\documentclass[12pt,letterpaper,oneside,openany]{report}
\usepackage[utf8]{inputenc}
\usepackage[T1]{fontenc}
\usepackage[spanish,es-tabla,es-noquoting]{babel}
\usepackage{lmodern}
\usepackage{textcomp,amsmath,amssymb}
\usepackage[paperwidth=612bp,paperheight=792bp,left=2.5cm,right=2.5cm,top=2.5cm,bottom=2.5cm]{geometry}
\usepackage{graphicx,tikz}
\usetikzlibrary{positioning,calc}
\usepackage{booktabs,longtable,array,fvextra,float,setspace}
\usepackage[hidelinks,unicode]{hyperref}
\usepackage{xurl}
\usepackage{bookmark}
\newcolumntype{P}[1]{>{\raggedright\arraybackslash}p{#1}}
\raggedbottom
\setstretch{1.12}
\setlength{\parskip}{0.35em}
\setlength{\emergencystretch}{2em}
\setcounter{tocdepth}{2}
\setcounter{secnumdepth}{2}
\newcommand{\archivo}[1]{\nolinkurl{#1}}
\pagestyle{empty}
\setlength{\parindent}{0pt}
% Una página usa el fondo vectorial sin texto y texto LaTeX editable.
\newenvironment{pagina}[1]{\ifcsname entrada#1\endcsname\csname entrada#1\endcsname\fi\noindent\begin{tikzpicture}[remember picture,overlay,x=1bp,y=1bp]
\node[anchor=south west,inner sep=0] at (current page.south west) {\includegraphics[page=#1,width=612bp,height=792bp]{assets/fondos.pdf}};
}{\end{tikzpicture}\null\newpage}
\newcommand{\texto}[5]{\node[anchor=base west,inner sep=0pt,outer sep=0pt] at ([xshift=#1bp,yshift=#2bp]current page.south west) {{#4\fontsize{#3bp}{#3bp}\selectfont #5}};}
\newcommand{\evidencia}[3]{\IfFileExists{assets/evidencias/#1}{\begin{figure}[htbp]\centering\includegraphics[width=0.95\linewidth,height=0.42\textheight,keepaspectratio]{assets/evidencias/#1}\caption{#2.}\end{figure}}{}}
\expandafter\def\csname entrada2\endcsname{\phantomsection\addcontentsline{toc}{chapter}{Dedicatoria}}
\expandafter\def\csname entrada3\endcsname{\phantomsection\addcontentsline{toc}{chapter}{Agradecimientos}}
\expandafter\def\csname entrada4\endcsname{\phantomsection\addcontentsline{toc}{chapter}{Resumen}}
\expandafter\def\csname entrada5\endcsname{\phantomsection\addcontentsline{toc}{chapter}{Abstract}}
\expandafter\def\csname entrada11\endcsname{\phantomsection\addcontentsline{toc}{chapter}{\protect\numberline{1}Formulación del proyecto}}
\expandafter\def\csname entrada12\endcsname{\phantomsection\addcontentsline{toc}{section}{\protect\numberline{1.1}Introducción}}
\expandafter\def\csname entrada13\endcsname{\phantomsection\addcontentsline{toc}{section}{\protect\numberline{1.2}Contexto del proyecto}}
\expandafter\def\csname entrada14\endcsname{\phantomsection\addcontentsline{toc}{section}{\protect\numberline{1.3}Diagnóstico de la situación actual}\phantomsection\addcontentsline{toc}{subsection}{\protect\numberline{1.3.1}Proceso actual de gestión de los operadores}\phantomsection\addcontentsline{toc}{subsection}{\protect\numberline{1.3.2}Personas y recursos involucrados}\phantomsection\addcontentsline{toc}{subsection}{\protect\numberline{1.3.3}Necesidades identificadas y principales hallazgos}}
\expandafter\def\csname entrada16\endcsname{\phantomsection\addcontentsline{toc}{section}{\protect\numberline{1.4}Planteamiento del problema}}
\expandafter\def\csname entrada17\endcsname{\phantomsection\addcontentsline{toc}{section}{\protect\numberline{1.5}Justificación}}
\expandafter\def\csname entrada18\endcsname{\phantomsection\addcontentsline{toc}{section}{\protect\numberline{1.6}Objetivos}\phantomsection\addcontentsline{toc}{subsection}{\protect\numberline{1.6.1}Objetivo general}\phantomsection\addcontentsline{toc}{subsection}{\protect\numberline{1.6.2}Objetivos específicos}}
\expandafter\def\csname entrada19\endcsname{\phantomsection\addcontentsline{toc}{section}{\protect\numberline{1.7}Alcance del proyecto}}
\expandafter\def\csname entrada20\endcsname{\phantomsection\addcontentsline{toc}{section}{\protect\numberline{1.8}Metodología de desarrollo}}
\expandafter\def\csname entrada21\endcsname{\phantomsection\addcontentsline{toc}{section}{\protect\numberline{1.9}Estructura del ecosistema}}
\expandafter\def\csname entrada22\endcsname{\phantomsection\addcontentsline{toc}{chapter}{\protect\numberline{2}Investigación documental}}
\expandafter\def\csname entrada23\endcsname{\phantomsection\addcontentsline{toc}{section}{\protect\numberline{2.1}Metodología de búsqueda}}
\expandafter\def\csname entrada24\endcsname{\phantomsection\addcontentsline{toc}{section}{\protect\numberline{2.2}Tendencias tecnológicas relacionadas}\phantomsection\addcontentsline{toc}{subsection}{\protect\numberline{2.2.1}Aplicaciones móviles híbridas}\phantomsection\addcontentsline{toc}{subsection}{\protect\numberline{2.2.2}Aplicaciones de escritorio y frameworks .NET}\phantomsection\addcontentsline{toc}{subsection}{\protect\numberline{2.2.3}Consumo de APIs REST y arquitecturas orientadas a servicios}}
\expandafter\def\csname entrada25\endcsname{\phantomsection\addcontentsline{toc}{subsection}{\protect\numberline{2.2.4}Control de acceso basado en roles (RBAC)}\phantomsection\addcontentsline{toc}{subsection}{\protect\numberline{2.2.5}Pasarelas de pago digital en Centroamérica}\phantomsection\addcontentsline{toc}{subsection}{\protect\numberline{2.2.6}Geolocalización y mapas}\phantomsection\addcontentsline{toc}{subsection}{\protect\numberline{2.2.7}Marketplaces turísticos y comisiones}\phantomsection\addcontentsline{toc}{subsection}{\protect\numberline{2.2.8}Autenticación y seguridad}}
\expandafter\def\csname entrada26\endcsname{\phantomsection\addcontentsline{toc}{subsection}{\protect\numberline{2.2.9}Sistemas de reseñas y verificación}\phantomsection\addcontentsline{toc}{subsection}{\protect\numberline{2.2.10}Notificaciones push}\phantomsection\addcontentsline{toc}{subsection}{\protect\numberline{2.2.11}Herramientas de modelado y metodologías Lean}}
\expandafter\def\csname entrada27\endcsname{\phantomsection\addcontentsline{toc}{section}{\protect\numberline{2.3}Estado del arte}}
\expandafter\def\csname entrada28\endcsname{\phantomsection\addcontentsline{toc}{section}{\protect\numberline{2.4}Tecnologías seleccionadas}\phantomsection\addcontentsline{toc}{subsection}{\protect\numberline{2.4.1}.NET 10 y C\#}\phantomsection\addcontentsline{toc}{subsection}{\protect\numberline{2.4.2}Windows Forms (WinForms)}\phantomsection\addcontentsline{toc}{subsection}{\protect\numberline{2.4.3}Arquitectura por capas}\phantomsection\addcontentsline{toc}{subsection}{\protect\numberline{2.4.4}React Native}\phantomsection\addcontentsline{toc}{subsection}{\protect\numberline{2.4.5}API REST, JWT y OAuth 2.0}}
\expandafter\def\csname entrada29\endcsname{\phantomsection\addcontentsline{toc}{subsection}{\protect\numberline{2.4.6}Pasarela de pago Stripe y SSL/TLS}\phantomsection\addcontentsline{toc}{subsection}{\protect\numberline{2.4.7}Base de datos y herramientas de soporte}}
\expandafter\def\csname entrada30\endcsname{\phantomsection\addcontentsline{toc}{section}{\protect\numberline{2.5}Análisis de las tecnologías y su pertinencia para el proyecto}\phantomsection\addcontentsline{toc}{subsection}{\protect\numberline{2.5.1}Comparación de alternativas para los portales de escritorio}\phantomsection\addcontentsline{toc}{subsection}{\protect\numberline{2.5.2}Justificación de la decisión}\phantomsection\addcontentsline{toc}{subsection}{\protect\numberline{2.5.3}Compatibilidad entre tecnologías}}
\expandafter\def\csname entrada32\endcsname{\phantomsection\addcontentsline{toc}{section}{\protect\numberline{2.6}Conclusiones de la investigación}}
\expandafter\def\csname entrada33\endcsname{\phantomsection\addcontentsline{toc}{chapter}{\protect\numberline{3}Emprendimiento digital}}
\expandafter\def\csname entrada34\endcsname{\phantomsection\addcontentsline{toc}{section}{\protect\numberline{3.1}Idea de negocio}}
\expandafter\def\csname entrada35\endcsname{\phantomsection\addcontentsline{toc}{section}{\protect\numberline{3.2}Problema o necesidad identificada}}
\expandafter\def\csname entrada36\endcsname{\phantomsection\addcontentsline{toc}{section}{\protect\numberline{3.3}Propuesta de valor}}
\expandafter\def\csname entrada37\endcsname{\phantomsection\addcontentsline{toc}{section}{\protect\numberline{3.4}Público objetivo}}
\expandafter\def\csname entrada38\endcsname{\phantomsection\addcontentsline{toc}{section}{\protect\numberline{3.5}Modelo de funcionamiento del emprendimiento}}
\expandafter\def\csname entrada39\endcsname{\phantomsection\addcontentsline{toc}{section}{\protect\numberline{3.6}Misión}}
\expandafter\def\csname entrada40\endcsname{\phantomsection\addcontentsline{toc}{section}{\protect\numberline{3.7}Visión}}
\expandafter\def\csname entrada41\endcsname{\phantomsection\addcontentsline{toc}{section}{\protect\numberline{3.8}Valores}}
\expandafter\def\csname entrada42\endcsname{\phantomsection\addcontentsline{toc}{section}{\protect\numberline{3.9}Identidad visual}}
\expandafter\def\csname entrada43\endcsname{\phantomsection\addcontentsline{toc}{section}{\protect\numberline{3.10}Análisis FODA}}
\expandafter\def\csname entrada44\endcsname{\phantomsection\addcontentsline{toc}{section}{\protect\numberline{3.11}Lean Canvas}}
\expandafter\def\csname entrada45\endcsname{\phantomsection\addcontentsline{toc}{section}{\protect\numberline{3.12}Consideraciones éticas, sociales e identidad nacional}}
\expandafter\def\csname entrada46\endcsname{\phantomsection\addcontentsline{toc}{chapter}{\protect\numberline{4}Modelado y automatización del proceso}}
\expandafter\def\csname entrada47\endcsname{\phantomsection\addcontentsline{toc}{section}{\protect\numberline{4.1}Descripción del proceso actual}}
\expandafter\def\csname entrada48\endcsname{\phantomsection\addcontentsline{toc}{section}{\protect\numberline{4.2}Actores y elementos del proceso}}
\expandafter\def\csname entrada49\endcsname{\phantomsection\addcontentsline{toc}{section}{\protect\numberline{4.3}BPMN del proceso actual}}
\expandafter\def\csname entrada50\endcsname{\phantomsection\addcontentsline{toc}{section}{\protect\numberline{4.4}Identificación de oportunidades de mejora}}
\expandafter\def\csname entrada51\endcsname{\phantomsection\addcontentsline{toc}{section}{\protect\numberline{4.5}Propuesta de automatización mediante TI}}
\expandafter\def\csname entrada52\endcsname{\phantomsection\addcontentsline{toc}{section}{\protect\numberline{4.6}BPMN del proceso propuesto}}
\expandafter\def\csname entrada53\endcsname{\phantomsection\addcontentsline{toc}{section}{\protect\numberline{4.7}BPMN de los procesos del prototipo de escritorio}\phantomsection\addcontentsline{toc}{subsection}{\protect\numberline{4.7.1}BPMN: Alta y verificación de operadores}}
\expandafter\def\csname entrada54\endcsname{\phantomsection\addcontentsline{toc}{subsection}{\protect\numberline{4.7.2}BPMN: Gestión del catálogo y disponibilidad}\phantomsection\addcontentsline{toc}{subsection}{\protect\numberline{4.7.3}BPMN: Reserva, pago y validación de voucher}\phantomsection\addcontentsline{toc}{subsection}{\protect\numberline{4.7.4}BPMN: Moderación de reseñas, reportes financieros y auditoría}}
\expandafter\def\csname entrada57\endcsname{\phantomsection\addcontentsline{toc}{section}{\protect\numberline{4.8}Descripción funcional del proceso propuesto}}
\expandafter\def\csname entrada58\endcsname{\phantomsection\addcontentsline{toc}{section}{\protect\numberline{4.9}Beneficios esperados}}

\hypersetup{pdftitle={Explora-Nicaragua: Proyecto Integrador II},pdfauthor={Equipo Explora-Nicaragua}}
\begin{document}
\newgeometry{margin=0pt}
\hypersetup{pageanchor=false}
\begingroup
\definecolor{PortadaAzul}{HTML}{09235B}
\definecolor{PortadaDorado}{HTML}{E3B23C}
\definecolor{PortadaGris}{HTML}{44474F}
\definecolor{PortadaFondo}{HTML}{EFF3FF}
\noindent\begin{tikzpicture}[remember picture,overlay,x=1bp,y=1bp]
\begin{scope}[shift={(current page.south west)}]
\node[anchor=south west,inner sep=0pt] at (0,0) {\includegraphics[width=612bp]{assets/fondo-portada-pptx.png}};

\fill[PortadaAzul] (48,624) rectangle (564,744);
\fill[PortadaDorado] (48,624) rectangle (564,628);
\fill[white,rounded corners=5bp] (60,654) rectangle (122,730);
\node[inner sep=0pt] at (91,692) {\includegraphics[width=48bp,height=59bp,keepaspectratio]{assets/escudo-unan.png}};

\node[anchor=north west,align=left,inner sep=0pt,text=white,font=\sffamily\bfseries\fontsize{10.5bp}{13bp}\selectfont] at (136,730) {UNIVERSIDAD NACIONAL AUTÓNOMA\\DE NICARAGUA, MANAGUA};
\node[anchor=base west,inner sep=0pt,text=white,font=\sffamily\bfseries\fontsize{14bp}{17bp}\selectfont] at (136,686) {Centro Universitario Regional de Carazo};
\node[anchor=base west,inner sep=0pt,text=white,font=\sffamily\bfseries\fontsize{14bp}{17bp}\selectfont] at (136,665) {``Miguel Ángel Ortez''};
\node[anchor=base west,inner sep=0pt,text=PortadaDorado,font=\sffamily\fontsize{10.7bp}{13bp}\selectfont] at (136,642) {Departamento de Ciencias Tecnológicas y Salud};

\node[anchor=base,inner sep=0pt,text=PortadaAzul,font=\sffamily\bfseries\fontsize{27bp}{32bp}\selectfont] at (306,578) {Proyecto Integrador II};
\node[anchor=base,inner sep=0pt,text=PortadaGris,font=\sffamily\fontsize{22bp}{26bp}\selectfont] at (306,543) {Explora-Nicaragua};
\fill[PortadaDorado] (91,500) rectangle (521,523);
\node[inner sep=0pt,text=PortadaAzul,font=\sffamily\bfseries\fontsize{9.7bp}{12bp}\selectfont] at (306,511.5) {INTEGRADOR II / SEGUNDO CORTE EVALUATIVO};

\node[anchor=base west,inner sep=0pt,text=PortadaAzul,font=\sffamily\bfseries\fontsize{14bp}{17bp}\selectfont] at (64,460) {Integrantes};
\fill[PortadaDorado] (64,448) rectangle (174,451);
\fill[PortadaFondo] (64,361) rectangle (548,438);
\node[anchor=base west,inner sep=0pt,text=PortadaGris,font=\sffamily\fontsize{13bp}{16bp}\selectfont] at (78,418) {Nohelia Nohemi Nicoya Bravo};
\node[anchor=base west,inner sep=0pt,text=PortadaGris,font=\sffamily\fontsize{13bp}{16bp}\selectfont] at (78,393) {Geoconda Nicole López Baltodano};
\node[anchor=base west,inner sep=0pt,text=PortadaGris,font=\sffamily\fontsize{13bp}{16bp}\selectfont] at (78,368) {Alexander José Cano Calero};

\node[anchor=base west,inner sep=0pt,text=PortadaAzul,font=\sffamily\bfseries\fontsize{14bp}{17bp}\selectfont] at (64,331) {Datos generales};
\fill[PortadaDorado] (64,319) rectangle (191,322);
\fill[PortadaFondo] (64,194) rectangle (548,309);
\fill[PortadaAzul] (64,194) rectangle (68,309);
\foreach \coverY/\coverLabel/\coverValue in {290/Carrera:/Ingeniería en Sistemas,263/Año:/I,236/Docente:/Mg. Martha Rodríguez,209/Fecha:/Agosto 2026}{
\node[anchor=base west,inner sep=0pt,text=PortadaAzul,font=\sffamily\bfseries\fontsize{11.5bp}{14bp}\selectfont] at (79,\coverY) {\coverLabel};
\node[anchor=base west,inner sep=0pt,text=PortadaGris,font=\sffamily\fontsize{12.5bp}{15bp}\selectfont] at (195,\coverY) {\coverValue};
}

\node[anchor=base west,inner sep=0pt,text=PortadaAzul,font=\sffamily\bfseries\fontsize{11.5bp}{14bp}\selectfont] at (64,162) {Componentes que contribuyen al proyecto};
\foreach \coverX/\coverText in {64/{HISTORIA E\\IDENTIDAD\\NACIONAL},162/{INTEGRADOR II},260/{TI EN LA\\EMPRESA},358/{MATEMÁTICA\\DISCRETA},456/{PROGRAMACIÓN\\ESTRUCTURADA}}{
\draw[PortadaAzul!45,dashed,line width=0.6bp] (\coverX,97) rectangle +(92,47);
\node[align=center,text width=87bp,inner sep=0pt,text=PortadaGris,font=\sffamily\fontsize{8.5bp}{10.5bp}\selectfont] at (\coverX+46,120.5) {\coverText};
}

\fill[PortadaDorado] (48,71) rectangle (564,74);
\fill[PortadaAzul] (48,47) rectangle (564,71);
\node[inner sep=0pt,text=white,font=\sffamily\fontsize{9bp}{11bp}\selectfont] at (306,59) {UNAN-MANAGUA / CUR-CARAZO};
\end{scope}
\end{tikzpicture}\null\newpage
\endgroup

\pagenumbering{roman}
% Texto editable; coordenadas en puntos PDF.
% Página original 2
\begin{pagina}{2}
\texto{70.8660}{625.4930}{24.7871}{\normalfont\bfseries}{Dedicatoria}
\texto{70.8660}{571.1960}{11.9552}{\normalfont}{A}
\texto{84.0150}{571.1960}{11.9552}{\normalfont\bfseries}{nuestros}
\texto{137.8062}{571.1960}{11.9552}{\normalfont\bfseries}{padres}
\texto{176.2230}{571.1960}{11.9552}{\normalfont}{,}
\texto{183.8385}{571.1960}{11.9552}{\normalfont}{por}
\texto{205.4463}{571.1960}{11.9552}{\normalfont}{su}
\texto{220.9426}{571.1960}{11.9552}{\normalfont}{amor}
\texto{251.3327}{571.1960}{11.9552}{\normalfont}{incondicional,}
\texto{325.9475}{571.1960}{11.9552}{\normalfont}{su}
\texto{341.4439}{571.1960}{11.9552}{\normalfont}{sacrificio}
\texto{390.1087}{571.1960}{11.9552}{\normalfont}{y}
\texto{400.6628}{571.1960}{11.9552}{\normalfont}{por}
\texto{422.2706}{571.1960}{11.9552}{\normalfont}{inculcarnos}
\texto{483.9307}{571.1960}{11.9552}{\normalfont}{el}
\texto{496.7610}{571.1960}{11.9552}{\normalfont}{valor}
\texto{526.1804}{571.1960}{11.9552}{\normalfont}{del}
\texto{70.8660}{556.7500}{11.9552}{\normalfont}{esfuerzo}
\texto{115.4756}{556.7500}{11.9552}{\normalfont}{y}
\texto{125.5515}{556.7500}{11.9552}{\normalfont}{la}
\texto{138.5659}{556.7500}{11.9552}{\normalfont}{perseverancia;}
\texto{214.3930}{556.7500}{11.9552}{\normalfont}{sin}
\texto{232.6629}{556.7500}{11.9552}{\normalfont}{su}
\texto{247.6930}{556.7500}{11.9552}{\normalfont}{apoyo}
\texto{281.5095}{556.7500}{11.9552}{\normalfont}{este}
\texto{304.9823}{556.7500}{11.9552}{\normalfont}{logro}
\texto{334.2558}{556.7500}{11.9552}{\normalfont}{no}
\texto{350.5101}{556.7500}{11.9552}{\normalfont}{habría}
\texto{386.9376}{556.7500}{11.9552}{\normalfont}{sido}
\texto{411.0608}{556.7500}{11.9552}{\normalfont}{posible.}
\texto{70.8660}{528.1310}{11.9552}{\normalfont}{A}
\texto{82.3860}{528.1310}{11.9552}{\normalfont\bfseries}{nuestros}
\texto{134.3003}{528.1310}{11.9552}{\normalfont\bfseries}{hijos}
\texto{161.6540}{528.1310}{11.9552}{\normalfont}{,}
\texto{167.6436}{528.1310}{11.9552}{\normalfont}{que}
\texto{188.2782}{528.1310}{11.9552}{\normalfont}{son}
\texto{208.0019}{528.1310}{11.9552}{\normalfont}{nuestra}
\texto{248.2025}{528.1310}{11.9552}{\normalfont}{mayor}
\texto{282.4876}{528.1310}{11.9552}{\normalfont}{motivación}
\texto{340.8481}{528.1310}{11.9552}{\normalfont}{y}
\texto{349.7643}{528.1310}{11.9552}{\normalfont}{la}
\texto{361.6191}{528.1310}{11.9552}{\normalfont}{razón}
\texto{392.3344}{528.1310}{11.9552}{\normalfont}{por}
\texto{412.3043}{528.1310}{11.9552}{\normalfont}{la}
\texto{424.1591}{528.1310}{11.9552}{\normalfont}{cual}
\texto{447.7204}{528.1310}{11.9552}{\normalfont}{cada}
\texto{473.8712}{528.1310}{11.9552}{\normalfont}{día}
\texto{492.2296}{528.1310}{11.9552}{\normalfont}{buscamos}
\texto{70.8660}{513.6850}{11.9552}{\normalfont}{ser}
\texto{89.5185}{513.6850}{11.9552}{\normalfont}{mejores;}
\texto{135.8114}{513.6850}{11.9552}{\normalfont}{que}
\texto{157.9883}{513.6850}{11.9552}{\normalfont}{este}
\texto{181.8437}{513.6850}{11.9552}{\normalfont}{trabajo}
\texto{223.2039}{513.6850}{11.9552}{\normalfont}{sea}
\texto{243.1572}{513.6850}{11.9552}{\normalfont}{un}
\texto{260.4444}{513.6850}{11.9552}{\normalfont}{ejemplo}
\texto{304.0833}{513.6850}{11.9552}{\normalfont}{de}
\texto{320.0698}{513.6850}{11.9552}{\normalfont}{que}
\texto{342.2347}{513.6850}{11.9552}{\normalfont}{con}
\texto{364.0864}{513.6850}{11.9552}{\normalfont}{dedicación}
\texto{421.6961}{513.6850}{11.9552}{\normalfont}{y}
\texto{432.1545}{513.6850}{11.9552}{\normalfont}{constancia}
\texto{489.8408}{513.6850}{11.9552}{\normalfont}{los}
\texto{507.8429}{513.6850}{11.9552}{\normalfont}{sueños}
\texto{70.8660}{499.2400}{11.9552}{\normalfont}{se}
\texto{84.5834}{499.2400}{11.9552}{\normalfont}{alcanzan.}
\texto{70.8660}{470.6210}{11.9552}{\normalfont}{Y}
\texto{82.6897}{470.6210}{11.9552}{\normalfont}{a}
\texto{91.7600}{470.6210}{11.9552}{\normalfont\bfseries}{todas}
\texto{126.9167}{470.6210}{11.9552}{\normalfont\bfseries}{aquellas}
\texto{176.9934}{470.6210}{11.9552}{\normalfont\bfseries}{personas}
\texto{230.8780}{470.6210}{11.9552}{\normalfont}{que}
\texto{251.8116}{470.6210}{11.9552}{\normalfont}{de}
\texto{266.5786}{470.6210}{11.9552}{\normalfont}{una}
\texto{288.4997}{470.6210}{11.9552}{\normalfont}{u}
\texto{298.0638}{470.6210}{11.9552}{\normalfont}{otra}
\texto{321.9360}{470.6210}{11.9552}{\normalfont}{forma}
\texto{354.5880}{470.6210}{11.9552}{\normalfont}{nos}
\texto{374.6225}{470.6210}{11.9552}{\normalfont}{brindaron}
\texto{427.7610}{470.6210}{11.9552}{\normalfont}{su}
\texto{441.9423}{470.6210}{11.9552}{\normalfont}{apoyo,}
\texto{478.1737}{470.6210}{11.9552}{\normalfont}{su}
\texto{492.3550}{470.6210}{11.9552}{\normalfont}{compañía}
\texto{70.8660}{456.1750}{11.9552}{\normalfont}{y}
\texto{80.9418}{456.1750}{11.9552}{\normalfont}{sus}
\texto{100.5771}{456.1750}{11.9552}{\normalfont}{palabras}
\texto{147.4749}{456.1750}{11.9552}{\normalfont}{de}
\texto{163.0788}{456.1750}{11.9552}{\normalfont}{aliento}
\texto{201.1227}{456.1750}{11.9552}{\normalfont}{a}
\texto{210.8853}{456.1750}{11.9552}{\normalfont}{lo}
\texto{223.8878}{456.1750}{11.9552}{\normalfont}{largo}
\texto{253.1493}{456.1750}{11.9552}{\normalfont}{de}
\texto{268.7652}{456.1750}{11.9552}{\normalfont}{este}
\texto{292.2380}{456.1750}{11.9552}{\normalfont}{camino.}
\texto{304.1930}{36.9930}{11.9552}{\normalfont\scshape}{i}
\end{pagina}
% Página original 3
\begin{pagina}{3}
\texto{70.8660}{625.4930}{24.7871}{\normalfont\bfseries}{Agradecimientos}
\texto{70.8660}{571.1960}{11.9552}{\normalfont}{En}
\texto{89.4707}{571.1960}{11.9552}{\normalfont}{primer}
\texto{127.4261}{571.1960}{11.9552}{\normalfont}{lugar,}
\texto{160.8289}{571.1960}{11.9552}{\normalfont}{agradecemos}
\texto{229.4159}{571.1960}{11.9552}{\normalfont}{a}
\texto{239.5560}{571.1960}{11.9552}{\normalfont\bfseries}{Dios}
\texto{269.7680}{571.1960}{11.9552}{\normalfont}{por}
\texto{291.1367}{571.1960}{11.9552}{\normalfont}{darnos}
\texto{329.1567}{571.1960}{11.9552}{\normalfont}{la}
\texto{342.3982}{571.1960}{11.9552}{\normalfont}{vida,}
\texto{371.5737}{571.1960}{11.9552}{\normalfont}{la}
\texto{384.8153}{571.1960}{11.9552}{\normalfont}{salud}
\texto{415.6812}{571.1960}{11.9552}{\normalfont}{y}
\texto{426.0081}{571.1960}{11.9552}{\normalfont}{la}
\texto{439.2497}{571.1960}{11.9552}{\normalfont}{sabiduría}
\texto{490.2769}{571.1960}{11.9552}{\normalfont}{necesarias}
\texto{70.8660}{556.7500}{11.9552}{\normalfont}{para}
\texto{97.5261}{556.7500}{11.9552}{\normalfont}{culminar}
\texto{146.2985}{556.7500}{11.9552}{\normalfont}{este}
\texto{169.7833}{556.7500}{11.9552}{\normalfont}{proyecto.}
\texto{70.8660}{528.1310}{11.9552}{\normalfont}{A}
\texto{83.2050}{528.1310}{11.9552}{\normalfont\bfseries}{nuestras}
\texto{135.8772}{528.1310}{11.9552}{\normalfont\bfseries}{familias}
\texto{180.7860}{528.1310}{11.9552}{\normalfont}{,}
\texto{187.6005}{528.1310}{11.9552}{\normalfont}{por}
\texto{208.3953}{528.1310}{11.9552}{\normalfont}{su}
\texto{223.0787}{528.1310}{11.9552}{\normalfont}{apoyo}
\texto{256.5604}{528.1310}{11.9552}{\normalfont}{incondicional,}
\texto{330.3623}{528.1310}{11.9552}{\normalfont}{su}
\texto{345.0457}{528.1310}{11.9552}{\normalfont}{paciencia}
\texto{395.4464}{528.1310}{11.9552}{\normalfont}{y}
\texto{405.1875}{528.1310}{11.9552}{\normalfont}{por}
\texto{425.9824}{528.1310}{11.9552}{\normalfont}{acompañarnos}
\texto{502.4502}{528.1310}{11.9552}{\normalfont}{en}
\texto{517.7194}{528.1310}{11.9552}{\normalfont}{cada}
\texto{70.8660}{513.6850}{11.9552}{\normalfont}{paso}
\texto{97.5907}{513.6850}{11.9552}{\normalfont}{de}
\texto{113.1946}{513.6850}{11.9552}{\normalfont}{este}
\texto{136.6794}{513.6850}{11.9552}{\normalfont}{camino}
\texto{176.9971}{513.6850}{11.9552}{\normalfont}{académico.}
\texto{70.8660}{485.0660}{11.9552}{\normalfont}{A}
\texto{83.3100}{485.0660}{11.9552}{\normalfont\bfseries}{nuestros}
\texto{136.2883}{485.0660}{11.9552}{\normalfont\bfseries}{docentes}
\texto{186.7480}{485.0660}{11.9552}{\normalfont}{,}
\texto{193.6581}{485.0660}{11.9552}{\normalfont}{por}
\texto{214.5725}{485.0660}{11.9552}{\normalfont}{compartir}
\texto{268.3207}{485.0660}{11.9552}{\normalfont}{sus}
\texto{287.7168}{485.0660}{11.9552}{\normalfont}{conocimientos,}
\texto{366.2434}{485.0660}{11.9552}{\normalfont}{su}
\texto{381.0343}{485.0660}{11.9552}{\normalfont}{tiempo}
\texto{420.1470}{485.0660}{11.9552}{\normalfont}{y}
\texto{429.9957}{485.0660}{11.9552}{\normalfont}{su}
\texto{444.7866}{485.0660}{11.9552}{\normalfont}{dedicación}
\texto{501.7866}{485.0660}{11.9552}{\normalfont}{durante}
\texto{70.8660}{470.6210}{11.9552}{\normalfont}{el}
\texto{83.4094}{470.6210}{11.9552}{\normalfont}{desarrollo}
\texto{136.9902}{470.6210}{11.9552}{\normalfont}{de}
\texto{152.7854}{470.6210}{11.9552}{\normalfont}{este}
\texto{176.4495}{470.6210}{11.9552}{\normalfont}{proyecto}
\texto{223.7921}{470.6210}{11.9552}{\normalfont}{y}
\texto{234.0593}{470.6210}{11.9552}{\normalfont}{a}
\texto{244.0012}{470.6210}{11.9552}{\normalfont}{lo}
\texto{257.1950}{470.6210}{11.9552}{\normalfont}{largo}
\texto{286.6478}{470.6210}{11.9552}{\normalfont}{de}
\texto{302.4550}{470.6210}{11.9552}{\normalfont}{toda}
\texto{329.6291}{470.6210}{11.9552}{\normalfont}{la}
\texto{342.8229}{470.6210}{11.9552}{\normalfont}{carrera;}
\texto{385.9333}{470.6210}{11.9552}{\normalfont}{en}
\texto{401.7286}{470.6210}{11.9552}{\normalfont}{especial,}
\texto{448.4901}{470.6210}{11.9552}{\normalfont}{por}
\texto{469.8110}{470.6210}{11.9552}{\normalfont}{guiarnos}
\texto{516.8882}{470.6210}{11.9552}{\normalfont}{en}
\texto{532.6834}{470.6210}{11.9552}{\normalfont}{el}
\texto{70.8660}{456.1750}{11.9552}{\normalfont}{componente}
\texto{135.8975}{456.1750}{11.9552}{\normalfont}{de}
\texto{151.5014}{456.1750}{11.9552}{\normalfont}{Integrador}
\texto{208.7262}{456.1750}{11.9552}{\normalfont}{II.}
\texto{70.8660}{427.5560}{11.9552}{\normalfont}{A}
\texto{83.2038}{427.5560}{11.9552}{\normalfont}{la}
\texto{96.0400}{427.5560}{11.9552}{\normalfont\bfseries}{Universidad}
\texto{170.5675}{427.5560}{11.9552}{\normalfont\bfseries}{Nacional}
\texto{225.9069}{427.5560}{11.9552}{\normalfont\bfseries}{Autónoma}
\texto{291.1752}{427.5560}{11.9552}{\normalfont\bfseries}{de}
\texto{308.8844}{427.5560}{11.9552}{\normalfont\bfseries}{Nicaragua}
\texto{368.4210}{427.5560}{11.9552}{\normalfont}{,}
\texto{375.2474}{427.5560}{11.9552}{\normalfont}{en}
\texto{390.5166}{427.5560}{11.9552}{\normalfont}{particular}
\texto{444.1691}{427.5560}{11.9552}{\normalfont}{al}
\texto{456.8488}{427.5560}{11.9552}{\normalfont}{Centro}
\texto{495.2202}{427.5560}{11.9552}{\normalfont}{Universi-}
\texto{70.8660}{413.1100}{11.9552}{\normalfont}{tario}
\texto{97.6911}{413.1100}{11.9552}{\normalfont}{Regional}
\texto{144.8376}{413.1100}{11.9552}{\normalfont}{de}
\texto{159.3058}{413.1100}{11.9552}{\normalfont}{Carazo}
\texto{197.8494}{413.1100}{11.9552}{\normalfont}{(CUR-Carazo),}
\texto{278.4848}{413.1100}{11.9552}{\normalfont}{por}
\texto{298.4906}{413.1100}{11.9552}{\normalfont}{brindarnos}
\texto{355.9473}{413.1100}{11.9552}{\normalfont}{la}
\texto{367.8140}{413.1100}{11.9552}{\normalfont}{formación}
\texto{420.9908}{413.1100}{11.9552}{\normalfont}{y}
\texto{429.9309}{413.1100}{11.9552}{\normalfont}{los}
\texto{446.4147}{413.1100}{11.9552}{\normalfont}{recursos}
\texto{490.2903}{413.1100}{11.9552}{\normalfont}{necesarios}
\texto{70.8660}{398.6640}{11.9552}{\normalfont}{para}
\texto{97.5261}{398.6640}{11.9552}{\normalfont}{nuestro}
\texto{138.8863}{398.6640}{11.9552}{\normalfont}{desarrollo}
\texto{192.2758}{398.6640}{11.9552}{\normalfont}{profesional.}
\texto{70.8660}{370.0450}{11.9552}{\normalfont}{Y}
\texto{83.2755}{370.0450}{11.9552}{\normalfont}{a}
\texto{92.7751}{370.0450}{11.9552}{\normalfont}{todas}
\texto{124.1121}{370.0450}{11.9552}{\normalfont}{las}
\texto{141.4860}{370.0450}{11.9552}{\normalfont\bfseries}{personas}
\texto{196.4895}{370.0450}{11.9552}{\normalfont\bfseries}{y}
\texto{207.7728}{370.0450}{11.9552}{\normalfont\bfseries}{operadores}
\texto{275.8254}{370.0450}{11.9552}{\normalfont\bfseries}{turísticos}
\texto{333.6790}{370.0450}{11.9552}{\normalfont}{que}
\texto{355.1984}{370.0450}{11.9552}{\normalfont}{colaboraron}
\texto{419.0009}{370.0450}{11.9552}{\normalfont}{con}
\texto{440.1950}{370.0450}{11.9552}{\normalfont}{información}
\texto{504.0000}{370.0450}{11.9552}{\normalfont}{y}
\texto{513.8247}{370.0450}{11.9552}{\normalfont}{expe-}
\texto{70.8660}{355.5990}{11.9552}{\normalfont}{riencias,}
\texto{116.4512}{355.5990}{11.9552}{\normalfont}{contribuyendo}
\texto{192.8688}{355.5990}{11.9552}{\normalfont}{de}
\texto{208.4727}{355.5990}{11.9552}{\normalfont}{manera}
\texto{250.1031}{355.5990}{11.9552}{\normalfont}{significativa}
\texto{314.2284}{355.5990}{11.9552}{\normalfont}{a}
\texto{323.9791}{355.5990}{11.9552}{\normalfont}{la}
\texto{336.9816}{355.5990}{11.9552}{\normalfont}{realización}
\texto{394.8710}{355.5990}{11.9552}{\normalfont}{de}
\texto{410.4750}{355.5990}{11.9552}{\normalfont}{este}
\texto{433.9598}{355.5990}{11.9552}{\normalfont}{trabajo.}
\texto{302.2540}{36.9930}{11.9552}{\normalfont\scshape}{ii}
\end{pagina}
% Página original 4
\begin{pagina}{4}
\texto{70.8660}{625.4930}{24.7871}{\normalfont\bfseries}{Resumen}
\texto{88.7990}{571.1960}{11.9552}{\normalfont}{El}
\texto{103.9247}{571.1960}{11.9552}{\normalfont}{presente}
\texto{149.8375}{571.1960}{11.9552}{\normalfont}{proyecto}
\texto{197.0007}{571.1960}{11.9552}{\normalfont}{de}
\texto{212.6166}{571.1960}{11.9552}{\normalfont}{Integrador}
\texto{269.8533}{571.1960}{11.9552}{\normalfont}{II}
\texto{282.5306}{571.1960}{11.9552}{\normalfont}{desarrolla}
\texto{335.9441}{571.1960}{11.9552}{\normalfont}{el}
\texto{348.3081}{571.1960}{11.9552}{\normalfont}{prototipo}
\texto{400.0168}{571.1960}{11.9552}{\normalfont}{de}
\texto{415.6326}{571.1960}{11.9552}{\normalfont}{las}
\texto{433.2640}{571.1960}{11.9552}{\normalfont\bfseries}{aplicaciones}
\texto{507.0981}{571.1960}{11.9552}{\normalfont\bfseries}{de}
\texto{525.2019}{571.1960}{11.9552}{\normalfont\bfseries}{es-}
\texto{70.8660}{556.7500}{11.9552}{\normalfont\bfseries}{critorio}
\texto{117.0070}{556.7500}{11.9552}{\normalfont}{del}
\texto{134.9900}{556.7500}{11.9552}{\normalfont}{ecosistema}
\texto{192.1239}{556.7500}{11.9552}{\normalfont}{tecnológico}
\texto{251.7270}{556.7500}{11.9552}{\normalfont\bfseries}{Explora-Nicaragua}
\texto{361.6340}{556.7500}{11.9552}{\normalfont}{,}
\texto{367.9105}{556.7500}{11.9552}{\normalfont}{orientado}
\texto{418.7273}{556.7500}{11.9552}{\normalfont}{a}
\texto{427.6052}{556.7500}{11.9552}{\normalfont}{la}
\texto{439.7469}{556.7500}{11.9552}{\normalfont}{promoción}
\texto{496.4241}{556.7500}{11.9552}{\normalfont}{y}
\texto{505.6272}{556.7500}{11.9552}{\normalfont}{reserva}
\texto{70.8660}{542.3040}{11.9552}{\normalfont}{de}
\texto{86.7329}{542.3040}{11.9552}{\normalfont}{tours}
\texto{116.9724}{542.3040}{11.9552}{\normalfont}{turísticos}
\texto{168.1000}{542.3040}{11.9552}{\normalfont}{en}
\texto{183.9670}{542.3040}{11.9552}{\normalfont}{Nicaragua.}
\texto{243.0902}{542.3040}{11.9552}{\normalfont}{El}
\texto{258.4670}{542.3040}{11.9552}{\normalfont}{alcance}
\texto{299.7100}{542.3040}{11.9552}{\normalfont}{inicial}
\texto{334.4375}{542.3040}{11.9552}{\normalfont}{se}
\texto{348.4179}{542.3040}{11.9552}{\normalfont}{limita}
\texto{382.5070}{542.3040}{11.9552}{\normalfont}{a}
\texto{392.5207}{542.3040}{11.9552}{\normalfont}{dos}
\texto{413.6670}{542.3040}{11.9552}{\normalfont}{aplicaciones}
\texto{478.3757}{542.3040}{11.9552}{\normalfont}{de}
\texto{494.2546}{542.3040}{11.9552}{\normalfont}{escritorio}
\texto{70.8660}{527.8590}{11.9552}{\normalfont}{que}
\texto{93.0549}{527.8590}{11.9552}{\normalfont}{trabajan}
\texto{140.9306}{527.8590}{11.9552}{\normalfont}{de}
\texto{156.9410}{527.8590}{11.9552}{\normalfont}{forma}
\texto{190.8484}{527.8590}{11.9552}{\normalfont}{coordinada}
\texto{251.4063}{527.8590}{11.9552}{\normalfont}{sobre}
\texto{282.4396}{527.8590}{11.9552}{\normalfont}{una}
\texto{305.6040}{527.8590}{11.9552}{\normalfont}{misma}
\texto{343.1529}{527.8590}{11.9552}{\normalfont}{API}
\texto{368.4190}{527.8590}{11.9552}{\normalfont}{REST:}
\texto{407.5125}{527.8590}{11.9552}{\normalfont}{un}
\texto{424.8236}{527.8590}{11.9552}{\normalfont}{backoffice}
\texto{478.2347}{527.8590}{11.9552}{\normalfont}{de}
\texto{494.2451}{527.8590}{11.9552}{\normalfont}{escritorio}
\texto{70.8660}{513.4130}{11.9552}{\normalfont}{para}
\texto{97.8967}{513.4130}{11.9552}{\normalfont}{las}
\texto{115.8869}{513.4130}{11.9552}{\normalfont}{agencias}
\texto{162.4930}{513.4130}{11.9552}{\normalfont}{y}
\texto{172.9395}{513.4130}{11.9552}{\normalfont}{operadores}
\texto{232.2253}{513.4130}{11.9552}{\normalfont}{turísticos}
\texto{283.4486}{513.4130}{11.9552}{\normalfont}{y}
\texto{293.9070}{513.4130}{11.9552}{\normalfont}{un}
\texto{311.1823}{513.4130}{11.9552}{\normalfont}{panel}
\texto{342.7655}{513.4130}{11.9552}{\normalfont}{de}
\texto{358.7400}{513.4130}{11.9552}{\normalfont}{administración}
\texto{438.5147}{513.4130}{11.9552}{\normalfont}{de}
\texto{454.4892}{513.4130}{11.9552}{\normalfont}{escritorio}
\texto{505.6479}{513.4130}{11.9552}{\normalfont}{para}
\texto{532.6906}{513.4130}{11.9552}{\normalfont}{el}
\texto{70.8660}{498.9670}{11.9552}{\normalfont}{equipo}
\texto{109.2494}{498.9670}{11.9552}{\normalfont}{gestor;}
\texto{147.6997}{498.9670}{11.9552}{\normalfont}{la}
\texto{161.3716}{498.9670}{11.9552}{\normalfont}{aplicación}
\texto{216.6668}{498.9670}{11.9552}{\normalfont}{móvil}
\texto{249.2017}{498.9670}{11.9552}{\normalfont}{para}
\texto{276.5313}{498.9670}{11.9552}{\normalfont}{el}
\texto{289.5529}{498.9670}{11.9552}{\normalfont}{turista}
\texto{328.0032}{498.9670}{11.9552}{\normalfont}{se}
\texto{342.3901}{498.9670}{11.9552}{\normalfont}{contempla}
\texto{399.3136}{498.9670}{11.9552}{\normalfont}{como}
\texto{430.5454}{498.9670}{11.9552}{\normalfont}{trabajo}
\texto{472.1925}{498.9670}{11.9552}{\normalfont}{futuro}
\texto{508.3020}{498.9670}{11.9552}{\normalfont}{dentro}
\texto{70.8660}{484.5210}{11.9552}{\normalfont}{de}
\texto{86.4699}{484.5210}{11.9552}{\normalfont}{la}
\texto{99.4724}{484.5210}{11.9552}{\normalfont}{misma}
\texto{136.6148}{484.5210}{11.9552}{\normalfont}{arquitectura.}
\texto{88.7990}{470.0750}{11.9552}{\normalfont}{La}
\texto{105.4454}{470.0750}{11.9552}{\normalfont}{investigación}
\texto{174.3671}{470.0750}{11.9552}{\normalfont}{aborda}
\texto{213.3005}{470.0750}{11.9552}{\normalfont}{el}
\texto{225.2461}{470.0750}{11.9552}{\normalfont}{diagnóstico}
\texto{286.0335}{470.0750}{11.9552}{\normalfont}{de}
\texto{301.2309}{470.0750}{11.9552}{\normalfont}{la}
\texto{313.8269}{470.0750}{11.9552}{\normalfont}{situación}
\texto{362.8959}{470.0750}{11.9552}{\normalfont}{actual}
\texto{397.6042}{470.0750}{11.9552}{\normalfont}{del}
\texto{416.0535}{470.0750}{11.9552}{\normalfont}{sector,}
\texto{452.7775}{470.0750}{11.9552}{\normalfont}{caracterizado}
\texto{523.9061}{470.0750}{11.9552}{\normalfont}{por}
\texto{70.8660}{455.6290}{11.9552}{\normalfont}{una}
\texto{94.3293}{455.6290}{11.9552}{\normalfont}{escasa}
\texto{130.2786}{455.6290}{11.9552}{\normalfont}{digitalización}
\texto{202.5191}{455.6290}{11.9552}{\normalfont}{de}
\texto{218.8283}{455.6290}{11.9552}{\normalfont}{la}
\texto{232.5362}{455.6290}{11.9552}{\normalfont}{gestión}
\texto{272.9734}{455.6290}{11.9552}{\normalfont}{de}
\texto{289.2827}{455.6290}{11.9552}{\normalfont}{los}
\texto{307.6077}{455.6290}{11.9552}{\normalfont}{operadores}
\texto{367.2402}{455.6290}{11.9552}{\normalfont}{y}
\texto{378.0214}{455.6290}{11.9552}{\normalfont}{por}
\texto{399.8564}{455.6290}{11.9552}{\normalfont}{la}
\texto{413.5642}{455.6290}{11.9552}{\normalfont}{dependencia}
\texto{480.9246}{455.6290}{11.9552}{\normalfont}{de}
\texto{497.2458}{455.6290}{11.9552}{\normalfont}{platafor-}
\texto{70.8660}{441.1840}{11.9552}{\normalfont}{mas}
\texto{94.7142}{441.1840}{11.9552}{\normalfont}{extranjeras}
\texto{154.9828}{441.1840}{11.9552}{\normalfont}{con}
\texto{176.1650}{441.1840}{11.9552}{\normalfont}{comisiones}
\texto{233.8967}{441.1840}{11.9552}{\normalfont}{elevadas}
\texto{279.5249}{441.1840}{11.9552}{\normalfont}{(15}
\texto{297.7805}{441.1840}{11.9552}{\normalfont}{\%--25}
\texto{327.0922}{441.1840}{11.9552}{\normalfont}{\%).}
\texto{348.2745}{441.1840}{11.9552}{\normalfont}{Se}
\texto{363.6034}{441.1840}{11.9552}{\normalfont}{propone}
\texto{408.5215}{441.1840}{11.9552}{\normalfont}{una}
\texto{431.0044}{441.1840}{11.9552}{\normalfont}{solución}
\texto{475.6643}{441.1840}{11.9552}{\normalfont}{con}
\texto{496.8465}{441.1840}{11.9552}{\normalfont}{comisión}
\texto{70.8660}{426.7380}{11.9552}{\normalfont}{reducida}
\texto{118.2899}{426.7380}{11.9552}{\normalfont}{del}
\texto{137.1098}{426.7380}{11.9552}{\normalfont}{8}
\texto{144.9476}{426.7380}{11.9552}{\normalfont}{\%,}
\texto{161.8164}{426.7380}{11.9552}{\normalfont}{verificación}
\texto{222.9003}{426.7380}{11.9552}{\normalfont}{documental}
\texto{285.9448}{426.7380}{11.9552}{\normalfont}{de}
\texto{301.5009}{426.7380}{11.9552}{\normalfont}{operadores,}
\texto{363.6321}{426.7380}{11.9552}{\normalfont}{reseñas}
\texto{404.0311}{426.7380}{11.9552}{\normalfont}{verificadas}
\texto{461.2894}{426.7380}{11.9552}{\normalfont}{y}
\texto{471.3174}{426.7380}{11.9552}{\normalfont}{control}
\texto{510.6261}{426.7380}{11.9552}{\normalfont}{de}
\texto{526.1822}{426.7380}{11.9552}{\normalfont}{ac-}
\texto{70.8660}{412.2920}{11.9552}{\normalfont}{ceso}
\texto{95.9026}{412.2920}{11.9552}{\normalfont}{basado}
\texto{135.2591}{412.2920}{11.9552}{\normalfont}{en}
\texto{151.1380}{412.2920}{11.9552}{\normalfont}{roles.}
\texto{182.0398}{412.2920}{11.9552}{\normalfont}{La}
\texto{199.3796}{412.2920}{11.9552}{\normalfont}{metodología}
\texto{265.6592}{412.2920}{11.9552}{\normalfont}{empleada}
\texto{317.9585}{412.2920}{11.9552}{\normalfont}{combina}
\texto{364.7320}{412.2920}{11.9552}{\normalfont}{la}
\texto{378.0094}{412.2920}{11.9552}{\normalfont}{investigación}
\texto{447.6007}{412.2920}{11.9552}{\normalfont}{documental}
\texto{510.9560}{412.2920}{11.9552}{\normalfont}{con}
\texto{532.6882}{412.2920}{11.9552}{\normalfont}{el}
\texto{70.8660}{397.8460}{11.9552}{\normalfont}{desarrollo}
\texto{123.9925}{397.8460}{11.9552}{\normalfont}{ágil}
\texto{145.8371}{397.8460}{11.9552}{\normalfont}{basado}
\texto{184.6556}{397.8460}{11.9552}{\normalfont}{en}
\texto{199.9965}{397.8460}{11.9552}{\normalfont}{Scrum.}
\texto{239.4008}{397.8460}{11.9552}{\normalfont}{Los}
\texto{260.8198}{397.8460}{11.9552}{\normalfont}{resultados}
\texto{315.9619}{397.8460}{11.9552}{\normalfont}{confirman}
\texto{370.3246}{397.8460}{11.9552}{\normalfont}{la}
\texto{383.0641}{397.8460}{11.9552}{\normalfont}{viabilidad}
\texto{437.1016}{397.8460}{11.9552}{\normalfont}{técnica}
\texto{476.5059}{397.8460}{11.9552}{\normalfont}{de}
\texto{491.8468}{397.8460}{11.9552}{\normalfont}{un}
\texto{508.4885}{397.8460}{11.9552}{\normalfont}{ecosis-}
\texto{70.8660}{383.4000}{11.9552}{\normalfont}{tema}
\texto{100.4623}{383.4000}{11.9552}{\normalfont}{coordinado}
\texto{160.9604}{383.4000}{11.9552}{\normalfont}{de}
\texto{176.9110}{383.4000}{11.9552}{\normalfont}{aplicaciones}
\texto{241.7034}{383.4000}{11.9552}{\normalfont}{sobre}
\texto{272.6650}{383.4000}{11.9552}{\normalfont}{una}
\texto{295.7696}{383.4000}{11.9552}{\normalfont}{arquitectura}
\texto{362.7737}{383.4000}{11.9552}{\normalfont}{orientada}
\texto{414.8099}{383.4000}{11.9552}{\normalfont}{a}
\texto{424.9073}{383.4000}{11.9552}{\normalfont}{servicios,}
\texto{475.1311}{383.4000}{11.9552}{\normalfont}{evidenciando}
\texto{70.8660}{368.9540}{11.9552}{\normalfont}{la}
\texto{84.6695}{368.9540}{11.9552}{\normalfont}{pertinencia}
\texto{146.2842}{368.9540}{11.9552}{\normalfont}{de}
\texto{162.7011}{368.9540}{11.9552}{\normalfont}{digitalizar}
\texto{218.7781}{368.9540}{11.9552}{\normalfont}{tanto}
\texto{250.4809}{368.9540}{11.9552}{\normalfont}{la}
\texto{264.2844}{368.9540}{11.9552}{\normalfont}{experiencia}
\texto{326.2243}{368.9540}{11.9552}{\normalfont}{del}
\texto{345.8930}{368.9540}{11.9552}{\normalfont}{turista}
\texto{384.4868}{368.9540}{11.9552}{\normalfont}{como}
\texto{415.8500}{368.9540}{11.9552}{\normalfont}{la}
\texto{429.6655}{368.9540}{11.9552}{\normalfont}{gestión}
\texto{470.2103}{368.9540}{11.9552}{\normalfont}{interna}
\texto{511.0063}{368.9540}{11.9552}{\normalfont}{de}
\texto{527.4231}{368.9540}{11.9552}{\normalfont}{los}
\texto{70.8660}{354.5090}{11.9552}{\normalfont}{operadores.}
\texto{88.7990}{340.0630}{11.9552}{\normalfont\bfseries}{Palabras}
\texto{143.2908}{340.0630}{11.9552}{\normalfont\bfseries}{clave:}
\texto{179.3170}{340.0630}{11.9552}{\normalfont}{turismo}
\texto{221.9421}{340.0630}{11.9552}{\normalfont}{digital,}
\texto{261.2508}{340.0630}{11.9552}{\normalfont}{marketplace,}
\texto{329.8282}{340.0630}{11.9552}{\normalfont}{backoffice}
\texto{382.4621}{340.0630}{11.9552}{\normalfont}{de}
\texto{397.7194}{340.0630}{11.9552}{\normalfont}{escritorio,}
\texto{451.4006}{340.0630}{11.9552}{\normalfont}{panel}
\texto{482.2546}{340.0630}{11.9552}{\normalfont}{de}
\texto{497.4999}{340.0630}{11.9552}{\normalfont}{adminis-}
\texto{70.8660}{325.6170}{11.9552}{\normalfont}{tración,}
\texto{113.7852}{325.6170}{11.9552}{\normalfont}{API}
\texto{138.6448}{325.6170}{11.9552}{\normalfont}{REST,}
\texto{177.3438}{325.6170}{11.9552}{\normalfont}{control}
\texto{216.6884}{325.6170}{11.9552}{\normalfont}{de}
\texto{232.2923}{325.6170}{11.9552}{\normalfont}{acceso}
\texto{268.1340}{325.6170}{11.9552}{\normalfont}{basado}
\texto{307.2155}{325.6170}{11.9552}{\normalfont}{en}
\texto{322.8194}{325.6170}{11.9552}{\normalfont}{roles.}
\texto{300.3150}{36.9930}{11.9552}{\normalfont\scshape}{iii}
\end{pagina}
% Página original 5
\begin{pagina}{5}
\texto{70.8660}{625.4930}{24.7871}{\normalfont\bfseries}{Abstract}
\texto{88.7990}{571.1960}{11.9552}{\normalfont}{This}
\texto{116.2171}{571.1960}{11.9552}{\normalfont}{Integrador}
\texto{174.1472}{571.1960}{11.9552}{\normalfont}{II}
\texto{187.5179}{571.1960}{11.9552}{\normalfont}{project}
\texto{227.8882}{571.1960}{11.9552}{\normalfont}{develops}
\texto{275.4818}{571.1960}{11.9552}{\normalfont}{the}
\texto{296.3317}{571.1960}{11.9552}{\normalfont}{prototype}
\texto{350.6872}{571.1960}{11.9552}{\normalfont}{of}
\texto{364.7202}{571.1960}{11.9552}{\normalfont}{the}
\texto{385.5810}{571.1960}{11.9552}{\normalfont\bfseries}{desktop}
\texto{436.3057}{571.1960}{11.9552}{\normalfont\bfseries}{applications}
\texto{510.8450}{571.1960}{11.9552}{\normalfont}{of}
\texto{524.8661}{571.1960}{11.9552}{\normalfont}{the}
\texto{70.8660}{556.7500}{11.9552}{\normalfont}{technological}
\texto{142.6307}{556.7500}{11.9552}{\normalfont}{ecosystem}
\texto{199.6350}{556.7500}{11.9552}{\normalfont\bfseries}{Explora-Nicaragua}
\texto{309.5420}{556.7500}{11.9552}{\normalfont}{,}
\texto{318.5562}{556.7500}{11.9552}{\normalfont}{focused}
\texto{362.1018}{556.7500}{11.9552}{\normalfont}{on}
\texto{380.2092}{556.7500}{11.9552}{\normalfont}{the}
\texto{402.2307}{556.7500}{11.9552}{\normalfont}{promotion}
\texto{460.6725}{556.7500}{11.9552}{\normalfont}{and}
\texto{485.2954}{556.7500}{11.9552}{\normalfont}{booking}
\texto{531.7007}{556.7500}{11.9552}{\normalfont}{of}
\texto{70.8660}{542.3040}{11.9552}{\normalfont}{tourist}
\texto{108.9457}{542.3040}{11.9552}{\normalfont}{tours}
\texto{139.2330}{542.3040}{11.9552}{\normalfont}{in}
\texto{153.1847}{542.3040}{11.9552}{\normalfont}{Nicaragua.}
\texto{212.3438}{542.3040}{11.9552}{\normalfont}{The}
\texto{236.7014}{542.3040}{11.9552}{\normalfont}{initial}
\texto{270.8263}{542.3040}{11.9552}{\normalfont}{scope}
\texto{302.7371}{542.3040}{11.9552}{\normalfont}{is}
\texto{314.8023}{542.3040}{11.9552}{\normalfont}{limited}
\texto{354.7805}{542.3040}{11.9552}{\normalfont}{to}
\texto{369.3826}{542.3040}{11.9552}{\normalfont}{two}
\texto{391.8057}{542.3040}{11.9552}{\normalfont}{desktop}
\texto{435.4135}{542.3040}{11.9552}{\normalfont}{applications}
\texto{500.8204}{542.3040}{11.9552}{\normalfont}{working}
\texto{70.8660}{527.8590}{11.9552}{\normalfont}{over}
\texto{95.5463}{527.8590}{11.9552}{\normalfont}{a}
\texto{104.9503}{527.8590}{11.9552}{\normalfont}{single}
\texto{137.1815}{527.8590}{11.9552}{\normalfont}{REST}
\texto{172.2700}{527.8590}{11.9552}{\normalfont}{API:}
\texto{200.0348}{527.8590}{11.9552}{\normalfont}{a}
\texto{209.4387}{527.8590}{11.9552}{\normalfont}{desktop}
\texto{252.4009}{527.8590}{11.9552}{\normalfont}{backoffice}
\texto{305.0468}{527.8590}{11.9552}{\normalfont}{for}
\texto{322.5803}{527.8590}{11.9552}{\normalfont}{tour}
\texto{347.5930}{527.8590}{11.9552}{\normalfont}{agencies}
\texto{392.8315}{527.8590}{11.9552}{\normalfont}{and}
\texto{415.2427}{527.8590}{11.9552}{\normalfont}{operators,}
\texto{469.9091}{527.8590}{11.9552}{\normalfont}{and}
\texto{492.3203}{527.8590}{11.9552}{\normalfont}{a}
\texto{501.7242}{527.8590}{11.9552}{\normalfont}{desktop}
\texto{70.8660}{513.4130}{11.9552}{\normalfont}{administration}
\texto{148.8187}{513.4130}{11.9552}{\normalfont}{panel}
\texto{179.2303}{513.4130}{11.9552}{\normalfont}{for}
\texto{196.3095}{513.4130}{11.9552}{\normalfont}{the}
\texto{215.6769}{513.4130}{11.9552}{\normalfont}{management}
\texto{283.4868}{513.4130}{11.9552}{\normalfont}{team;}
\texto{315.1992}{513.4130}{11.9552}{\normalfont}{the}
\texto{334.5666}{513.4130}{11.9552}{\normalfont}{mobile}
\texto{371.4819}{513.4130}{11.9552}{\normalfont}{application}
\texto{431.1599}{513.4130}{11.9552}{\normalfont}{for}
\texto{448.2510}{513.4130}{11.9552}{\normalfont}{tourists}
\texto{489.8479}{513.4130}{11.9552}{\normalfont}{is}
\texto{500.8132}{513.4130}{11.9552}{\normalfont}{planned}
\texto{70.8660}{498.9670}{11.9552}{\normalfont}{as}
\texto{85.2338}{498.9670}{11.9552}{\normalfont}{future}
\texto{120.0234}{498.9670}{11.9552}{\normalfont}{work}
\texto{148.6370}{498.9670}{11.9552}{\normalfont}{within}
\texto{185.0645}{498.9670}{11.9552}{\normalfont}{the}
\texto{205.2209}{498.9670}{11.9552}{\normalfont}{same}
\texto{234.5590}{498.9670}{11.9552}{\normalfont}{architecture.}
\texto{88.7990}{484.5210}{11.9552}{\normalfont}{The}
\texto{112.8696}{484.5210}{11.9552}{\normalfont}{research}
\texto{158.1439}{484.5210}{11.9552}{\normalfont}{addresses}
\texto{209.7354}{484.5210}{11.9552}{\normalfont}{the}
\texto{229.9038}{484.5210}{11.9552}{\normalfont}{diagnosis}
\texto{280.1300}{484.5210}{11.9552}{\normalfont}{of}
\texto{293.4696}{484.5210}{11.9552}{\normalfont}{the}
\texto{313.6500}{484.5210}{11.9552}{\normalfont}{current}
\texto{354.3073}{484.5210}{11.9552}{\normalfont}{situation}
\texto{403.1562}{484.5210}{11.9552}{\normalfont}{of}
\texto{416.5078}{484.5210}{11.9552}{\normalfont}{the}
\texto{436.6762}{484.5210}{11.9552}{\normalfont}{sector,}
\texto{473.8306}{484.5210}{11.9552}{\normalfont}{characterized}
\texto{70.8660}{470.0750}{11.9552}{\normalfont}{by}
\texto{88.8203}{470.0750}{11.9552}{\normalfont}{a}
\texto{100.2806}{470.0750}{11.9552}{\normalfont}{low}
\texto{123.1126}{470.0750}{11.9552}{\normalfont}{level}
\texto{151.4847}{470.0750}{11.9552}{\normalfont}{of}
\texto{166.5100}{470.0750}{11.9552}{\normalfont}{digitalization}
\texto{239.1044}{470.0750}{11.9552}{\normalfont}{in}
\texto{254.4668}{470.0750}{11.9552}{\normalfont}{operators'}
\texto{311.1894}{470.0750}{11.9552}{\normalfont}{management}
\texto{381.4980}{470.0750}{11.9552}{\normalfont}{and}
\texto{405.9655}{470.0750}{11.9552}{\normalfont}{by}
\texto{423.9198}{470.0750}{11.9552}{\normalfont}{dependence}
\texto{488.3822}{470.0750}{11.9552}{\normalfont}{on}
\texto{506.3461}{470.0750}{11.9552}{\normalfont}{foreign}
\texto{70.8660}{455.6290}{11.9552}{\normalfont}{platforms}
\texto{124.3440}{455.6290}{11.9552}{\normalfont}{with}
\texto{152.0801}{455.6290}{11.9552}{\normalfont}{high}
\texto{179.1658}{455.6290}{11.9552}{\normalfont}{commissions}
\texto{247.4180}{455.6290}{11.9552}{\normalfont}{(15}
\texto{265.6616}{455.6290}{11.9552}{\normalfont}{\%--25}
\texto{294.9734}{455.6290}{11.9552}{\normalfont}{\%).}
\texto{317.5066}{455.6290}{11.9552}{\normalfont}{A}
\texto{331.2550}{455.6290}{11.9552}{\normalfont}{solution}
\texto{376.6155}{455.6290}{11.9552}{\normalfont}{is}
\texto{389.4458}{455.6290}{11.9552}{\normalfont}{proposed}
\texto{440.3319}{455.6290}{11.9552}{\normalfont}{with}
\texto{468.0680}{455.6290}{11.9552}{\normalfont}{a}
\texto{478.8946}{455.6290}{11.9552}{\normalfont}{reduced}
\texto{523.5401}{455.6290}{11.9552}{\normalfont}{8}
\texto{531.3779}{455.6290}{11.9552}{\normalfont}{\%}
\texto{70.8660}{441.1840}{11.9552}{\normalfont}{commission,}
\texto{138.2192}{441.1840}{11.9552}{\normalfont}{documentary}
\texto{210.3330}{441.1840}{11.9552}{\normalfont}{verification}
\texto{272.3685}{441.1840}{11.9552}{\normalfont}{of}
\texto{287.2503}{441.1840}{11.9552}{\normalfont}{operators,}
\texto{343.8176}{441.1840}{11.9552}{\normalfont}{verified}
\texto{386.3303}{441.1840}{11.9552}{\normalfont}{reviews}
\texto{429.2422}{441.1840}{11.9552}{\normalfont}{and}
\texto{453.5543}{441.1840}{11.9552}{\normalfont}{role-based}
\texto{510.4491}{441.1840}{11.9552}{\normalfont}{access}
\texto{70.8660}{426.7380}{11.9552}{\normalfont}{control.}
\texto{113.5461}{426.7380}{11.9552}{\normalfont}{The}
\texto{137.6884}{426.7380}{11.9552}{\normalfont}{methodology}
\texto{207.3657}{426.7380}{11.9552}{\normalfont}{combines}
\texto{257.9147}{426.7380}{11.9552}{\normalfont}{documentary}
\texto{328.5699}{426.7380}{11.9552}{\normalfont}{research}
\texto{373.9160}{426.7380}{11.9552}{\normalfont}{with}
\texto{400.6717}{426.7380}{11.9552}{\normalfont}{agile}
\texto{428.0659}{426.7380}{11.9552}{\normalfont}{development}
\texto{496.1125}{426.7380}{11.9552}{\normalfont}{based}
\texto{528.7860}{426.7380}{11.9552}{\normalfont}{on}
\texto{70.8660}{412.2920}{11.9552}{\normalfont}{Scrum.}
\texto{109.5411}{412.2920}{11.9552}{\normalfont}{The}
\texto{132.5955}{412.2920}{11.9552}{\normalfont}{results}
\texto{168.7982}{412.2920}{11.9552}{\normalfont}{confirm}
\texto{210.0628}{412.2920}{11.9552}{\normalfont}{the}
\texto{229.2270}{412.2920}{11.9552}{\normalfont}{technical}
\texto{277.3227}{412.2920}{11.9552}{\normalfont}{feasibility}
\texto{329.3852}{412.2920}{11.9552}{\normalfont}{of}
\texto{341.7206}{412.2920}{11.9552}{\normalfont}{a}
\texto{350.4670}{412.2920}{11.9552}{\normalfont}{coordinated}
\texto{413.5283}{412.2920}{11.9552}{\normalfont}{ecosystem}
\texto{467.6160}{412.2920}{11.9552}{\normalfont}{of}
\texto{479.9395}{412.2920}{11.9552}{\normalfont}{applications}
\texto{70.8660}{397.8460}{11.9552}{\normalfont}{over}
\texto{95.9528}{397.8460}{11.9552}{\normalfont}{a}
\texto{105.7632}{397.8460}{11.9552}{\normalfont}{service-oriented}
\texto{189.1316}{397.8460}{11.9552}{\normalfont}{architecture,}
\texto{257.1400}{397.8460}{11.9552}{\normalfont}{evidencing}
\texto{314.7521}{397.8460}{11.9552}{\normalfont}{the}
\texto{334.9683}{397.8460}{11.9552}{\normalfont}{relevance}
\texto{385.4312}{397.8460}{11.9552}{\normalfont}{of}
\texto{398.8187}{397.8460}{11.9552}{\normalfont}{digitalizing}
\texto{459.3574}{397.8460}{11.9552}{\normalfont}{both}
\texto{487.0504}{397.8460}{11.9552}{\normalfont}{the}
\texto{507.2667}{397.8460}{11.9552}{\normalfont}{tourist}
\texto{70.8660}{383.4000}{11.9552}{\normalfont}{experience}
\texto{128.0908}{383.4000}{11.9552}{\normalfont}{and}
\texto{150.8487}{383.4000}{11.9552}{\normalfont}{the}
\texto{171.0171}{383.4000}{11.9552}{\normalfont}{internal}
\texto{214.2638}{383.4000}{11.9552}{\normalfont}{management}
\texto{282.8747}{383.4000}{11.9552}{\normalfont}{of}
\texto{296.2143}{383.4000}{11.9552}{\normalfont}{operators.}
\texto{88.7990}{368.9540}{11.9552}{\normalfont\bfseries}{Keywords:}
\texto{153.3740}{368.9540}{11.9552}{\normalfont}{digital}
\texto{188.6179}{368.9540}{11.9552}{\normalfont}{tourism,}
\texto{233.6938}{368.9540}{11.9552}{\normalfont}{marketplace,}
\texto{301.4702}{368.9540}{11.9552}{\normalfont}{desktop}
\texto{343.6195}{368.9540}{11.9552}{\normalfont}{backoffice,}
\texto{398.7043}{368.9540}{11.9552}{\normalfont}{administration}
\texto{476.2983}{368.9540}{11.9552}{\normalfont}{panel,}
\texto{509.6031}{368.9540}{11.9552}{\normalfont}{REST}
\texto{70.8660}{354.5090}{11.9552}{\normalfont}{API,}
\texto{98.9775}{354.5090}{11.9552}{\normalfont}{role-based}
\texto{154.3181}{354.5090}{11.9552}{\normalfont}{access}
\texto{188.9236}{354.5090}{11.9552}{\normalfont}{control.}
\texto{300.5270}{36.9930}{11.9552}{\normalfont\scshape}{iv}
\end{pagina}

\clearpage\restoregeometry
\pagestyle{plain}
\hypersetup{pageanchor=true}
\begingroup\small\setstretch{1.0}\tableofcontents\endgroup
\clearpage\newgeometry{margin=0pt}
\pagenumbering{arabic}
\pagestyle{empty}
% Texto editable; coordenadas en puntos PDF.
% Página original 11
\begin{pagina}{11}
\texto{70.8660}{625.4930}{24.7871}{\normalfont\bfseries}{Capítulo}
\texto{185.1197}{625.4930}{24.7871}{\normalfont\bfseries}{1}
\texto{70.8660}{575.6790}{24.7871}{\normalfont\bfseries}{F}
\texto{86.0902}{575.6790}{24.7871}{\normalfont\bfseries}{ormulación}
\texto{229.5307}{575.6790}{24.7871}{\normalfont\bfseries}{del}
\texto{274.7870}{575.6790}{24.7871}{\normalfont\bfseries}{proyecto}
\texto{303.0730}{36.9930}{11.9552}{\normalfont}{1}
\end{pagina}
% Página original 12
\begin{pagina}{12}
\texto{70.8660}{746.3890}{11.9552}{\normalfont}{2}
\texto{266.5910}{746.3890}{11.9552}{\normalfont}{CAPÍTULO}
\texto{333.5545}{746.3890}{11.9552}{\normalfont}{1.}
\texto{350.4663}{746.3890}{11.9552}{\normalfont}{FORMULACIÓN}
\texto{445.2256}{746.3890}{11.9552}{\normalfont}{DEL}
\texto{473.3418}{746.3890}{11.9552}{\normalfont}{PROYECTO}
\texto{70.8660}{705.1940}{17.2154}{\normalfont\bfseries}{1.1.}
\texto{120.3603}{705.1940}{17.2154}{\normalfont\bfseries}{Introducción}
\texto{88.7990}{678.9090}{11.9552}{\normalfont}{El}
\texto{103.9965}{678.9090}{11.9552}{\normalfont}{turismo}
\texto{147.0758}{678.9090}{11.9552}{\normalfont}{constituye}
\texto{203.1648}{678.9090}{11.9552}{\normalfont}{una}
\texto{226.0064}{678.9090}{11.9552}{\normalfont}{actividad}
\texto{277.1508}{678.9090}{11.9552}{\normalfont}{económica}
\texto{333.8232}{678.9090}{11.9552}{\normalfont}{de}
\texto{349.5228}{678.9090}{11.9552}{\normalfont}{creciente}
\texto{398.0561}{678.9090}{11.9552}{\normalfont}{relevancia}
\texto{452.4571}{678.9090}{11.9552}{\normalfont}{para}
\texto{479.2008}{678.9090}{11.9552}{\normalfont}{el}
\texto{491.6486}{678.9090}{11.9552}{\normalfont}{desarrollo}
\texto{70.8660}{664.4630}{11.9552}{\normalfont}{sostenible}
\texto{123.6626}{664.4630}{11.9552}{\normalfont}{de}
\texto{138.6089}{664.4630}{11.9552}{\normalfont}{Centroamérica,}
\texto{219.5696}{664.4630}{11.9552}{\normalfont}{y}
\texto{228.9879}{664.4630}{11.9552}{\normalfont}{Nicaragua}
\texto{283.9268}{664.4630}{11.9552}{\normalfont}{no}
\texto{299.5235}{664.4630}{11.9552}{\normalfont}{es}
\texto{312.5834}{664.4630}{11.9552}{\normalfont}{una}
\texto{334.6838}{664.4630}{11.9552}{\normalfont}{excepción.}
\texto{390.6126}{664.4630}{11.9552}{\normalfont}{El}
\texto{405.0688}{664.4630}{11.9552}{\normalfont}{país}
\texto{428.5345}{664.4630}{11.9552}{\normalfont}{destaca}
\texto{469.5600}{664.4630}{11.9552}{\normalfont}{por}
\texto{490.0320}{664.4630}{11.9552}{\normalfont}{su}
\texto{504.3926}{664.4630}{11.9552}{\normalfont}{riqueza}
\texto{70.8660}{650.0170}{11.9552}{\normalfont}{natural}
\texto{112.9937}{650.0170}{11.9552}{\normalfont}{y}
\texto{124.2412}{650.0170}{11.9552}{\normalfont}{cultural:}
\texto{172.2341}{650.0170}{11.9552}{\normalfont}{la}
\texto{186.3963}{650.0170}{11.9552}{\normalfont}{isla}
\texto{208.4393}{650.0170}{11.9552}{\normalfont}{de}
\texto{225.2028}{650.0170}{11.9552}{\normalfont}{Ometepe,}
\texto{279.3719}{650.0170}{11.9552}{\normalfont}{el}
\texto{292.8956}{650.0170}{11.9552}{\normalfont}{Complejo}
\texto{346.4047}{650.0170}{11.9552}{\normalfont}{Volcánico}
\texto{400.2413}{650.0170}{11.9552}{\normalfont}{de}
\texto{417.0169}{650.0170}{11.9552}{\normalfont}{Momotombo,}
\texto{490.0440}{650.0170}{11.9552}{\normalfont}{el}
\texto{503.5558}{650.0170}{11.9552}{\normalfont}{río}
\texto{522.2824}{650.0170}{11.9552}{\normalfont}{San}
\texto{70.8660}{635.5710}{11.9552}{\normalfont}{Juan,}
\texto{103.9054}{635.5710}{11.9552}{\normalfont}{las}
\texto{122.5412}{635.5710}{11.9552}{\normalfont}{ruinas}
\texto{158.7487}{635.5710}{11.9552}{\normalfont}{prehispánicas}
\texto{232.0795}{635.5710}{11.9552}{\normalfont}{de}
\texto{248.7116}{635.5710}{11.9552}{\normalfont}{Solentiname}
\texto{315.7372}{635.5710}{11.9552}{\normalfont}{y}
\texto{326.8292}{635.5710}{11.9552}{\normalfont}{las}
\texto{345.4769}{635.5710}{11.9552}{\normalfont}{playas}
\texto{382.0025}{635.5710}{11.9552}{\normalfont}{del}
\texto{401.8744}{635.5710}{11.9552}{\normalfont}{Pacífico}
\texto{446.2975}{635.5710}{11.9552}{\normalfont}{son}
\texto{468.1851}{635.5710}{11.9552}{\normalfont}{atractivos}
\texto{523.2556}{635.5710}{11.9552}{\normalfont}{que}
\texto{70.8660}{621.1260}{11.9552}{\normalfont}{atraen,}
\texto{110.5214}{621.1260}{11.9552}{\normalfont}{cada}
\texto{137.8199}{621.1260}{11.9552}{\normalfont}{año,}
\texto{163.1673}{621.1260}{11.9552}{\normalfont}{a}
\texto{172.9060}{621.1260}{11.9552}{\normalfont}{una}
\texto{195.6520}{621.1260}{11.9552}{\normalfont}{cantidad}
\texto{243.4393}{621.1260}{11.9552}{\normalfont}{creciente}
\texto{291.8770}{621.1260}{11.9552}{\normalfont}{de}
\texto{307.4809}{621.1260}{11.9552}{\normalfont}{visitantes}
\texto{359.6247}{621.1260}{11.9552}{\normalfont}{internacionales}
\texto{439.3444}{621.1260}{11.9552}{\normalfont}{y}
\texto{449.4083}{621.1260}{11.9552}{\normalfont}{nacionales}
\texto{505.3873}{621.1260}{11.9552}{\normalfont}{[1].}
\texto{524.8741}{621.1260}{11.9552}{\normalfont}{Sin}
\texto{70.8660}{606.6800}{11.9552}{\normalfont}{embargo,}
\texto{122.1514}{606.6800}{11.9552}{\normalfont}{la}
\texto{136.0386}{606.6800}{11.9552}{\normalfont}{digitalización}
\texto{208.4704}{606.6800}{11.9552}{\normalfont}{de}
\texto{224.9590}{606.6800}{11.9552}{\normalfont}{esta}
\texto{249.9788}{606.6800}{11.9552}{\normalfont}{oferta}
\texto{284.3524}{606.6800}{11.9552}{\normalfont}{(es}
\texto{303.5190}{606.6800}{11.9552}{\normalfont}{decir,}
\texto{336.2667}{606.6800}{11.9552}{\normalfont}{la}
\texto{350.1658}{606.6800}{11.9552}{\normalfont}{posibilidad}
\texto{410.6160}{606.6800}{11.9552}{\normalfont}{de}
\texto{427.1166}{606.6800}{11.9552}{\normalfont}{descubrir,}
\texto{482.0531}{606.6800}{11.9552}{\normalfont}{comparar}
\texto{534.9621}{606.6800}{11.9552}{\normalfont}{y}
\texto{70.8660}{592.2340}{11.9552}{\normalfont}{reservar}
\texto{114.2921}{592.2340}{11.9552}{\normalfont}{experiencias}
\texto{179.5101}{592.2340}{11.9552}{\normalfont}{turísticas}
\texto{229.8367}{592.2340}{11.9552}{\normalfont}{mediante}
\texto{279.7114}{592.2340}{11.9552}{\normalfont}{herramientas}
\texto{349.1496}{592.2340}{11.9552}{\normalfont}{tecnológicas)}
\texto{418.2722}{592.2340}{11.9552}{\normalfont}{aún}
\texto{440.5040}{592.2340}{11.9552}{\normalfont}{se}
\texto{453.6954}{592.2340}{11.9552}{\normalfont}{encuentra}
\texto{506.8219}{592.2340}{11.9552}{\normalfont}{en}
\texto{521.8879}{592.2340}{11.9552}{\normalfont}{fase}
\texto{70.8660}{577.7880}{11.9552}{\normalfont}{incipiente,}
\texto{126.7016}{577.7880}{11.9552}{\normalfont}{tanto}
\texto{157.1849}{577.7880}{11.9552}{\normalfont}{en}
\texto{172.3824}{577.7880}{11.9552}{\normalfont}{el}
\texto{184.3280}{577.7880}{11.9552}{\normalfont}{punto}
\texto{217.4128}{577.7880}{11.9552}{\normalfont}{de}
\texto{232.6103}{577.7880}{11.9552}{\normalfont}{venta}
\texto{263.7464}{577.7880}{11.9552}{\normalfont}{hacia}
\texto{293.9022}{577.7880}{11.9552}{\normalfont}{el}
\texto{305.8478}{577.7880}{11.9552}{\normalfont}{turista}
\texto{343.2222}{577.7880}{11.9552}{\normalfont}{como}
\texto{373.3899}{577.7880}{11.9552}{\normalfont}{en}
\texto{388.5874}{577.7880}{11.9552}{\normalfont}{la}
\texto{401.1834}{577.7880}{11.9552}{\normalfont}{gestión}
\texto{440.5088}{577.7880}{11.9552}{\normalfont}{interna}
\texto{480.0973}{577.7880}{11.9552}{\normalfont}{que}
\texto{501.4732}{577.7880}{11.9552}{\normalfont}{realizan}
\texto{70.8660}{563.3420}{11.9552}{\normalfont}{los}
\texto{88.4856}{563.3420}{11.9552}{\normalfont}{operadores.}
\texto{88.7990}{548.8960}{11.9552}{\normalfont}{Este}
\texto{115.4519}{548.8960}{11.9552}{\normalfont}{proyecto}
\texto{163.0336}{548.8960}{11.9552}{\normalfont}{propone}
\texto{208.6571}{548.8960}{11.9552}{\normalfont}{el}
\texto{221.4396}{548.8960}{11.9552}{\normalfont}{desarrollo}
\texto{275.2595}{548.8960}{11.9552}{\normalfont}{de}
\texto{291.2980}{548.8960}{11.9552}{\normalfont\bfseries}{Explora-Nicaragua}
\texto{401.2050}{548.8960}{11.9552}{\normalfont}{,}
\texto{408.7846}{548.8960}{11.9552}{\normalfont}{un}
\texto{426.1196}{548.8960}{11.9552}{\normalfont}{ecosistema}
\texto{484.5567}{548.8960}{11.9552}{\normalfont}{tecnológico}
\texto{70.8660}{534.4510}{11.9552}{\normalfont}{integral}
\texto{114.1438}{534.4510}{11.9552}{\normalfont}{orientado}
\texto{166.5387}{534.4510}{11.9552}{\normalfont}{a}
\texto{176.9947}{534.4510}{11.9552}{\normalfont}{la}
\texto{190.6906}{534.4510}{11.9552}{\normalfont}{promoción}
\texto{248.9339}{534.4510}{11.9552}{\normalfont}{y}
\texto{259.7032}{534.4510}{11.9552}{\normalfont}{reserva}
\texto{299.8081}{534.4510}{11.9552}{\normalfont}{de}
\texto{316.1054}{534.4510}{11.9552}{\normalfont}{tours}
\texto{346.7753}{534.4510}{11.9552}{\normalfont}{turísticos}
\texto{398.3213}{534.4510}{11.9552}{\normalfont}{en}
\texto{414.6306}{534.4510}{11.9552}{\normalfont}{Nicaragua.}
\texto{474.1723}{534.4510}{11.9552}{\normalfont}{La}
\texto{491.9305}{534.4510}{11.9552}{\normalfont}{visión}
\texto{526.1774}{534.4510}{11.9552}{\normalfont}{del}
\texto{70.8660}{520.0050}{11.9552}{\normalfont}{ecosistema}
\texto{130.2594}{520.0050}{11.9552}{\normalfont}{contempla}
\texto{187.9122}{520.0050}{11.9552}{\normalfont}{tres}
\texto{212.1215}{520.0050}{11.9552}{\normalfont}{aplicaciones}
\texto{277.9659}{520.0050}{11.9552}{\normalfont}{que}
\texto{301.1471}{520.0050}{11.9552}{\normalfont}{trabajan}
\texto{350.0032}{520.0050}{11.9552}{\normalfont}{de}
\texto{367.0059}{520.0050}{11.9552}{\normalfont}{forma}
\texto{401.8935}{520.0050}{11.9552}{\normalfont}{coordinada}
\texto{463.4437}{520.0050}{11.9552}{\normalfont}{sobre}
\texto{495.4573}{520.0050}{11.9552}{\normalfont}{una}
\texto{519.6140}{520.0050}{11.9552}{\normalfont}{mis-}
\texto{70.8660}{505.5590}{11.9552}{\normalfont}{ma}
\texto{91.4720}{505.5590}{11.9552}{\normalfont}{API}
\texto{117.4315}{505.5590}{11.9552}{\normalfont}{REST:}
\texto{157.2304}{505.5590}{11.9552}{\normalfont}{una}
\texto{181.0882}{505.5590}{11.9552}{\normalfont}{aplicación}
\texto{236.8137}{505.5590}{11.9552}{\normalfont}{móvil}
\texto{269.7910}{505.5590}{11.9552}{\normalfont}{transaccional}
\texto{341.8402}{505.5590}{11.9552}{\normalfont}{para}
\texto{369.6001}{505.5590}{11.9552}{\normalfont}{que}
\texto{392.4944}{505.5590}{11.9552}{\normalfont}{el}
\texto{405.9463}{505.5590}{11.9552}{\normalfont}{turista}
\texto{444.8270}{505.5590}{11.9552}{\normalfont}{descubra}
\texto{494.7759}{505.5590}{11.9552}{\normalfont}{y}
\texto{505.9516}{505.5590}{11.9552}{\normalfont}{reserve}
\texto{70.8660}{491.1130}{11.9552}{\normalfont}{experiencias,}
\texto{140.9737}{491.1130}{11.9552}{\normalfont}{un}
\texto{158.9902}{491.1130}{11.9552}{\normalfont}{panel}
\texto{191.3146}{491.1130}{11.9552}{\normalfont}{de}
\texto{208.0184}{491.1130}{11.9552}{\normalfont}{administración}
\texto{288.5343}{491.1130}{11.9552}{\normalfont}{web}
\texto{313.3820}{491.1130}{11.9552}{\normalfont}{para}
\texto{341.1539}{491.1130}{11.9552}{\normalfont}{que}
\texto{364.0482}{491.1130}{11.9552}{\normalfont}{el}
\texto{377.5121}{491.1130}{11.9552}{\normalfont}{equipo}
\texto{416.3378}{491.1130}{11.9552}{\normalfont}{gestor}
\texto{451.9786}{491.1130}{11.9552}{\normalfont}{supervise}
\texto{503.9407}{491.1130}{11.9552}{\normalfont}{toda}
\texto{532.0354}{491.1130}{11.9552}{\normalfont}{la}
\texto{70.8660}{476.6670}{11.9552}{\normalfont}{plataforma,}
\texto{133.7527}{476.6670}{11.9552}{\normalfont}{y}
\texto{143.9720}{476.6670}{11.9552}{\normalfont}{un}
\texto{161.0082}{476.6670}{11.9552}{\normalfont}{backoffice}
\texto{214.1443}{476.6670}{11.9552}{\normalfont}{de}
\texto{229.8917}{476.6670}{11.9552}{\normalfont}{escritorio}
\texto{280.8232}{476.6670}{11.9552}{\normalfont}{(}
\texto{285.3700}{476.6670}{11.9552}{\normalfont\itshape}{desktop}
\texto{323.0500}{476.6670}{11.9552}{\normalfont}{)}
\texto{331.6314}{476.6670}{11.9552}{\normalfont}{para}
\texto{358.4350}{476.6670}{11.9552}{\normalfont}{que}
\texto{380.3489}{476.6670}{11.9552}{\normalfont}{las}
\texto{398.1119}{476.6670}{11.9552}{\normalfont}{agencias}
\texto{444.4790}{476.6670}{11.9552}{\normalfont}{y}
\texto{454.6983}{476.6670}{11.9552}{\normalfont}{operadores}
\texto{513.7570}{476.6670}{11.9552}{\normalfont}{turís-}
\texto{70.8660}{462.2210}{11.9552}{\normalfont}{ticos}
\texto{98.2291}{462.2210}{11.9552}{\normalfont}{gestionen}
\texto{149.6675}{462.2210}{11.9552}{\normalfont}{sus}
\texto{169.2908}{462.2210}{11.9552}{\normalfont}{propios}
\texto{210.3234}{462.2210}{11.9552}{\normalfont}{tours,}
\texto{243.5397}{462.2210}{11.9552}{\normalfont}{disponibilidad}
\texto{319.3644}{462.2210}{11.9552}{\normalfont}{y}
\texto{329.4402}{462.2210}{11.9552}{\normalfont}{reservas.}
\texto{376.7070}{462.2210}{11.9552}{\normalfont\bfseries}{El}
\texto{393.7420}{462.2210}{11.9552}{\normalfont\bfseries}{alcance}
\texto{440.5896}{462.2210}{11.9552}{\normalfont\bfseries}{de}
\texto{458.6695}{462.2210}{11.9552}{\normalfont\bfseries}{este}
\texto{485.9500}{462.2210}{11.9552}{\normalfont\bfseries}{prototipo}
\texto{70.8660}{447.7760}{11.9552}{\normalfont\bfseries}{inicial}
\texto{110.9625}{447.7760}{11.9552}{\normalfont\bfseries}{se}
\texto{127.5815}{447.7760}{11.9552}{\normalfont\bfseries}{limita}
\texto{166.9428}{447.7760}{11.9552}{\normalfont\bfseries}{a}
\texto{178.6577}{447.7760}{11.9552}{\normalfont\bfseries}{las}
\texto{199.4023}{447.7760}{11.9552}{\normalfont\bfseries}{dos}
\texto{224.0815}{447.7760}{11.9552}{\normalfont\bfseries}{aplicaciones}
\texto{298.6090}{447.7760}{11.9552}{\normalfont\bfseries}{de}
\texto{317.3942}{447.7760}{11.9552}{\normalfont\bfseries}{escritorio}
\texto{376.4550}{447.7760}{11.9552}{\normalfont}{---}
\texto{392.6567}{447.7760}{11.9552}{\normalfont}{el}
\texto{405.6185}{447.7760}{11.9552}{\normalfont}{backoffice}
\texto{459.2208}{447.7760}{11.9552}{\normalfont}{de}
\texto{475.4345}{447.7760}{11.9552}{\normalfont}{operadores}
\texto{534.9594}{447.7760}{11.9552}{\normalfont}{y}
\texto{70.8660}{433.3300}{11.9552}{\normalfont}{el}
\texto{83.0388}{433.3300}{11.9552}{\normalfont}{panel}
\texto{114.0721}{433.3300}{11.9552}{\normalfont}{de}
\texto{129.5086}{433.3300}{11.9552}{\normalfont}{administración---,}
\texto{223.6917}{433.3300}{11.9552}{\normalfont}{en}
\texto{239.1163}{433.3300}{11.9552}{\normalfont}{tanto}
\texto{269.8268}{433.3300}{11.9552}{\normalfont}{la}
\texto{282.6619}{433.3300}{11.9552}{\normalfont}{aplicación}
\texto{337.1083}{433.3300}{11.9552}{\normalfont}{móvil}
\texto{368.7944}{433.3300}{11.9552}{\normalfont}{se}
\texto{382.3444}{433.3300}{11.9552}{\normalfont}{contempla}
\texto{438.4191}{433.3300}{11.9552}{\normalfont}{como}
\texto{468.8020}{433.3300}{11.9552}{\normalfont}{trabajo}
\texto{509.6003}{433.3300}{11.9552}{\normalfont}{futuro}
\texto{70.8660}{418.8840}{11.9552}{\normalfont}{mediante}
\texto{121.2667}{418.8840}{11.9552}{\normalfont}{integraciones}
\texto{191.2429}{418.8840}{11.9552}{\normalfont}{externas.}
\texto{88.7990}{404.4380}{11.9552}{\normalfont}{El}
\texto{102.7890}{404.4380}{11.9552}{\normalfont}{informe}
\texto{144.2592}{404.4380}{11.9552}{\normalfont}{está}
\texto{167.2586}{404.4380}{11.9552}{\normalfont}{organizado}
\texto{225.3130}{404.4380}{11.9552}{\normalfont}{en}
\texto{239.7932}{404.4380}{11.9552}{\normalfont}{cuatro}
\texto{275.0849}{404.4380}{11.9552}{\normalfont}{capítulos.}
\texto{326.7123}{404.4380}{11.9552}{\normalfont}{El}
\texto{340.7510}{404.4380}{11.9552}{\normalfont\bfseries}{Capítulo}
\texto{394.5542}{404.4380}{11.9552}{\normalfont\bfseries}{I}
\texto{402.3470}{404.4380}{11.9552}{\normalfont}{formula}
\texto{444.1328}{404.4380}{11.9552}{\normalfont}{el}
\texto{455.3731}{404.4380}{11.9552}{\normalfont}{proyecto:}
\texto{504.6524}{404.4380}{11.9552}{\normalfont}{presen-}
\texto{70.8660}{389.9920}{11.9552}{\normalfont}{ta}
\texto{84.1769}{389.9920}{11.9552}{\normalfont}{el}
\texto{95.5487}{389.9920}{11.9552}{\normalfont}{contexto}
\texto{142.0305}{389.9920}{11.9552}{\normalfont}{donde}
\texto{175.5146}{389.9920}{11.9552}{\normalfont}{surge}
\texto{205.1612}{389.9920}{11.9552}{\normalfont}{la}
\texto{217.1713}{389.9920}{11.9552}{\normalfont}{propuesta,}
\texto{273.4827}{389.9920}{11.9552}{\normalfont}{el}
\texto{284.8545}{389.9920}{11.9552}{\normalfont}{diagnóstico}
\texto{345.0561}{389.9920}{11.9552}{\normalfont}{de}
\texto{359.6797}{389.9920}{11.9552}{\normalfont}{la}
\texto{371.7019}{389.9920}{11.9552}{\normalfont}{situación}
\texto{420.1969}{389.9920}{11.9552}{\normalfont}{actual,}
\texto{457.5832}{389.9920}{11.9552}{\normalfont}{el}
\texto{468.9550}{389.9920}{11.9552}{\normalfont}{planteamiento}
\texto{70.8660}{375.5460}{11.9552}{\normalfont}{del}
\texto{89.3511}{375.5460}{11.9552}{\normalfont}{problema,}
\texto{143.6062}{375.5460}{11.9552}{\normalfont}{la}
\texto{156.2500}{375.5460}{11.9552}{\normalfont}{justificación,}
\texto{223.9021}{375.5460}{11.9552}{\normalfont}{los}
\texto{241.1511}{375.5460}{11.9552}{\normalfont}{objetivos,}
\texto{293.5316}{375.5460}{11.9552}{\normalfont}{el}
\texto{305.5131}{375.5460}{11.9552}{\normalfont}{alcance}
\texto{346.1106}{375.5460}{11.9552}{\normalfont}{y}
\texto{355.8278}{375.5460}{11.9552}{\normalfont}{la}
\texto{368.4596}{375.5460}{11.9552}{\normalfont}{metodología}
\texto{434.0937}{375.5460}{11.9552}{\normalfont}{de}
\texto{449.3389}{375.5460}{11.9552}{\normalfont}{desarrollo.}
\texto{505.6097}{375.5460}{11.9552}{\normalfont}{El}
\texto{520.4000}{375.5460}{11.9552}{\normalfont\bfseries}{Ca-}
\texto{70.8660}{361.1010}{11.9552}{\normalfont\bfseries}{pítulo}
\texto{110.2823}{361.1010}{11.9552}{\normalfont\bfseries}{II}
\texto{125.0860}{361.1010}{11.9552}{\normalfont}{documenta}
\texto{185.4167}{361.1010}{11.9552}{\normalfont}{la}
\texto{198.9213}{361.1010}{11.9552}{\normalfont}{investigación}
\texto{268.7397}{361.1010}{11.9552}{\normalfont}{bibliográfica}
\texto{335.5860}{361.1010}{11.9552}{\normalfont}{que}
\texto{357.8705}{361.1010}{11.9552}{\normalfont}{fundamenta}
\texto{422.7561}{361.1010}{11.9552}{\normalfont}{el}
\texto{435.6104}{361.1010}{11.9552}{\normalfont}{proyecto,}
\texto{486.5037}{361.1010}{11.9552}{\normalfont}{incluyendo}
\texto{70.8660}{346.6550}{11.9552}{\normalfont}{las}
\texto{87.9954}{346.6550}{11.9552}{\normalfont}{tendencias}
\texto{144.7970}{346.6550}{11.9552}{\normalfont}{tecnológicas,}
\texto{212.6547}{346.6550}{11.9552}{\normalfont}{el}
\texto{224.5166}{346.6550}{11.9552}{\normalfont}{estado}
\texto{260.5066}{346.6550}{11.9552}{\normalfont}{del}
\texto{278.8721}{346.6550}{11.9552}{\normalfont}{arte}
\texto{302.4406}{346.6550}{11.9552}{\normalfont}{y}
\texto{312.0263}{346.6550}{11.9552}{\normalfont}{el}
\texto{323.8882}{346.6550}{11.9552}{\normalfont}{análisis}
\texto{364.4953}{346.6550}{11.9552}{\normalfont}{de}
\texto{379.6090}{346.6550}{11.9552}{\normalfont}{las}
\texto{396.7385}{346.6550}{11.9552}{\normalfont}{tecnologías}
\texto{456.1414}{346.6550}{11.9552}{\normalfont}{seleccionadas.}
\texto{529.9289}{346.6550}{11.9552}{\normalfont}{El}
\texto{70.8660}{332.2090}{11.9552}{\normalfont\bfseries}{Capítulo}
\texto{126.1875}{332.2090}{11.9552}{\normalfont\bfseries}{III}
\texto{146.0650}{332.2090}{11.9552}{\normalfont}{describe}
\texto{191.5139}{332.2090}{11.9552}{\normalfont}{el}
\texto{204.0692}{332.2090}{11.9552}{\normalfont}{emprendimiento}
\texto{290.4312}{332.2090}{11.9552}{\normalfont}{digital:}
\texto{330.3018}{332.2090}{11.9552}{\normalfont}{la}
\texto{343.5075}{332.2090}{11.9552}{\normalfont}{idea}
\texto{368.4078}{332.2090}{11.9552}{\normalfont}{de}
\texto{384.2150}{332.2090}{11.9552}{\normalfont}{negocio,}
\texto{429.6113}{332.2090}{11.9552}{\normalfont}{la}
\texto{442.8050}{332.2090}{11.9552}{\normalfont}{propuesta}
\texto{497.0482}{332.2090}{11.9552}{\normalfont}{de}
\texto{512.8434}{332.2090}{11.9552}{\normalfont}{valor,}
\texto{70.8660}{317.7630}{11.9552}{\normalfont}{el}
\texto{84.4615}{317.7630}{11.9552}{\normalfont}{público}
\texto{126.6729}{317.7630}{11.9552}{\normalfont}{objetivo,}
\texto{176.0383}{317.7630}{11.9552}{\normalfont}{el}
\texto{189.6457}{317.7630}{11.9552}{\normalfont}{modelo}
\texto{231.5295}{317.7630}{11.9552}{\normalfont}{de}
\texto{248.3768}{317.7630}{11.9552}{\normalfont}{funcionamiento,}
\texto{334.8153}{317.7630}{11.9552}{\normalfont}{la}
\texto{349.0731}{317.7630}{11.9552}{\normalfont}{identidad}
\texto{402.0179}{317.7630}{11.9552}{\normalfont}{visual,}
\texto{440.0665}{317.7630}{11.9552}{\normalfont}{el}
\texto{453.6619}{317.7630}{11.9552}{\normalfont}{análisis}
\texto{496.0025}{317.7630}{11.9552}{\normalfont}{FODA}
\texto{534.9573}{317.7630}{11.9552}{\normalfont}{y}
\texto{70.8660}{303.3170}{11.9552}{\normalfont}{el}
\texto{83.7800}{303.3170}{11.9552}{\normalfont}{Lean}
\texto{113.1133}{303.3170}{11.9552}{\normalfont}{Canvas.}
\texto{157.3045}{303.3170}{11.9552}{\normalfont}{El}
\texto{173.0330}{303.3170}{11.9552}{\normalfont\bfseries}{Capítulo}
\texto{228.7729}{303.3170}{11.9552}{\normalfont\bfseries}{IV}
\texto{248.4430}{303.3170}{11.9552}{\normalfont}{modela}
\texto{289.6454}{303.3170}{11.9552}{\normalfont}{y}
\texto{300.2831}{303.3170}{11.9552}{\normalfont}{automatiza}
\texto{361.9744}{303.3170}{11.9552}{\normalfont}{el}
\texto{374.8884}{303.3170}{11.9552}{\normalfont}{proceso}
\texto{417.4561}{303.3170}{11.9552}{\normalfont}{de}
\texto{433.6219}{303.3170}{11.9552}{\normalfont}{promoción,}
\texto{494.9855}{303.3170}{11.9552}{\normalfont}{reserva}
\texto{534.9589}{303.3170}{11.9552}{\normalfont}{y}
\texto{70.8660}{288.8710}{11.9552}{\normalfont}{gestión}
\texto{111.0522}{288.8710}{11.9552}{\normalfont}{de}
\texto{127.1224}{288.8710}{11.9552}{\normalfont}{tours}
\texto{157.5651}{288.8710}{11.9552}{\normalfont}{mediante}
\texto{208.4321}{288.8710}{11.9552}{\normalfont}{la}
\texto{221.9008}{288.8710}{11.9552}{\normalfont}{notación}
\texto{269.8388}{288.8710}{11.9552}{\normalfont}{BPMN,}
\texto{313.2099}{288.8710}{11.9552}{\normalfont}{tanto}
\texto{344.5660}{288.8710}{11.9552}{\normalfont}{en}
\texto{360.6361}{288.8710}{11.9552}{\normalfont}{su}
\texto{376.1205}{288.8710}{11.9552}{\normalfont}{estado}
\texto{413.0549}{288.8710}{11.9552}{\normalfont}{actual}
\texto{448.6360}{288.8710}{11.9552}{\normalfont}{como}
\texto{479.6645}{288.8710}{11.9552}{\normalfont}{en}
\texto{495.7347}{288.8710}{11.9552}{\normalfont}{el}
\texto{508.5530}{288.8710}{11.9552}{\normalfont}{estado}
\texto{70.8660}{274.4260}{11.9552}{\normalfont}{propuesto.}
\texto{129.5445}{274.4260}{11.9552}{\normalfont}{La}
\texto{148.0081}{274.4260}{11.9552}{\normalfont}{lectura}
\texto{188.4119}{274.4260}{11.9552}{\normalfont}{del}
\texto{208.6545}{274.4260}{11.9552}{\normalfont}{informe}
\texto{252.6353}{274.4260}{11.9552}{\normalfont}{permite}
\texto{297.2760}{274.4260}{11.9552}{\normalfont}{al}
\texto{311.6652}{274.4260}{11.9552}{\normalfont}{lector}
\texto{345.5654}{274.4260}{11.9552}{\normalfont}{comprender}
\texto{410.6830}{274.4260}{11.9552}{\normalfont}{qué}
\texto{433.8641}{274.4260}{11.9552}{\normalfont}{se}
\texto{448.9683}{274.4260}{11.9552}{\normalfont}{propone}
\texto{495.5482}{274.4260}{11.9552}{\normalfont}{construir}
\texto{70.8660}{259.9800}{11.9552}{\normalfont}{(las}
\texto{93.9826}{259.9800}{11.9552}{\normalfont}{aplicaciones}
\texto{159.3847}{259.9800}{11.9552}{\normalfont}{de}
\texto{175.9450}{259.9800}{11.9552}{\normalfont}{escritorio}
\texto{227.6895}{259.9800}{11.9552}{\normalfont}{del}
\texto{247.5017}{259.9800}{11.9552}{\normalfont}{ecosistema}
\texto{306.4647}{259.9800}{11.9552}{\normalfont}{Explora-Nicaragua),}
\texto{414.8817}{259.9800}{11.9552}{\normalfont}{por}
\texto{436.9677}{259.9800}{11.9552}{\normalfont}{qué}
\texto{459.7065}{259.9800}{11.9552}{\normalfont}{surge}
\texto{491.2897}{259.9800}{11.9552}{\normalfont}{(la}
\texto{509.7892}{259.9800}{11.9552}{\normalfont}{escasa}
\texto{70.8660}{245.5340}{11.9552}{\normalfont}{digitalización}
\texto{142.1262}{245.5340}{11.9552}{\normalfont}{de}
\texto{157.4671}{245.5340}{11.9552}{\normalfont}{la}
\texto{170.1946}{245.5340}{11.9552}{\normalfont}{oferta}
\texto{203.4205}{245.5340}{11.9552}{\normalfont}{turística}
\texto{249.3810}{245.5340}{11.9552}{\normalfont}{y}
\texto{259.1939}{245.5340}{11.9552}{\normalfont}{de}
\texto{274.5348}{245.5340}{11.9552}{\normalfont}{la}
\texto{287.2623}{245.5340}{11.9552}{\normalfont}{gestión}
\texto{326.7312}{245.5340}{11.9552}{\normalfont}{de}
\texto{342.0601}{245.5340}{11.9552}{\normalfont}{los}
\texto{359.4167}{245.5340}{11.9552}{\normalfont}{operadores),}
\texto{425.8733}{245.5340}{11.9552}{\normalfont}{a}
\texto{435.3490}{245.5340}{11.9552}{\normalfont}{quién}
\texto{466.6238}{245.5340}{11.9552}{\normalfont}{beneficia}
\texto{514.8056}{245.5340}{11.9552}{\normalfont}{(ope-}
\texto{70.8660}{231.0880}{11.9552}{\normalfont}{radores,}
\texto{114.9831}{231.0880}{11.9552}{\normalfont}{guías,}
\texto{148.0440}{231.0880}{11.9552}{\normalfont}{cooperativas}
\texto{215.2609}{231.0880}{11.9552}{\normalfont}{comunitarias}
\texto{284.4193}{231.0880}{11.9552}{\normalfont}{y}
\texto{294.3278}{231.0880}{11.9552}{\normalfont}{el}
\texto{306.5245}{231.0880}{11.9552}{\normalfont}{equipo}
\texto{344.0710}{231.0880}{11.9552}{\normalfont}{gestor)}
\texto{382.9971}{231.0880}{11.9552}{\normalfont}{y}
\texto{392.9056}{231.0880}{11.9552}{\normalfont}{cómo}
\texto{423.3005}{231.0880}{11.9552}{\normalfont}{se}
\texto{436.8625}{231.0880}{11.9552}{\normalfont}{desarrolla}
\texto{490.0847}{231.0880}{11.9552}{\normalfont}{(mediante}
\texto{70.8660}{216.6420}{11.9552}{\normalfont}{una}
\texto{93.6239}{216.6420}{11.9552}{\normalfont}{metodología}
\texto{159.6286}{216.6420}{11.9552}{\normalfont}{ágil}
\texto{181.7481}{216.6420}{11.9552}{\normalfont}{basada}
\texto{220.8296}{216.6420}{11.9552}{\normalfont}{en}
\texto{236.4455}{216.6420}{11.9552}{\normalfont}{Scrum).}
\end{pagina}
% Página original 13
\begin{pagina}{13}
\texto{70.8660}{746.3890}{11.9552}{\normalfont}{1.2.}
\texto{96.8710}{746.3890}{11.9552}{\normalfont}{CONTEXTO}
\texto{169.8694}{746.3890}{11.9552}{\normalfont}{DEL}
\texto{197.9856}{746.3890}{11.9552}{\normalfont}{PROYECTO}
\texto{535.2810}{746.3890}{11.9552}{\normalfont}{3}
\texto{70.8660}{705.1940}{17.2154}{\normalfont\bfseries}{1.2.}
\texto{120.3603}{705.1940}{17.2154}{\normalfont\bfseries}{Contexto}
\texto{204.5023}{705.1940}{17.2154}{\normalfont\bfseries}{del}
\texto{235.9341}{705.1940}{17.2154}{\normalfont\bfseries}{proyecto}
\texto{88.7990}{678.9090}{11.9552}{\normalfont}{El}
\texto{105.1561}{678.9090}{11.9552}{\normalfont}{turismo}
\texto{149.3832}{678.9090}{11.9552}{\normalfont}{nicaragüense}
\texto{219.6367}{678.9090}{11.9552}{\normalfont}{genera}
\texto{257.9460}{678.9090}{11.9552}{\normalfont}{aproximadamente}
\texto{353.8171}{678.9090}{11.9552}{\normalfont}{el}
\texto{367.4245}{678.9090}{11.9552}{\normalfont}{5.2}
\texto{384.3674}{678.9090}{11.9552}{\normalfont}{\%}
\texto{399.2755}{678.9090}{11.9552}{\normalfont}{del}
\texto{419.3746}{678.9090}{11.9552}{\normalfont}{PIB}
\texto{445.0042}{678.9090}{11.9552}{\normalfont}{nacional}
\texto{492.4185}{678.9090}{11.9552}{\normalfont}{e}
\texto{502.7621}{678.9090}{11.9552}{\normalfont}{indirec-}
\texto{70.8660}{664.4630}{11.9552}{\normalfont}{tamente}
\texto{116.1834}{664.4630}{11.9552}{\normalfont}{mantiene}
\texto{166.7037}{664.4630}{11.9552}{\normalfont}{empleada}
\texto{218.8475}{664.4630}{11.9552}{\normalfont}{a}
\texto{228.7296}{664.4630}{11.9552}{\normalfont}{una}
\texto{251.6071}{664.4630}{11.9552}{\normalfont}{población}
\texto{304.7360}{664.4630}{11.9552}{\normalfont}{significativa}
\texto{368.9809}{664.4630}{11.9552}{\normalfont}{en}
\texto{384.7043}{664.4630}{11.9552}{\normalfont}{sectores}
\texto{428.5345}{664.4630}{11.9552}{\normalfont}{como}
\texto{459.2163}{664.4630}{11.9552}{\normalfont}{el}
\texto{471.6880}{664.4630}{11.9552}{\normalfont}{transporte,}
\texto{532.0354}{664.4630}{11.9552}{\normalfont}{la}
\texto{70.8660}{650.0170}{11.9552}{\normalfont}{gastronomía}
\texto{136.4977}{650.0170}{11.9552}{\normalfont}{y}
\texto{145.8084}{650.0170}{11.9552}{\normalfont}{el}
\texto{157.3953}{650.0170}{11.9552}{\normalfont}{alojamiento}
\texto{220.0478}{650.0170}{11.9552}{\normalfont}{[2].}
\texto{238.7879}{650.0170}{11.9552}{\normalfont}{A}
\texto{250.6953}{650.0170}{11.9552}{\normalfont}{pesar}
\texto{280.8917}{650.0170}{11.9552}{\normalfont}{de}
\texto{295.7305}{650.0170}{11.9552}{\normalfont}{esta}
\texto{319.1005}{650.0170}{11.9552}{\normalfont}{importancia,}
\texto{386.9415}{650.0170}{11.9552}{\normalfont}{la}
\texto{399.1908}{650.0170}{11.9552}{\normalfont}{oferta}
\texto{431.9146}{650.0170}{11.9552}{\normalfont}{turística}
\texto{477.3850}{650.0170}{11.9552}{\normalfont}{se}
\texto{490.3492}{650.0170}{11.9552}{\normalfont}{distribuye}
\texto{70.8660}{635.5710}{11.9552}{\normalfont}{en}
\texto{85.7885}{635.5710}{11.9552}{\normalfont}{gran}
\texto{111.7791}{635.5710}{11.9552}{\normalfont}{parte}
\texto{141.6719}{635.5710}{11.9552}{\normalfont}{a}
\texto{150.7411}{635.5710}{11.9552}{\normalfont}{través}
\texto{184.2802}{635.5710}{11.9552}{\normalfont}{de}
\texto{199.2027}{635.5710}{11.9552}{\normalfont}{canales}
\texto{238.9155}{635.5710}{11.9552}{\normalfont}{tradicionales:}
\texto{309.8457}{635.5710}{11.9552}{\normalfont}{agencias}
\texto{355.4117}{635.5710}{11.9552}{\normalfont}{de}
\texto{370.3342}{635.5710}{11.9552}{\normalfont}{viajes}
\texto{402.5654}{635.5710}{11.9552}{\normalfont}{físicas,}
\texto{439.4161}{635.5710}{11.9552}{\normalfont}{propaganda}
\texto{502.4774}{635.5710}{11.9552}{\normalfont}{radial}
\texto{534.9597}{635.5710}{11.9552}{\normalfont}{y}
\texto{70.8660}{621.1260}{11.9552}{\normalfont}{redes}
\texto{100.1730}{621.1260}{11.9552}{\normalfont}{sociales}
\texto{141.5858}{621.1260}{11.9552}{\normalfont}{sin}
\texto{159.1982}{621.1260}{11.9552}{\normalfont}{funcionalidad}
\texto{231.0513}{621.1260}{11.9552}{\normalfont}{transaccional.}
\texto{304.5949}{621.1260}{11.9552}{\normalfont}{Esta}
\texto{330.8223}{621.1260}{11.9552}{\normalfont}{falta}
\texto{357.1500}{621.1260}{11.9552}{\normalfont}{de}
\texto{372.0964}{621.1260}{11.9552}{\normalfont}{digitalización}
\texto{442.9740}{621.1260}{11.9552}{\normalfont}{dificulta}
\texto{487.8371}{621.1260}{11.9552}{\normalfont}{la}
\texto{500.1820}{621.1260}{11.9552}{\normalfont}{toma}
\texto{529.4364}{621.1260}{11.9552}{\normalfont}{de}
\texto{70.8660}{606.6800}{11.9552}{\normalfont}{decisiones}
\texto{124.4564}{606.6800}{11.9552}{\normalfont}{informada}
\texto{179.5436}{606.6800}{11.9552}{\normalfont}{de}
\texto{194.6334}{606.6800}{11.9552}{\normalfont}{los}
\texto{211.7509}{606.6800}{11.9552}{\normalfont}{turistas,}
\texto{256.8865}{606.6800}{11.9552}{\normalfont}{limita}
\texto{290.1985}{606.6800}{11.9552}{\normalfont}{la}
\texto{302.6869}{606.6800}{11.9552}{\normalfont}{visibilidad}
\texto{358.4890}{606.6800}{11.9552}{\normalfont}{de}
\texto{373.5788}{606.6800}{11.9552}{\normalfont}{los}
\texto{390.6963}{606.6800}{11.9552}{\normalfont}{operadores}
\texto{449.0974}{606.6800}{11.9552}{\normalfont}{locales}
\texto{486.0486}{606.6800}{11.9552}{\normalfont}{y}
\texto{495.6103}{606.6800}{11.9552}{\normalfont}{perpetúa}
\texto{70.8660}{592.2340}{11.9552}{\normalfont}{la}
\texto{83.8685}{592.2340}{11.9552}{\normalfont}{dependencia}
\texto{150.5235}{592.2340}{11.9552}{\normalfont}{de}
\texto{166.1394}{592.2340}{11.9552}{\normalfont}{plataformas}
\texto{230.2599}{592.2340}{11.9552}{\normalfont}{extranjeras}
\texto{290.8154}{592.2340}{11.9552}{\normalfont}{que}
\texto{312.5978}{592.2340}{11.9552}{\normalfont}{retienen}
\texto{357.4680}{592.2340}{11.9552}{\normalfont}{importantes}
\texto{422.5760}{592.2340}{11.9552}{\normalfont}{márgenes}
\texto{474.0145}{592.2340}{11.9552}{\normalfont}{de}
\texto{489.6304}{592.2340}{11.9552}{\normalfont}{comisión.}
\texto{88.7990}{577.7880}{11.9552}{\normalfont}{El}
\texto{103.9247}{577.7880}{11.9552}{\normalfont}{ecosistema}
\texto{161.9433}{577.7880}{11.9552}{\normalfont}{se}
\texto{175.6727}{577.7880}{11.9552}{\normalfont}{desarrolla}
\texto{229.0861}{577.7880}{11.9552}{\normalfont}{en}
\texto{244.7020}{577.7880}{11.9552}{\normalfont}{el}
\texto{257.0660}{577.7880}{11.9552}{\normalfont}{contexto}
\texto{304.5521}{577.7880}{11.9552}{\normalfont}{de}
\texto{320.1680}{577.7880}{11.9552}{\normalfont}{un}
\texto{337.0965}{577.7880}{11.9552}{\normalfont}{mercado}
\texto{383.9298}{577.7880}{11.9552}{\normalfont}{turístico}
\texto{430.1773}{577.7880}{11.9552}{\normalfont}{nicaragüense}
\texto{499.1875}{577.7880}{11.9552}{\normalfont}{formado}
\texto{70.8660}{563.3420}{11.9552}{\normalfont}{por}
\texto{91.2783}{563.3420}{11.9552}{\normalfont}{tres}
\texto{113.3835}{563.3420}{11.9552}{\normalfont}{tipos}
\texto{141.6766}{563.3420}{11.9552}{\normalfont}{de}
\texto{156.5633}{563.3420}{11.9552}{\normalfont}{actores}
\texto{195.5779}{563.3420}{11.9552}{\normalfont}{principales.}
\texto{256.7168}{563.3420}{11.9552}{\normalfont}{En}
\texto{274.3650}{563.3420}{11.9552}{\normalfont}{primer}
\texto{311.3640}{563.3420}{11.9552}{\normalfont}{lugar,}
\texto{343.8223}{563.3420}{11.9552}{\normalfont}{los}
\texto{360.7230}{563.3420}{11.9552}{\normalfont\bfseries}{turistas}
\texto{408.2170}{563.3420}{11.9552}{\normalfont}{(locales}
\texto{449.5055}{563.3420}{11.9552}{\normalfont}{y}
\texto{458.8640}{563.3420}{11.9552}{\normalfont}{extranjeros)}
\texto{523.2428}{563.3420}{11.9552}{\normalfont}{que}
\texto{70.8660}{548.8960}{11.9552}{\normalfont}{demandan}
\texto{128.5905}{548.8960}{11.9552}{\normalfont}{experiencias}
\texto{195.4822}{548.8960}{11.9552}{\normalfont}{y}
\texto{206.7058}{548.8960}{11.9552}{\normalfont}{buscan}
\texto{246.9350}{548.8960}{11.9552}{\normalfont}{información}
\texto{312.1506}{548.8960}{11.9552}{\normalfont}{confiable,}
\texto{365.3226}{548.8960}{11.9552}{\normalfont}{precios}
\texto{405.5518}{548.8960}{11.9552}{\normalfont}{comparables}
\texto{473.7467}{548.8960}{11.9552}{\normalfont}{y}
\texto{484.9702}{548.8960}{11.9552}{\normalfont}{un}
\texto{503.0345}{548.8960}{11.9552}{\normalfont}{proceso}
\texto{70.8660}{534.4510}{11.9552}{\normalfont}{de}
\texto{87.7013}{534.4510}{11.9552}{\normalfont}{reserva}
\texto{128.3323}{534.4510}{11.9552}{\normalfont}{seguro.}
\texto{169.3075}{534.4510}{11.9552}{\normalfont}{En}
\texto{188.9045}{534.4510}{11.9552}{\normalfont}{segundo}
\texto{235.0707}{534.4510}{11.9552}{\normalfont}{lugar,}
\texto{269.4777}{534.4510}{11.9552}{\normalfont}{las}
\texto{288.3280}{534.4510}{11.9552}{\normalfont\bfseries}{agencias}
\texto{342.6512}{534.4510}{11.9552}{\normalfont\bfseries}{de}
\texto{362.1538}{534.4510}{11.9552}{\normalfont\bfseries}{viajes,}
\texto{405.0801}{534.4510}{11.9552}{\normalfont\bfseries}{guías}
\texto{440.7508}{534.4510}{11.9552}{\normalfont\bfseries}{independientes}
\texto{534.0313}{534.4510}{11.9552}{\normalfont\bfseries}{y}
\texto{70.8660}{520.0050}{11.9552}{\normalfont\bfseries}{cooperativas}
\texto{148.6489}{520.0050}{11.9552}{\normalfont\bfseries}{comunitarias}
\texto{227.8980}{520.0050}{11.9552}{\normalfont}{que}
\texto{249.9673}{520.0050}{11.9552}{\normalfont}{ofrecen}
\texto{290.2468}{520.0050}{11.9552}{\normalfont}{experiencias}
\texto{356.2897}{520.0050}{11.9552}{\normalfont}{pero}
\texto{382.9091}{520.0050}{11.9552}{\normalfont}{que,}
\texto{408.2302}{520.0050}{11.9552}{\normalfont}{en}
\texto{424.1330}{520.0050}{11.9552}{\normalfont}{su}
\texto{439.4381}{520.0050}{11.9552}{\normalfont}{mayoría,}
\texto{487.5267}{520.0050}{11.9552}{\normalfont}{carecen}
\texto{529.4321}{520.0050}{11.9552}{\normalfont}{de}
\texto{70.8660}{505.5590}{11.9552}{\normalfont}{herramientas}
\texto{140.6150}{505.5590}{11.9552}{\normalfont}{digitales}
\texto{186.6354}{505.5590}{11.9552}{\normalfont}{para}
\texto{213.0683}{505.5590}{11.9552}{\normalfont}{gestionar}
\texto{262.9908}{505.5590}{11.9552}{\normalfont}{su}
\texto{277.7818}{505.5590}{11.9552}{\normalfont}{catálogo,}
\texto{326.9894}{505.5590}{11.9552}{\normalfont}{su}
\texto{341.7924}{505.5590}{11.9552}{\normalfont}{disponibilidad}
\texto{417.4018}{505.5590}{11.9552}{\normalfont}{y}
\texto{427.2505}{505.5590}{11.9552}{\normalfont}{sus}
\texto{446.6705}{505.5590}{11.9552}{\normalfont}{reservas.}
\texto{493.7238}{505.5590}{11.9552}{\normalfont}{En}
\texto{511.8742}{505.5590}{11.9552}{\normalfont}{tercer}
\texto{70.8660}{491.1130}{11.9552}{\normalfont}{lugar,}
\texto{103.2646}{491.1130}{11.9552}{\normalfont}{el}
\texto{114.9060}{491.1130}{11.9552}{\normalfont\bfseries}{equipo}
\texto{157.5155}{491.1130}{11.9552}{\normalfont\bfseries}{gestor}
\texto{196.7417}{491.1130}{11.9552}{\normalfont\bfseries}{de}
\texto{213.9489}{491.1130}{11.9552}{\normalfont\bfseries}{la}
\texto{227.8336}{491.1130}{11.9552}{\normalfont\bfseries}{plataforma}
\texto{291.4240}{491.1130}{11.9552}{\normalfont}{,}
\texto{297.8081}{491.1130}{11.9552}{\normalfont}{encargado}
\texto{352.3190}{491.1130}{11.9552}{\normalfont}{de}
\texto{367.1698}{491.1130}{11.9552}{\normalfont}{verificar}
\texto{411.2773}{491.1130}{11.9552}{\normalfont}{a}
\texto{420.2628}{491.1130}{11.9552}{\normalfont}{los}
\texto{437.1292}{491.1130}{11.9552}{\normalfont}{operadores,}
\texto{498.5431}{491.1130}{11.9552}{\normalfont}{moderar}
\texto{70.8660}{476.6670}{11.9552}{\normalfont}{la}
\texto{83.8685}{476.6670}{11.9552}{\normalfont}{oferta}
\texto{117.3574}{476.6670}{11.9552}{\normalfont}{y}
\texto{127.4452}{476.6670}{11.9552}{\normalfont}{supervisar}
\texto{183.4983}{476.6670}{11.9552}{\normalfont}{la}
\texto{196.5128}{476.6670}{11.9552}{\normalfont}{actividad}
\texto{247.5615}{476.6670}{11.9552}{\normalfont}{económica}
\texto{304.1383}{476.6670}{11.9552}{\normalfont}{de}
\texto{319.7541}{476.6670}{11.9552}{\normalfont}{la}
\texto{332.7566}{476.6670}{11.9552}{\normalfont}{plataforma.}
\texto{88.7990}{462.2210}{11.9552}{\normalfont\bfseries}{Explora-Nicaragua}
\texto{203.2950}{462.2210}{11.9552}{\normalfont}{nace}
\texto{230.3735}{462.2210}{11.9552}{\normalfont}{como}
\texto{261.3662}{462.2210}{11.9552}{\normalfont}{respuesta}
\texto{313.2996}{462.2210}{11.9552}{\normalfont}{a}
\texto{323.4806}{462.2210}{11.9552}{\normalfont}{esta}
\texto{348.0342}{462.2210}{11.9552}{\normalfont}{problemática,}
\texto{421.9508}{462.2210}{11.9552}{\normalfont}{atendiendo}
\texto{482.2098}{462.2210}{11.9552}{\normalfont}{simultánea-}
\texto{70.8660}{447.7760}{11.9552}{\normalfont}{mente}
\texto{105.9808}{447.7760}{11.9552}{\normalfont}{la}
\texto{119.3061}{447.7760}{11.9552}{\normalfont}{experiencia}
\texto{180.7678}{447.7760}{11.9552}{\normalfont}{del}
\texto{199.9463}{447.7760}{11.9552}{\normalfont}{turista}
\texto{238.0619}{447.7760}{11.9552}{\normalfont}{y}
\texto{248.4605}{447.7760}{11.9552}{\normalfont}{la}
\texto{261.7977}{447.7760}{11.9552}{\normalfont}{eficiencia}
\texto{312.1937}{447.7760}{11.9552}{\normalfont}{operativa}
\texto{363.8927}{447.7760}{11.9552}{\normalfont}{de}
\texto{379.8314}{447.7760}{11.9552}{\normalfont}{quienes}
\texto{421.5120}{447.7760}{11.9552}{\normalfont}{ofrecen}
\texto{461.8393}{447.7760}{11.9552}{\normalfont}{los}
\texto{479.7817}{447.7760}{11.9552}{\normalfont}{tours.}
\texto{513.3447}{447.7760}{11.9552}{\normalfont}{En}
\texto{532.0331}{447.7760}{11.9552}{\normalfont}{la}
\texto{70.8660}{433.3300}{11.9552}{\normalfont}{visión}
\texto{105.0411}{433.3300}{11.9552}{\normalfont}{completa}
\texto{155.7479}{433.3300}{11.9552}{\normalfont}{del}
\texto{175.2373}{433.3300}{11.9552}{\normalfont}{ecosistema,}
\texto{237.1294}{433.3300}{11.9552}{\normalfont}{la}
\texto{250.7535}{433.3300}{11.9552}{\normalfont}{app}
\texto{274.1450}{433.3300}{11.9552}{\normalfont}{móvil}
\texto{306.6441}{433.3300}{11.9552}{\normalfont}{permite}
\texto{350.5077}{433.3300}{11.9552}{\normalfont}{al}
\texto{364.1438}{433.3300}{11.9552}{\normalfont}{turista}
\texto{402.5582}{433.3300}{11.9552}{\normalfont}{explorar}
\texto{449.0377}{433.3300}{11.9552}{\normalfont}{destinos}
\texto{494.6587}{433.3300}{11.9552}{\normalfont}{por}
\texto{516.4220}{433.3300}{11.9552}{\normalfont}{cate-}
\texto{70.8660}{418.8840}{11.9552}{\normalfont}{goría,}
\texto{103.5587}{418.8840}{11.9552}{\normalfont}{leer}
\texto{125.8456}{418.8840}{11.9552}{\normalfont}{reseñas}
\texto{166.4717}{418.8840}{11.9552}{\normalfont}{verificadas,}
\texto{227.1970}{418.8840}{11.9552}{\normalfont}{comparar}
\texto{279.4125}{418.8840}{11.9552}{\normalfont}{opciones}
\texto{326.8005}{418.8840}{11.9552}{\normalfont}{y}
\texto{337.0677}{418.8840}{11.9552}{\normalfont}{confirmar}
\texto{389.9216}{418.8840}{11.9552}{\normalfont}{reservas}
\texto{434.1295}{418.8840}{11.9552}{\normalfont}{mediante}
\texto{484.7096}{418.8840}{11.9552}{\normalfont}{un}
\texto{501.7936}{418.8840}{11.9552}{\normalfont}{voucher}
\texto{70.8660}{404.4380}{11.9552}{\normalfont}{digital;}
\texto{110.5214}{404.4380}{11.9552}{\normalfont}{el}
\texto{122.9270}{404.4380}{11.9552}{\normalfont\bfseries}{backoffice}
\texto{183.7790}{404.4380}{11.9552}{\normalfont\bfseries}{desktop}
\texto{233.1240}{404.4380}{11.9552}{\normalfont}{de}
\texto{248.7160}{404.4380}{11.9552}{\normalfont}{operadores}
\texto{307.6312}{404.4380}{11.9552}{\normalfont}{permite}
\texto{350.8851}{404.4380}{11.9552}{\normalfont}{a}
\texto{360.6238}{404.4380}{11.9552}{\normalfont}{las}
\texto{378.2434}{404.4380}{11.9552}{\normalfont}{agencias}
\texto{424.4789}{404.4380}{11.9552}{\normalfont}{registrar}
\texto{471.3648}{404.4380}{11.9552}{\normalfont}{sus}
\texto{491.0001}{404.4380}{11.9552}{\normalfont}{tours,}
\texto{524.2283}{404.4380}{11.9552}{\normalfont}{pu-}
\texto{70.8660}{389.9920}{11.9552}{\normalfont}{blicar}
\texto{103.1881}{389.9920}{11.9552}{\normalfont}{descripciones}
\texto{173.6951}{389.9920}{11.9552}{\normalfont}{con}
\texto{194.9729}{389.9920}{11.9552}{\normalfont}{fotografías,}
\texto{255.3252}{389.9920}{11.9552}{\normalfont}{establecer}
\texto{309.1857}{389.9920}{11.9552}{\normalfont}{precios}
\texto{348.0760}{389.9920}{11.9552}{\normalfont}{y}
\texto{357.9725}{389.9920}{11.9552}{\normalfont}{disponibilidad}
\texto{433.6179}{389.9920}{11.9552}{\normalfont}{y}
\texto{443.5144}{389.9920}{11.9552}{\normalfont}{recibir}
\texto{479.7386}{389.9920}{11.9552}{\normalfont}{reservas}
\texto{523.5759}{389.9920}{11.9552}{\normalfont}{con}
\texto{70.8660}{375.5460}{11.9552}{\normalfont}{pagos}
\texto{103.8863}{375.5460}{11.9552}{\normalfont}{en}
\texto{119.9445}{375.5460}{11.9552}{\normalfont}{línea;}
\texto{151.6114}{375.5460}{11.9552}{\normalfont}{y}
\texto{162.1416}{375.5460}{11.9552}{\normalfont}{el}
\texto{174.9910}{375.5460}{11.9552}{\normalfont\bfseries}{panel}
\texto{211.3312}{375.5460}{11.9552}{\normalfont\bfseries}{de}
\texto{229.9371}{375.5460}{11.9552}{\normalfont\bfseries}{administración}
\texto{320.9700}{375.5460}{11.9552}{\normalfont}{permite}
\texto{364.6663}{375.5460}{11.9552}{\normalfont}{al}
\texto{378.1111}{375.5460}{11.9552}{\normalfont}{equipo}
\texto{416.2792}{375.5460}{11.9552}{\normalfont}{gestor}
\texto{451.2626}{375.5460}{11.9552}{\normalfont}{verificar}
\texto{496.5895}{375.5460}{11.9552}{\normalfont}{operado-}
\texto{70.8660}{361.1010}{11.9552}{\normalfont}{res,}
\texto{92.3041}{361.1010}{11.9552}{\normalfont}{moderar}
\texto{138.7261}{361.1010}{11.9552}{\normalfont}{la}
\texto{151.6569}{361.1010}{11.9552}{\normalfont}{oferta,}
\texto{188.3258}{361.1010}{11.9552}{\normalfont}{gestionar}
\texto{238.3918}{361.1010}{11.9552}{\normalfont}{usuarios}
\texto{283.9698}{361.1010}{11.9552}{\normalfont}{y}
\texto{293.9739}{361.1010}{11.9552}{\normalfont}{consultar}
\texto{344.6903}{361.1010}{11.9552}{\normalfont}{reportes}
\texto{389.8762}{361.1010}{11.9552}{\normalfont}{financieros.}
\texto{450.9920}{361.1010}{11.9552}{\normalfont\bfseries}{En}
\texto{471.6793}{361.1010}{11.9552}{\normalfont\bfseries}{el}
\texto{485.9514}{361.1010}{11.9552}{\normalfont\bfseries}{prototipo}
\texto{70.8660}{346.6550}{11.9552}{\normalfont\bfseries}{inicial}
\texto{110.4843}{346.6550}{11.9552}{\normalfont\bfseries}{de}
\texto{128.7913}{346.6550}{11.9552}{\normalfont\bfseries}{esta}
\texto{156.6888}{346.6550}{11.9552}{\normalfont\bfseries}{etapa,}
\texto{197.0388}{346.6550}{11.9552}{\normalfont\bfseries}{el}
\texto{211.5978}{346.6550}{11.9552}{\normalfont\bfseries}{alcance}
\texto{258.6726}{346.6550}{11.9552}{\normalfont\bfseries}{del}
\texto{280.7156}{346.6550}{11.9552}{\normalfont\bfseries}{desarrollo}
\texto{342.7631}{346.6550}{11.9552}{\normalfont\bfseries}{se}
\texto{358.9038}{346.6550}{11.9552}{\normalfont\bfseries}{limita}
\texto{397.7750}{346.6550}{11.9552}{\normalfont\bfseries}{a}
\texto{409.0117}{346.6550}{11.9552}{\normalfont\bfseries}{las}
\texto{429.2781}{346.6550}{11.9552}{\normalfont\bfseries}{dos}
\texto{453.4790}{346.6550}{11.9552}{\normalfont\bfseries}{aplicaciones}
\texto{527.5283}{346.6550}{11.9552}{\normalfont\bfseries}{de}
\texto{70.8660}{332.2090}{11.9552}{\normalfont\bfseries}{escritorio}
\texto{130.5570}{332.2090}{11.9552}{\normalfont}{(backoffice}
\texto{189.3335}{332.2090}{11.9552}{\normalfont}{de}
\texto{206.1808}{332.2090}{11.9552}{\normalfont}{operadores}
\texto{266.3274}{332.2090}{11.9552}{\normalfont}{y}
\texto{277.6466}{332.2090}{11.9552}{\normalfont}{panel}
\texto{310.0906}{332.2090}{11.9552}{\normalfont}{de}
\texto{326.9259}{332.2090}{11.9552}{\normalfont}{administración),}
\texto{415.3777}{332.2090}{11.9552}{\normalfont}{que}
\texto{438.3914}{332.2090}{11.9552}{\normalfont}{consumen}
\texto{493.6627}{332.2090}{11.9552}{\normalfont}{la}
\texto{507.9085}{332.2090}{11.9552}{\normalfont}{misma}
\texto{70.8660}{317.7630}{11.9552}{\normalfont}{API}
\texto{96.1560}{317.7630}{11.9552}{\normalfont}{REST}
\texto{132.0216}{317.7630}{11.9552}{\normalfont}{y}
\texto{142.5279}{317.7630}{11.9552}{\normalfont}{trabajan}
\texto{190.4395}{317.7630}{11.9552}{\normalfont}{sobre}
\texto{221.4967}{317.7630}{11.9552}{\normalfont}{una}
\texto{244.6970}{317.7630}{11.9552}{\normalfont}{base}
\texto{271.2017}{317.7630}{11.9552}{\normalfont}{de}
\texto{287.2360}{317.7630}{11.9552}{\normalfont}{datos}
\texto{318.9436}{317.7630}{11.9552}{\normalfont}{centralizada,}
\texto{387.9968}{317.7630}{11.9552}{\normalfont}{garantizando}
\texto{458.3388}{317.7630}{11.9552}{\normalfont}{la}
\texto{471.7717}{317.7630}{11.9552}{\normalfont}{coherencia}
\texto{529.4292}{317.7630}{11.9552}{\normalfont}{de}
\texto{70.8660}{303.3170}{11.9552}{\normalfont}{la}
\texto{83.8685}{303.3170}{11.9552}{\normalfont}{información}
\texto{147.9244}{303.3170}{11.9552}{\normalfont}{en}
\texto{163.5403}{303.3170}{11.9552}{\normalfont}{tiempo}
\texto{202.8801}{303.3170}{11.9552}{\normalfont}{real.}
\end{pagina}
% Página original 14
\begin{pagina}{14}
\texto{70.8660}{746.3890}{11.9552}{\normalfont}{4}
\texto{266.5910}{746.3890}{11.9552}{\normalfont}{CAPÍTULO}
\texto{333.5545}{746.3890}{11.9552}{\normalfont}{1.}
\texto{350.4663}{746.3890}{11.9552}{\normalfont}{FORMULACIÓN}
\texto{445.2256}{746.3890}{11.9552}{\normalfont}{DEL}
\texto{473.3418}{746.3890}{11.9552}{\normalfont}{PROYECTO}
\texto{70.8660}{705.1940}{17.2154}{\normalfont\bfseries}{1.3.}
\texto{120.3603}{705.1940}{17.2154}{\normalfont\bfseries}{Diagnóstico}
\texto{225.4259}{705.1940}{17.2154}{\normalfont\bfseries}{de}
\texto{251.4779}{705.1940}{17.2154}{\normalfont\bfseries}{la}
\texto{272.7286}{705.1940}{17.2154}{\normalfont\bfseries}{situación}
\texto{354.3417}{705.1940}{17.2154}{\normalfont\bfseries}{actual}
\texto{88.7990}{678.5380}{11.9552}{\normalfont}{Para}
\texto{116.2864}{678.5380}{11.9552}{\normalfont}{diagnosticar}
\texto{181.7459}{678.5380}{11.9552}{\normalfont}{la}
\texto{194.4495}{678.5380}{11.9552}{\normalfont}{situación}
\texto{243.6380}{678.5380}{11.9552}{\normalfont}{actual}
\texto{278.4539}{678.5380}{11.9552}{\normalfont}{se}
\texto{291.8724}{678.5380}{11.9552}{\normalfont}{aplicaron}
\texto{342.2971}{678.5380}{11.9552}{\normalfont}{técnicas}
\texto{386.2826}{678.5380}{11.9552}{\normalfont}{de}
\texto{401.5877}{678.5380}{11.9552}{\normalfont}{recopilación}
\texto{465.6819}{678.5380}{11.9552}{\normalfont}{de}
\texto{480.9869}{678.5380}{11.9552}{\normalfont}{información}
\texto{70.8660}{664.0920}{11.9552}{\normalfont}{como}
\texto{102.4086}{664.0920}{11.9552}{\normalfont}{entrevistas}
\texto{162.0364}{664.0920}{11.9552}{\normalfont}{semiestructuradas}
\texto{258.8089}{664.0920}{11.9552}{\normalfont}{a}
\texto{269.5519}{664.0920}{11.9552}{\normalfont}{operadores}
\texto{329.4474}{664.0920}{11.9552}{\normalfont}{turísticos,}
\texto{384.5322}{664.0920}{11.9552}{\normalfont}{revisión}
\texto{428.8214}{664.0920}{11.9552}{\normalfont}{documental}
\texto{492.8822}{664.0920}{11.9552}{\normalfont}{y}
\texto{503.9383}{664.0920}{11.9552}{\normalfont}{análisis}
\texto{70.8660}{649.6460}{11.9552}{\normalfont}{de}
\texto{87.0318}{649.6460}{11.9552}{\normalfont}{competencia.}
\texto{158.1748}{649.6460}{11.9552}{\normalfont}{El}
\texto{173.8505}{649.6460}{11.9552}{\normalfont}{diagnóstico}
\texto{235.6182}{649.6460}{11.9552}{\normalfont}{se}
\texto{249.9095}{649.6460}{11.9552}{\normalfont}{centró}
\texto{285.9137}{649.6460}{11.9552}{\normalfont}{tanto}
\texto{317.3774}{649.6460}{11.9552}{\normalfont}{en}
\texto{333.5552}{649.6460}{11.9552}{\normalfont}{el}
\texto{346.4692}{649.6460}{11.9552}{\normalfont}{punto}
\texto{380.5344}{649.6460}{11.9552}{\normalfont}{de}
\texto{396.7121}{649.6460}{11.9552}{\normalfont}{venta}
\texto{428.8166}{649.6460}{11.9552}{\normalfont}{hacia}
\texto{459.9528}{649.6460}{11.9552}{\normalfont}{el}
\texto{472.8787}{649.6460}{11.9552}{\normalfont}{turista}
\texto{511.2214}{649.6460}{11.9552}{\normalfont}{como,}
\texto{70.8660}{635.2000}{11.9552}{\normalfont}{de}
\texto{86.7090}{635.2000}{11.9552}{\normalfont}{manera}
\texto{128.5785}{635.2000}{11.9552}{\normalfont}{destacada,}
\texto{186.1093}{635.2000}{11.9552}{\normalfont}{en}
\texto{201.9643}{635.2000}{11.9552}{\normalfont}{la}
\texto{215.2179}{635.2000}{11.9552}{\normalfont}{gestión}
\texto{255.1889}{635.2000}{11.9552}{\normalfont}{interna}
\texto{295.4349}{635.2000}{11.9552}{\normalfont}{que}
\texto{317.4563}{635.2000}{11.9552}{\normalfont}{realizan}
\texto{361.2769}{635.2000}{11.9552}{\normalfont}{los}
\texto{379.1476}{635.2000}{11.9552}{\normalfont}{operadores,}
\texto{441.5657}{635.2000}{11.9552}{\normalfont}{ya}
\texto{457.4111}{635.2000}{11.9552}{\normalfont}{que}
\texto{479.4445}{635.2000}{11.9552}{\normalfont}{el}
\texto{492.0357}{635.2000}{11.9552}{\normalfont}{backoffice}
\texto{70.8660}{620.7540}{11.9552}{\normalfont}{desktop}
\texto{114.1749}{620.7540}{11.9552}{\normalfont}{constituye}
\texto{170.1683}{620.7540}{11.9552}{\normalfont}{el}
\texto{182.5204}{620.7540}{11.9552}{\normalfont}{foco}
\texto{207.2390}{620.7540}{11.9552}{\normalfont}{principal}
\texto{256.0114}{620.7540}{11.9552}{\normalfont}{del}
\texto{274.8791}{620.7540}{11.9552}{\normalfont}{prototipo}
\texto{326.5758}{620.7540}{11.9552}{\normalfont}{de}
\texto{342.1917}{620.7540}{11.9552}{\normalfont}{esta}
\texto{366.3149}{620.7540}{11.9552}{\normalfont}{etapa.}
\texto{70.8660}{584.9660}{14.3462}{\normalfont\bfseries}{1.3.1.}
\texto{124.6642}{584.9660}{14.3462}{\normalfont\bfseries}{Proceso}
\texto{185.1493}{584.9660}{14.3462}{\normalfont\bfseries}{actual}
\texto{233.1201}{584.9660}{14.3462}{\normalfont\bfseries}{de}
\texto{254.8302}{584.9660}{14.3462}{\normalfont\bfseries}{gestión}
\texto{309.8063}{584.9660}{14.3462}{\normalfont\bfseries}{de}
\texto{331.5164}{584.9660}{14.3462}{\normalfont\bfseries}{los}
\texto{355.8160}{584.9660}{14.3462}{\normalfont\bfseries}{operadores}
\texto{88.7990}{562.4280}{11.9552}{\normalfont}{La}
\texto{107.3702}{562.4280}{11.9552}{\normalfont}{gestión}
\texto{148.6085}{562.4280}{11.9552}{\normalfont}{de}
\texto{165.7188}{562.4280}{11.9552}{\normalfont}{tours}
\texto{197.2135}{562.4280}{11.9552}{\normalfont}{por}
\texto{219.8495}{562.4280}{11.9552}{\normalfont}{parte}
\texto{251.9301}{562.4280}{11.9552}{\normalfont}{de}
\texto{269.0404}{562.4280}{11.9552}{\normalfont}{los}
\texto{288.1663}{562.4280}{11.9552}{\normalfont}{operadores}
\texto{348.5998}{562.4280}{11.9552}{\normalfont}{nicaragüenses}
\texto{423.7215}{562.4280}{11.9552}{\normalfont}{se}
\texto{438.9453}{562.4280}{11.9552}{\normalfont}{caracteriza}
\texto{499.6418}{562.4280}{11.9552}{\normalfont}{por}
\texto{522.2778}{562.4280}{11.9552}{\normalfont}{una}
\texto{70.8660}{547.9820}{11.9552}{\normalfont}{fragmentación}
\texto{148.7613}{547.9820}{11.9552}{\normalfont}{de}
\texto{166.1824}{547.9820}{11.9552}{\normalfont}{la}
\texto{181.0140}{547.9820}{11.9552}{\normalfont}{oferta}
\texto{216.3201}{547.9820}{11.9552}{\normalfont}{y}
\texto{228.2132}{547.9820}{11.9552}{\normalfont}{una}
\texto{252.7883}{547.9820}{11.9552}{\normalfont}{escasa}
\texto{289.8494}{547.9820}{11.9552}{\normalfont}{digitalización.}
\texto{366.4655}{547.9820}{11.9552}{\normalfont}{Las}
\texto{389.9646}{547.9820}{11.9552}{\normalfont}{principales}
\texto{450.3743}{547.9820}{11.9552}{\normalfont}{formas}
\texto{490.2975}{547.9820}{11.9552}{\normalfont}{en}
\texto{507.7186}{547.9820}{11.9552}{\normalfont}{que}
\texto{531.3181}{547.9820}{11.9552}{\normalfont}{se}
\texto{70.8660}{533.5370}{11.9552}{\normalfont}{gestiona}
\texto{116.4512}{533.5370}{11.9552}{\normalfont}{la}
\texto{129.4537}{533.5370}{11.9552}{\normalfont}{oferta}
\texto{162.9545}{533.5370}{11.9552}{\normalfont}{hoy}
\texto{185.0645}{533.5370}{11.9552}{\normalfont}{son:}
\texto{100.7540}{504.8610}{11.9552}{\normalfont\bfseries}{Registro}
\texto{155.0509}{504.8610}{11.9552}{\normalfont\bfseries}{y}
\texto{167.0755}{504.8610}{11.9552}{\normalfont\bfseries}{actualización}
\texto{247.6559}{504.8610}{11.9552}{\normalfont\bfseries}{de}
\texto{266.1901}{504.8610}{11.9552}{\normalfont\bfseries}{tours}
\texto{301.3419}{504.8610}{11.9552}{\normalfont\bfseries}{en}
\texto{319.8761}{504.8610}{11.9552}{\normalfont\bfseries}{papel}
\texto{356.5391}{504.8610}{11.9552}{\normalfont\bfseries}{u}
\texto{368.9366}{504.8610}{11.9552}{\normalfont\bfseries}{hojas}
\texto{404.3838}{504.8610}{11.9552}{\normalfont\bfseries}{de}
\texto{422.9179}{504.8610}{11.9552}{\normalfont\bfseries}{cálculo:}
\texto{471.1030}{504.8610}{11.9552}{\normalfont}{la}
\texto{484.4880}{504.8610}{11.9552}{\normalfont}{mayoría}
\texto{529.4205}{504.8610}{11.9552}{\normalfont}{de}
\texto{100.7540}{490.4150}{11.9552}{\normalfont}{las}
\texto{119.1985}{490.4150}{11.9552}{\normalfont}{agencias}
\texto{166.2709}{490.4150}{11.9552}{\normalfont}{y}
\texto{177.1836}{490.4150}{11.9552}{\normalfont}{guías}
\texto{207.9969}{490.4150}{11.9552}{\normalfont}{independientes}
\texto{288.2378}{490.4150}{11.9552}{\normalfont}{mantienen}
\texto{345.9791}{490.4150}{11.9552}{\normalfont}{su}
\texto{361.8340}{490.4150}{11.9552}{\normalfont}{catálogo}
\texto{408.8419}{490.4150}{11.9552}{\normalfont}{de}
\texto{425.2827}{490.4150}{11.9552}{\normalfont}{tours,}
\texto{459.3478}{490.4150}{11.9552}{\normalfont}{precios}
\texto{499.2662}{490.4150}{11.9552}{\normalfont}{y}
\texto{510.1790}{490.4150}{11.9552}{\normalfont}{dispo-}
\texto{100.7540}{475.9690}{11.9552}{\normalfont}{nibilidad}
\texto{149.8612}{475.9690}{11.9552}{\normalfont}{en}
\texto{165.8118}{475.9690}{11.9552}{\normalfont}{documentos}
\texto{230.5923}{475.9690}{11.9552}{\normalfont}{locales}
\texto{268.3922}{475.9690}{11.9552}{\normalfont}{(Word,}
\texto{308.3967}{475.9690}{11.9552}{\normalfont}{Excel)}
\texto{344.9820}{475.9690}{11.9552}{\normalfont}{o}
\texto{355.0794}{475.9690}{11.9552}{\normalfont}{en}
\texto{371.0300}{475.9690}{11.9552}{\normalfont}{libretas}
\texto{413.0477}{475.9690}{11.9552}{\normalfont}{físicas,}
\texto{450.9146}{475.9690}{11.9552}{\normalfont}{sin}
\texto{469.5193}{475.9690}{11.9552}{\normalfont}{un}
\texto{486.7707}{475.9690}{11.9552}{\normalfont}{repositorio}
\texto{100.7540}{461.5230}{11.9552}{\normalfont}{centralizado.}
\texto{100.7540}{436.2020}{11.9552}{\normalfont\bfseries}{Reservas}
\texto{156.4353}{436.2020}{11.9552}{\normalfont\bfseries}{por}
\texto{180.8957}{436.2020}{11.9552}{\normalfont\bfseries}{canales}
\texto{226.9878}{436.2020}{11.9552}{\normalfont\bfseries}{informales:}
\texto{295.0060}{436.2020}{11.9552}{\normalfont}{WhatsApp,}
\texto{357.4074}{436.2020}{11.9552}{\normalfont}{llamadas}
\texto{406.1726}{436.2020}{11.9552}{\normalfont}{telefónicas}
\texto{463.0555}{436.2020}{11.9552}{\normalfont}{y}
\texto{473.0596}{436.2020}{11.9552}{\normalfont}{redes}
\texto{502.9524}{436.2020}{11.9552}{\normalfont}{sociales}
\texto{100.7540}{421.7560}{11.9552}{\normalfont}{son}
\texto{121.0396}{421.7560}{11.9552}{\normalfont}{los}
\texto{138.0734}{421.7560}{11.9552}{\normalfont}{canales}
\texto{177.8698}{421.7560}{11.9552}{\normalfont}{más}
\texto{201.4072}{421.7560}{11.9552}{\normalfont}{frecuentes}
\texto{255.5021}{421.7560}{11.9552}{\normalfont}{para}
\texto{281.5764}{421.7560}{11.9552}{\normalfont}{recibir}
\texto{317.4061}{421.7560}{11.9552}{\normalfont}{reservas.}
\texto{364.1008}{421.7560}{11.9552}{\normalfont}{Estas}
\texto{395.0169}{421.7560}{11.9552}{\normalfont}{reservas}
\texto{438.4597}{421.7560}{11.9552}{\normalfont}{se}
\texto{451.5913}{421.7560}{11.9552}{\normalfont}{registran}
\texto{499.8425}{421.7560}{11.9552}{\normalfont}{manual-}
\texto{100.7540}{407.3100}{11.9552}{\normalfont}{mente}
\texto{135.5460}{407.3100}{11.9552}{\normalfont}{y}
\texto{145.6219}{407.3100}{11.9552}{\normalfont}{no}
\texto{161.8762}{407.3100}{11.9552}{\normalfont}{siempre}
\texto{204.8718}{407.3100}{11.9552}{\normalfont}{se}
\texto{218.5892}{407.3100}{11.9552}{\normalfont}{confirma}
\texto{266.7233}{407.3100}{11.9552}{\normalfont}{al}
\texto{279.7257}{407.3100}{11.9552}{\normalfont}{turista}
\texto{317.5066}{407.3100}{11.9552}{\normalfont}{de}
\texto{333.1224}{407.3100}{11.9552}{\normalfont}{manera}
\texto{374.7409}{407.3100}{11.9552}{\normalfont}{formal.}
\texto{100.7540}{381.9890}{11.9552}{\normalfont\bfseries}{Gestión}
\texto{151.0005}{381.9890}{11.9552}{\normalfont\bfseries}{manual}
\texto{198.6635}{381.9890}{11.9552}{\normalfont\bfseries}{de}
\texto{217.3411}{381.9890}{11.9552}{\normalfont\bfseries}{disponibilidad:}
\texto{307.9580}{381.9890}{11.9552}{\normalfont}{la}
\texto{321.4745}{381.9890}{11.9552}{\normalfont}{disponibilidad}
\texto{397.8133}{381.9890}{11.9552}{\normalfont}{de}
\texto{413.9432}{381.9890}{11.9552}{\normalfont}{un}
\texto{431.3620}{381.9890}{11.9552}{\normalfont}{tour}
\texto{457.2474}{381.9890}{11.9552}{\normalfont}{(fechas,}
\texto{500.0972}{381.9890}{11.9552}{\normalfont}{horarios}
\texto{100.7540}{367.5430}{11.9552}{\normalfont}{y}
\texto{111.8580}{367.5430}{11.9552}{\normalfont}{cupos)}
\texto{150.3394}{367.5430}{11.9552}{\normalfont}{se}
\texto{165.0969}{367.5430}{11.9552}{\normalfont}{maneja}
\texto{206.7679}{367.5430}{11.9552}{\normalfont}{de}
\texto{223.4120}{367.5430}{11.9552}{\normalfont}{forma}
\texto{257.9290}{367.5430}{11.9552}{\normalfont}{manual,}
\texto{303.5046}{367.5430}{11.9552}{\normalfont}{lo}
\texto{317.5472}{367.5430}{11.9552}{\normalfont}{que}
\texto{340.3577}{367.5430}{11.9552}{\normalfont}{provoca}
\texto{384.9578}{367.5430}{11.9552}{\normalfont}{dobles}
\texto{421.8276}{367.5430}{11.9552}{\normalfont}{reservas,}
\texto{470.1362}{367.5430}{11.9552}{\normalfont}{sobreventa}
\texto{529.4364}{367.5430}{11.9552}{\normalfont}{de}
\texto{100.7540}{353.0970}{11.9552}{\normalfont}{cupos}
\texto{133.6547}{353.0970}{11.9552}{\normalfont}{y}
\texto{143.7306}{353.0970}{11.9552}{\normalfont}{dificultad}
\texto{195.7668}{353.0970}{11.9552}{\normalfont}{para}
\texto{222.4269}{353.0970}{11.9552}{\normalfont}{actualizar}
\texto{276.4141}{353.0970}{11.9552}{\normalfont}{la}
\texto{289.4166}{353.0970}{11.9552}{\normalfont}{información}
\texto{353.4726}{353.0970}{11.9552}{\normalfont}{en}
\texto{369.0885}{353.0970}{11.9552}{\normalfont}{tiempo}
\texto{408.4282}{353.0970}{11.9552}{\normalfont}{real.}
\texto{100.7540}{327.7750}{11.9552}{\normalfont\bfseries}{Control}
\texto{149.1486}{327.7750}{11.9552}{\normalfont\bfseries}{sin}
\texto{169.3446}{327.7750}{11.9552}{\normalfont\bfseries}{trazabilidad:}
\texto{246.2310}{327.7750}{11.9552}{\normalfont}{la}
\texto{258.5281}{327.7750}{11.9552}{\normalfont}{ausencia}
\texto{304.7206}{327.7750}{11.9552}{\normalfont}{de}
\texto{319.6311}{327.7750}{11.9552}{\normalfont}{un}
\texto{335.8424}{327.7750}{11.9552}{\normalfont}{voucher}
\texto{378.3981}{327.7750}{11.9552}{\normalfont}{digital}
\texto{414.1203}{327.7750}{11.9552}{\normalfont}{y}
\texto{423.5027}{327.7750}{11.9552}{\normalfont}{de}
\texto{438.4132}{327.7750}{11.9552}{\normalfont}{registro}
\texto{480.0532}{327.7750}{11.9552}{\normalfont}{de}
\texto{494.9637}{327.7750}{11.9552}{\normalfont}{auditoría}
\texto{100.7540}{313.3300}{11.9552}{\normalfont}{dificulta}
\texto{146.2746}{313.3300}{11.9552}{\normalfont}{la}
\texto{159.2771}{313.3300}{11.9552}{\normalfont}{validación}
\texto{214.2328}{313.3300}{11.9552}{\normalfont}{posterior}
\texto{263.4045}{313.3300}{11.9552}{\normalfont}{de}
\texto{279.0084}{313.3300}{11.9552}{\normalfont}{los}
\texto{296.6400}{313.3300}{11.9552}{\normalfont}{servicios}
\texto{343.2652}{313.3300}{11.9552}{\normalfont}{y}
\texto{353.3411}{313.3300}{11.9552}{\normalfont}{la}
\texto{366.3555}{313.3300}{11.9552}{\normalfont}{gestión}
\texto{406.0874}{313.3300}{11.9552}{\normalfont}{de}
\texto{421.7033}{313.3300}{11.9552}{\normalfont}{incidencias}
\texto{480.9461}{313.3300}{11.9552}{\normalfont}{o}
\texto{490.6968}{313.3300}{11.9552}{\normalfont}{reclamos.}
\texto{70.8660}{277.5420}{14.3462}{\normalfont\bfseries}{1.3.2.}
\texto{124.6642}{277.5420}{14.3462}{\normalfont\bfseries}{Personas}
\texto{192.1818}{277.5420}{14.3462}{\normalfont\bfseries}{y}
\texto{206.0804}{277.5420}{14.3462}{\normalfont\bfseries}{recursos}
\texto{268.9512}{277.5420}{14.3462}{\normalfont\bfseries}{involucrados}
\texto{88.7990}{255.0040}{11.9552}{\normalfont}{En}
\texto{107.1168}{255.0040}{11.9552}{\normalfont}{el}
\texto{119.4330}{255.0040}{11.9552}{\normalfont}{proceso}
\texto{161.4029}{255.0040}{11.9552}{\normalfont}{actual}
\texto{196.4819}{255.0040}{11.9552}{\normalfont}{intervienen}
\texto{256.9154}{255.0040}{11.9552}{\normalfont}{el}
\texto{269.2790}{255.0040}{11.9552}{\normalfont\bfseries}{operador}
\texto{325.5670}{255.0040}{11.9552}{\normalfont}{o}
\texto{335.2698}{255.0040}{11.9552}{\normalfont}{su}
\texto{350.2521}{255.0040}{11.9552}{\normalfont}{personal}
\texto{396.7746}{255.0040}{11.9552}{\normalfont}{administrativo}
\texto{474.8444}{255.0040}{11.9552}{\normalfont}{(quien}
\texto{510.8869}{255.0040}{11.9552}{\normalfont}{recibe}
\texto{70.8660}{240.5580}{11.9552}{\normalfont}{las}
\texto{88.0791}{240.5580}{11.9552}{\normalfont}{reservas}
\texto{131.7132}{240.5580}{11.9552}{\normalfont}{manualmente),}
\texto{211.3013}{240.5580}{11.9552}{\normalfont}{el}
\texto{223.3040}{240.5580}{11.9552}{\normalfont\bfseries}{turista}
\texto{266.0710}{240.5580}{11.9552}{\normalfont}{(quien}
\texto{301.7549}{240.5580}{11.9552}{\normalfont}{consulta}
\texto{347.5959}{240.5580}{11.9552}{\normalfont}{y}
\texto{357.2772}{240.5580}{11.9552}{\normalfont}{reserva}
\texto{396.2823}{240.5580}{11.9552}{\normalfont}{por}
\texto{417.0174}{240.5580}{11.9552}{\normalfont}{canales}
\texto{457.0051}{240.5580}{11.9552}{\normalfont}{informales)}
\texto{517.4793}{240.5580}{11.9552}{\normalfont}{y,}
\texto{529.4321}{240.5580}{11.9552}{\normalfont}{en}
\texto{70.8660}{226.1120}{11.9552}{\normalfont}{algunos}
\texto{113.5819}{226.1120}{11.9552}{\normalfont}{casos,}
\texto{147.2690}{226.1120}{11.9552}{\normalfont\bfseries}{plataformas}
\texto{221.0780}{226.1120}{11.9552}{\normalfont\bfseries}{extranjeras}
\texto{290.9240}{226.1120}{11.9552}{\normalfont}{como}
\texto{321.8688}{226.1120}{11.9552}{\normalfont}{intermediarias}
\texto{398.7312}{226.1120}{11.9552}{\normalfont}{que}
\texto{420.9081}{226.1120}{11.9552}{\normalfont}{ofrecen}
\texto{461.2832}{226.1120}{11.9552}{\normalfont}{reserva}
\texto{501.0773}{226.1120}{11.9552}{\normalfont}{en}
\texto{517.0757}{226.1120}{11.9552}{\normalfont}{línea}
\texto{70.8660}{211.6660}{11.9552}{\normalfont}{pero}
\texto{96.3856}{211.6660}{11.9552}{\normalfont}{con}
\texto{117.0298}{211.6660}{11.9552}{\normalfont}{comisiones}
\texto{174.2354}{211.6660}{11.9552}{\normalfont}{elevadas}
\texto{219.3257}{211.6660}{11.9552}{\normalfont}{(15}
\texto{237.5813}{211.6660}{11.9552}{\normalfont}{\%--25}
\texto{266.8930}{211.6660}{11.9552}{\normalfont}{\%)}
\texto{284.2854}{211.6660}{11.9552}{\normalfont}{[6].}
\texto{302.9791}{211.6660}{11.9552}{\normalfont}{Los}
\texto{323.8481}{211.6660}{11.9552}{\normalfont}{recursos}
\texto{368.0464}{211.6660}{11.9552}{\normalfont}{utilizados}
\texto{419.9727}{211.6660}{11.9552}{\normalfont}{son}
\texto{440.0311}{211.6660}{11.9552}{\normalfont}{limitados:}
\texto{493.2700}{211.6660}{11.9552}{\normalfont}{teléfonos,}
\texto{70.8660}{197.2200}{11.9552}{\normalfont}{hojas}
\texto{101.2752}{197.2200}{11.9552}{\normalfont}{de}
\texto{116.6640}{197.2200}{11.9552}{\normalfont}{cálculo,}
\texto{158.7176}{197.2200}{11.9552}{\normalfont}{WhatsApp}
\texto{217.7356}{197.2200}{11.9552}{\normalfont}{y,}
\texto{229.8678}{197.2200}{11.9552}{\normalfont}{en}
\texto{245.2565}{197.2200}{11.9552}{\normalfont}{el}
\texto{257.3934}{197.2200}{11.9552}{\normalfont}{mejor}
\texto{290.0168}{197.2200}{11.9552}{\normalfont}{de}
\texto{305.4055}{197.2200}{11.9552}{\normalfont}{los}
\texto{322.8099}{197.2200}{11.9552}{\normalfont}{casos,}
\texto{355.8875}{197.2200}{11.9552}{\normalfont}{plataformas}
\texto{419.7929}{197.2200}{11.9552}{\normalfont}{de}
\texto{435.1816}{197.2200}{11.9552}{\normalfont}{terceros}
\texto{478.6005}{197.2200}{11.9552}{\normalfont}{cuyo}
\texto{505.6981}{197.2200}{11.9552}{\normalfont}{control}
\texto{70.8660}{182.7740}{11.9552}{\normalfont}{de}
\texto{86.4699}{182.7740}{11.9552}{\normalfont}{la}
\texto{99.4724}{182.7740}{11.9552}{\normalfont}{oferta}
\texto{132.9733}{182.7740}{11.9552}{\normalfont}{escapa}
\texto{170.1037}{182.7740}{11.9552}{\normalfont}{al}
\texto{183.1062}{182.7740}{11.9552}{\normalfont}{operador}
\texto{232.2134}{182.7740}{11.9552}{\normalfont}{local.}
\texto{70.8660}{146.9870}{14.3462}{\normalfont\bfseries}{1.3.3.}
\texto{124.6642}{146.9870}{14.3462}{\normalfont\bfseries}{Necesidades}
\texto{214.9234}{146.9870}{14.3462}{\normalfont\bfseries}{identificadas}
\texto{307.5482}{146.9870}{14.3462}{\normalfont\bfseries}{y}
\texto{321.4468}{146.9870}{14.3462}{\normalfont\bfseries}{principales}
\texto{402.5171}{146.9870}{14.3462}{\normalfont\bfseries}{hallazgos}
\texto{88.7990}{124.4480}{11.9552}{\normalfont}{El}
\texto{103.9128}{124.4480}{11.9552}{\normalfont}{diagnóstico}
\texto{165.1067}{124.4480}{11.9552}{\normalfont}{permitió}
\texto{212.2627}{124.4480}{11.9552}{\normalfont}{identificar}
\texto{267.2160}{124.4480}{11.9552}{\normalfont}{las}
\texto{284.8356}{124.4480}{11.9552}{\normalfont}{siguientes}
\texto{337.9788}{124.4480}{11.9552}{\normalfont}{necesidades}
\texto{400.5380}{124.4480}{11.9552}{\normalfont}{y}
\texto{410.6139}{124.4480}{11.9552}{\normalfont}{hallazgos:}
\texto{100.7540}{95.7730}{11.9552}{\normalfont\bfseries}{Baja}
\texto{133.4802}{95.7730}{11.9552}{\normalfont\bfseries}{penetración}
\texto{207.8594}{95.7730}{11.9552}{\normalfont\bfseries}{digital}
\texto{250.6399}{95.7730}{11.9552}{\normalfont\bfseries}{en}
\texto{269.8675}{95.7730}{11.9552}{\normalfont\bfseries}{la}
\texto{285.7488}{95.7730}{11.9552}{\normalfont\bfseries}{gestión:}
\texto{335.7120}{95.7730}{11.9552}{\normalfont}{la}
\texto{349.6948}{95.7730}{11.9552}{\normalfont}{mayoría}
\texto{395.2250}{95.7730}{11.9552}{\normalfont}{de}
\texto{411.8212}{95.7730}{11.9552}{\normalfont}{los}
\texto{430.4330}{95.7730}{11.9552}{\normalfont}{operadores}
\texto{490.3286}{95.7730}{11.9552}{\normalfont}{locales}
\texto{528.7741}{95.7730}{11.9552}{\normalfont}{no}
\texto{100.7540}{81.3270}{11.9552}{\normalfont}{disponen}
\texto{150.2844}{81.3270}{11.9552}{\normalfont}{de}
\texto{166.2709}{81.3270}{11.9552}{\normalfont}{sistemas}
\texto{213.0062}{81.3270}{11.9552}{\normalfont}{propios}
\texto{254.4214}{81.3270}{11.9552}{\normalfont}{de}
\texto{270.3959}{81.3270}{11.9552}{\normalfont}{reserva}
\texto{310.1780}{81.3270}{11.9552}{\normalfont}{en}
\texto{326.1526}{81.3270}{11.9552}{\normalfont}{línea}
\texto{354.4959}{81.3270}{11.9552}{\normalfont}{ni}
\texto{368.5194}{81.3270}{11.9552}{\normalfont}{de}
\texto{384.5059}{81.3270}{11.9552}{\normalfont}{herramientas}
\texto{454.8527}{81.3270}{11.9552}{\normalfont}{para}
\texto{481.8953}{81.3270}{11.9552}{\normalfont}{administrar}
\texto{100.7540}{66.8810}{11.9552}{\normalfont}{su}
\texto{115.7721}{66.8810}{11.9552}{\normalfont}{catálogo}
\texto{161.9431}{66.8810}{11.9552}{\normalfont}{y}
\texto{172.0309}{66.8810}{11.9552}{\normalfont}{disponibilidad.}
\end{pagina}
% Página original 15
\begin{pagina}{15}
\texto{70.8660}{746.3890}{11.9552}{\normalfont}{1.3.}
\texto{96.8710}{746.3890}{11.9552}{\normalfont}{DIAGNÓSTICO}
\texto{186.1894}{746.3890}{11.9552}{\normalfont}{DE}
\texto{207.0034}{746.3890}{11.9552}{\normalfont}{LA}
\texto{226.9902}{746.3890}{11.9552}{\normalfont}{SITUACIÓN}
\texto{297.8534}{746.3890}{11.9552}{\normalfont}{ACTUAL}
\texto{535.2810}{746.3890}{11.9552}{\normalfont}{5}
\texto{100.7540}{705.1940}{11.9552}{\normalfont\bfseries}{Errores}
\texto{150.1708}{705.1940}{11.9552}{\normalfont\bfseries}{de}
\texto{169.7212}{705.1940}{11.9552}{\normalfont\bfseries}{disponibilidad}
\texto{258.1263}{705.1940}{11.9552}{\normalfont\bfseries}{y}
\texto{271.1790}{705.1940}{11.9552}{\normalfont\bfseries}{doble}
\texto{308.6621}{705.1940}{11.9552}{\normalfont\bfseries}{reserva:}
\texto{359.0400}{705.1940}{11.9552}{\normalfont}{la}
\texto{373.3097}{705.1940}{11.9552}{\normalfont}{gestión}
\texto{414.3208}{705.1940}{11.9552}{\normalfont}{manual}
\texto{456.8957}{705.1940}{11.9552}{\normalfont}{de}
\texto{473.7788}{705.1940}{11.9552}{\normalfont}{cupos}
\texto{507.9588}{705.1940}{11.9552}{\normalfont}{genera}
\texto{100.7540}{690.7480}{11.9552}{\normalfont}{pérdida}
\texto{143.3456}{690.7480}{11.9552}{\normalfont}{de}
\texto{158.9495}{690.7480}{11.9552}{\normalfont}{ingresos}
\texto{203.3105}{690.7480}{11.9552}{\normalfont}{y}
\texto{213.3863}{690.7480}{11.9552}{\normalfont}{deterioro}
\texto{262.8211}{690.7480}{11.9552}{\normalfont}{de}
\texto{278.4250}{690.7480}{11.9552}{\normalfont}{la}
\texto{291.4275}{690.7480}{11.9552}{\normalfont}{confianza}
\texto{342.8133}{690.7480}{11.9552}{\normalfont}{del}
\texto{361.6691}{690.7480}{11.9552}{\normalfont}{turista.}
\texto{100.7540}{666.3390}{11.9552}{\normalfont\bfseries}{Comisiones}
\texto{171.4559}{666.3390}{11.9552}{\normalfont\bfseries}{elevadas}
\texto{224.3158}{666.3390}{11.9552}{\normalfont\bfseries}{de}
\texto{242.5750}{666.3390}{11.9552}{\normalfont\bfseries}{plataformas}
\texto{316.1090}{666.3390}{11.9552}{\normalfont\bfseries}{extranjeras:}
\texto{389.4430}{666.3390}{11.9552}{\normalfont}{servicios}
\texto{436.2117}{666.3390}{11.9552}{\normalfont}{como}
\texto{466.9190}{666.3390}{11.9552}{\normalfont\itshape}{Civitatis}
\texto{514.4700}{666.3390}{11.9552}{\normalfont}{redu-}
\texto{100.7540}{651.8940}{11.9552}{\normalfont}{cen}
\texto{121.5608}{651.8940}{11.9552}{\normalfont}{significativamente}
\texto{216.5808}{651.8940}{11.9552}{\normalfont}{los}
\texto{234.2003}{651.8940}{11.9552}{\normalfont}{ingresos}
\texto{278.5493}{651.8940}{11.9552}{\normalfont}{de}
\texto{294.1652}{651.8940}{11.9552}{\normalfont}{los}
\texto{311.7848}{651.8940}{11.9552}{\normalfont}{proveedores}
\texto{375.2597}{651.8940}{11.9552}{\normalfont}{nicaragüenses}
\texto{448.8870}{651.8940}{11.9552}{\normalfont}{[6].}
\texto{100.7540}{627.4850}{11.9552}{\normalfont\bfseries}{F}
\texto{108.0969}{627.4850}{11.9552}{\normalfont\bfseries}{alta}
\texto{135.9979}{627.4850}{11.9552}{\normalfont\bfseries}{de}
\texto{155.4646}{627.4850}{11.9552}{\normalfont\bfseries}{confianza}
\texto{215.5072}{627.4850}{11.9552}{\normalfont\bfseries}{en}
\texto{234.9858}{627.4850}{11.9552}{\normalfont\bfseries}{pagos}
\texto{273.6214}{627.4850}{11.9552}{\normalfont\bfseries}{en}
\texto{293.1000}{627.4850}{11.9552}{\normalfont\bfseries}{línea:}
\texto{329.5670}{627.4850}{11.9552}{\normalfont}{el}
\texto{343.1266}{627.4850}{11.9552}{\normalfont}{turista}
\texto{382.1149}{627.4850}{11.9552}{\normalfont}{expresa}
\texto{425.3305}{627.4850}{11.9552}{\normalfont}{reticencia}
\texto{479.2126}{627.4850}{11.9552}{\normalfont}{a}
\texto{490.1827}{627.4850}{11.9552}{\normalfont}{pagar}
\texto{523.9036}{627.4850}{11.9552}{\normalfont}{por}
\texto{100.7540}{613.0390}{11.9552}{\normalfont}{servicios}
\texto{147.3793}{613.0390}{11.9552}{\normalfont}{de}
\texto{162.9832}{613.0390}{11.9552}{\normalfont}{proveedores}
\texto{226.4581}{613.0390}{11.9552}{\normalfont}{que}
\texto{248.2525}{613.0390}{11.9552}{\normalfont}{no}
\texto{264.5068}{613.0390}{11.9552}{\normalfont}{conoce}
\texto{302.5578}{613.0390}{11.9552}{\normalfont}{previamente.}
\texto{100.7540}{588.6310}{11.9552}{\normalfont\bfseries}{F}
\texto{108.0969}{588.6310}{11.9552}{\normalfont\bfseries}{alta}
\texto{134.6231}{588.6310}{11.9552}{\normalfont\bfseries}{de}
\texto{152.7388}{588.6310}{11.9552}{\normalfont\bfseries}{un}
\texto{172.1780}{588.6310}{11.9552}{\normalfont\bfseries}{canal}
\texto{206.9353}{588.6310}{11.9552}{\normalfont\bfseries}{de}
\texto{225.0510}{588.6310}{11.9552}{\normalfont\bfseries}{administración}
\texto{316.1867}{588.6310}{11.9552}{\normalfont\bfseries}{y}
\texto{327.7928}{588.6310}{11.9552}{\normalfont\bfseries}{supervisión:}
\texto{401.7710}{588.6310}{11.9552}{\normalfont}{no}
\texto{418.0372}{588.6310}{11.9552}{\normalfont}{existe}
\texto{450.9643}{588.6310}{11.9552}{\normalfont}{una}
\texto{473.7341}{588.6310}{11.9552}{\normalfont}{instancia}
\texto{523.2454}{588.6310}{11.9552}{\normalfont}{que}
\texto{100.7540}{574.1850}{11.9552}{\normalfont}{verifique}
\texto{147.0661}{574.1850}{11.9552}{\normalfont}{a}
\texto{155.9799}{574.1850}{11.9552}{\normalfont}{los}
\texto{172.7626}{574.1850}{11.9552}{\normalfont}{operadores,}
\texto{234.1047}{574.1850}{11.9552}{\normalfont}{modere}
\texto{274.5587}{574.1850}{11.9552}{\normalfont}{la}
\texto{286.7363}{574.1850}{11.9552}{\normalfont}{oferta}
\texto{319.3883}{574.1850}{11.9552}{\normalfont}{y}
\texto{328.6273}{574.1850}{11.9552}{\normalfont}{supervise}
\texto{378.6526}{574.1850}{11.9552}{\normalfont}{la}
\texto{390.8182}{574.1850}{11.9552}{\normalfont}{actividad}
\texto{441.0301}{574.1850}{11.9552}{\normalfont}{económica}
\texto{496.7820}{574.1850}{11.9552}{\normalfont}{de}
\texto{511.5490}{574.1850}{11.9552}{\normalfont}{forma}
\texto{100.7540}{559.7390}{11.9552}{\normalfont}{centralizada.}
\texto{88.7990}{532.3420}{11.9552}{\normalfont}{Estos}
\texto{120.7433}{532.3420}{11.9552}{\normalfont}{hallazgos}
\texto{171.3353}{532.3420}{11.9552}{\normalfont}{sustentan}
\texto{224.2586}{532.3420}{11.9552}{\normalfont}{la}
\texto{237.7154}{532.3420}{11.9552}{\normalfont}{necesidad}
\texto{290.8969}{532.3420}{11.9552}{\normalfont}{de}
\texto{306.9551}{532.3420}{11.9552}{\normalfont}{una}
\texto{330.1554}{532.3420}{11.9552}{\normalfont}{solución}
\texto{375.5445}{532.3420}{11.9552}{\normalfont}{que}
\texto{397.7811}{532.3420}{11.9552}{\normalfont}{brinde}
\texto{434.6390}{532.3420}{11.9552}{\normalfont}{a}
\texto{444.8440}{532.3420}{11.9552}{\normalfont}{los}
\texto{462.9059}{532.3420}{11.9552}{\normalfont}{operadores}
\texto{522.2754}{532.3420}{11.9552}{\normalfont}{una}
\texto{70.8660}{517.8960}{11.9552}{\normalfont}{herramienta}
\texto{135.9979}{517.8960}{11.9552}{\normalfont}{de}
\texto{151.3747}{517.8960}{11.9552}{\normalfont}{escritorio}
\texto{201.9476}{517.8960}{11.9552}{\normalfont}{para}
\texto{228.3805}{517.8960}{11.9552}{\normalfont}{gestionar}
\texto{278.2911}{517.8960}{11.9552}{\normalfont}{eficientemente}
\texto{354.1564}{517.8960}{11.9552}{\normalfont}{su}
\texto{368.9593}{517.8960}{11.9552}{\normalfont}{negocio}
\texto{410.6734}{517.8960}{11.9552}{\normalfont}{(el}
\texto{427.3629}{517.8960}{11.9552}{\normalfont}{backoffice)}
\texto{484.6929}{517.8960}{11.9552}{\normalfont}{y}
\texto{494.5416}{517.8960}{11.9552}{\normalfont}{al}
\texto{507.3288}{517.8960}{11.9552}{\normalfont}{equipo}
\texto{70.8660}{503.4500}{11.9552}{\normalfont}{gestor}
\texto{106.9731}{503.4500}{11.9552}{\normalfont}{un}
\texto{125.4558}{503.4500}{11.9552}{\normalfont}{panel}
\texto{158.2585}{503.4500}{11.9552}{\normalfont}{de}
\texto{175.4405}{503.4500}{11.9552}{\normalfont}{escritorio}
\texto{227.8067}{503.4500}{11.9552}{\normalfont}{para}
\texto{256.0568}{503.4500}{11.9552}{\normalfont}{supervisar}
\texto{313.6881}{503.4500}{11.9552}{\normalfont}{la}
\texto{328.2806}{503.4500}{11.9552}{\normalfont}{plataforma}
\texto{389.3621}{503.4500}{11.9552}{\normalfont}{(el}
\texto{407.8448}{503.4500}{11.9552}{\normalfont}{panel}
\texto{440.6475}{503.4500}{11.9552}{\normalfont}{de}
\texto{457.8295}{503.4500}{11.9552}{\normalfont}{administración).}
\texto{70.8660}{489.0040}{11.9552}{\normalfont}{Ambas}
\texto{109.0724}{489.0040}{11.9552}{\normalfont}{aplicaciones}
\texto{172.6574}{489.0040}{11.9552}{\normalfont}{de}
\texto{187.4005}{489.0040}{11.9552}{\normalfont}{escritorio}
\texto{237.3278}{489.0040}{11.9552}{\normalfont}{constituyen}
\texto{298.9640}{489.0040}{11.9552}{\normalfont}{el}
\texto{310.4554}{489.0040}{11.9552}{\normalfont}{foco}
\texto{334.3012}{489.0040}{11.9552}{\normalfont}{del}
\texto{352.2962}{489.0040}{11.9552}{\normalfont}{prototipo}
\texto{403.1321}{489.0040}{11.9552}{\normalfont}{de}
\texto{417.8752}{489.0040}{11.9552}{\normalfont}{esta}
\texto{441.1377}{489.0040}{11.9552}{\normalfont}{etapa,}
\texto{475.3917}{489.0040}{11.9552}{\normalfont}{consumiendo}
\texto{70.8660}{474.5580}{11.9552}{\normalfont}{una}
\texto{93.1457}{474.5580}{11.9552}{\normalfont}{API}
\texto{117.5271}{474.5580}{11.9552}{\normalfont}{REST}
\texto{152.4842}{474.5580}{11.9552}{\normalfont}{común}
\texto{189.3994}{474.5580}{11.9552}{\normalfont}{que}
\texto{210.7155}{474.5580}{11.9552}{\normalfont}{en}
\texto{225.8413}{474.5580}{11.9552}{\normalfont}{el}
\texto{237.7152}{474.5580}{11.9552}{\normalfont}{futuro}
\texto{272.6889}{474.5580}{11.9552}{\normalfont}{también}
\texto{317.4085}{474.5580}{11.9552}{\normalfont}{alimentará}
\texto{375.1354}{474.5580}{11.9552}{\normalfont}{la}
\texto{387.6597}{474.5580}{11.9552}{\normalfont}{aplicación}
\texto{441.8072}{474.5580}{11.9552}{\normalfont}{móvil}
\texto{473.1943}{474.5580}{11.9552}{\normalfont}{para}
\texto{499.3882}{474.5580}{11.9552}{\normalfont}{turistas.}
\end{pagina}
% Página original 16
\begin{pagina}{16}
\texto{70.8660}{746.3890}{11.9552}{\normalfont}{6}
\texto{266.5910}{746.3890}{11.9552}{\normalfont}{CAPÍTULO}
\texto{333.5545}{746.3890}{11.9552}{\normalfont}{1.}
\texto{350.4663}{746.3890}{11.9552}{\normalfont}{FORMULACIÓN}
\texto{445.2256}{746.3890}{11.9552}{\normalfont}{DEL}
\texto{473.3418}{746.3890}{11.9552}{\normalfont}{PROYECTO}
\texto{70.8660}{705.1940}{17.2154}{\normalfont\bfseries}{1.4.}
\texto{120.3603}{705.1940}{17.2154}{\normalfont\bfseries}{Planteamiento}
\texto{248.6305}{705.1940}{17.2154}{\normalfont\bfseries}{del}
\texto{280.0624}{705.1940}{17.2154}{\normalfont\bfseries}{problema}
\texto{88.7990}{678.9090}{11.9552}{\normalfont}{£Cómo}
\texto{131.9621}{678.9090}{11.9552}{\normalfont}{pueden}
\texto{172.8536}{678.9090}{11.9552}{\normalfont}{las}
\texto{191.0470}{678.9090}{11.9552}{\normalfont}{aplicaciones}
\texto{256.0666}{678.9090}{11.9552}{\normalfont}{de}
\texto{272.2563}{678.9090}{11.9552}{\normalfont}{escritorio}
\texto{323.6183}{678.9090}{11.9552}{\normalfont}{del}
\texto{343.0478}{678.9090}{11.9552}{\normalfont}{ecosistema}
\texto{401.7780}{678.9090}{11.9552}{\normalfont\bfseries}{Explora-Nicaragua}
\texto{516.4210}{678.9090}{11.9552}{\normalfont}{---un}
\texto{70.8660}{664.4630}{11.9552}{\normalfont}{backoffice}
\texto{124.5759}{664.4630}{11.9552}{\normalfont}{para}
\texto{151.9533}{664.4630}{11.9552}{\normalfont}{operadores}
\texto{211.5859}{664.4630}{11.9552}{\normalfont}{turísticos}
\texto{263.1558}{664.4630}{11.9552}{\normalfont}{y}
\texto{273.9370}{664.4630}{11.9552}{\normalfont}{un}
\texto{291.5590}{664.4630}{11.9552}{\normalfont}{panel}
\texto{323.4889}{664.4630}{11.9552}{\normalfont}{de}
\texto{339.8102}{664.4630}{11.9552}{\normalfont}{administración}
\texto{419.9315}{664.4630}{11.9552}{\normalfont}{para}
\texto{447.3089}{664.4630}{11.9552}{\normalfont}{el}
\texto{460.3664}{664.4630}{11.9552}{\normalfont}{equipo}
\texto{498.7976}{664.4630}{11.9552}{\normalfont}{gestor---}
\texto{70.8660}{650.0170}{11.9552}{\normalfont}{centralizar}
\texto{127.7034}{650.0170}{11.9552}{\normalfont}{la}
\texto{139.9886}{650.0170}{11.9552}{\normalfont}{promoción}
\texto{196.8212}{650.0170}{11.9552}{\normalfont}{y}
\texto{206.1797}{650.0170}{11.9552}{\normalfont}{reserva}
\texto{244.8739}{650.0170}{11.9552}{\normalfont}{de}
\texto{259.7606}{650.0170}{11.9552}{\normalfont}{tours}
\texto{289.0197}{650.0170}{11.9552}{\normalfont}{turísticos,}
\texto{342.4188}{650.0170}{11.9552}{\normalfont}{facilitar}
\texto{384.9458}{650.0170}{11.9552}{\normalfont}{la}
\texto{397.2310}{650.0170}{11.9552}{\normalfont}{gestión}
\texto{436.2456}{650.0170}{11.9552}{\normalfont}{eficiente}
\texto{480.7261}{650.0170}{11.9552}{\normalfont}{del}
\texto{498.8646}{650.0170}{11.9552}{\normalfont}{catálogo}
\texto{70.8660}{635.5710}{11.9552}{\normalfont}{y}
\texto{81.4081}{635.5710}{11.9552}{\normalfont}{la}
\texto{94.8888}{635.5710}{11.9552}{\normalfont}{disponibilidad}
\texto{171.1916}{635.5710}{11.9552}{\normalfont}{por}
\texto{192.7995}{635.5710}{11.9552}{\normalfont}{parte}
\texto{223.8280}{635.5710}{11.9552}{\normalfont}{de}
\texto{239.9101}{635.5710}{11.9552}{\normalfont}{los}
\texto{258.0079}{635.5710}{11.9552}{\normalfont}{operadores,}
\texto{320.6412}{635.5710}{11.9552}{\normalfont}{y}
\texto{331.1953}{635.5710}{11.9552}{\normalfont}{permitir}
\texto{377.5169}{635.5710}{11.9552}{\normalfont}{la}
\texto{390.9976}{635.5710}{11.9552}{\normalfont}{supervisión}
\texto{452.7199}{635.5710}{11.9552}{\normalfont}{de}
\texto{468.8020}{635.5710}{11.9552}{\normalfont}{la}
\texto{482.2827}{635.5710}{11.9552}{\normalfont}{plataforma,}
\texto{70.8660}{621.1260}{11.9552}{\normalfont}{reduciendo}
\texto{130.0442}{621.1260}{11.9552}{\normalfont}{al}
\texto{143.0467}{621.1260}{11.9552}{\normalfont}{mismo}
\texto{180.1891}{621.1260}{11.9552}{\normalfont}{tiempo}
\texto{219.5289}{621.1260}{11.9552}{\normalfont}{las}
\texto{237.1604}{621.1260}{11.9552}{\normalfont}{comisiones}
\texto{295.1671}{621.1260}{11.9552}{\normalfont}{asociadas}
\texto{347.0048}{621.1260}{11.9552}{\normalfont}{al}
\texto{360.0073}{621.1260}{11.9552}{\normalfont}{uso}
\texto{380.8787}{621.1260}{11.9552}{\normalfont}{de}
\texto{396.4946}{621.1260}{11.9552}{\normalfont}{plataformas}
\texto{460.6151}{621.1260}{11.9552}{\normalfont}{extranjeras?}
\texto{88.7990}{606.6800}{11.9552}{\normalfont}{Este}
\texto{114.6151}{606.6800}{11.9552}{\normalfont}{planteamiento}
\texto{190.2771}{606.6800}{11.9552}{\normalfont}{integra}
\texto{229.2033}{606.6800}{11.9552}{\normalfont}{tres}
\texto{251.6073}{606.6800}{11.9552}{\normalfont}{dimensiones}
\texto{316.3495}{606.6800}{11.9552}{\normalfont}{del}
\texto{334.7868}{606.6800}{11.9552}{\normalfont}{problema.}
\texto{389.0060}{606.6800}{11.9552}{\normalfont}{La}
\texto{405.6524}{606.6800}{11.9552}{\normalfont}{primera}
\texto{448.8035}{606.6800}{11.9552}{\normalfont}{es}
\texto{462.1025}{606.6800}{11.9552}{\normalfont}{la}
\texto{474.8270}{606.6800}{11.9552}{\normalfont\bfseries}{experiencia}
\texto{70.8660}{592.2340}{11.9552}{\normalfont\bfseries}{del}
\texto{91.9884}{592.2340}{11.9552}{\normalfont\bfseries}{turista}
\texto{131.2590}{592.2340}{11.9552}{\normalfont}{:}
\texto{137.7985}{592.2340}{11.9552}{\normalfont}{descubrir,}
\texto{191.2287}{592.2340}{11.9552}{\normalfont}{comparar}
\texto{242.6552}{592.2340}{11.9552}{\normalfont}{y}
\texto{252.1213}{592.2340}{11.9552}{\normalfont}{reservar}
\texto{295.4756}{592.2340}{11.9552}{\normalfont}{tours}
\texto{324.8424}{592.2340}{11.9552}{\normalfont}{con}
\texto{345.6899}{592.2340}{11.9552}{\normalfont}{un}
\texto{361.9848}{592.2340}{11.9552}{\normalfont}{pago}
\texto{389.3479}{592.2340}{11.9552}{\normalfont}{en}
\texto{404.3421}{592.2340}{11.9552}{\normalfont}{línea}
\texto{431.6932}{592.2340}{11.9552}{\normalfont}{seguro}
\texto{467.5755}{592.2340}{11.9552}{\normalfont}{y}
\texto{477.0416}{592.2340}{11.9552}{\normalfont}{con}
\texto{497.8891}{592.2340}{11.9552}{\normalfont}{informa-}
\texto{70.8660}{577.7880}{11.9552}{\normalfont}{ción}
\texto{95.9337}{577.7880}{11.9552}{\normalfont}{confiable}
\texto{145.0648}{577.7880}{11.9552}{\normalfont}{---dimensión}
\texto{212.4706}{577.7880}{11.9552}{\normalfont}{que}
\texto{234.6116}{577.7880}{11.9552}{\normalfont}{se}
\texto{248.6876}{577.7880}{11.9552}{\normalfont}{contempla}
\texto{305.3003}{577.7880}{11.9552}{\normalfont}{en}
\texto{321.2629}{577.7880}{11.9552}{\normalfont}{la}
\texto{334.6240}{577.7880}{11.9552}{\normalfont}{visión}
\texto{368.5481}{577.7880}{11.9552}{\normalfont}{completa}
\texto{418.9799}{577.7880}{11.9552}{\normalfont}{del}
\texto{438.1943}{577.7880}{11.9552}{\normalfont}{ecosistema}
\texto{496.5596}{577.7880}{11.9552}{\normalfont}{pero}
\texto{523.2508}{577.7880}{11.9552}{\normalfont}{que}
\texto{70.8660}{563.3420}{11.9552}{\normalfont}{no}
\texto{87.4670}{563.3420}{11.9552}{\normalfont}{forma}
\texto{121.3026}{563.3420}{11.9552}{\normalfont}{parte}
\texto{152.2116}{563.3420}{11.9552}{\normalfont}{del}
\texto{171.4021}{563.3420}{11.9552}{\normalfont}{prototipo}
\texto{223.4574}{563.3420}{11.9552}{\normalfont}{inicial}
\texto{258.2685}{563.3420}{11.9552}{\normalfont}{de}
\texto{274.2192}{563.3420}{11.9552}{\normalfont}{escritorio.}
\texto{328.6058}{563.3420}{11.9552}{\normalfont}{La}
\texto{346.0173}{563.3420}{11.9552}{\normalfont}{segunda}
\texto{391.2988}{563.3420}{11.9552}{\normalfont}{es}
\texto{405.3629}{563.3420}{11.9552}{\normalfont}{la}
\texto{418.8530}{563.3420}{11.9552}{\normalfont\bfseries}{gestión}
\texto{465.0491}{563.3420}{11.9552}{\normalfont\bfseries}{operativa}
\texto{519.9160}{563.3420}{11.9552}{\normalfont}{:}
\texto{527.4119}{563.3420}{11.9552}{\normalfont}{los}
\texto{70.8660}{548.8960}{11.9552}{\normalfont}{operadores}
\texto{129.9366}{548.8960}{11.9552}{\normalfont}{requieren}
\texto{181.1527}{548.8960}{11.9552}{\normalfont}{herramientas}
\texto{251.2843}{548.8960}{11.9552}{\normalfont}{de}
\texto{267.0556}{548.8960}{11.9552}{\normalfont}{escritorio}
\texto{317.9991}{548.8960}{11.9552}{\normalfont}{eficientes}
\texto{367.9695}{548.8960}{11.9552}{\normalfont}{para}
\texto{394.7969}{548.8960}{11.9552}{\normalfont}{administrar}
\texto{458.0973}{548.8960}{11.9552}{\normalfont}{su}
\texto{473.2828}{548.8960}{11.9552}{\normalfont}{catálogo,}
\texto{522.8610}{548.8960}{11.9552}{\normalfont}{dis-}
\texto{70.8660}{534.4510}{11.9552}{\normalfont}{ponibilidad}
\texto{132.4855}{534.4510}{11.9552}{\normalfont}{y}
\texto{142.7407}{534.4510}{11.9552}{\normalfont}{reservas,}
\texto{190.1885}{534.4510}{11.9552}{\normalfont}{superando}
\texto{246.6816}{534.4510}{11.9552}{\normalfont}{la}
\texto{259.8634}{534.4510}{11.9552}{\normalfont}{gestión}
\texto{299.7746}{534.4510}{11.9552}{\normalfont}{manual}
\texto{341.2376}{534.4510}{11.9552}{\normalfont}{informal.}
\texto{391.1649}{534.4510}{11.9552}{\normalfont}{La}
\texto{408.3972}{534.4510}{11.9552}{\normalfont}{tercera}
\texto{447.5935}{534.4510}{11.9552}{\normalfont}{es}
\texto{461.4902}{534.4510}{11.9552}{\normalfont}{la}
\texto{474.8090}{534.4510}{11.9552}{\normalfont\bfseries}{supervisión}
\texto{70.8660}{520.0050}{11.9552}{\normalfont\bfseries}{de}
\texto{89.5436}{520.0050}{11.9552}{\normalfont\bfseries}{la}
\texto{104.8989}{520.0050}{11.9552}{\normalfont\bfseries}{plataforma}
\texto{168.7590}{520.0050}{11.9552}{\normalfont}{:}
\texto{176.4223}{520.0050}{11.9552}{\normalfont}{el}
\texto{189.3004}{520.0050}{11.9552}{\normalfont}{equipo}
\texto{227.5403}{520.0050}{11.9552}{\normalfont}{gestor}
\texto{262.5954}{520.0050}{11.9552}{\normalfont}{necesita}
\texto{307.4058}{520.0050}{11.9552}{\normalfont}{un}
\texto{324.8365}{520.0050}{11.9552}{\normalfont}{panel}
\texto{356.5752}{520.0050}{11.9552}{\normalfont}{de}
\texto{372.7051}{520.0050}{11.9552}{\normalfont}{escritorio}
\texto{424.0073}{520.0050}{11.9552}{\normalfont}{para}
\texto{451.1934}{520.0050}{11.9552}{\normalfont}{verificar}
\texto{496.5921}{520.0050}{11.9552}{\normalfont}{operado-}
\texto{70.8660}{505.5590}{11.9552}{\normalfont}{res,}
\texto{92.9855}{505.5590}{11.9552}{\normalfont}{moderar}
\texto{140.0770}{505.5590}{11.9552}{\normalfont}{la}
\texto{153.6892}{505.5590}{11.9552}{\normalfont}{oferta}
\texto{187.7759}{505.5590}{11.9552}{\normalfont}{y}
\texto{198.4615}{505.5590}{11.9552}{\normalfont}{consultar}
\texto{249.8473}{505.5590}{11.9552}{\normalfont}{reportes.}
\texto{298.9664}{505.5590}{11.9552}{\normalfont}{Las}
\texto{321.2461}{505.5590}{11.9552}{\normalfont}{dos}
\texto{342.7153}{505.5590}{11.9552}{\normalfont}{últimas}
\texto{385.0080}{505.5590}{11.9552}{\normalfont}{dimensiones}
\texto{450.7664}{505.5590}{11.9552}{\normalfont}{se}
\texto{465.0935}{505.5590}{11.9552}{\normalfont}{resuelven}
\texto{516.4817}{505.5590}{11.9552}{\normalfont}{en}
\texto{532.6834}{505.5590}{11.9552}{\normalfont}{el}
\texto{70.8660}{491.1130}{11.9552}{\normalfont}{prototipo}
\texto{122.4073}{491.1130}{11.9552}{\normalfont}{de}
\texto{137.8558}{491.1130}{11.9552}{\normalfont}{esta}
\texto{161.8116}{491.1130}{11.9552}{\normalfont}{etapa}
\texto{193.5192}{491.1130}{11.9552}{\normalfont}{mediante}
\texto{243.7645}{491.1130}{11.9552}{\normalfont}{dos}
\texto{264.4805}{491.1130}{11.9552}{\normalfont}{aplicaciones}
\texto{328.7707}{491.1130}{11.9552}{\normalfont}{de}
\texto{344.2193}{491.1130}{11.9552}{\normalfont}{escritorio}
\texto{394.8519}{491.1130}{11.9552}{\normalfont}{que}
\texto{416.4789}{491.1130}{11.9552}{\normalfont}{consumen}
\texto{470.3514}{491.1130}{11.9552}{\normalfont}{la}
\texto{483.1985}{491.1130}{11.9552}{\normalfont}{misma}
\texto{520.1735}{491.1130}{11.9552}{\normalfont}{API}
\texto{70.8660}{476.6670}{11.9552}{\normalfont}{REST,}
\texto{109.5530}{476.6670}{11.9552}{\normalfont}{garantizando}
\texto{179.4646}{476.6670}{11.9552}{\normalfont}{coherencia}
\texto{236.6918}{476.6670}{11.9552}{\normalfont}{de}
\texto{252.3077}{476.6670}{11.9552}{\normalfont}{datos}
\texto{283.5849}{476.6670}{11.9552}{\normalfont}{y}
\texto{293.6727}{476.6670}{11.9552}{\normalfont}{control}
\texto{333.0172}{476.6670}{11.9552}{\normalfont}{de}
\texto{348.6212}{476.6670}{11.9552}{\normalfont}{acceso}
\texto{384.4629}{476.6670}{11.9552}{\normalfont}{basado}
\texto{423.5444}{476.6670}{11.9552}{\normalfont}{en}
\texto{439.1483}{476.6670}{11.9552}{\normalfont}{roles}
\texto{466.5353}{476.6670}{11.9552}{\normalfont}{(RBAC).}
\end{pagina}
% Página original 17
\begin{pagina}{17}
\texto{70.8660}{746.3890}{11.9552}{\normalfont}{1.5.}
\texto{96.8710}{746.3890}{11.9552}{\normalfont}{JUSTIFICACIÓN}
\texto{535.2810}{746.3890}{11.9552}{\normalfont}{7}
\texto{70.8660}{705.1940}{17.2154}{\normalfont\bfseries}{1.5.}
\texto{120.3603}{705.1940}{17.2154}{\normalfont\bfseries}{Justificación}
\texto{88.7990}{678.9090}{11.9552}{\normalfont}{El}
\texto{104.8692}{678.9090}{11.9552}{\normalfont}{turismo}
\texto{148.8093}{678.9090}{11.9552}{\normalfont}{es}
\texto{163.4831}{678.9090}{11.9552}{\normalfont}{una}
\texto{187.2094}{678.9090}{11.9552}{\normalfont}{fuente}
\texto{223.2830}{678.9090}{11.9552}{\normalfont}{estratégica}
\texto{282.8319}{678.9090}{11.9552}{\normalfont}{de}
\texto{299.3922}{678.9090}{11.9552}{\normalfont}{divisas}
\texto{338.5312}{678.9090}{11.9552}{\normalfont}{para}
\texto{366.1477}{678.9090}{11.9552}{\normalfont}{Nicaragua:}
\texto{425.9524}{678.9090}{11.9552}{\normalfont}{el}
\texto{439.2729}{678.9090}{11.9552}{\normalfont}{ingreso}
\texto{479.9612}{678.9090}{11.9552}{\normalfont}{por}
\texto{502.0472}{678.9090}{11.9552}{\normalfont}{turismo}
\texto{70.8660}{664.4630}{11.9552}{\normalfont}{representó}
\texto{125.9651}{664.4630}{11.9552}{\normalfont}{aproximadamente}
\texto{219.3854}{664.4630}{11.9552}{\normalfont}{el}
\texto{230.5301}{664.4630}{11.9552}{\normalfont}{5.2}
\texto{247.4730}{664.4630}{11.9552}{\normalfont}{\%}
\texto{259.9184}{664.4630}{11.9552}{\normalfont}{del}
\texto{277.5666}{664.4630}{11.9552}{\normalfont}{PIB}
\texto{300.7334}{664.4630}{11.9552}{\normalfont}{en}
\texto{315.1299}{664.4630}{11.9552}{\normalfont}{2023}
\texto{341.2328}{664.4630}{11.9552}{\normalfont}{[2].}
\texto{359.5275}{664.4630}{11.9552}{\normalfont}{Sin}
\texto{378.4765}{664.4630}{11.9552}{\normalfont}{embargo,}
\texto{427.6698}{664.4630}{11.9552}{\normalfont}{la}
\texto{439.4648}{664.4630}{11.9552}{\normalfont}{escasa}
\texto{473.5013}{664.4630}{11.9552}{\normalfont}{digitalización}
\texto{70.8660}{650.0170}{11.9552}{\normalfont}{de}
\texto{85.5972}{650.0170}{11.9552}{\normalfont}{la}
\texto{97.7389}{650.0170}{11.9552}{\normalfont}{oferta}
\texto{130.3670}{650.0170}{11.9552}{\normalfont}{limita}
\texto{163.3203}{650.0170}{11.9552}{\normalfont}{la}
\texto{175.4501}{650.0170}{11.9552}{\normalfont}{capacidad}
\texto{229.2150}{650.0170}{11.9552}{\normalfont}{del}
\texto{247.2100}{650.0170}{11.9552}{\normalfont}{país}
\texto{270.4724}{650.0170}{11.9552}{\normalfont}{para}
\texto{296.2717}{650.0170}{11.9552}{\normalfont}{captar}
\texto{331.8145}{650.0170}{11.9552}{\normalfont}{y}
\texto{341.0296}{650.0170}{11.9552}{\normalfont}{retener}
\texto{379.8362}{650.0170}{11.9552}{\normalfont}{valor}
\texto{407.9046}{650.0170}{11.9552}{\normalfont}{en}
\texto{422.6478}{650.0170}{11.9552}{\normalfont}{la}
\texto{434.7895}{650.0170}{11.9552}{\normalfont}{cadena}
\texto{472.9457}{650.0170}{11.9552}{\normalfont}{de}
\texto{487.6888}{650.0170}{11.9552}{\normalfont}{suministro}
\texto{70.8660}{635.5710}{11.9552}{\normalfont}{turística.}
\texto{120.3534}{635.5710}{11.9552}{\normalfont}{Un}
\texto{139.5295}{635.5710}{11.9552}{\normalfont}{ecosistema}
\texto{197.5481}{635.5710}{11.9552}{\normalfont}{como}
\texto{228.1130}{635.5710}{11.9552}{\normalfont\bfseries}{Explora-Nicaragua}
\texto{342.1840}{635.5710}{11.9552}{\normalfont}{permitiría:}
\texto{100.7540}{608.1740}{11.9552}{\normalfont\bfseries}{Digitalizar}
\texto{166.9571}{608.1740}{11.9552}{\normalfont\bfseries}{la}
\texto{181.6907}{608.1740}{11.9552}{\normalfont\bfseries}{gestión}
\texto{227.4803}{608.1740}{11.9552}{\normalfont\bfseries}{de}
\texto{245.5482}{608.1740}{11.9552}{\normalfont\bfseries}{los}
\texto{265.7740}{608.1740}{11.9552}{\normalfont\bfseries}{operadores:}
\texto{337.2590}{608.1740}{11.9552}{\normalfont}{al}
\texto{350.2376}{608.1740}{11.9552}{\normalfont}{brindar}
\texto{391.8321}{608.1740}{11.9552}{\normalfont}{un}
\texto{408.7248}{608.1740}{11.9552}{\normalfont}{backoffice}
\texto{461.7055}{608.1740}{11.9552}{\normalfont}{de}
\texto{477.2974}{608.1740}{11.9552}{\normalfont}{escritorio,}
\texto{531.3134}{608.1740}{11.9552}{\normalfont}{se}
\texto{100.7540}{593.7280}{11.9552}{\normalfont}{automatizan}
\texto{168.6380}{593.7280}{11.9552}{\normalfont}{procesos}
\texto{215.5120}{593.7280}{11.9552}{\normalfont}{que}
\texto{237.5454}{593.7280}{11.9552}{\normalfont}{hoy}
\texto{259.9064}{593.7280}{11.9552}{\normalfont}{se}
\texto{273.8749}{593.7280}{11.9552}{\normalfont}{realizan}
\texto{317.6955}{593.7280}{11.9552}{\normalfont}{de}
\texto{333.5504}{593.7280}{11.9552}{\normalfont}{forma}
\texto{367.2904}{593.7280}{11.9552}{\normalfont}{manual}
\texto{408.8371}{593.7280}{11.9552}{\normalfont}{(registro}
\texto{455.9741}{593.7280}{11.9552}{\normalfont}{de}
\texto{471.8291}{593.7280}{11.9552}{\normalfont}{tours,}
\texto{505.3084}{593.7280}{11.9552}{\normalfont}{gestión}
\texto{100.7540}{579.2820}{11.9552}{\normalfont}{de}
\texto{117.0035}{579.2820}{11.9552}{\normalfont}{disponibilidad,}
\texto{196.7256}{579.2820}{11.9552}{\normalfont}{control}
\texto{236.7157}{579.2820}{11.9552}{\normalfont}{de}
\texto{252.9652}{579.2820}{11.9552}{\normalfont}{reservas),}
\texto{305.4438}{579.2820}{11.9552}{\normalfont}{reduciendo}
\texto{365.2676}{579.2820}{11.9552}{\normalfont}{errores}
\texto{404.3443}{579.2820}{11.9552}{\normalfont}{como}
\texto{435.5522}{579.2820}{11.9552}{\normalfont}{las}
\texto{453.8174}{579.2820}{11.9552}{\normalfont}{dobles}
\texto{490.2927}{579.2820}{11.9552}{\normalfont}{reservas}
\texto{534.9549}{579.2820}{11.9552}{\normalfont}{y}
\texto{100.7540}{564.8370}{11.9552}{\normalfont}{mejorando}
\texto{158.3063}{564.8370}{11.9552}{\normalfont}{la}
\texto{171.3088}{564.8370}{11.9552}{\normalfont}{eficiencia}
\texto{221.3939}{564.8370}{11.9552}{\normalfont}{operativa.}
\texto{100.7540}{540.4280}{11.9552}{\normalfont\bfseries}{Supervisar}
\texto{167.4796}{540.4280}{11.9552}{\normalfont\bfseries}{la}
\texto{181.8904}{540.4280}{11.9552}{\normalfont\bfseries}{plataforma}
\texto{249.6166}{540.4280}{11.9552}{\normalfont\bfseries}{desde}
\texto{286.2760}{540.4280}{11.9552}{\normalfont\bfseries}{una}
\texto{311.8948}{540.4280}{11.9552}{\normalfont\bfseries}{aplicación}
\texto{373.9399}{540.4280}{11.9552}{\normalfont\bfseries}{de}
\texto{391.6850}{540.4280}{11.9552}{\normalfont\bfseries}{escritorio:}
\texto{453.5890}{540.4280}{11.9552}{\normalfont}{el}
\texto{465.6422}{540.4280}{11.9552}{\normalfont}{panel}
\texto{496.5560}{540.4280}{11.9552}{\normalfont}{de}
\texto{511.8610}{540.4280}{11.9552}{\normalfont}{admi-}
\texto{100.7540}{525.9820}{11.9552}{\normalfont}{nistración}
\texto{154.3635}{525.9820}{11.9552}{\normalfont}{de}
\texto{169.5490}{525.9820}{11.9552}{\normalfont}{escritorio}
\texto{219.9067}{525.9820}{11.9552}{\normalfont}{permite}
\texto{262.7302}{525.9820}{11.9552}{\normalfont}{al}
\texto{275.3143}{525.9820}{11.9552}{\normalfont}{equipo}
\texto{312.5978}{525.9820}{11.9552}{\normalfont}{gestor}
\texto{346.7083}{525.9820}{11.9552}{\normalfont}{verificar}
\texto{391.1506}{525.9820}{11.9552}{\normalfont}{operadores,}
\texto{452.8992}{525.9820}{11.9552}{\normalfont}{moderar}
\texto{498.9745}{525.9820}{11.9552}{\normalfont}{la}
\texto{511.5466}{525.9820}{11.9552}{\normalfont}{oferta}
\texto{100.7540}{511.5360}{11.9552}{\normalfont}{y}
\texto{110.7701}{511.5360}{11.9552}{\normalfont}{consultar}
\texto{161.4984}{511.5360}{11.9552}{\normalfont}{reportes}
\texto{206.6962}{511.5360}{11.9552}{\normalfont}{sin}
\texto{224.9064}{511.5360}{11.9552}{\normalfont}{depender}
\texto{275.2545}{511.5360}{11.9552}{\normalfont}{de}
\texto{290.7986}{511.5360}{11.9552}{\normalfont}{un}
\texto{307.6435}{511.5360}{11.9552}{\normalfont}{navegador}
\texto{363.1898}{511.5360}{11.9552}{\normalfont}{web,}
\texto{390.1177}{511.5360}{11.9552}{\normalfont}{ofreciendo}
\texto{445.6591}{511.5360}{11.9552}{\normalfont}{una}
\texto{468.3573}{511.5360}{11.9552}{\normalfont}{experiencia}
\texto{529.4364}{511.5360}{11.9552}{\normalfont}{de}
\texto{100.7540}{497.0910}{11.9552}{\normalfont}{gestión}
\texto{140.4859}{497.0910}{11.9552}{\normalfont}{más}
\texto{164.6091}{497.0910}{11.9552}{\normalfont}{fluida}
\texto{197.1344}{497.0910}{11.9552}{\normalfont}{y}
\texto{207.2103}{497.0910}{11.9552}{\normalfont}{productiva.}
\texto{100.7540}{472.6820}{11.9552}{\normalfont\bfseries}{Generar}
\texto{153.6534}{472.6820}{11.9552}{\normalfont\bfseries}{divisas:}
\texto{200.9990}{472.6820}{11.9552}{\normalfont}{al}
\texto{214.5036}{472.6820}{11.9552}{\normalfont}{conectar}
\texto{261.8390}{472.6820}{11.9552}{\normalfont}{a}
\texto{272.1037}{472.6820}{11.9552}{\normalfont}{turistas}
\texto{315.0157}{472.6820}{11.9552}{\normalfont}{extranjeros}
\texto{376.0614}{472.6820}{11.9552}{\normalfont}{con}
\texto{398.0327}{472.6820}{11.9552}{\normalfont}{operadores}
\texto{457.4620}{472.6820}{11.9552}{\normalfont}{locales}
\texto{495.4293}{472.6820}{11.9552}{\normalfont}{de}
\texto{511.5473}{472.6820}{11.9552}{\normalfont}{forma}
\texto{100.7540}{458.2360}{11.9552}{\normalfont}{directa,}
\texto{143.0228}{458.2360}{11.9552}{\normalfont}{se}
\texto{156.7402}{458.2360}{11.9552}{\normalfont}{favorece}
\texto{201.6152}{458.2360}{11.9552}{\normalfont}{la}
\texto{214.6297}{458.2360}{11.9552}{\normalfont}{entrada}
\texto{257.2260}{458.2360}{11.9552}{\normalfont}{de}
\texto{272.8300}{458.2360}{11.9552}{\normalfont}{divisas}
\texto{311.0125}{458.2360}{11.9552}{\normalfont}{al}
\texto{324.0150}{458.2360}{11.9552}{\normalfont}{país.}
\texto{100.7540}{433.8280}{11.9552}{\normalfont\bfseries}{Crear}
\texto{139.5689}{433.8280}{11.9552}{\normalfont\bfseries}{empleo:}
\texto{189.4740}{433.8280}{11.9552}{\normalfont}{la}
\texto{203.3133}{433.8280}{11.9552}{\normalfont}{formalización}
\texto{276.0225}{433.8280}{11.9552}{\normalfont}{y}
\texto{286.9471}{433.8280}{11.9552}{\normalfont}{digitalización}
\texto{359.3191}{433.8280}{11.9552}{\normalfont}{de}
\texto{375.7719}{433.8280}{11.9552}{\normalfont}{la}
\texto{389.6232}{433.8280}{11.9552}{\normalfont}{oferta}
\texto{423.9489}{433.8280}{11.9552}{\normalfont}{turística}
\texto{471.0333}{433.8280}{11.9552}{\normalfont}{incrementa}
\texto{532.0383}{433.8280}{11.9552}{\normalfont}{la}
\texto{100.7540}{419.3820}{11.9552}{\normalfont}{demanda}
\texto{150.8272}{419.3820}{11.9552}{\normalfont}{de}
\texto{166.4311}{419.3820}{11.9552}{\normalfont}{servicios}
\texto{213.0683}{419.3820}{11.9552}{\normalfont}{y}
\texto{223.1442}{419.3820}{11.9552}{\normalfont}{de}
\texto{238.7481}{419.3820}{11.9552}{\normalfont}{profesionales}
\texto{307.4977}{419.3820}{11.9552}{\normalfont}{del}
\texto{326.3534}{419.3820}{11.9552}{\normalfont}{sector.}
\texto{100.7540}{394.9740}{11.9552}{\normalfont\bfseries}{F}
\texto{108.0969}{394.9740}{11.9552}{\normalfont\bfseries}{omentar}
\texto{161.4314}{394.9740}{11.9552}{\normalfont\bfseries}{la}
\texto{176.6193}{394.9740}{11.9552}{\normalfont\bfseries}{confianza:}
\texto{238.8080}{394.9740}{11.9552}{\normalfont}{un}
\texto{256.0833}{394.9740}{11.9552}{\normalfont}{sistema}
\texto{298.2134}{394.9740}{11.9552}{\normalfont}{de}
\texto{314.1999}{394.9740}{11.9552}{\normalfont}{reseñas}
\texto{355.0293}{394.9740}{11.9552}{\normalfont}{verificadas}
\texto{412.6940}{394.9740}{11.9552}{\normalfont}{y}
\texto{423.1524}{394.9740}{11.9552}{\normalfont}{pagos}
\texto{456.1129}{394.9740}{11.9552}{\normalfont}{en}
\texto{472.0994}{394.9740}{11.9552}{\normalfont}{línea}
\texto{500.4428}{394.9740}{11.9552}{\normalfont}{con}
\texto{522.2706}{394.9740}{11.9552}{\normalfont}{cer-}
\texto{100.7540}{380.5280}{11.9552}{\normalfont}{tificación}
\texto{151.1500}{380.5280}{11.9552}{\normalfont}{SSL/TLS,}
\texto{207.0692}{380.5280}{11.9552}{\normalfont}{sumado}
\texto{250.3877}{380.5280}{11.9552}{\normalfont}{a}
\texto{260.4611}{380.5280}{11.9552}{\normalfont}{la}
\texto{273.7864}{380.5280}{11.9552}{\normalfont}{verificación}
\texto{335.2529}{380.5280}{11.9552}{\normalfont}{documental}
\texto{398.6561}{380.5280}{11.9552}{\normalfont}{de}
\texto{414.5828}{380.5280}{11.9552}{\normalfont}{operadores}
\texto{473.8328}{380.5280}{11.9552}{\normalfont}{por}
\texto{495.2971}{380.5280}{11.9552}{\normalfont}{parte}
\texto{526.1822}{380.5280}{11.9552}{\normalfont}{del}
\texto{100.7540}{366.0820}{11.9552}{\normalfont}{panel}
\texto{131.9666}{366.0820}{11.9552}{\normalfont}{de}
\texto{147.5706}{366.0820}{11.9552}{\normalfont}{administración,}
\texto{230.2384}{366.0820}{11.9552}{\normalfont}{mejora}
\texto{268.9302}{366.0820}{11.9552}{\normalfont}{la}
\texto{281.9327}{366.0820}{11.9552}{\normalfont}{percepción}
\texto{340.4677}{366.0820}{11.9552}{\normalfont}{de}
\texto{356.0836}{366.0820}{11.9552}{\normalfont}{seguridad}
\texto{408.8228}{366.0820}{11.9552}{\normalfont}{del}
\texto{427.6785}{366.0820}{11.9552}{\normalfont}{turista.}
\texto{100.7540}{341.6730}{11.9552}{\normalfont\bfseries}{Reducción}
\texto{165.0658}{341.6730}{11.9552}{\normalfont\bfseries}{de}
\texto{181.9860}{341.6730}{11.9552}{\normalfont\bfseries}{comisiones:}
\texto{250.9360}{341.6730}{11.9552}{\normalfont}{al}
\texto{262.9223}{341.6730}{11.9552}{\normalfont}{operar}
\texto{298.6444}{341.6730}{11.9552}{\normalfont}{como}
\texto{328.1905}{341.6730}{11.9552}{\normalfont}{plataforma}
\texto{386.6777}{341.6730}{11.9552}{\normalfont}{nicaragüense}
\texto{454.6598}{341.6730}{11.9552}{\normalfont}{con}
\texto{475.1008}{341.6730}{11.9552}{\normalfont}{una}
\texto{496.8425}{341.6730}{11.9552}{\normalfont}{comisión}
\texto{100.7540}{327.2280}{11.9552}{\normalfont}{menor}
\texto{136.1964}{327.2280}{11.9552}{\normalfont}{al}
\texto{148.8880}{327.2280}{11.9552}{\normalfont}{10}
\texto{162.5791}{327.2280}{11.9552}{\normalfont}{\%,}
\texto{179.1729}{327.2280}{11.9552}{\normalfont}{se}
\texto{192.5675}{327.2280}{11.9552}{\normalfont}{incrementa}
\texto{252.4129}{327.2280}{11.9552}{\normalfont}{el}
\texto{264.4422}{327.2280}{11.9552}{\normalfont}{margen}
\texto{305.7498}{327.2280}{11.9552}{\normalfont}{de}
\texto{321.0429}{327.2280}{11.9552}{\normalfont}{los}
\texto{338.3397}{327.2280}{11.9552}{\normalfont}{proveedores}
\texto{401.4918}{327.2280}{11.9552}{\normalfont}{locales}
\texto{438.6343}{327.2280}{11.9552}{\normalfont}{frente}
\texto{471.4776}{327.2280}{11.9552}{\normalfont}{a}
\texto{480.9174}{327.2280}{11.9552}{\normalfont}{plataformas}
\texto{100.7540}{312.7820}{11.9552}{\normalfont}{extranjeras}
\texto{161.2975}{312.7820}{11.9552}{\normalfont}{[6].}
\texto{100.7540}{288.3730}{11.9552}{\normalfont\bfseries}{Aporte}
\texto{148.3751}{288.3730}{11.9552}{\normalfont\bfseries}{tecnológico}
\texto{219.5899}{288.3730}{11.9552}{\normalfont\bfseries}{y}
\texto{232.7263}{288.3730}{11.9552}{\normalfont\bfseries}{académico:}
\texto{302.0440}{288.3730}{11.9552}{\normalfont}{el}
\texto{315.7590}{288.3730}{11.9552}{\normalfont}{proyecto}
\texto{364.2613}{288.3730}{11.9552}{\normalfont}{integra}
\texto{404.9687}{288.3730}{11.9552}{\normalfont}{arquitecturas}
\texto{477.6061}{288.3730}{11.9552}{\normalfont}{orientadas}
\texto{535.2876}{288.3730}{11.9552}{\normalfont}{a}
\texto{100.7540}{273.9280}{11.9552}{\normalfont}{servicios,}
\texto{150.9419}{273.9280}{11.9552}{\normalfont}{aplicaciones}
\texto{215.6985}{273.9280}{11.9552}{\normalfont}{de}
\texto{231.6252}{273.9280}{11.9552}{\normalfont}{escritorio}
\texto{282.7241}{273.9280}{11.9552}{\normalfont}{en}
\texto{298.6508}{273.9280}{11.9552}{\normalfont}{.NET,}
\texto{334.5571}{273.9280}{11.9552}{\normalfont}{control}
\texto{374.2125}{273.9280}{11.9552}{\normalfont}{de}
\texto{390.1392}{273.9280}{11.9552}{\normalfont}{acceso}
\texto{426.2797}{273.9280}{11.9552}{\normalfont}{basado}
\texto{465.6721}{273.9280}{11.9552}{\normalfont}{en}
\texto{481.5989}{273.9280}{11.9552}{\normalfont}{roles}
\texto{509.2847}{273.9280}{11.9552}{\normalfont}{y}
\texto{519.6833}{273.9280}{11.9552}{\normalfont}{con-}
\texto{100.7540}{259.4820}{11.9552}{\normalfont}{sumo}
\texto{131.6797}{259.4820}{11.9552}{\normalfont}{de}
\texto{147.5945}{259.4820}{11.9552}{\normalfont}{API}
\texto{172.7650}{259.4820}{11.9552}{\normalfont}{REST,}
\texto{211.7628}{259.4820}{11.9552}{\normalfont}{constituyendo}
\texto{286.9275}{259.4820}{11.9552}{\normalfont}{un}
\texto{304.1430}{259.4820}{11.9552}{\normalfont}{caso}
\texto{329.8778}{259.4820}{11.9552}{\normalfont}{de}
\texto{345.7926}{259.4820}{11.9552}{\normalfont}{aplicación}
\texto{400.7291}{259.4820}{11.9552}{\normalfont}{relevante}
\texto{450.4579}{259.4820}{11.9552}{\normalfont}{para}
\texto{477.4289}{259.4820}{11.9552}{\normalfont}{la}
\texto{490.7422}{259.4820}{11.9552}{\normalfont}{Ingeniería}
\texto{100.7540}{245.0360}{11.9552}{\normalfont}{en}
\texto{116.3579}{245.0360}{11.9552}{\normalfont}{Sistemas.}
\texto{88.7990}{217.6390}{11.9552}{\normalfont}{La}
\texto{107.0713}{217.6390}{11.9552}{\normalfont}{relevancia}
\texto{162.5841}{217.6390}{11.9552}{\normalfont}{social,}
\texto{199.2937}{217.6390}{11.9552}{\normalfont}{económica}
\texto{257.0780}{217.6390}{11.9552}{\normalfont}{y}
\texto{268.3733}{217.6390}{11.9552}{\normalfont}{tecnológica}
\texto{330.0597}{217.6390}{11.9552}{\normalfont}{del}
\texto{350.1229}{217.6390}{11.9552}{\normalfont}{proyecto}
\texto{398.4817}{217.6390}{11.9552}{\normalfont}{justifica}
\texto{443.6485}{217.6390}{11.9552}{\normalfont}{su}
\texto{459.8741}{217.6390}{11.9552}{\normalfont}{desarrollo}
\texto{514.4711}{217.6390}{11.9552}{\normalfont}{como}
\texto{70.8660}{203.1930}{11.9552}{\normalfont}{aporte}
\texto{107.6043}{203.1930}{11.9552}{\normalfont}{al}
\texto{120.6068}{203.1930}{11.9552}{\normalfont}{proceso}
\texto{162.6246}{203.1930}{11.9552}{\normalfont}{de}
\texto{178.2285}{203.1930}{11.9552}{\normalfont}{digitalización}
\texto{249.7756}{203.1930}{11.9552}{\normalfont}{del}
\texto{268.6313}{203.1930}{11.9552}{\normalfont}{turismo}
\texto{311.6270}{203.1930}{11.9552}{\normalfont}{nicaragüense.}
\end{pagina}
% Página original 18
\begin{pagina}{18}
\texto{70.8660}{746.3890}{11.9552}{\normalfont}{8}
\texto{266.5910}{746.3890}{11.9552}{\normalfont}{CAPÍTULO}
\texto{333.5545}{746.3890}{11.9552}{\normalfont}{1.}
\texto{350.4663}{746.3890}{11.9552}{\normalfont}{FORMULACIÓN}
\texto{445.2256}{746.3890}{11.9552}{\normalfont}{DEL}
\texto{473.3418}{746.3890}{11.9552}{\normalfont}{PROYECTO}
\texto{70.8660}{705.1940}{17.2154}{\normalfont\bfseries}{1.6.}
\texto{120.3603}{705.1940}{17.2154}{\normalfont\bfseries}{Objetivos}
\texto{70.8660}{675.4220}{14.3462}{\normalfont\bfseries}{1.6.1.}
\texto{124.6642}{675.4220}{14.3462}{\normalfont\bfseries}{Objetivo}
\texto{190.7758}{675.4220}{14.3462}{\normalfont\bfseries}{general}
\texto{88.7990}{653.2550}{11.9552}{\normalfont}{Desarrollar}
\texto{148.8308}{653.2550}{11.9552}{\normalfont}{el}
\texto{160.8482}{653.2550}{11.9552}{\normalfont}{prototipo}
\texto{212.2101}{653.2550}{11.9552}{\normalfont}{funcional}
\texto{262.2618}{653.2550}{11.9552}{\normalfont}{de}
\texto{277.5310}{653.2550}{11.9552}{\normalfont}{las}
\texto{294.8158}{653.2550}{11.9552}{\normalfont}{aplicaciones}
\texto{358.9148}{653.2550}{11.9552}{\normalfont}{de}
\texto{374.1840}{653.2550}{11.9552}{\normalfont}{escritorio}
\texto{424.6373}{653.2550}{11.9552}{\normalfont}{del}
\texto{443.1463}{653.2550}{11.9552}{\normalfont}{ecosistema}
\texto{500.8182}{653.2550}{11.9552}{\normalfont}{tecnoló-}
\texto{70.8660}{638.8090}{11.9552}{\normalfont}{gico}
\texto{93.9563}{638.8090}{11.9552}{\normalfont}{denominado}
\texto{158.6880}{638.8090}{11.9552}{\normalfont\bfseries}{Explora-Nicaragua}
\texto{268.5950}{638.8090}{11.9552}{\normalfont}{,}
\texto{274.7758}{638.8090}{11.9552}{\normalfont}{orientado}
\texto{325.5209}{638.8090}{11.9552}{\normalfont}{a}
\texto{334.3151}{638.8090}{11.9552}{\normalfont}{la}
\texto{346.3612}{638.8090}{11.9552}{\normalfont}{promoción}
\texto{402.9547}{638.8090}{11.9552}{\normalfont}{y}
\texto{412.0741}{638.8090}{11.9552}{\normalfont}{reserva}
\texto{450.5173}{638.8090}{11.9552}{\normalfont}{de}
\texto{465.1648}{638.8090}{11.9552}{\normalfont}{tours}
\texto{494.1848}{638.8090}{11.9552}{\normalfont}{turísticos}
\texto{70.8660}{624.3630}{11.9552}{\normalfont}{en}
\texto{86.0874}{624.3630}{11.9552}{\normalfont}{Nicaragua.}
\texto{144.5650}{624.3630}{11.9552}{\normalfont}{El}
\texto{159.2962}{624.3630}{11.9552}{\normalfont}{prototipo}
\texto{210.6223}{624.3630}{11.9552}{\normalfont}{comprende}
\texto{269.4299}{624.3630}{11.9552}{\normalfont}{el}
\texto{281.3995}{624.3630}{11.9552}{\normalfont}{backoffice}
\texto{334.0215}{624.3630}{11.9552}{\normalfont}{de}
\texto{349.2548}{624.3630}{11.9552}{\normalfont}{escritorio}
\texto{399.6603}{624.3630}{11.9552}{\normalfont}{para}
\texto{425.9498}{624.3630}{11.9552}{\normalfont}{operadores}
\texto{484.4824}{624.3630}{11.9552}{\normalfont}{turísticos}
\texto{534.9645}{624.3630}{11.9552}{\normalfont}{y}
\texto{70.8660}{609.9180}{11.9552}{\normalfont}{el}
\texto{83.2898}{609.9180}{11.9552}{\normalfont}{panel}
\texto{114.5742}{609.9180}{11.9552}{\normalfont}{de}
\texto{130.2499}{609.9180}{11.9552}{\normalfont}{administración}
\texto{209.7256}{609.9180}{11.9552}{\normalfont}{de}
\texto{225.4013}{609.9180}{11.9552}{\normalfont}{escritorio}
\texto{276.2611}{609.9180}{11.9552}{\normalfont}{para}
\texto{303.0049}{609.9180}{11.9552}{\normalfont}{el}
\texto{315.4287}{609.9180}{11.9552}{\normalfont}{equipo}
\texto{353.2143}{609.9180}{11.9552}{\normalfont}{gestor,}
\texto{391.0669}{609.9180}{11.9552}{\normalfont}{ambos}
\texto{427.6187}{609.9180}{11.9552}{\normalfont}{consumiendo}
\texto{497.3391}{609.9180}{11.9552}{\normalfont}{una}
\texto{520.1807}{609.9180}{11.9552}{\normalfont}{API}
\texto{70.8660}{595.4720}{11.9552}{\normalfont}{REST}
\texto{105.5839}{595.4720}{11.9552}{\normalfont}{común,}
\texto{145.5238}{595.4720}{11.9552}{\normalfont}{con}
\texto{166.2757}{595.4720}{11.9552}{\normalfont}{el}
\texto{177.9224}{595.4720}{11.9552}{\normalfont}{fin}
\texto{194.1217}{595.4720}{11.9552}{\normalfont}{de}
\texto{209.0203}{595.4720}{11.9552}{\normalfont}{centralizar}
\texto{265.8697}{595.4720}{11.9552}{\normalfont}{la}
\texto{278.1668}{595.4720}{11.9552}{\normalfont}{gestión}
\texto{317.1933}{595.4720}{11.9552}{\normalfont}{del}
\texto{335.3437}{595.4720}{11.9552}{\normalfont}{catálogo,}
\texto{384.0612}{595.4720}{11.9552}{\normalfont}{la}
\texto{396.3583}{595.4720}{11.9552}{\normalfont}{disponibilidad}
\texto{471.4776}{595.4720}{11.9552}{\normalfont}{y}
\texto{480.8481}{595.4720}{11.9552}{\normalfont}{las}
\texto{497.7623}{595.4720}{11.9552}{\normalfont}{reservas.}
\texto{70.8660}{560.8100}{14.3462}{\normalfont\bfseries}{1.6.2.}
\texto{124.6642}{560.8100}{14.3462}{\normalfont\bfseries}{Objetivos}
\texto{197.1427}{560.8100}{14.3462}{\normalfont\bfseries}{específicos}
\texto{100.7540}{538.6430}{11.9552}{\normalfont\bfseries}{Investigar}
\texto{164.4870}{538.6430}{11.9552}{\normalfont}{la}
\texto{179.3186}{538.6430}{11.9552}{\normalfont}{oferta}
\texto{214.6486}{538.6430}{11.9552}{\normalfont}{turística}
\texto{262.7253}{538.6430}{11.9552}{\normalfont}{nicaragüense}
\texto{333.5646}{538.6430}{11.9552}{\normalfont}{y}
\texto{345.4696}{538.6430}{11.9552}{\normalfont}{las}
\texto{364.9303}{538.6430}{11.9552}{\normalfont}{tendencias}
\texto{424.0631}{538.6430}{11.9552}{\normalfont}{tecnológicas}
\texto{491.0002}{538.6430}{11.9552}{\normalfont}{relevantes}
\texto{100.7540}{524.1970}{11.9552}{\normalfont}{mediante}
\texto{151.3341}{524.1970}{11.9552}{\normalfont}{una}
\texto{174.2833}{524.1970}{11.9552}{\normalfont}{investigación}
\texto{243.8028}{524.1970}{11.9552}{\normalfont}{documental}
\texto{307.0745}{524.1970}{11.9552}{\normalfont}{que}
\texto{329.0481}{524.1970}{11.9552}{\normalfont}{sustente}
\texto{374.5663}{524.1970}{11.9552}{\normalfont}{el}
\texto{387.1097}{524.1970}{11.9552}{\normalfont}{desarrollo}
\texto{440.6905}{524.1970}{11.9552}{\normalfont}{del}
\texto{459.7376}{524.1970}{11.9552}{\normalfont}{prototipo}
\texto{511.6255}{524.1970}{11.9552}{\normalfont}{de}
\texto{527.4207}{524.1970}{11.9552}{\normalfont}{las}
\texto{100.7540}{509.7520}{11.9552}{\normalfont}{aplicaciones}
\texto{165.1997}{509.7520}{11.9552}{\normalfont}{de}
\texto{180.8036}{509.7520}{11.9552}{\normalfont}{escritorio.}
\texto{100.7540}{485.3430}{11.9552}{\normalfont\bfseries}{F}
\texto{108.0969}{485.3430}{11.9552}{\normalfont\bfseries}{ormular}
\texto{157.7960}{485.3430}{11.9552}{\normalfont}{el}
\texto{169.6579}{485.3430}{11.9552}{\normalfont}{emprendimiento}
\texto{255.3504}{485.3430}{11.9552}{\normalfont}{digital}
\texto{291.3300}{485.3430}{11.9552}{\normalfont\bfseries}{Explora-Nicaragua}
\texto{401.2370}{485.3430}{11.9552}{\normalfont}{,}
\texto{407.8960}{485.3430}{11.9552}{\normalfont}{definiendo}
\texto{463.3443}{485.3430}{11.9552}{\normalfont}{la}
\texto{475.8685}{485.3430}{11.9552}{\normalfont}{propuesta}
\texto{529.4302}{485.3430}{11.9552}{\normalfont}{de}
\texto{100.7540}{470.8970}{11.9552}{\normalfont}{valor,}
\texto{133.0904}{470.8970}{11.9552}{\normalfont}{el}
\texto{145.6099}{470.8970}{11.9552}{\normalfont}{público}
\texto{186.7334}{470.8970}{11.9552}{\normalfont}{objetivo,}
\texto{235.0109}{470.8970}{11.9552}{\normalfont}{el}
\texto{247.5304}{470.8970}{11.9552}{\normalfont}{modelo}
\texto{288.3263}{470.8970}{11.9552}{\normalfont}{de}
\texto{304.0857}{470.8970}{11.9552}{\normalfont}{funcionamiento,}
\texto{389.4482}{470.8970}{11.9552}{\normalfont}{la}
\texto{402.6061}{470.8970}{11.9552}{\normalfont}{identidad}
\texto{454.4629}{470.8970}{11.9552}{\normalfont}{visual,}
\texto{491.4236}{470.8970}{11.9552}{\normalfont}{el}
\texto{503.9431}{470.8970}{11.9552}{\normalfont}{análisis}
\texto{100.7540}{456.4510}{11.9552}{\normalfont}{FODA}
\texto{138.4535}{456.4510}{11.9552}{\normalfont}{y}
\texto{148.5413}{456.4510}{11.9552}{\normalfont}{el}
\texto{160.8934}{456.4510}{11.9552}{\normalfont}{Lean}
\texto{189.6768}{456.4510}{11.9552}{\normalfont}{Canvas.}
\texto{100.7540}{432.0430}{11.9552}{\normalfont\bfseries}{Modelar}
\texto{153.1240}{432.0430}{11.9552}{\normalfont}{el}
\texto{164.7110}{432.0430}{11.9552}{\normalfont}{proceso}
\texto{205.9516}{432.0430}{11.9552}{\normalfont}{de}
\texto{220.8024}{432.0430}{11.9552}{\normalfont}{promoción,}
\texto{280.8390}{432.0430}{11.9552}{\normalfont}{reserva}
\texto{319.4854}{432.0430}{11.9552}{\normalfont}{y}
\texto{328.8081}{432.0430}{11.9552}{\normalfont}{gestión}
\texto{367.7748}{432.0430}{11.9552}{\normalfont}{de}
\texto{382.6256}{432.0430}{11.9552}{\normalfont}{tours}
\texto{411.8369}{432.0430}{11.9552}{\normalfont}{en}
\texto{426.6757}{432.0430}{11.9552}{\normalfont}{su}
\texto{440.9407}{432.0430}{11.9552}{\normalfont}{estado}
\texto{476.6556}{432.0430}{11.9552}{\normalfont}{actual}
\texto{511.0173}{432.0430}{11.9552}{\normalfont}{y}
\texto{520.3280}{432.0430}{11.9552}{\normalfont}{pro-}
\texto{100.7540}{417.5970}{11.9552}{\normalfont}{puesto}
\texto{137.0476}{417.5970}{11.9552}{\normalfont}{mediante}
\texto{186.5995}{417.5970}{11.9552}{\normalfont}{la}
\texto{198.7771}{417.5970}{11.9552}{\normalfont}{notación}
\texto{245.4119}{417.5970}{11.9552}{\normalfont}{BPMN,}
\texto{287.4799}{417.5970}{11.9552}{\normalfont}{identificando}
\texto{355.9043}{417.5970}{11.9552}{\normalfont}{las}
\texto{372.6870}{417.5970}{11.9552}{\normalfont}{oportunidades}
\texto{448.3251}{417.5970}{11.9552}{\normalfont}{de}
\texto{463.0922}{417.5970}{11.9552}{\normalfont}{automatización}
\texto{100.7540}{403.1510}{11.9552}{\normalfont}{que}
\texto{122.5364}{403.1510}{11.9552}{\normalfont}{serán}
\texto{153.1632}{403.1510}{11.9552}{\normalfont}{cubiertas}
\texto{203.3129}{403.1510}{11.9552}{\normalfont}{por}
\texto{224.4425}{403.1510}{11.9552}{\normalfont}{el}
\texto{236.8066}{403.1510}{11.9552}{\normalfont}{prototipo.}
\texto{100.7540}{378.7430}{11.9552}{\normalfont\bfseries}{Diseñar}
\texto{151.0930}{378.7430}{11.9552}{\normalfont}{la}
\texto{165.5421}{378.7430}{11.9552}{\normalfont}{arquitectura}
\texto{233.6460}{378.7430}{11.9552}{\normalfont}{del}
\texto{253.9603}{378.7430}{11.9552}{\normalfont}{sistema,}
\texto{300.4063}{378.7430}{11.9552}{\normalfont}{definiendo}
\texto{357.7912}{378.7430}{11.9552}{\normalfont}{los}
\texto{376.8574}{378.7430}{11.9552}{\normalfont}{módulos}
\texto{424.8743}{378.7430}{11.9552}{\normalfont}{de}
\texto{441.9248}{378.7430}{11.9552}{\normalfont}{turista,}
\texto{484.4040}{378.7430}{11.9552}{\normalfont}{operador}
\texto{534.9577}{378.7430}{11.9552}{\normalfont}{y}
\texto{100.7540}{364.2970}{11.9552}{\normalfont}{administrador,}
\texto{181.3249}{364.2970}{11.9552}{\normalfont}{el}
\texto{195.4942}{364.2970}{11.9552}{\normalfont}{modelo}
\texto{237.9519}{364.2970}{11.9552}{\normalfont}{de}
\texto{255.3730}{364.2970}{11.9552}{\normalfont}{datos}
\texto{288.4674}{364.2970}{11.9552}{\normalfont}{y}
\texto{300.3604}{364.2970}{11.9552}{\normalfont}{las}
\texto{319.7972}{364.2970}{11.9552}{\normalfont}{interfaces}
\texto{373.7056}{364.2970}{11.9552}{\normalfont}{de}
\texto{391.1267}{364.2970}{11.9552}{\normalfont}{las}
\texto{410.5634}{364.2970}{11.9552}{\normalfont}{aplicaciones}
\texto{476.8263}{364.2970}{11.9552}{\normalfont}{de}
\texto{494.2475}{364.2970}{11.9552}{\normalfont}{escritorio}
\texto{100.7540}{349.8510}{11.9552}{\normalfont}{(backoffice}
\texto{158.2992}{349.8510}{11.9552}{\normalfont}{de}
\texto{173.9150}{349.8510}{11.9552}{\normalfont}{operadores}
\texto{232.8303}{349.8510}{11.9552}{\normalfont}{y}
\texto{242.9181}{349.8510}{11.9552}{\normalfont}{panel}
\texto{274.1307}{349.8510}{11.9552}{\normalfont}{de}
\texto{289.7466}{349.8510}{11.9552}{\normalfont}{administración).}
\texto{100.7540}{325.4430}{11.9552}{\normalfont\bfseries}{Desarrollar}
\texto{170.9540}{325.4430}{11.9552}{\normalfont}{el}
\texto{184.0952}{325.4430}{11.9552}{\normalfont}{prototipo}
\texto{236.5928}{325.4430}{11.9552}{\normalfont}{funcional}
\texto{287.7922}{325.4430}{11.9552}{\normalfont}{de}
\texto{304.1971}{325.4430}{11.9552}{\normalfont}{las}
\texto{322.6210}{325.4430}{11.9552}{\normalfont\bfseries}{aplicaciones}
\texto{397.3637}{325.4430}{11.9552}{\normalfont\bfseries}{de}
\texto{416.3641}{325.4430}{11.9552}{\normalfont\bfseries}{escritorio}
\texto{470.9320}{325.4430}{11.9552}{\normalfont}{:}
\texto{478.8822}{325.4430}{11.9552}{\normalfont}{el}
\texto{492.0234}{325.4430}{11.9552}{\normalfont}{backoffice}
\texto{100.7540}{310.9970}{11.9552}{\normalfont}{para}
\texto{127.4500}{310.9970}{11.9552}{\normalfont}{operadores,}
\texto{189.6529}{310.9970}{11.9552}{\normalfont}{permitiendo}
\texto{255.0550}{310.9970}{11.9552}{\normalfont}{la}
\texto{268.0933}{310.9970}{11.9552}{\normalfont}{gestión}
\texto{307.8731}{310.9970}{11.9552}{\normalfont}{de}
\texto{323.5129}{310.9970}{11.9552}{\normalfont}{tours,}
\texto{356.7890}{310.9970}{11.9552}{\normalfont}{disponibilidad,}
\texto{435.9013}{310.9970}{11.9552}{\normalfont}{reservas}
\texto{479.9658}{310.9970}{11.9552}{\normalfont}{y}
\texto{490.0894}{310.9970}{11.9552}{\normalfont}{validación}
\texto{100.7540}{296.5510}{11.9552}{\normalfont}{de}
\texto{116.3938}{296.5510}{11.9552}{\normalfont}{vouchers;}
\texto{167.5477}{296.5510}{11.9552}{\normalfont}{y}
\texto{177.6714}{296.5510}{11.9552}{\normalfont}{el}
\texto{190.0593}{296.5510}{11.9552}{\normalfont}{panel}
\texto{221.3198}{296.5510}{11.9552}{\normalfont}{de}
\texto{236.9596}{296.5510}{11.9552}{\normalfont}{administración}
\texto{316.4115}{296.5510}{11.9552}{\normalfont}{para}
\texto{343.1074}{296.5510}{11.9552}{\normalfont}{el}
\texto{355.5074}{296.5510}{11.9552}{\normalfont}{equipo}
\texto{393.2691}{296.5510}{11.9552}{\normalfont}{gestor,}
\texto{431.0857}{296.5510}{11.9552}{\normalfont}{permitiendo}
\texto{496.4879}{296.5510}{11.9552}{\normalfont}{la}
\texto{509.5262}{296.5510}{11.9552}{\normalfont}{super-}
\texto{100.7540}{282.1050}{11.9552}{\normalfont}{visión}
\texto{134.7737}{282.1050}{11.9552}{\normalfont}{de}
\texto{150.8439}{282.1050}{11.9552}{\normalfont}{la}
\texto{164.3126}{282.1050}{11.9552}{\normalfont}{plataforma,}
\texto{227.5341}{282.1050}{11.9552}{\normalfont}{la}
\texto{241.0028}{282.1050}{11.9552}{\normalfont}{verificación}
\texto{302.6008}{282.1050}{11.9552}{\normalfont}{de}
\texto{318.6710}{282.1050}{11.9552}{\normalfont}{operadores,}
\texto{381.3043}{282.1050}{11.9552}{\normalfont}{la}
\texto{394.7730}{282.1050}{11.9552}{\normalfont}{moderación}
\texto{457.9921}{282.1050}{11.9552}{\normalfont}{de}
\texto{474.0623}{282.1050}{11.9552}{\normalfont}{la}
\texto{487.5310}{282.1050}{11.9552}{\normalfont}{oferta}
\texto{521.4981}{282.1050}{11.9552}{\normalfont}{y}
\texto{532.0402}{282.1050}{11.9552}{\normalfont}{la}
\texto{100.7540}{267.6590}{11.9552}{\normalfont}{consulta}
\texto{146.9895}{267.6590}{11.9552}{\normalfont}{de}
\texto{162.5935}{267.6590}{11.9552}{\normalfont}{reportes.}
\texto{100.7540}{243.2510}{11.9552}{\normalfont\bfseries}{V}
\texto{109.7981}{243.2510}{11.9552}{\normalfont\bfseries}{alidar}
\texto{148.3830}{243.2510}{11.9552}{\normalfont}{el}
\texto{161.9187}{243.2510}{11.9552}{\normalfont}{prototipo}
\texto{214.7989}{243.2510}{11.9552}{\normalfont}{mediante}
\texto{266.3832}{243.2510}{11.9552}{\normalfont}{pruebas}
\texto{311.2009}{243.2510}{11.9552}{\normalfont}{funcionales}
\texto{372.6028}{243.2510}{11.9552}{\normalfont}{de}
\texto{389.3903}{243.2510}{11.9552}{\normalfont}{los}
\texto{408.1934}{243.2510}{11.9552}{\normalfont}{procesos}
\texto{455.9999}{243.2510}{11.9552}{\normalfont}{críticos,}
\texto{500.8175}{243.2510}{11.9552}{\normalfont}{verifica-}
\texto{100.7540}{228.8050}{11.9552}{\normalfont}{ción}
\texto{126.1445}{228.8050}{11.9552}{\normalfont}{del}
\texto{145.6936}{228.8050}{11.9552}{\normalfont}{control}
\texto{185.7196}{228.8050}{11.9552}{\normalfont}{de}
\texto{202.0050}{228.8050}{11.9552}{\normalfont}{acceso}
\texto{238.5281}{228.8050}{11.9552}{\normalfont}{basado}
\texto{278.2911}{228.8050}{11.9552}{\normalfont}{en}
\texto{294.5884}{228.8050}{11.9552}{\normalfont}{roles}
\texto{322.6569}{228.8050}{11.9552}{\normalfont}{y}
\texto{333.4141}{228.8050}{11.9552}{\normalfont}{evaluación}
\texto{391.0023}{228.8050}{11.9552}{\normalfont}{de}
\texto{407.2997}{228.8050}{11.9552}{\normalfont}{usabilidad}
\texto{463.9721}{228.8050}{11.9552}{\normalfont}{con}
\texto{486.1227}{228.8050}{11.9552}{\normalfont}{operadores}
\texto{100.7540}{214.3590}{11.9552}{\normalfont}{turísticos}
\texto{151.6066}{214.3590}{11.9552}{\normalfont}{piloto}
\texto{184.7704}{214.3590}{11.9552}{\normalfont}{en}
\texto{200.3862}{214.3590}{11.9552}{\normalfont}{los}
\texto{218.0058}{214.3590}{11.9552}{\normalfont}{destinos}
\texto{263.0052}{214.3590}{11.9552}{\normalfont}{de}
\texto{278.6211}{214.3590}{11.9552}{\normalfont}{Ometepe,}
\texto{331.6185}{214.3590}{11.9552}{\normalfont}{Somoto}
\texto{373.8992}{214.3590}{11.9552}{\normalfont}{y}
\texto{383.9751}{214.3590}{11.9552}{\normalfont}{San}
\texto{406.7450}{214.3590}{11.9552}{\normalfont}{Juan}
\texto{435.5163}{214.3590}{11.9552}{\normalfont}{del}
\texto{454.3721}{214.3590}{11.9552}{\normalfont}{Sur.}
\end{pagina}
% Página original 19
\begin{pagina}{19}
\texto{70.8660}{746.3890}{11.9552}{\normalfont}{1.7.}
\texto{96.8710}{746.3890}{11.9552}{\normalfont}{ALCANCE}
\texto{159.2938}{746.3890}{11.9552}{\normalfont}{DEL}
\texto{187.4101}{746.3890}{11.9552}{\normalfont}{PROYECTO}
\texto{535.2810}{746.3890}{11.9552}{\normalfont}{9}
\texto{70.8660}{705.1940}{17.2154}{\normalfont\bfseries}{1.7.}
\texto{120.3603}{705.1940}{17.2154}{\normalfont\bfseries}{Alcance}
\texto{193.0471}{705.1940}{17.2154}{\normalfont\bfseries}{del}
\texto{224.4790}{705.1940}{17.2154}{\normalfont\bfseries}{proyecto}
\texto{88.7990}{678.9090}{11.9552}{\normalfont}{El}
\texto{103.5541}{678.9090}{11.9552}{\normalfont}{proyecto}
\texto{150.3560}{678.9090}{11.9552}{\normalfont\bfseries}{Explora-Nicaragua}
\texto{264.0710}{678.9090}{11.9552}{\normalfont}{contempla}
\texto{319.9663}{678.9090}{11.9552}{\normalfont}{la}
\texto{332.6221}{678.9090}{11.9552}{\normalfont}{visión}
\texto{365.8169}{678.9090}{11.9552}{\normalfont}{de}
\texto{381.0742}{678.9090}{11.9552}{\normalfont}{un}
\texto{397.6321}{678.9090}{11.9552}{\normalfont}{ecosistema}
\texto{455.2801}{678.9090}{11.9552}{\normalfont}{de}
\texto{470.5373}{678.9090}{11.9552}{\normalfont}{tres}
\texto{493.0131}{678.9090}{11.9552}{\normalfont}{aplicacio-}
\texto{70.8660}{664.4630}{11.9552}{\normalfont}{nes}
\texto{91.1588}{664.4630}{11.9552}{\normalfont}{coordinadas}
\texto{156.0109}{664.4630}{11.9552}{\normalfont}{para}
\texto{182.7547}{664.4630}{11.9552}{\normalfont}{la}
\texto{195.8409}{664.4630}{11.9552}{\normalfont}{promoción}
\texto{253.4745}{664.4630}{11.9552}{\normalfont}{y}
\texto{263.6221}{664.4630}{11.9552}{\normalfont}{reserva}
\texto{303.1173}{664.4630}{11.9552}{\normalfont}{de}
\texto{318.7929}{664.4630}{11.9552}{\normalfont}{tours}
\texto{348.8531}{664.4630}{11.9552}{\normalfont}{turísticos}
\texto{399.7894}{664.4630}{11.9552}{\normalfont}{en}
\texto{415.4770}{664.4630}{11.9552}{\normalfont}{Nicaragua.}
\texto{474.3970}{664.4630}{11.9552}{\normalfont}{Sin}
\texto{494.6372}{664.4630}{11.9552}{\normalfont}{embargo,}
\texto{70.8660}{650.0170}{11.9552}{\normalfont\bfseries}{el}
\texto{85.7957}{650.0170}{11.9552}{\normalfont\bfseries}{alcance}
\texto{133.2291}{650.0170}{11.9552}{\normalfont\bfseries}{de}
\texto{151.8948}{650.0170}{11.9552}{\normalfont\bfseries}{este}
\texto{179.7611}{650.0170}{11.9552}{\normalfont\bfseries}{prototipo}
\texto{239.9974}{650.0170}{11.9552}{\normalfont\bfseries}{inicial}
\texto{279.9864}{650.0170}{11.9552}{\normalfont\bfseries}{se}
\texto{296.4857}{650.0170}{11.9552}{\normalfont\bfseries}{limita}
\texto{335.7275}{650.0170}{11.9552}{\normalfont\bfseries}{a}
\texto{347.3228}{650.0170}{11.9552}{\normalfont\bfseries}{las}
\texto{367.9599}{650.0170}{11.9552}{\normalfont\bfseries}{aplicaciones}
\texto{442.3679}{650.0170}{11.9552}{\normalfont\bfseries}{de}
\texto{461.0335}{650.0170}{11.9552}{\normalfont\bfseries}{escritorio}
\texto{515.5960}{650.0170}{11.9552}{\normalfont}{,}
\texto{523.2473}{650.0170}{11.9552}{\normalfont}{que}
\texto{70.8660}{635.5710}{11.9552}{\normalfont}{constituyen}
\texto{133.3630}{635.5710}{11.9552}{\normalfont}{el}
\texto{145.7151}{635.5710}{11.9552}{\normalfont}{foco}
\texto{170.4337}{635.5710}{11.9552}{\normalfont}{del}
\texto{189.2894}{635.5710}{11.9552}{\normalfont}{desarrollo}
\texto{242.6790}{635.5710}{11.9552}{\normalfont}{en}
\texto{258.2948}{635.5710}{11.9552}{\normalfont}{esta}
\texto{282.4180}{635.5710}{11.9552}{\normalfont}{etapa.}
\texto{70.8660}{600.9100}{14.3462}{\normalfont\bfseries}{Alcance}
\texto{131.4385}{600.9100}{14.3462}{\normalfont\bfseries}{del}
\texto{157.6318}{600.9100}{14.3462}{\normalfont\bfseries}{prototipo}
\texto{229.2265}{600.9100}{14.3462}{\normalfont\bfseries}{(aplicaciones}
\texto{324.1037}{600.9100}{14.3462}{\normalfont\bfseries}{de}
\texto{345.8138}{600.9100}{14.3462}{\normalfont\bfseries}{escritorio)}
\texto{70.8660}{578.7430}{11.9552}{\normalfont\bfseries}{Backoffice}
\texto{133.8197}{578.7430}{11.9552}{\normalfont\bfseries}{de}
\texto{151.9115}{578.7430}{11.9552}{\normalfont\bfseries}{escritorio}
\texto{210.9534}{578.7430}{11.9552}{\normalfont\bfseries}{(operador)}
\texto{88.7990}{556.5760}{11.9552}{\normalfont}{Registro}
\texto{136.7011}{556.5760}{11.9552}{\normalfont}{y}
\texto{148.2952}{556.5760}{11.9552}{\normalfont}{actualización}
\texto{220.0480}{556.5760}{11.9552}{\normalfont}{de}
\texto{237.1702}{556.5760}{11.9552}{\normalfont}{tours}
\texto{268.6650}{556.5760}{11.9552}{\normalfont}{con}
\texto{291.6405}{556.5760}{11.9552}{\normalfont}{descripciones,}
\texto{367.0969}{556.5760}{11.9552}{\normalfont}{fotografías,}
\texto{429.1588}{556.5760}{11.9552}{\normalfont}{precios}
\texto{469.7586}{556.5760}{11.9552}{\normalfont}{y}
\texto{481.3528}{556.5760}{11.9552}{\normalfont}{políticas}
\texto{529.4294}{556.5760}{11.9552}{\normalfont}{de}
\texto{70.8660}{542.1300}{11.9552}{\normalfont}{cancelación;}
\texto{136.4235}{542.1300}{11.9552}{\normalfont}{gestión}
\texto{176.6815}{542.1300}{11.9552}{\normalfont}{de}
\texto{192.8234}{542.1300}{11.9552}{\normalfont}{disponibilidad}
\texto{269.1860}{542.1300}{11.9552}{\normalfont}{(fechas,}
\texto{312.0478}{542.1300}{11.9552}{\normalfont}{horarios}
\texto{357.5087}{542.1300}{11.9552}{\normalfont}{y}
\texto{368.1225}{542.1300}{11.9552}{\normalfont}{cupos);}
\texto{409.3536}{542.1300}{11.9552}{\normalfont}{visualización}
\texto{478.5646}{542.1300}{11.9552}{\normalfont}{y}
\texto{489.1665}{542.1300}{11.9552}{\normalfont}{gestión}
\texto{529.4364}{542.1300}{11.9552}{\normalfont}{de}
\texto{70.8660}{527.6840}{11.9552}{\normalfont}{reservas;}
\texto{119.6647}{527.6840}{11.9552}{\normalfont}{validación}
\texto{176.1387}{527.6840}{11.9552}{\normalfont}{de}
\texto{193.2729}{527.6840}{11.9552}{\normalfont}{vouchers;}
\texto{245.9093}{527.6840}{11.9552}{\normalfont}{consulta}
\texto{293.6751}{527.6840}{11.9552}{\normalfont}{de}
\texto{310.8093}{527.6840}{11.9552}{\normalfont}{reportes}
\texto{357.5971}{527.6840}{11.9552}{\normalfont}{financieros}
\texto{417.0575}{527.6840}{11.9552}{\normalfont}{(bruto,}
\texto{458.2551}{527.6840}{11.9552}{\normalfont}{comisión,}
\texto{511.2238}{527.6840}{11.9552}{\normalfont}{neto);}
\texto{70.8660}{513.2380}{11.9552}{\normalfont}{gestión}
\texto{110.5979}{513.2380}{11.9552}{\normalfont}{del}
\texto{129.4537}{513.2380}{11.9552}{\normalfont}{equipo}
\texto{167.1795}{513.2380}{11.9552}{\normalfont}{de}
\texto{182.7834}{513.2380}{11.9552}{\normalfont}{la}
\texto{195.7978}{513.2380}{11.9552}{\normalfont}{agencia}
\texto{237.4163}{513.2380}{11.9552}{\normalfont}{(rol}
\texto{259.5238}{513.2380}{11.9552}{\normalfont}{propietario)}
\texto{323.9169}{513.2380}{11.9552}{\normalfont}{y}
\texto{333.9928}{513.2380}{11.9552}{\normalfont}{gestión}
\texto{373.7366}{513.2380}{11.9552}{\normalfont}{de}
\texto{389.3406}{513.2380}{11.9552}{\normalfont}{la}
\texto{402.3430}{513.2380}{11.9552}{\normalfont}{suscripción}
\texto{462.6355}{513.2380}{11.9552}{\normalfont}{y}
\texto{472.7114}{513.2380}{11.9552}{\normalfont}{del}
\texto{491.5791}{513.2380}{11.9552}{\normalfont}{perfil.}
\texto{70.8660}{482.0640}{11.9552}{\normalfont\bfseries}{Panel}
\texto{108.0371}{482.0640}{11.9552}{\normalfont\bfseries}{de}
\texto{126.1289}{482.0640}{11.9552}{\normalfont\bfseries}{administración}
\texto{217.2526}{482.0640}{11.9552}{\normalfont\bfseries}{de}
\texto{235.3444}{482.0640}{11.9552}{\normalfont\bfseries}{escritorio}
\texto{294.3864}{482.0640}{11.9552}{\normalfont\bfseries}{(equipo}
\texto{343.1230}{482.0640}{11.9552}{\normalfont\bfseries}{gestor)}
\texto{88.7990}{459.8970}{11.9552}{\normalfont}{Dashboard}
\texto{146.8248}{459.8970}{11.9552}{\normalfont}{con}
\texto{167.2418}{459.8970}{11.9552}{\normalfont}{KPIs}
\texto{195.9917}{459.8970}{11.9552}{\normalfont}{de}
\texto{210.5555}{459.8970}{11.9552}{\normalfont}{la}
\texto{222.5179}{459.8970}{11.9552}{\normalfont}{plataforma;}
\texto{284.2330}{459.8970}{11.9552}{\normalfont}{verificación}
\texto{344.3127}{459.8970}{11.9552}{\normalfont}{y}
\texto{353.3484}{459.8970}{11.9552}{\normalfont}{gestión}
\texto{392.0283}{459.8970}{11.9552}{\normalfont}{de}
\texto{406.5921}{459.8970}{11.9552}{\normalfont}{operadores;}
\texto{467.7191}{459.8970}{11.9552}{\normalfont}{moderación}
\texto{529.4318}{459.8970}{11.9552}{\normalfont}{de}
\texto{70.8660}{445.4510}{11.9552}{\normalfont}{tours}
\texto{100.7827}{445.4510}{11.9552}{\normalfont}{y}
\texto{110.8107}{445.4510}{11.9552}{\normalfont}{reseñas;}
\texto{154.4615}{445.4510}{11.9552}{\normalfont}{gestión}
\texto{194.1456}{445.4510}{11.9552}{\normalfont}{de}
\texto{209.7017}{445.4510}{11.9552}{\normalfont}{usuarios;}
\texto{258.5555}{445.4510}{11.9552}{\normalfont}{supervisión}
\texto{319.7637}{445.4510}{11.9552}{\normalfont}{de}
\texto{335.3079}{445.4510}{11.9552}{\normalfont}{reservas}
\texto{379.2886}{445.4510}{11.9552}{\normalfont}{y}
\texto{389.3167}{445.4510}{11.9552}{\normalfont}{pagos;}
\texto{425.0866}{445.4510}{11.9552}{\normalfont}{consulta}
\texto{471.2743}{445.4510}{11.9552}{\normalfont}{y}
\texto{481.3024}{445.4510}{11.9552}{\normalfont}{exportación}
\texto{70.8660}{431.0050}{11.9552}{\normalfont}{de}
\texto{88.2632}{431.0050}{11.9552}{\normalfont}{reportes}
\texto{135.3380}{431.0050}{11.9552}{\normalfont}{financieros;}
\texto{198.3371}{431.0050}{11.9552}{\normalfont}{configuración}
\texto{271.6775}{431.0050}{11.9552}{\normalfont}{de}
\texto{289.0867}{431.0050}{11.9552}{\normalfont}{la}
\texto{303.8824}{431.0050}{11.9552}{\normalfont}{plataforma}
\texto{365.1911}{431.0050}{11.9552}{\normalfont}{(comisión}
\texto{419.7355}{431.0050}{11.9552}{\normalfont}{y}
\texto{431.6166}{431.0050}{11.9552}{\normalfont}{planes)}
\texto{473.8041}{431.0050}{11.9552}{\normalfont}{y}
\texto{485.6851}{431.0050}{11.9552}{\normalfont}{módulo}
\texto{529.4316}{431.0050}{11.9552}{\normalfont}{de}
\texto{70.8660}{416.5590}{11.9552}{\normalfont}{auditoría.}
\texto{70.8660}{381.8980}{14.3462}{\normalfont\bfseries}{F}
\texto{79.6774}{381.8980}{14.3462}{\normalfont\bfseries}{uera}
\texto{115.8112}{381.8980}{14.3462}{\normalfont\bfseries}{de}
\texto{137.5213}{381.8980}{14.3462}{\normalfont\bfseries}{alcance}
\texto{193.7527}{381.8980}{14.3462}{\normalfont\bfseries}{de}
\texto{215.4628}{381.8980}{14.3462}{\normalfont\bfseries}{este}
\texto{248.2137}{381.8980}{14.3462}{\normalfont\bfseries}{prototipo}
\texto{100.7540}{359.7310}{11.9552}{\normalfont\bfseries}{Aplicación}
\texto{167.3014}{359.7310}{11.9552}{\normalfont\bfseries}{móvil}
\texto{204.4558}{359.7310}{11.9552}{\normalfont\bfseries}{para}
\texto{235.5309}{359.7310}{11.9552}{\normalfont\bfseries}{turistas:}
\texto{287.9570}{359.7310}{11.9552}{\normalfont}{la}
\texto{301.4257}{359.7310}{11.9552}{\normalfont}{app}
\texto{324.6619}{359.7310}{11.9552}{\normalfont}{móvil}
\texto{357.0055}{359.7310}{11.9552}{\normalfont}{(exploración}
\texto{424.1411}{359.7310}{11.9552}{\normalfont}{de}
\texto{440.2232}{359.7310}{11.9552}{\normalfont}{tours,}
\texto{473.9297}{359.7310}{11.9552}{\normalfont}{reserva,}
\texto{517.0712}{359.7310}{11.9552}{\normalfont}{pago}
\texto{100.7540}{345.2850}{11.9552}{\normalfont}{en}
\texto{117.1828}{345.2850}{11.9552}{\normalfont}{línea}
\texto{145.9686}{345.2850}{11.9552}{\normalfont}{y}
\texto{156.8813}{345.2850}{11.9552}{\normalfont}{voucher}
\texto{200.9553}{345.2850}{11.9552}{\normalfont}{digital)}
\texto{242.7483}{345.2850}{11.9552}{\normalfont}{se}
\texto{257.2906}{345.2850}{11.9552}{\normalfont}{contempla}
\texto{314.3815}{345.2850}{11.9552}{\normalfont}{como}
\texto{345.7687}{345.2850}{11.9552}{\normalfont}{trabajo}
\texto{387.5712}{345.2850}{11.9552}{\normalfont}{futuro}
\texto{423.8361}{345.2850}{11.9552}{\normalfont}{mediante}
\texto{475.0617}{345.2850}{11.9552}{\normalfont}{integraciones}
\texto{100.7540}{330.8390}{11.9552}{\normalfont}{externas}
\texto{147.3147}{330.8390}{11.9552}{\normalfont}{con}
\texto{168.7719}{330.8390}{11.9552}{\normalfont}{la}
\texto{181.7863}{330.8390}{11.9552}{\normalfont}{misma}
\texto{218.9168}{330.8390}{11.9552}{\normalfont}{API}
\texto{243.7764}{330.8390}{11.9552}{\normalfont}{REST.}
\texto{100.7540}{306.4310}{11.9552}{\normalfont\bfseries}{Reserva}
\texto{152.1936}{306.4310}{11.9552}{\normalfont\bfseries}{de}
\texto{171.2419}{306.4310}{11.9552}{\normalfont\bfseries}{hospedaje,}
\texto{238.5389}{306.4310}{11.9552}{\normalfont\bfseries}{transporte}
\texto{305.4581}{306.4310}{11.9552}{\normalfont\bfseries}{terrestre}
\texto{361.5543}{306.4310}{11.9552}{\normalfont\bfseries}{o}
\texto{373.7307}{306.4310}{11.9552}{\normalfont\bfseries}{aéreo:}
\texto{413.2430}{306.4310}{11.9552}{\normalfont}{el}
\texto{426.4320}{306.4310}{11.9552}{\normalfont}{proyecto}
\texto{474.4201}{306.4310}{11.9552}{\normalfont}{no}
\texto{491.5233}{306.4310}{11.9552}{\normalfont}{incluye}
\texto{532.0299}{306.4310}{11.9552}{\normalfont}{la}
\texto{100.7540}{291.9850}{11.9552}{\normalfont}{gestión}
\texto{140.4859}{291.9850}{11.9552}{\normalfont}{de}
\texto{156.0898}{291.9850}{11.9552}{\normalfont}{alojamiento}
\texto{219.5074}{291.9850}{11.9552}{\normalfont}{ni}
\texto{233.1602}{291.9850}{11.9552}{\normalfont}{transporte.}
\texto{100.7540}{267.5760}{11.9552}{\normalfont\bfseries}{F}
\texto{108.0969}{267.5760}{11.9552}{\normalfont\bfseries}{uncionalidades}
\texto{198.5487}{267.5760}{11.9552}{\normalfont\bfseries}{avanzadas}
\texto{261.5909}{267.5760}{11.9552}{\normalfont\bfseries}{de}
\texto{279.8859}{267.5760}{11.9552}{\normalfont\bfseries}{la}
\texto{294.8467}{267.5760}{11.9552}{\normalfont\bfseries}{app}
\texto{321.0035}{267.5760}{11.9552}{\normalfont\bfseries}{móvil:}
\texto{360.9480}{267.5760}{11.9552}{\normalfont}{reseñas,}
\texto{404.8140}{267.5760}{11.9552}{\normalfont}{notificaciones}
\texto{477.2434}{267.5760}{11.9552}{\normalfont}{push}
\texto{505.4481}{267.5760}{11.9552}{\normalfont}{y}
\texto{515.7033}{267.5760}{11.9552}{\normalfont}{otras}
\texto{100.7540}{253.1310}{11.9552}{\normalfont}{funcionalidades}
\texto{183.0847}{253.1310}{11.9552}{\normalfont}{transaccionales}
\texto{263.8540}{253.1310}{11.9552}{\normalfont}{del}
\texto{282.7217}{253.1310}{11.9552}{\normalfont}{turista.}
\texto{70.8660}{218.4690}{14.3462}{\normalfont\bfseries}{Restricciones}
\texto{168.9624}{218.4690}{14.3462}{\normalfont\bfseries}{y}
\texto{182.8610}{218.4690}{14.3462}{\normalfont\bfseries}{limitaciones}
\texto{88.7990}{196.3020}{11.9552}{\normalfont}{El}
\texto{103.1237}{196.3020}{11.9552}{\normalfont}{desarrollo}
\texto{155.7362}{196.3020}{11.9552}{\normalfont}{se}
\texto{168.6645}{196.3020}{11.9552}{\normalfont}{limita}
\texto{201.7015}{196.3020}{11.9552}{\normalfont}{a}
\texto{210.6751}{196.3020}{11.9552}{\normalfont}{un}
\texto{226.7907}{196.3020}{11.9552}{\normalfont}{prototipo}
\texto{277.7103}{196.3020}{11.9552}{\normalfont}{funcional}
\texto{327.3315}{196.3020}{11.9552}{\normalfont}{de}
\texto{342.1584}{196.3020}{11.9552}{\normalfont}{las}
\texto{358.9889}{196.3020}{11.9552}{\normalfont}{aplicaciones}
\texto{422.6575}{196.3020}{11.9552}{\normalfont}{de}
\texto{437.4844}{196.3020}{11.9552}{\normalfont}{escritorio}
\texto{487.4834}{196.3020}{11.9552}{\normalfont}{(backoffice}
\texto{70.8660}{181.8560}{11.9552}{\normalfont}{de}
\texto{88.2273}{181.8560}{11.9552}{\normalfont}{operadores}
\texto{148.9119}{181.8560}{11.9552}{\normalfont}{y}
\texto{160.7571}{181.8560}{11.9552}{\normalfont}{panel}
\texto{193.7392}{181.8560}{11.9552}{\normalfont}{de}
\texto{211.1125}{181.8560}{11.9552}{\normalfont}{administración)}
\texto{296.8384}{181.8560}{11.9552}{\normalfont}{sobre}
\texto{329.2346}{181.8560}{11.9552}{\normalfont}{.NET}
\texto{363.3356}{181.8560}{11.9552}{\normalfont}{Windows}
\texto{415.8907}{181.8560}{11.9552}{\normalfont}{Forms,}
\texto{456.2467}{181.8560}{11.9552}{\normalfont}{con}
\texto{479.4732}{181.8560}{11.9552}{\normalfont}{consumo}
\texto{529.4292}{181.8560}{11.9552}{\normalfont}{de}
\texto{70.8660}{167.4100}{11.9552}{\normalfont}{una}
\texto{95.2618}{167.4100}{11.9552}{\normalfont}{API}
\texto{121.7593}{167.4100}{11.9552}{\normalfont}{REST.}
\texto{162.0842}{167.4100}{11.9552}{\normalfont}{El}
\texto{178.8478}{167.4100}{11.9552}{\normalfont}{alcance}
\texto{221.4537}{167.4100}{11.9552}{\normalfont}{temporal}
\texto{272.8371}{167.4100}{11.9552}{\normalfont}{y}
\texto{284.5628}{167.4100}{11.9552}{\normalfont}{de}
\texto{301.8046}{167.4100}{11.9552}{\normalfont}{recursos}
\texto{348.4418}{167.4100}{11.9552}{\normalfont}{disponible}
\texto{405.7551}{167.4100}{11.9552}{\normalfont}{en}
\texto{422.9969}{167.4100}{11.9552}{\normalfont}{la}
\texto{437.6372}{167.4100}{11.9552}{\normalfont}{asignatura}
\texto{496.5787}{167.4100}{11.9552}{\normalfont}{limita}
\texto{532.0307}{167.4100}{11.9552}{\normalfont}{la}
\texto{70.8660}{152.9640}{11.9552}{\normalfont}{implementación}
\texto{155.0617}{152.9640}{11.9552}{\normalfont}{completa}
\texto{205.1109}{152.9640}{11.9552}{\normalfont}{de}
\texto{220.6910}{152.9640}{11.9552}{\normalfont}{todos}
\texto{252.2790}{152.9640}{11.9552}{\normalfont}{los}
\texto{269.8747}{152.9640}{11.9552}{\normalfont}{módulos}
\texto{316.4091}{152.9640}{11.9552}{\normalfont}{previstos,}
\texto{368.8757}{152.9640}{11.9552}{\normalfont}{priorizándose}
\texto{440.4514}{152.9640}{11.9552}{\normalfont}{los}
\texto{458.0471}{152.9640}{11.9552}{\normalfont}{procesos}
\texto{504.6580}{152.9640}{11.9552}{\normalfont}{críticos}
\texto{70.8660}{138.5190}{11.9552}{\normalfont}{de}
\texto{86.4699}{138.5190}{11.9552}{\normalfont}{gestión}
\texto{126.2018}{138.5190}{11.9552}{\normalfont}{de}
\texto{141.8177}{138.5190}{11.9552}{\normalfont}{tours,}
\texto{175.0460}{138.5190}{11.9552}{\normalfont}{disponibilidad}
\texto{250.8826}{138.5190}{11.9552}{\normalfont}{y}
\texto{260.9585}{138.5190}{11.9552}{\normalfont}{reservas.}
\end{pagina}
% Página original 20
\begin{pagina}{20}
\texto{70.8660}{746.3890}{11.9552}{\normalfont}{10}
\texto{266.5910}{746.3890}{11.9552}{\normalfont}{CAPÍTULO}
\texto{333.5545}{746.3890}{11.9552}{\normalfont}{1.}
\texto{350.4663}{746.3890}{11.9552}{\normalfont}{FORMULACIÓN}
\texto{445.2256}{746.3890}{11.9552}{\normalfont}{DEL}
\texto{473.3418}{746.3890}{11.9552}{\normalfont}{PROYECTO}
\texto{70.8660}{705.1940}{17.2154}{\normalfont\bfseries}{1.8.}
\texto{120.3603}{705.1940}{17.2154}{\normalfont\bfseries}{Metodología}
\texto{231.7577}{705.1940}{17.2154}{\normalfont\bfseries}{de}
\texto{257.8098}{705.1940}{17.2154}{\normalfont\bfseries}{desarrollo}
\texto{88.7990}{678.9090}{11.9552}{\normalfont}{El}
\texto{105.2039}{678.9090}{11.9552}{\normalfont}{desarrollo}
\texto{159.8966}{678.9090}{11.9552}{\normalfont}{de}
\texto{176.7990}{678.9090}{11.9552}{\normalfont\bfseries}{Explora-Nicaragua}
\texto{292.1630}{678.9090}{11.9552}{\normalfont}{se}
\texto{307.1716}{678.9090}{11.9552}{\normalfont}{sustenta}
\texto{354.4400}{678.9090}{11.9552}{\normalfont}{en}
\texto{371.3471}{678.9090}{11.9552}{\normalfont}{un}
\texto{389.5548}{678.9090}{11.9552}{\normalfont}{enfoque}
\texto{434.0880}{678.9090}{11.9552}{\normalfont}{metodológico}
\texto{506.5986}{678.9090}{11.9552}{\normalfont}{estruc-}
\texto{70.8660}{664.4630}{11.9552}{\normalfont}{turado}
\texto{109.3833}{664.4630}{11.9552}{\normalfont}{que}
\texto{131.9666}{664.4630}{11.9552}{\normalfont}{combina}
\texto{179.2781}{664.4630}{11.9552}{\normalfont}{investigación,}
\texto{252.6592}{664.4630}{11.9552}{\normalfont}{diseño,}
\texto{292.5417}{664.4630}{11.9552}{\normalfont}{ingeniería}
\texto{346.6796}{664.4630}{11.9552}{\normalfont}{de}
\texto{363.0846}{664.4630}{11.9552}{\normalfont}{software}
\texto{409.8007}{664.4630}{11.9552}{\normalfont}{y}
\texto{420.6775}{664.4630}{11.9552}{\normalfont}{validación}
\texto{476.4342}{664.4630}{11.9552}{\normalfont}{continua.}
\texto{527.9755}{664.4630}{11.9552}{\normalfont}{La}
\texto{70.8660}{650.0170}{11.9552}{\normalfont}{metodología}
\texto{137.3010}{650.0170}{11.9552}{\normalfont}{seleccionada}
\texto{204.1282}{650.0170}{11.9552}{\normalfont}{es}
\texto{218.2750}{650.0170}{11.9552}{\normalfont\bfseries}{Scrum}
\texto{255.8980}{650.0170}{11.9552}{\normalfont}{,}
\texto{263.4776}{650.0170}{11.9552}{\normalfont}{elegida}
\texto{302.9250}{650.0170}{11.9552}{\normalfont}{por}
\texto{324.4850}{650.0170}{11.9552}{\normalfont}{su}
\texto{339.9335}{650.0170}{11.9552}{\normalfont}{carácter}
\texto{385.2341}{650.0170}{11.9552}{\normalfont}{iterativo}
\texto{432.4883}{650.0170}{11.9552}{\normalfont}{e}
\texto{442.0189}{650.0170}{11.9552}{\normalfont}{incremental,}
\texto{509.1091}{650.0170}{11.9552}{\normalfont}{que}
\texto{531.3219}{650.0170}{11.9552}{\normalfont}{se}
\texto{70.8660}{635.5710}{11.9552}{\normalfont}{adapta}
\texto{110.1101}{635.5710}{11.9552}{\normalfont}{a}
\texto{120.0999}{635.5710}{11.9552}{\normalfont}{las}
\texto{137.9586}{635.5710}{11.9552}{\normalfont}{entregas}
\texto{184.1104}{635.5710}{11.9552}{\normalfont}{parciales}
\texto{232.5362}{635.5710}{11.9552}{\normalfont}{del}
\texto{251.6310}{635.5710}{11.9552}{\normalfont}{componente}
\texto{316.9016}{635.5710}{11.9552}{\normalfont}{académico}
\texto{373.7175}{635.5710}{11.9552}{\normalfont}{y}
\texto{384.0325}{635.5710}{11.9552}{\normalfont}{permite}
\texto{427.5135}{635.5710}{11.9552}{\normalfont}{incorporar}
\texto{484.6522}{635.5710}{11.9552}{\normalfont}{el}
\texto{497.2434}{635.5710}{11.9552}{\normalfont}{feedback}
\texto{70.8660}{621.1260}{11.9552}{\normalfont}{de}
\texto{86.4699}{621.1260}{11.9552}{\normalfont}{los}
\texto{104.0895}{621.1260}{11.9552}{\normalfont}{operadores}
\texto{163.0167}{621.1260}{11.9552}{\normalfont}{piloto}
\texto{196.1804}{621.1260}{11.9552}{\normalfont}{en}
\texto{211.7963}{621.1260}{11.9552}{\normalfont}{ciclos}
\texto{243.0735}{621.1260}{11.9552}{\normalfont}{cortos.}
\texto{70.8660}{586.4640}{14.3462}{\normalfont\bfseries}{Equipo}
\texto{126.2882}{586.4640}{14.3462}{\normalfont\bfseries}{y}
\texto{140.1868}{586.4640}{14.3462}{\normalfont\bfseries}{organización}
\texto{88.7990}{564.2970}{11.9552}{\normalfont}{El}
\texto{104.1040}{564.2970}{11.9552}{\normalfont}{equipo}
\texto{142.0092}{564.2970}{11.9552}{\normalfont}{se}
\texto{155.9179}{564.2970}{11.9552}{\normalfont}{organiza}
\texto{202.9305}{564.2970}{11.9552}{\normalfont}{de}
\texto{218.7257}{564.2970}{11.9552}{\normalfont}{forma}
\texto{252.4059}{564.2970}{11.9552}{\normalfont}{multidisciplinaria,}
\texto{348.5879}{564.2970}{11.9552}{\normalfont}{cubriendo}
\texto{402.7545}{564.2970}{11.9552}{\normalfont}{desarrollo}
\texto{456.3353}{564.2970}{11.9552}{\normalfont}{móvil,}
\texto{491.6438}{564.2970}{11.9552}{\normalfont}{desarrollo}
\texto{70.8660}{549.8510}{11.9552}{\normalfont}{de}
\texto{87.8567}{549.8510}{11.9552}{\normalfont}{escritorio,}
\texto{143.2834}{549.8510}{11.9552}{\normalfont}{desarrollo}
\texto{198.0717}{549.8510}{11.9552}{\normalfont}{web}
\texto{223.1944}{549.8510}{11.9552}{\normalfont}{y}
\texto{234.6570}{549.8510}{11.9552}{\normalfont}{backend.}
\texto{284.5078}{549.8510}{11.9552}{\normalfont}{La}
\texto{302.9595}{549.8510}{11.9552}{\normalfont}{comunicación}
\texto{376.8593}{549.8510}{11.9552}{\normalfont}{se}
\texto{391.9755}{549.8510}{11.9552}{\normalfont}{coordina}
\texto{441.1568}{549.8510}{11.9552}{\normalfont}{mediante}
\texto{492.9443}{549.8510}{11.9552}{\normalfont}{reuniones}
\texto{70.8660}{535.4050}{11.9552}{\normalfont}{periódicas}
\texto{126.2018}{535.4050}{11.9552}{\normalfont}{y}
\texto{136.2896}{535.4050}{11.9552}{\normalfont}{un}
\texto{153.1943}{535.4050}{11.9552}{\normalfont}{tablero}
\texto{192.8616}{535.4050}{11.9552}{\normalfont}{de}
\texto{208.4775}{535.4050}{11.9552}{\normalfont}{trabajo}
\texto{249.4432}{535.4050}{11.9552}{\normalfont}{donde}
\texto{283.9196}{535.4050}{11.9552}{\normalfont}{se}
\texto{297.6370}{535.4050}{11.9552}{\normalfont}{prioriza}
\texto{340.5681}{535.4050}{11.9552}{\normalfont}{el}
\texto{352.9203}{535.4050}{11.9552}{\normalfont}{backlog}
\texto{395.1914}{535.4050}{11.9552}{\normalfont}{del}
\texto{414.0591}{535.4050}{11.9552}{\normalfont}{sprint.}
\texto{70.8660}{500.7440}{14.3462}{\normalfont\bfseries}{F}
\texto{79.6774}{500.7440}{14.3462}{\normalfont\bfseries}{ases}
\texto{112.9864}{500.7440}{14.3462}{\normalfont\bfseries}{y}
\texto{126.8850}{500.7440}{14.3462}{\normalfont\bfseries}{actividades}
\texto{85.6720}{478.5770}{11.9552}{\normalfont}{1.}
\texto{100.7540}{478.5770}{11.9552}{\normalfont\bfseries}{Investigación}
\texto{181.3595}{478.5770}{11.9552}{\normalfont\bfseries}{documental:}
\texto{256.7240}{478.5770}{11.9552}{\normalfont}{revisión}
\texto{299.7818}{478.5770}{11.9552}{\normalfont}{de}
\texto{315.1467}{478.5770}{11.9552}{\normalfont}{literatura}
\texto{366.9199}{478.5770}{11.9552}{\normalfont}{sobre}
\texto{397.3076}{478.5770}{11.9552}{\normalfont}{turismo}
\texto{440.0403}{478.5770}{11.9552}{\normalfont}{digital,}
\texto{479.4685}{478.5770}{11.9552}{\normalfont}{pasarelas}
\texto{529.4317}{478.5770}{11.9552}{\normalfont}{de}
\texto{100.7540}{464.1310}{11.9552}{\normalfont}{pago,}
\texto{131.9666}{464.1310}{11.9552}{\normalfont}{aplicaciones}
\texto{196.4123}{464.1310}{11.9552}{\normalfont}{de}
\texto{212.0282}{464.1310}{11.9552}{\normalfont}{escritorio}
\texto{262.8163}{464.1310}{11.9552}{\normalfont}{en}
\texto{278.4202}{464.1310}{11.9552}{\normalfont}{.NET}
\texto{310.7758}{464.1310}{11.9552}{\normalfont}{y}
\texto{320.8516}{464.1310}{11.9552}{\normalfont}{control}
\texto{360.1962}{464.1310}{11.9552}{\normalfont}{de}
\texto{375.8121}{464.1310}{11.9552}{\normalfont}{acceso}
\texto{411.6418}{464.1310}{11.9552}{\normalfont}{basado}
\texto{450.7353}{464.1310}{11.9552}{\normalfont}{en}
\texto{466.3392}{464.1310}{11.9552}{\normalfont}{roles.}
\texto{85.6720}{439.7220}{11.9552}{\normalfont}{2.}
\texto{100.7540}{439.7220}{11.9552}{\normalfont\bfseries}{Análisis}
\texto{151.1631}{439.7220}{11.9552}{\normalfont\bfseries}{de}
\texto{169.1951}{439.7220}{11.9552}{\normalfont\bfseries}{requerimientos:}
\texto{263.7970}{439.7220}{11.9552}{\normalfont}{identificación}
\texto{334.9615}{439.7220}{11.9552}{\normalfont}{de}
\texto{350.5176}{439.7220}{11.9552}{\normalfont}{necesidades}
\texto{413.0290}{439.7220}{11.9552}{\normalfont}{de}
\texto{428.5970}{439.7220}{11.9552}{\normalfont}{operadores}
\texto{487.4644}{439.7220}{11.9552}{\normalfont}{y}
\texto{497.4925}{439.7220}{11.9552}{\normalfont}{adminis-}
\texto{100.7540}{425.2770}{11.9552}{\normalfont}{tradores}
\texto{146.3392}{425.2770}{11.9552}{\normalfont}{mediante}
\texto{196.7399}{425.2770}{11.9552}{\normalfont}{entrevistas}
\texto{255.3993}{425.2770}{11.9552}{\normalfont}{y}
\texto{265.4871}{425.2770}{11.9552}{\normalfont}{encuestas.}
\texto{85.6720}{400.8680}{11.9552}{\normalfont}{3.}
\texto{100.7540}{400.8680}{11.9552}{\normalfont\bfseries}{Diseño}
\texto{146.4157}{400.8680}{11.9552}{\normalfont\bfseries}{de}
\texto{166.0019}{400.8680}{11.9552}{\normalfont\bfseries}{la}
\texto{182.2538}{400.8680}{11.9552}{\normalfont\bfseries}{arquitectura:}
\texto{263.6200}{400.8680}{11.9552}{\normalfont}{definición}
\texto{317.5977}{400.8680}{11.9552}{\normalfont}{de}
\texto{334.5048}{400.8680}{11.9552}{\normalfont}{la}
\texto{348.8223}{400.8680}{11.9552}{\normalfont}{arquitectura}
\texto{416.7828}{400.8680}{11.9552}{\normalfont}{orientada}
\texto{469.7874}{400.8680}{11.9552}{\normalfont}{a}
\texto{480.8412}{400.8680}{11.9552}{\normalfont}{servicios,}
\texto{532.0334}{400.8680}{11.9552}{\normalfont}{la}
\texto{100.7540}{386.4220}{11.9552}{\normalfont}{base}
\texto{126.8283}{386.4220}{11.9552}{\normalfont}{de}
\texto{142.4322}{386.4220}{11.9552}{\normalfont}{datos,}
\texto{176.9732}{386.4220}{11.9552}{\normalfont}{la}
\texto{189.9757}{386.4220}{11.9552}{\normalfont}{API}
\texto{214.8353}{386.4220}{11.9552}{\normalfont}{REST}
\texto{250.2825}{386.4220}{11.9552}{\normalfont}{y}
\texto{260.3583}{386.4220}{11.9552}{\normalfont}{el}
\texto{272.7224}{386.4220}{11.9552}{\normalfont}{modelo}
\texto{313.3629}{386.4220}{11.9552}{\normalfont}{de}
\texto{328.9788}{386.4220}{11.9552}{\normalfont}{datos}
\texto{360.2560}{386.4220}{11.9552}{\normalfont}{compartido.}
\texto{85.6720}{362.0140}{11.9552}{\normalfont}{4.}
\texto{100.7540}{362.0140}{11.9552}{\normalfont\bfseries}{Desarrollo}
\texto{164.5744}{362.0140}{11.9552}{\normalfont\bfseries}{ágil}
\texto{188.9439}{362.0140}{11.9552}{\normalfont\bfseries}{(Scrum):}
\texto{243.9260}{362.0140}{11.9552}{\normalfont}{implementación}
\texto{327.4044}{362.0140}{11.9552}{\normalfont}{iterativa}
\texto{373.1521}{362.0140}{11.9552}{\normalfont}{en}
\texto{388.0268}{362.0140}{11.9552}{\normalfont}{sprints}
\texto{425.4585}{362.0140}{11.9552}{\normalfont}{de}
\texto{440.3212}{362.0140}{11.9552}{\normalfont}{dos}
\texto{460.4514}{362.0140}{11.9552}{\normalfont}{semanas,}
\texto{509.2621}{362.0140}{11.9552}{\normalfont}{priori-}
\texto{100.7540}{347.5680}{11.9552}{\normalfont}{zando}
\texto{133.4204}{347.5680}{11.9552}{\normalfont}{en}
\texto{147.8766}{347.5680}{11.9552}{\normalfont}{esta}
\texto{170.8521}{347.5680}{11.9552}{\normalfont}{etapa}
\texto{201.5674}{347.5680}{11.9552}{\normalfont}{el}
\texto{212.7838}{347.5680}{11.9552}{\normalfont}{backoffice}
\texto{264.6407}{347.5680}{11.9552}{\normalfont}{de}
\texto{279.0969}{347.5680}{11.9552}{\normalfont}{escritorio}
\texto{328.7373}{347.5680}{11.9552}{\normalfont}{de}
\texto{343.1935}{347.5680}{11.9552}{\normalfont}{operadores}
\texto{400.9730}{347.5680}{11.9552}{\normalfont}{y}
\texto{409.9011}{347.5680}{11.9552}{\normalfont}{el}
\texto{421.1055}{347.5680}{11.9552}{\normalfont}{panel}
\texto{451.1705}{347.5680}{11.9552}{\normalfont}{de}
\texto{465.6387}{347.5680}{11.9552}{\normalfont}{administración}
\texto{100.7540}{333.1220}{11.9552}{\normalfont}{de}
\texto{116.3579}{333.1220}{11.9552}{\normalfont}{escritorio.}
\texto{85.6720}{308.7140}{11.9552}{\normalfont}{5.}
\texto{100.7540}{308.7140}{11.9552}{\normalfont\bfseries}{Integración}
\texto{171.8002}{308.7140}{11.9552}{\normalfont\bfseries}{de}
\texto{190.3224}{308.7140}{11.9552}{\normalfont\bfseries}{componentes:}
\texto{274.1940}{308.7140}{11.9552}{\normalfont}{conexión}
\texto{323.0119}{308.7140}{11.9552}{\normalfont}{de}
\texto{338.9864}{308.7140}{11.9552}{\normalfont}{las}
\texto{356.9885}{308.7140}{11.9552}{\normalfont}{dos}
\texto{378.2305}{308.7140}{11.9552}{\normalfont}{aplicaciones}
\texto{443.0588}{308.7140}{11.9552}{\normalfont}{de}
\texto{459.0333}{308.7140}{11.9552}{\normalfont}{escritorio}
\texto{510.1920}{308.7140}{11.9552}{\normalfont}{con}
\texto{532.0318}{308.7140}{11.9552}{\normalfont}{la}
\texto{100.7540}{294.2680}{11.9552}{\normalfont}{API}
\texto{125.6136}{294.2680}{11.9552}{\normalfont}{REST,}
\texto{164.3007}{294.2680}{11.9552}{\normalfont}{autenticación}
\texto{236.1634}{294.2680}{11.9552}{\normalfont}{JWT}
\texto{266.5678}{294.2680}{11.9552}{\normalfont}{y}
\texto{276.6437}{294.2680}{11.9552}{\normalfont}{control}
\texto{315.9882}{294.2680}{11.9552}{\normalfont}{de}
\texto{331.6041}{294.2680}{11.9552}{\normalfont}{acceso}
\texto{367.4339}{294.2680}{11.9552}{\normalfont}{basado}
\texto{406.5274}{294.2680}{11.9552}{\normalfont}{en}
\texto{422.1313}{294.2680}{11.9552}{\normalfont}{roles.}
\texto{85.6720}{269.8590}{11.9552}{\normalfont}{6.}
\texto{100.7540}{269.8590}{11.9552}{\normalfont\bfseries}{Pruebas}
\texto{151.7907}{269.8590}{11.9552}{\normalfont\bfseries}{de}
\texto{168.8305}{269.8590}{11.9552}{\normalfont\bfseries}{usabilidad:}
\texto{235.0410}{269.8590}{11.9552}{\normalfont}{evaluación}
\texto{291.0272}{269.8590}{11.9552}{\normalfont}{con}
\texto{311.5758}{269.8590}{11.9552}{\normalfont}{operadores}
\texto{369.5824}{269.8590}{11.9552}{\normalfont}{piloto}
\texto{401.8376}{269.8590}{11.9552}{\normalfont}{en}
\texto{416.5329}{269.8590}{11.9552}{\normalfont}{los}
\texto{433.2439}{269.8590}{11.9552}{\normalfont}{destinos}
\texto{477.3346}{269.8590}{11.9552}{\normalfont}{de}
\texto{492.0300}{269.8590}{11.9552}{\normalfont}{Ometepe,}
\texto{100.7540}{255.4140}{11.9552}{\normalfont}{Somoto}
\texto{143.0228}{255.4140}{11.9552}{\normalfont}{y}
\texto{153.0986}{255.4140}{11.9552}{\normalfont}{San}
\texto{175.8685}{255.4140}{11.9552}{\normalfont}{Juan}
\texto{204.6399}{255.4140}{11.9552}{\normalfont}{del}
\texto{223.4956}{255.4140}{11.9552}{\normalfont}{Sur.}
\texto{85.6720}{231.0050}{11.9552}{\normalfont}{7.}
\texto{100.7540}{231.0050}{11.9552}{\normalfont\bfseries}{Despliegue}
\texto{170.1121}{231.0050}{11.9552}{\normalfont\bfseries}{piloto:}
\texto{212.8750}{231.0050}{11.9552}{\normalfont}{lanzamiento}
\texto{279.7165}{231.0050}{11.9552}{\normalfont}{controlado}
\texto{338.7656}{231.0050}{11.9552}{\normalfont}{de}
\texto{355.8759}{231.0050}{11.9552}{\normalfont}{las}
\texto{374.9899}{231.0050}{11.9552}{\normalfont}{aplicaciones}
\texto{440.9300}{231.0050}{11.9552}{\normalfont}{de}
\texto{458.0283}{231.0050}{11.9552}{\normalfont}{escritorio}
\texto{510.3108}{231.0050}{11.9552}{\normalfont}{en}
\texto{527.4091}{231.0050}{11.9552}{\normalfont}{los}
\texto{100.7540}{216.5590}{11.9552}{\normalfont}{destinos}
\texto{145.7534}{216.5590}{11.9552}{\normalfont}{seleccionados.}
\texto{88.7990}{189.1620}{11.9552}{\normalfont}{Como}
\texto{123.2108}{189.1620}{11.9552}{\normalfont}{evidencias}
\texto{179.4744}{189.1620}{11.9552}{\normalfont}{del}
\texto{198.9399}{189.1620}{11.9552}{\normalfont}{proceso}
\texto{241.5434}{189.1620}{11.9552}{\normalfont}{se}
\texto{255.8586}{189.1620}{11.9552}{\normalfont}{contemplan}
\texto{319.2259}{189.1620}{11.9552}{\normalfont}{el}
\texto{332.2210}{189.1620}{11.9552}{\normalfont\bfseries}{cronograma}
\texto{405.8770}{189.1620}{11.9552}{\normalfont}{de}
\texto{422.0787}{189.1620}{11.9552}{\normalfont}{actividades,}
\texto{486.7970}{189.1620}{11.9552}{\normalfont}{el}
\texto{499.8030}{189.1620}{11.9552}{\normalfont\bfseries}{tablero}
\texto{70.8660}{174.7160}{11.9552}{\normalfont\bfseries}{de}
\texto{88.9578}{174.7160}{11.9552}{\normalfont\bfseries}{trabajo}
\texto{135.3400}{174.7160}{11.9552}{\normalfont}{(backlog,}
\texto{185.4156}{174.7160}{11.9552}{\normalfont}{sprint}
\texto{218.9714}{174.7160}{11.9552}{\normalfont}{y}
\texto{229.0472}{174.7160}{11.9552}{\normalfont}{reuniones)}
\texto{285.7006}{174.7160}{11.9552}{\normalfont}{y}
\texto{295.7764}{174.7160}{11.9552}{\normalfont}{el}
\texto{308.1285}{174.7160}{11.9552}{\normalfont}{registro}
\texto{350.4738}{174.7160}{11.9552}{\normalfont}{del}
\texto{369.3296}{174.7160}{11.9552}{\normalfont}{avance}
\texto{407.0530}{174.7160}{11.9552}{\normalfont}{incremental}
\texto{470.4610}{174.7160}{11.9552}{\normalfont}{por}
\texto{491.6026}{174.7160}{11.9552}{\normalfont}{sprint.}
\end{pagina}
% Página original 21
\begin{pagina}{21}
\texto{70.8660}{746.3890}{11.9552}{\normalfont}{1.9.}
\texto{96.8710}{746.3890}{11.9552}{\normalfont}{ESTRUCTURA}
\texto{183.8331}{746.3890}{11.9552}{\normalfont}{DEL}
\texto{211.9613}{746.3890}{11.9552}{\normalfont}{ECOSISTEMA}
\texto{529.4280}{746.3890}{11.9552}{\normalfont}{11}
\texto{70.8660}{705.1940}{17.2154}{\normalfont\bfseries}{1.9.}
\texto{120.3603}{705.1940}{17.2154}{\normalfont\bfseries}{Estructura}
\texto{217.5946}{705.1940}{17.2154}{\normalfont\bfseries}{del}
\texto{249.0265}{705.1940}{17.2154}{\normalfont\bfseries}{ecosistema}
\texto{88.7990}{678.9090}{11.9552}{\normalfont}{Para}
\texto{116.8602}{678.9090}{11.9552}{\normalfont}{clarificar}
\texto{165.2573}{678.9090}{11.9552}{\normalfont}{la}
\texto{178.5347}{678.9090}{11.9552}{\normalfont}{coordinación}
\texto{247.4278}{678.9090}{11.9552}{\normalfont}{entre}
\texto{277.2799}{678.9090}{11.9552}{\normalfont}{los}
\texto{295.1744}{678.9090}{11.9552}{\normalfont}{componentes,}
\texto{368.3498}{678.9090}{11.9552}{\normalfont}{la}
\texto{381.6273}{678.9090}{11.9552}{\normalfont}{figura}
\texto{415.0680}{678.9090}{11.9552}{\normalfont}{1.1}
\texto{434.1980}{678.9090}{11.9552}{\normalfont}{presenta}
\texto{481.0241}{678.9090}{11.9552}{\normalfont}{el}
\texto{493.6512}{678.9090}{11.9552}{\normalfont}{diagrama}
\texto{70.8660}{664.4630}{11.9552}{\normalfont}{de}
\texto{88.2154}{664.4630}{11.9552}{\normalfont}{alto}
\texto{113.3691}{664.4630}{11.9552}{\normalfont}{nivel}
\texto{143.0778}{664.4630}{11.9552}{\normalfont}{del}
\texto{163.6790}{664.4630}{11.9552}{\normalfont}{ecosistema}
\texto{223.5840}{664.4630}{11.9552}{\normalfont\bfseries}{Explora-Nicaragua}
\texto{339.3980}{664.4630}{11.9552}{\normalfont}{con}
\texto{362.6007}{664.4630}{11.9552}{\normalfont}{la}
\texto{377.3486}{664.4630}{11.9552}{\normalfont}{representación}
\texto{456.2242}{664.4630}{11.9552}{\normalfont}{completa}
\texto{508.0428}{664.4630}{11.9552}{\normalfont}{de}
\texto{525.3922}{664.4630}{11.9552}{\normalfont}{sus}
\texto{70.8660}{650.0170}{11.9552}{\normalfont}{aplicaciones}
\texto{134.8096}{650.0170}{11.9552}{\normalfont}{y}
\texto{144.3833}{650.0170}{11.9552}{\normalfont}{servicios.}
\texto{193.7583}{650.0170}{11.9552}{\normalfont}{Si}
\texto{206.9210}{650.0170}{11.9552}{\normalfont}{bien}
\texto{231.7782}{650.0170}{11.9552}{\normalfont}{la}
\texto{244.2786}{650.0170}{11.9552}{\normalfont}{visión}
\texto{277.3299}{650.0170}{11.9552}{\normalfont}{completa}
\texto{326.9010}{650.0170}{11.9552}{\normalfont}{contempla}
\texto{382.6528}{650.0170}{11.9552}{\normalfont}{tres}
\texto{404.9851}{650.0170}{11.9552}{\normalfont}{aplicaciones,}
\texto{472.1760}{650.0170}{11.9552}{\normalfont\bfseries}{el}
\texto{485.9460}{650.0170}{11.9552}{\normalfont\bfseries}{prototipo}
\texto{70.8660}{635.5710}{11.9552}{\normalfont\bfseries}{inicial}
\texto{110.4007}{635.5710}{11.9552}{\normalfont\bfseries}{de}
\texto{128.6120}{635.5710}{11.9552}{\normalfont\bfseries}{este}
\texto{156.0241}{635.5710}{11.9552}{\normalfont\bfseries}{proyecto}
\texto{210.7442}{635.5710}{11.9552}{\normalfont\bfseries}{se}
\texto{226.7893}{635.5710}{11.9552}{\normalfont\bfseries}{limita}
\texto{265.5887}{635.5710}{11.9552}{\normalfont\bfseries}{a}
\texto{276.7298}{635.5710}{11.9552}{\normalfont\bfseries}{las}
\texto{296.9125}{635.5710}{11.9552}{\normalfont\bfseries}{dos}
\texto{321.0298}{635.5710}{11.9552}{\normalfont\bfseries}{aplicaciones}
\texto{394.9834}{635.5710}{11.9552}{\normalfont\bfseries}{de}
\texto{413.1948}{635.5710}{11.9552}{\normalfont\bfseries}{escritorio}
\texto{471.7660}{635.5710}{11.9552}{\normalfont}{(backoffice}
\texto{529.4188}{635.5710}{11.9552}{\normalfont}{de}
\texto{70.8660}{621.1260}{11.9552}{\normalfont}{operadores}
\texto{129.7812}{621.1260}{11.9552}{\normalfont}{y}
\texto{139.8571}{621.1260}{11.9552}{\normalfont}{panel}
\texto{171.0817}{621.1260}{11.9552}{\normalfont}{de}
\texto{186.6856}{621.1260}{11.9552}{\normalfont}{administración),}
\texto{273.9059}{621.1260}{11.9552}{\normalfont}{que}
\texto{295.6883}{621.1260}{11.9552}{\normalfont}{se}
\texto{309.4057}{621.1260}{11.9552}{\normalfont}{indican}
\texto{350.3857}{621.1260}{11.9552}{\normalfont}{con}
\texto{371.8429}{621.1260}{11.9552}{\normalfont}{color}
\texto{400.4661}{621.1260}{11.9552}{\normalfont}{destacado}
\texto{454.5060}{621.1260}{11.9552}{\normalfont}{en}
\texto{470.1099}{621.1260}{11.9552}{\normalfont}{la}
\texto{483.1243}{621.1260}{11.9552}{\normalfont}{figura.}
\texto{145.7480}{594.5980}{7.9701}{\normalfont}{App}
\texto{164.3215}{594.5980}{7.9701}{\normalfont}{móvil}
\texto{149.5980}{585.1340}{7.9701}{\normalfont}{(turista)}
\texto{150.4450}{575.6690}{7.9701}{\normalfont}{(futuro)}
\texto{149.2160}{531.2980}{7.9701}{\normalfont}{Panel}
\texto{172.6122}{531.2980}{7.9701}{\normalfont}{de}
\texto{137.8370}{521.8340}{7.9701}{\normalfont}{administración}
\texto{142.5380}{512.3690}{7.9701}{\normalfont}{de}
\texto{153.8276}{512.3690}{7.9701}{\normalfont}{escritorio}
\texto{137.1310}{502.9050}{7.9701}{\normalfont}{(equipo}
\texto{167.7059}{502.9050}{7.9701}{\normalfont}{gestor)}
\texto{146.7460}{458.5330}{7.9701}{\normalfont}{Backoffice}
\texto{142.5380}{449.0690}{7.9701}{\normalfont}{de}
\texto{153.8276}{449.0690}{7.9701}{\normalfont}{escritorio}
\texto{145.5020}{439.6040}{7.9701}{\normalfont}{(operador)}
\texto{297.8320}{526.3440}{7.9701}{\normalfont}{API}
\texto{315.8046}{526.3440}{7.9701}{\normalfont}{REST}
\texto{281.7140}{516.8800}{7.9701}{\normalfont}{Autenticación}
\texto{335.5720}{516.8800}{7.9701}{\normalfont}{JWT}
\texto{302.3540}{507.4150}{7.9701}{\normalfont}{y}
\texto{309.6450}{507.4150}{7.9701}{\normalfont}{RBAC}
\texto{417.3230}{543.5140}{7.9701}{\normalfont}{Base}
\texto{437.4770}{543.5140}{7.9701}{\normalfont}{de}
\texto{448.7666}{543.5140}{7.9701}{\normalfont}{datos}
\texto{425.3040}{534.0500}{7.9701}{\normalfont}{relacional}
\texto{421.8130}{487.8170}{7.9701}{\normalfont}{Pasarela}
\texto{455.6142}{487.8170}{7.9701}{\normalfont}{de}
\texto{418.3650}{478.3530}{7.9701}{\normalfont}{pago}
\texto{438.5947}{478.3530}{7.9701}{\normalfont}{(Stripe)}
\texto{70.8660}{409.1040}{11.9552}{\normalfont}{Figura}
\texto{108.9792}{409.1040}{11.9552}{\normalfont}{1.1:}
\texto{131.6730}{409.1040}{11.9552}{\normalfont}{Ecosistema}
\texto{193.0032}{409.1040}{11.9552}{\normalfont}{Explora-Nicaragua:}
\texto{296.4970}{409.1040}{11.9552}{\normalfont}{dos}
\texto{317.9422}{409.1040}{11.9552}{\normalfont}{aplicaciones}
\texto{382.9498}{409.1040}{11.9552}{\normalfont}{de}
\texto{399.1276}{409.1040}{11.9552}{\normalfont}{escritorio}
\texto{450.4895}{409.1040}{11.9552}{\normalfont}{(prototipo}
\texto{507.3245}{409.1040}{11.9552}{\normalfont}{inicial,}
\texto{70.8660}{394.6580}{11.9552}{\normalfont}{sombreado}
\texto{129.9606}{394.6580}{11.9552}{\normalfont}{verde)}
\texto{166.5650}{394.6580}{11.9552}{\normalfont}{y}
\texto{177.4657}{394.6580}{11.9552}{\normalfont}{una}
\texto{201.0605}{394.6580}{11.9552}{\normalfont}{app}
\texto{224.6553}{394.6580}{11.9552}{\normalfont}{móvil}
\texto{257.3576}{394.6580}{11.9552}{\normalfont}{(futuro)}
\texto{302.7275}{394.6580}{11.9552}{\normalfont}{coordinadas}
\texto{368.3329}{394.6580}{11.9552}{\normalfont}{sobre}
\texto{399.7966}{394.6580}{11.9552}{\normalfont}{una}
\texto{423.3914}{394.6580}{11.9552}{\normalfont}{API}
\texto{449.0879}{394.6580}{11.9552}{\normalfont}{REST}
\texto{485.3480}{394.6580}{11.9552}{\normalfont}{común}
\texto{523.5783}{394.6580}{11.9552}{\normalfont}{con}
\texto{70.8660}{380.2120}{11.9552}{\normalfont}{autenticación}
\texto{142.7287}{380.2120}{11.9552}{\normalfont}{JWT/RBAC,}
\texto{216.0380}{380.2120}{11.9552}{\normalfont}{base}
\texto{242.1123}{380.2120}{11.9552}{\normalfont}{de}
\texto{257.7282}{380.2120}{11.9552}{\normalfont}{datos}
\texto{289.0054}{380.2120}{11.9552}{\normalfont}{relacional}
\texto{341.6800}{380.2120}{11.9552}{\normalfont}{y}
\texto{351.7678}{380.2120}{11.9552}{\normalfont}{pasarela}
\texto{397.3529}{380.2120}{11.9552}{\normalfont}{de}
\texto{412.9688}{380.2120}{11.9552}{\normalfont}{pago}
\texto{440.9297}{380.2120}{11.9552}{\normalfont}{Stripe}
\texto{88.7990}{351.8190}{11.9552}{\normalfont}{La}
\texto{106.9398}{351.8190}{11.9552}{\normalfont}{API}
\texto{132.8874}{351.8190}{11.9552}{\normalfont}{REST}
\texto{169.3986}{351.8190}{11.9552}{\normalfont}{centraliza}
\texto{223.4767}{351.8190}{11.9552}{\normalfont}{la}
\texto{237.5671}{351.8190}{11.9552}{\normalfont}{lógica}
\texto{271.8188}{351.8190}{11.9552}{\normalfont}{de}
\texto{288.4987}{351.8190}{11.9552}{\normalfont}{negocio}
\texto{331.5278}{351.8190}{11.9552}{\normalfont}{y}
\texto{342.6796}{351.8190}{11.9552}{\normalfont}{expone}
\texto{383.4325}{351.8190}{11.9552}{\normalfont}{los}
\texto{402.1400}{351.8190}{11.9552}{\normalfont}{mismos}
\texto{444.9635}{351.8190}{11.9552}{\normalfont}{servicios}
\texto{492.6767}{351.8190}{11.9552}{\normalfont}{a}
\texto{503.5153}{351.8190}{11.9552}{\normalfont}{las}
\texto{522.2109}{351.8190}{11.9552}{\normalfont}{tres}
\texto{70.8660}{337.3730}{11.9552}{\normalfont}{aplicaciones;}
\texto{139.2569}{337.3730}{11.9552}{\normalfont}{la}
\texto{152.9647}{337.3730}{11.9552}{\normalfont}{autenticación}
\texto{225.5209}{337.3730}{11.9552}{\normalfont}{mediante}
\texto{276.6150}{337.3730}{11.9552}{\normalfont}{tokens}
\texto{313.7909}{337.3730}{11.9552}{\normalfont}{JWT}
\texto{344.8887}{337.3730}{11.9552}{\normalfont}{y}
\texto{355.6580}{337.3730}{11.9552}{\normalfont}{el}
\texto{368.7035}{337.3730}{11.9552}{\normalfont}{control}
\texto{408.7534}{337.3730}{11.9552}{\normalfont}{de}
\texto{425.0508}{337.3730}{11.9552}{\normalfont}{acceso}
\texto{461.5858}{337.3730}{11.9552}{\normalfont}{basado}
\texto{501.3608}{337.3730}{11.9552}{\normalfont}{en}
\texto{517.6581}{337.3730}{11.9552}{\normalfont}{roles}
\texto{70.8660}{322.9270}{11.9552}{\normalfont}{(RBAC)}
\texto{119.2224}{322.9270}{11.9552}{\normalfont}{garantizan}
\texto{178.3313}{322.9270}{11.9552}{\normalfont}{que}
\texto{201.6678}{322.9270}{11.9552}{\normalfont}{cada}
\texto{230.5325}{322.9270}{11.9552}{\normalfont}{actor}
\texto{261.9986}{322.9270}{11.9552}{\normalfont}{acceda}
\texto{301.2690}{322.9270}{11.9552}{\normalfont}{únicamente}
\texto{364.9304}{322.9270}{11.9552}{\normalfont}{a}
\texto{376.2353}{322.9270}{11.9552}{\normalfont}{las}
\texto{395.4090}{322.9270}{11.9552}{\normalfont}{funcionalidades}
\texto{479.2939}{322.9270}{11.9552}{\normalfont}{autorizadas.}
\texto{70.8660}{308.4810}{11.9552}{\normalfont}{La}
\texto{89.1981}{308.4810}{11.9552}{\normalfont}{base}
\texto{116.5396}{308.4810}{11.9552}{\normalfont}{de}
\texto{133.4108}{308.4810}{11.9552}{\normalfont}{datos}
\texto{165.9553}{308.4810}{11.9552}{\normalfont}{relacional}
\texto{219.8971}{308.4810}{11.9552}{\normalfont}{almacena}
\texto{272.5383}{308.4810}{11.9552}{\normalfont}{el}
\texto{286.1696}{308.4810}{11.9552}{\normalfont}{catálogo,}
\texto{336.8596}{308.4810}{11.9552}{\normalfont}{las}
\texto{355.7465}{308.4810}{11.9552}{\normalfont}{reservas,}
\texto{404.2941}{308.4810}{11.9552}{\normalfont}{los}
\texto{423.1810}{308.4810}{11.9552}{\normalfont}{pagos}
\texto{457.0261}{308.4810}{11.9552}{\normalfont}{y}
\texto{468.3692}{308.4810}{11.9552}{\normalfont}{la}
\texto{482.6390}{308.4810}{11.9552}{\normalfont}{bitácora}
\texto{529.4388}{308.4810}{11.9552}{\normalfont}{de}
\texto{70.8660}{294.0360}{11.9552}{\normalfont}{auditoría,}
\texto{124.3703}{294.0360}{11.9552}{\normalfont}{y}
\texto{134.6374}{294.0360}{11.9552}{\normalfont}{la}
\texto{147.8312}{294.0360}{11.9552}{\normalfont}{pasarela}
\texto{193.6076}{294.0360}{11.9552}{\normalfont}{Stripe}
\texto{228.5862}{294.0360}{11.9552}{\normalfont}{procesa}
\texto{270.7832}{294.0360}{11.9552}{\normalfont}{los}
\texto{288.5941}{294.0360}{11.9552}{\normalfont}{pagos}
\texto{321.3633}{294.0360}{11.9552}{\normalfont}{con}
\texto{342.9998}{294.0360}{11.9552}{\normalfont}{certificación}
\texto{408.2226}{294.0360}{11.9552}{\normalfont}{PCI}
\texto{432.9531}{294.0360}{11.9552}{\normalfont}{DSS.}
\texto{462.2410}{294.0360}{11.9552}{\normalfont}{En}
\texto{480.7979}{294.0360}{11.9552}{\normalfont}{el}
\texto{493.3412}{294.0360}{11.9552}{\normalfont}{prototipo}
\texto{70.8660}{279.5900}{11.9552}{\normalfont}{inicial,}
\texto{109.0127}{279.5900}{11.9552}{\normalfont}{los}
\texto{127.0746}{279.5900}{11.9552}{\normalfont}{servicios}
\texto{174.1422}{279.5900}{11.9552}{\normalfont}{de}
\texto{190.1885}{279.5900}{11.9552}{\normalfont}{la}
\texto{203.6333}{279.5900}{11.9552}{\normalfont}{API}
\texto{228.9353}{279.5900}{11.9552}{\normalfont}{REST}
\texto{264.8009}{279.5900}{11.9552}{\normalfont}{son}
\texto{286.1146}{279.5900}{11.9552}{\normalfont}{consumidos}
\texto{349.1161}{279.5900}{11.9552}{\normalfont}{por}
\texto{370.6881}{279.5900}{11.9552}{\normalfont}{las}
\texto{388.7500}{279.5900}{11.9552}{\normalfont}{dos}
\texto{410.0637}{279.5900}{11.9552}{\normalfont}{aplicaciones}
\texto{474.9518}{279.5900}{11.9552}{\normalfont}{de}
\texto{490.9980}{279.5900}{11.9552}{\normalfont}{escritorio,}
\texto{70.8660}{265.1440}{11.9552}{\normalfont}{mientras}
\texto{119.3754}{265.1440}{11.9552}{\normalfont}{que}
\texto{141.8034}{265.1440}{11.9552}{\normalfont}{la}
\texto{155.4634}{265.1440}{11.9552}{\normalfont}{app}
\texto{178.8669}{265.1440}{11.9552}{\normalfont}{móvil}
\texto{211.3779}{265.1440}{11.9552}{\normalfont}{permanece}
\texto{270.2357}{265.1440}{11.9552}{\normalfont}{como}
\texto{301.4436}{265.1440}{11.9552}{\normalfont}{trabajo}
\texto{343.0668}{265.1440}{11.9552}{\normalfont}{futuro.}
\texto{382.4042}{265.1440}{11.9552}{\normalfont}{Al}
\texto{398.9860}{265.1440}{11.9552}{\normalfont}{compartir}
\texto{453.6069}{265.1440}{11.9552}{\normalfont}{el}
\texto{466.6166}{265.1440}{11.9552}{\normalfont}{mismo}
\texto{504.3926}{265.1440}{11.9552}{\normalfont}{modelo}
\texto{70.8660}{250.6980}{11.9552}{\normalfont}{de}
\texto{86.0993}{250.6980}{11.9552}{\normalfont}{datos,}
\texto{120.2697}{250.6980}{11.9552}{\normalfont}{una}
\texto{142.6689}{250.6980}{11.9552}{\normalfont}{acción}
\texto{178.0755}{250.6980}{11.9552}{\normalfont}{realizada}
\texto{227.1396}{250.6980}{11.9552}{\normalfont}{en}
\texto{242.3849}{250.6980}{11.9552}{\normalfont}{cualquiera}
\texto{298.2778}{250.6980}{11.9552}{\normalfont}{de}
\texto{313.5231}{250.6980}{11.9552}{\normalfont}{las}
\texto{330.7840}{250.6980}{11.9552}{\normalfont}{aplicaciones}
\texto{394.8710}{250.6980}{11.9552}{\normalfont}{se}
\texto{408.2298}{250.6980}{11.9552}{\normalfont}{refleja}
\texto{442.6608}{250.6980}{11.9552}{\normalfont}{de}
\texto{457.9060}{250.6980}{11.9552}{\normalfont}{inmediato}
\texto{512.1731}{250.6980}{11.9552}{\normalfont}{en}
\texto{527.4184}{250.6980}{11.9552}{\normalfont}{las}
\texto{70.8660}{236.2520}{11.9552}{\normalfont}{demás.}
\end{pagina}

% Texto editable; coordenadas en puntos PDF.
% Página original 22
\begin{pagina}{22}
\texto{70.8660}{625.4930}{24.7871}{\normalfont\bfseries}{Capítulo}
\texto{185.1197}{625.4930}{24.7871}{\normalfont\bfseries}{2}
\texto{70.8660}{575.6790}{24.7871}{\normalfont\bfseries}{Investigación}
\texto{238.5830}{575.6790}{24.7871}{\normalfont\bfseries}{documental}
\texto{300.1470}{36.9930}{11.9552}{\normalfont}{12}
\end{pagina}
% Página original 23
\begin{pagina}{23}
\texto{70.8660}{746.3890}{11.9552}{\normalfont}{2.1.}
\texto{96.8710}{746.3890}{11.9552}{\normalfont}{METODOLOGÍA}
\texto{193.6782}{746.3890}{11.9552}{\normalfont}{DE}
\texto{214.4802}{746.3890}{11.9552}{\normalfont}{BÚSQUEDA}
\texto{529.4280}{746.3890}{11.9552}{\normalfont}{13}
\texto{70.8660}{705.1940}{17.2154}{\normalfont\bfseries}{2.1.}
\texto{120.3603}{705.1940}{17.2154}{\normalfont\bfseries}{Metodología}
\texto{231.7577}{705.1940}{17.2154}{\normalfont\bfseries}{de}
\texto{257.8098}{705.1940}{17.2154}{\normalfont\bfseries}{búsqueda}
\texto{88.7990}{678.9090}{11.9552}{\normalfont}{La}
\texto{105.6367}{678.9090}{11.9552}{\normalfont}{investigación}
\texto{174.7378}{678.9090}{11.9552}{\normalfont}{documental}
\texto{237.5910}{678.9090}{11.9552}{\normalfont}{que}
\texto{259.1463}{678.9090}{11.9552}{\normalfont}{sustenta}
\texto{304.8964}{678.9090}{11.9552}{\normalfont}{el}
\texto{317.0333}{678.9090}{11.9552}{\normalfont}{presente}
\texto{362.7189}{678.9090}{11.9552}{\normalfont}{proyecto}
\texto{409.6431}{678.9090}{11.9552}{\normalfont}{se}
\texto{423.1333}{678.9090}{11.9552}{\normalfont}{enmarca}
\texto{469.7275}{678.9090}{11.9552}{\normalfont}{en}
\texto{485.1163}{678.9090}{11.9552}{\normalfont}{un}
\texto{501.7938}{678.9090}{11.9552}{\normalfont}{enfoque}
\texto{70.8660}{664.4630}{11.9552}{\normalfont}{sistemático}
\texto{131.3784}{664.4630}{11.9552}{\normalfont}{de}
\texto{146.8987}{664.4630}{11.9552}{\normalfont}{revisión}
\texto{190.1119}{664.4630}{11.9552}{\normalfont}{de}
\texto{205.6322}{664.4630}{11.9552}{\normalfont}{fuentes,}
\texto{248.5227}{664.4630}{11.9552}{\normalfont}{cuyo}
\texto{275.7399}{664.4630}{11.9552}{\normalfont}{objetivo}
\texto{320.5264}{664.4630}{11.9552}{\normalfont}{es}
\texto{334.1482}{664.4630}{11.9552}{\normalfont}{identificar}
\texto{389.0058}{664.4630}{11.9552}{\normalfont}{el}
\texto{401.2743}{664.4630}{11.9552}{\normalfont}{marco}
\texto{436.2934}{664.4630}{11.9552}{\normalfont}{teórico}
\texto{474.5764}{664.4630}{11.9552}{\normalfont}{y}
\texto{484.5566}{664.4630}{11.9552}{\normalfont}{tecnológico}
\texto{70.8660}{650.0170}{11.9552}{\normalfont}{pertinente}
\texto{127.9401}{650.0170}{11.9552}{\normalfont}{para}
\texto{155.7479}{650.0170}{11.9552}{\normalfont}{el}
\texto{169.2477}{650.0170}{11.9552}{\normalfont}{diseño}
\texto{206.2371}{650.0170}{11.9552}{\normalfont}{del}
\texto{226.2406}{650.0170}{11.9552}{\normalfont}{ecosistema}
\texto{285.5450}{650.0170}{11.9552}{\normalfont\bfseries}{Explora-Nicaragua}
\texto{395.4520}{650.0170}{11.9552}{\normalfont}{.}
\texto{403.7489}{650.0170}{11.9552}{\normalfont}{La}
\texto{421.9734}{650.0170}{11.9552}{\normalfont}{estrategia}
\texto{476.5106}{650.0170}{11.9552}{\normalfont}{de}
\texto{493.2623}{650.0170}{11.9552}{\normalfont}{búsqueda}
\texto{70.8660}{635.5710}{11.9552}{\normalfont}{consideró}
\texto{122.3044}{635.5710}{11.9552}{\normalfont}{las}
\texto{139.9240}{635.5710}{11.9552}{\normalfont}{siguientes}
\texto{193.0553}{635.5710}{11.9552}{\normalfont}{fuentes:}
\texto{100.7540}{608.1740}{11.9552}{\normalfont\bfseries}{F}
\texto{108.0969}{608.1740}{11.9552}{\normalfont\bfseries}{uentes}
\texto{148.8988}{608.1740}{11.9552}{\normalfont\bfseries}{académicas:}
\texto{221.0540}{608.1740}{11.9552}{\normalfont}{artículos}
\texto{267.6697}{608.1740}{11.9552}{\normalfont}{indexados}
\texto{321.1262}{608.1740}{11.9552}{\normalfont}{en}
\texto{335.8215}{608.1740}{11.9552}{\normalfont}{bases}
\texto{365.6043}{608.1740}{11.9552}{\normalfont}{de}
\texto{380.2997}{608.1740}{11.9552}{\normalfont}{datos}
\texto{410.6683}{608.1740}{11.9552}{\normalfont}{como}
\texto{440.3219}{608.1740}{11.9552}{\normalfont}{Scopus,}
\texto{481.7467}{608.1740}{11.9552}{\normalfont}{IEEE}
\texto{512.8517}{608.1740}{11.9552}{\normalfont}{Xplo-}
\texto{100.7540}{593.7280}{11.9552}{\normalfont}{re}
\texto{114.9926}{593.7280}{11.9552}{\normalfont}{y}
\texto{125.6662}{593.7280}{11.9552}{\normalfont}{Google}
\texto{165.6958}{593.7280}{11.9552}{\normalfont}{Scholar,}
\texto{210.8291}{593.7280}{11.9552}{\normalfont}{enfocados}
\texto{264.8260}{593.7280}{11.9552}{\normalfont}{en}
\texto{281.0277}{593.7280}{11.9552}{\normalfont}{turismo}
\texto{324.6091}{593.7280}{11.9552}{\normalfont}{digital,}
\texto{364.8743}{593.7280}{11.9552}{\normalfont}{desarrollo}
\texto{418.8615}{593.7280}{11.9552}{\normalfont}{de}
\texto{435.0632}{593.7280}{11.9552}{\normalfont}{aplicaciones}
\texto{500.1067}{593.7280}{11.9552}{\normalfont}{móviles,}
\texto{100.7540}{579.2820}{11.9552}{\normalfont}{aplicaciones}
\texto{165.1997}{579.2820}{11.9552}{\normalfont}{de}
\texto{180.8036}{579.2820}{11.9552}{\normalfont}{escritorio}
\texto{231.6037}{579.2820}{11.9552}{\normalfont}{y}
\texto{241.6795}{579.2820}{11.9552}{\normalfont}{arquitecturas}
\texto{312.9540}{579.2820}{11.9552}{\normalfont}{orientadas}
\texto{369.2845}{579.2820}{11.9552}{\normalfont}{a}
\texto{379.0352}{579.2820}{11.9552}{\normalfont}{servicios.}
\texto{100.7540}{554.8740}{11.9552}{\normalfont\bfseries}{Repositorios}
\texto{177.1178}{554.8740}{11.9552}{\normalfont\bfseries}{de}
\texto{194.1098}{554.8740}{11.9552}{\normalfont\bfseries}{tesis:}
\texto{226.5080}{554.8740}{11.9552}{\normalfont}{consultas}
\texto{276.4042}{554.8740}{11.9552}{\normalfont}{en}
\texto{291.0398}{554.8740}{11.9552}{\normalfont}{el}
\texto{302.4355}{554.8740}{11.9552}{\normalfont}{repositorio}
\texto{359.7559}{554.8740}{11.9552}{\normalfont}{institucional}
\texto{425.8466}{554.8740}{11.9552}{\normalfont}{de}
\texto{440.4941}{554.8740}{11.9552}{\normalfont}{la}
\texto{452.5282}{554.8740}{11.9552}{\normalfont}{UNAN}
\texto{490.5697}{554.8740}{11.9552}{\normalfont}{y}
\texto{499.6891}{554.8740}{11.9552}{\normalfont}{en}
\texto{514.3366}{554.8740}{11.9552}{\normalfont}{bases}
\texto{100.7540}{540.4280}{11.9552}{\normalfont}{de}
\texto{116.5253}{540.4280}{11.9552}{\normalfont}{datos}
\texto{147.9818}{540.4280}{11.9552}{\normalfont}{como}
\texto{178.7234}{540.4280}{11.9552}{\normalfont}{TDX}
\texto{208.6353}{540.4280}{11.9552}{\normalfont}{para}
\texto{235.4748}{540.4280}{11.9552}{\normalfont}{identificar}
\texto{290.5954}{540.4280}{11.9552}{\normalfont}{trabajos}
\texto{336.3575}{540.4280}{11.9552}{\normalfont}{previos}
\texto{376.5939}{540.4280}{11.9552}{\normalfont}{sobre}
\texto{407.3881}{540.4280}{11.9552}{\normalfont}{turismo}
\texto{450.5512}{540.4280}{11.9552}{\normalfont}{y}
\texto{460.8064}{540.4280}{11.9552}{\normalfont}{tecnología}
\texto{516.2498}{540.4280}{11.9552}{\normalfont}{en}
\texto{532.0331}{540.4280}{11.9552}{\normalfont}{la}
\texto{100.7540}{525.9820}{11.9552}{\normalfont}{región.}
\texto{100.7540}{501.5740}{11.9552}{\normalfont\bfseries}{Legislación}
\texto{168.0187}{501.5740}{11.9552}{\normalfont\bfseries}{y}
\texto{178.2261}{501.5740}{11.9552}{\normalfont\bfseries}{normativa:}
\texto{243.9700}{501.5740}{11.9552}{\normalfont}{documentos}
\texto{307.2202}{501.5740}{11.9552}{\normalfont}{oficiales}
\texto{349.6588}{501.5740}{11.9552}{\normalfont}{del}
\texto{367.3190}{501.5740}{11.9552}{\normalfont}{Instituto}
\texto{414.6305}{501.5740}{11.9552}{\normalfont}{Nicaragüense}
\texto{484.7166}{501.5740}{11.9552}{\normalfont}{de}
\texto{499.1251}{501.5740}{11.9552}{\normalfont}{Turismo}
\texto{100.7540}{487.1280}{11.9552}{\normalfont}{(INTUR)}
\texto{153.1130}{487.1280}{11.9552}{\normalfont}{y}
\texto{163.7029}{487.1280}{11.9552}{\normalfont}{del}
\texto{183.0847}{487.1280}{11.9552}{\normalfont}{Banco}
\texto{219.1989}{487.1280}{11.9552}{\normalfont}{Central}
\texto{261.6710}{487.1280}{11.9552}{\normalfont}{de}
\texto{277.7890}{487.1280}{11.9552}{\normalfont}{Nicaragua}
\texto{333.8995}{487.1280}{11.9552}{\normalfont}{relacionados}
\texto{400.8223}{487.1280}{11.9552}{\normalfont}{con}
\texto{422.7936}{487.1280}{11.9552}{\normalfont}{el}
\texto{435.6718}{487.1280}{11.9552}{\normalfont}{sector}
\texto{470.0645}{487.1280}{11.9552}{\normalfont}{turístico}
\texto{516.8260}{487.1280}{11.9552}{\normalfont}{y}
\texto{527.4160}{487.1280}{11.9552}{\normalfont}{los}
\texto{100.7540}{472.6820}{11.9552}{\normalfont}{servicios}
\texto{147.3793}{472.6820}{11.9552}{\normalfont}{de}
\texto{162.9832}{472.6820}{11.9552}{\normalfont}{pago}
\texto{190.9560}{472.6820}{11.9552}{\normalfont}{digital.}
\texto{100.7540}{448.2740}{11.9552}{\normalfont\bfseries}{Documentación}
\texto{195.8935}{448.2740}{11.9552}{\normalfont\bfseries}{técnica:}
\texto{244.3310}{448.2740}{11.9552}{\normalfont}{manuales}
\texto{295.1717}{448.2740}{11.9552}{\normalfont}{de}
\texto{310.5006}{448.2740}{11.9552}{\normalfont}{desarrollo}
\texto{363.6152}{448.2740}{11.9552}{\normalfont}{de}
\texto{378.9561}{448.2740}{11.9552}{\normalfont}{React}
\texto{412.0051}{448.2740}{11.9552}{\normalfont}{Native,}
\texto{452.3706}{448.2740}{11.9552}{\normalfont}{.NET}
\texto{484.4392}{448.2740}{11.9552}{\normalfont}{y}
\texto{494.2521}{448.2740}{11.9552}{\normalfont}{Windows}
\texto{100.7540}{433.8280}{11.9552}{\normalfont}{Forms,}
\texto{139.3406}{433.8280}{11.9552}{\normalfont}{así}
\texto{156.9721}{433.8280}{11.9552}{\normalfont}{como}
\texto{187.5344}{433.8280}{11.9552}{\normalfont}{la}
\texto{200.5369}{433.8280}{11.9552}{\normalfont}{documentación}
\texto{281.1891}{433.8280}{11.9552}{\normalfont}{de}
\texto{296.7930}{433.8280}{11.9552}{\normalfont}{las}
\texto{314.4245}{433.8280}{11.9552}{\normalfont}{APIs}
\texto{343.9012}{433.8280}{11.9552}{\normalfont}{de}
\texto{359.5052}{433.8280}{11.9552}{\normalfont}{Google}
\texto{398.9370}{433.8280}{11.9552}{\normalfont}{Maps,}
\texto{433.7864}{433.8280}{11.9552}{\normalfont}{Mapbox}
\texto{479.3142}{433.8280}{11.9552}{\normalfont}{y}
\texto{489.3901}{433.8280}{11.9552}{\normalfont}{Stripe.}
\texto{88.7990}{406.4310}{11.9552}{\normalfont}{Los}
\texto{110.4810}{406.4310}{11.9552}{\normalfont}{criterios}
\texto{155.4158}{406.4310}{11.9552}{\normalfont}{de}
\texto{171.0316}{406.4310}{11.9552}{\normalfont}{inclusión}
\texto{219.8686}{406.4310}{11.9552}{\normalfont}{aplicados}
\texto{270.6567}{406.4310}{11.9552}{\normalfont}{a}
\texto{280.4193}{406.4310}{11.9552}{\normalfont}{la}
\texto{293.4218}{406.4310}{11.9552}{\normalfont}{selección}
\texto{341.6204}{406.4310}{11.9552}{\normalfont}{de}
\texto{357.2243}{406.4310}{11.9552}{\normalfont}{las}
\texto{374.8439}{406.4310}{11.9552}{\normalfont}{fuentes}
\texto{414.5902}{406.4310}{11.9552}{\normalfont}{fueron}
\texto{450.6805}{406.4310}{11.9552}{\normalfont}{los}
\texto{468.3001}{406.4310}{11.9552}{\normalfont}{siguientes:}
\texto{85.6720}{379.0330}{11.9552}{\normalfont}{1.}
\texto{100.7540}{379.0330}{11.9552}{\normalfont}{Idioma:}
\texto{143.3432}{379.0330}{11.9552}{\normalfont}{español}
\texto{185.0262}{379.0330}{11.9552}{\normalfont}{o}
\texto{194.7888}{379.0330}{11.9552}{\normalfont}{inglés.}
\texto{85.6720}{354.6250}{11.9552}{\normalfont}{2.}
\texto{100.7540}{354.6250}{11.9552}{\normalfont}{Actualidad:}
\texto{163.8320}{354.6250}{11.9552}{\normalfont}{publicaciones}
\texto{235.4317}{354.6250}{11.9552}{\normalfont}{de}
\texto{251.0356}{354.6250}{11.9552}{\normalfont}{los}
\texto{268.6672}{354.6250}{11.9552}{\normalfont}{últimos}
\texto{310.3502}{354.6250}{11.9552}{\normalfont}{cinco}
\texto{340.2740}{354.6250}{11.9552}{\normalfont}{años}
\texto{366.9987}{354.6250}{11.9552}{\normalfont}{(2021--2026).}
\texto{85.6720}{330.2160}{11.9552}{\normalfont}{3.}
\texto{100.7540}{330.2160}{11.9552}{\normalfont}{Relevancia}
\texto{159.2651}{330.2160}{11.9552}{\normalfont}{directa:}
\texto{201.6894}{330.2160}{11.9552}{\normalfont}{relación}
\texto{245.4143}{330.2160}{11.9552}{\normalfont}{con}
\texto{267.0269}{330.2160}{11.9552}{\normalfont}{al}
\texto{280.1848}{330.2160}{11.9552}{\normalfont}{menos}
\texto{316.1700}{330.2160}{11.9552}{\normalfont}{uno}
\texto{339.0833}{330.2160}{11.9552}{\normalfont}{de}
\texto{354.8426}{330.2160}{11.9552}{\normalfont}{los}
\texto{372.6176}{330.2160}{11.9552}{\normalfont}{temas}
\texto{406.6517}{330.2160}{11.9552}{\normalfont}{centrales}
\texto{455.3213}{330.2160}{11.9552}{\normalfont}{del}
\texto{474.3325}{330.2160}{11.9552}{\normalfont}{proyecto}
\texto{521.6273}{330.2160}{11.9552}{\normalfont}{(tu-}
\texto{100.7540}{315.7710}{11.9552}{\normalfont}{rismo}
\texto{132.3468}{315.7710}{11.9552}{\normalfont}{digital,}
\texto{171.6914}{315.7710}{11.9552}{\normalfont}{pasarelas}
\texto{221.5709}{315.7710}{11.9552}{\normalfont}{de}
\texto{236.8520}{315.7710}{11.9552}{\normalfont}{pago,}
\texto{267.7418}{315.7710}{11.9552}{\normalfont}{aplicaciones}
\texto{331.8648}{315.7710}{11.9552}{\normalfont}{de}
\texto{347.1459}{315.7710}{11.9552}{\normalfont}{escritorio,}
\texto{400.8630}{315.7710}{11.9552}{\normalfont}{control}
\texto{439.8848}{315.7710}{11.9552}{\normalfont}{de}
\texto{455.1659}{315.7710}{11.9552}{\normalfont}{acceso}
\texto{490.6728}{315.7710}{11.9552}{\normalfont}{basado}
\texto{529.4316}{315.7710}{11.9552}{\normalfont}{en}
\texto{100.7540}{301.3250}{11.9552}{\normalfont}{roles,}
\texto{131.3808}{301.3250}{11.9552}{\normalfont}{marketplaces).}
\texto{85.6720}{276.9160}{11.9552}{\normalfont}{4.}
\texto{100.7540}{276.9160}{11.9552}{\normalfont}{Calidad:}
\texto{147.3243}{276.9160}{11.9552}{\normalfont}{publicaciones}
\texto{218.6729}{276.9160}{11.9552}{\normalfont}{revisadas}
\texto{268.9613}{276.9160}{11.9552}{\normalfont}{por}
\texto{289.8518}{276.9160}{11.9552}{\normalfont}{pares}
\texto{320.2276}{276.9160}{11.9552}{\normalfont}{o}
\texto{329.7391}{276.9160}{11.9552}{\normalfont}{documentación}
\texto{410.1283}{276.9160}{11.9552}{\normalfont}{oficial}
\texto{443.6913}{276.9160}{11.9552}{\normalfont}{de}
\texto{459.0561}{276.9160}{11.9552}{\normalfont}{proveedores}
\texto{522.2800}{276.9160}{11.9552}{\normalfont}{tec-}
\texto{100.7540}{262.4700}{11.9552}{\normalfont}{nológicos.}
\end{pagina}
% Página original 24
\begin{pagina}{24}
\texto{70.8660}{746.3890}{11.9552}{\normalfont}{14}
\texto{269.9360}{746.3890}{11.9552}{\normalfont}{CAPÍTULO}
\texto{336.8995}{746.3890}{11.9552}{\normalfont}{2.}
\texto{353.8113}{746.3890}{11.9552}{\normalfont}{INVESTIGACIÓN}
\texto{454.8220}{746.3890}{11.9552}{\normalfont}{DOCUMENTAL}
\texto{70.8660}{705.1940}{17.2154}{\normalfont\bfseries}{2.2.}
\texto{120.3603}{705.1940}{17.2154}{\normalfont\bfseries}{T}
\texto{132.2286}{705.1940}{17.2154}{\normalfont\bfseries}{endencias}
\texto{219.6621}{705.1940}{17.2154}{\normalfont\bfseries}{tecnológicas}
\texto{327.3273}{705.1940}{17.2154}{\normalfont\bfseries}{relacionadas}
\texto{88.7990}{677.5970}{11.9552}{\normalfont}{A}
\texto{101.6508}{677.5970}{11.9552}{\normalfont}{continuación,}
\texto{173.3701}{677.5970}{11.9552}{\normalfont}{se}
\texto{187.2668}{677.5970}{11.9552}{\normalfont}{describen}
\texto{239.2193}{677.5970}{11.9552}{\normalfont}{las}
\texto{257.0302}{677.5970}{11.9552}{\normalfont}{tendencias}
\texto{314.5012}{677.5970}{11.9552}{\normalfont}{tecnológicas}
\texto{379.7766}{677.5970}{11.9552}{\normalfont}{que}
\texto{401.7503}{677.5970}{11.9552}{\normalfont}{influyen}
\texto{446.1519}{677.5970}{11.9552}{\normalfont}{directamente}
\texto{516.2428}{677.5970}{11.9552}{\normalfont}{en}
\texto{532.0380}{677.5970}{11.9552}{\normalfont}{la}
\texto{70.8660}{663.1510}{11.9552}{\normalfont}{arquitectura}
\texto{137.7147}{663.1510}{11.9552}{\normalfont}{propuesta}
\texto{191.9578}{663.1510}{11.9552}{\normalfont}{para}
\texto{218.9620}{663.1510}{11.9552}{\normalfont\bfseries}{Explora-Nicaragua}
\texto{328.8690}{663.1510}{11.9552}{\normalfont}{,}
\texto{336.2095}{663.1510}{11.9552}{\normalfont}{estableciendo}
\texto{408.0124}{663.1510}{11.9552}{\normalfont}{la}
\texto{421.2062}{663.1510}{11.9552}{\normalfont}{relación}
\texto{464.9790}{663.1510}{11.9552}{\normalfont}{entre}
\texto{494.7594}{663.1510}{11.9552}{\normalfont}{cada}
\texto{522.2731}{663.1510}{11.9552}{\normalfont}{una}
\texto{70.8660}{648.7050}{11.9552}{\normalfont}{de}
\texto{86.4699}{648.7050}{11.9552}{\normalfont}{ellas}
\texto{112.5442}{648.7050}{11.9552}{\normalfont}{y}
\texto{122.6320}{648.7050}{11.9552}{\normalfont}{los}
\texto{140.2516}{648.7050}{11.9552}{\normalfont}{componentes}
\texto{209.9002}{648.7050}{11.9552}{\normalfont}{del}
\texto{228.7679}{648.7050}{11.9552}{\normalfont}{sistema.}
\texto{70.8660}{610.0630}{14.3462}{\normalfont\bfseries}{2.2.1.}
\texto{124.6642}{610.0630}{14.3462}{\normalfont\bfseries}{Aplicaciones}
\texto{217.6061}{610.0630}{14.3462}{\normalfont\bfseries}{móviles}
\texto{275.2620}{610.0630}{14.3462}{\normalfont\bfseries}{híbridas}
\texto{88.7990}{586.5840}{11.9552}{\normalfont}{El}
\texto{103.2552}{586.5840}{11.9552}{\normalfont}{desarrollo}
\texto{155.9872}{586.5840}{11.9552}{\normalfont}{de}
\texto{170.9456}{586.5840}{11.9552}{\normalfont}{aplicaciones}
\texto{234.7337}{586.5840}{11.9552}{\normalfont}{móviles}
\texto{275.7735}{586.5840}{11.9552}{\normalfont}{multiplataforma}
\texto{361.6119}{586.5840}{11.9552}{\normalfont}{(Android}
\texto{411.3479}{586.5840}{11.9552}{\normalfont}{e}
\texto{419.8026}{586.5840}{11.9552}{\normalfont}{iOS)}
\texto{446.4555}{586.5840}{11.9552}{\normalfont}{se}
\texto{459.5274}{586.5840}{11.9552}{\normalfont}{ha}
\texto{475.1241}{586.5840}{11.9552}{\normalfont}{popularizado}
\texto{70.8660}{572.1380}{11.9552}{\normalfont}{mediante}
\texto{121.6852}{572.1380}{11.9552}{\normalfont}{frameworks}
\texto{184.2754}{572.1380}{11.9552}{\normalfont}{como}
\texto{215.2700}{572.1380}{11.9552}{\normalfont\bfseries}{React}
\texto{254.1459}{572.1380}{11.9552}{\normalfont\bfseries}{Native}
\texto{297.3580}{572.1380}{11.9552}{\normalfont}{y}
\texto{307.9670}{572.1380}{11.9552}{\normalfont\bfseries}{Flutter}
\texto{349.7200}{572.1380}{11.9552}{\normalfont}{.}
\texto{357.2876}{572.1380}{11.9552}{\normalfont}{React}
\texto{391.0420}{572.1380}{11.9552}{\normalfont}{Native,}
\texto{432.1009}{572.1380}{11.9552}{\normalfont}{basado}
\texto{471.6128}{572.1380}{11.9552}{\normalfont}{en}
\texto{487.6471}{572.1380}{11.9552}{\normalfont}{JavaScript}
\texto{70.8660}{557.6920}{11.9552}{\normalfont}{y}
\texto{80.9418}{557.6920}{11.9552}{\normalfont}{React,}
\texto{117.5295}{557.6920}{11.9552}{\normalfont}{permite}
\texto{160.7835}{557.6920}{11.9552}{\normalfont}{compartir}
\texto{214.7707}{557.6920}{11.9552}{\normalfont}{hasta}
\texto{246.0599}{557.6920}{11.9552}{\normalfont}{el}
\texto{258.4240}{557.6920}{11.9552}{\normalfont}{90}
\texto{272.1151}{557.6920}{11.9552}{\normalfont}{\%}
\texto{285.7798}{557.6920}{11.9552}{\normalfont}{del}
\texto{304.6475}{557.6920}{11.9552}{\normalfont}{código}
\texto{341.3978}{557.6920}{11.9552}{\normalfont}{fuente}
\texto{376.5270}{557.6920}{11.9552}{\normalfont}{entre}
\texto{406.1161}{557.6920}{11.9552}{\normalfont}{plataformas,}
\texto{473.5004}{557.6920}{11.9552}{\normalfont}{reduciendo}
\texto{532.6906}{557.6920}{11.9552}{\normalfont}{el}
\texto{70.8660}{543.2460}{11.9552}{\normalfont}{tiempo}
\texto{110.8275}{543.2460}{11.9552}{\normalfont}{y}
\texto{121.5250}{543.2460}{11.9552}{\normalfont}{costo}
\texto{152.1231}{543.2460}{11.9552}{\normalfont}{de}
\texto{168.3487}{543.2460}{11.9552}{\normalfont}{desarrollo}
\texto{222.3599}{543.2460}{11.9552}{\normalfont}{[5].}
\texto{242.4927}{543.2460}{11.9552}{\normalfont}{Flutter,}
\texto{286.5189}{543.2460}{11.9552}{\normalfont}{escrito}
\texto{324.2710}{543.2460}{11.9552}{\normalfont}{en}
\texto{340.4966}{543.2460}{11.9552}{\normalfont}{Dart,}
\texto{372.1779}{543.2460}{11.9552}{\normalfont}{ofrece}
\texto{406.2885}{543.2460}{11.9552}{\normalfont}{un}
\texto{423.8148}{543.2460}{11.9552}{\normalfont}{rendimiento}
\texto{489.1452}{543.2460}{11.9552}{\normalfont}{cercano}
\texto{532.0357}{543.2460}{11.9552}{\normalfont}{al}
\texto{70.8660}{528.8000}{11.9552}{\normalfont}{nativo}
\texto{106.6336}{528.8000}{11.9552}{\normalfont}{gracias}
\texto{145.7151}{528.8000}{11.9552}{\normalfont}{a}
\texto{155.4658}{528.8000}{11.9552}{\normalfont}{su}
\texto{170.4959}{528.8000}{11.9552}{\normalfont}{motor}
\texto{204.9603}{528.8000}{11.9552}{\normalfont}{de}
\texto{220.5762}{528.8000}{11.9552}{\normalfont}{renderizado}
\texto{283.6566}{528.8000}{11.9552}{\normalfont}{propio.}
\texto{88.7990}{513.7090}{11.9552}{\normalfont}{Para}
\texto{116.1360}{513.7090}{11.9552}{\normalfont\bfseries}{Explora-Nicaragua}
\texto{226.0430}{513.7090}{11.9552}{\normalfont}{,}
\texto{232.5825}{513.7090}{11.9552}{\normalfont}{la}
\texto{244.9753}{513.7090}{11.9552}{\normalfont}{app}
\texto{267.1354}{513.7090}{11.9552}{\normalfont}{móvil}
\texto{298.3911}{513.7090}{11.9552}{\normalfont}{del}
\texto{316.6491}{513.7090}{11.9552}{\normalfont}{turista}
\texto{353.8202}{513.7090}{11.9552}{\normalfont}{se}
\texto{366.9279}{513.7090}{11.9552}{\normalfont}{construye}
\texto{419.0717}{513.7090}{11.9552}{\normalfont}{con}
\texto{439.9191}{513.7090}{11.9552}{\normalfont}{un}
\texto{456.2260}{513.7090}{11.9552}{\normalfont}{framework}
\texto{513.1710}{513.7090}{11.9552}{\normalfont}{híbri-}
\texto{70.8660}{499.2630}{11.9552}{\normalfont}{do,}
\texto{89.1168}{499.2630}{11.9552}{\normalfont}{aprovechando}
\texto{161.3549}{499.2630}{11.9552}{\normalfont}{la}
\texto{173.1140}{499.2630}{11.9552}{\normalfont}{capacidad}
\texto{226.4844}{499.2630}{11.9552}{\normalfont}{de}
\texto{240.8450}{499.2630}{11.9552}{\normalfont}{lanzar}
\texto{274.7046}{499.2630}{11.9552}{\normalfont}{la}
\texto{286.4637}{499.2630}{11.9552}{\normalfont}{aplicación}
\texto{339.8460}{499.2630}{11.9552}{\normalfont}{simultáneamente}
\texto{428.4054}{499.2630}{11.9552}{\normalfont}{en}
\texto{442.7540}{499.2630}{11.9552}{\normalfont}{ambas}
\texto{477.6680}{499.2630}{11.9552}{\normalfont}{plataformas.}
\texto{70.8660}{484.8180}{11.9552}{\normalfont}{Esta}
\texto{96.8781}{484.8180}{11.9552}{\normalfont}{aplicación}
\texto{150.6430}{484.8180}{11.9552}{\normalfont}{forma}
\texto{183.2712}{484.8180}{11.9552}{\normalfont}{parte}
\texto{212.9727}{484.8180}{11.9552}{\normalfont}{de}
\texto{227.7158}{484.8180}{11.9552}{\normalfont}{la}
\texto{239.8575}{484.8180}{11.9552}{\normalfont}{visión}
\texto{272.5502}{484.8180}{11.9552}{\normalfont}{completa}
\texto{321.7626}{484.8180}{11.9552}{\normalfont}{del}
\texto{339.7576}{484.8180}{11.9552}{\normalfont}{ecosistema}
\texto{396.9034}{484.8180}{11.9552}{\normalfont}{y}
\texto{406.1185}{484.8180}{11.9552}{\normalfont}{se}
\texto{418.9751}{484.8180}{11.9552}{\normalfont}{contempla}
\texto{474.3683}{484.8180}{11.9552}{\normalfont}{como}
\texto{504.0698}{484.8180}{11.9552}{\normalfont}{trabajo}
\texto{70.8660}{470.3720}{11.9552}{\normalfont}{futuro;}
\texto{110.1556}{470.3720}{11.9552}{\normalfont}{su}
\texto{125.7834}{470.3720}{11.9552}{\normalfont}{arquitectura}
\texto{193.0386}{470.3720}{11.9552}{\normalfont}{separa}
\texto{230.1164}{470.3720}{11.9552}{\normalfont}{la}
\texto{243.7286}{470.3720}{11.9552}{\normalfont}{lógica}
\texto{277.4901}{470.3720}{11.9552}{\normalfont}{de}
\texto{293.7038}{470.3720}{11.9552}{\normalfont}{negocio}
\texto{336.2547}{470.3720}{11.9552}{\normalfont}{de}
\texto{352.4564}{470.3720}{11.9552}{\normalfont}{la}
\texto{366.0566}{470.3720}{11.9552}{\normalfont}{presentación,}
\texto{437.2929}{470.3720}{11.9552}{\normalfont}{de}
\texto{453.4946}{470.3720}{11.9552}{\normalfont}{modo}
\texto{486.2901}{470.3720}{11.9552}{\normalfont}{que}
\texto{508.6702}{470.3720}{11.9552}{\normalfont}{la}
\texto{522.2824}{470.3720}{11.9552}{\normalfont}{app}
\texto{70.8660}{455.9260}{11.9552}{\normalfont}{pueda}
\texto{105.3305}{455.9260}{11.9552}{\normalfont}{consumir}
\texto{155.4682}{455.9260}{11.9552}{\normalfont}{la}
\texto{168.4826}{455.9260}{11.9552}{\normalfont}{misma}
\texto{205.6131}{455.9260}{11.9552}{\normalfont}{API}
\texto{230.4727}{455.9260}{11.9552}{\normalfont}{REST}
\texto{265.9199}{455.9260}{11.9552}{\normalfont}{que}
\texto{287.7022}{455.9260}{11.9552}{\normalfont}{los}
\texto{305.3338}{455.9260}{11.9552}{\normalfont}{portales}
\texto{349.9410}{455.9260}{11.9552}{\normalfont}{de}
\texto{365.5569}{455.9260}{11.9552}{\normalfont}{escritorio}
\texto{416.3450}{455.9260}{11.9552}{\normalfont}{del}
\texto{435.2007}{455.9260}{11.9552}{\normalfont}{prototipo}
\texto{486.9094}{455.9260}{11.9552}{\normalfont}{inicial.}
\texto{70.8660}{417.2840}{14.3462}{\normalfont\bfseries}{2.2.2.}
\texto{124.6642}{417.2840}{14.3462}{\normalfont\bfseries}{Aplicaciones}
\texto{217.6061}{417.2840}{14.3462}{\normalfont\bfseries}{de}
\texto{239.3162}{417.2840}{14.3462}{\normalfont\bfseries}{escritorio}
\texto{310.1664}{417.2840}{14.3462}{\normalfont\bfseries}{y}
\texto{324.0650}{417.2840}{14.3462}{\normalfont\bfseries}{frameworks}
\texto{410.3731}{417.2840}{14.3462}{\normalfont\bfseries}{.NET}
\texto{88.7990}{393.8050}{11.9552}{\normalfont}{Las}
\texto{109.7875}{393.8050}{11.9552}{\normalfont}{aplicaciones}
\texto{173.5398}{393.8050}{11.9552}{\normalfont}{de}
\texto{188.4504}{393.8050}{11.9552}{\normalfont}{escritorio}
\texto{238.5451}{393.8050}{11.9552}{\normalfont}{siguen}
\texto{273.6814}{393.8050}{11.9552}{\normalfont}{siendo}
\texto{308.8177}{393.8050}{11.9552}{\normalfont}{relevantes}
\texto{362.1714}{393.8050}{11.9552}{\normalfont}{para}
\texto{388.1381}{393.8050}{11.9552}{\normalfont}{tareas}
\texto{421.9737}{393.8050}{11.9552}{\normalfont}{administrativas}
\texto{503.6803}{393.8050}{11.9552}{\normalfont}{intensi-}
\texto{70.8660}{379.3590}{11.9552}{\normalfont}{vas}
\texto{90.1928}{379.3590}{11.9552}{\normalfont}{en}
\texto{105.2468}{379.3590}{11.9552}{\normalfont}{datos,}
\texto{139.2139}{379.3590}{11.9552}{\normalfont}{como}
\texto{169.2143}{379.3590}{11.9552}{\normalfont}{la}
\texto{181.6668}{379.3590}{11.9552}{\normalfont}{gestión}
\texto{220.8368}{379.3590}{11.9552}{\normalfont}{de}
\texto{235.8788}{379.3590}{11.9552}{\normalfont}{catálogos,}
\texto{289.3688}{379.3590}{11.9552}{\normalfont}{disponibilidad}
\texto{364.6316}{379.3590}{11.9552}{\normalfont}{y}
\texto{374.1575}{379.3590}{11.9552}{\normalfont}{reportes.}
\texto{422.1050}{379.3590}{11.9552}{\normalfont}{El}
\texto{436.6688}{379.3590}{11.9552}{\normalfont}{framework}
\texto{493.6570}{379.3590}{11.9552}{\normalfont\bfseries}{.NET}
\texto{529.4290}{379.3590}{11.9552}{\normalfont}{de}
\texto{70.8660}{364.9130}{11.9552}{\normalfont}{Microsoft}
\texto{123.3087}{364.9130}{11.9552}{\normalfont}{y}
\texto{133.7551}{364.9130}{11.9552}{\normalfont}{su}
\texto{149.1319}{364.9130}{11.9552}{\normalfont}{tecnología}
\texto{204.9290}{364.9130}{11.9552}{\normalfont\bfseries}{Windows}
\texto{263.7641}{364.9130}{11.9552}{\normalfont\bfseries}{F}
\texto{271.1070}{364.9130}{11.9552}{\normalfont\bfseries}{orms}
\texto{304.7286}{364.9130}{11.9552}{\normalfont\bfseries}{(WinF}
\texto{342.4018}{364.9130}{11.9552}{\normalfont\bfseries}{orms)}
\texto{380.6220}{364.9130}{11.9552}{\normalfont}{permiten}
\texto{430.7262}{364.9130}{11.9552}{\normalfont}{construir}
\texto{480.5842}{364.9130}{11.9552}{\normalfont}{aplicaciones}
\texto{70.8660}{350.4670}{11.9552}{\normalfont}{de}
\texto{86.6373}{350.4670}{11.9552}{\normalfont}{escritorio}
\texto{137.5928}{350.4670}{11.9552}{\normalfont}{para}
\texto{164.4322}{350.4670}{11.9552}{\normalfont}{Windows}
\texto{215.3852}{350.4670}{11.9552}{\normalfont}{con}
\texto{237.0218}{350.4670}{11.9552}{\normalfont}{controles}
\texto{286.3537}{350.4670}{11.9552}{\normalfont}{ricos,}
\texto{317.1479}{350.4670}{11.9552}{\normalfont}{enlazado}
\texto{365.4493}{350.4670}{11.9552}{\normalfont}{de}
\texto{381.2206}{350.4670}{11.9552}{\normalfont}{datos}
\texto{412.6771}{350.4670}{11.9552}{\normalfont}{(data}
\texto{444.0571}{350.4670}{11.9552}{\normalfont}{binding)}
\texto{491.0459}{350.4670}{11.9552}{\normalfont}{y}
\texto{501.3010}{350.4670}{11.9552}{\normalfont}{un}
\texto{518.3730}{350.4670}{11.9552}{\normalfont}{ciclo}
\texto{70.8660}{336.0220}{11.9552}{\normalfont}{de}
\texto{87.1992}{336.0220}{11.9552}{\normalfont}{desarrollo}
\texto{141.3180}{336.0220}{11.9552}{\normalfont}{maduro}
\texto{184.9664}{336.0220}{11.9552}{\normalfont}{y}
\texto{195.7715}{336.0220}{11.9552}{\normalfont}{documentado.}
\texto{271.9381}{336.0220}{11.9552}{\normalfont}{WinForms}
\texto{329.7965}{336.0220}{11.9552}{\normalfont}{es}
\texto{344.2432}{336.0220}{11.9552}{\normalfont}{especialmente}
\texto{419.1736}{336.0220}{11.9552}{\normalfont}{adecuado}
\texto{471.2767}{336.0220}{11.9552}{\normalfont}{para}
\texto{498.6661}{336.0220}{11.9552}{\normalfont}{sistemas}
\texto{70.8660}{321.5760}{11.9552}{\normalfont}{de}
\texto{86.5177}{321.5760}{11.9552}{\normalfont}{gestión}
\texto{126.3094}{321.5760}{11.9552}{\normalfont}{internos}
\texto{170.9693}{321.5760}{11.9552}{\normalfont}{donde}
\texto{205.4935}{321.5760}{11.9552}{\normalfont}{se}
\texto{219.2587}{321.5760}{11.9552}{\normalfont}{valora}
\texto{254.1129}{321.5760}{11.9552}{\normalfont}{la}
\texto{267.1632}{321.5760}{11.9552}{\normalfont}{productividad}
\texto{342.6579}{321.5760}{11.9552}{\normalfont}{del}
\texto{361.5734}{321.5760}{11.9552}{\normalfont}{desarrollador}
\texto{431.9202}{321.5760}{11.9552}{\normalfont}{y}
\texto{442.0558}{321.5760}{11.9552}{\normalfont}{la}
\texto{455.1061}{321.5760}{11.9552}{\normalfont}{estabilidad}
\texto{514.4087}{321.5760}{11.9552}{\normalfont}{sobre}
\texto{70.8660}{307.1300}{11.9552}{\normalfont}{la}
\texto{83.8685}{307.1300}{11.9552}{\normalfont}{novedad}
\texto{129.7191}{307.1300}{11.9552}{\normalfont}{visual.}
\texto{88.7990}{292.0390}{11.9552}{\normalfont}{Para}
\texto{117.1760}{292.0390}{11.9552}{\normalfont\bfseries}{Explora-Nicaragua}
\texto{227.0830}{292.0390}{11.9552}{\normalfont}{,}
\texto{234.6626}{292.0390}{11.9552}{\normalfont}{el}
\texto{247.5010}{292.0390}{11.9552}{\normalfont\bfseries}{backoffice}
\texto{308.8551}{292.0390}{11.9552}{\normalfont\bfseries}{de}
\texto{327.4490}{292.0390}{11.9552}{\normalfont\bfseries}{operadores}
\texto{395.6570}{292.0390}{11.9552}{\normalfont}{y}
\texto{406.1632}{292.0390}{11.9552}{\normalfont}{el}
\texto{419.0010}{292.0390}{11.9552}{\normalfont\bfseries}{panel}
\texto{455.3293}{292.0390}{11.9552}{\normalfont\bfseries}{de}
\texto{473.9112}{292.0390}{11.9552}{\normalfont\bfseries}{administra-}
\texto{70.8660}{277.5930}{11.9552}{\normalfont}{\parbox[t]{470.268bp}{\fontsize{11.9552bp}{14.446bp}\selectfont ción utiliza .NET 10 y Windows Forms con una arquitectura por capas. Los formularios coordinan la interacción del usuario y consumen los servicios de la API REST mediante HttpClient e IHttpClientFactory. El servidor concentra las reglas de negocio y la persistencia de los datos.}}
\texto{70.8660}{195.6130}{14.3462}{\normalfont\bfseries}{2.2.3.}
\texto{124.6642}{195.6130}{14.3462}{\normalfont\bfseries}{Consumo}
\texto{196.2633}{195.6130}{14.3462}{\normalfont\bfseries}{de}
\texto{218.6620}{195.6130}{14.3462}{\normalfont\bfseries}{APIs}
\texto{260.3205}{195.6130}{14.3462}{\normalfont\bfseries}{REST}
\texto{309.2109}{195.6130}{14.3462}{\normalfont\bfseries}{y}
\texto{323.7838}{195.6130}{14.3462}{\normalfont\bfseries}{arquitecturas}
\texto{423.1197}{195.6130}{14.3462}{\normalfont\bfseries}{orientadas}
\texto{501.5203}{195.6130}{14.3462}{\normalfont\bfseries}{a}
\texto{515.4347}{195.6130}{14.3462}{\normalfont\bfseries}{ser-}
\texto{124.6640}{177.6810}{14.3462}{\normalfont\bfseries}{vicios}
\texto{88.7990}{154.2010}{11.9552}{\normalfont}{Las}
\texto{110.1582}{154.2010}{11.9552}{\normalfont}{arquitecturas}
\texto{181.1218}{154.2010}{11.9552}{\normalfont}{orientadas}
\texto{237.1176}{154.2010}{11.9552}{\normalfont}{a}
\texto{246.5574}{154.2010}{11.9552}{\normalfont}{servicios,}
\texto{296.1237}{154.2010}{11.9552}{\normalfont}{materializadas}
\texto{373.2658}{154.2010}{11.9552}{\normalfont}{en}
\texto{388.6070}{154.2010}{11.9552}{\normalfont\bfseries}{APIs}
\texto{422.3780}{154.2010}{11.9552}{\normalfont\bfseries}{REST}
\texto{458.0670}{154.2010}{11.9552}{\normalfont}{,}
\texto{464.8934}{154.2010}{11.9552}{\normalfont}{son}
\texto{485.4540}{154.2010}{11.9552}{\normalfont}{el}
\texto{497.4952}{154.2010}{11.9552}{\normalfont}{estándar}
\texto{70.8660}{139.7560}{11.9552}{\normalfont}{para}
\texto{99.0444}{139.7560}{11.9552}{\normalfont}{integrar}
\texto{144.4718}{139.7560}{11.9552}{\normalfont}{múltiples}
\texto{196.4673}{139.7560}{11.9552}{\normalfont}{clientes}
\texto{239.3578}{139.7560}{11.9552}{\normalfont}{sobre}
\texto{271.5149}{139.7560}{11.9552}{\normalfont}{una}
\texto{295.7911}{139.7560}{11.9552}{\normalfont}{misma}
\texto{334.4519}{139.7560}{11.9552}{\normalfont}{lógica}
\texto{369.1458}{139.7560}{11.9552}{\normalfont}{de}
\texto{386.2800}{139.7560}{11.9552}{\normalfont}{negocio}
\texto{429.7515}{139.7560}{11.9552}{\normalfont}{[14].}
\texto{456.6423}{139.7560}{11.9552}{\normalfont}{Una}
\texto{483.2020}{139.7560}{11.9552}{\normalfont}{API}
\texto{509.5919}{139.7560}{11.9552}{\normalfont}{REST}
\texto{70.8660}{125.3100}{11.9552}{\normalfont}{expone}
\texto{110.2321}{125.3100}{11.9552}{\normalfont}{recursos}
\texto{154.9326}{125.3100}{11.9552}{\normalfont}{mediante}
\texto{205.0225}{125.3100}{11.9552}{\normalfont}{los}
\texto{222.3432}{125.3100}{11.9552}{\normalfont}{métodos}
\texto{268.6026}{125.3100}{11.9552}{\normalfont}{HTTP}
\texto{305.8502}{125.3100}{11.9552}{\normalfont}{(GET,}
\texto{342.8576}{125.3100}{11.9552}{\normalfont}{POST,}
\texto{381.7359}{125.3100}{11.9552}{\normalfont}{PUT,}
\texto{413.7806}{125.3100}{11.9552}{\normalfont}{DELETE)}
\texto{470.5343}{125.3100}{11.9552}{\normalfont}{y}
\texto{480.3113}{125.3100}{11.9552}{\normalfont}{permite}
\texto{523.2544}{125.3100}{11.9552}{\normalfont}{que}
\texto{70.8660}{110.8640}{11.9552}{\normalfont}{aplicaciones}
\texto{135.3117}{110.8640}{11.9552}{\normalfont}{distintas}
\texto{182.5897}{110.8640}{11.9552}{\normalfont}{(móvil,}
\texto{222.2595}{110.8640}{11.9552}{\normalfont}{web}
\texto{245.9953}{110.8640}{11.9552}{\normalfont}{y}
\texto{256.0831}{110.8640}{11.9552}{\normalfont}{escritorio)}
\texto{311.4238}{110.8640}{11.9552}{\normalfont}{consuman}
\texto{366.1260}{110.8640}{11.9552}{\normalfont}{los}
\texto{383.7455}{110.8640}{11.9552}{\normalfont}{mismos}
\texto{425.4931}{110.8640}{11.9552}{\normalfont}{servicios.}
\texto{88.7990}{95.7730}{11.9552}{\normalfont}{Para}
\texto{117.4300}{95.7730}{11.9552}{\normalfont\bfseries}{Explora-Nicaragua}
\texto{227.3370}{95.7730}{11.9552}{\normalfont}{,}
\texto{235.1677}{95.7730}{11.9552}{\normalfont}{la}
\texto{248.8635}{95.7730}{11.9552}{\normalfont}{API}
\texto{274.4166}{95.7730}{11.9552}{\normalfont}{REST}
\texto{310.5332}{95.7730}{11.9552}{\normalfont}{es}
\texto{324.9440}{95.7730}{11.9552}{\normalfont}{el}
\texto{337.9895}{95.7730}{11.9552}{\normalfont}{elemento}
\texto{387.7710}{95.7730}{11.9552}{\normalfont}{central}
\texto{427.1586}{95.7730}{11.9552}{\normalfont}{que}
\texto{449.6224}{95.7730}{11.9552}{\normalfont}{coordina}
\texto{498.1103}{95.7730}{11.9552}{\normalfont}{los}
\texto{516.4233}{95.7730}{11.9552}{\normalfont}{com-}
\texto{70.8660}{81.3270}{11.9552}{\normalfont}{ponentes}
\texto{121.6158}{81.3270}{11.9552}{\normalfont}{del}
\texto{142.3844}{81.3270}{11.9552}{\normalfont}{ecosistema:}
\texto{205.5676}{81.3270}{11.9552}{\normalfont}{en}
\texto{223.0963}{81.3270}{11.9552}{\normalfont}{el}
\texto{237.3613}{81.3270}{11.9552}{\normalfont}{prototipo}
\texto{290.9828}{81.3270}{11.9552}{\normalfont}{inicial,}
\texto{330.6238}{81.3270}{11.9552}{\normalfont}{el}
\texto{344.8887}{81.3270}{11.9552}{\normalfont}{backoffice}
\texto{399.8062}{81.3270}{11.9552}{\normalfont}{de}
\texto{417.3349}{81.3270}{11.9552}{\normalfont}{escritorio}
\texto{470.0358}{81.3270}{11.9552}{\normalfont}{y}
\texto{482.0364}{81.3270}{11.9552}{\normalfont}{el}
\texto{496.3014}{81.3270}{11.9552}{\normalfont}{panel}
\texto{529.4388}{81.3270}{11.9552}{\normalfont}{de}
\texto{70.8660}{66.8810}{11.9552}{\normalfont}{administración}
\texto{149.7918}{66.8810}{11.9552}{\normalfont}{de}
\texto{164.9176}{66.8810}{11.9552}{\normalfont}{escritorio;}
\texto{218.4912}{66.8810}{11.9552}{\normalfont}{en}
\texto{233.6169}{66.8810}{11.9552}{\normalfont}{la}
\texto{246.1531}{66.8810}{11.9552}{\normalfont}{visión}
\texto{279.2284}{66.8810}{11.9552}{\normalfont}{completa,}
\texto{332.0871}{66.8810}{11.9552}{\normalfont}{también}
\texto{376.8067}{66.8810}{11.9552}{\normalfont}{la}
\texto{389.3310}{66.8810}{11.9552}{\normalfont}{app}
\texto{411.6227}{66.8810}{11.9552}{\normalfont}{móvil.}
\texto{446.2617}{66.8810}{11.9552}{\normalfont}{Esto}
\texto{472.6803}{66.8810}{11.9552}{\normalfont}{garantiza}
\texto{523.2532}{66.8810}{11.9552}{\normalfont}{que}
\end{pagina}
% Página original 25
\begin{pagina}{25}
\texto{70.8660}{746.3890}{11.9552}{\normalfont}{2.2.}
\texto{96.8710}{746.3890}{11.9552}{\normalfont}{TENDENCIAS}
\texto{179.6105}{746.3890}{11.9552}{\normalfont}{TECNOLÓGICAS}
\texto{279.8226}{746.3890}{11.9552}{\normalfont}{RELACIONADAS}
\texto{529.4280}{746.3890}{11.9552}{\normalfont}{15}
\texto{70.8660}{705.1940}{11.9552}{\normalfont}{el}
\texto{83.3138}{705.1940}{11.9552}{\normalfont}{modelo}
\texto{124.0619}{705.1940}{11.9552}{\normalfont}{de}
\texto{139.7614}{705.1940}{11.9552}{\normalfont}{datos}
\texto{171.1343}{705.1940}{11.9552}{\normalfont}{sea}
\texto{190.8125}{705.1940}{11.9552}{\normalfont}{único}
\texto{222.1208}{705.1940}{11.9552}{\normalfont}{y}
\texto{232.3042}{705.1940}{11.9552}{\normalfont}{consistente}
\texto{292.0348}{705.1940}{11.9552}{\normalfont}{en}
\texto{307.7344}{705.1940}{11.9552}{\normalfont}{tiempo}
\texto{347.1818}{705.1940}{11.9552}{\normalfont}{real,}
\texto{373.2871}{705.1940}{11.9552}{\normalfont}{evitando}
\texto{421.1916}{705.1940}{11.9552}{\normalfont}{la}
\texto{434.2897}{705.1940}{11.9552}{\normalfont}{duplicación}
\texto{496.1651}{705.1940}{11.9552}{\normalfont}{de}
\texto{511.8766}{705.1940}{11.9552}{\normalfont}{lógica}
\texto{70.8660}{690.7480}{11.9552}{\normalfont}{de}
\texto{86.4699}{690.7480}{11.9552}{\normalfont}{negocio}
\texto{128.4112}{690.7480}{11.9552}{\normalfont}{en}
\texto{144.0270}{690.7480}{11.9552}{\normalfont}{cada}
\texto{171.3375}{690.7480}{11.9552}{\normalfont}{cliente.}
\texto{70.8660}{652.8680}{14.3462}{\normalfont\bfseries}{2.2.4.}
\texto{124.6642}{652.8680}{14.3462}{\normalfont\bfseries}{Control}
\texto{183.6989}{652.8680}{14.3462}{\normalfont\bfseries}{de}
\texto{205.4090}{652.8680}{14.3462}{\normalfont\bfseries}{acceso}
\texto{254.7814}{652.8680}{14.3462}{\normalfont\bfseries}{basado}
\texto{308.2224}{652.8680}{14.3462}{\normalfont\bfseries}{en}
\texto{329.9325}{652.8680}{14.3462}{\normalfont\bfseries}{roles}
\texto{368.1881}{652.8680}{14.3462}{\normalfont\bfseries}{(RBAC)}
\texto{88.7990}{629.6400}{11.9552}{\normalfont}{El}
\texto{105.7550}{629.6400}{11.9552}{\normalfont\bfseries}{control}
\texto{153.2829}{629.6400}{11.9552}{\normalfont\bfseries}{de}
\texto{173.4310}{629.6400}{11.9552}{\normalfont\bfseries}{acceso}
\texto{216.6311}{629.6400}{11.9552}{\normalfont\bfseries}{basado}
\texto{263.2217}{629.6400}{11.9552}{\normalfont\bfseries}{en}
\texto{283.3818}{629.6400}{11.9552}{\normalfont\bfseries}{roles}
\texto{316.4750}{629.6400}{11.9552}{\normalfont}{(}
\texto{321.0270}{629.6400}{11.9552}{\normalfont\itshape}{Role-Based}
\texto{383.0601}{629.6400}{11.9552}{\normalfont\itshape}{Access}
\texto{421.4949}{629.6400}{11.9552}{\normalfont\itshape}{Control}
\texto{459.6100}{629.6400}{11.9552}{\normalfont}{,}
\texto{468.5525}{629.6400}{11.9552}{\normalfont}{RBAC)}
\texto{512.6074}{629.6400}{11.9552}{\normalfont}{es}
\texto{528.1181}{629.6400}{11.9552}{\normalfont}{un}
\texto{70.8660}{615.1950}{11.9552}{\normalfont}{modelo}
\texto{110.3349}{615.1950}{11.9552}{\normalfont}{de}
\texto{124.7792}{615.1950}{11.9552}{\normalfont}{seguridad}
\texto{176.3467}{615.1950}{11.9552}{\normalfont}{que}
\texto{196.9694}{615.1950}{11.9552}{\normalfont}{asigna}
\texto{231.6276}{615.1950}{11.9552}{\normalfont}{permisos}
\texto{279.0419}{615.1950}{11.9552}{\normalfont}{a}
\texto{287.6329}{615.1950}{11.9552}{\normalfont}{los}
\texto{304.0809}{615.1950}{11.9552}{\normalfont}{usuarios}
\texto{348.5709}{615.1950}{11.9552}{\normalfont}{según}
\texto{379.9773}{615.1950}{11.9552}{\normalfont}{su}
\texto{393.8357}{615.1950}{11.9552}{\normalfont}{rol}
\texto{410.2191}{615.1950}{11.9552}{\normalfont}{dentro}
\texto{445.7906}{615.1950}{11.9552}{\normalfont}{del}
\texto{463.4867}{615.1950}{11.9552}{\normalfont}{sistema.}
\texto{507.3264}{615.1950}{11.9552}{\normalfont}{RBAC}
\texto{70.8660}{600.7490}{11.9552}{\normalfont}{reduce}
\texto{107.9319}{600.7490}{11.9552}{\normalfont}{la}
\texto{120.9463}{600.7490}{11.9552}{\normalfont}{complejidad}
\texto{186.3150}{600.7490}{11.9552}{\normalfont}{administrativa}
\texto{264.1099}{600.7490}{11.9552}{\normalfont}{y}
\texto{274.1977}{600.7490}{11.9552}{\normalfont}{minimiza}
\texto{324.9212}{600.7490}{11.9552}{\normalfont}{riesgos}
\texto{362.7785}{600.7490}{11.9552}{\normalfont}{al}
\texto{375.7929}{600.7490}{11.9552}{\normalfont}{garantizar}
\texto{431.4085}{600.7490}{11.9552}{\normalfont}{que}
\texto{453.2029}{600.7490}{11.9552}{\normalfont}{cada}
\texto{480.5253}{600.7490}{11.9552}{\normalfont}{usuario}
\texto{521.5699}{600.7490}{11.9552}{\normalfont}{solo}
\texto{70.8660}{586.3030}{11.9552}{\normalfont}{puede}
\texto{104.6801}{586.3030}{11.9552}{\normalfont}{realizar}
\texto{146.2985}{586.3030}{11.9552}{\normalfont}{las}
\texto{163.9301}{586.3030}{11.9552}{\normalfont}{acciones}
\texto{209.5152}{586.3030}{11.9552}{\normalfont}{que}
\texto{231.2976}{586.3030}{11.9552}{\normalfont}{su}
\texto{246.3277}{586.3030}{11.9552}{\normalfont}{rol}
\texto{263.8827}{586.3030}{11.9552}{\normalfont}{autoriza}
\texto{309.4153}{586.3030}{11.9552}{\normalfont}{[7].}
\texto{88.7990}{571.3350}{11.9552}{\normalfont}{Para}
\texto{116.9980}{571.3350}{11.9552}{\normalfont\bfseries}{Explora-Nicaragua}
\texto{226.9050}{571.3350}{11.9552}{\normalfont}{,}
\texto{234.3053}{571.3350}{11.9552}{\normalfont}{el}
\texto{246.9204}{571.3350}{11.9552}{\normalfont}{RBAC}
\texto{284.8805}{571.3350}{11.9552}{\normalfont}{se}
\texto{298.8610}{571.3350}{11.9552}{\normalfont}{implementa}
\texto{362.5200}{571.3350}{11.9552}{\normalfont}{en}
\texto{378.3869}{571.3350}{11.9552}{\normalfont}{dos}
\texto{399.5094}{571.3350}{11.9552}{\normalfont}{capas:}
\texto{434.9518}{571.3350}{11.9552}{\normalfont}{el}
\texto{447.5669}{571.3350}{11.9552}{\normalfont}{backend}
\texto{493.0182}{571.3350}{11.9552}{\normalfont}{valida}
\texto{527.4133}{571.3350}{11.9552}{\normalfont}{los}
\texto{70.8660}{556.8900}{11.9552}{\normalfont}{permisos}
\texto{118.8422}{556.8900}{11.9552}{\normalfont}{en}
\texto{133.8603}{556.8900}{11.9552}{\normalfont}{cada}
\texto{160.5850}{556.8900}{11.9552}{\normalfont}{endpoint}
\texto{208.7716}{556.8900}{11.9552}{\normalfont}{(fuente}
\texto{247.8436}{556.8900}{11.9552}{\normalfont}{de}
\texto{262.8617}{556.8900}{11.9552}{\normalfont}{verdad)}
\texto{305.1975}{556.8900}{11.9552}{\normalfont}{y}
\texto{314.6756}{556.8900}{11.9552}{\normalfont}{la}
\texto{327.0922}{556.8900}{11.9552}{\normalfont}{interfaz}
\texto{368.7776}{556.8900}{11.9552}{\normalfont}{de}
\texto{383.7838}{556.8900}{11.9552}{\normalfont}{cada}
\texto{410.5084}{556.8900}{11.9552}{\normalfont}{aplicación}
\texto{464.5483}{556.8900}{11.9552}{\normalfont}{oculta}
\texto{499.4001}{556.8900}{11.9552}{\normalfont}{o}
\texto{508.5650}{556.8900}{11.9552}{\normalfont}{desha-}
\texto{70.8660}{542.4440}{11.9552}{\normalfont}{bilita}
\texto{100.4001}{542.4440}{11.9552}{\normalfont}{elementos}
\texto{153.1010}{542.4440}{11.9552}{\normalfont}{según}
\texto{184.6628}{542.4440}{11.9552}{\normalfont}{el}
\texto{195.9987}{542.4440}{11.9552}{\normalfont}{rol.}
\texto{215.7893}{542.4440}{11.9552}{\normalfont}{En}
\texto{233.1387}{542.4440}{11.9552}{\normalfont}{el}
\texto{244.4746}{542.4440}{11.9552}{\normalfont}{panel}
\texto{274.6711}{542.4440}{11.9552}{\normalfont}{de}
\texto{289.2588}{542.4440}{11.9552}{\normalfont}{administración}
\texto{367.6467}{542.4440}{11.9552}{\normalfont}{se}
\texto{380.3479}{542.4440}{11.9552}{\normalfont}{definen}
\texto{419.6494}{542.4440}{11.9552}{\normalfont}{roles}
\texto{446.0082}{542.4440}{11.9552}{\normalfont}{como}
\texto{475.5450}{542.4440}{11.9552}{\normalfont\itshape}{SuperAdmin}
\texto{537.8830}{542.4440}{11.9552}{\normalfont}{,}
\texto{70.8660}{527.9980}{11.9552}{\normalfont\itshape}{Soporte}
\texto{108.9730}{527.9980}{11.9552}{\normalfont}{,}
\texto{115.4830}{527.9980}{11.9552}{\normalfont\itshape}{Analista}
\texto{161.0990}{527.9980}{11.9552}{\normalfont}{y}
\texto{170.5350}{527.9980}{11.9552}{\normalfont\itshape}{Moderador}
\texto{226.0970}{527.9980}{11.9552}{\normalfont}{;}
\texto{232.6006}{527.9980}{11.9552}{\normalfont}{en}
\texto{247.5709}{527.9980}{11.9552}{\normalfont}{el}
\texto{259.2775}{527.9980}{11.9552}{\normalfont}{backoffice}
\texto{311.6365}{527.9980}{11.9552}{\normalfont}{de}
\texto{326.6068}{527.9980}{11.9552}{\normalfont}{operadores,}
\texto{388.1282}{527.9980}{11.9552}{\normalfont}{roles}
\texto{414.8696}{527.9980}{11.9552}{\normalfont}{como}
\texto{444.7890}{527.9980}{11.9552}{\normalfont\itshape}{Propietario}
\texto{502.3350}{527.9980}{11.9552}{\normalfont}{,}
\texto{508.8450}{527.9980}{11.9552}{\normalfont\itshape}{Editor}
\texto{70.8660}{513.5520}{11.9552}{\normalfont}{y}
\texto{80.9460}{513.5520}{11.9552}{\normalfont\itshape}{Solo}
\texto{106.6497}{513.5520}{11.9552}{\normalfont\itshape}{lectura}
\texto{140.2740}{513.5520}{11.9552}{\normalfont}{.}
\texto{70.8660}{475.6730}{14.3462}{\normalfont\bfseries}{2.2.5.}
\texto{124.6642}{475.6730}{14.3462}{\normalfont\bfseries}{Pasarelas}
\texto{195.3207}{475.6730}{14.3462}{\normalfont\bfseries}{de}
\texto{217.0308}{475.6730}{14.3462}{\normalfont\bfseries}{pago}
\texto{255.3624}{475.6730}{14.3462}{\normalfont\bfseries}{digital}
\texto{305.3503}{475.6730}{14.3462}{\normalfont\bfseries}{en}
\texto{327.0604}{475.6730}{14.3462}{\normalfont\bfseries}{Centroamérica}
\texto{88.7990}{452.4450}{11.9552}{\normalfont}{El}
\texto{103.3150}{452.4450}{11.9552}{\normalfont}{ecosistema}
\texto{160.7239}{452.4450}{11.9552}{\normalfont}{de}
\texto{175.7300}{452.4450}{11.9552}{\normalfont}{pagos}
\texto{207.7102}{452.4450}{11.9552}{\normalfont}{en}
\texto{222.7164}{452.4450}{11.9552}{\normalfont}{línea}
\texto{250.0794}{452.4450}{11.9552}{\normalfont}{en}
\texto{265.0856}{452.4450}{11.9552}{\normalfont}{Centroamérica}
\texto{342.8542}{452.4450}{11.9552}{\normalfont}{ha}
\texto{358.5227}{452.4450}{11.9552}{\normalfont}{evolucionado}
\texto{427.4994}{452.4450}{11.9552}{\normalfont}{hacia}
\texto{457.4758}{452.4450}{11.9552}{\normalfont}{la}
\texto{469.8806}{452.4450}{11.9552}{\normalfont}{integración}
\texto{529.4390}{452.4450}{11.9552}{\normalfont}{de}
\texto{70.8660}{437.9990}{11.9552}{\normalfont}{pasarelas}
\texto{120.5422}{437.9990}{11.9552}{\normalfont}{internacionales}
\texto{199.7478}{437.9990}{11.9552}{\normalfont}{y}
\texto{209.2977}{437.9990}{11.9552}{\normalfont}{locales.}
\texto{249.4870}{437.9990}{11.9552}{\normalfont\bfseries}{Stripe}
\texto{288.7750}{437.9990}{11.9552}{\normalfont}{es}
\texto{301.9664}{437.9990}{11.9552}{\normalfont}{ampliamente}
\texto{370.7016}{437.9990}{11.9552}{\normalfont}{adoptada}
\texto{421.5495}{437.9990}{11.9552}{\normalfont}{por}
\texto{442.1530}{437.9990}{11.9552}{\normalfont}{su}
\texto{456.6571}{437.9990}{11.9552}{\normalfont}{API}
\texto{480.9907}{437.9990}{11.9552}{\normalfont}{flexible}
\texto{520.4572}{437.9990}{11.9552}{\normalfont}{y}
\texto{530.0070}{437.9990}{11.9552}{\normalfont}{su}
\texto{70.8660}{423.5530}{11.9552}{\normalfont}{soporte}
\texto{111.3009}{423.5530}{11.9552}{\normalfont}{para}
\texto{137.0404}{423.5530}{11.9552}{\normalfont}{múltiples}
\texto{186.5852}{423.5530}{11.9552}{\normalfont}{métodos}
\texto{232.2229}{423.5530}{11.9552}{\normalfont}{de}
\texto{246.9183}{423.5530}{11.9552}{\normalfont}{pago}
\texto{273.9585}{423.5530}{11.9552}{\normalfont}{[6].}
\texto{292.5435}{423.5530}{11.9552}{\normalfont}{La}
\texto{308.6878}{423.5530}{11.9552}{\normalfont}{integración}
\texto{367.9234}{423.5530}{11.9552}{\normalfont}{de}
\texto{382.6068}{423.5530}{11.9552}{\normalfont}{Stripe}
\texto{416.4854}{423.5530}{11.9552}{\normalfont}{en}
\texto{431.2280}{423.5530}{11.9552}{\normalfont\bfseries}{Explora-Nicaragua}
\texto{70.8660}{409.1070}{11.9552}{\normalfont}{se}
\texto{85.6952}{409.1070}{11.9552}{\normalfont}{sustenta}
\texto{132.7963}{409.1070}{11.9552}{\normalfont}{en}
\texto{149.5240}{409.1070}{11.9552}{\normalfont}{la}
\texto{163.6503}{409.1070}{11.9552}{\normalfont}{certificación}
\texto{229.8660}{409.1070}{11.9552}{\normalfont\bfseries}{PCI}
\texto{259.5388}{409.1070}{11.9552}{\normalfont\bfseries}{DSS}
\texto{289.8160}{409.1070}{11.9552}{\normalfont}{y}
\texto{301.0037}{409.1070}{11.9552}{\normalfont}{en}
\texto{317.7314}{409.1070}{11.9552}{\normalfont}{el}
\texto{331.2073}{409.1070}{11.9552}{\normalfont}{uso}
\texto{353.2025}{409.1070}{11.9552}{\normalfont}{de}
\texto{369.9380}{409.1070}{11.9552}{\normalfont\bfseries}{SSL/TLS}
\texto{429.6140}{409.1070}{11.9552}{\normalfont}{para}
\texto{457.3859}{409.1070}{11.9552}{\normalfont}{el}
\texto{470.8618}{409.1070}{11.9552}{\normalfont}{cifrado}
\texto{510.6774}{409.1070}{11.9552}{\normalfont}{de}
\texto{527.4051}{409.1070}{11.9552}{\normalfont}{las}
\texto{70.8660}{394.6610}{11.9552}{\normalfont}{transacciones.}
\texto{70.8660}{356.7820}{14.3462}{\normalfont\bfseries}{2.2.6.}
\texto{124.6642}{356.7820}{14.3462}{\normalfont\bfseries}{Geolocalización}
\texto{238.8586}{356.7820}{14.3462}{\normalfont\bfseries}{y}
\texto{252.7572}{356.7820}{14.3462}{\normalfont\bfseries}{mapas}
\texto{88.7990}{333.5540}{11.9552}{\normalfont}{Las}
\texto{110.0984}{333.5540}{11.9552}{\normalfont}{APIs}
\texto{139.1926}{333.5540}{11.9552}{\normalfont}{de}
\texto{154.4139}{333.5540}{11.9552}{\normalfont}{geolocalización,}
\texto{236.9454}{333.5540}{11.9552}{\normalfont}{como}
\texto{267.1260}{333.5540}{11.9552}{\normalfont\bfseries}{Google}
\texto{312.1565}{333.5540}{11.9552}{\normalfont\bfseries}{Maps}
\texto{348.2588}{333.5540}{11.9552}{\normalfont\bfseries}{Platform}
\texto{404.0510}{333.5540}{11.9552}{\normalfont}{y}
\texto{413.8520}{333.5540}{11.9552}{\normalfont\bfseries}{Mapbox}
\texto{461.9160}{333.5540}{11.9552}{\normalfont}{,}
\texto{468.6826}{333.5540}{11.9552}{\normalfont}{permiten}
\texto{518.0457}{333.5540}{11.9552}{\normalfont}{inte-}
\texto{70.8660}{319.1080}{11.9552}{\normalfont}{grar}
\texto{95.0251}{319.1080}{11.9552}{\normalfont}{mapas}
\texto{130.9672}{319.1080}{11.9552}{\normalfont}{interactivos}
\texto{193.2538}{319.1080}{11.9552}{\normalfont}{y}
\texto{202.7916}{319.1080}{11.9552}{\normalfont}{cálculo}
\texto{241.2587}{319.1080}{11.9552}{\normalfont}{de}
\texto{256.3246}{319.1080}{11.9552}{\normalfont}{rutas.}
\texto{289.0269}{319.1080}{11.9552}{\normalfont}{Para}
\texto{316.2752}{319.1080}{11.9552}{\normalfont}{el}
\texto{328.0893}{319.1080}{11.9552}{\normalfont}{ecosistema}
\texto{385.5579}{319.1080}{11.9552}{\normalfont}{turístico,}
\texto{434.5073}{319.1080}{11.9552}{\normalfont}{estas}
\texto{462.7096}{319.1080}{11.9552}{\normalfont}{funcionalidades}
\texto{70.8660}{304.6620}{11.9552}{\normalfont}{son}
\texto{91.7374}{304.6620}{11.9552}{\normalfont}{esenciales}
\texto{144.5411}{304.6620}{11.9552}{\normalfont}{para}
\texto{171.2132}{304.6620}{11.9552}{\normalfont}{que}
\texto{192.9955}{304.6620}{11.9552}{\normalfont}{el}
\texto{205.3477}{304.6620}{11.9552}{\normalfont}{turista}
\texto{243.1404}{304.6620}{11.9552}{\normalfont}{identifique}
\texto{300.3700}{304.6620}{11.9552}{\normalfont}{la}
\texto{313.3724}{304.6620}{11.9552}{\normalfont}{ubicación}
\texto{365.4087}{304.6620}{11.9552}{\normalfont}{exacta}
\texto{402.1494}{304.6620}{11.9552}{\normalfont}{de}
\texto{417.7653}{304.6620}{11.9552}{\normalfont}{los}
\texto{435.3848}{304.6620}{11.9552}{\normalfont}{tours.}
\texto{70.8660}{266.7830}{14.3462}{\normalfont\bfseries}{2.2.7.}
\texto{124.6642}{266.7830}{14.3462}{\normalfont\bfseries}{Marketplaces}
\texto{223.6903}{266.7830}{14.3462}{\normalfont\bfseries}{turísticos}
\texto{294.1244}{266.7830}{14.3462}{\normalfont\bfseries}{y}
\texto{308.0230}{266.7830}{14.3462}{\normalfont\bfseries}{comisiones}
\texto{88.7990}{243.5550}{11.9552}{\normalfont}{El}
\texto{102.9085}{243.5550}{11.9552}{\normalfont}{modelo}
\texto{142.5568}{243.5550}{11.9552}{\normalfont}{de}
\texto{157.1564}{243.5550}{11.9552}{\normalfont}{negocio}
\texto{198.1054}{243.5550}{11.9552}{\normalfont}{de}
\texto{212.7051}{243.5550}{11.9552}{\normalfont}{los}
\texto{229.3204}{243.5550}{11.9552}{\normalfont}{marketplaces}
\texto{298.6295}{243.5550}{11.9552}{\normalfont}{turísticos}
\texto{348.4779}{243.5550}{11.9552}{\normalfont}{se}
\texto{361.2030}{243.5550}{11.9552}{\normalfont}{sustenta}
\texto{406.1761}{243.5550}{11.9552}{\normalfont}{en}
\texto{420.7758}{243.5550}{11.9552}{\normalfont}{la}
\texto{432.7860}{243.5550}{11.9552}{\normalfont}{comisión}
\texto{479.9684}{243.5550}{11.9552}{\normalfont}{por}
\texto{500.1057}{243.5550}{11.9552}{\normalfont}{transac-}
\texto{70.8660}{229.1090}{11.9552}{\normalfont}{ción.}
\texto{98.4682}{229.1090}{11.9552}{\normalfont}{Plataformas}
\texto{163.7029}{229.1090}{11.9552}{\normalfont}{como}
\texto{193.9150}{229.1090}{11.9552}{\normalfont\itshape}{Civitatis}
\texto{240.9690}{229.1090}{11.9552}{\normalfont}{cobran}
\texto{278.9770}{229.1090}{11.9552}{\normalfont}{entre}
\texto{308.2074}{229.1090}{11.9552}{\normalfont}{el}
\texto{320.2128}{229.1090}{11.9552}{\normalfont}{15}
\texto{333.9159}{229.1090}{11.9552}{\normalfont}{\%}
\texto{347.2101}{229.1090}{11.9552}{\normalfont}{y}
\texto{356.9392}{229.1090}{11.9552}{\normalfont}{el}
\texto{368.9327}{229.1090}{11.9552}{\normalfont}{25}
\texto{382.6357}{229.1090}{11.9552}{\normalfont}{\%}
\texto{395.9419}{229.1090}{11.9552}{\normalfont}{del}
\texto{414.4390}{229.1090}{11.9552}{\normalfont}{valor}
\texto{443.0214}{229.1090}{11.9552}{\normalfont}{total,}
\texto{473.8874}{229.1090}{11.9552}{\normalfont}{lo}
\texto{486.5432}{229.1090}{11.9552}{\normalfont}{que}
\texto{507.9669}{229.1090}{11.9552}{\normalfont}{reduce}
\texto{70.8660}{214.6630}{11.9552}{\normalfont}{los}
\texto{89.7528}{214.6630}{11.9552}{\normalfont}{ingresos}
\texto{135.3810}{214.6630}{11.9552}{\normalfont}{de}
\texto{152.2642}{214.6630}{11.9552}{\normalfont}{los}
\texto{171.1510}{214.6630}{11.9552}{\normalfont}{proveedores}
\texto{235.9051}{214.6630}{11.9552}{\normalfont}{locales}
\texto{274.6376}{214.6630}{11.9552}{\normalfont}{[6].}
\texto{295.4230}{214.6630}{11.9552}{\normalfont\bfseries}{Explora-Nicaragua}
\texto{410.7640}{214.6630}{11.9552}{\normalfont}{propone}
\texto{457.2243}{214.6630}{11.9552}{\normalfont}{una}
\texto{481.2614}{214.6630}{11.9552}{\normalfont}{comisión}
\texto{530.7273}{214.6630}{11.9552}{\normalfont}{fi-}
\texto{70.8660}{200.2180}{11.9552}{\normalfont}{ja}
\texto{84.6958}{200.2180}{11.9552}{\normalfont}{del}
\texto{104.0656}{200.2180}{11.9552}{\normalfont}{8}
\texto{111.9154}{200.2180}{11.9552}{\normalfont}{\%}
\texto{126.0823}{200.2180}{11.9552}{\normalfont}{para}
\texto{153.2565}{200.2180}{11.9552}{\normalfont}{operadores,}
\texto{215.9376}{200.2180}{11.9552}{\normalfont}{junto}
\texto{247.0044}{200.2180}{11.9552}{\normalfont}{con}
\texto{268.9756}{200.2180}{11.9552}{\normalfont}{un}
\texto{286.3943}{200.2180}{11.9552}{\normalfont}{sistema}
\texto{328.6560}{200.2180}{11.9552}{\normalfont}{de}
\texto{344.7740}{200.2180}{11.9552}{\normalfont}{verificación}
\texto{406.4078}{200.2180}{11.9552}{\normalfont}{documental}
\texto{470.0023}{200.2180}{11.9552}{\normalfont}{de}
\texto{486.1203}{200.2180}{11.9552}{\normalfont}{operadores}
\texto{70.8660}{185.7720}{11.9552}{\normalfont}{gestionado}
\texto{128.8081}{185.7720}{11.9552}{\normalfont}{desde}
\texto{160.7356}{185.7720}{11.9552}{\normalfont}{el}
\texto{173.0997}{185.7720}{11.9552}{\normalfont}{panel}
\texto{204.3123}{185.7720}{11.9552}{\normalfont}{de}
\texto{219.9163}{185.7720}{11.9552}{\normalfont}{administración.}
\texto{70.8660}{147.8920}{14.3462}{\normalfont\bfseries}{2.2.8.}
\texto{124.6642}{147.8920}{14.3462}{\normalfont\bfseries}{Autenticación}
\texto{227.3716}{147.8920}{14.3462}{\normalfont\bfseries}{y}
\texto{241.2702}{147.8920}{14.3462}{\normalfont\bfseries}{seguridad}
\texto{88.7990}{124.6640}{11.9552}{\normalfont}{La}
\texto{107.5854}{124.6640}{11.9552}{\normalfont}{seguridad}
\texto{162.0461}{124.6640}{11.9552}{\normalfont}{en}
\texto{179.3716}{124.6640}{11.9552}{\normalfont}{aplicaciones}
\texto{245.5508}{124.6640}{11.9552}{\normalfont}{transaccionales}
\texto{328.0417}{124.6640}{11.9552}{\normalfont}{exige}
\texto{359.3500}{124.6640}{11.9552}{\normalfont}{el}
\texto{373.4236}{124.6640}{11.9552}{\normalfont}{uso}
\texto{396.0285}{124.6640}{11.9552}{\normalfont}{de}
\texto{413.3540}{124.6640}{11.9552}{\normalfont}{protocolos}
\texto{471.3893}{124.6640}{11.9552}{\normalfont}{modernos}
\texto{526.1848}{124.6640}{11.9552}{\normalfont}{co-}
\texto{70.8660}{110.2190}{11.9552}{\normalfont}{mo}
\texto{90.9580}{110.2190}{11.9552}{\normalfont\bfseries}{OAuth}
\texto{135.7888}{110.2190}{11.9552}{\normalfont\bfseries}{2.0}
\texto{157.4620}{110.2190}{11.9552}{\normalfont}{y}
\texto{168.2310}{110.2190}{11.9552}{\normalfont\bfseries}{JWT}
\texto{202.9050}{110.2190}{11.9552}{\normalfont}{(JSON}
\texto{242.3261}{110.2190}{11.9552}{\normalfont}{Web}
\texto{269.5744}{110.2190}{11.9552}{\normalfont}{Tokens)}
\texto{314.1171}{110.2190}{11.9552}{\normalfont}{para}
\texto{341.3630}{110.2190}{11.9552}{\normalfont}{la}
\texto{354.9512}{110.2190}{11.9552}{\normalfont}{autenticación,}
\texto{430.6516}{110.2190}{11.9552}{\normalfont}{junto}
\texto{461.8020}{110.2190}{11.9552}{\normalfont}{con}
\texto{483.8450}{110.2190}{11.9552}{\normalfont}{certificados}
\texto{70.8660}{95.7730}{11.9552}{\normalfont\bfseries}{SSL/TLS}
\texto{129.3050}{95.7730}{11.9552}{\normalfont}{para}
\texto{155.8455}{95.7730}{11.9552}{\normalfont}{el}
\texto{168.0781}{95.7730}{11.9552}{\normalfont}{cifrado}
\texto{206.6623}{95.7730}{11.9552}{\normalfont}{de}
\texto{222.1467}{95.7730}{11.9552}{\normalfont}{las}
\texto{239.6467}{95.7730}{11.9552}{\normalfont}{comunicaciones}
\texto{321.8722}{95.7730}{11.9552}{\normalfont}{[7].}
\texto{341.2545}{95.7730}{11.9552}{\normalfont}{En}
\texto{359.5680}{95.7730}{11.9552}{\normalfont\bfseries}{Explora-Nicaragua}
\texto{469.4750}{95.7730}{11.9552}{\normalfont}{,}
\texto{476.5047}{95.7730}{11.9552}{\normalfont}{el}
\texto{488.7372}{95.7730}{11.9552}{\normalfont}{JWT}
\texto{519.0221}{95.7730}{11.9552}{\normalfont}{emi-}
\texto{70.8660}{81.3270}{11.9552}{\normalfont}{tido}
\texto{94.4823}{81.3270}{11.9552}{\normalfont}{por}
\texto{115.1815}{81.3270}{11.9552}{\normalfont}{el}
\texto{127.1033}{81.3270}{11.9552}{\normalfont}{backend}
\texto{171.8731}{81.3270}{11.9552}{\normalfont}{incluye}
\texto{211.1005}{81.3270}{11.9552}{\normalfont}{los}
\texto{228.2950}{81.3270}{11.9552}{\normalfont\itshape}{claims}
\texto{264.5090}{81.3270}{11.9552}{\normalfont}{de}
\texto{279.6706}{81.3270}{11.9552}{\normalfont}{rol,}
\texto{300.0470}{81.3270}{11.9552}{\normalfont}{y}
\texto{309.6925}{81.3270}{11.9552}{\normalfont}{todos}
\texto{340.8621}{81.3270}{11.9552}{\normalfont}{los}
\texto{358.0513}{81.3270}{11.9552}{\normalfont}{clientes}
\texto{398.9811}{81.3270}{11.9552}{\normalfont}{lo}
\texto{411.5532}{81.3270}{11.9552}{\normalfont}{utilizan}
\texto{453.3796}{81.3270}{11.9552}{\normalfont}{para}
\texto{479.6093}{81.3270}{11.9552}{\normalfont}{autenticarse}
\texto{70.8660}{66.8810}{11.9552}{\normalfont}{e}
\texto{79.9663}{66.8810}{11.9552}{\normalfont}{informar}
\texto{127.7632}{66.8810}{11.9552}{\normalfont}{la}
\texto{140.7776}{66.8810}{11.9552}{\normalfont}{interfaz}
\texto{183.0488}{66.8810}{11.9552}{\normalfont}{según}
\texto{215.6267}{66.8810}{11.9552}{\normalfont}{el}
\texto{227.9908}{66.8810}{11.9552}{\normalfont}{rol.}
\end{pagina}
% Página original 26
\begin{pagina}{26}
\texto{70.8660}{746.3890}{11.9552}{\normalfont}{16}
\texto{269.9360}{746.3890}{11.9552}{\normalfont}{CAPÍTULO}
\texto{336.8995}{746.3890}{11.9552}{\normalfont}{2.}
\texto{353.8113}{746.3890}{11.9552}{\normalfont}{INVESTIGACIÓN}
\texto{454.8220}{746.3890}{11.9552}{\normalfont}{DOCUMENTAL}
\texto{70.8660}{705.1940}{14.3462}{\normalfont\bfseries}{2.2.9.}
\texto{124.6642}{705.1940}{14.3462}{\normalfont\bfseries}{Sistemas}
\texto{191.1632}{705.1940}{14.3462}{\normalfont\bfseries}{de}
\texto{212.8733}{705.1940}{14.3462}{\normalfont\bfseries}{reseñas}
\texto{269.1190}{705.1940}{14.3462}{\normalfont\bfseries}{y}
\texto{283.0176}{705.1940}{14.3462}{\normalfont\bfseries}{verificación}
\texto{88.7990}{683.0270}{11.9552}{\normalfont}{Los}
\texto{111.4971}{683.0270}{11.9552}{\normalfont}{sistemas}
\texto{158.8899}{683.0270}{11.9552}{\normalfont}{de}
\texto{175.5220}{683.0270}{11.9552}{\normalfont}{reseñas}
\texto{216.9970}{683.0270}{11.9552}{\normalfont}{verificadas}
\texto{275.3192}{683.0270}{11.9552}{\normalfont}{son}
\texto{297.2188}{683.0270}{11.9552}{\normalfont}{un}
\texto{315.1516}{683.0270}{11.9552}{\normalfont}{componente}
\texto{381.2112}{683.0270}{11.9552}{\normalfont}{crítico}
\texto{417.9926}{683.0270}{11.9552}{\normalfont}{en}
\texto{434.6247}{683.0270}{11.9552}{\normalfont}{marketplaces}
\texto{505.9542}{683.0270}{11.9552}{\normalfont}{turísti-}
\texto{70.8660}{668.5810}{11.9552}{\normalfont}{cos}
\texto{90.4845}{668.5810}{11.9552}{\normalfont}{[8].}
\texto{110.0470}{668.5810}{11.9552}{\normalfont\bfseries}{Explora-Nicaragua}
\texto{224.1660}{668.5810}{11.9552}{\normalfont}{implementa}
\texto{287.6218}{668.5810}{11.9552}{\normalfont}{un}
\texto{304.5743}{668.5810}{11.9552}{\normalfont}{mecanismo}
\texto{364.5153}{668.5810}{11.9552}{\normalfont}{de}
\texto{380.1790}{668.5810}{11.9552}{\normalfont}{verificación}
\texto{441.3585}{668.5810}{11.9552}{\normalfont}{basado}
\texto{480.4879}{668.5810}{11.9552}{\normalfont}{en}
\texto{496.1516}{668.5810}{11.9552}{\normalfont}{el}
\texto{508.5515}{668.5810}{11.9552}{\normalfont}{estado}
\texto{70.8660}{654.1350}{11.9552}{\normalfont}{del}
\texto{88.7773}{654.1350}{11.9552}{\normalfont}{voucher:}
\texto{134.3457}{654.1350}{11.9552}{\normalfont}{únicamente}
\texto{195.5085}{654.1350}{11.9552}{\normalfont}{los}
\texto{212.1956}{654.1350}{11.9552}{\normalfont}{turistas}
\texto{253.6610}{654.1350}{11.9552}{\normalfont}{que}
\texto{274.5109}{654.1350}{11.9552}{\normalfont}{presentan}
\texto{326.6332}{654.1350}{11.9552}{\normalfont}{el}
\texto{338.0528}{654.1350}{11.9552}{\normalfont}{voucher}
\texto{380.3574}{654.1350}{11.9552}{\normalfont}{digital}
\texto{415.8405}{654.1350}{11.9552}{\normalfont}{al}
\texto{427.9104}{654.1350}{11.9552}{\normalfont}{operador}
\texto{476.0732}{654.1350}{11.9552}{\normalfont}{pueden}
\texto{515.4464}{654.1350}{11.9552}{\normalfont}{dejar}
\texto{70.8660}{639.6890}{11.9552}{\normalfont}{una}
\texto{93.6239}{639.6890}{11.9552}{\normalfont}{reseña.}
\texto{132.7055}{639.6890}{11.9552}{\normalfont}{La}
\texto{149.7823}{639.6890}{11.9552}{\normalfont}{moderación}
\texto{212.5351}{639.6890}{11.9552}{\normalfont}{de}
\texto{228.1510}{639.6890}{11.9552}{\normalfont}{reseñas}
\texto{268.5978}{639.6890}{11.9552}{\normalfont}{se}
\texto{282.3152}{639.6890}{11.9552}{\normalfont}{realiza}
\texto{319.3931}{639.6890}{11.9552}{\normalfont}{desde}
\texto{351.3206}{639.6890}{11.9552}{\normalfont}{el}
\texto{363.6847}{639.6890}{11.9552}{\normalfont}{panel}
\texto{394.8973}{639.6890}{11.9552}{\normalfont}{de}
\texto{410.5013}{639.6890}{11.9552}{\normalfont}{administración.}
\texto{70.8660}{605.0280}{14.3462}{\normalfont\bfseries}{2.2.10.}
\texto{132.7340}{605.0280}{14.3462}{\normalfont\bfseries}{Notificaciones}
\texto{235.9707}{605.0280}{14.3462}{\normalfont\bfseries}{push}
\texto{88.7990}{582.8610}{11.9552}{\normalfont}{Las}
\texto{110.6350}{582.8610}{11.9552}{\normalfont\bfseries}{notificaciones}
\texto{193.7894}{582.8610}{11.9552}{\normalfont\bfseries}{push}
\texto{221.5140}{582.8610}{11.9552}{\normalfont}{,}
\texto{228.8067}{582.8610}{11.9552}{\normalfont}{soportadas}
\texto{287.9419}{582.8610}{11.9552}{\normalfont}{por}
\texto{309.2389}{582.8610}{11.9552}{\normalfont}{servicios}
\texto{356.0076}{582.8610}{11.9552}{\normalfont}{como}
\texto{386.7320}{582.8610}{11.9552}{\normalfont\bfseries}{Firebase}
\texto{440.6559}{582.8610}{11.9552}{\normalfont\bfseries}{Cloud}
\texto{480.4249}{582.8610}{11.9552}{\normalfont\bfseries}{Messaging}
\texto{70.8660}{568.4150}{11.9552}{\normalfont\bfseries}{(FCM)}
\texto{111.8790}{568.4150}{11.9552}{\normalfont}{,}
\texto{118.8011}{568.4150}{11.9552}{\normalfont}{permiten}
\texto{168.3315}{568.4150}{11.9552}{\normalfont}{comunicar}
\texto{224.3703}{568.4150}{11.9552}{\normalfont}{actualizaciones}
\texto{304.2095}{568.4150}{11.9552}{\normalfont}{en}
\texto{319.5982}{568.4150}{11.9552}{\normalfont}{tiempo}
\texto{358.7228}{568.4150}{11.9552}{\normalfont}{real}
\texto{381.2536}{568.4150}{11.9552}{\normalfont}{al}
\texto{394.0408}{568.4150}{11.9552}{\normalfont}{turista}
\texto{431.6065}{568.4150}{11.9552}{\normalfont}{[10].}
\texto{456.7339}{568.4150}{11.9552}{\normalfont}{En}
\texto{474.8843}{568.4150}{11.9552}{\normalfont}{el}
\texto{487.0212}{568.4150}{11.9552}{\normalfont}{ecosistema}
\texto{70.8660}{553.9690}{11.9552}{\normalfont}{se}
\texto{84.5834}{553.9690}{11.9552}{\normalfont}{segmentan}
\texto{142.2027}{553.9690}{11.9552}{\normalfont}{por}
\texto{163.3443}{553.9690}{11.9552}{\normalfont}{tipo}
\texto{187.7257}{553.9690}{11.9552}{\normalfont}{de}
\texto{203.3416}{553.9690}{11.9552}{\normalfont}{usuario}
\texto{244.3742}{553.9690}{11.9552}{\normalfont}{y}
\texto{254.4501}{553.9690}{11.9552}{\normalfont}{contexto}
\texto{301.9361}{553.9690}{11.9552}{\normalfont}{(confirmación,}
\texto{378.6741}{553.9690}{11.9552}{\normalfont}{recordatorio,}
\texto{447.6078}{553.9690}{11.9552}{\normalfont}{oferta).}
\texto{70.8660}{519.3070}{14.3462}{\normalfont\bfseries}{2.2.11.}
\texto{132.7340}{519.3070}{14.3462}{\normalfont\bfseries}{Herramientas}
\texto{233.4199}{519.3070}{14.3462}{\normalfont\bfseries}{de}
\texto{255.1300}{519.3070}{14.3462}{\normalfont\bfseries}{modelado}
\texto{328.1837}{519.3070}{14.3462}{\normalfont\bfseries}{y}
\texto{342.0823}{519.3070}{14.3462}{\normalfont\bfseries}{metodologías}
\texto{439.4213}{519.3070}{14.3462}{\normalfont\bfseries}{Lean}
\texto{88.7990}{497.1400}{11.9552}{\normalfont}{El}
\texto{105.0900}{497.1400}{11.9552}{\normalfont\bfseries}{BPMN}
\texto{152.8869}{497.1400}{11.9552}{\normalfont\bfseries}{2.0}
\texto{175.1060}{497.1400}{11.9552}{\normalfont}{es}
\texto{189.9472}{497.1400}{11.9552}{\normalfont}{el}
\texto{203.4350}{497.1400}{11.9552}{\normalfont}{estándar}
\texto{252.1071}{497.1400}{11.9552}{\normalfont}{del}
\texto{272.0866}{497.1400}{11.9552}{\normalfont}{OMG}
\texto{306.1362}{497.1400}{11.9552}{\normalfont}{para}
\texto{333.9320}{497.1400}{11.9552}{\normalfont}{modelar}
\texto{380.2608}{497.1400}{11.9552}{\normalfont}{procesos}
\texto{428.0195}{497.1400}{11.9552}{\normalfont}{de}
\texto{444.7472}{497.1400}{11.9552}{\normalfont}{negocio}
\texto{487.8241}{497.1400}{11.9552}{\normalfont}{[11],}
\texto{514.3088}{497.1400}{11.9552}{\normalfont}{y}
\texto{525.5204}{497.1400}{11.9552}{\normalfont}{he-}
\texto{70.8660}{482.6950}{11.9552}{\normalfont}{rramientas}
\texto{129.0280}{482.6950}{11.9552}{\normalfont}{como}
\texto{159.5020}{482.6950}{11.9552}{\normalfont\bfseries}{Camunda}
\texto{219.9092}{482.6950}{11.9552}{\normalfont\bfseries}{Modeler}
\texto{272.5440}{482.6950}{11.9552}{\normalfont}{permiten}
\texto{322.1939}{482.6950}{11.9552}{\normalfont}{diseñarlos}
\texto{376.2028}{482.6950}{11.9552}{\normalfont}{[12].}
\texto{401.4674}{482.6950}{11.9552}{\normalfont}{El}
\texto{416.4855}{482.6950}{11.9552}{\normalfont}{enfoque}
\texto{459.6340}{482.6950}{11.9552}{\normalfont\itshape}{Lean}
\texto{487.7933}{482.6950}{11.9552}{\normalfont\itshape}{Startup}
\texto{529.4280}{482.6950}{11.9552}{\normalfont}{de}
\texto{70.8660}{468.2490}{11.9552}{\normalfont}{Eric}
\texto{96.0939}{468.2490}{11.9552}{\normalfont}{Ries}
\texto{122.0486}{468.2490}{11.9552}{\normalfont}{[3]}
\texto{138.6619}{468.2490}{11.9552}{\normalfont}{y}
\texto{149.1083}{468.2490}{11.9552}{\normalfont}{el}
\texto{161.8260}{468.2490}{11.9552}{\normalfont\itshape}{Lean}
\texto{190.4037}{468.2490}{11.9552}{\normalfont\itshape}{Canvas}
\texto{232.2060}{468.2490}{11.9552}{\normalfont}{de}
\texto{248.1686}{468.2490}{11.9552}{\normalfont}{Ash}
\texto{272.3324}{468.2490}{11.9552}{\normalfont}{Maurya}
\texto{315.9330}{468.2490}{11.9552}{\normalfont}{[13]}
\texto{338.4039}{468.2490}{11.9552}{\normalfont}{complementan}
\texto{416.4785}{468.2490}{11.9552}{\normalfont}{la}
\texto{429.8516}{468.2490}{11.9552}{\normalfont}{validación}
\texto{485.1659}{468.2490}{11.9552}{\normalfont}{del}
\texto{504.3923}{468.2490}{11.9552}{\normalfont}{modelo}
\texto{70.8660}{453.8030}{11.9552}{\normalfont}{de}
\texto{87.3905}{453.8030}{11.9552}{\normalfont}{negocio}
\texto{130.2523}{453.8030}{11.9552}{\normalfont}{mediante}
\texto{181.5735}{453.8030}{11.9552}{\normalfont}{un}
\texto{199.3987}{453.8030}{11.9552}{\normalfont}{producto}
\texto{250.0649}{453.8030}{11.9552}{\normalfont}{mínimo}
\texto{293.2542}{453.8030}{11.9552}{\normalfont}{viable}
\texto{328.3141}{453.8030}{11.9552}{\normalfont}{y}
\texto{339.3224}{453.8030}{11.9552}{\normalfont}{el}
\texto{352.5951}{453.8030}{11.9552}{\normalfont}{ciclo}
\texto{380.1720}{453.8030}{11.9552}{\normalfont\bfseries}{Build--Measure--Learn}
\texto{509.1600}{453.8030}{11.9552}{\normalfont}{.}
\texto{517.2298}{453.8030}{11.9552}{\normalfont}{Para}
\texto{70.8660}{439.3570}{11.9552}{\normalfont\bfseries}{Explora-Nicaragua}
\texto{180.7730}{439.3570}{11.9552}{\normalfont}{,}
\texto{187.6472}{439.3570}{11.9552}{\normalfont}{BPMN}
\texto{227.0372}{439.3570}{11.9552}{\normalfont}{se}
\texto{240.4916}{439.3570}{11.9552}{\normalfont}{aplica}
\texto{274.0427}{439.3570}{11.9552}{\normalfont}{a}
\texto{283.5303}{439.3570}{11.9552}{\normalfont}{procesos}
\texto{329.8902}{439.3570}{11.9552}{\normalfont}{como}
\texto{360.1895}{439.3570}{11.9552}{\normalfont}{la}
\texto{372.9289}{439.3570}{11.9552}{\normalfont}{reserva,}
\texto{415.3292}{439.3570}{11.9552}{\normalfont}{la}
\texto{428.0687}{439.3570}{11.9552}{\normalfont}{gestión}
\texto{467.5376}{439.3570}{11.9552}{\normalfont}{del}
\texto{486.1184}{439.3570}{11.9552}{\normalfont}{operador}
\texto{534.9625}{439.3570}{11.9552}{\normalfont}{y}
\texto{70.8660}{424.9110}{11.9552}{\normalfont}{la}
\texto{84.2750}{424.9110}{11.9552}{\normalfont}{verificación}
\texto{145.8251}{424.9110}{11.9552}{\normalfont}{de}
\texto{161.8475}{424.9110}{11.9552}{\normalfont}{operadores}
\texto{221.1811}{424.9110}{11.9552}{\normalfont}{por}
\texto{242.7411}{424.9110}{11.9552}{\normalfont}{el}
\texto{255.5117}{424.9110}{11.9552}{\normalfont}{administrador;}
\texto{334.6838}{424.9110}{11.9552}{\normalfont}{mientras}
\texto{382.9541}{424.9110}{11.9552}{\normalfont}{que}
\texto{405.1549}{424.9110}{11.9552}{\normalfont}{las}
\texto{423.1929}{424.9110}{11.9552}{\normalfont}{métricas}
\texto{470.4973}{424.9110}{11.9552}{\normalfont}{Lean}
\texto{499.6871}{424.9110}{11.9552}{\normalfont}{(tasa}
\texto{529.4316}{424.9110}{11.9552}{\normalfont}{de}
\texto{70.8660}{410.4660}{11.9552}{\normalfont}{conversión,}
\texto{131.0772}{410.4660}{11.9552}{\normalfont}{NPS)}
\texto{162.7824}{410.4660}{11.9552}{\normalfont}{miden}
\texto{197.8972}{410.4660}{11.9552}{\normalfont}{la}
\texto{210.9116}{410.4660}{11.9552}{\normalfont}{validación}
\texto{265.8673}{410.4660}{11.9552}{\normalfont}{del}
\texto{284.7230}{410.4660}{11.9552}{\normalfont}{MVP.}
\end{pagina}
% Página original 27
\begin{pagina}{27}
\texto{70.8660}{746.3890}{11.9552}{\normalfont}{2.3.}
\texto{96.8710}{746.3890}{11.9552}{\normalfont}{ESTADO}
\texto{149.5432}{746.3890}{11.9552}{\normalfont}{DEL}
\texto{177.6594}{746.3890}{11.9552}{\normalfont}{ARTE}
\texto{529.4280}{746.3890}{11.9552}{\normalfont}{17}
\texto{70.8660}{705.1940}{17.2154}{\normalfont\bfseries}{2.3.}
\texto{120.3603}{705.1940}{17.2154}{\normalfont\bfseries}{Estado}
\texto{184.5582}{705.1940}{17.2154}{\normalfont\bfseries}{del}
\texto{215.9901}{705.1940}{17.2154}{\normalfont\bfseries}{arte}
\texto{88.7990}{678.9090}{11.9552}{\normalfont}{Se}
\texto{103.6258}{678.9090}{11.9552}{\normalfont}{revisaron}
\texto{153.3117}{678.9090}{11.9552}{\normalfont}{investigaciones,}
\texto{234.9346}{678.9090}{11.9552}{\normalfont}{tesis}
\texto{260.2963}{678.9090}{11.9552}{\normalfont}{y}
\texto{269.6071}{678.9090}{11.9552}{\normalfont}{soluciones}
\texto{323.5848}{678.9090}{11.9552}{\normalfont}{existentes}
\texto{376.2642}{678.9090}{11.9552}{\normalfont}{relacionadas}
\texto{441.8839}{678.9090}{11.9552}{\normalfont}{con}
\texto{462.5759}{678.9090}{11.9552}{\normalfont}{la}
\texto{474.8013}{678.9090}{11.9552}{\normalfont}{problemática}
\texto{70.8660}{664.4630}{11.9552}{\normalfont}{estudiada,}
\texto{126.8570}{664.4630}{11.9552}{\normalfont}{tanto}
\texto{157.7468}{664.4630}{11.9552}{\normalfont}{a}
\texto{167.4975}{664.4630}{11.9552}{\normalfont}{nivel}
\texto{195.4727}{664.4630}{11.9552}{\normalfont}{nacional}
\texto{241.6436}{664.4630}{11.9552}{\normalfont}{como}
\texto{272.2059}{664.4630}{11.9552}{\normalfont}{internacional.}
\texto{70.8660}{629.8010}{14.3462}{\normalfont\bfseries}{Investigaciones}
\texto{181.6674}{629.8010}{14.3462}{\normalfont\bfseries}{internacionales}
\texto{291.8491}{629.8010}{14.3462}{\normalfont\bfseries}{y}
\texto{305.7477}{629.8010}{14.3462}{\normalfont\bfseries}{sistemas}
\texto{369.6471}{629.8010}{14.3462}{\normalfont\bfseries}{similares}
\texto{88.7990}{607.6340}{11.9552}{\normalfont}{Existen}
\texto{130.6207}{607.6340}{11.9552}{\normalfont}{plataformas}
\texto{194.3945}{607.6340}{11.9552}{\normalfont}{internacionales}
\texto{273.7794}{607.6340}{11.9552}{\normalfont}{consolidadas}
\texto{341.2067}{607.6340}{11.9552}{\normalfont}{como}
\texto{371.4260}{607.6340}{11.9552}{\normalfont\itshape}{Civitatis}
\texto{414.4640}{607.6340}{11.9552}{\normalfont}{,}
\texto{421.2720}{607.6340}{11.9552}{\normalfont\itshape}{GetYourGuide}
\texto{498.8010}{607.6340}{11.9552}{\normalfont}{y}
\texto{508.5350}{607.6340}{11.9552}{\normalfont\itshape}{Viator}
\texto{70.8660}{593.1890}{11.9552}{\normalfont}{que}
\texto{92.4571}{593.1890}{11.9552}{\normalfont}{permiten}
\texto{142.0114}{593.1890}{11.9552}{\normalfont}{reservar}
\texto{185.7842}{593.1890}{11.9552}{\normalfont}{tours}
\texto{215.5813}{593.1890}{11.9552}{\normalfont}{en}
\texto{230.9939}{593.1890}{11.9552}{\normalfont}{línea.}
\texto{262.0273}{593.1890}{11.9552}{\normalfont}{Estos}
\texto{293.3379}{593.1890}{11.9552}{\normalfont}{sistemas}
\texto{339.5113}{593.1890}{11.9552}{\normalfont}{comparten}
\texto{397.2095}{593.1890}{11.9552}{\normalfont}{características}
\texto{473.2350}{593.1890}{11.9552}{\normalfont}{como}
\texto{503.6179}{593.1890}{11.9552}{\normalfont}{el}
\texto{515.7788}{593.1890}{11.9552}{\normalfont}{catá-}
\texto{70.8660}{578.7430}{11.9552}{\normalfont}{logo}
\texto{95.2881}{578.7430}{11.9552}{\normalfont}{por}
\texto{116.1308}{578.7430}{11.9552}{\normalfont}{destino,}
\texto{159.4899}{578.7430}{11.9552}{\normalfont}{el}
\texto{171.5551}{578.7430}{11.9552}{\normalfont}{proceso}
\texto{213.2859}{578.7430}{11.9552}{\normalfont}{de}
\texto{228.6029}{578.7430}{11.9552}{\normalfont}{reserva}
\texto{267.7275}{578.7430}{11.9552}{\normalfont}{y}
\texto{277.5164}{578.7430}{11.9552}{\normalfont}{las}
\texto{294.8610}{578.7430}{11.9552}{\normalfont}{comisiones}
\texto{352.5807}{578.7430}{11.9552}{\normalfont}{por}
\texto{373.4234}{578.7430}{11.9552}{\normalfont}{transacción;}
\texto{438.2445}{578.7430}{11.9552}{\normalfont}{sin}
\texto{456.2275}{578.7430}{11.9552}{\normalfont}{embargo,}
\texto{506.3413}{578.7430}{11.9552}{\normalfont}{operan}
\texto{70.8660}{564.2970}{11.9552}{\normalfont}{con}
\texto{91.7254}{564.2970}{11.9552}{\normalfont}{comisiones}
\texto{149.1463}{564.2970}{11.9552}{\normalfont}{elevadas}
\texto{194.4756}{564.2970}{11.9552}{\normalfont}{(15}
\texto{212.7312}{564.2970}{11.9552}{\normalfont}{\%--25}
\texto{242.0310}{564.2970}{11.9552}{\normalfont}{\%)}
\texto{259.6506}{564.2970}{11.9552}{\normalfont}{[6]}
\texto{275.3094}{564.2970}{11.9552}{\normalfont}{y}
\texto{284.7995}{564.2970}{11.9552}{\normalfont}{generalmente}
\texto{355.4260}{564.2970}{11.9552}{\normalfont}{no}
\texto{371.0945}{564.2970}{11.9552}{\normalfont}{ofrecen}
\texto{410.5012}{564.2970}{11.9552}{\normalfont}{un}
\texto{426.8201}{564.2970}{11.9552}{\normalfont}{backoffice}
\texto{479.2269}{564.2970}{11.9552}{\normalfont}{de}
\texto{494.2450}{564.2970}{11.9552}{\normalfont}{escritorio}
\texto{70.8660}{549.8510}{11.9552}{\normalfont}{adaptado}
\texto{122.2399}{549.8510}{11.9552}{\normalfont}{a}
\texto{131.9905}{549.8510}{11.9552}{\normalfont}{las}
\texto{149.6221}{549.8510}{11.9552}{\normalfont}{necesidades}
\texto{212.1812}{549.8510}{11.9552}{\normalfont}{locales}
\texto{249.6465}{549.8510}{11.9552}{\normalfont}{de}
\texto{265.2504}{549.8510}{11.9552}{\normalfont}{gestión}
\texto{304.9823}{549.8510}{11.9552}{\normalfont}{de}
\texto{320.5982}{549.8510}{11.9552}{\normalfont}{los}
\texto{338.2177}{549.8510}{11.9552}{\normalfont}{operadores,}
\texto{400.3967}{549.8510}{11.9552}{\normalfont}{sino}
\texto{424.5199}{549.8510}{11.9552}{\normalfont}{un}
\texto{441.4366}{549.8510}{11.9552}{\normalfont}{panel}
\texto{472.6492}{549.8510}{11.9552}{\normalfont}{web}
\texto{496.3850}{549.8510}{11.9552}{\normalfont}{básico.}
\texto{70.8660}{515.1890}{14.3462}{\normalfont\bfseries}{Investigaciones}
\texto{181.6674}{515.1890}{14.3462}{\normalfont\bfseries}{nacionales}
\texto{88.7990}{493.0230}{11.9552}{\normalfont}{A}
\texto{102.3203}{493.0230}{11.9552}{\normalfont}{nivel}
\texto{131.1443}{493.0230}{11.9552}{\normalfont}{nacional,}
\texto{181.4279}{493.0230}{11.9552}{\normalfont}{las}
\texto{199.8963}{493.0230}{11.9552}{\normalfont}{propuestas}
\texto{259.4141}{493.0230}{11.9552}{\normalfont}{académicas}
\texto{321.4687}{493.0230}{11.9552}{\normalfont}{y}
\texto{332.4053}{493.0230}{11.9552}{\normalfont}{comerciales}
\texto{395.0984}{493.0230}{11.9552}{\normalfont}{para}
\texto{422.6193}{493.0230}{11.9552}{\normalfont}{la}
\texto{436.4825}{493.0230}{11.9552}{\normalfont}{gestión}
\texto{477.0752}{493.0230}{11.9552}{\normalfont}{turística}
\texto{524.1596}{493.0230}{11.9552}{\normalfont}{son}
\texto{70.8660}{478.5770}{11.9552}{\normalfont}{limitadas}
\texto{122.4790}{478.5770}{11.9552}{\normalfont}{y}
\texto{133.3797}{478.5770}{11.9552}{\normalfont}{en}
\texto{149.8205}{478.5770}{11.9552}{\normalfont}{su}
\texto{165.6636}{478.5770}{11.9552}{\normalfont}{mayoría}
\texto{211.0383}{478.5770}{11.9552}{\normalfont}{se}
\texto{225.5926}{478.5770}{11.9552}{\normalfont}{centran}
\texto{268.3635}{478.5770}{11.9552}{\normalfont}{en}
\texto{284.8043}{478.5770}{11.9552}{\normalfont}{portales}
\texto{330.2365}{478.5770}{11.9552}{\normalfont}{informativos}
\texto{398.4480}{478.5770}{11.9552}{\normalfont}{del}
\texto{418.1287}{478.5770}{11.9552}{\normalfont}{catálogo,}
\texto{468.3884}{478.5770}{11.9552}{\normalfont}{sin}
\texto{487.4832}{478.5770}{11.9552}{\normalfont}{funcionali-}
\texto{70.8660}{464.1310}{11.9552}{\normalfont}{dad}
\texto{93.6239}{464.1310}{11.9552}{\normalfont}{transaccional}
\texto{164.5732}{464.1310}{11.9552}{\normalfont}{ni}
\texto{178.2261}{464.1310}{11.9552}{\normalfont}{de}
\texto{193.8300}{464.1310}{11.9552}{\normalfont}{gestión}
\texto{233.5739}{464.1310}{11.9552}{\normalfont}{de}
\texto{249.1778}{464.1310}{11.9552}{\normalfont}{disponibilidad.}
\texto{328.2543}{464.1310}{11.9552}{\normalfont}{El}
\texto{343.3680}{464.1310}{11.9552}{\normalfont}{directorio}
\texto{396.0546}{464.1310}{11.9552}{\normalfont}{de}
\texto{411.6585}{464.1310}{11.9552}{\normalfont}{operadores}
\texto{470.5738}{464.1310}{11.9552}{\normalfont}{publicado}
\texto{523.8987}{464.1310}{11.9552}{\normalfont}{por}
\texto{70.8660}{449.6850}{11.9552}{\normalfont}{INTUR}
\texto{113.6058}{449.6850}{11.9552}{\normalfont}{[1]}
\texto{129.8642}{449.6850}{11.9552}{\normalfont}{es}
\texto{143.5816}{449.6850}{11.9552}{\normalfont}{un}
\texto{160.4982}{449.6850}{11.9552}{\normalfont}{ejemplo}
\texto{203.7426}{449.6850}{11.9552}{\normalfont}{de}
\texto{219.3465}{449.6850}{11.9552}{\normalfont}{solución}
\texto{264.2933}{449.6850}{11.9552}{\normalfont}{informativa}
\texto{326.7281}{449.6850}{11.9552}{\normalfont}{sin}
\texto{344.9980}{449.6850}{11.9552}{\normalfont}{reserva}
\texto{384.4095}{449.6850}{11.9552}{\normalfont}{ni}
\texto{398.0743}{449.6850}{11.9552}{\normalfont}{pago}
\texto{426.0352}{449.6850}{11.9552}{\normalfont}{en}
\texto{441.6391}{449.6850}{11.9552}{\normalfont}{línea.}
\texto{70.8660}{415.0230}{14.3462}{\normalfont\bfseries}{Similitudes}
\texto{154.5345}{415.0230}{14.3462}{\normalfont\bfseries}{y}
\texto{168.4331}{415.0230}{14.3462}{\normalfont\bfseries}{diferencias}
\texto{248.3500}{415.0230}{14.3462}{\normalfont\bfseries}{con}
\texto{277.9391}{415.0230}{14.3462}{\normalfont\bfseries}{la}
\texto{295.6480}{415.0230}{14.3462}{\normalfont\bfseries}{propuesta}
\texto{88.7990}{392.8570}{11.9552}{\normalfont}{La}
\texto{106.4138}{392.8570}{11.9552}{\normalfont}{propuesta}
\texto{161.1530}{392.8570}{11.9552}{\normalfont\bfseries}{Explora-Nicaragua}
\texto{275.7720}{392.8570}{11.9552}{\normalfont}{se}
\texto{290.0393}{392.8570}{11.9552}{\normalfont}{diferencia}
\texto{343.5891}{392.8570}{11.9552}{\normalfont}{de}
\texto{359.7429}{392.8570}{11.9552}{\normalfont}{las}
\texto{377.9124}{392.8570}{11.9552}{\normalfont}{soluciones}
\texto{433.2292}{392.8570}{11.9552}{\normalfont}{existentes}
\texto{487.2356}{392.8570}{11.9552}{\normalfont}{en}
\texto{503.3894}{392.8570}{11.9552}{\normalfont}{tres}
\texto{526.7619}{392.8570}{11.9552}{\normalfont}{as-}
\texto{70.8660}{378.4110}{11.9552}{\normalfont}{pectos:}
\texto{111.6930}{378.4110}{11.9552}{\normalfont}{(1)}
\texto{131.9714}{378.4110}{11.9552}{\normalfont}{integra}
\texto{172.7267}{378.4110}{11.9552}{\normalfont}{de}
\texto{189.7533}{378.4110}{11.9552}{\normalfont}{forma}
\texto{224.6649}{378.4110}{11.9552}{\normalfont}{coordinada}
\texto{286.2389}{378.4110}{11.9552}{\normalfont}{un}
\texto{304.5662}{378.4110}{11.9552}{\normalfont}{backoffice}
\texto{358.9815}{378.4110}{11.9552}{\normalfont}{de}
\texto{376.0081}{378.4110}{11.9552}{\normalfont}{escritorio}
\texto{428.2189}{378.4110}{11.9552}{\normalfont}{para}
\texto{456.3016}{378.4110}{11.9552}{\normalfont}{operadores}
\texto{516.6395}{378.4110}{11.9552}{\normalfont}{y}
\texto{528.1381}{378.4110}{11.9552}{\normalfont}{un}
\texto{70.8660}{363.9650}{11.9552}{\normalfont}{panel}
\texto{101.6841}{363.9650}{11.9552}{\normalfont}{de}
\texto{116.8935}{363.9650}{11.9552}{\normalfont}{administración}
\texto{195.9030}{363.9650}{11.9552}{\normalfont}{de}
\texto{211.1244}{363.9650}{11.9552}{\normalfont}{escritorio}
\texto{261.5180}{363.9650}{11.9552}{\normalfont}{para}
\texto{287.7835}{363.9650}{11.9552}{\normalfont}{el}
\texto{299.7531}{363.9650}{11.9552}{\normalfont}{equipo}
\texto{337.0724}{363.9650}{11.9552}{\normalfont}{gestor}
\texto{371.2069}{363.9650}{11.9552}{\normalfont}{sobre}
\texto{401.4512}{363.9650}{11.9552}{\normalfont}{una}
\texto{423.8146}{363.9650}{11.9552}{\normalfont}{API}
\texto{448.2797}{363.9650}{11.9552}{\normalfont}{REST}
\texto{483.3324}{363.9650}{11.9552}{\normalfont}{común,}
\texto{523.5831}{363.9650}{11.9552}{\normalfont}{con}
\texto{70.8660}{349.5190}{11.9552}{\normalfont}{la}
\texto{84.3467}{349.5190}{11.9552}{\normalfont}{app}
\texto{107.5828}{349.5190}{11.9552}{\normalfont}{móvil}
\texto{139.9264}{349.5190}{11.9552}{\normalfont}{contemplada}
\texto{209.0155}{349.5190}{11.9552}{\normalfont}{como}
\texto{240.0560}{349.5190}{11.9552}{\normalfont}{trabajo}
\texto{281.4999}{349.5190}{11.9552}{\normalfont}{futuro}
\texto{317.4300}{349.5190}{11.9552}{\normalfont}{dentro}
\texto{354.6514}{349.5190}{11.9552}{\normalfont}{de}
\texto{370.7335}{349.5190}{11.9552}{\normalfont}{la}
\texto{384.2142}{349.5190}{11.9552}{\normalfont}{misma}
\texto{421.8229}{349.5190}{11.9552}{\normalfont}{arquitectura;}
\texto{492.2103}{349.5190}{11.9552}{\normalfont}{(2)}
\texto{511.5442}{349.5190}{11.9552}{\normalfont}{ofrece}
\texto{70.8660}{335.0730}{11.9552}{\normalfont}{a}
\texto{81.1427}{335.0730}{11.9552}{\normalfont}{los}
\texto{99.2883}{335.0730}{11.9552}{\normalfont}{operadores}
\texto{158.7415}{335.0730}{11.9552}{\normalfont}{una}
\texto{182.0254}{335.0730}{11.9552}{\normalfont}{herramienta}
\texto{247.9106}{335.0730}{11.9552}{\normalfont}{de}
\texto{264.0525}{335.0730}{11.9552}{\normalfont}{gestión}
\texto{304.3104}{335.0730}{11.9552}{\normalfont}{de}
\texto{320.4523}{335.0730}{11.9552}{\normalfont}{escritorio}
\texto{371.7664}{335.0730}{11.9552}{\normalfont}{eficiente}
\texto{417.4903}{335.0730}{11.9552}{\normalfont}{para}
\texto{444.6764}{335.0730}{11.9552}{\normalfont}{su}
\texto{460.2325}{335.0730}{11.9552}{\normalfont}{catálogo,}
\texto{510.1813}{335.0730}{11.9552}{\normalfont}{dispo-}
\texto{70.8660}{320.6270}{11.9552}{\normalfont}{nibilidad}
\texto{120.2362}{320.6270}{11.9552}{\normalfont}{y}
\texto{130.9218}{320.6270}{11.9552}{\normalfont}{reservas,}
\texto{178.7999}{320.6270}{11.9552}{\normalfont}{un}
\texto{196.3143}{320.6270}{11.9552}{\normalfont}{aspecto}
\texto{238.9298}{320.6270}{11.9552}{\normalfont}{poco}
\texto{267.4956}{320.6270}{11.9552}{\normalfont}{cubierto}
\texto{313.6139}{320.6270}{11.9552}{\normalfont}{por}
\texto{335.3533}{320.6270}{11.9552}{\normalfont}{las}
\texto{353.5826}{320.6270}{11.9552}{\normalfont}{plataformas}
\texto{418.3008}{320.6270}{11.9552}{\normalfont}{extranjeras;}
\texto{482.7059}{320.6270}{11.9552}{\normalfont}{y}
\texto{493.3915}{320.6270}{11.9552}{\normalfont}{(3)}
\texto{512.8450}{320.6270}{11.9552}{\normalfont}{opera}
\texto{70.8660}{306.1820}{11.9552}{\normalfont}{con}
\texto{91.9884}{306.1820}{11.9552}{\normalfont}{una}
\texto{114.4116}{306.1820}{11.9552}{\normalfont}{comisión}
\texto{162.2635}{306.1820}{11.9552}{\normalfont}{reducida}
\texto{209.4005}{306.1820}{11.9552}{\normalfont}{del}
\texto{227.9334}{306.1820}{11.9552}{\normalfont}{8}
\texto{235.7712}{306.1820}{11.9552}{\normalfont}{\%}
\texto{249.0893}{306.1820}{11.9552}{\normalfont}{frente}
\texto{281.9207}{306.1820}{11.9552}{\normalfont}{al}
\texto{294.6004}{306.1820}{11.9552}{\normalfont}{15}
\texto{308.2915}{306.1820}{11.9552}{\normalfont}{\%--25}
\texto{337.6032}{306.1820}{11.9552}{\normalfont}{\%}
\texto{350.9213}{306.1820}{11.9552}{\normalfont}{de}
\texto{366.1905}{306.1820}{11.9552}{\normalfont}{las}
\texto{383.4873}{306.1820}{11.9552}{\normalfont}{plataformas}
\texto{447.2731}{306.1820}{11.9552}{\normalfont}{internacionales.}
\texto{529.9218}{306.1820}{11.9552}{\normalfont}{El}
\texto{70.8660}{291.7360}{11.9552}{\normalfont}{aporte}
\texto{107.2337}{291.7360}{11.9552}{\normalfont}{de}
\texto{122.4790}{291.7360}{11.9552}{\normalfont}{las}
\texto{139.7280}{291.7360}{11.9552}{\normalfont}{investigaciones}
\texto{218.5055}{291.7360}{11.9552}{\normalfont}{consultadas}
\texto{281.3564}{291.7360}{11.9552}{\normalfont}{radica}
\texto{316.1006}{291.7360}{11.9552}{\normalfont}{en}
\texto{331.3459}{291.7360}{11.9552}{\normalfont}{la}
\texto{343.9778}{291.7360}{11.9552}{\normalfont}{validación}
\texto{398.5628}{291.7360}{11.9552}{\normalfont}{del}
\texto{417.0599}{291.7360}{11.9552}{\normalfont}{valor}
\texto{445.6304}{291.7360}{11.9552}{\normalfont}{de}
\texto{460.8638}{291.7360}{11.9552}{\normalfont}{la}
\texto{473.5076}{291.7360}{11.9552}{\normalfont}{digitalización}
\texto{70.8660}{277.2900}{11.9552}{\normalfont}{turística}
\texto{117.1015}{277.2900}{11.9552}{\normalfont}{y}
\texto{127.1774}{277.2900}{11.9552}{\normalfont}{de}
\texto{142.7933}{277.2900}{11.9552}{\normalfont}{los}
\texto{160.4128}{277.2900}{11.9552}{\normalfont}{patrones}
\texto{207.9491}{277.2900}{11.9552}{\normalfont}{de}
\texto{223.5650}{277.2900}{11.9552}{\normalfont}{seguridad}
\texto{276.3042}{277.2900}{11.9552}{\normalfont}{y}
\texto{286.3920}{277.2900}{11.9552}{\normalfont}{arquitectura}
\texto{353.0494}{277.2900}{11.9552}{\normalfont}{aplicados.}
\end{pagina}
% Página original 28
\begin{pagina}{28}
\texto{70.8660}{746.3890}{11.9552}{\normalfont}{18}
\texto{269.9360}{746.3890}{11.9552}{\normalfont}{CAPÍTULO}
\texto{336.8995}{746.3890}{11.9552}{\normalfont}{2.}
\texto{353.8113}{746.3890}{11.9552}{\normalfont}{INVESTIGACIÓN}
\texto{454.8220}{746.3890}{11.9552}{\normalfont}{DOCUMENTAL}
\texto{70.8660}{705.1940}{17.2154}{\normalfont\bfseries}{2.4.}
\texto{120.3603}{705.1940}{17.2154}{\normalfont\bfseries}{T}
\texto{132.2286}{705.1940}{17.2154}{\normalfont\bfseries}{ecnologías}
\texto{223.7370}{705.1940}{17.2154}{\normalfont\bfseries}{seleccionadas}
\texto{88.7990}{676.3850}{11.9552}{\normalfont}{A}
\texto{100.8977}{676.3850}{11.9552}{\normalfont}{continuación}
\texto{168.6119}{676.3850}{11.9552}{\normalfont}{se}
\texto{181.7674}{676.3850}{11.9552}{\normalfont}{describen}
\texto{232.9668}{676.3850}{11.9552}{\normalfont}{las}
\texto{250.0125}{676.3850}{11.9552}{\normalfont}{tecnologías}
\texto{309.3437}{676.3850}{11.9552}{\normalfont}{seleccionadas}
\texto{379.7957}{676.3850}{11.9552}{\normalfont}{para}
\texto{405.8820}{676.3850}{11.9552}{\normalfont}{desarrollar}
\texto{463.2622}{676.3850}{11.9552}{\normalfont}{el}
\texto{475.0524}{676.3850}{11.9552}{\normalfont}{prototipo}
\texto{526.1872}{676.3850}{11.9552}{\normalfont}{del}
\texto{70.8660}{661.9390}{11.9552}{\normalfont}{ecosistema,}
\texto{132.1244}{661.9390}{11.9552}{\normalfont}{justificando}
\texto{194.9442}{661.9390}{11.9552}{\normalfont}{su}
\texto{209.9743}{661.9390}{11.9552}{\normalfont}{elección.}
\texto{70.8660}{619.6230}{14.3462}{\normalfont\bfseries}{2.4.1.}
\texto{124.6642}{619.6230}{14.3462}{\normalfont\bfseries}{.NET}
\texto{168.9639}{619.6230}{14.3462}{\normalfont\bfseries}{10}
\texto{189.5865}{619.6230}{14.3462}{\normalfont\bfseries}{y}
\texto{203.4851}{619.6230}{14.3462}{\normalfont\bfseries}{C\#}
\texto{88.7990}{594.9330}{11.9552}{\normalfont\bfseries}{.NET}
\texto{124.5880}{594.9330}{11.9552}{\normalfont}{es}
\texto{137.7555}{594.9330}{11.9552}{\normalfont}{un}
\texto{154.1221}{594.9330}{11.9552}{\normalfont}{framework}
\texto{211.1389}{594.9330}{11.9552}{\normalfont}{de}
\texto{226.1929}{594.9330}{11.9552}{\normalfont}{desarrollo}
\texto{279.0444}{594.9330}{11.9552}{\normalfont}{multiplataforma}
\texto{365.0023}{594.9330}{11.9552}{\normalfont}{mantenido}
\texto{422.0071}{594.9330}{11.9552}{\normalfont}{por}
\texto{442.5987}{594.9330}{11.9552}{\normalfont}{Microsoft,}
\texto{497.3966}{594.9330}{11.9552}{\normalfont}{y}
\texto{507.0370}{594.9330}{11.9552}{\normalfont\bfseries}{C\#}
\texto{531.3150}{594.9330}{11.9552}{\normalfont}{es}
\texto{70.8660}{580.4870}{11.9552}{\normalfont}{su}
\texto{84.9397}{580.4870}{11.9552}{\normalfont}{lenguaje}
\texto{130.1877}{580.4870}{11.9552}{\normalfont}{principal.}
\texto{181.2794}{580.4870}{11.9552}{\normalfont}{Se}
\texto{195.9509}{580.4870}{11.9552}{\normalfont}{utiliza}
\texto{230.7835}{580.4870}{11.9552}{\normalfont}{para}
\texto{256.5111}{580.4870}{11.9552}{\normalfont}{construir}
\texto{305.0660}{580.4870}{11.9552}{\normalfont}{aplicaciones}
\texto{368.5792}{580.4870}{11.9552}{\normalfont}{de}
\texto{383.2386}{580.4870}{11.9552}{\normalfont}{escritorio,}
\texto{436.3460}{580.4870}{11.9552}{\normalfont}{web}
\texto{459.1494}{580.4870}{11.9552}{\normalfont}{y}
\texto{468.2927}{580.4870}{11.9552}{\normalfont}{servicios.}
\texto{517.2373}{580.4870}{11.9552}{\normalfont}{Para}
\texto{70.8660}{566.0410}{11.9552}{\normalfont}{el}
\texto{83.4333}{566.0410}{11.9552}{\normalfont}{desarrollo}
\texto{137.0500}{566.0410}{11.9552}{\normalfont}{de}
\texto{152.8811}{566.0410}{11.9552}{\normalfont}{los}
\texto{170.7278}{566.0410}{11.9552}{\normalfont}{portales}
\texto{215.5502}{566.0410}{11.9552}{\normalfont}{de}
\texto{231.3813}{566.0410}{11.9552}{\normalfont}{escritorio}
\texto{282.3965}{566.0410}{11.9552}{\normalfont}{de}
\texto{298.2240}{566.0410}{11.9552}{\normalfont\bfseries}{Explora-Nicaragua}
\texto{412.5140}{566.0410}{11.9552}{\normalfont}{(backoffice}
\texto{470.2863}{566.0410}{11.9552}{\normalfont}{de}
\texto{486.1174}{566.0410}{11.9552}{\normalfont}{operadores}
\texto{70.8660}{551.5950}{11.9552}{\normalfont}{y}
\texto{80.9538}{551.5950}{11.9552}{\normalfont}{panel}
\texto{112.1784}{551.5950}{11.9552}{\normalfont}{de}
\texto{127.8062}{551.5950}{11.9552}{\normalfont}{administración)}
\texto{211.7748}{551.5950}{11.9552}{\normalfont}{se}
\texto{225.5161}{551.5950}{11.9552}{\normalfont}{selecciona}
\texto{279.5679}{551.5950}{11.9552}{\normalfont}{.NET}
\texto{311.9354}{551.5950}{11.9552}{\normalfont}{10}
\texto{327.6757}{551.5950}{11.9552}{\normalfont}{por}
\texto{348.8292}{551.5950}{11.9552}{\normalfont}{su}
\texto{363.8593}{551.5950}{11.9552}{\normalfont}{madurez,}
\texto{414.5947}{551.5950}{11.9552}{\normalfont}{su}
\texto{429.6368}{551.5950}{11.9552}{\normalfont}{rendimiento,}
\texto{497.6093}{551.5950}{11.9552}{\normalfont}{su}
\texto{512.6513}{551.5950}{11.9552}{\normalfont}{amplia}
\texto{70.8660}{537.1490}{11.9552}{\normalfont}{documentación}
\texto{151.5062}{537.1490}{11.9552}{\normalfont}{y}
\texto{161.5821}{537.1490}{11.9552}{\normalfont}{su}
\texto{176.6121}{537.1490}{11.9552}{\normalfont}{integración}
\texto{236.7683}{537.1490}{11.9552}{\normalfont}{con}
\texto{258.2255}{537.1490}{11.9552}{\normalfont}{controles}
\texto{307.3901}{537.1490}{11.9552}{\normalfont}{ricos}
\texto{334.7770}{537.1490}{11.9552}{\normalfont}{para}
\texto{361.4371}{537.1490}{11.9552}{\normalfont}{aplicaciones}
\texto{425.8948}{537.1490}{11.9552}{\normalfont}{de}
\texto{441.4987}{537.1490}{11.9552}{\normalfont}{gestión.}
\texto{70.8660}{494.8330}{14.3462}{\normalfont\bfseries}{2.4.2.}
\texto{124.6642}{494.8330}{14.3462}{\normalfont\bfseries}{Windows}
\texto{194.7641}{494.8330}{14.3462}{\normalfont\bfseries}{F}
\texto{203.5899}{494.8330}{14.3462}{\normalfont\bfseries}{orms}
\texto{243.4336}{494.8330}{14.3462}{\normalfont\bfseries}{(WinF}
\texto{288.6413}{494.8330}{14.3462}{\normalfont\bfseries}{orms)}
\texto{88.7990}{470.1430}{11.9552}{\normalfont\bfseries}{Windows}
\texto{145.8528}{470.1430}{11.9552}{\normalfont\bfseries}{F}
\texto{153.1957}{470.1430}{11.9552}{\normalfont\bfseries}{orms}
\texto{184.6310}{470.1430}{11.9552}{\normalfont}{es}
\texto{197.1529}{470.1430}{11.9552}{\normalfont}{la}
\texto{208.9718}{470.1430}{11.9552}{\normalfont}{tecnología}
\texto{263.0643}{470.1430}{11.9552}{\normalfont}{de}
\texto{277.4846}{470.1430}{11.9552}{\normalfont}{Microsoft}
\texto{328.3732}{470.1430}{11.9552}{\normalfont}{para}
\texto{353.8497}{470.1430}{11.9552}{\normalfont}{construir}
\texto{402.1535}{470.1430}{11.9552}{\normalfont}{aplicaciones}
\texto{465.4156}{470.1430}{11.9552}{\normalfont}{de}
\texto{479.8360}{470.1430}{11.9552}{\normalfont}{escritorio}
\texto{529.4285}{470.1430}{11.9552}{\normalfont}{de}
\texto{70.8660}{455.6970}{11.9552}{\normalfont}{Windows}
\texto{120.3486}{455.6970}{11.9552}{\normalfont}{mediante}
\texto{169.4462}{455.6970}{11.9552}{\normalfont}{el}
\texto{180.5071}{455.6970}{11.9552}{\normalfont}{framework}
\texto{236.7587}{455.6970}{11.9552}{\normalfont}{.NET.}
\texto{271.0630}{455.6970}{11.9552}{\normalfont}{Ofrece}
\texto{306.5006}{455.6970}{11.9552}{\normalfont}{controles}
\texto{354.3740}{455.6970}{11.9552}{\normalfont}{ricos}
\texto{380.4459}{455.6970}{11.9552}{\normalfont}{(}
\texto{385.0000}{455.6970}{11.9552}{\normalfont\ttfamily}{DataGridView}
\texto{458.8240}{455.6970}{11.9552}{\normalfont}{,}
\texto{464.6780}{455.6970}{11.9552}{\normalfont\ttfamily}{BindingSource}
\texto{544.6530}{455.6970}{11.9552}{\normalfont}{),}
\texto{70.8660}{441.2510}{11.9552}{\normalfont}{diseño}
\texto{107.0305}{441.2510}{11.9552}{\normalfont}{visual}
\texto{140.9187}{441.2510}{11.9552}{\normalfont}{asistido}
\texto{183.6633}{441.2510}{11.9552}{\normalfont}{y}
\texto{194.0739}{441.2510}{11.9552}{\normalfont}{un}
\texto{211.3133}{441.2510}{11.9552}{\normalfont}{ciclo}
\texto{238.3201}{441.2510}{11.9552}{\normalfont}{de}
\texto{254.2588}{441.2510}{11.9552}{\normalfont}{desarrollo}
\texto{307.9830}{441.2510}{11.9552}{\normalfont}{rápido.}
\texto{347.9971}{441.2510}{11.9552}{\normalfont}{Se}
\texto{363.9358}{441.2510}{11.9552}{\normalfont}{seleccionó}
\texto{418.3104}{441.2510}{11.9552}{\normalfont}{para}
\texto{445.3172}{441.2510}{11.9552}{\normalfont}{los}
\texto{463.2715}{441.2510}{11.9552}{\normalfont}{dos}
\texto{484.4777}{441.2510}{11.9552}{\normalfont}{portales}
\texto{529.4316}{441.2510}{11.9552}{\normalfont}{de}
\texto{70.8660}{426.8050}{11.9552}{\normalfont}{escritorio}
\texto{121.7139}{426.8050}{11.9552}{\normalfont}{por}
\texto{142.9033}{426.8050}{11.9552}{\normalfont}{su}
\texto{157.9812}{426.8050}{11.9552}{\normalfont}{idoneidad}
\texto{211.3659}{426.8050}{11.9552}{\normalfont}{para}
\texto{238.0858}{426.8050}{11.9552}{\normalfont}{sistemas}
\texto{284.5102}{426.8050}{11.9552}{\normalfont}{de}
\texto{300.1739}{426.8050}{11.9552}{\normalfont}{gestión}
\texto{339.9656}{426.8050}{11.9552}{\normalfont}{interna}
\texto{380.0203}{426.8050}{11.9552}{\normalfont}{intensivos}
\texto{433.8522}{426.8050}{11.9552}{\normalfont}{en}
\texto{449.5278}{426.8050}{11.9552}{\normalfont}{datos}
\texto{480.8648}{426.8050}{11.9552}{\normalfont}{(grillas,}
\texto{523.2579}{426.8050}{11.9552}{\normalfont}{for-}
\texto{70.8660}{412.3590}{11.9552}{\normalfont}{mularios}
\texto{117.7447}{412.3590}{11.9552}{\normalfont}{CRUD,}
\texto{159.0332}{412.3590}{11.9552}{\normalfont}{reportes)}
\texto{208.5206}{412.3590}{11.9552}{\normalfont}{y}
\texto{218.2736}{412.3590}{11.9552}{\normalfont}{por}
\texto{239.0804}{412.3590}{11.9552}{\normalfont}{la}
\texto{251.7601}{412.3590}{11.9552}{\normalfont}{productividad}
\texto{326.8723}{412.3590}{11.9552}{\normalfont}{que}
\texto{348.3318}{412.3590}{11.9552}{\normalfont}{aporta}
\texto{385.3977}{412.3590}{11.9552}{\normalfont}{al}
\texto{398.0774}{412.3590}{11.9552}{\normalfont}{equipo,}
\texto{438.7203}{412.3590}{11.9552}{\normalfont}{en}
\texto{454.0015}{412.3590}{11.9552}{\normalfont}{línea}
\texto{481.6395}{412.3590}{11.9552}{\normalfont}{con}
\texto{502.7739}{412.3590}{11.9552}{\normalfont}{el}
\texto{514.8032}{412.3590}{11.9552}{\normalfont}{perfil}
\texto{70.8660}{397.9130}{11.9552}{\normalfont}{administrativo}
\texto{148.9717}{397.9130}{11.9552}{\normalfont}{de}
\texto{164.5756}{397.9130}{11.9552}{\normalfont}{los}
\texto{182.1952}{397.9130}{11.9552}{\normalfont}{procesos}
\texto{228.8301}{397.9130}{11.9552}{\normalfont}{de}
\texto{244.4340}{397.9130}{11.9552}{\normalfont}{gestión}
\texto{284.1779}{397.9130}{11.9552}{\normalfont}{de}
\texto{299.7818}{397.9130}{11.9552}{\normalfont}{tours}
\texto{329.7702}{397.9130}{11.9552}{\normalfont}{y}
\texto{339.8460}{397.9130}{11.9552}{\normalfont}{reservas.}
\texto{70.8660}{355.5980}{14.3462}{\normalfont\bfseries}{2.4.3.}
\texto{124.6642}{355.5980}{14.3462}{\normalfont\bfseries}{\parbox[t]{470.268bp}{\fontsize{14.3462bp}{14.446bp}\selectfont Arquitectura por capas}}
\texto{70.8660}{330.9070}{11.9552}{\normalfont}{\parbox[t]{470.268bp}{\fontsize{11.9552bp}{14.446bp}\selectfont La arquitectura separa dominio, aplicación, infraestructura y presentación. El dominio contiene las entidades; aplicación define contratos y validadores; infraestructura implementa servicios y persistencia; presentación incorpora la API y los clientes Windows Forms. La inyección de dependencias conecta las interfaces con sus implementaciones.}}
\texto{70.8660}{245.2530}{14.3462}{\normalfont\bfseries}{2.4.4.}
\texto{124.6642}{245.2530}{14.3462}{\normalfont\bfseries}{React}
\texto{170.7414}{245.2530}{14.3462}{\normalfont\bfseries}{Native}
\texto{88.7990}{220.5630}{11.9552}{\normalfont\bfseries}{React}
\texto{128.0934}{220.5630}{11.9552}{\normalfont\bfseries}{Native}
\texto{171.6730}{220.5630}{11.9552}{\normalfont}{es}
\texto{186.1794}{220.5630}{11.9552}{\normalfont}{un}
\texto{203.8731}{220.5630}{11.9552}{\normalfont}{framework}
\texto{262.2169}{220.5630}{11.9552}{\normalfont}{de}
\texto{278.6099}{220.5630}{11.9552}{\normalfont}{código}
\texto{316.1492}{220.5630}{11.9552}{\normalfont}{abierto}
\texto{356.6056}{220.5630}{11.9552}{\normalfont}{para}
\texto{384.0547}{220.5630}{11.9552}{\normalfont}{construir}
\texto{434.3311}{220.5630}{11.9552}{\normalfont}{apps}
\texto{462.5072}{220.5630}{11.9552}{\normalfont}{móviles}
\texto{504.9816}{220.5630}{11.9552}{\normalfont}{nativas}
\texto{70.8660}{206.1170}{11.9552}{\normalfont}{usando}
\texto{111.7815}{206.1170}{11.9552}{\normalfont}{JavaScript}
\texto{170.3476}{206.1170}{11.9552}{\normalfont}{y}
\texto{181.6070}{206.1170}{11.9552}{\normalfont}{React}
\texto{216.1145}{206.1170}{11.9552}{\normalfont}{[5].}
\texto{236.8076}{206.1170}{11.9552}{\normalfont}{Se}
\texto{253.5951}{206.1170}{11.9552}{\normalfont}{selecciona}
\texto{308.8185}{206.1170}{11.9552}{\normalfont}{para}
\texto{336.6622}{206.1170}{11.9552}{\normalfont}{la}
\texto{350.8482}{206.1170}{11.9552}{\normalfont}{app}
\texto{374.7897}{206.1170}{11.9552}{\normalfont}{móvil}
\texto{407.8387}{206.1170}{11.9552}{\normalfont}{del}
\texto{427.8780}{206.1170}{11.9552}{\normalfont}{turista}
\texto{466.8424}{206.1170}{11.9552}{\normalfont}{por}
\texto{489.1675}{206.1170}{11.9552}{\normalfont}{su}
\texto{505.3692}{206.1170}{11.9552}{\normalfont}{capaci-}
\texto{70.8660}{191.6710}{11.9552}{\normalfont}{dad}
\texto{94.0782}{191.6710}{11.9552}{\normalfont}{multiplataforma;}
\texto{184.2802}{191.6710}{11.9552}{\normalfont}{esta}
\texto{208.8697}{191.6710}{11.9552}{\normalfont}{aplicación}
\texto{263.9497}{191.6710}{11.9552}{\normalfont}{forma}
\texto{297.9048}{191.6710}{11.9552}{\normalfont}{parte}
\texto{328.9214}{191.6710}{11.9552}{\normalfont}{de}
\texto{344.9916}{191.6710}{11.9552}{\normalfont}{la}
\texto{358.4483}{191.6710}{11.9552}{\normalfont}{visión}
\texto{392.4681}{191.6710}{11.9552}{\normalfont}{completa}
\texto{442.9955}{191.6710}{11.9552}{\normalfont}{del}
\texto{462.3175}{191.6710}{11.9552}{\normalfont}{ecosistema}
\texto{520.7784}{191.6710}{11.9552}{\normalfont}{y}
\texto{531.3205}{191.6710}{11.9552}{\normalfont}{se}
\texto{70.8660}{177.2250}{11.9552}{\normalfont}{contempla}
\texto{127.1200}{177.2250}{11.9552}{\normalfont}{como}
\texto{157.6823}{177.2250}{11.9552}{\normalfont}{trabajo}
\texto{198.6599}{177.2250}{11.9552}{\normalfont}{futuro}
\texto{234.0999}{177.2250}{11.9552}{\normalfont}{del}
\texto{252.9557}{177.2250}{11.9552}{\normalfont}{proyecto.}
\texto{70.8660}{134.9090}{14.3462}{\normalfont\bfseries}{2.4.5.}
\texto{124.6642}{134.9090}{14.3462}{\normalfont\bfseries}{API}
\texto{159.2673}{134.9090}{14.3462}{\normalfont\bfseries}{REST,}
\texto{211.9523}{134.9090}{14.3462}{\normalfont\bfseries}{JWT}
\texto{253.5620}{134.9090}{14.3462}{\normalfont\bfseries}{y}
\texto{267.4606}{134.9090}{14.3462}{\normalfont\bfseries}{OAuth}
\texto{320.4540}{134.9090}{14.3462}{\normalfont\bfseries}{2.0}
\texto{88.7990}{110.2190}{11.9552}{\normalfont}{La}
\texto{104.8238}{110.2190}{11.9552}{\normalfont}{arquitectura}
\texto{170.4530}{110.2190}{11.9552}{\normalfont}{orientada}
\texto{221.1263}{110.2190}{11.9552}{\normalfont}{a}
\texto{229.8369}{110.2190}{11.9552}{\normalfont}{servicios}
\texto{275.4340}{110.2190}{11.9552}{\normalfont}{se}
\texto{288.1233}{110.2190}{11.9552}{\normalfont}{materializa}
\texto{347.5741}{110.2190}{11.9552}{\normalfont}{en}
\texto{362.1379}{110.2190}{11.9552}{\normalfont}{una}
\texto{384.0110}{110.2190}{11.9552}{\normalfont\bfseries}{API}
\texto{411.6514}{110.2190}{11.9552}{\normalfont\bfseries}{REST}
\texto{450.2080}{110.2190}{11.9552}{\normalfont}{[14]}
\texto{471.2751}{110.2190}{11.9552}{\normalfont}{que}
\texto{492.0293}{110.2190}{11.9552}{\normalfont}{centraliza}
\texto{70.8660}{95.7730}{11.9552}{\normalfont}{la}
\texto{84.6575}{95.7730}{11.9552}{\normalfont}{lógica}
\texto{118.6103}{95.7730}{11.9552}{\normalfont}{de}
\texto{135.0033}{95.7730}{11.9552}{\normalfont}{negocio.}
\texto{180.9973}{95.7730}{11.9552}{\normalfont}{La}
\texto{198.8512}{95.7730}{11.9552}{\normalfont}{autenticación}
\texto{271.5029}{95.7730}{11.9552}{\normalfont}{se}
\texto{286.0094}{95.7730}{11.9552}{\normalfont}{basa}
\texto{313.5350}{95.7730}{11.9552}{\normalfont}{en}
\texto{329.9820}{95.7730}{11.9552}{\normalfont\bfseries}{JWT}
\texto{364.8620}{95.7730}{11.9552}{\normalfont}{y}
\texto{375.8360}{95.7730}{11.9552}{\normalfont\bfseries}{OAuth}
\texto{420.9059}{95.7730}{11.9552}{\normalfont\bfseries}{2.0}
\texto{442.7820}{95.7730}{11.9552}{\normalfont}{[7],}
\texto{463.0771}{95.7730}{11.9552}{\normalfont}{emitidos}
\texto{510.7520}{95.7730}{11.9552}{\normalfont}{por}
\texto{532.6707}{95.7730}{11.9552}{\normalfont}{el}
\texto{70.8660}{81.3270}{11.9552}{\normalfont}{backend}
\texto{115.3609}{81.3270}{11.9552}{\normalfont}{e}
\texto{123.7678}{81.3270}{11.9552}{\normalfont}{incluidos}
\texto{171.9233}{81.3270}{11.9552}{\normalfont}{en}
\texto{186.8338}{81.3270}{11.9552}{\normalfont}{las}
\texto{203.7600}{81.3270}{11.9552}{\normalfont}{peticiones}
\texto{257.4293}{81.3270}{11.9552}{\normalfont}{de}
\texto{272.3398}{81.3270}{11.9552}{\normalfont}{todos}
\texto{303.2584}{81.3270}{11.9552}{\normalfont}{los}
\texto{320.1845}{81.3270}{11.9552}{\normalfont}{clientes.}
\texto{364.1031}{81.3270}{11.9552}{\normalfont}{Asimismo,}
\texto{420.4360}{81.3270}{11.9552}{\normalfont}{se}
\texto{433.4600}{81.3270}{11.9552}{\normalfont}{define}
\texto{466.5927}{81.3270}{11.9552}{\normalfont}{un}
\texto{482.8570}{81.3270}{11.9552}{\normalfont\bfseries}{control}
\texto{527.5156}{81.3270}{11.9552}{\normalfont\bfseries}{de}
\texto{70.8660}{66.8810}{11.9552}{\normalfont\bfseries}{acceso}
\texto{112.0098}{66.8810}{11.9552}{\normalfont\bfseries}{basado}
\texto{156.5441}{66.8810}{11.9552}{\normalfont\bfseries}{en}
\texto{174.6359}{66.8810}{11.9552}{\normalfont\bfseries}{roles}
\texto{206.5157}{66.8810}{11.9552}{\normalfont\bfseries}{(RBAC)}
\texto{259.9660}{66.8810}{11.9552}{\normalfont}{para}
\texto{286.6261}{66.8810}{11.9552}{\normalfont}{garantizar}
\texto{342.2297}{66.8810}{11.9552}{\normalfont}{permisos}
\texto{390.8157}{66.8810}{11.9552}{\normalfont}{según}
\texto{423.3936}{66.8810}{11.9552}{\normalfont}{el}
\texto{435.7457}{66.8810}{11.9552}{\normalfont}{rol}
\texto{453.3127}{66.8810}{11.9552}{\normalfont}{de}
\texto{468.9166}{66.8810}{11.9552}{\normalfont}{cada}
\texto{496.2390}{66.8810}{11.9552}{\normalfont}{usuario.}
\end{pagina}
% Página original 29
\begin{pagina}{29}
\texto{70.8660}{746.3890}{11.9552}{\normalfont}{2.4.}
\texto{96.8710}{746.3890}{11.9552}{\normalfont}{TECNOLOGÍAS}
\texto{188.6403}{746.3890}{11.9552}{\normalfont}{SELECCIONADAS}
\texto{529.4280}{746.3890}{11.9552}{\normalfont}{19}
\texto{70.8660}{705.1940}{14.3462}{\normalfont\bfseries}{2.4.6.}
\texto{124.6642}{705.1940}{14.3462}{\normalfont\bfseries}{Pasarela}
\texto{188.9539}{705.1940}{14.3462}{\normalfont\bfseries}{de}
\texto{210.6640}{705.1940}{14.3462}{\normalfont\bfseries}{pago}
\texto{248.9956}{705.1940}{14.3462}{\normalfont\bfseries}{Stripe}
\texto{297.4829}{705.1940}{14.3462}{\normalfont\bfseries}{y}
\texto{311.3815}{705.1940}{14.3462}{\normalfont\bfseries}{SSL/TLS}
\texto{88.7990}{683.0270}{11.9552}{\normalfont\bfseries}{Stripe}
\texto{129.9770}{683.0270}{11.9552}{\normalfont}{se}
\texto{145.0573}{683.0270}{11.9552}{\normalfont}{selecciona}
\texto{200.4601}{683.0270}{11.9552}{\normalfont}{como}
\texto{232.3972}{683.0270}{11.9552}{\normalfont}{pasarela}
\texto{279.3453}{683.0270}{11.9552}{\normalfont}{de}
\texto{296.3121}{683.0270}{11.9552}{\normalfont}{pago}
\texto{325.6358}{683.0270}{11.9552}{\normalfont}{por}
\texto{348.1403}{683.0270}{11.9552}{\normalfont}{su}
\texto{364.5213}{683.0270}{11.9552}{\normalfont}{certificación}
\texto{430.9276}{683.0270}{11.9552}{\normalfont}{PCI}
\texto{456.8298}{683.0270}{11.9552}{\normalfont}{DSS,}
\texto{487.2892}{683.0270}{11.9552}{\normalfont}{su}
\texto{503.6822}{683.0270}{11.9552}{\normalfont}{soporte}
\texto{70.8660}{668.5810}{11.9552}{\normalfont}{para}
\texto{98.1119}{668.5810}{11.9552}{\normalfont}{múltiples}
\texto{149.1510}{668.5810}{11.9552}{\normalfont}{métodos}
\texto{196.3071}{668.5810}{11.9552}{\normalfont}{de}
\texto{212.4969}{668.5810}{11.9552}{\normalfont}{pago}
\texto{241.0435}{668.5810}{11.9552}{\normalfont}{y}
\texto{251.7051}{668.5810}{11.9552}{\normalfont}{su}
\texto{267.3091}{668.5810}{11.9552}{\normalfont}{facilidad}
\texto{315.0414}{668.5810}{11.9552}{\normalfont}{de}
\texto{331.2311}{668.5810}{11.9552}{\normalfont}{integración}
\texto{391.9731}{668.5810}{11.9552}{\normalfont}{vía}
\texto{411.7398}{668.5810}{11.9552}{\normalfont}{API.}
\texto{440.4371}{668.5810}{11.9552}{\normalfont}{Las}
\texto{462.7049}{668.5810}{11.9552}{\normalfont}{comunicaciones}
\texto{70.8660}{654.1350}{11.9552}{\normalfont}{se}
\texto{84.5834}{654.1350}{11.9552}{\normalfont}{aseguran}
\texto{133.4204}{654.1350}{11.9552}{\normalfont}{mediante}
\texto{183.8300}{654.1350}{11.9552}{\normalfont\bfseries}{SSL/TLS}
\texto{238.4860}{654.1350}{11.9552}{\normalfont}{.}
\texto{70.8660}{619.4730}{14.3462}{\normalfont\bfseries}{2.4.7.}
\texto{124.6642}{619.4730}{14.3462}{\normalfont\bfseries}{Base}
\texto{163.0949}{619.4730}{14.3462}{\normalfont\bfseries}{de}
\texto{184.8050}{619.4730}{14.3462}{\normalfont\bfseries}{datos}
\texto{227.7101}{619.4730}{14.3462}{\normalfont\bfseries}{y}
\texto{241.6087}{619.4730}{14.3462}{\normalfont\bfseries}{herramientas}
\texto{338.6421}{619.4730}{14.3462}{\normalfont\bfseries}{de}
\texto{360.3522}{619.4730}{14.3462}{\normalfont\bfseries}{soporte}
\texto{88.7990}{597.3070}{11.9552}{\normalfont}{Para}
\texto{117.5776}{597.3070}{11.9552}{\normalfont}{el}
\texto{130.9339}{597.3070}{11.9552}{\normalfont}{modelo}
\texto{172.5667}{597.3070}{11.9552}{\normalfont}{de}
\texto{189.1629}{597.3070}{11.9552}{\normalfont}{datos}
\texto{221.4324}{597.3070}{11.9552}{\normalfont}{centralizado}
\texto{287.7837}{597.3070}{11.9552}{\normalfont}{se}
\texto{302.4934}{597.3070}{11.9552}{\normalfont}{emplea}
\texto{343.1531}{597.3070}{11.9552}{\normalfont}{un}
\texto{361.0500}{597.3070}{11.9552}{\normalfont}{gestor}
\texto{396.5713}{597.3070}{11.9552}{\normalfont}{de}
\texto{413.1794}{597.3070}{11.9552}{\normalfont}{base}
\texto{440.2460}{597.3070}{11.9552}{\normalfont}{de}
\texto{456.8422}{597.3070}{11.9552}{\normalfont}{datos}
\texto{489.1117}{597.3070}{11.9552}{\normalfont}{relacional.}
\texto{70.8660}{582.8610}{11.9552}{\normalfont}{Como}
\texto{105.8995}{582.8610}{11.9552}{\normalfont}{herramientas}
\texto{177.1071}{582.8610}{11.9552}{\normalfont}{de}
\texto{193.9424}{582.8610}{11.9552}{\normalfont}{desarrollo}
\texto{248.5633}{582.8610}{11.9552}{\normalfont}{se}
\texto{263.5121}{582.8610}{11.9552}{\normalfont}{utilizan}
\texto{307.0123}{582.8610}{11.9552}{\normalfont}{un}
\texto{325.1483}{582.8610}{11.9552}{\normalfont}{IDE}
\texto{351.4043}{582.8610}{11.9552}{\normalfont}{moderno,}
\texto{404.3324}{582.8610}{11.9552}{\normalfont}{herramientas}
\texto{475.5400}{582.8610}{11.9552}{\normalfont}{de}
\texto{492.3753}{582.8610}{11.9552}{\normalfont}{diseño}
\texto{529.4364}{582.8610}{11.9552}{\normalfont}{de}
\texto{70.8660}{568.4150}{11.9552}{\normalfont}{interfaces}
\texto{124.2723}{568.4150}{11.9552}{\normalfont}{y}
\texto{135.6632}{568.4150}{11.9552}{\normalfont}{herramientas}
\texto{206.9544}{568.4150}{11.9552}{\normalfont}{de}
\texto{223.8734}{568.4150}{11.9552}{\normalfont}{modelado}
\texto{278.1979}{568.4150}{11.9552}{\normalfont}{de}
\texto{295.1169}{568.4150}{11.9552}{\normalfont}{procesos}
\texto{343.0668}{568.4150}{11.9552}{\normalfont}{(BPMN),}
\texto{396.3917}{568.4150}{11.9552}{\normalfont}{complementando}
\texto{487.7797}{568.4150}{11.9552}{\normalfont}{el}
\texto{501.4588}{568.4150}{11.9552}{\normalfont}{ciclo}
\texto{529.4340}{568.4150}{11.9552}{\normalfont}{de}
\texto{70.8660}{553.9690}{11.9552}{\normalfont}{desarrollo.}
\end{pagina}
% Página original 30
\begin{pagina}{30}
\texto{70.8660}{746.3890}{11.9552}{\normalfont}{20}
\texto{269.9360}{746.3890}{11.9552}{\normalfont}{CAPÍTULO}
\texto{336.8995}{746.3890}{11.9552}{\normalfont}{2.}
\texto{353.8113}{746.3890}{11.9552}{\normalfont}{INVESTIGACIÓN}
\texto{454.8220}{746.3890}{11.9552}{\normalfont}{DOCUMENTAL}
\texto{70.8660}{705.1940}{17.2154}{\normalfont\bfseries}{2.5.}
\texto{120.3603}{705.1940}{17.2154}{\normalfont\bfseries}{Análisis}
\texto{192.1743}{705.1940}{17.2154}{\normalfont\bfseries}{de}
\texto{217.3828}{705.1940}{17.2154}{\normalfont\bfseries}{las}
\texto{245.4302}{705.1940}{17.2154}{\normalfont\bfseries}{tecnologías}
\texto{343.6440}{705.1940}{17.2154}{\normalfont\bfseries}{y}
\texto{359.4787}{705.1940}{17.2154}{\normalfont\bfseries}{su}
\texto{383.4736}{705.1940}{17.2154}{\normalfont\bfseries}{pertinencia}
\texto{483.7963}{705.1940}{17.2154}{\normalfont\bfseries}{para}
\texto{526.9088}{705.1940}{17.2154}{\normalfont\bfseries}{el}
\texto{120.3600}{683.2760}{17.2154}{\normalfont\bfseries}{proyecto}
\texto{88.7990}{656.3720}{11.9552}{\normalfont}{Con}
\texto{114.3688}{656.3720}{11.9552}{\normalfont}{el}
\texto{127.5817}{656.3720}{11.9552}{\normalfont}{fin}
\texto{145.3471}{656.3720}{11.9552}{\normalfont}{de}
\texto{161.8118}{656.3720}{11.9552}{\normalfont}{seleccionar}
\texto{221.2650}{656.3720}{11.9552}{\normalfont}{la}
\texto{235.1283}{656.3720}{11.9552}{\normalfont}{tecnología}
\texto{291.2651}{656.3720}{11.9552}{\normalfont}{más}
\texto{316.2491}{656.3720}{11.9552}{\normalfont}{adecuada}
\texto{368.4837}{656.3720}{11.9552}{\normalfont}{para}
\texto{396.0046}{656.3720}{11.9552}{\normalfont}{los}
\texto{414.4849}{656.3720}{11.9552}{\normalfont}{portales}
\texto{459.9530}{656.3720}{11.9552}{\normalfont}{de}
\texto{476.4177}{656.3720}{11.9552}{\normalfont}{escritorio,}
\texto{531.3183}{656.3720}{11.9552}{\normalfont}{se}
\texto{70.8660}{641.9260}{11.9552}{\normalfont}{evaluaron}
\texto{125.2335}{641.9260}{11.9552}{\normalfont}{alternativas}
\texto{190.0713}{641.9260}{11.9552}{\normalfont}{considerando}
\texto{261.7451}{641.9260}{11.9552}{\normalfont}{las}
\texto{280.7276}{641.9260}{11.9552}{\normalfont}{características}
\texto{358.3192}{641.9260}{11.9552}{\normalfont}{del}
\texto{378.5498}{641.9260}{11.9552}{\normalfont}{proyecto,}
\texto{430.3158}{641.9260}{11.9552}{\normalfont}{los}
\texto{449.3102}{641.9260}{11.9552}{\normalfont}{conocimientos}
\texto{526.1870}{641.9260}{11.9552}{\normalfont}{del}
\texto{70.8660}{627.4800}{11.9552}{\normalfont}{equipo,}
\texto{112.4773}{627.4800}{11.9552}{\normalfont}{la}
\texto{126.1253}{627.4800}{11.9552}{\normalfont}{facilidad}
\texto{173.9055}{627.4800}{11.9552}{\normalfont}{de}
\texto{190.1550}{627.4800}{11.9552}{\normalfont}{desarrollo,}
\texto{247.4419}{627.4800}{11.9552}{\normalfont}{la}
\texto{261.0900}{627.4800}{11.9552}{\normalfont}{disponibilidad}
\texto{337.5602}{627.4800}{11.9552}{\normalfont}{de}
\texto{353.8097}{627.4800}{11.9552}{\normalfont}{herramientas,}
\texto{427.6833}{627.4800}{11.9552}{\normalfont}{la}
\texto{441.3313}{627.4800}{11.9552}{\normalfont}{compatibilidad}
\texto{521.9668}{627.4800}{11.9552}{\normalfont}{y}
\texto{532.6882}{627.4800}{11.9552}{\normalfont}{el}
\texto{70.8660}{613.0340}{11.9552}{\normalfont}{tiempo}
\texto{110.2058}{613.0340}{11.9552}{\normalfont}{disponible.}
\texto{70.8660}{576.4950}{14.3462}{\normalfont\bfseries}{2.5.1.}
\texto{124.6642}{576.4950}{14.3462}{\normalfont\bfseries}{Comparación}
\texto{223.0189}{576.4950}{14.3462}{\normalfont\bfseries}{de}
\texto{244.5856}{576.4950}{14.3462}{\normalfont\bfseries}{alternativas}
\texto{331.7976}{576.4950}{14.3462}{\normalfont\bfseries}{para}
\texto{368.2842}{576.4950}{14.3462}{\normalfont\bfseries}{los}
\texto{392.4404}{576.4950}{14.3462}{\normalfont\bfseries}{portales}
\texto{454.0860}{576.4950}{14.3462}{\normalfont\bfseries}{de}
\texto{475.6670}{576.4950}{14.3462}{\normalfont\bfseries}{escritorio}
\texto{124.9140}{532.3050}{11.9552}{\normalfont}{Cuadro}
\texto{166.5324}{532.3050}{11.9552}{\normalfont}{2.1:}
\texto{188.6470}{532.3050}{11.9552}{\normalfont}{Comparación}
\texto{260.1821}{532.3050}{11.9552}{\normalfont}{de}
\texto{275.7861}{532.3050}{11.9552}{\normalfont}{tecnologías}
\texto{335.6912}{532.3050}{11.9552}{\normalfont}{para}
\texto{362.3513}{532.3050}{11.9552}{\normalfont}{los}
\texto{379.9708}{532.3050}{11.9552}{\normalfont}{portales}
\texto{424.5900}{532.3050}{11.9552}{\normalfont}{de}
\texto{440.1940}{532.3050}{11.9552}{\normalfont}{escritorio}
\texto{89.6420}{514.5900}{11.9552}{\normalfont\bfseries}{Criterio}
\texto{192.3060}{514.5900}{11.9552}{\normalfont\bfseries}{.NET}
\texto{229.2105}{514.5900}{11.9552}{\normalfont\bfseries}{WinF}
\texto{261.6653}{514.5900}{11.9552}{\normalfont\bfseries}{orms}
\texto{306.2976}{514.5900}{11.9552}{\normalfont\bfseries}{Java}
\texto{336.7906}{514.5900}{11.9552}{\normalfont\bfseries}{Swing}
\texto{420.3110}{514.5900}{11.9552}{\normalfont\bfseries}{Python}
\texto{467.9680}{514.5900}{11.9552}{\normalfont\bfseries}{Tkinter}
\texto{89.6420}{494.1540}{11.9552}{\normalfont}{Madurez}
\texto{140.0332}{494.1540}{11.9552}{\normalfont}{y}
\texto{152.0577}{494.1540}{11.9552}{\normalfont}{docu-}
\texto{89.6420}{479.7090}{11.9552}{\normalfont}{mentación}
\texto{192.3060}{494.1540}{11.9552}{\normalfont}{Alta}
\texto{306.3080}{494.1540}{11.9552}{\normalfont}{Alta}
\texto{420.3110}{494.1540}{11.9552}{\normalfont}{Media}
\texto{89.6420}{465.2630}{11.9552}{\normalfont}{Controles}
\texto{156.8709}{465.2630}{11.9552}{\normalfont}{ricos}
\texto{89.6420}{450.8170}{11.9552}{\normalfont}{para}
\texto{127.9704}{450.8170}{11.9552}{\normalfont}{grillas}
\texto{174.1677}{450.8170}{11.9552}{\normalfont}{y}
\texto{89.6420}{436.3710}{11.9552}{\normalfont}{CRUD}
\texto{192.3060}{465.2630}{11.9552}{\normalfont}{Muy}
\texto{219.6117}{465.2630}{11.9552}{\normalfont}{buenos}
\texto{306.3080}{465.2630}{11.9552}{\normalfont}{Buenos}
\texto{420.3110}{465.2630}{11.9552}{\normalfont}{Básicos}
\texto{89.6420}{421.9250}{11.9552}{\normalfont}{Curva}
\texto{123.4441}{421.9250}{11.9552}{\normalfont}{de}
\texto{138.0677}{421.9250}{11.9552}{\normalfont}{aprendi-}
\texto{89.6420}{407.4800}{11.9552}{\normalfont}{zaje}
\texto{113.6983}{407.4800}{11.9552}{\normalfont}{del}
\texto{132.5540}{407.4800}{11.9552}{\normalfont}{equipo}
\texto{192.3060}{421.9250}{11.9552}{\normalfont}{Adecuada}
\texto{246.2789}{421.9250}{11.9552}{\normalfont}{(C\#)}
\texto{306.3036}{421.9250}{11.9552}{\normalfont}{Media}
\texto{341.7388}{421.9250}{11.9552}{\normalfont}{(Java)}
\texto{420.3110}{421.9250}{11.9552}{\normalfont}{Baja}
\texto{448.1044}{421.9250}{11.9552}{\normalfont}{(Python)}
\texto{89.6420}{393.0340}{11.9552}{\normalfont}{Rendimiento}
\texto{192.3060}{393.0340}{11.9552}{\normalfont}{Alto}
\texto{306.3080}{393.0340}{11.9552}{\normalfont}{Alto}
\texto{420.3110}{393.0340}{11.9552}{\normalfont}{Medio}
\texto{89.6420}{378.5880}{11.9552}{\normalfont}{Soporte}
\texto{132.6688}{378.5880}{11.9552}{\normalfont}{de}
\texto{148.0455}{378.5880}{11.9552}{\normalfont}{la}
\texto{160.8328}{378.5880}{11.9552}{\normalfont}{pla-}
\texto{89.6420}{364.1420}{11.9552}{\normalfont}{taforma}
\texto{192.3060}{378.5880}{11.9552}{\normalfont}{Windows}
\texto{243.0917}{378.5880}{11.9552}{\normalfont}{nativo}
\texto{306.3084}{378.5880}{11.9552}{\normalfont}{Multiplataforma}
\texto{420.3036}{378.5880}{11.9552}{\normalfont}{Multiplataforma}
\texto{88.7990}{329.9910}{11.9552}{\normalfont}{Fuente:}
\texto{130.2501}{329.9910}{11.9552}{\normalfont}{Elaboración}
\texto{195.1262}{329.9910}{11.9552}{\normalfont}{propia}
\texto{231.5417}{329.9910}{11.9552}{\normalfont}{(2026).}
\texto{70.8660}{293.4510}{14.3462}{\normalfont\bfseries}{2.5.2.}
\texto{124.6642}{293.4510}{14.3462}{\normalfont\bfseries}{Justificación}
\texto{217.1499}{293.4510}{14.3462}{\normalfont\bfseries}{de}
\texto{238.8600}{293.4510}{14.3462}{\normalfont\bfseries}{la}
\texto{256.5690}{293.4510}{14.3462}{\normalfont\bfseries}{decisión}
\texto{88.7990}{270.6650}{11.9552}{\normalfont}{La}
\texto{105.1465}{270.6650}{11.9552}{\normalfont}{tabla}
\texto{134.3430}{270.6650}{11.9552}{\normalfont}{2.1}
\texto{152.4830}{270.6650}{11.9552}{\normalfont}{compara}
\texto{199.2374}{270.6650}{11.9552}{\normalfont}{las}
\texto{216.1397}{270.6650}{11.9552}{\normalfont}{alternativas}
\texto{278.8853}{270.6650}{11.9552}{\normalfont}{consideradas}
\texto{346.5924}{270.6650}{11.9552}{\normalfont}{para}
\texto{372.5352}{270.6650}{11.9552}{\normalfont}{el}
\texto{384.1700}{270.6650}{11.9552}{\normalfont}{desarrollo}
\texto{436.8422}{270.6650}{11.9552}{\normalfont}{de}
\texto{451.7288}{270.6650}{11.9552}{\normalfont}{los}
\texto{468.6311}{270.6650}{11.9552}{\normalfont}{portales}
\texto{512.5330}{270.6650}{11.9552}{\normalfont}{de}
\texto{527.4196}{270.6650}{11.9552}{\normalfont}{es-}
\texto{70.8660}{256.2190}{11.9552}{\normalfont}{critorio.}
\texto{114.3566}{256.2190}{11.9552}{\normalfont}{Se}
\texto{129.2432}{256.2190}{11.9552}{\normalfont}{selecciona}
\texto{182.5630}{256.2190}{11.9552}{\normalfont\bfseries}{.NET}
\texto{218.6426}{256.2190}{11.9552}{\normalfont\bfseries}{10}
\texto{234.9913}{256.2190}{11.9552}{\normalfont\bfseries}{con}
\texto{258.8121}{256.2190}{11.9552}{\normalfont\bfseries}{Windows}
\texto{316.4039}{256.2190}{11.9552}{\normalfont\bfseries}{F}
\texto{323.7467}{256.2190}{11.9552}{\normalfont\bfseries}{orms}
\texto{355.6456}{256.2190}{11.9552}{\normalfont}{porque}
\texto{393.9309}{256.2190}{11.9552}{\normalfont}{ofrece}
\texto{426.7025}{256.2190}{11.9552}{\normalfont}{el}
\texto{438.3254}{256.2190}{11.9552}{\normalfont}{mejor}
\texto{470.4466}{256.2190}{11.9552}{\normalfont}{equilibrio}
\texto{521.4284}{256.2190}{11.9552}{\normalfont}{entre}
\texto{70.8660}{241.7730}{11.9552}{\normalfont}{madurez,}
\texto{122.3068}{241.7730}{11.9552}{\normalfont}{controles}
\texto{172.1887}{241.7730}{11.9552}{\normalfont}{ricos}
\texto{200.2810}{241.7730}{11.9552}{\normalfont}{orientados}
\texto{257.3289}{241.7730}{11.9552}{\normalfont}{a}
\texto{267.7968}{241.7730}{11.9552}{\normalfont}{sistemas}
\texto{314.8788}{241.7730}{11.9552}{\normalfont}{de}
\texto{331.2120}{241.7730}{11.9552}{\normalfont}{gestión}
\texto{371.6612}{241.7730}{11.9552}{\normalfont}{(}
\texto{376.2130}{241.7730}{11.9552}{\normalfont\ttfamily}{DataGridView}
\texto{450.0360}{241.7730}{11.9552}{\normalfont}{,}
\texto{457.9070}{241.7730}{11.9552}{\normalfont\ttfamily}{BindingSource}
\texto{537.8830}{241.7730}{11.9552}{\normalfont}{,}
\texto{70.8660}{227.3280}{11.9552}{\normalfont}{patrones}
\texto{119.8488}{227.3280}{11.9552}{\normalfont}{de}
\texto{136.9113}{227.3280}{11.9552}{\normalfont}{CRUD),}
\texto{184.5337}{227.3280}{11.9552}{\normalfont}{documentación}
\texto{266.6324}{227.3280}{11.9552}{\normalfont}{abundante}
\texto{325.6337}{227.3280}{11.9552}{\normalfont}{y}
\texto{337.1681}{227.3280}{11.9552}{\normalfont}{productividad}
\texto{414.0615}{227.3280}{11.9552}{\normalfont}{para}
\texto{442.1802}{227.3280}{11.9552}{\normalfont}{el}
\texto{455.9908}{227.3280}{11.9552}{\normalfont}{equipo,}
\texto{498.4150}{227.3280}{11.9552}{\normalfont}{aspectos}
\texto{70.8660}{212.8820}{11.9552}{\normalfont}{críticos}
\texto{112.3123}{212.8820}{11.9552}{\normalfont}{dado}
\texto{141.9994}{212.8820}{11.9552}{\normalfont}{el}
\texto{155.4275}{212.8820}{11.9552}{\normalfont}{tiempo}
\texto{195.8433}{212.8820}{11.9552}{\normalfont}{disponible}
\texto{252.5707}{212.8820}{11.9552}{\normalfont}{en}
\texto{269.2506}{212.8820}{11.9552}{\normalfont}{el}
\texto{282.6787}{212.8820}{11.9552}{\normalfont}{componente}
\texto{348.7861}{212.8820}{11.9552}{\normalfont}{académico.}
\texto{409.6907}{212.8820}{11.9552}{\normalfont}{Además,}
\texto{458.3006}{212.8820}{11.9552}{\normalfont}{al}
\texto{472.3670}{212.8820}{11.9552}{\normalfont}{compartir}
\texto{527.4184}{212.8820}{11.9552}{\normalfont}{los}
\texto{70.8660}{198.4360}{11.9552}{\normalfont}{portales}
\texto{116.1069}{198.4360}{11.9552}{\normalfont}{el}
\texto{129.0926}{198.4360}{11.9552}{\normalfont}{modelo}
\texto{170.3787}{198.4360}{11.9552}{\normalfont}{de}
\texto{186.6162}{198.4360}{11.9552}{\normalfont}{datos}
\texto{218.5390}{198.4360}{11.9552}{\normalfont}{y}
\texto{229.2485}{198.4360}{11.9552}{\normalfont}{la}
\texto{242.8965}{198.4360}{11.9552}{\normalfont}{API}
\texto{268.3898}{198.4360}{11.9552}{\normalfont}{REST,}
\texto{307.7105}{198.4360}{11.9552}{\normalfont}{la}
\texto{321.3585}{198.4360}{11.9552}{\normalfont}{elección}
\texto{365.5617}{198.4360}{11.9552}{\normalfont}{de}
\texto{381.7992}{198.4360}{11.9552}{\normalfont}{una}
\texto{405.2027}{198.4360}{11.9552}{\normalfont}{tecnología}
\texto{461.1124}{198.4360}{11.9552}{\normalfont}{de}
\texto{477.3619}{198.4360}{11.9552}{\normalfont}{escritorio}
\texto{528.7836}{198.4360}{11.9552}{\normalfont}{no}
\texto{70.8660}{183.9900}{11.9552}{\normalfont}{condiciona}
\texto{129.1500}{183.9900}{11.9552}{\normalfont}{la}
\texto{142.5709}{183.9900}{11.9552}{\normalfont}{futura}
\texto{178.4293}{183.9900}{11.9552}{\normalfont}{integración}
\texto{238.9920}{183.9900}{11.9552}{\normalfont}{con}
\texto{260.8676}{183.9900}{11.9552}{\normalfont}{la}
\texto{274.2885}{183.9900}{11.9552}{\normalfont}{app}
\texto{297.4529}{183.9900}{11.9552}{\normalfont}{móvil:}
\texto{332.9885}{183.9900}{11.9552}{\normalfont}{todos}
\texto{365.0070}{183.9900}{11.9552}{\normalfont}{los}
\texto{383.0450}{183.9900}{11.9552}{\normalfont}{componentes}
\texto{453.1000}{183.9900}{11.9552}{\normalfont}{se}
\texto{467.2359}{183.9900}{11.9552}{\normalfont}{comunican}
\texto{525.8594}{183.9900}{11.9552}{\normalfont}{vía}
\texto{70.8660}{169.5440}{11.9552}{\normalfont}{HTTP,}
\texto{110.6840}{169.5440}{11.9552}{\normalfont}{manteniendo}
\texto{179.9452}{169.5440}{11.9552}{\normalfont}{independencia}
\texto{256.3677}{169.5440}{11.9552}{\normalfont}{tecnológica}
\texto{316.8466}{169.5440}{11.9552}{\normalfont}{entre}
\texto{346.4477}{169.5440}{11.9552}{\normalfont}{cliente}
\texto{383.1908}{169.5440}{11.9552}{\normalfont}{y}
\texto{393.2667}{169.5440}{11.9552}{\normalfont}{servidor.}
\texto{70.8660}{133.0050}{14.3462}{\normalfont\bfseries}{2.5.3.}
\texto{124.6642}{133.0050}{14.3462}{\normalfont\bfseries}{Compatibilidad}
\texto{238.9862}{133.0050}{14.3462}{\normalfont\bfseries}{entre}
\texto{280.4697}{133.0050}{14.3462}{\normalfont\bfseries}{tecnologías}
\texto{88.7990}{110.2190}{11.9552}{\normalfont}{Las}
\texto{111.0189}{110.2190}{11.9552}{\normalfont}{tecnologías}
\texto{171.4501}{110.2190}{11.9552}{\normalfont}{seleccionadas}
\texto{243.0020}{110.2190}{11.9552}{\normalfont}{son}
\texto{264.4113}{110.2190}{11.9552}{\normalfont}{compatibles}
\texto{329.3950}{110.2190}{11.9552}{\normalfont}{entre}
\texto{359.5221}{110.2190}{11.9552}{\normalfont}{sí:}
\texto{375.0782}{110.2190}{11.9552}{\normalfont}{la}
\texto{388.6187}{110.2190}{11.9552}{\normalfont}{API}
\texto{414.0163}{110.2190}{11.9552}{\normalfont}{REST}
\texto{449.9895}{110.2190}{11.9552}{\normalfont}{actúa}
\texto{482.3905}{110.2190}{11.9552}{\normalfont}{como}
\texto{513.5027}{110.2190}{11.9552}{\normalfont}{inter-}
\texto{70.8660}{95.7730}{11.9552}{\normalfont}{mediario}
\texto{118.9403}{95.7730}{11.9552}{\normalfont}{neutral}
\texto{159.2101}{95.7730}{11.9552}{\normalfont}{que}
\texto{180.9447}{95.7730}{11.9552}{\normalfont}{permite}
\texto{224.1508}{95.7730}{11.9552}{\normalfont}{que}
\texto{245.8853}{95.7730}{11.9552}{\normalfont}{el}
\texto{258.1896}{95.7730}{11.9552}{\normalfont}{backoffice}
\texto{311.1464}{95.7730}{11.9552}{\normalfont}{de}
\texto{326.7025}{95.7730}{11.9552}{\normalfont}{escritorio}
\texto{377.4428}{95.7730}{11.9552}{\normalfont}{(.NET}
\texto{414.3030}{95.7730}{11.9552}{\normalfont}{WinForms)}
\texto{475.9249}{95.7730}{11.9552}{\normalfont}{y}
\texto{485.9649}{95.7730}{11.9552}{\normalfont}{el}
\texto{498.2692}{95.7730}{11.9552}{\normalfont}{panel}
\texto{529.4340}{95.7730}{11.9552}{\normalfont}{de}
\texto{70.8660}{81.3270}{11.9552}{\normalfont}{administración}
\texto{150.2342}{81.3270}{11.9552}{\normalfont}{de}
\texto{165.8022}{81.3270}{11.9552}{\normalfont}{escritorio}
\texto{216.5545}{81.3270}{11.9552}{\normalfont}{(WinForms)}
\texto{282.7408}{81.3270}{11.9552}{\normalfont}{consuman}
\texto{337.3952}{81.3270}{11.9552}{\normalfont}{los}
\texto{354.9789}{81.3270}{11.9552}{\normalfont}{mismos}
\texto{396.6906}{81.3270}{11.9552}{\normalfont}{servicios}
\texto{443.2920}{81.3270}{11.9552}{\normalfont}{y}
\texto{453.3320}{81.3270}{11.9552}{\normalfont}{la}
\texto{466.2986}{81.3270}{11.9552}{\normalfont}{misma}
\texto{503.3932}{81.3270}{11.9552}{\normalfont}{base}
\texto{529.4316}{81.3270}{11.9552}{\normalfont}{de}
\texto{70.8660}{66.8810}{11.9552}{\normalfont}{datos,}
\texto{106.0645}{66.8810}{11.9552}{\normalfont}{y}
\texto{116.8098}{66.8810}{11.9552}{\normalfont}{que,}
\texto{142.5135}{66.8810}{11.9552}{\normalfont}{en}
\texto{158.7869}{66.8810}{11.9552}{\normalfont}{la}
\texto{172.4589}{66.8810}{11.9552}{\normalfont}{visión}
\texto{206.6819}{66.8810}{11.9552}{\normalfont}{completa,}
\texto{260.6763}{66.8810}{11.9552}{\normalfont}{la}
\texto{274.3483}{66.8810}{11.9552}{\normalfont}{app}
\texto{297.7757}{66.8810}{11.9552}{\normalfont}{móvil}
\texto{330.3106}{66.8810}{11.9552}{\normalfont}{(React}
\texto{368.8565}{66.8810}{11.9552}{\normalfont}{Native)}
\texto{411.4673}{66.8810}{11.9552}{\normalfont}{se}
\texto{425.8541}{66.8810}{11.9552}{\normalfont}{integre}
\texto{465.2178}{66.8810}{11.9552}{\normalfont}{sin}
\texto{484.1573}{66.8810}{11.9552}{\normalfont}{cambios}
\texto{529.4268}{66.8810}{11.9552}{\normalfont}{de}
\end{pagina}
% Página original 31
\begin{pagina}{31}
\texto{70.8660}{746.3890}{11.9552}{\normalfont}{2.5.}
\texto{96.8710}{746.3890}{11.9552}{\normalfont}{ANÁLISIS}
\texto{155.8723}{746.3890}{11.9552}{\normalfont}{DE}
\texto{176.6743}{746.3890}{11.9552}{\normalfont}{LAS}
\texto{203.1646}{746.3890}{11.9552}{\normalfont}{TECNOLOGÍAS}
\texto{294.9339}{746.3890}{11.9552}{\normalfont}{Y}
\texto{307.6065}{746.3890}{11.9552}{\normalfont}{SU}
\texto{326.7946}{746.3890}{11.9552}{\normalfont}{PERTINENCIA}
\texto{413.9001}{746.3890}{11.9552}{\normalfont}{PARA}
\texto{450.9589}{746.3890}{11.9552}{\normalfont}{EL}
\texto{470.1470}{746.3890}{11.9552}{\normalfont}{PROYECTO}
\texto{529.4400}{746.3890}{11.9552}{\normalfont}{21}
\texto{70.8660}{705.1940}{11.9552}{\normalfont}{arquitectura.}
\texto{141.1578}{705.1940}{11.9552}{\normalfont}{La}
\texto{158.6052}{705.1940}{11.9552}{\normalfont}{autenticación}
\texto{230.8505}{705.1940}{11.9552}{\normalfont}{JWT}
\texto{261.6136}{705.1940}{11.9552}{\normalfont}{y}
\texto{272.0840}{705.1940}{11.9552}{\normalfont}{el}
\texto{284.8187}{705.1940}{11.9552}{\normalfont}{RBAC}
\texto{322.8984}{705.1940}{11.9552}{\normalfont}{son}
\texto{344.1523}{705.1940}{11.9552}{\normalfont}{transversales}
\texto{413.6646}{705.1940}{11.9552}{\normalfont}{y}
\texto{424.1230}{705.1940}{11.9552}{\normalfont}{se}
\texto{438.2230}{705.1940}{11.9552}{\normalfont}{aplican}
\texto{478.9233}{705.1940}{11.9552}{\normalfont}{por}
\texto{500.4355}{705.1940}{11.9552}{\normalfont}{igual}
\texto{529.4292}{705.1940}{11.9552}{\normalfont}{en}
\texto{70.8660}{690.7480}{11.9552}{\normalfont}{todos}
\texto{102.4660}{690.7480}{11.9552}{\normalfont}{los}
\texto{120.0856}{690.7480}{11.9552}{\normalfont}{clientes,}
\texto{164.7095}{690.7480}{11.9552}{\normalfont}{garantizando}
\texto{234.6212}{690.7480}{11.9552}{\normalfont}{coherencia}
\texto{291.8483}{690.7480}{11.9552}{\normalfont}{y}
\texto{301.9361}{690.7480}{11.9552}{\normalfont}{seguridad.}
\end{pagina}
% Página original 32
\begin{pagina}{32}
\texto{70.8660}{746.3890}{11.9552}{\normalfont}{22}
\texto{269.9360}{746.3890}{11.9552}{\normalfont}{CAPÍTULO}
\texto{336.8995}{746.3890}{11.9552}{\normalfont}{2.}
\texto{353.8113}{746.3890}{11.9552}{\normalfont}{INVESTIGACIÓN}
\texto{454.8220}{746.3890}{11.9552}{\normalfont}{DOCUMENTAL}
\texto{70.8660}{705.1940}{17.2154}{\normalfont\bfseries}{2.6.}
\texto{120.3603}{705.1940}{17.2154}{\normalfont\bfseries}{Conclusiones}
\texto{235.9341}{705.1940}{17.2154}{\normalfont\bfseries}{de}
\texto{261.9862}{705.1940}{17.2154}{\normalfont\bfseries}{la}
\texto{283.2369}{705.1940}{17.2154}{\normalfont\bfseries}{investigación}
\texto{88.7990}{678.9090}{11.9552}{\normalfont}{La}
\texto{105.4693}{678.9090}{11.9552}{\normalfont}{investigación}
\texto{174.4030}{678.9090}{11.9552}{\normalfont}{documental}
\texto{237.1009}{678.9090}{11.9552}{\normalfont}{desarrollada}
\texto{302.4647}{678.9090}{11.9552}{\normalfont}{aporta}
\texto{339.4709}{678.9090}{11.9552}{\normalfont}{los}
\texto{356.6959}{678.9090}{11.9552}{\normalfont}{siguientes}
\texto{409.4446}{678.9090}{11.9552}{\normalfont}{hallazgos}
\texto{459.1998}{678.9090}{11.9552}{\normalfont}{que}
\texto{480.5876}{678.9090}{11.9552}{\normalfont}{sustentan}
\texto{532.6860}{678.9090}{11.9552}{\normalfont}{el}
\texto{70.8660}{664.4630}{11.9552}{\normalfont}{proyecto}
\texto{118.0150}{664.4630}{11.9552}{\normalfont\bfseries}{Explora-Nicaragua}
\texto{228.1840}{664.4630}{11.9552}{\normalfont}{:}
\texto{100.7540}{637.0660}{11.9552}{\normalfont\bfseries}{T}
\texto{108.9959}{637.0660}{11.9552}{\normalfont\bfseries}{endencias}
\texto{170.4074}{637.0660}{11.9552}{\normalfont\bfseries}{relevantes:}
\texto{236.8230}{637.0660}{11.9552}{\normalfont}{la}
\texto{250.4352}{637.0660}{11.9552}{\normalfont}{digitalización}
\texto{322.5800}{637.0660}{11.9552}{\normalfont}{turística,}
\texto{372.6771}{637.0660}{11.9552}{\normalfont}{las}
\texto{390.9064}{637.0660}{11.9552}{\normalfont}{aplicaciones}
\texto{455.9618}{637.0660}{11.9552}{\normalfont}{de}
\texto{472.1755}{637.0660}{11.9552}{\normalfont}{escritorio}
\texto{523.5733}{637.0660}{11.9552}{\normalfont}{con}
\texto{100.7540}{622.6200}{11.9552}{\normalfont}{.NET,}
\texto{136.0505}{622.6200}{11.9552}{\normalfont}{las}
\texto{153.3832}{622.6200}{11.9552}{\normalfont}{APIs}
\texto{182.5730}{622.6200}{11.9552}{\normalfont}{REST,}
\texto{220.9731}{622.6200}{11.9552}{\normalfont}{el}
\texto{233.0383}{622.6200}{11.9552}{\normalfont}{control}
\texto{272.0959}{622.6200}{11.9552}{\normalfont}{de}
\texto{287.4010}{622.6200}{11.9552}{\normalfont}{acceso}
\texto{322.9438}{622.6200}{11.9552}{\normalfont}{basado}
\texto{361.7384}{622.6200}{11.9552}{\normalfont}{en}
\texto{377.0554}{622.6200}{11.9552}{\normalfont}{roles}
\texto{404.1435}{622.6200}{11.9552}{\normalfont}{y}
\texto{413.9324}{622.6200}{11.9552}{\normalfont}{las}
\texto{431.2651}{622.6200}{11.9552}{\normalfont}{pasarelas}
\texto{481.1685}{622.6200}{11.9552}{\normalfont}{de}
\texto{496.4855}{622.6200}{11.9552}{\normalfont}{pago}
\texto{524.1594}{622.6200}{11.9552}{\normalfont}{son}
\texto{100.7540}{608.1740}{11.9552}{\normalfont}{tendencias}
\texto{158.0457}{608.1740}{11.9552}{\normalfont}{maduras}
\texto{205.5820}{608.1740}{11.9552}{\normalfont}{y}
\texto{215.6698}{608.1740}{11.9552}{\normalfont}{aplicables}
\texto{269.0593}{608.1740}{11.9552}{\normalfont}{a}
\texto{278.8100}{608.1740}{11.9552}{\normalfont}{la}
\texto{291.8244}{608.1740}{11.9552}{\normalfont}{solución}
\texto{336.7592}{608.1740}{11.9552}{\normalfont}{propuesta.}
\texto{100.7540}{583.7660}{11.9552}{\normalfont\bfseries}{Soluciones}
\texto{167.8956}{583.7660}{11.9552}{\normalfont\bfseries}{existentes:}
\texto{234.6290}{583.7660}{11.9552}{\normalfont}{existen}
\texto{275.6975}{583.7660}{11.9552}{\normalfont}{plataformas}
\texto{341.4918}{583.7660}{11.9552}{\normalfont}{internacionales}
\texto{422.8851}{583.7660}{11.9552}{\normalfont}{consolidadas,}
\texto{495.5728}{583.7660}{11.9552}{\normalfont}{pero}
\texto{523.5790}{583.7660}{11.9552}{\normalfont}{con}
\texto{100.7540}{569.3200}{11.9552}{\normalfont}{comisiones}
\texto{158.8324}{569.3200}{11.9552}{\normalfont}{elevadas}
\texto{204.8073}{569.3200}{11.9552}{\normalfont}{y}
\texto{214.9549}{569.3200}{11.9552}{\normalfont}{sin}
\texto{233.2965}{569.3200}{11.9552}{\normalfont}{un}
\texto{250.2729}{569.3200}{11.9552}{\normalfont}{backoffice}
\texto{303.3492}{569.3200}{11.9552}{\normalfont}{de}
\texto{319.0249}{569.3200}{11.9552}{\normalfont}{escritorio}
\texto{369.8847}{569.3200}{11.9552}{\normalfont}{adaptado}
\texto{421.3303}{569.3200}{11.9552}{\normalfont}{a}
\texto{431.1527}{569.3200}{11.9552}{\normalfont}{la}
\texto{444.2269}{569.3200}{11.9552}{\normalfont}{gestión}
\texto{484.0305}{569.3200}{11.9552}{\normalfont}{local}
\texto{511.7355}{569.3200}{11.9552}{\normalfont}{de}
\texto{527.4112}{569.3200}{11.9552}{\normalfont}{los}
\texto{100.7540}{554.8740}{11.9552}{\normalfont}{operadores,}
\texto{162.9210}{554.8740}{11.9552}{\normalfont}{lo}
\texto{175.9235}{554.8740}{11.9552}{\normalfont}{que}
\texto{197.7178}{554.8740}{11.9552}{\normalfont}{constituye}
\texto{253.7112}{554.8740}{11.9552}{\normalfont}{una}
\texto{276.4691}{554.8740}{11.9552}{\normalfont}{oportunidad}
\texto{343.1361}{554.8740}{11.9552}{\normalfont}{para}
\texto{369.7962}{554.8740}{11.9552}{\normalfont}{la}
\texto{382.8106}{554.8740}{11.9552}{\normalfont}{propuesta.}
\texto{100.7540}{530.4660}{11.9552}{\normalfont\bfseries}{T}
\texto{108.9959}{530.4660}{11.9552}{\normalfont\bfseries}{ecnologías}
\texto{173.9664}{530.4660}{11.9552}{\normalfont\bfseries}{pertinentes:}
\texto{248.6870}{530.4660}{11.9552}{\normalfont}{la}
\texto{262.9328}{530.4660}{11.9552}{\normalfont}{combinación}
\texto{331.4863}{530.4660}{11.9552}{\normalfont}{de}
\texto{348.3455}{530.4660}{11.9552}{\normalfont}{.NET}
\texto{381.9325}{530.4660}{11.9552}{\normalfont}{WinForms}
\texto{440.3049}{530.4660}{11.9552}{\normalfont}{(aplicaciones}
\texto{510.5585}{530.4660}{11.9552}{\normalfont}{de}
\texto{527.4177}{530.4660}{11.9552}{\normalfont}{es-}
\texto{100.7540}{516.0200}{11.9552}{\normalfont}{critorio}
\texto{143.1926}{516.0200}{11.9552}{\normalfont}{del}
\texto{163.5308}{516.0200}{11.9552}{\normalfont}{prototipo),}
\texto{224.5142}{516.0200}{11.9552}{\normalfont}{React}
\texto{259.3086}{516.0200}{11.9552}{\normalfont}{Native}
\texto{298.1798}{516.0200}{11.9552}{\normalfont}{(app}
\texto{326.9607}{516.0200}{11.9552}{\normalfont}{móvil,}
\texto{363.5604}{516.0200}{11.9552}{\normalfont}{trabajo}
\texto{406.0085}{516.0200}{11.9552}{\normalfont}{futuro)}
\texto{447.4715}{516.0200}{11.9552}{\normalfont}{y}
\texto{459.0298}{516.0200}{11.9552}{\normalfont}{una}
\texto{483.2702}{516.0200}{11.9552}{\normalfont}{API}
\texto{509.6003}{516.0200}{11.9552}{\normalfont}{REST}
\texto{100.7540}{501.5740}{11.9552}{\normalfont}{común}
\texto{139.3789}{501.5740}{11.9552}{\normalfont}{con}
\texto{162.0794}{501.5740}{11.9552}{\normalfont}{autenticación}
\texto{235.1854}{501.5740}{11.9552}{\normalfont}{JWT}
\texto{266.8213}{501.5740}{11.9552}{\normalfont}{y}
\texto{278.1285}{501.5740}{11.9552}{\normalfont}{RBAC}
\texto{317.0810}{501.5740}{11.9552}{\normalfont}{resulta}
\texto{356.7555}{501.5740}{11.9552}{\normalfont}{pertinente}
\texto{413.9252}{501.5740}{11.9552}{\normalfont}{y}
\texto{425.2444}{501.5740}{11.9552}{\normalfont}{viable}
\texto{460.6270}{501.5740}{11.9552}{\normalfont}{para}
\texto{488.5185}{501.5740}{11.9552}{\normalfont}{atender}
\texto{532.0307}{501.5740}{11.9552}{\normalfont}{la}
\texto{100.7540}{487.1280}{11.9552}{\normalfont}{problemática}
\texto{170.9884}{487.1280}{11.9552}{\normalfont}{identificada.}
\texto{100.7540}{462.7200}{11.9552}{\normalfont\bfseries}{Aporte}
\texto{148.3512}{462.7200}{11.9552}{\normalfont\bfseries}{de}
\texto{167.9733}{462.7200}{11.9552}{\normalfont\bfseries}{la}
\texto{184.2611}{462.7200}{11.9552}{\normalfont\bfseries}{investigación:}
\texto{268.3840}{462.7200}{11.9552}{\normalfont}{la}
\texto{282.7135}{462.7200}{11.9552}{\normalfont}{investigación}
\texto{353.3807}{462.7200}{11.9552}{\normalfont}{confirma}
\texto{402.8417}{462.7200}{11.9552}{\normalfont}{la}
\texto{417.1832}{462.7200}{11.9552}{\normalfont}{relevancia}
\texto{472.8155}{462.7200}{11.9552}{\normalfont}{de}
\texto{489.7584}{462.7200}{11.9552}{\normalfont}{digitalizar}
\texto{100.7540}{448.2740}{11.9552}{\normalfont}{tanto}
\texto{133.1382}{448.2740}{11.9552}{\normalfont}{el}
\texto{146.9967}{448.2740}{11.9552}{\normalfont}{punto}
\texto{181.9824}{448.2740}{11.9552}{\normalfont}{de}
\texto{199.0927}{448.2740}{11.9552}{\normalfont}{venta}
\texto{232.1297}{448.2740}{11.9552}{\normalfont}{hacia}
\texto{264.1983}{448.2740}{11.9552}{\normalfont}{el}
\texto{278.0568}{448.2740}{11.9552}{\normalfont}{turista}
\texto{317.3320}{448.2740}{11.9552}{\normalfont}{como}
\texto{349.4006}{448.2740}{11.9552}{\normalfont}{la}
\texto{363.9095}{448.2740}{11.9552}{\normalfont}{gestión}
\texto{405.1477}{448.2740}{11.9552}{\normalfont}{interna}
\texto{446.6371}{448.2740}{11.9552}{\normalfont}{de}
\texto{463.7473}{448.2740}{11.9552}{\normalfont}{los}
\texto{482.8733}{448.2740}{11.9552}{\normalfont}{operadores,}
\texto{100.7540}{433.8280}{11.9552}{\normalfont}{sustentando}
\texto{166.2972}{433.8280}{11.9552}{\normalfont}{el}
\texto{179.3666}{433.8280}{11.9552}{\normalfont}{desarrollo}
\texto{233.4735}{433.8280}{11.9552}{\normalfont}{de}
\texto{249.7947}{433.8280}{11.9552}{\normalfont}{las}
\texto{268.1316}{433.8280}{11.9552}{\normalfont}{aplicaciones}
\texto{333.2946}{433.8280}{11.9552}{\normalfont}{de}
\texto{349.6158}{433.8280}{11.9552}{\normalfont}{escritorio}
\texto{401.1212}{433.8280}{11.9552}{\normalfont}{del}
\texto{420.6943}{433.8280}{11.9552}{\normalfont}{prototipo}
\texto{473.1202}{433.8280}{11.9552}{\normalfont}{inicial}
\texto{508.3020}{433.8280}{11.9552}{\normalfont}{dentro}
\texto{100.7540}{419.3820}{11.9552}{\normalfont}{de}
\texto{116.3579}{419.3820}{11.9552}{\normalfont}{la}
\texto{129.3604}{419.3820}{11.9552}{\normalfont}{visión}
\texto{162.9258}{419.3820}{11.9552}{\normalfont}{integral}
\texto{205.5222}{419.3820}{11.9552}{\normalfont}{del}
\texto{224.3779}{419.3820}{11.9552}{\normalfont}{ecosistema.}
\texto{88.7990}{391.9850}{11.9552}{\normalfont}{Estos}
\texto{120.4564}{391.9850}{11.9552}{\normalfont}{hallazgos}
\texto{170.7495}{391.9850}{11.9552}{\normalfont}{respaldan}
\texto{223.6560}{391.9850}{11.9552}{\normalfont}{la}
\texto{236.8139}{391.9850}{11.9552}{\normalfont}{viabilidad}
\texto{291.2699}{391.9850}{11.9552}{\normalfont}{técnica}
\texto{331.1046}{391.9850}{11.9552}{\normalfont}{y}
\texto{341.3359}{391.9850}{11.9552}{\normalfont}{la}
\texto{354.5057}{391.9850}{11.9552}{\normalfont}{pertinencia}
\texto{415.4629}{391.9850}{11.9552}{\normalfont}{de}
\texto{431.2280}{391.9850}{11.9552}{\normalfont\bfseries}{Explora-Nicaragua}
\texto{70.8660}{377.5390}{11.9552}{\normalfont}{como}
\texto{102.6357}{377.5390}{11.9552}{\normalfont}{una}
\texto{126.6011}{377.5390}{11.9552}{\normalfont}{solución}
\texto{172.7434}{377.5390}{11.9552}{\normalfont}{integral}
\texto{216.5592}{377.5390}{11.9552}{\normalfont}{para}
\texto{244.4268}{377.5390}{11.9552}{\normalfont}{la}
\texto{258.6368}{377.5390}{11.9552}{\normalfont}{promoción}
\texto{317.4061}{377.5390}{11.9552}{\normalfont}{y}
\texto{328.6895}{377.5390}{11.9552}{\normalfont}{reserva}
\texto{369.3084}{377.5390}{11.9552}{\normalfont}{de}
\texto{386.1198}{377.5390}{11.9552}{\normalfont}{tours}
\texto{417.3157}{377.5390}{11.9552}{\normalfont}{turísticos}
\texto{469.3759}{377.5390}{11.9552}{\normalfont}{en}
\texto{486.1873}{377.5390}{11.9552}{\normalfont}{Nicaragua,}
\texto{70.8660}{363.0930}{11.9552}{\normalfont}{con}
\texto{92.4786}{363.0930}{11.9552}{\normalfont}{énfasis}
\texto{130.1542}{363.0930}{11.9552}{\normalfont}{en}
\texto{145.9136}{363.0930}{11.9552}{\normalfont}{el}
\texto{158.4211}{363.0930}{11.9552}{\normalfont}{desarrollo}
\texto{211.9661}{363.0930}{11.9552}{\normalfont}{de}
\texto{227.7254}{363.0930}{11.9552}{\normalfont}{las}
\texto{245.5004}{363.0930}{11.9552}{\normalfont}{aplicaciones}
\texto{310.1135}{363.0930}{11.9552}{\normalfont}{de}
\texto{325.8728}{363.0930}{11.9552}{\normalfont}{escritorio}
\texto{376.8163}{363.0930}{11.9552}{\normalfont}{del}
\texto{395.8275}{363.0930}{11.9552}{\normalfont}{prototipo}
\texto{447.6796}{363.0930}{11.9552}{\normalfont}{inicial:}
\texto{485.5512}{363.0930}{11.9552}{\normalfont}{el}
\texto{498.0588}{363.0930}{11.9552}{\normalfont}{panel}
\texto{529.4388}{363.0930}{11.9552}{\normalfont}{de}
\texto{70.8660}{348.6470}{11.9552}{\normalfont}{administración}
\texto{150.2700}{348.6470}{11.9552}{\normalfont}{y}
\texto{160.3459}{348.6470}{11.9552}{\normalfont}{el}
\texto{172.7100}{348.6470}{11.9552}{\normalfont}{backoffice}
\texto{225.7145}{348.6470}{11.9552}{\normalfont}{de}
\texto{241.3185}{348.6470}{11.9552}{\normalfont}{operadores.}
\end{pagina}

% Texto editable; coordenadas en puntos PDF.
% Página original 33
\begin{pagina}{33}
\texto{70.8660}{625.4930}{24.7871}{\normalfont\bfseries}{Capítulo}
\texto{185.1197}{625.4930}{24.7871}{\normalfont\bfseries}{3}
\texto{70.8660}{575.6790}{24.7871}{\normalfont\bfseries}{Emprendimiento}
\texto{283.2295}{575.6790}{24.7871}{\normalfont\bfseries}{digital}
\texto{88.7990}{521.3830}{11.9552}{\normalfont}{El}
\texto{103.7932}{521.3830}{11.9552}{\normalfont}{presente}
\texto{149.5984}{521.3830}{11.9552}{\normalfont}{capítulo}
\texto{194.3610}{521.3830}{11.9552}{\normalfont}{describe}
\texto{239.5110}{521.3830}{11.9552}{\normalfont}{el}
\texto{251.7555}{521.3830}{11.9552}{\normalfont}{modelo}
\texto{292.2885}{521.3830}{11.9552}{\normalfont}{de}
\texto{307.7728}{521.3830}{11.9552}{\normalfont}{negocio}
\texto{349.6065}{521.3830}{11.9552}{\normalfont}{del}
\texto{368.3546}{521.3830}{11.9552}{\normalfont}{emprendimiento}
\texto{454.4177}{521.3830}{11.9552}{\normalfont}{digital}
\texto{490.7610}{521.3830}{11.9552}{\normalfont\bfseries}{Explora-}
\texto{70.8660}{506.9370}{11.9552}{\normalfont\bfseries}{Nicaragua}
\texto{130.3990}{506.9370}{11.9552}{\normalfont}{,}
\texto{139.2219}{506.9370}{11.9552}{\normalfont}{incluyendo}
\texto{199.4259}{506.9370}{11.9552}{\normalfont}{la}
\texto{214.1141}{506.9370}{11.9552}{\normalfont}{idea}
\texto{240.4968}{506.9370}{11.9552}{\normalfont}{de}
\texto{257.7745}{506.9370}{11.9552}{\normalfont}{negocio,}
\texto{304.6532}{506.9370}{11.9552}{\normalfont}{la}
\texto{319.3414}{506.9370}{11.9552}{\normalfont}{propuesta}
\texto{375.0550}{506.9370}{11.9552}{\normalfont}{de}
\texto{392.3326}{506.9370}{11.9552}{\normalfont}{valor,}
\texto{426.1993}{506.9370}{11.9552}{\normalfont}{el}
\texto{440.2371}{506.9370}{11.9552}{\normalfont}{público}
\texto{482.8789}{506.9370}{11.9552}{\normalfont}{objetivo,}
\texto{532.6867}{506.9370}{11.9552}{\normalfont}{el}
\texto{70.8660}{492.4910}{11.9552}{\normalfont}{modelo}
\texto{112.3553}{492.4910}{11.9552}{\normalfont}{de}
\texto{128.8200}{492.4910}{11.9552}{\normalfont}{funcionamiento,}
\texto{214.8640}{492.4910}{11.9552}{\normalfont}{la}
\texto{228.7272}{492.4910}{11.9552}{\normalfont}{identidad}
\texto{281.2895}{492.4910}{11.9552}{\normalfont}{visual,}
\texto{318.9436}{492.4910}{11.9552}{\normalfont}{el}
\texto{332.1565}{492.4910}{11.9552}{\normalfont}{análisis}
\texto{374.1144}{492.4910}{11.9552}{\normalfont}{estratégico}
\texto{433.5557}{492.4910}{11.9552}{\normalfont}{y}
\texto{444.4923}{492.4910}{11.9552}{\normalfont}{las}
\texto{462.9607}{492.4910}{11.9552}{\normalfont}{consideraciones}
\texto{70.8660}{478.0450}{11.9552}{\normalfont}{éticas}
\texto{103.4439}{478.0450}{11.9552}{\normalfont}{que}
\texto{125.2263}{478.0450}{11.9552}{\normalfont}{enmarcan}
\texto{178.5632}{478.0450}{11.9552}{\normalfont}{el}
\texto{190.9153}{478.0450}{11.9552}{\normalfont}{proyecto.}
\texto{300.1470}{36.9930}{11.9552}{\normalfont}{23}
\end{pagina}
% Página original 34
\begin{pagina}{34}
\texto{70.8660}{746.3890}{11.9552}{\normalfont}{24}
\texto{288.7940}{746.3890}{11.9552}{\normalfont}{CAPÍTULO}
\texto{355.7575}{746.3890}{11.9552}{\normalfont}{3.}
\texto{372.6693}{746.3890}{11.9552}{\normalfont}{EMPRENDIMIENTO}
\texto{490.9875}{746.3890}{11.9552}{\normalfont}{DIGITAL}
\texto{70.8660}{705.1940}{17.2154}{\normalfont\bfseries}{3.1.}
\texto{120.3603}{705.1940}{17.2154}{\normalfont\bfseries}{Idea}
\texto{163.0407}{705.1940}{17.2154}{\normalfont\bfseries}{de}
\texto{189.0928}{705.1940}{17.2154}{\normalfont\bfseries}{negocio}
\texto{88.7990}{678.9090}{11.9552}{\normalfont}{La}
\texto{106.1030}{678.9090}{11.9552}{\normalfont}{idea}
\texto{131.0630}{678.9090}{11.9552}{\normalfont}{de}
\texto{146.9180}{678.9090}{11.9552}{\normalfont}{negocio}
\texto{189.1223}{678.9090}{11.9552}{\normalfont}{consiste}
\texto{233.0720}{678.9090}{11.9552}{\normalfont}{en}
\texto{248.9269}{678.9090}{11.9552}{\normalfont}{el}
\texto{261.5301}{678.9090}{11.9552}{\normalfont}{desarrollo}
\texto{315.1707}{678.9090}{11.9552}{\normalfont}{de}
\texto{331.0257}{678.9090}{11.9552}{\normalfont}{un}
\texto{348.2340}{678.9090}{11.9552}{\normalfont\bfseries}{marketplace}
\texto{424.1878}{678.9090}{11.9552}{\normalfont\bfseries}{digital}
\texto{466.1314}{678.9090}{11.9552}{\normalfont\bfseries}{nicaragüense}
\texto{70.8660}{664.4630}{11.9552}{\normalfont}{que}
\texto{93.7124}{664.4630}{11.9552}{\normalfont}{conecte}
\texto{136.3829}{664.4630}{11.9552}{\normalfont}{directamente}
\texto{207.3585}{664.4630}{11.9552}{\normalfont}{a}
\texto{218.1732}{664.4630}{11.9552}{\normalfont}{turistas}
\texto{261.6351}{664.4630}{11.9552}{\normalfont}{(locales}
\texto{304.7049}{664.4630}{11.9552}{\normalfont}{y}
\texto{315.8448}{664.4630}{11.9552}{\normalfont}{extranjeros)}
\texto{382.0049}{664.4630}{11.9552}{\normalfont}{con}
\texto{404.5261}{664.4630}{11.9552}{\normalfont}{operadores}
\texto{464.5173}{664.4630}{11.9552}{\normalfont}{turísticos}
\texto{516.4339}{664.4630}{11.9552}{\normalfont}{nica-}
\texto{70.8660}{650.0170}{11.9552}{\normalfont}{ragüenses}
\texto{125.5108}{650.0170}{11.9552}{\normalfont}{(agencias}
\texto{178.1400}{650.0170}{11.9552}{\normalfont}{de}
\texto{195.5970}{650.0170}{11.9552}{\normalfont}{viajes,}
\texto{233.5906}{650.0170}{11.9552}{\normalfont}{guías}
\texto{265.4201}{650.0170}{11.9552}{\normalfont}{independientes}
\texto{346.6653}{650.0170}{11.9552}{\normalfont}{y}
\texto{358.5942}{650.0170}{11.9552}{\normalfont}{cooperativas}
\texto{427.8076}{650.0170}{11.9552}{\normalfont}{comunitarias)}
\texto{503.5271}{650.0170}{11.9552}{\normalfont}{para}
\texto{532.0402}{650.0170}{11.9552}{\normalfont}{la}
\texto{70.8660}{635.5710}{11.9552}{\normalfont}{promoción,}
\texto{131.6678}{635.5710}{11.9552}{\normalfont}{comparación}
\texto{199.9511}{635.5710}{11.9552}{\normalfont}{y}
\texto{210.0389}{635.5710}{11.9552}{\normalfont}{reserva}
\texto{249.4504}{635.5710}{11.9552}{\normalfont}{de}
\texto{265.0543}{635.5710}{11.9552}{\normalfont}{tours}
\texto{295.0427}{635.5710}{11.9552}{\normalfont}{turísticos}
\texto{345.8954}{635.5710}{11.9552}{\normalfont}{en}
\texto{361.4993}{635.5710}{11.9552}{\normalfont}{todo}
\texto{388.4941}{635.5710}{11.9552}{\normalfont}{el}
\texto{400.8463}{635.5710}{11.9552}{\normalfont}{país.}
\texto{88.7990}{621.1260}{11.9552}{\normalfont}{A}
\texto{101.1607}{621.1260}{11.9552}{\normalfont}{diferencia}
\texto{153.8496}{621.1260}{11.9552}{\normalfont}{de}
\texto{169.1547}{621.1260}{11.9552}{\normalfont}{las}
\texto{186.4634}{621.1260}{11.9552}{\normalfont}{plataformas}
\texto{250.2851}{621.1260}{11.9552}{\normalfont}{internacionales}
\texto{329.7058}{621.1260}{11.9552}{\normalfont}{existentes}
\texto{382.8515}{621.1260}{11.9552}{\normalfont}{(Civitatis,}
\texto{438.2184}{621.1260}{11.9552}{\normalfont}{GetYourGuide,}
\texto{519.3607}{621.1260}{11.9552}{\normalfont}{Via-}
\texto{70.8660}{606.6800}{11.9552}{\normalfont}{tor),}
\texto{96.3425}{606.6800}{11.9552}{\normalfont}{que}
\texto{116.9413}{606.6800}{11.9552}{\normalfont}{cobran}
\texto{154.1244}{606.6800}{11.9552}{\normalfont}{comisiones}
\texto{210.9475}{606.6800}{11.9552}{\normalfont}{del}
\texto{228.6196}{606.6800}{11.9552}{\normalfont}{15}
\texto{242.3227}{606.6800}{11.9552}{\normalfont}{\%}
\texto{254.7920}{606.6800}{11.9552}{\normalfont}{al}
\texto{266.6109}{606.6800}{11.9552}{\normalfont}{25}
\texto{280.3139}{606.6800}{11.9552}{\normalfont}{\%}
\texto{292.7832}{606.6800}{11.9552}{\normalfont}{[6],}
\texto{311.1070}{606.6800}{11.9552}{\normalfont\bfseries}{Explora-Nicaragua}
\texto{423.9910}{606.6800}{11.9552}{\normalfont}{opera}
\texto{454.9932}{606.6800}{11.9552}{\normalfont}{con}
\texto{475.2669}{606.6800}{11.9552}{\normalfont}{una}
\texto{496.8412}{606.6800}{11.9552}{\normalfont}{comisión}
\texto{70.8660}{592.2340}{11.9552}{\normalfont}{fija}
\texto{91.1755}{592.2340}{11.9552}{\normalfont}{reducida}
\texto{139.1254}{592.2340}{11.9552}{\normalfont}{del}
\texto{158.4594}{592.2340}{11.9552}{\normalfont}{8}
\texto{166.2972}{592.2340}{11.9552}{\normalfont}{\%,}
\texto{183.6800}{592.2340}{11.9552}{\normalfont}{favoreciendo}
\texto{251.1456}{592.2340}{11.9552}{\normalfont}{directamente}
\texto{321.5355}{592.2340}{11.9552}{\normalfont}{a}
\texto{331.7643}{592.2340}{11.9552}{\normalfont}{los}
\texto{349.8621}{592.2340}{11.9552}{\normalfont}{proveedores}
\texto{413.8153}{592.2340}{11.9552}{\normalfont}{locales}
\texto{451.7467}{592.2340}{11.9552}{\normalfont}{y}
\texto{462.3008}{592.2340}{11.9552}{\normalfont}{fortaleciendo}
\texto{532.0378}{592.2340}{11.9552}{\normalfont}{la}
\texto{70.8660}{577.7880}{11.9552}{\normalfont}{economía}
\texto{122.2399}{577.7880}{11.9552}{\normalfont}{turística}
\texto{168.4754}{577.7880}{11.9552}{\normalfont}{del}
\texto{187.3431}{577.7880}{11.9552}{\normalfont}{país.}
\texto{88.7990}{563.3420}{11.9552}{\normalfont}{El}
\texto{103.7932}{563.3420}{11.9552}{\normalfont}{ecosistema}
\texto{161.6803}{563.3420}{11.9552}{\normalfont}{se}
\texto{175.2781}{563.3420}{11.9552}{\normalfont}{compone}
\texto{224.2658}{563.3420}{11.9552}{\normalfont}{de}
\texto{239.7501}{563.3420}{11.9552}{\normalfont}{tres}
\texto{262.4531}{563.3420}{11.9552}{\normalfont}{aplicaciones}
\texto{326.7792}{563.3420}{11.9552}{\normalfont}{coordinadas}
\texto{391.4401}{563.3420}{11.9552}{\normalfont}{sobre}
\texto{421.9474}{563.3420}{11.9552}{\normalfont}{una}
\texto{444.5858}{563.3420}{11.9552}{\normalfont}{misma}
\texto{481.6086}{563.3420}{11.9552}{\normalfont}{API}
\texto{506.3487}{563.3420}{11.9552}{\normalfont}{REST,}
\texto{70.8660}{548.8960}{11.9552}{\normalfont}{de}
\texto{87.2948}{548.8960}{11.9552}{\normalfont}{las}
\texto{105.7513}{548.8960}{11.9552}{\normalfont}{cuales}
\texto{141.1210}{548.8960}{11.9552}{\normalfont\bfseries}{las}
\texto{162.1287}{548.8960}{11.9552}{\normalfont\bfseries}{dos}
\texto{187.0708}{548.8960}{11.9552}{\normalfont\bfseries}{aplicaciones}
\texto{261.8613}{548.8960}{11.9552}{\normalfont\bfseries}{de}
\texto{280.9096}{548.8960}{11.9552}{\normalfont\bfseries}{escritorio}
\texto{340.8960}{548.8960}{11.9552}{\normalfont\bfseries}{constituyen}
\texto{413.8215}{548.8960}{11.9552}{\normalfont\bfseries}{el}
\texto{429.1337}{548.8960}{11.9552}{\normalfont\bfseries}{prototipo}
\texto{489.7526}{548.8960}{11.9552}{\normalfont\bfseries}{inicial}
\texto{529.4280}{548.8960}{11.9552}{\normalfont}{de}
\texto{70.8660}{534.4510}{11.9552}{\normalfont}{este}
\texto{94.3388}{534.4510}{11.9552}{\normalfont}{proyecto:}
\texto{100.7540}{510.0420}{11.9552}{\normalfont\bfseries}{Aplicación}
\texto{166.2015}{510.0420}{11.9552}{\normalfont\bfseries}{móvil}
\texto{202.2560}{510.0420}{11.9552}{\normalfont\bfseries}{(turista):}
\texto{258.8780}{510.0420}{11.9552}{\normalfont}{permite}
\texto{301.6298}{510.0420}{11.9552}{\normalfont}{explorar}
\texto{346.9974}{510.0420}{11.9552}{\normalfont}{destinos,}
\texto{394.7704}{510.0420}{11.9552}{\normalfont}{leer}
\texto{416.3997}{510.0420}{11.9552}{\normalfont}{reseñas}
\texto{456.3564}{510.0420}{11.9552}{\normalfont}{verificadas,}
\texto{516.4241}{510.0420}{11.9552}{\normalfont}{com-}
\texto{100.7540}{495.5960}{11.9552}{\normalfont}{parar}
\texto{132.1101}{495.5960}{11.9552}{\normalfont}{opciones}
\texto{179.4742}{495.5960}{11.9552}{\normalfont}{y}
\texto{189.7055}{495.5960}{11.9552}{\normalfont}{confirmar}
\texto{242.5355}{495.5960}{11.9552}{\normalfont}{reservas}
\texto{286.7076}{495.5960}{11.9552}{\normalfont}{con}
\texto{308.3202}{495.5960}{11.9552}{\normalfont}{pago}
\texto{336.4364}{495.5960}{11.9552}{\normalfont}{en}
\texto{352.1958}{495.5960}{11.9552}{\normalfont}{línea}
\texto{380.3120}{495.5960}{11.9552}{\normalfont}{y}
\texto{390.5313}{495.5960}{11.9552}{\normalfont}{voucher}
\texto{433.9359}{495.5960}{11.9552}{\normalfont}{digital.}
\texto{473.7500}{495.5960}{11.9552}{\normalfont\itshape}{Se}
\texto{490.0210}{495.5960}{11.9552}{\normalfont\itshape}{contempla}
\texto{100.7540}{481.1500}{11.9552}{\normalfont\itshape}{como}
\texto{131.2398}{481.1500}{11.9552}{\normalfont\itshape}{trabajo}
\texto{169.9447}{481.1500}{11.9552}{\normalfont\itshape}{futuro}
\texto{204.4534}{481.1500}{11.9552}{\normalfont\itshape}{en}
\texto{220.5929}{481.1500}{11.9552}{\normalfont\itshape}{la}
\texto{233.7436}{481.1500}{11.9552}{\normalfont\itshape}{visión}
\texto{267.8160}{481.1500}{11.9552}{\normalfont\itshape}{del}
\texto{286.3465}{481.1500}{11.9552}{\normalfont\itshape}{ecosistema.}
\texto{100.7540}{456.7420}{11.9552}{\normalfont\bfseries}{Backoffice}
\texto{163.5523}{456.7420}{11.9552}{\normalfont\bfseries}{de}
\texto{181.5006}{456.7420}{11.9552}{\normalfont\bfseries}{escritorio}
\texto{240.3871}{456.7420}{11.9552}{\normalfont\bfseries}{(operador):}
\texto{310.7900}{456.7420}{11.9552}{\normalfont}{herramienta}
\texto{376.0176}{456.7420}{11.9552}{\normalfont}{administrativa}
\texto{453.6690}{456.7420}{11.9552}{\normalfont}{para}
\texto{480.1976}{456.7420}{11.9552}{\normalfont}{que}
\texto{501.8484}{456.7420}{11.9552}{\normalfont}{los}
\texto{519.3485}{456.7420}{11.9552}{\normalfont}{ope-}
\texto{100.7540}{442.2960}{11.9552}{\normalfont}{radores}
\texto{142.0018}{442.2960}{11.9552}{\normalfont}{gestionen}
\texto{193.6555}{442.2960}{11.9552}{\normalfont}{su}
\texto{208.8888}{442.2960}{11.9552}{\normalfont}{catálogo}
\texto{255.2750}{442.2960}{11.9552}{\normalfont}{de}
\texto{271.0941}{442.2960}{11.9552}{\normalfont}{tours,}
\texto{304.5376}{442.2960}{11.9552}{\normalfont}{disponibilidad,}
\texto{383.8292}{442.2960}{11.9552}{\normalfont}{reservas}
\texto{428.0730}{442.2960}{11.9552}{\normalfont}{y}
\texto{438.3641}{442.2960}{11.9552}{\normalfont}{reportes}
\texto{483.8369}{442.2960}{11.9552}{\normalfont}{financieros.}
\texto{100.7540}{427.8500}{11.9552}{\normalfont\itshape}{Incluido}
\texto{146.2065}{427.8500}{11.9552}{\normalfont\itshape}{en}
\texto{162.3460}{427.8500}{11.9552}{\normalfont\itshape}{el}
\texto{174.8990}{427.8500}{11.9552}{\normalfont\itshape}{prototipo}
\texto{224.0528}{427.8500}{11.9552}{\normalfont\itshape}{inicial.}
\texto{100.7540}{403.4420}{11.9552}{\normalfont\bfseries}{Panel}
\texto{137.0046}{403.4420}{11.9552}{\normalfont\bfseries}{de}
\texto{154.1758}{403.4420}{11.9552}{\normalfont\bfseries}{administración}
\texto{244.3670}{403.4420}{11.9552}{\normalfont\bfseries}{de}
\texto{261.5383}{403.4420}{11.9552}{\normalfont\bfseries}{escritorio}
\texto{319.6597}{403.4420}{11.9552}{\normalfont\bfseries}{(equipo}
\texto{367.4638}{403.4420}{11.9552}{\normalfont\bfseries}{gestor):}
\texto{415.1460}{403.4420}{11.9552}{\normalfont}{aplicación}
\texto{468.9707}{403.4420}{11.9552}{\normalfont}{de}
\texto{483.7736}{403.4420}{11.9552}{\normalfont}{supervisión}
\texto{100.7540}{388.9960}{11.9552}{\normalfont}{para}
\texto{128.0836}{388.9960}{11.9552}{\normalfont}{verificar}
\texto{173.6377}{388.9960}{11.9552}{\normalfont}{operadores,}
\texto{236.4862}{388.9960}{11.9552}{\normalfont}{moderar}
\texto{283.6614}{388.9960}{11.9552}{\normalfont}{la}
\texto{297.3453}{388.9960}{11.9552}{\normalfont}{oferta}
\texto{331.5157}{388.9960}{11.9552}{\normalfont}{y}
\texto{342.2730}{388.9960}{11.9552}{\normalfont}{gestionar}
\texto{393.0921}{388.9960}{11.9552}{\normalfont}{la}
\texto{406.7641}{388.9960}{11.9552}{\normalfont}{actividad}
\texto{458.4942}{388.9960}{11.9552}{\normalfont}{económica}
\texto{515.7525}{388.9960}{11.9552}{\normalfont}{de}
\texto{532.0378}{388.9960}{11.9552}{\normalfont}{la}
\texto{100.7540}{374.5500}{11.9552}{\normalfont}{plataforma.}
\texto{163.5120}{374.5500}{11.9552}{\normalfont\itshape}{Incluido}
\texto{208.9645}{374.5500}{11.9552}{\normalfont\itshape}{en}
\texto{225.1040}{374.5500}{11.9552}{\normalfont\itshape}{el}
\texto{237.6570}{374.5500}{11.9552}{\normalfont\itshape}{prototipo}
\texto{286.8227}{374.5500}{11.9552}{\normalfont\itshape}{inicial.}
\end{pagina}
% Página original 35
\begin{pagina}{35}
\texto{70.8660}{746.3890}{11.9552}{\normalfont}{3.2.}
\texto{96.8710}{746.3890}{11.9552}{\normalfont}{PROBLEMA}
\texto{169.1999}{746.3890}{11.9552}{\normalfont}{O}
\texto{182.2143}{746.3890}{11.9552}{\normalfont}{NECESIDAD}
\texto{256.3294}{746.3890}{11.9552}{\normalfont}{IDENTIFICADA}
\texto{529.4280}{746.3890}{11.9552}{\normalfont}{25}
\texto{70.8660}{705.1940}{17.2154}{\normalfont\bfseries}{3.2.}
\texto{120.3603}{705.1940}{17.2154}{\normalfont\bfseries}{Problema}
\texto{208.1709}{705.1940}{17.2154}{\normalfont\bfseries}{o}
\texto{224.3103}{705.1940}{17.2154}{\normalfont\bfseries}{necesidad}
\texto{311.7611}{705.1940}{17.2154}{\normalfont\bfseries}{identificada}
\texto{88.7990}{678.9090}{11.9552}{\normalfont}{El}
\texto{103.7813}{678.9090}{11.9552}{\normalfont}{diagnóstico}
\texto{164.8436}{678.9090}{11.9552}{\normalfont}{realizado}
\texto{214.1469}{678.9090}{11.9552}{\normalfont}{en}
\texto{229.6193}{678.9090}{11.9552}{\normalfont}{el}
\texto{241.8519}{678.9090}{11.9552}{\normalfont}{Capítulo}
\texto{289.8424}{678.9090}{11.9552}{\normalfont}{I}
\texto{297.8428}{678.9090}{11.9552}{\normalfont}{permitió}
\texto{344.8674}{678.9090}{11.9552}{\normalfont}{identificar}
\texto{399.6892}{678.9090}{11.9552}{\normalfont}{las}
\texto{417.1773}{678.9090}{11.9552}{\normalfont}{siguientes}
\texto{470.1890}{678.9090}{11.9552}{\normalfont}{problemáticas}
\texto{70.8660}{664.4630}{11.9552}{\normalfont}{que}
\texto{92.6484}{664.4630}{11.9552}{\normalfont}{dan}
\texto{115.4063}{664.4630}{11.9552}{\normalfont}{origen}
\texto{150.5331}{664.4630}{11.9552}{\normalfont}{al}
\texto{163.5355}{664.4630}{11.9552}{\normalfont}{emprendimiento:}
\texto{100.7540}{637.0660}{11.9552}{\normalfont\bfseries}{Escasa}
\texto{144.0055}{637.0660}{11.9552}{\normalfont\bfseries}{digitalización}
\texto{226.0995}{637.0660}{11.9552}{\normalfont\bfseries}{de}
\texto{244.4663}{637.0660}{11.9552}{\normalfont\bfseries}{la}
\texto{259.4987}{637.0660}{11.9552}{\normalfont\bfseries}{oferta}
\texto{298.5026}{637.0660}{11.9552}{\normalfont\bfseries}{turística:}
\texto{355.1000}{637.0660}{11.9552}{\normalfont}{la}
\texto{368.3416}{637.0660}{11.9552}{\normalfont}{mayoría}
\texto{413.1305}{637.0660}{11.9552}{\normalfont}{de}
\texto{428.9855}{637.0660}{11.9552}{\normalfont}{los}
\texto{446.8442}{637.0660}{11.9552}{\normalfont}{operadores}
\texto{506.0105}{637.0660}{11.9552}{\normalfont}{nicara-}
\texto{100.7540}{622.6200}{11.9552}{\normalfont}{güenses}
\texto{143.3432}{622.6200}{11.9552}{\normalfont}{gestionan}
\texto{195.6352}{622.6200}{11.9552}{\normalfont}{sus}
\texto{215.4618}{622.6200}{11.9552}{\normalfont}{tours}
\texto{245.6415}{622.6200}{11.9552}{\normalfont}{mediante}
\texto{296.2335}{622.6200}{11.9552}{\normalfont}{canales}
\texto{336.8070}{622.6200}{11.9552}{\normalfont}{informales}
\texto{393.3264}{622.6200}{11.9552}{\normalfont}{(WhatsApp,}
\texto{460.5553}{622.6200}{11.9552}{\normalfont}{llamadas}
\texto{509.5955}{622.6200}{11.9552}{\normalfont}{telefó-}
\texto{100.7540}{608.1740}{11.9552}{\normalfont}{nicas,}
\texto{133.3319}{608.1740}{11.9552}{\normalfont}{hojas}
\texto{163.9564}{608.1740}{11.9552}{\normalfont}{de}
\texto{179.5722}{608.1740}{11.9552}{\normalfont}{cálculo),}
\texto{226.3936}{608.1740}{11.9552}{\normalfont}{sin}
\texto{244.6755}{608.1740}{11.9552}{\normalfont}{un}
\texto{261.5801}{608.1740}{11.9552}{\normalfont}{sistema}
\texto{303.3277}{608.1740}{11.9552}{\normalfont}{centralizado}
\texto{368.6987}{608.1740}{11.9552}{\normalfont}{de}
\texto{384.3027}{608.1740}{11.9552}{\normalfont}{reserva}
\texto{423.7142}{608.1740}{11.9552}{\normalfont}{en}
\texto{439.3181}{608.1740}{11.9552}{\normalfont}{línea.}
\texto{100.7540}{583.7660}{11.9552}{\normalfont\bfseries}{Comisiones}
\texto{173.4046}{583.7660}{11.9552}{\normalfont\bfseries}{elevadas}
\texto{228.2132}{583.7660}{11.9552}{\normalfont\bfseries}{de}
\texto{248.4091}{583.7660}{11.9552}{\normalfont\bfseries}{plataformas}
\texto{323.8918}{583.7660}{11.9552}{\normalfont\bfseries}{extranjeras:}
\texto{398.9280}{583.7660}{11.9552}{\normalfont}{los}
\texto{418.3767}{583.7660}{11.9552}{\normalfont}{operadores}
\texto{479.1330}{583.7660}{11.9552}{\normalfont}{que}
\texto{502.7565}{583.7660}{11.9552}{\normalfont}{utilizan}
\texto{100.7540}{569.3200}{11.9552}{\normalfont}{plataformas}
\texto{165.8429}{569.3200}{11.9552}{\normalfont}{como}
\texto{197.3735}{569.3200}{11.9552}{\normalfont}{Civitatis}
\texto{246.2034}{569.3200}{11.9552}{\normalfont}{o}
\texto{256.9343}{569.3200}{11.9552}{\normalfont}{GetYourGuide}
\texto{336.0921}{569.3200}{11.9552}{\normalfont}{enfrentan}
\texto{389.0871}{569.3200}{11.9552}{\normalfont}{comisiones}
\texto{448.0741}{569.3200}{11.9552}{\normalfont}{del}
\texto{467.8982}{569.3200}{11.9552}{\normalfont}{15}
\texto{481.6012}{569.3200}{11.9552}{\normalfont}{\%--25}
\texto{510.9010}{569.3200}{11.9552}{\normalfont}{\%}
\texto{525.5342}{569.3200}{11.9552}{\normalfont}{[6],}
\texto{100.7540}{554.8740}{11.9552}{\normalfont}{reduciendo}
\texto{159.9322}{554.8740}{11.9552}{\normalfont}{significativamente}
\texto{254.9522}{554.8740}{11.9552}{\normalfont}{sus}
\texto{274.5874}{554.8740}{11.9552}{\normalfont}{ingresos.}
\texto{100.7540}{530.4660}{11.9552}{\normalfont\bfseries}{Gestión}
\texto{149.3387}{530.4660}{11.9552}{\normalfont\bfseries}{manual}
\texto{195.3519}{530.4660}{11.9552}{\normalfont\bfseries}{de}
\texto{212.3797}{530.4660}{11.9552}{\normalfont\bfseries}{disponibilidad:}
\texto{301.5450}{530.4660}{11.9552}{\normalfont}{la}
\texto{313.6150}{530.4660}{11.9552}{\normalfont}{ausencia}
\texto{359.5803}{530.4660}{11.9552}{\normalfont}{de}
\texto{374.2637}{530.4660}{11.9552}{\normalfont}{un}
\texto{390.2358}{530.4660}{11.9552}{\normalfont}{sistema}
\texto{431.0629}{530.4660}{11.9552}{\normalfont}{digital}
\texto{466.5578}{530.4660}{11.9552}{\normalfont}{provoca}
\texto{509.1973}{530.4660}{11.9552}{\normalfont}{dobles}
\texto{100.7540}{516.0200}{11.9552}{\normalfont}{reservas,}
\texto{148.0225}{516.0200}{11.9552}{\normalfont}{sobreventa}
\texto{206.2945}{516.0200}{11.9552}{\normalfont}{de}
\texto{221.9104}{516.0200}{11.9552}{\normalfont}{cupos}
\texto{254.8111}{516.0200}{11.9552}{\normalfont}{y}
\texto{264.8989}{516.0200}{11.9552}{\normalfont}{pérdida}
\texto{307.4905}{516.0200}{11.9552}{\normalfont}{de}
\texto{323.1064}{516.0200}{11.9552}{\normalfont}{ingresos.}
\texto{100.7540}{491.6110}{11.9552}{\normalfont\bfseries}{F}
\texto{108.0969}{491.6110}{11.9552}{\normalfont\bfseries}{alta}
\texto{135.6871}{491.6110}{11.9552}{\normalfont\bfseries}{de}
\texto{154.8429}{491.6110}{11.9552}{\normalfont\bfseries}{confianza}
\texto{214.5747}{491.6110}{11.9552}{\normalfont\bfseries}{del}
\texto{237.4784}{491.6110}{11.9552}{\normalfont\bfseries}{turista:}
\texto{285.0610}{491.6110}{11.9552}{\normalfont}{la}
\texto{298.9960}{491.6110}{11.9552}{\normalfont}{inexistencia}
\texto{363.4106}{491.6110}{11.9552}{\normalfont}{de}
\texto{379.9470}{491.6110}{11.9552}{\normalfont}{un}
\texto{397.7961}{491.6110}{11.9552}{\normalfont}{canal}
\texto{429.2909}{491.6110}{11.9552}{\normalfont}{verificado}
\texto{482.9004}{491.6110}{11.9552}{\normalfont}{genera}
\texto{520.9108}{491.6110}{11.9552}{\normalfont}{des-}
\texto{100.7540}{477.1650}{11.9552}{\normalfont}{confianza}
\texto{152.6659}{477.1650}{11.9552}{\normalfont}{hacia}
\texto{183.7661}{477.1650}{11.9552}{\normalfont}{los}
\texto{201.9237}{477.1650}{11.9552}{\normalfont}{operadores}
\texto{261.3769}{477.1650}{11.9552}{\normalfont}{locales,}
\texto{302.6319}{477.1650}{11.9552}{\normalfont}{impulsando}
\texto{365.6645}{477.1650}{11.9552}{\normalfont}{al}
\texto{379.2050}{477.1650}{11.9552}{\normalfont}{turista}
\texto{417.5238}{477.1650}{11.9552}{\normalfont}{a}
\texto{427.8124}{477.1650}{11.9552}{\normalfont}{recurrir}
\texto{470.6312}{477.1650}{11.9552}{\normalfont}{a}
\texto{480.9198}{477.1650}{11.9552}{\normalfont}{plataformas}
\texto{100.7540}{462.7200}{11.9552}{\normalfont}{extranjeras.}
\texto{100.7540}{438.3110}{11.9552}{\normalfont\bfseries}{Ausencia}
\texto{157.9118}{438.3110}{11.9552}{\normalfont\bfseries}{de}
\texto{176.2666}{438.3110}{11.9552}{\normalfont\bfseries}{supervisión}
\texto{247.3487}{438.3110}{11.9552}{\normalfont\bfseries}{centralizada:}
\texto{325.6920}{438.3110}{11.9552}{\normalfont}{no}
\texto{342.1734}{438.3110}{11.9552}{\normalfont}{existe}
\texto{375.3156}{438.3110}{11.9552}{\normalfont}{una}
\texto{398.3007}{438.3110}{11.9552}{\normalfont}{instancia}
\texto{448.0272}{438.3110}{11.9552}{\normalfont}{que}
\texto{470.0486}{438.3110}{11.9552}{\normalfont}{verifique}
\texto{517.4247}{438.3110}{11.9552}{\normalfont}{a}
\texto{527.4145}{438.3110}{11.9552}{\normalfont}{los}
\texto{100.7540}{423.8650}{11.9552}{\normalfont}{operadores,}
\texto{162.9210}{423.8650}{11.9552}{\normalfont}{modere}
\texto{204.2239}{423.8650}{11.9552}{\normalfont}{la}
\texto{217.2263}{423.8650}{11.9552}{\normalfont}{oferta}
\texto{250.7152}{423.8650}{11.9552}{\normalfont}{y}
\texto{260.8030}{423.8650}{11.9552}{\normalfont}{supervise}
\texto{311.6533}{423.8650}{11.9552}{\normalfont}{la}
\texto{324.6677}{423.8650}{11.9552}{\normalfont}{actividad}
\texto{375.7164}{423.8650}{11.9552}{\normalfont}{económica}
\texto{432.3052}{423.8650}{11.9552}{\normalfont}{de}
\texto{447.9091}{423.8650}{11.9552}{\normalfont}{la}
\texto{460.9116}{423.8650}{11.9552}{\normalfont}{plataforma.}
\end{pagina}
% Página original 36
\begin{pagina}{36}
\texto{70.8660}{746.3890}{11.9552}{\normalfont}{26}
\texto{288.7940}{746.3890}{11.9552}{\normalfont}{CAPÍTULO}
\texto{355.7575}{746.3890}{11.9552}{\normalfont}{3.}
\texto{372.6693}{746.3890}{11.9552}{\normalfont}{EMPRENDIMIENTO}
\texto{490.9875}{746.3890}{11.9552}{\normalfont}{DIGITAL}
\texto{70.8660}{705.1940}{17.2154}{\normalfont\bfseries}{3.3.}
\texto{120.3603}{705.1940}{17.2154}{\normalfont\bfseries}{Propuesta}
\texto{212.5832}{705.1940}{17.2154}{\normalfont\bfseries}{de}
\texto{238.6352}{705.1940}{17.2154}{\normalfont\bfseries}{v}
\texto{247.7904}{705.1940}{17.2154}{\normalfont\bfseries}{alor}
\texto{88.7990}{678.9090}{11.9552}{\normalfont}{La}
\texto{106.8203}{678.9090}{11.9552}{\normalfont}{propuesta}
\texto{161.8166}{678.9090}{11.9552}{\normalfont}{de}
\texto{178.3889}{678.9090}{11.9552}{\normalfont}{valor}
\texto{208.2864}{678.9090}{11.9552}{\normalfont}{de}
\texto{224.8540}{678.9090}{11.9552}{\normalfont\bfseries}{Explora-Nicaragua}
\texto{339.8820}{678.9090}{11.9552}{\normalfont}{se}
\texto{354.5558}{678.9090}{11.9552}{\normalfont}{sintetiza}
\texto{402.0753}{678.9090}{11.9552}{\normalfont}{en}
\texto{418.6476}{678.9090}{11.9552}{\normalfont}{los}
\texto{437.2236}{678.9090}{11.9552}{\normalfont}{siguientes}
\texto{491.3113}{678.9090}{11.9552}{\normalfont}{elementos}
\texto{70.8660}{664.4630}{11.9552}{\normalfont}{diferenciadores:}
\texto{100.7540}{637.0660}{11.9552}{\normalfont\bfseries}{Comisión}
\texto{159.6310}{637.0660}{11.9552}{\normalfont\bfseries}{reducida:}
\texto{217.3800}{637.0660}{11.9552}{\normalfont}{8}
\texto{225.2178}{637.0660}{11.9552}{\normalfont}{\%}
\texto{238.6794}{637.0660}{11.9552}{\normalfont}{por}
\texto{259.6297}{637.0660}{11.9552}{\normalfont}{reserva,}
\texto{302.1017}{637.0660}{11.9552}{\normalfont}{frente}
\texto{335.0766}{637.0660}{11.9552}{\normalfont}{al}
\texto{347.8877}{637.0660}{11.9552}{\normalfont}{15}
\texto{361.5908}{637.0660}{11.9552}{\normalfont}{\%--25}
\texto{390.8906}{637.0660}{11.9552}{\normalfont}{\%}
\texto{404.3641}{637.0660}{11.9552}{\normalfont}{de}
\texto{419.7768}{637.0660}{11.9552}{\normalfont}{las}
\texto{437.2050}{637.0660}{11.9552}{\normalfont}{plataformas}
\texto{501.1343}{637.0660}{11.9552}{\normalfont}{interna-}
\texto{100.7540}{622.6200}{11.9552}{\normalfont}{cionales,}
\texto{147.6399}{622.6200}{11.9552}{\normalfont}{incrementando}
\texto{226.6566}{622.6200}{11.9552}{\normalfont}{el}
\texto{239.0087}{622.6200}{11.9552}{\normalfont}{margen}
\texto{280.6391}{622.6200}{11.9552}{\normalfont}{de}
\texto{296.2430}{622.6200}{11.9552}{\normalfont}{los}
\texto{313.8746}{622.6200}{11.9552}{\normalfont}{operadores}
\texto{372.7898}{622.6200}{11.9552}{\normalfont}{locales.}
\texto{100.7540}{598.2110}{11.9552}{\normalfont\bfseries}{V}
\texto{109.7981}{598.2110}{11.9552}{\normalfont\bfseries}{erificación}
\texto{173.2587}{598.2110}{11.9552}{\normalfont\bfseries}{de}
\texto{191.0875}{598.2110}{11.9552}{\normalfont\bfseries}{operadores:}
\texto{262.3690}{598.2110}{11.9552}{\normalfont}{proceso}
\texto{304.1476}{598.2110}{11.9552}{\normalfont}{documental}
\texto{366.9890}{598.2110}{11.9552}{\normalfont}{de}
\texto{382.3657}{598.2110}{11.9552}{\normalfont}{verificación}
\texto{443.2703}{598.2110}{11.9552}{\normalfont}{gestionado}
\texto{500.9852}{598.2110}{11.9552}{\normalfont}{desde}
\texto{532.6856}{598.2110}{11.9552}{\normalfont}{el}
\texto{100.7540}{583.7660}{11.9552}{\normalfont}{panel}
\texto{131.9666}{583.7660}{11.9552}{\normalfont}{de}
\texto{147.5706}{583.7660}{11.9552}{\normalfont}{administración}
\texto{226.9866}{583.7660}{11.9552}{\normalfont}{de}
\texto{242.5905}{583.7660}{11.9552}{\normalfont}{escritorio,}
\texto{296.6304}{583.7660}{11.9552}{\normalfont}{que}
\texto{318.4247}{583.7660}{11.9552}{\normalfont}{genera}
\texto{355.4906}{583.7660}{11.9552}{\normalfont}{confianza}
\texto{406.8765}{583.7660}{11.9552}{\normalfont}{en}
\texto{422.4804}{583.7660}{11.9552}{\normalfont}{el}
\texto{434.8325}{583.7660}{11.9552}{\normalfont}{turista.}
\texto{100.7540}{559.3570}{11.9552}{\normalfont\bfseries}{Reseñas}
\texto{150.7889}{559.3570}{11.9552}{\normalfont\bfseries}{verificadas:}
\texto{218.6450}{559.3570}{11.9552}{\normalfont}{únicamente}
\texto{279.5567}{559.3570}{11.9552}{\normalfont}{los}
\texto{295.9928}{559.3570}{11.9552}{\normalfont}{turistas}
\texto{337.2071}{559.3570}{11.9552}{\normalfont}{que}
\texto{357.8059}{559.3570}{11.9552}{\normalfont}{presentaron}
\texto{420.0949}{559.3570}{11.9552}{\normalfont}{el}
\texto{431.2634}{559.3570}{11.9552}{\normalfont}{voucher}
\texto{473.3290}{559.3570}{11.9552}{\normalfont}{digital}
\texto{508.5490}{559.3570}{11.9552}{\normalfont}{(reser-}
\texto{100.7540}{544.9110}{11.9552}{\normalfont}{va}
\texto{115.3083}{544.9110}{11.9552}{\normalfont}{completada)}
\texto{181.5855}{544.9110}{11.9552}{\normalfont}{pueden}
\texto{221.1859}{544.9110}{11.9552}{\normalfont}{dejar}
\texto{250.0553}{544.9110}{11.9552}{\normalfont}{una}
\texto{272.1079}{544.9110}{11.9552}{\normalfont}{reseña,}
\texto{310.4721}{544.9110}{11.9552}{\normalfont}{garantizando}
\texto{379.6664}{544.9110}{11.9552}{\normalfont}{la}
\texto{391.9516}{544.9110}{11.9552}{\normalfont}{autenticidad}
\texto{458.5564}{544.9110}{11.9552}{\normalfont}{de}
\texto{473.4430}{544.9110}{11.9552}{\normalfont}{las}
\texto{490.3453}{544.9110}{11.9552}{\normalfont}{opiniones.}
\texto{100.7540}{520.5030}{11.9552}{\normalfont\bfseries}{Herramienta}
\texto{179.9632}{520.5030}{11.9552}{\normalfont\bfseries}{de}
\texto{198.6527}{520.5030}{11.9552}{\normalfont\bfseries}{gestión}
\texto{245.0760}{520.5030}{11.9552}{\normalfont\bfseries}{de}
\texto{263.7655}{520.5030}{11.9552}{\normalfont\bfseries}{escritorio:}
\texto{326.4950}{520.5030}{11.9552}{\normalfont}{el}
\texto{339.3731}{520.5030}{11.9552}{\normalfont}{backoffice}
\texto{392.9037}{520.5030}{11.9552}{\normalfont}{desktop}
\texto{436.7387}{520.5030}{11.9552}{\normalfont}{permite}
\texto{480.5067}{520.5030}{11.9552}{\normalfont}{a}
\texto{490.7953}{520.5030}{11.9552}{\normalfont}{los}
\texto{508.9409}{520.5030}{11.9552}{\normalfont}{opera-}
\texto{100.7540}{506.0570}{11.9552}{\normalfont}{dores}
\texto{130.8428}{506.0570}{11.9552}{\normalfont}{administrar}
\texto{193.4618}{506.0570}{11.9552}{\normalfont}{tours,}
\texto{226.1521}{506.0570}{11.9552}{\normalfont}{disponibilidad}
\texto{301.4507}{506.0570}{11.9552}{\normalfont}{y}
\texto{311.0005}{506.0570}{11.9552}{\normalfont}{reservas}
\texto{354.4912}{506.0570}{11.9552}{\normalfont}{con}
\texto{375.4223}{506.0570}{11.9552}{\normalfont}{una}
\texto{397.6542}{506.0570}{11.9552}{\normalfont}{interfaz}
\texto{439.3874}{506.0570}{11.9552}{\normalfont}{eficiente}
\texto{484.0473}{506.0570}{11.9552}{\normalfont}{y}
\texto{493.5971}{506.0570}{11.9552}{\normalfont}{accesible,}
\texto{100.7540}{491.6110}{11.9552}{\normalfont}{sin}
\texto{119.0239}{491.6110}{11.9552}{\normalfont}{depender}
\texto{169.4199}{491.6110}{11.9552}{\normalfont}{de}
\texto{185.0358}{491.6110}{11.9552}{\normalfont}{un}
\texto{201.9404}{491.6110}{11.9552}{\normalfont}{navegador}
\texto{257.5464}{491.6110}{11.9552}{\normalfont}{web.}
\texto{100.7540}{467.2030}{11.9552}{\normalfont\bfseries}{Supervisión}
\texto{173.5002}{467.2030}{11.9552}{\normalfont\bfseries}{centralizada}
\texto{248.2441}{467.2030}{11.9552}{\normalfont\bfseries}{desde}
\texto{285.0231}{467.2030}{11.9552}{\normalfont\bfseries}{el}
\texto{299.1637}{467.2030}{11.9552}{\normalfont\bfseries}{escritorio:}
\texto{361.1700}{467.2030}{11.9552}{\normalfont}{el}
\texto{373.3308}{467.2030}{11.9552}{\normalfont}{panel}
\texto{404.3522}{467.2030}{11.9552}{\normalfont}{de}
\texto{419.7648}{467.2030}{11.9552}{\normalfont}{administración}
\texto{498.9895}{467.2030}{11.9552}{\normalfont}{de}
\texto{514.4022}{467.2030}{11.9552}{\normalfont}{escri-}
\texto{100.7540}{452.7570}{11.9552}{\normalfont}{torio}
\texto{128.6670}{452.7570}{11.9552}{\normalfont}{brinda}
\texto{165.6970}{452.7570}{11.9552}{\normalfont}{al}
\texto{178.6636}{452.7570}{11.9552}{\normalfont}{equipo}
\texto{216.3417}{452.7570}{11.9552}{\normalfont}{gestor}
\texto{250.8348}{452.7570}{11.9552}{\normalfont}{herramientas}
\texto{320.7751}{452.7570}{11.9552}{\normalfont}{de}
\texto{336.3432}{452.7570}{11.9552}{\normalfont}{verificación,}
\texto{400.6789}{452.7570}{11.9552}{\normalfont}{moderación,}
\texto{466.6477}{452.7570}{11.9552}{\normalfont}{reportes}
\texto{511.8694}{452.7570}{11.9552}{\normalfont}{finan-}
\texto{100.7540}{438.3110}{11.9552}{\normalfont}{cieros}
\texto{133.3319}{438.3110}{11.9552}{\normalfont}{y}
\texto{143.4078}{438.3110}{11.9552}{\normalfont}{auditoría.}
\texto{100.7540}{413.9030}{11.9552}{\normalfont\bfseries}{Ecosistema}
\texto{170.2161}{413.9030}{11.9552}{\normalfont\bfseries}{coordinado:}
\texto{242.5540}{413.9030}{11.9552}{\normalfont}{dos}
\texto{263.4134}{413.9030}{11.9552}{\normalfont}{aplicaciones}
\texto{327.8591}{413.9030}{11.9552}{\normalfont}{de}
\texto{343.4511}{413.9030}{11.9552}{\normalfont}{escritorio}
\texto{394.2392}{413.9030}{11.9552}{\normalfont}{sobre}
\texto{424.8660}{413.9030}{11.9552}{\normalfont}{una}
\texto{447.6120}{413.9030}{11.9552}{\normalfont}{misma}
\texto{484.7424}{413.9030}{11.9552}{\normalfont}{API}
\texto{509.5901}{413.9030}{11.9552}{\normalfont}{REST}
\texto{100.7540}{399.4570}{11.9552}{\normalfont}{en}
\texto{117.1828}{399.4570}{11.9552}{\normalfont}{el}
\texto{130.3718}{399.4570}{11.9552}{\normalfont}{prototipo}
\texto{182.9054}{399.4570}{11.9552}{\normalfont}{inicial,}
\texto{221.4465}{399.4570}{11.9552}{\normalfont}{ampliable}
\texto{275.6084}{399.4570}{11.9552}{\normalfont}{a}
\texto{286.1959}{399.4570}{11.9552}{\normalfont}{la}
\texto{300.0233}{399.4570}{11.9552}{\normalfont}{app}
\texto{323.6181}{399.4570}{11.9552}{\normalfont}{móvil}
\texto{356.3084}{399.4570}{11.9552}{\normalfont}{en}
\texto{372.7491}{399.4570}{11.9552}{\normalfont}{la}
\texto{386.5765}{399.4570}{11.9552}{\normalfont}{visión}
\texto{420.9669}{399.4570}{11.9552}{\normalfont}{completa,}
\texto{475.1287}{399.4570}{11.9552}{\normalfont}{garantizando}
\texto{100.7540}{385.0110}{11.9552}{\normalfont}{coherencia}
\texto{157.9812}{385.0110}{11.9552}{\normalfont}{de}
\texto{173.5851}{385.0110}{11.9552}{\normalfont}{datos}
\texto{204.8742}{385.0110}{11.9552}{\normalfont}{en}
\texto{220.4782}{385.0110}{11.9552}{\normalfont}{tiempo}
\texto{259.8299}{385.0110}{11.9552}{\normalfont}{real.}
\end{pagina}
% Página original 37
\begin{pagina}{37}
\texto{70.8660}{746.3890}{11.9552}{\normalfont}{3.4.}
\texto{96.8710}{746.3890}{11.9552}{\normalfont}{PÚBLICO}
\texto{154.9063}{746.3890}{11.9552}{\normalfont}{OBJETIVO}
\texto{529.4280}{746.3890}{11.9552}{\normalfont}{27}
\texto{70.8660}{705.1940}{17.2154}{\normalfont\bfseries}{3.4.}
\texto{120.3603}{705.1940}{17.2154}{\normalfont\bfseries}{Público}
\texto{190.6163}{705.1940}{17.2154}{\normalfont\bfseries}{objetivo}
\texto{88.7990}{678.9090}{11.9552}{\normalfont}{El}
\texto{103.9128}{678.9090}{11.9552}{\normalfont}{emprendimiento}
\texto{190.0835}{678.9090}{11.9552}{\normalfont}{atiende}
\texto{231.0515}{678.9090}{11.9552}{\normalfont}{a}
\texto{240.8142}{678.9090}{11.9552}{\normalfont}{tres}
\texto{263.6366}{678.9090}{11.9552}{\normalfont}{segmentos}
\texto{319.3694}{678.9090}{11.9552}{\normalfont}{principales:}
\texto{85.6720}{651.5120}{11.9552}{\normalfont}{1.}
\texto{100.7540}{651.5120}{11.9552}{\normalfont\bfseries}{T}
\texto{108.9959}{651.5120}{11.9552}{\normalfont\bfseries}{uristas}
\texto{153.3019}{651.5120}{11.9552}{\normalfont\bfseries}{(consumidores}
\texto{242.7973}{651.5120}{11.9552}{\normalfont\bfseries}{finales):}
\texto{292.9890}{651.5120}{11.9552}{\normalfont}{viajeros}
\texto{336.9531}{651.5120}{11.9552}{\normalfont}{locales}
\texto{375.0758}{651.5120}{11.9552}{\normalfont}{y}
\texto{385.8092}{651.5120}{11.9552}{\normalfont}{extranjeros}
\texto{447.0102}{651.5120}{11.9552}{\normalfont}{que}
\texto{469.4621}{651.5120}{11.9552}{\normalfont}{buscan}
\texto{509.2012}{651.5120}{11.9552}{\normalfont}{descu-}
\texto{100.7540}{637.0660}{11.9552}{\normalfont}{brir,}
\texto{127.0387}{637.0660}{11.9552}{\normalfont}{comparar}
\texto{179.3499}{637.0660}{11.9552}{\normalfont}{y}
\texto{189.7126}{637.0660}{11.9552}{\normalfont}{reservar}
\texto{233.9517}{637.0660}{11.9552}{\normalfont}{experiencias}
\texto{299.9826}{637.0660}{11.9552}{\normalfont}{turísticas}
\texto{351.1222}{637.0660}{11.9552}{\normalfont}{en}
\texto{367.0130}{637.0660}{11.9552}{\normalfont}{Nicaragua}
\texto{422.8964}{637.0660}{11.9552}{\normalfont}{de}
\texto{438.7753}{637.0660}{11.9552}{\normalfont}{forma}
\texto{472.5512}{637.0660}{11.9552}{\normalfont}{segura}
\texto{509.3182}{637.0660}{11.9552}{\normalfont}{y}
\texto{519.6809}{637.0660}{11.9552}{\normalfont}{con-}
\texto{100.7540}{622.6200}{11.9552}{\normalfont}{fiable.}
\texto{134.9793}{622.6200}{11.9552}{\normalfont}{Perfil:}
\texto{169.0421}{622.6200}{11.9552}{\normalfont}{usuarios}
\texto{214.4647}{622.6200}{11.9552}{\normalfont}{de}
\texto{229.8295}{622.6200}{11.9552}{\normalfont}{smartphones}
\texto{298.0148}{622.6200}{11.9552}{\normalfont}{con}
\texto{319.2329}{622.6200}{11.9552}{\normalfont}{acceso}
\texto{354.8355}{622.6200}{11.9552}{\normalfont}{a}
\texto{364.3470}{622.6200}{11.9552}{\normalfont}{internet,}
\texto{411.2568}{622.6200}{11.9552}{\normalfont}{interesados}
\texto{471.3150}{622.6200}{11.9552}{\normalfont}{en}
\texto{486.6798}{622.6200}{11.9552}{\normalfont}{turismo}
\texto{529.4364}{622.6200}{11.9552}{\normalfont}{de}
\texto{100.7540}{608.1740}{11.9552}{\normalfont}{aventura,}
\texto{152.1231}{608.1740}{11.9552}{\normalfont}{cultural}
\texto{195.7046}{608.1740}{11.9552}{\normalfont}{y}
\texto{205.7804}{608.1740}{11.9552}{\normalfont}{ecológico.}
\texto{85.6720}{583.7660}{11.9552}{\normalfont}{2.}
\texto{100.7540}{583.7660}{11.9552}{\normalfont\bfseries}{Operadores}
\texto{173.7010}{583.7660}{11.9552}{\normalfont\bfseries}{turísticos}
\texto{233.6157}{583.7660}{11.9552}{\normalfont\bfseries}{(proveedores):}
\texto{322.2390}{583.7660}{11.9552}{\normalfont}{agencias}
\texto{369.5386}{583.7660}{11.9552}{\normalfont}{de}
\texto{386.2065}{583.7660}{11.9552}{\normalfont}{viajes,}
\texto{423.4350}{583.7660}{11.9552}{\normalfont}{guías}
\texto{454.4755}{583.7660}{11.9552}{\normalfont}{independientes}
\texto{534.9555}{583.7660}{11.9552}{\normalfont}{y}
\texto{100.7540}{569.3200}{11.9552}{\normalfont}{cooperativas}
\texto{168.4611}{569.3200}{11.9552}{\normalfont}{comunitarias}
\texto{238.1216}{569.3200}{11.9552}{\normalfont}{que}
\texto{260.2507}{569.3200}{11.9552}{\normalfont}{ofrecen}
\texto{300.5780}{569.3200}{11.9552}{\normalfont}{tours}
\texto{330.8892}{569.3200}{11.9552}{\normalfont}{en}
\texto{346.8398}{569.3200}{11.9552}{\normalfont}{destinos}
\texto{392.1740}{569.3200}{11.9552}{\normalfont}{como}
\texto{423.0829}{569.3200}{11.9552}{\normalfont}{Ometepe,}
\texto{476.4151}{569.3200}{11.9552}{\normalfont}{Somoto,}
\texto{522.2824}{569.3200}{11.9552}{\normalfont}{San}
\texto{100.7540}{554.8740}{11.9552}{\normalfont}{Juan}
\texto{129.7764}{554.8740}{11.9552}{\normalfont}{del}
\texto{148.8832}{554.8740}{11.9552}{\normalfont}{Sur,}
\texto{173.8433}{554.8740}{11.9552}{\normalfont}{Solentiname}
\texto{240.1038}{554.8740}{11.9552}{\normalfont}{y}
\texto{250.4307}{554.8740}{11.9552}{\normalfont}{el}
\texto{263.0339}{554.8740}{11.9552}{\normalfont}{Río}
\texto{284.9023}{554.8740}{11.9552}{\normalfont}{San}
\texto{307.9113}{554.8740}{11.9552}{\normalfont}{Juan.}
\texto{340.1975}{554.8740}{11.9552}{\normalfont}{Perfil:}
\texto{374.7504}{554.8740}{11.9552}{\normalfont}{pequeños}
\texto{425.7872}{554.8740}{11.9552}{\normalfont}{y}
\texto{436.1141}{554.8740}{11.9552}{\normalfont}{medianos}
\texto{487.8036}{554.8740}{11.9552}{\normalfont}{emprende-}
\texto{100.7540}{540.4280}{11.9552}{\normalfont}{dores}
\texto{131.8471}{540.4280}{11.9552}{\normalfont}{con}
\texto{153.7825}{540.4280}{11.9552}{\normalfont}{presencia}
\texto{205.0488}{540.4280}{11.9552}{\normalfont}{limitada}
\texto{251.6860}{540.4280}{11.9552}{\normalfont}{en}
\texto{267.7681}{540.4280}{11.9552}{\normalfont}{canales}
\texto{308.6286}{540.4280}{11.9552}{\normalfont}{digitales,}
\texto{358.5822}{540.4280}{11.9552}{\normalfont}{que}
\texto{380.8428}{540.4280}{11.9552}{\normalfont}{buscan}
\texto{420.4026}{540.4280}{11.9552}{\normalfont}{aumentar}
\texto{473.2206}{540.4280}{11.9552}{\normalfont}{su}
\texto{488.7170}{540.4280}{11.9552}{\normalfont}{visibilidad}
\texto{100.7540}{525.9820}{11.9552}{\normalfont}{y}
\texto{110.8298}{525.9820}{11.9552}{\normalfont}{eficiencia}
\texto{160.9030}{525.9820}{11.9552}{\normalfont}{operativa.}
\texto{85.6720}{501.5740}{11.9552}{\normalfont}{3.}
\texto{100.7540}{501.5740}{11.9552}{\normalfont\bfseries}{Equipo}
\texto{145.4449}{501.5740}{11.9552}{\normalfont\bfseries}{gestor}
\texto{184.0495}{501.5740}{11.9552}{\normalfont\bfseries}{de}
\texto{200.6469}{501.5740}{11.9552}{\normalfont\bfseries}{la}
\texto{213.9100}{501.5740}{11.9552}{\normalfont\bfseries}{plataforma}
\texto{280.4885}{501.5740}{11.9552}{\normalfont\bfseries}{(administradores):}
\texto{391.1520}{501.5740}{11.9552}{\normalfont}{personal}
\texto{436.4072}{501.5740}{11.9552}{\normalfont}{de}
\texto{450.7230}{501.5740}{11.9552}{\normalfont\bfseries}{Explora-Nicaragua}
\texto{100.7540}{487.1280}{11.9552}{\normalfont}{encargado}
\texto{155.5877}{487.1280}{11.9552}{\normalfont}{de}
\texto{170.7613}{487.1280}{11.9552}{\normalfont}{verificar}
\texto{215.1916}{487.1280}{11.9552}{\normalfont}{operadores,}
\texto{276.9282}{487.1280}{11.9552}{\normalfont}{moderar}
\texto{322.9916}{487.1280}{11.9552}{\normalfont}{la}
\texto{335.5517}{487.1280}{11.9552}{\normalfont}{oferta,}
\texto{371.8621}{487.1280}{11.9552}{\normalfont}{gestionar}
\texto{421.5574}{487.1280}{11.9552}{\normalfont}{usuarios}
\texto{466.7768}{487.1280}{11.9552}{\normalfont}{y}
\texto{476.4103}{487.1280}{11.9552}{\normalfont}{supervisar}
\texto{532.0331}{487.1280}{11.9552}{\normalfont}{la}
\texto{100.7540}{472.6820}{11.9552}{\normalfont}{actividad}
\texto{151.8027}{472.6820}{11.9552}{\normalfont}{económica}
\texto{208.3795}{472.6820}{11.9552}{\normalfont}{y}
\texto{218.4673}{472.6820}{11.9552}{\normalfont}{financiera}
\texto{271.7923}{472.6820}{11.9552}{\normalfont}{de}
\texto{287.3962}{472.6820}{11.9552}{\normalfont}{la}
\texto{300.4106}{472.6820}{11.9552}{\normalfont}{plataforma.}
\end{pagina}
% Página original 38
\begin{pagina}{38}
\texto{70.8660}{746.3890}{11.9552}{\normalfont}{28}
\texto{288.7940}{746.3890}{11.9552}{\normalfont}{CAPÍTULO}
\texto{355.7575}{746.3890}{11.9552}{\normalfont}{3.}
\texto{372.6693}{746.3890}{11.9552}{\normalfont}{EMPRENDIMIENTO}
\texto{490.9875}{746.3890}{11.9552}{\normalfont}{DIGITAL}
\texto{70.8660}{705.1940}{17.2154}{\normalfont\bfseries}{3.5.}
\texto{120.3603}{705.1940}{17.2154}{\normalfont\bfseries}{Modelo}
\texto{190.0637}{705.1940}{17.2154}{\normalfont\bfseries}{de}
\texto{216.1158}{705.1940}{17.2154}{\normalfont\bfseries}{funcionamiento}
\texto{351.6526}{705.1940}{17.2154}{\normalfont\bfseries}{del}
\texto{383.0845}{705.1940}{17.2154}{\normalfont\bfseries}{emprendimiento}
\texto{88.7990}{678.9090}{11.9552}{\normalfont}{El}
\texto{103.0161}{678.9090}{11.9552}{\normalfont}{modelo}
\texto{142.7600}{678.9090}{11.9552}{\normalfont}{de}
\texto{157.4792}{678.9090}{11.9552}{\normalfont}{funcionamiento}
\texto{238.5259}{678.9090}{11.9552}{\normalfont}{de}
\texto{253.2370}{678.9090}{11.9552}{\normalfont\bfseries}{Explora-Nicaragua}
\texto{366.4110}{678.9090}{11.9552}{\normalfont}{se}
\texto{379.2318}{678.9090}{11.9552}{\normalfont}{basa}
\texto{405.0598}{678.9090}{11.9552}{\normalfont}{en}
\texto{419.7671}{678.9090}{11.9552}{\normalfont}{un}
\texto{435.7870}{678.9090}{11.9552}{\normalfont}{modelo}
\texto{475.5309}{678.9090}{11.9552}{\normalfont}{de}
\texto{490.2500}{678.9090}{11.9552}{\normalfont\bfseries}{comisión}
\texto{70.8660}{664.4630}{11.9552}{\normalfont\bfseries}{por}
\texto{94.6808}{664.4630}{11.9552}{\normalfont\bfseries}{transacción}
\texto{164.4580}{664.4630}{11.9552}{\normalfont}{complementado}
\texto{247.3936}{664.4630}{11.9552}{\normalfont}{con}
\texto{268.2291}{664.4630}{11.9552}{\normalfont}{una}
\texto{290.3534}{664.4630}{11.9552}{\normalfont}{suscripción}
\texto{350.0123}{664.4630}{11.9552}{\normalfont}{opcional.}
\texto{399.1362}{664.4630}{11.9552}{\normalfont}{El}
\texto{413.6283}{664.4630}{11.9552}{\normalfont}{flujo}
\texto{439.3415}{664.4630}{11.9552}{\normalfont}{completo}
\texto{488.7811}{664.4630}{11.9552}{\normalfont}{contempla}
\texto{70.8660}{650.0170}{11.9552}{\normalfont}{la}
\texto{84.6695}{650.0170}{11.9552}{\normalfont}{interacción}
\texto{144.9882}{650.0170}{11.9552}{\normalfont}{del}
\texto{164.6569}{650.0170}{11.9552}{\normalfont}{turista}
\texto{203.2507}{650.0170}{11.9552}{\normalfont}{a}
\texto{213.8143}{650.0170}{11.9552}{\normalfont}{través}
\texto{248.8239}{650.0170}{11.9552}{\normalfont}{de}
\texto{265.2408}{650.0170}{11.9552}{\normalfont}{una}
\texto{288.8117}{650.0170}{11.9552}{\normalfont}{aplicación}
\texto{344.2503}{650.0170}{11.9552}{\normalfont}{móvil,}
\texto{380.1805}{650.0170}{11.9552}{\normalfont}{que}
\texto{402.7639}{650.0170}{11.9552}{\normalfont}{en}
\texto{419.1808}{650.0170}{11.9552}{\normalfont}{la}
\texto{432.9962}{650.0170}{11.9552}{\normalfont}{visión}
\texto{467.3626}{650.0170}{11.9552}{\normalfont}{del}
\texto{487.0313}{650.0170}{11.9552}{\normalfont}{ecosistema}
\texto{70.8660}{635.5710}{11.9552}{\normalfont}{consumirá}
\texto{127.1798}{635.5710}{11.9552}{\normalfont}{la}
\texto{140.5050}{635.5710}{11.9552}{\normalfont}{misma}
\texto{177.9583}{635.5710}{11.9552}{\normalfont}{API}
\texto{203.1407}{635.5710}{11.9552}{\normalfont}{REST}
\texto{238.8987}{635.5710}{11.9552}{\normalfont}{que}
\texto{261.0039}{635.5710}{11.9552}{\normalfont}{las}
\texto{278.9463}{635.5710}{11.9552}{\normalfont}{aplicaciones}
\texto{343.7267}{635.5710}{11.9552}{\normalfont}{de}
\texto{359.6534}{635.5710}{11.9552}{\normalfont}{escritorio;}
\texto{414.0161}{635.5710}{11.9552}{\normalfont}{en}
\texto{429.9428}{635.5710}{11.9552}{\normalfont}{el}
\texto{442.6177}{635.5710}{11.9552}{\normalfont}{prototipo}
\texto{494.6372}{635.5710}{11.9552}{\normalfont}{inicial,}
\texto{532.6882}{635.5710}{11.9552}{\normalfont}{el}
\texto{70.8660}{621.1260}{11.9552}{\normalfont}{flujo}
\texto{97.1292}{621.1260}{11.9552}{\normalfont}{se}
\texto{110.7868}{621.1260}{11.9552}{\normalfont}{automatiza}
\texto{171.8444}{621.1260}{11.9552}{\normalfont}{en}
\texto{187.3886}{621.1260}{11.9552}{\normalfont}{las}
\texto{204.9484}{621.1260}{11.9552}{\normalfont}{dos}
\texto{225.7480}{621.1260}{11.9552}{\normalfont}{aplicaciones}
\texto{290.1339}{621.1260}{11.9552}{\normalfont}{de}
\texto{305.6661}{621.1260}{11.9552}{\normalfont}{escritorio}
\texto{356.3944}{621.1260}{11.9552}{\normalfont}{(backoffice}
\texto{413.8798}{621.1260}{11.9552}{\normalfont}{de}
\texto{429.4240}{621.1260}{11.9552}{\normalfont}{operadores}
\texto{488.2794}{621.1260}{11.9552}{\normalfont}{y}
\texto{498.2835}{621.1260}{11.9552}{\normalfont}{panel}
\texto{529.4364}{621.1260}{11.9552}{\normalfont}{de}
\texto{70.8660}{606.6800}{11.9552}{\normalfont}{administración):}
\texto{85.6720}{579.2820}{11.9552}{\normalfont}{1.}
\texto{100.7540}{579.2820}{11.9552}{\normalfont\bfseries}{Registro}
\texto{155.9237}{579.2820}{11.9552}{\normalfont\bfseries}{y}
\texto{168.8209}{579.2820}{11.9552}{\normalfont\bfseries}{verificación:}
\texto{243.6070}{579.2820}{11.9552}{\normalfont}{el}
\texto{257.1068}{579.2820}{11.9552}{\normalfont}{operador}
\texto{307.3497}{579.2820}{11.9552}{\normalfont}{se}
\texto{322.2148}{579.2820}{11.9552}{\normalfont}{registra}
\texto{365.6959}{579.2820}{11.9552}{\normalfont}{en}
\texto{382.4475}{579.2820}{11.9552}{\normalfont}{la}
\texto{396.5977}{579.2820}{11.9552}{\normalfont}{plataforma}
\texto{457.2488}{579.2820}{11.9552}{\normalfont}{y}
\texto{468.4843}{579.2820}{11.9552}{\normalfont}{sube}
\texto{496.6795}{579.2820}{11.9552}{\normalfont}{su}
\texto{512.8453}{579.2820}{11.9552}{\normalfont}{docu-}
\texto{100.7540}{564.8370}{11.9552}{\normalfont}{mentación}
\texto{157.8329}{564.8370}{11.9552}{\normalfont}{oficial.}
\texto{195.7357}{564.8370}{11.9552}{\normalfont}{El}
\texto{211.6863}{564.8370}{11.9552}{\normalfont}{equipo}
\texto{250.2370}{564.8370}{11.9552}{\normalfont}{gestor}
\texto{285.6029}{564.8370}{11.9552}{\normalfont}{verifica}
\texto{326.7479}{564.8370}{11.9552}{\normalfont}{la}
\texto{340.5873}{564.8370}{11.9552}{\normalfont}{documentación}
\texto{422.0643}{564.8370}{11.9552}{\normalfont}{y}
\texto{432.9651}{564.8370}{11.9552}{\normalfont}{aprueba}
\texto{478.6722}{564.8370}{11.9552}{\normalfont}{o}
\texto{489.2597}{564.8370}{11.9552}{\normalfont}{rechaza}
\texto{532.0307}{564.8370}{11.9552}{\normalfont}{la}
\texto{100.7540}{550.3910}{11.9552}{\normalfont}{solicitud}
\texto{147.6399}{550.3910}{11.9552}{\normalfont}{desde}
\texto{179.5675}{550.3910}{11.9552}{\normalfont}{el}
\texto{191.9315}{550.3910}{11.9552}{\normalfont}{panel}
\texto{223.1442}{550.3910}{11.9552}{\normalfont}{de}
\texto{238.7481}{550.3910}{11.9552}{\normalfont}{administración}
\texto{318.1641}{550.3910}{11.9552}{\normalfont}{de}
\texto{333.7680}{550.3910}{11.9552}{\normalfont}{escritorio.}
\texto{85.6720}{525.9820}{11.9552}{\normalfont}{2.}
\texto{100.7540}{525.9820}{11.9552}{\normalfont\bfseries}{Publicación}
\texto{175.0604}{525.9820}{11.9552}{\normalfont\bfseries}{de}
\texto{194.9574}{525.9820}{11.9552}{\normalfont\bfseries}{tours:}
\texto{234.4080}{525.9820}{11.9552}{\normalfont}{una}
\texto{258.7320}{525.9820}{11.9552}{\normalfont}{vez}
\texto{280.4690}{525.9820}{11.9552}{\normalfont}{verificado,}
\texto{337.9759}{525.9820}{11.9552}{\normalfont}{el}
\texto{351.8941}{525.9820}{11.9552}{\normalfont}{operador}
\texto{402.5675}{525.9820}{11.9552}{\normalfont}{publica}
\texto{445.1136}{525.9820}{11.9552}{\normalfont}{sus}
\texto{466.3269}{525.9820}{11.9552}{\normalfont}{tours}
\texto{497.8815}{525.9820}{11.9552}{\normalfont}{con}
\texto{520.9048}{525.9820}{11.9552}{\normalfont}{des-}
\texto{100.7540}{511.5370}{11.9552}{\normalfont}{cripciones,}
\texto{159.6118}{511.5370}{11.9552}{\normalfont}{fotografías,}
\texto{221.3987}{511.5370}{11.9552}{\normalfont}{precios,}
\texto{264.9754}{511.5370}{11.9552}{\normalfont}{disponibilidad}
\texto{342.0434}{511.5370}{11.9552}{\normalfont}{y}
\texto{353.3626}{511.5370}{11.9552}{\normalfont}{políticas}
\texto{401.1762}{511.5370}{11.9552}{\normalfont}{de}
\texto{418.0235}{511.5370}{11.9552}{\normalfont}{cancelación}
\texto{481.0465}{511.5370}{11.9552}{\normalfont}{mediante}
\texto{532.6786}{511.5370}{11.9552}{\normalfont}{el}
\texto{100.7540}{497.0910}{11.9552}{\normalfont}{backoffice}
\texto{153.7466}{497.0910}{11.9552}{\normalfont}{de}
\texto{169.3625}{497.0910}{11.9552}{\normalfont}{escritorio.}
\texto{85.6720}{472.6820}{11.9552}{\normalfont}{3.}
\texto{100.7540}{472.6820}{11.9552}{\normalfont\bfseries}{Descubrimiento}
\texto{196.8822}{472.6820}{11.9552}{\normalfont\bfseries}{y}
\texto{208.0818}{472.6820}{11.9552}{\normalfont\bfseries}{reserva:}
\texto{256.8400}{472.6820}{11.9552}{\normalfont}{el}
\texto{268.8574}{472.6820}{11.9552}{\normalfont}{turista}
\texto{306.3034}{472.6820}{11.9552}{\normalfont}{explora}
\texto{347.2620}{472.6820}{11.9552}{\normalfont}{tours}
\texto{376.9156}{472.6820}{11.9552}{\normalfont}{y}
\texto{386.6567}{472.6820}{11.9552}{\normalfont}{confirma}
\texto{434.4441}{472.6820}{11.9552}{\normalfont}{una}
\texto{456.8672}{472.6820}{11.9552}{\normalfont}{reserva}
\texto{495.9440}{472.6820}{11.9552}{\normalfont}{con}
\texto{517.0784}{472.6820}{11.9552}{\normalfont}{pago}
\texto{100.7540}{458.2360}{11.9552}{\normalfont}{en}
\texto{117.5176}{458.2360}{11.9552}{\normalfont}{línea}
\texto{146.6381}{458.2360}{11.9552}{\normalfont}{mediante}
\texto{198.2104}{458.2360}{11.9552}{\normalfont}{Stripe.}
\texto{237.4211}{458.2360}{11.9552}{\normalfont}{En}
\texto{256.9463}{458.2360}{11.9552}{\normalfont}{el}
\texto{270.4700}{458.2360}{11.9552}{\normalfont}{prototipo}
\texto{323.3263}{458.2360}{11.9552}{\normalfont}{inicial,}
\texto{362.2142}{458.2360}{11.9552}{\normalfont}{las}
\texto{380.9935}{458.2360}{11.9552}{\normalfont}{reservas}
\texto{426.1817}{458.2360}{11.9552}{\normalfont}{se}
\texto{441.0707}{458.2360}{11.9552}{\normalfont}{gestionan}
\texto{494.3192}{458.2360}{11.9552}{\normalfont}{desde}
\texto{527.4184}{458.2360}{11.9552}{\normalfont}{las}
\texto{100.7540}{443.7910}{11.9552}{\normalfont}{aplicaciones}
\texto{165.8931}{443.7910}{11.9552}{\normalfont}{de}
\texto{182.2024}{443.7910}{11.9552}{\normalfont}{escritorio;}
\texto{236.9476}{443.7910}{11.9552}{\normalfont}{la}
\texto{250.6555}{443.7910}{11.9552}{\normalfont}{exposición}
\texto{307.9998}{443.7910}{11.9552}{\normalfont}{directa}
\texto{347.7221}{443.7910}{11.9552}{\normalfont}{al}
\texto{361.4300}{443.7910}{11.9552}{\normalfont}{turista}
\texto{399.9161}{443.7910}{11.9552}{\normalfont}{a}
\texto{410.3722}{443.7910}{11.9552}{\normalfont}{través}
\texto{445.2742}{443.7910}{11.9552}{\normalfont}{de}
\texto{461.5835}{443.7910}{11.9552}{\normalfont}{la}
\texto{475.2913}{443.7910}{11.9552}{\normalfont}{app}
\texto{498.7546}{443.7910}{11.9552}{\normalfont}{móvil}
\texto{531.3253}{443.7910}{11.9552}{\normalfont}{se}
\texto{100.7540}{429.3450}{11.9552}{\normalfont}{contempla}
\texto{157.0080}{429.3450}{11.9552}{\normalfont}{como}
\texto{187.5703}{429.3450}{11.9552}{\normalfont}{trabajo}
\texto{228.5479}{429.3450}{11.9552}{\normalfont}{futuro.}
\texto{85.6720}{404.9360}{11.9552}{\normalfont}{4.}
\texto{100.7540}{404.9360}{11.9552}{\normalfont\bfseries}{Generación}
\texto{170.3966}{404.9360}{11.9552}{\normalfont\bfseries}{de}
\texto{187.4005}{404.9360}{11.9552}{\normalfont\bfseries}{voucher:}
\texto{239.7200}{404.9360}{11.9552}{\normalfont}{al}
\texto{251.7661}{404.9360}{11.9552}{\normalfont}{confirmar}
\texto{303.4962}{404.9360}{11.9552}{\normalfont}{el}
\texto{314.9039}{404.9360}{11.9552}{\normalfont}{pago,}
\texto{345.1720}{404.9360}{11.9552}{\normalfont}{se}
\texto{357.9450}{404.9360}{11.9552}{\normalfont}{genera}
\texto{394.0545}{404.9360}{11.9552}{\normalfont}{un}
\texto{410.0147}{404.9360}{11.9552}{\normalfont}{voucher}
\texto{452.3193}{404.9360}{11.9552}{\normalfont}{digital}
\texto{487.7784}{404.9360}{11.9552}{\normalfont}{con}
\texto{508.2912}{404.9360}{11.9552}{\normalfont}{código}
\texto{100.7540}{390.4900}{11.9552}{\normalfont}{QR}
\texto{122.3714}{390.4900}{11.9552}{\normalfont}{que}
\texto{144.1538}{390.4900}{11.9552}{\normalfont}{el}
\texto{156.5178}{390.4900}{11.9552}{\normalfont}{turista}
\texto{194.2987}{390.4900}{11.9552}{\normalfont}{presenta}
\texto{240.8618}{390.4900}{11.9552}{\normalfont}{al}
\texto{253.8762}{390.4900}{11.9552}{\normalfont}{operador}
\texto{302.9714}{390.4900}{11.9552}{\normalfont}{el}
\texto{315.3355}{390.4900}{11.9552}{\normalfont}{día}
\texto{334.8416}{390.4900}{11.9552}{\normalfont}{del}
\texto{353.6973}{390.4900}{11.9552}{\normalfont}{tour.}
\texto{85.6720}{366.0820}{11.9552}{\normalfont}{5.}
\texto{100.7540}{366.0820}{11.9552}{\normalfont\bfseries}{V}
\texto{109.7981}{366.0820}{11.9552}{\normalfont\bfseries}{alidación}
\texto{165.9254}{366.0820}{11.9552}{\normalfont\bfseries}{y}
\texto{177.2326}{366.0820}{11.9552}{\normalfont\bfseries}{cierre:}
\texto{217.6170}{366.0820}{11.9552}{\normalfont}{el}
\texto{229.7300}{366.0820}{11.9552}{\normalfont}{operador}
\texto{278.5981}{366.0820}{11.9552}{\normalfont}{valida}
\texto{312.5030}{366.0820}{11.9552}{\normalfont}{el}
\texto{324.6280}{366.0820}{11.9552}{\normalfont}{voucher}
\texto{367.6380}{366.0820}{11.9552}{\normalfont}{en}
\texto{383.0028}{366.0820}{11.9552}{\normalfont}{el}
\texto{395.1278}{366.0820}{11.9552}{\normalfont}{backoffice}
\texto{447.8933}{366.0820}{11.9552}{\normalfont}{de}
\texto{463.2701}{366.0820}{11.9552}{\normalfont}{escritorio,}
\texto{517.0708}{366.0820}{11.9552}{\normalfont}{mar-}
\texto{100.7540}{351.6360}{11.9552}{\normalfont}{ca}
\texto{116.2934}{351.6360}{11.9552}{\normalfont}{el}
\texto{129.2432}{351.6360}{11.9552}{\normalfont}{tour}
\texto{155.1884}{351.6360}{11.9552}{\normalfont}{como}
\texto{186.3484}{351.6360}{11.9552}{\normalfont}{completado}
\texto{249.3763}{351.6360}{11.9552}{\normalfont}{y}
\texto{260.0379}{351.6360}{11.9552}{\normalfont}{el}
\texto{272.9878}{351.6360}{11.9552}{\normalfont}{sistema}
\texto{315.3211}{351.6360}{11.9552}{\normalfont}{genera}
\texto{352.9848}{351.6360}{11.9552}{\normalfont}{el}
\texto{365.9347}{351.6360}{11.9552}{\normalfont}{reporte}
\texto{407.1610}{351.6360}{11.9552}{\normalfont}{financiero}
\texto{461.0837}{351.6360}{11.9552}{\normalfont}{correspondiente}
\texto{100.7540}{337.1900}{11.9552}{\normalfont}{(monto}
\texto{141.3993}{337.1900}{11.9552}{\normalfont}{bruto,}
\texto{176.5141}{337.1900}{11.9552}{\normalfont}{comisión}
\texto{224.7007}{337.1900}{11.9552}{\normalfont}{del}
\texto{243.5684}{337.1900}{11.9552}{\normalfont}{8}
\texto{251.4182}{337.1900}{11.9552}{\normalfont}{\%,}
\texto{268.3229}{337.1900}{11.9552}{\normalfont}{monto}
\texto{304.4156}{337.1900}{11.9552}{\normalfont}{neto).}
\texto{85.6720}{312.7820}{11.9552}{\normalfont}{6.}
\texto{100.7540}{312.7820}{11.9552}{\normalfont\bfseries}{Reseña}
\texto{147.6112}{312.7820}{11.9552}{\normalfont\bfseries}{verificada:}
\texto{212.0100}{312.7820}{11.9552}{\normalfont}{después}
\texto{255.7158}{312.7820}{11.9552}{\normalfont}{de}
\texto{271.9892}{312.7820}{11.9552}{\normalfont}{completar}
\texto{327.2844}{312.7820}{11.9552}{\normalfont}{el}
\texto{340.3060}{312.7820}{11.9552}{\normalfont}{tour,}
\texto{369.5748}{312.7820}{11.9552}{\normalfont}{el}
\texto{382.5964}{312.7820}{11.9552}{\normalfont}{turista}
\texto{421.0467}{312.7820}{11.9552}{\normalfont}{puede}
\texto{455.5303}{312.7820}{11.9552}{\normalfont}{dejar}
\texto{485.7745}{312.7820}{11.9552}{\normalfont}{una}
\texto{509.2019}{312.7820}{11.9552}{\normalfont}{reseña}
\texto{100.7540}{298.3360}{11.9552}{\normalfont}{verificada}
\texto{153.4310}{298.3360}{11.9552}{\normalfont}{en}
\texto{169.0349}{298.3360}{11.9552}{\normalfont}{la}
\texto{182.0374}{298.3360}{11.9552}{\normalfont}{plataforma.}
\texto{85.6720}{273.9280}{11.9552}{\normalfont}{7.}
\texto{100.7540}{273.9280}{11.9552}{\normalfont\bfseries}{Comisión}
\texto{161.7949}{273.9280}{11.9552}{\normalfont\bfseries}{y}
\texto{175.3258}{273.9280}{11.9552}{\normalfont\bfseries}{pago:}
\texto{212.1260}{273.9280}{11.9552}{\normalfont}{la}
\texto{226.8261}{273.9280}{11.9552}{\normalfont}{plataforma}
\texto{288.0272}{273.9280}{11.9552}{\normalfont}{retiene}
\texto{328.0914}{273.9280}{11.9552}{\normalfont}{el}
\texto{342.1412}{273.9280}{11.9552}{\normalfont}{8}
\texto{349.9790}{273.9280}{11.9552}{\normalfont}{\%}
\texto{365.3295}{273.9280}{11.9552}{\normalfont}{de}
\texto{382.6311}{273.9280}{11.9552}{\normalfont}{comisión}
\texto{432.5153}{273.9280}{11.9552}{\normalfont}{sobre}
\texto{464.8398}{273.9280}{11.9552}{\normalfont}{cada}
\texto{493.8479}{273.9280}{11.9552}{\normalfont}{reserva}
\texto{534.9570}{273.9280}{11.9552}{\normalfont}{y}
\texto{100.7540}{259.4820}{11.9552}{\normalfont}{transfiere}
\texto{152.1924}{259.4820}{11.9552}{\normalfont}{el}
\texto{164.5446}{259.4820}{11.9552}{\normalfont}{monto}
\texto{200.6373}{259.4820}{11.9552}{\normalfont}{neto}
\texto{226.6590}{259.4820}{11.9552}{\normalfont}{al}
\texto{239.6615}{259.4820}{11.9552}{\normalfont}{operador.}
\end{pagina}
% Página original 39
\begin{pagina}{39}
\texto{70.8660}{746.3890}{11.9552}{\normalfont}{3.6.}
\texto{96.8710}{746.3890}{11.9552}{\normalfont}{MISIÓN}
\texto{529.4280}{746.3890}{11.9552}{\normalfont}{29}
\texto{70.8660}{705.1940}{17.2154}{\normalfont\bfseries}{3.6.}
\texto{120.3603}{705.1940}{17.2154}{\normalfont\bfseries}{Misión}
\texto{88.7990}{678.9090}{11.9552}{\normalfont\bfseries}{Misión:}
\texto{136.4630}{678.9090}{11.9552}{\normalfont}{Conectar}
\texto{186.8231}{678.9090}{11.9552}{\normalfont}{a}
\texto{196.8726}{678.9090}{11.9552}{\normalfont}{los}
\texto{214.7911}{678.9090}{11.9552}{\normalfont}{turistas}
\texto{257.4879}{678.9090}{11.9552}{\normalfont}{con}
\texto{279.2440}{678.9090}{11.9552}{\normalfont}{los}
\texto{297.1624}{678.9090}{11.9552}{\normalfont}{mejores}
\texto{340.1198}{678.9090}{11.9552}{\normalfont}{operadores}
\texto{399.3339}{678.9090}{11.9552}{\normalfont}{turísticos}
\texto{450.4855}{678.9090}{11.9552}{\normalfont}{de}
\texto{466.3883}{678.9090}{11.9552}{\normalfont}{Nicaragua}
\texto{522.2836}{678.9090}{11.9552}{\normalfont}{me-}
\texto{70.8660}{664.4630}{11.9552}{\normalfont}{diante}
\texto{107.2409}{664.4630}{11.9552}{\normalfont}{una}
\texto{130.9433}{664.4630}{11.9552}{\normalfont}{plataforma}
\texto{191.3912}{664.4630}{11.9552}{\normalfont}{digital}
\texto{228.7631}{664.4630}{11.9552}{\normalfont}{segura,}
\texto{269.4395}{664.4630}{11.9552}{\normalfont}{confiable}
\texto{319.1564}{664.4630}{11.9552}{\normalfont}{y}
\texto{330.1767}{664.4630}{11.9552}{\normalfont}{accesible,}
\texto{382.5596}{664.4630}{11.9552}{\normalfont}{que}
\texto{405.2864}{664.4630}{11.9552}{\normalfont}{promueva}
\texto{459.5511}{664.4630}{11.9552}{\normalfont}{la}
\texto{473.4980}{664.4630}{11.9552}{\normalfont}{digitalización}
\texto{70.8660}{650.0170}{11.9552}{\normalfont}{del}
\texto{89.7935}{650.0170}{11.9552}{\normalfont}{sector,}
\texto{127.0076}{650.0170}{11.9552}{\normalfont}{fortalezca}
\texto{180.0911}{650.0170}{11.9552}{\normalfont}{la}
\texto{193.1653}{650.0170}{11.9552}{\normalfont}{economía}
\texto{244.6229}{650.0170}{11.9552}{\normalfont}{local}
\texto{272.3398}{650.0170}{11.9552}{\normalfont}{y}
\texto{282.4993}{650.0170}{11.9552}{\normalfont}{brinde}
\texto{318.9986}{650.0170}{11.9552}{\normalfont}{experiencias}
\texto{384.8143}{650.0170}{11.9552}{\normalfont}{auténticas}
\texto{440.5662}{650.0170}{11.9552}{\normalfont}{a}
\texto{450.4006}{650.0170}{11.9552}{\normalfont}{los}
\texto{468.0919}{650.0170}{11.9552}{\normalfont}{visitantes,}
\texto{523.5831}{650.0170}{11.9552}{\normalfont}{con}
\texto{70.8660}{635.5710}{11.9552}{\normalfont}{comisiones}
\texto{128.8726}{635.5710}{11.9552}{\normalfont}{justas}
\texto{162.4907}{635.5710}{11.9552}{\normalfont}{que}
\texto{184.2850}{635.5710}{11.9552}{\normalfont}{beneficien}
\texto{238.5831}{635.5710}{11.9552}{\normalfont}{directamente}
\texto{308.4947}{635.5710}{11.9552}{\normalfont}{a}
\texto{318.2573}{635.5710}{11.9552}{\normalfont}{los}
\texto{335.8769}{635.5710}{11.9552}{\normalfont}{proveedores}
\texto{399.3519}{635.5710}{11.9552}{\normalfont}{nicaragüenses.}
\end{pagina}
% Página original 40
\begin{pagina}{40}
\texto{70.8660}{746.3890}{11.9552}{\normalfont}{30}
\texto{288.7940}{746.3890}{11.9552}{\normalfont}{CAPÍTULO}
\texto{355.7575}{746.3890}{11.9552}{\normalfont}{3.}
\texto{372.6693}{746.3890}{11.9552}{\normalfont}{EMPRENDIMIENTO}
\texto{490.9875}{746.3890}{11.9552}{\normalfont}{DIGITAL}
\texto{70.8660}{705.1940}{17.2154}{\normalfont\bfseries}{3.7.}
\texto{120.3603}{705.1940}{17.2154}{\normalfont\bfseries}{Visión}
\texto{88.7990}{678.9090}{11.9552}{\normalfont\bfseries}{Visión:}
\texto{133.3960}{678.9090}{11.9552}{\normalfont}{Ser}
\texto{153.3851}{678.9090}{11.9552}{\normalfont}{el}
\texto{165.5698}{678.9090}{11.9552}{\normalfont}{marketplace}
\texto{231.0867}{678.9090}{11.9552}{\normalfont}{turístico}
\texto{277.1549}{678.9090}{11.9552}{\normalfont}{líder}
\texto{303.6476}{678.9090}{11.9552}{\normalfont}{en}
\texto{319.0842}{678.9090}{11.9552}{\normalfont}{Nicaragua}
\texto{374.5133}{678.9090}{11.9552}{\normalfont}{y}
\texto{384.4217}{678.9090}{11.9552}{\normalfont}{Centroamérica,}
\texto{465.8725}{678.9090}{11.9552}{\normalfont}{reconocido}
\texto{523.9054}{678.9090}{11.9552}{\normalfont}{por}
\texto{70.8660}{664.4630}{11.9552}{\normalfont}{la}
\texto{84.1674}{664.4630}{11.9552}{\normalfont}{calidad}
\texto{124.7840}{664.4630}{11.9552}{\normalfont}{de}
\texto{140.6868}{664.4630}{11.9552}{\normalfont}{su}
\texto{156.0038}{664.4630}{11.9552}{\normalfont}{verificación,}
\texto{220.6862}{664.4630}{11.9552}{\normalfont}{la}
\texto{233.9875}{664.4630}{11.9552}{\normalfont}{confianza}
\texto{285.6603}{664.4630}{11.9552}{\normalfont}{de}
\texto{301.5631}{664.4630}{11.9552}{\normalfont}{sus}
\texto{321.4972}{664.4630}{11.9552}{\normalfont}{usuarios}
\texto{367.4458}{664.4630}{11.9552}{\normalfont}{y}
\texto{377.8205}{664.4630}{11.9552}{\normalfont}{el}
\texto{390.4715}{664.4630}{11.9552}{\normalfont}{apoyo}
\texto{424.5869}{664.4630}{11.9552}{\normalfont}{al}
\texto{437.8882}{664.4630}{11.9552}{\normalfont}{desarrollo}
\texto{491.5767}{664.4630}{11.9552}{\normalfont}{sostenible}
\texto{70.8660}{650.0170}{11.9552}{\normalfont}{del}
\texto{90.0804}{650.0170}{11.9552}{\normalfont}{turismo}
\texto{133.4347}{650.0170}{11.9552}{\normalfont}{local,}
\texto{164.6904}{650.0170}{11.9552}{\normalfont}{contribuyendo}
\texto{241.4787}{650.0170}{11.9552}{\normalfont}{a}
\texto{251.5880}{650.0170}{11.9552}{\normalfont}{que}
\texto{273.7410}{650.0170}{11.9552}{\normalfont}{cada}
\texto{301.4220}{650.0170}{11.9552}{\normalfont}{operador}
\texto{350.8879}{650.0170}{11.9552}{\normalfont}{turístico}
\texto{397.4940}{650.0170}{11.9552}{\normalfont}{nicaragüense}
\texto{466.8629}{650.0170}{11.9552}{\normalfont}{tenga}
\texto{499.0965}{650.0170}{11.9552}{\normalfont}{acceso}
\texto{535.2849}{650.0170}{11.9552}{\normalfont}{a}
\texto{70.8660}{635.5710}{11.9552}{\normalfont}{herramientas}
\texto{140.8422}{635.5710}{11.9552}{\normalfont}{digitales}
\texto{187.0777}{635.5710}{11.9552}{\normalfont}{que}
\texto{208.8601}{635.5710}{11.9552}{\normalfont}{impulsen}
\texto{258.3594}{635.5710}{11.9552}{\normalfont}{su}
\texto{273.3775}{635.5710}{11.9552}{\normalfont}{crecimiento.}
\end{pagina}
% Página original 41
\begin{pagina}{41}
\texto{70.8660}{746.3890}{11.9552}{\normalfont}{3.8.}
\texto{96.8710}{746.3890}{11.9552}{\normalfont}{V}
\texto{104.3549}{746.3890}{11.9552}{\normalfont}{ALORES}
\texto{529.4280}{746.3890}{11.9552}{\normalfont}{31}
\texto{70.8660}{705.1940}{17.2154}{\normalfont\bfseries}{3.8.}
\texto{120.3603}{705.1940}{17.2154}{\normalfont\bfseries}{V}
\texto{133.3837}{705.1940}{17.2154}{\normalfont\bfseries}{alores}
\texto{88.7990}{678.9090}{11.9552}{\normalfont}{Los}
\texto{110.4810}{678.9090}{11.9552}{\normalfont}{valores}
\texto{149.2421}{678.9090}{11.9552}{\normalfont}{que}
\texto{171.0245}{678.9090}{11.9552}{\normalfont}{orientan}
\texto{216.8727}{678.9090}{11.9552}{\normalfont}{el}
\texto{229.2367}{678.9090}{11.9552}{\normalfont}{emprendimiento}
\texto{315.4074}{678.9090}{11.9552}{\normalfont}{son:}
\texto{100.7540}{651.5120}{11.9552}{\normalfont\bfseries}{Confianza:}
\texto{165.4610}{651.5120}{11.9552}{\normalfont}{cada}
\texto{191.9346}{651.5120}{11.9552}{\normalfont}{operador}
\texto{240.1930}{651.5120}{11.9552}{\normalfont}{es}
\texto{253.0735}{651.5120}{11.9552}{\normalfont}{verificado}
\texto{304.9136}{651.5120}{11.9552}{\normalfont}{documentalmente}
\texto{398.0399}{651.5120}{11.9552}{\normalfont}{y}
\texto{407.2788}{651.5120}{11.9552}{\normalfont}{cada}
\texto{433.7524}{651.5120}{11.9552}{\normalfont}{reseña}
\texto{468.7453}{651.5120}{11.9552}{\normalfont}{es}
\texto{481.6378}{651.5120}{11.9552}{\normalfont}{autenticada}
\texto{100.7540}{637.0660}{11.9552}{\normalfont}{mediante}
\texto{151.2265}{637.0660}{11.9552}{\normalfont}{voucher,}
\texto{197.7991}{637.0660}{11.9552}{\normalfont}{construyendo}
\texto{269.4849}{637.0660}{11.9552}{\normalfont}{un}
\texto{286.4732}{637.0660}{11.9552}{\normalfont}{ecosistema}
\texto{344.5636}{637.0660}{11.9552}{\normalfont}{transparente}
\texto{412.6604}{637.0660}{11.9552}{\normalfont}{para}
\texto{439.4042}{637.0660}{11.9552}{\normalfont}{turistas}
\texto{481.8858}{637.0660}{11.9552}{\normalfont}{y}
\texto{492.0333}{637.0660}{11.9552}{\normalfont}{proveedo-}
\texto{100.7540}{622.6200}{11.9552}{\normalfont}{res.}
\texto{100.7540}{598.2110}{11.9552}{\normalfont\bfseries}{Equidad:}
\texto{158.8730}{598.2110}{11.9552}{\normalfont}{comisiones}
\texto{218.7327}{598.2110}{11.9552}{\normalfont}{reducidas}
\texto{272.6865}{598.2110}{11.9552}{\normalfont}{que}
\texto{296.3339}{598.2110}{11.9552}{\normalfont}{favorecen}
\texto{349.5776}{598.2110}{11.9552}{\normalfont}{a}
\texto{361.1932}{598.2110}{11.9552}{\normalfont}{los}
\texto{380.6778}{598.2110}{11.9552}{\normalfont}{operadores}
\texto{441.4581}{598.2110}{11.9552}{\normalfont}{locales}
\texto{480.7763}{598.2110}{11.9552}{\normalfont}{frente}
\texto{515.7955}{598.2110}{11.9552}{\normalfont}{a}
\texto{527.4112}{598.2110}{11.9552}{\normalfont}{las}
\texto{100.7540}{583.7660}{11.9552}{\normalfont}{plataformas}
\texto{164.8745}{583.7660}{11.9552}{\normalfont}{internacionales,}
\texto{247.8580}{583.7660}{11.9552}{\normalfont}{promoviendo}
\texto{317.4444}{583.7660}{11.9552}{\normalfont}{la}
\texto{330.4588}{583.7660}{11.9552}{\normalfont}{retención}
\texto{381.1823}{583.7660}{11.9552}{\normalfont}{de}
\texto{396.7863}{583.7660}{11.9552}{\normalfont}{valor}
\texto{425.7274}{583.7660}{11.9552}{\normalfont}{en}
\texto{441.3433}{583.7660}{11.9552}{\normalfont}{Nicaragua.}
\texto{100.7540}{559.3570}{11.9552}{\normalfont\bfseries}{Accesibilidad:}
\texto{185.1150}{559.3570}{11.9552}{\normalfont}{una}
\texto{207.5143}{559.3570}{11.9552}{\normalfont}{plataforma}
\texto{266.6710}{559.3570}{11.9552}{\normalfont}{fácil}
\texto{291.3465}{559.3570}{11.9552}{\normalfont}{de}
\texto{306.6037}{559.3570}{11.9552}{\normalfont}{usar}
\texto{331.6810}{559.3570}{11.9552}{\normalfont}{tanto}
\texto{362.2122}{559.3570}{11.9552}{\normalfont}{para}
\texto{388.5256}{559.3570}{11.9552}{\normalfont}{el}
\texto{400.5190}{559.3570}{11.9552}{\normalfont}{turista}
\texto{437.9531}{559.3570}{11.9552}{\normalfont}{que}
\texto{459.3888}{559.3570}{11.9552}{\normalfont}{reserva}
\texto{498.4417}{559.3570}{11.9552}{\normalfont}{desde}
\texto{530.0225}{559.3570}{11.9552}{\normalfont}{su}
\texto{100.7540}{544.9110}{11.9552}{\normalfont}{teléfono}
\texto{144.6487}{544.9110}{11.9552}{\normalfont}{como}
\texto{175.2110}{544.9110}{11.9552}{\normalfont}{para}
\texto{201.8830}{544.9110}{11.9552}{\normalfont}{el}
\texto{214.2351}{544.9110}{11.9552}{\normalfont}{operador}
\texto{263.3423}{544.9110}{11.9552}{\normalfont}{que}
\texto{285.1247}{544.9110}{11.9552}{\normalfont}{gestiona}
\texto{330.7099}{544.9110}{11.9552}{\normalfont}{desde}
\texto{362.6494}{544.9110}{11.9552}{\normalfont}{su}
\texto{377.6675}{544.9110}{11.9552}{\normalfont}{escritorio.}
\texto{100.7540}{520.5030}{11.9552}{\normalfont\bfseries}{Innovación}
\texto{168.6141}{520.5030}{11.9552}{\normalfont\bfseries}{tecnológica:}
\texto{241.4800}{520.5030}{11.9552}{\normalfont}{aprovechar}
\texto{300.9045}{520.5030}{11.9552}{\normalfont}{las}
\texto{318.7632}{520.5030}{11.9552}{\normalfont}{tendencias}
\texto{376.3060}{520.5030}{11.9552}{\normalfont}{actuales}
\texto{421.4918}{520.5030}{11.9552}{\normalfont}{(API}
\texto{451.1431}{520.5030}{11.9552}{\normalfont}{REST,}
\texto{490.0812}{520.5030}{11.9552}{\normalfont}{autentica-}
\texto{100.7540}{506.0570}{11.9552}{\normalfont}{ción}
\texto{124.4109}{506.0570}{11.9552}{\normalfont}{JWT,}
\texto{157.0032}{506.0570}{11.9552}{\normalfont}{control}
\texto{195.2957}{506.0570}{11.9552}{\normalfont}{de}
\texto{209.8476}{506.0570}{11.9552}{\normalfont}{acceso}
\texto{244.6253}{506.0570}{11.9552}{\normalfont}{basado}
\texto{282.6667}{506.0570}{11.9552}{\normalfont}{en}
\texto{297.2186}{506.0570}{11.9552}{\normalfont}{roles,}
\texto{326.7934}{506.0570}{11.9552}{\normalfont}{pagos}
\texto{358.3312}{506.0570}{11.9552}{\normalfont}{en}
\texto{372.8830}{506.0570}{11.9552}{\normalfont}{línea)}
\texto{404.3443}{506.0570}{11.9552}{\normalfont}{para}
\texto{429.9524}{506.0570}{11.9552}{\normalfont}{construir}
\texto{478.3877}{506.0570}{11.9552}{\normalfont}{una}
\texto{500.1055}{506.0570}{11.9552}{\normalfont}{solución}
\texto{100.7540}{491.6110}{11.9552}{\normalfont}{moderna}
\texto{149.1989}{491.6110}{11.9552}{\normalfont}{y}
\texto{159.2747}{491.6110}{11.9552}{\normalfont}{escalable.}
\texto{100.7540}{467.2030}{11.9552}{\normalfont\bfseries}{Responsabilidad}
\texto{201.2948}{467.2030}{11.9552}{\normalfont\bfseries}{social:}
\texto{241.8020}{467.2030}{11.9552}{\normalfont}{contribuir}
\texto{296.5712}{467.2030}{11.9552}{\normalfont}{a}
\texto{306.8000}{467.2030}{11.9552}{\normalfont}{la}
\texto{320.2807}{467.2030}{11.9552}{\normalfont}{digitalización}
\texto{392.2821}{467.2030}{11.9552}{\normalfont}{del}
\texto{411.6161}{467.2030}{11.9552}{\normalfont}{turismo}
\texto{455.0780}{467.2030}{11.9552}{\normalfont}{nicaragüense,}
\texto{527.8062}{467.2030}{11.9552}{\normalfont}{fo-}
\texto{100.7540}{452.7570}{11.9552}{\normalfont}{mentando}
\texto{155.0569}{452.7570}{11.9552}{\normalfont}{el}
\texto{167.4090}{452.7570}{11.9552}{\normalfont}{empleo}
\texto{207.0764}{452.7570}{11.9552}{\normalfont}{formal}
\texto{243.8291}{452.7570}{11.9552}{\normalfont}{y}
\texto{253.9049}{452.7570}{11.9552}{\normalfont}{el}
\texto{266.2690}{452.7570}{11.9552}{\normalfont}{desarrollo}
\texto{319.6585}{452.7570}{11.9552}{\normalfont}{sostenible}
\texto{373.1126}{452.7570}{11.9552}{\normalfont}{de}
\texto{388.7285}{452.7570}{11.9552}{\normalfont}{las}
\texto{406.3480}{452.7570}{11.9552}{\normalfont}{comunidades}
\texto{475.6739}{452.7570}{11.9552}{\normalfont}{turísticas.}
\end{pagina}
% Página original 42
\begin{pagina}{42}
\texto{70.8660}{746.3890}{11.9552}{\normalfont}{32}
\texto{288.7940}{746.3890}{11.9552}{\normalfont}{CAPÍTULO}
\texto{355.7575}{746.3890}{11.9552}{\normalfont}{3.}
\texto{372.6693}{746.3890}{11.9552}{\normalfont}{EMPRENDIMIENTO}
\texto{490.9875}{746.3890}{11.9552}{\normalfont}{DIGITAL}
\texto{70.8660}{705.1940}{17.2154}{\normalfont\bfseries}{3.9.}
\texto{120.3603}{705.1940}{17.2154}{\normalfont\bfseries}{Identidad}
\texto{207.6802}{705.1940}{17.2154}{\normalfont\bfseries}{visual}
\texto{88.7990}{678.9090}{11.9552}{\normalfont}{La}
\texto{105.8639}{678.9090}{11.9552}{\normalfont}{identidad}
\texto{157.5653}{678.9090}{11.9552}{\normalfont}{visual}
\texto{191.1188}{678.9090}{11.9552}{\normalfont}{de}
\texto{206.7370}{678.9090}{11.9552}{\normalfont\bfseries}{Explora-Nicaragua}
\texto{320.8070}{678.9090}{11.9552}{\normalfont}{se}
\texto{334.5244}{678.9090}{11.9552}{\normalfont}{fundamenta}
\texto{398.9079}{678.9090}{11.9552}{\normalfont}{en}
\texto{414.5119}{678.9090}{11.9552}{\normalfont}{los}
\texto{432.1434}{678.9090}{11.9552}{\normalfont}{siguientes}
\texto{485.2747}{678.9090}{11.9552}{\normalfont}{elementos:}
\texto{100.7540}{651.5120}{11.9552}{\normalfont\bfseries}{Nombre}
\texto{152.7280}{651.5120}{11.9552}{\normalfont\bfseries}{de}
\texto{171.1426}{651.5120}{11.9552}{\normalfont\bfseries}{marca:}
\texto{214.8180}{651.5120}{11.9552}{\normalfont}{Explora-Nicaragua,}
\texto{318.0009}{651.5120}{11.9552}{\normalfont}{que}
\texto{340.0702}{651.5120}{11.9552}{\normalfont}{evoca}
\texto{372.5453}{651.5120}{11.9552}{\normalfont}{la}
\texto{385.8228}{651.5120}{11.9552}{\normalfont}{acción}
\texto{421.8749}{651.5120}{11.9552}{\normalfont}{de}
\texto{437.7538}{651.5120}{11.9552}{\normalfont}{descubrir}
\texto{488.8288}{651.5120}{11.9552}{\normalfont}{y}
\texto{499.1916}{651.5120}{11.9552}{\normalfont}{explorar}
\texto{100.7540}{637.0660}{11.9552}{\normalfont}{los}
\texto{118.3736}{637.0660}{11.9552}{\normalfont}{destinos}
\texto{163.3729}{637.0660}{11.9552}{\normalfont}{del}
\texto{182.2406}{637.0660}{11.9552}{\normalfont}{país.}
\texto{100.7540}{612.6570}{11.9552}{\normalfont\bfseries}{Logotipo:}
\texto{161.6350}{612.6570}{11.9552}{\normalfont}{diseño}
\texto{198.9113}{612.6570}{11.9552}{\normalfont}{que}
\texto{222.1403}{612.6570}{11.9552}{\normalfont}{integra}
\texto{262.9314}{612.6570}{11.9552}{\normalfont}{elementos}
\texto{318.1070}{612.6570}{11.9552}{\normalfont}{que}
\texto{341.3360}{612.6570}{11.9552}{\normalfont}{representan}
\texto{405.6048}{612.6570}{11.9552}{\normalfont}{la}
\texto{420.0658}{612.6570}{11.9552}{\normalfont}{naturaleza}
\texto{478.7395}{612.6570}{11.9552}{\normalfont}{y}
\texto{490.2619}{612.6570}{11.9552}{\normalfont}{la}
\texto{504.7229}{612.6570}{11.9552}{\normalfont}{cultura}
\texto{100.7540}{598.2110}{11.9552}{\normalfont}{nicaragüense,}
\texto{173.0041}{598.2110}{11.9552}{\normalfont}{combinando}
\texto{238.3631}{598.2110}{11.9552}{\normalfont}{colores}
\texto{276.7943}{598.2110}{11.9552}{\normalfont}{vibrantes}
\texto{327.5967}{598.2110}{11.9552}{\normalfont}{asociados}
\texto{379.4225}{598.2110}{11.9552}{\normalfont}{al}
\texto{392.4370}{598.2110}{11.9552}{\normalfont}{turismo}
\texto{435.4207}{598.2110}{11.9552}{\normalfont}{tropical.}
\texto{100.7540}{573.8030}{11.9552}{\normalfont\bfseries}{Paleta}
\texto{143.7401}{573.8030}{11.9552}{\normalfont\bfseries}{de}
\texto{163.3383}{573.8030}{11.9552}{\normalfont\bfseries}{colores:}
\texto{212.4010}{573.8030}{11.9552}{\normalfont}{tonos}
\texto{244.9933}{573.8030}{11.9552}{\normalfont}{de}
\texto{261.9123}{573.8030}{11.9552}{\normalfont}{verde}
\texto{294.4424}{573.8030}{11.9552}{\normalfont}{(naturaleza}
\texto{357.5491}{573.8030}{11.9552}{\normalfont}{y}
\texto{368.9400}{573.8030}{11.9552}{\normalfont}{sostenibilidad),}
\texto{451.6867}{573.8030}{11.9552}{\normalfont}{azul}
\texto{477.7108}{573.8030}{11.9552}{\normalfont}{(confianza}
\texto{534.9643}{573.8030}{11.9552}{\normalfont}{y}
\texto{100.7540}{559.3570}{11.9552}{\normalfont}{tecnología)}
\texto{160.5826}{559.3570}{11.9552}{\normalfont}{y}
\texto{170.6584}{559.3570}{11.9552}{\normalfont}{acentos}
\texto{212.0187}{559.3570}{11.9552}{\normalfont}{de}
\texto{227.6345}{559.3570}{11.9552}{\normalfont}{color}
\texto{256.2457}{559.3570}{11.9552}{\normalfont}{cálido}
\texto{290.0718}{559.3570}{11.9552}{\normalfont}{(energía}
\texto{334.9420}{559.3570}{11.9552}{\normalfont}{y}
\texto{345.0179}{559.3570}{11.9552}{\normalfont}{hospitalidad).}
\texto{100.7540}{534.9490}{11.9552}{\normalfont\bfseries}{Tipografía:}
\texto{169.7510}{534.9490}{11.9552}{\normalfont}{fuentes}
\texto{210.0352}{534.9490}{11.9552}{\normalfont}{modernas}
\texto{263.6711}{534.9490}{11.9552}{\normalfont}{y}
\texto{274.3088}{534.9490}{11.9552}{\normalfont}{legibles}
\texto{315.9033}{534.9490}{11.9552}{\normalfont}{que}
\texto{338.2476}{534.9490}{11.9552}{\normalfont}{transmiten}
\texto{398.0523}{534.9490}{11.9552}{\normalfont}{profesionalismo}
\texto{481.0094}{534.9490}{11.9552}{\normalfont}{y}
\texto{491.6472}{534.9490}{11.9552}{\normalfont}{accesibili-}
\texto{100.7540}{520.5030}{11.9552}{\normalfont}{dad.}
\texto{100.7540}{496.0940}{11.9552}{\normalfont\bfseries}{Aplicación:}
\texto{170.0860}{496.0940}{11.9552}{\normalfont}{la}
\texto{183.2558}{496.0940}{11.9552}{\normalfont}{identidad}
\texto{235.1247}{496.0940}{11.9552}{\normalfont}{visual}
\texto{268.8575}{496.0940}{11.9552}{\normalfont}{se}
\texto{282.7422}{496.0940}{11.9552}{\normalfont}{aplica}
\texto{316.7357}{496.0940}{11.9552}{\normalfont}{de}
\texto{332.5070}{496.0940}{11.9552}{\normalfont}{forma}
\texto{366.1632}{496.0940}{11.9552}{\normalfont}{consistente}
\texto{425.9775}{496.0940}{11.9552}{\normalfont}{en}
\texto{441.7488}{496.0940}{11.9552}{\normalfont}{las}
\texto{459.5477}{496.0940}{11.9552}{\normalfont}{dos}
\texto{480.5865}{496.0940}{11.9552}{\normalfont}{aplicaciones}
\texto{100.7540}{481.6490}{11.9552}{\normalfont}{de}
\texto{116.0471}{481.6490}{11.9552}{\normalfont}{escritorio}
\texto{166.5363}{481.6490}{11.9552}{\normalfont}{del}
\texto{185.0932}{481.6490}{11.9552}{\normalfont}{prototipo}
\texto{236.4909}{481.6490}{11.9552}{\normalfont}{inicial}
\texto{270.6565}{481.6490}{11.9552}{\normalfont}{(backoffice}
\texto{327.9148}{481.6490}{11.9552}{\normalfont}{de}
\texto{343.2198}{481.6490}{11.9552}{\normalfont}{operadores}
\texto{401.8361}{481.6490}{11.9552}{\normalfont}{y}
\texto{411.6131}{481.6490}{11.9552}{\normalfont}{panel}
\texto{442.5269}{481.6490}{11.9552}{\normalfont}{de}
\texto{457.8319}{481.6490}{11.9552}{\normalfont}{administración),}
\texto{100.7540}{467.2030}{11.9552}{\normalfont}{así}
\texto{118.8876}{467.2030}{11.9552}{\normalfont}{como}
\texto{149.9640}{467.2030}{11.9552}{\normalfont}{en}
\texto{166.0820}{467.2030}{11.9552}{\normalfont}{los}
\texto{184.2276}{467.2030}{11.9552}{\normalfont}{materiales}
\texto{240.7327}{467.2030}{11.9552}{\normalfont}{promocionales}
\texto{317.7218}{467.2030}{11.9552}{\normalfont}{y}
\texto{328.3236}{467.2030}{11.9552}{\normalfont}{la}
\texto{341.8402}{467.2030}{11.9552}{\normalfont}{documentación}
\texto{422.9945}{467.2030}{11.9552}{\normalfont}{de}
\texto{439.1244}{467.2030}{11.9552}{\normalfont}{la}
\texto{452.6410}{467.2030}{11.9552}{\normalfont}{plataforma;}
\texto{515.9103}{467.2030}{11.9552}{\normalfont}{en}
\texto{532.0283}{467.2030}{11.9552}{\normalfont}{la}
\texto{100.7540}{452.7570}{11.9552}{\normalfont}{visión}
\texto{134.3075}{452.7570}{11.9552}{\normalfont}{completa}
\texto{184.3806}{452.7570}{11.9552}{\normalfont}{se}
\texto{198.1100}{452.7570}{11.9552}{\normalfont}{extenderá}
\texto{251.7601}{452.7570}{11.9552}{\normalfont}{a}
\texto{261.5108}{452.7570}{11.9552}{\normalfont}{la}
\texto{274.5252}{452.7570}{11.9552}{\normalfont}{aplicación}
\texto{329.1509}{452.7570}{11.9552}{\normalfont}{móvil.}
\texto{88.7990}{425.3600}{11.9552}{\normalfont}{La}
\texto{105.9954}{425.3600}{11.9552}{\normalfont}{coherencia}
\texto{163.3660}{425.3600}{11.9552}{\normalfont}{visual}
\texto{197.0629}{425.3600}{11.9552}{\normalfont}{entre}
\texto{226.7835}{425.3600}{11.9552}{\normalfont}{los}
\texto{244.5466}{425.3600}{11.9552}{\normalfont}{componentes}
\texto{314.3386}{425.3600}{11.9552}{\normalfont}{del}
\texto{333.3259}{425.3600}{11.9552}{\normalfont}{ecosistema}
\texto{391.4760}{425.3600}{11.9552}{\normalfont}{refuerza}
\texto{436.1645}{425.3600}{11.9552}{\normalfont}{el}
\texto{448.6601}{425.3600}{11.9552}{\normalfont}{reconocimiento}
\texto{529.4318}{425.3600}{11.9552}{\normalfont}{de}
\texto{70.8660}{410.9140}{11.9552}{\normalfont}{marca}
\texto{105.9808}{410.9140}{11.9552}{\normalfont}{y}
\texto{116.0567}{410.9140}{11.9552}{\normalfont}{genera}
\texto{153.1345}{410.9140}{11.9552}{\normalfont}{una}
\texto{175.8924}{410.9140}{11.9552}{\normalfont}{experiencia}
\texto{237.0313}{410.9140}{11.9552}{\normalfont}{unificada}
\texto{287.1045}{410.9140}{11.9552}{\normalfont}{para}
\texto{313.7646}{410.9140}{11.9552}{\normalfont}{todos}
\texto{345.3765}{410.9140}{11.9552}{\normalfont}{los}
\texto{362.9961}{410.9140}{11.9552}{\normalfont}{actores}
\texto{402.7400}{410.9140}{11.9552}{\normalfont}{del}
\texto{421.5957}{410.9140}{11.9552}{\normalfont}{marketplace.}
\end{pagina}
% Página original 43
\begin{pagina}{43}
\texto{70.8660}{746.3890}{11.9552}{\normalfont}{3.10.}
\texto{102.7242}{746.3890}{11.9552}{\normalfont}{ANÁLISIS}
\texto{161.7255}{746.3890}{11.9552}{\normalfont}{FODA}
\texto{529.4280}{746.3890}{11.9552}{\normalfont}{33}
\texto{70.8660}{705.1940}{17.2154}{\normalfont\bfseries}{3.10.}
\texto{130.0439}{705.1940}{17.2154}{\normalfont\bfseries}{Análisis}
\texto{202.7187}{705.1940}{17.2154}{\normalfont\bfseries}{FODA}
\texto{88.7990}{678.9090}{11.9552}{\normalfont}{El}
\texto{104.4866}{678.9090}{11.9552}{\normalfont}{análisis}
\texto{146.1696}{678.9090}{11.9552}{\normalfont}{FODA}
\texto{184.4549}{678.9090}{11.9552}{\normalfont}{del}
\texto{203.8965}{678.9090}{11.9552}{\normalfont}{emprendimiento}
\texto{290.6550}{678.9090}{11.9552}{\normalfont\bfseries}{Explora-Nicaragua}
\texto{405.3050}{678.9090}{11.9552}{\normalfont}{se}
\texto{419.5962}{678.9090}{11.9552}{\normalfont}{presenta}
\texto{466.7452}{678.9090}{11.9552}{\normalfont}{en}
\texto{482.9349}{678.9090}{11.9552}{\normalfont}{la}
\texto{496.5232}{678.9090}{11.9552}{\normalfont}{siguiente}
\texto{70.8660}{664.4630}{11.9552}{\normalfont}{tabla:}
\texto{177.5840}{632.3170}{11.9552}{\normalfont}{Cuadro}
\texto{219.2024}{632.3170}{11.9552}{\normalfont}{3.1:}
\texto{241.3170}{632.3170}{11.9552}{\normalfont}{Análisis}
\texto{285.3360}{632.3170}{11.9552}{\normalfont}{FODA}
\texto{323.0475}{632.3170}{11.9552}{\normalfont}{de}
\texto{338.6515}{632.3170}{11.9552}{\normalfont}{Explora-Nicaragua}
\texto{115.7700}{614.6020}{11.9552}{\normalfont\bfseries}{F}
\texto{123.1129}{614.6020}{11.9552}{\normalfont\bfseries}{actores}
\texto{168.9910}{614.6020}{11.9552}{\normalfont\bfseries}{Internos}
\texto{311.9780}{614.6020}{11.9552}{\normalfont\bfseries}{F}
\texto{319.3209}{614.6020}{11.9552}{\normalfont\bfseries}{actores}
\texto{365.1990}{614.6020}{11.9552}{\normalfont\bfseries}{Externos}
\texto{115.7700}{594.1660}{11.9552}{\normalfont\bfseries}{F}
\texto{123.1129}{594.1660}{11.9552}{\normalfont\bfseries}{ortalezas}
\texto{311.9780}{594.1660}{11.9552}{\normalfont\bfseries}{Oportunidades}
\texto{127.9250}{570.7540}{11.9552}{\normalfont}{Comisión}
\texto{179.3395}{570.7540}{11.9552}{\normalfont}{reducida}
\texto{226.7993}{570.7540}{11.9552}{\normalfont}{(8}
\texto{239.1897}{570.7540}{11.9552}{\normalfont}{\%)}
\texto{257.3831}{570.7540}{11.9552}{\normalfont}{vs.}
\texto{275.3039}{570.7540}{11.9552}{\normalfont}{com-}
\texto{127.9250}{556.3080}{11.9552}{\normalfont}{petencia}
\texto{174.4188}{556.3080}{11.9552}{\normalfont}{(15}
\texto{192.6744}{556.3080}{11.9552}{\normalfont}{\%--25}
\texto{221.9742}{556.3080}{11.9552}{\normalfont}{\%)}
\texto{127.9250}{541.8620}{11.9552}{\normalfont}{Ecosistema}
\texto{188.2988}{541.8620}{11.9552}{\normalfont}{integrado}
\texto{239.6057}{541.8620}{11.9552}{\normalfont}{de}
\texto{254.8271}{541.8620}{11.9552}{\normalfont}{tres}
\texto{277.2550}{541.8620}{11.9552}{\normalfont}{apli-}
\texto{127.9250}{527.4160}{11.9552}{\normalfont}{caciones}
\texto{173.4982}{527.4160}{11.9552}{\normalfont}{(prototipo}
\texto{229.7355}{527.4160}{11.9552}{\normalfont}{inicial:}
\texto{267.4517}{527.4160}{11.9552}{\normalfont}{dos}
\texto{288.3112}{527.4160}{11.9552}{\normalfont}{de}
\texto{127.9250}{512.9710}{11.9552}{\normalfont}{escritorio)}
\texto{127.9250}{498.5250}{11.9552}{\normalfont}{Verificación}
\texto{190.3742}{498.5250}{11.9552}{\normalfont}{documental}
\texto{252.8449}{498.5250}{11.9552}{\normalfont}{de}
\texto{267.8271}{498.5250}{11.9552}{\normalfont}{opera-}
\texto{127.9250}{484.0790}{11.9552}{\normalfont}{dores}
\texto{127.9250}{469.6330}{11.9552}{\normalfont}{Reseñas}
\texto{172.4342}{469.6330}{11.9552}{\normalfont}{verificadas}
\texto{229.7283}{469.6330}{11.9552}{\normalfont}{por}
\texto{250.8699}{469.6330}{11.9552}{\normalfont}{voucher}
\texto{127.9250}{455.1870}{11.9552}{\normalfont}{Arquitectura}
\texto{206.6500}{455.1870}{11.9552}{\normalfont}{técnica}
\texto{255.4750}{455.1870}{11.9552}{\normalfont}{moderna}
\texto{127.9250}{440.7420}{11.9552}{\normalfont}{(API}
\texto{157.3372}{440.7420}{11.9552}{\normalfont}{REST,}
\texto{196.0242}{440.7420}{11.9552}{\normalfont}{JWT,}
\texto{229.6805}{440.7420}{11.9552}{\normalfont}{RBAC)}
\texto{127.9250}{426.2960}{11.9552}{\normalfont}{Backoffice}
\texto{183.8753}{426.2960}{11.9552}{\normalfont}{de}
\texto{200.6509}{426.2960}{11.9552}{\normalfont}{escritorio}
\texto{252.6106}{426.2960}{11.9552}{\normalfont}{para}
\texto{280.4423}{426.2960}{11.9552}{\normalfont}{ges-}
\texto{127.9250}{411.8500}{11.9552}{\normalfont}{tión}
\texto{151.9836}{411.8500}{11.9552}{\normalfont}{eficiente}
\texto{324.1320}{571.4850}{11.9552}{\normalfont}{Crecimiento}
\texto{390.0649}{571.4850}{11.9552}{\normalfont}{del}
\texto{409.4945}{571.4850}{11.9552}{\normalfont}{turismo}
\texto{453.0640}{571.4850}{11.9552}{\normalfont}{en}
\texto{469.2418}{571.4850}{11.9552}{\normalfont}{Nica-}
\texto{324.1320}{557.0390}{11.9552}{\normalfont}{ragua}
\texto{356.6454}{557.0390}{11.9552}{\normalfont}{(5.2}
\texto{378.1528}{557.0390}{11.9552}{\normalfont}{\%}
\texto{391.8056}{557.0390}{11.9552}{\normalfont}{del}
\texto{410.6613}{557.0390}{11.9552}{\normalfont}{PIB)}
\texto{324.1320}{542.5930}{11.9552}{\normalfont}{Escasa}
\texto{363.2255}{542.5930}{11.9552}{\normalfont}{digitalización}
\texto{435.8486}{542.5930}{11.9552}{\normalfont}{de}
\texto{452.5404}{542.5930}{11.9552}{\normalfont}{la}
\texto{466.6428}{542.5930}{11.9552}{\normalfont}{oferta}
\texto{324.1320}{528.1470}{11.9552}{\normalfont}{local}
\texto{324.1320}{513.7010}{11.9552}{\normalfont}{Demanda}
\texto{380.5151}{513.7010}{11.9552}{\normalfont}{creciente}
\texto{432.8502}{513.7010}{11.9552}{\normalfont}{de}
\texto{452.3276}{513.7010}{11.9552}{\normalfont}{platafor-}
\texto{324.1320}{499.2560}{11.9552}{\normalfont}{mas}
\texto{348.2552}{499.2560}{11.9552}{\normalfont}{de}
\texto{363.8591}{499.2560}{11.9552}{\normalfont}{reserva}
\texto{403.2706}{499.2560}{11.9552}{\normalfont}{en}
\texto{418.8865}{499.2560}{11.9552}{\normalfont}{línea}
\texto{324.1320}{484.8100}{11.9552}{\normalfont}{Apoyo}
\texto{361.9583}{484.8100}{11.9552}{\normalfont}{gubernamental}
\texto{443.0480}{484.8100}{11.9552}{\normalfont}{al}
\texto{457.1384}{484.8100}{11.9552}{\normalfont}{turismo}
\texto{324.1320}{470.3640}{11.9552}{\normalfont}{sostenible}
\texto{324.1320}{455.9180}{11.9552}{\normalfont}{Integración}
\texto{385.1633}{455.9180}{11.9552}{\normalfont}{con}
\texto{406.5368}{455.9180}{11.9552}{\normalfont}{pasarelas}
\texto{456.6434}{455.9180}{11.9552}{\normalfont}{de}
\texto{472.1637}{455.9180}{11.9552}{\normalfont}{pago}
\texto{324.1320}{441.4720}{11.9552}{\normalfont}{internacionales}
\texto{324.1320}{427.0260}{11.9552}{\normalfont}{Posibilidad}
\texto{384.5129}{427.0260}{11.9552}{\normalfont}{de}
\texto{400.1169}{427.0260}{11.9552}{\normalfont}{expansión}
\texto{454.4819}{427.0260}{11.9552}{\normalfont}{regional}
\texto{115.7700}{391.4140}{11.9552}{\normalfont\bfseries}{Debilidades}
\texto{311.9780}{391.4140}{11.9552}{\normalfont\bfseries}{Amenazas}
\texto{127.9250}{368.0020}{11.9552}{\normalfont}{Prototipo}
\texto{182.6726}{368.0020}{11.9552}{\normalfont}{en}
\texto{199.8666}{368.0020}{11.9552}{\normalfont}{fase}
\texto{224.6163}{368.0020}{11.9552}{\normalfont}{inicial}
\texto{260.6708}{368.0020}{11.9552}{\normalfont}{(funcio-}
\texto{127.9250}{353.5560}{11.9552}{\normalfont}{nalidad}
\texto{169.5434}{353.5560}{11.9552}{\normalfont}{parcial)}
\texto{127.9250}{339.1100}{11.9552}{\normalfont}{Dependencia}
\texto{202.8028}{339.1100}{11.9552}{\normalfont}{de}
\texto{224.1931}{339.1100}{11.9552}{\normalfont}{infraestructura}
\texto{127.9250}{324.6640}{11.9552}{\normalfont}{tecnológica}
\texto{188.4040}{324.6640}{11.9552}{\normalfont}{del}
\texto{207.2597}{324.6640}{11.9552}{\normalfont}{país}
\texto{127.9250}{310.2190}{11.9552}{\normalfont}{Necesidad}
\texto{182.2064}{310.2190}{11.9552}{\normalfont}{de}
\texto{197.0930}{310.2190}{11.9552}{\normalfont}{capacitación}
\texto{263.3583}{310.2190}{11.9552}{\normalfont}{de}
\texto{278.2329}{310.2190}{11.9552}{\normalfont}{ope-}
\texto{127.9250}{295.7730}{11.9552}{\normalfont}{radores}
\texto{127.9250}{281.3270}{11.9552}{\normalfont}{Alcance}
\texto{170.6433}{281.3270}{11.9552}{\normalfont}{limitado}
\texto{215.6546}{281.3270}{11.9552}{\normalfont}{a}
\texto{224.2457}{281.3270}{11.9552}{\normalfont}{tours}
\texto{253.0625}{281.3270}{11.9552}{\normalfont}{turísticos}
\texto{127.9250}{266.8810}{11.9552}{\normalfont}{(sin}
\texto{150.7475}{266.8810}{11.9552}{\normalfont}{hospedaje}
\texto{205.1197}{266.8810}{11.9552}{\normalfont}{ni}
\texto{218.7726}{266.8810}{11.9552}{\normalfont}{transporte)}
\texto{127.9250}{252.4350}{11.9552}{\normalfont}{Tiempo}
\texto{170.2823}{252.4350}{11.9552}{\normalfont}{de}
\texto{185.0015}{252.4350}{11.9552}{\normalfont}{desarrollo}
\texto{237.5064}{252.4350}{11.9552}{\normalfont}{limitado}
\texto{282.7807}{252.4350}{11.9552}{\normalfont}{por}
\texto{127.9250}{237.9890}{11.9552}{\normalfont}{el}
\texto{140.2771}{237.9890}{11.9552}{\normalfont}{componente}
\texto{205.3086}{237.9890}{11.9552}{\normalfont}{académico}
\texto{324.1320}{368.7330}{11.9552}{\normalfont}{Competencia}
\texto{397.4030}{368.7330}{11.9552}{\normalfont}{de}
\texto{415.7327}{368.7330}{11.9552}{\normalfont}{plataformas}
\texto{482.5671}{368.7330}{11.9552}{\normalfont}{in-}
\texto{324.1320}{354.2870}{11.9552}{\normalfont}{ternacionales}
\texto{394.4310}{354.2870}{11.9552}{\normalfont}{consolidadas}
\texto{324.1320}{339.8410}{11.9552}{\normalfont}{Cambios}
\texto{371.0753}{339.8410}{11.9552}{\normalfont}{en}
\texto{385.7587}{339.8410}{11.9552}{\normalfont}{la}
\texto{397.8525}{339.8410}{11.9552}{\normalfont}{regulación}
\texto{452.8584}{339.8410}{11.9552}{\normalfont}{de}
\texto{467.5418}{339.8410}{11.9552}{\normalfont}{pagos}
\texto{324.1320}{325.3950}{11.9552}{\normalfont}{digitales}
\texto{324.1320}{310.9490}{11.9552}{\normalfont}{Riesgos}
\texto{366.6256}{310.9490}{11.9552}{\normalfont}{de}
\texto{382.8153}{310.9490}{11.9552}{\normalfont}{ciberseguridad}
\texto{461.1768}{310.9490}{11.9552}{\normalfont}{y}
\texto{471.8385}{310.9490}{11.9552}{\normalfont}{frau-}
\texto{324.1320}{296.5040}{11.9552}{\normalfont}{de}
\texto{324.1320}{282.0580}{11.9552}{\normalfont}{Fluctuaciones}
\texto{399.2680}{282.0580}{11.9552}{\normalfont}{en}
\texto{415.9838}{282.0580}{11.9552}{\normalfont}{la}
\texto{430.0862}{282.0580}{11.9552}{\normalfont}{demanda}
\texto{481.2711}{282.0580}{11.9552}{\normalfont}{tu-}
\texto{324.1320}{267.6120}{11.9552}{\normalfont}{rística}
\texto{324.1320}{253.1660}{11.9552}{\normalfont}{Barreras}
\texto{372.4358}{253.1660}{11.9552}{\normalfont}{de}
\texto{388.9722}{253.1660}{11.9552}{\normalfont}{adopción}
\texto{439.6503}{253.1660}{11.9552}{\normalfont}{tecnológica}
\texto{324.1320}{238.7200}{11.9552}{\normalfont}{en}
\texto{339.7359}{238.7200}{11.9552}{\normalfont}{operadores}
\texto{398.6512}{238.7200}{11.9552}{\normalfont}{tradicionales}
\texto{88.7990}{202.2670}{11.9552}{\normalfont}{Fuente:}
\texto{130.2501}{202.2670}{11.9552}{\normalfont}{Elaboración}
\texto{195.1262}{202.2670}{11.9552}{\normalfont}{propia}
\texto{231.5417}{202.2670}{11.9552}{\normalfont}{(2026).}
\end{pagina}
% Página original 44
\begin{pagina}{44}
\texto{70.8660}{746.3890}{11.9552}{\normalfont}{34}
\texto{288.7940}{746.3890}{11.9552}{\normalfont}{CAPÍTULO}
\texto{355.7575}{746.3890}{11.9552}{\normalfont}{3.}
\texto{372.6693}{746.3890}{11.9552}{\normalfont}{EMPRENDIMIENTO}
\texto{490.9875}{746.3890}{11.9552}{\normalfont}{DIGITAL}
\texto{70.8660}{705.1940}{17.2154}{\normalfont\bfseries}{3.11.}
\texto{130.0439}{705.1940}{17.2154}{\normalfont\bfseries}{Lean}
\texto{177.1470}{705.1940}{17.2154}{\normalfont\bfseries}{Canv}
\texto{219.9307}{705.1940}{17.2154}{\normalfont\bfseries}{as}
\texto{88.7990}{678.9090}{11.9552}{\normalfont}{El}
\texto{103.7693}{678.9090}{11.9552}{\normalfont}{modelo}
\texto{144.2700}{678.9090}{11.9552}{\normalfont\itshape}{Lean}
\texto{172.3815}{678.9090}{11.9552}{\normalfont\itshape}{Canvas}
\texto{213.6760}{678.9090}{11.9552}{\normalfont}{[13]}
\texto{235.6397}{678.9090}{11.9552}{\normalfont}{de}
\texto{251.1020}{678.9090}{11.9552}{\normalfont\bfseries}{Explora-Nicaragua}
\texto{365.0250}{678.9090}{11.9552}{\normalfont}{estructura}
\texto{420.8725}{678.9090}{11.9552}{\normalfont}{la}
\texto{433.7315}{678.9090}{11.9552}{\normalfont}{propuesta}
\texto{487.6280}{678.9090}{11.9552}{\normalfont}{de}
\texto{503.0884}{678.9090}{11.9552}{\normalfont}{negocio}
\texto{70.8660}{664.4630}{11.9552}{\normalfont}{en}
\texto{86.4699}{664.4630}{11.9552}{\normalfont}{los}
\texto{104.0895}{664.4630}{11.9552}{\normalfont}{siguientes}
\texto{157.2208}{664.4630}{11.9552}{\normalfont}{bloques:}
\texto{183.8710}{629.9920}{11.9552}{\normalfont}{Cuadro}
\texto{225.4894}{629.9920}{11.9552}{\normalfont}{3.2:}
\texto{247.6040}{629.9920}{11.9552}{\normalfont}{Lean}
\texto{276.3754}{629.9920}{11.9552}{\normalfont}{Canvas}
\texto{316.7648}{629.9920}{11.9552}{\normalfont}{de}
\texto{332.3688}{629.9920}{11.9552}{\normalfont}{Explora-Nicaragua}
\texto{76.8440}{612.2770}{11.9552}{\normalfont\bfseries}{Problema}
\texto{188.0110}{612.2770}{11.9552}{\normalfont\bfseries}{Solución}
\texto{329.7440}{612.2770}{11.9552}{\normalfont\bfseries}{Métricas}
\texto{385.3907}{612.2770}{11.9552}{\normalfont\bfseries}{clave}
\texto{440.9046}{612.2770}{11.9552}{\normalfont\bfseries}{V}
\texto{449.9488}{612.2770}{11.9552}{\normalfont\bfseries}{entaja}
\texto{490.4458}{612.2770}{11.9552}{\normalfont\bfseries}{injusta}
\texto{88.9980}{583.6060}{11.9552}{\normalfont}{Gestión}
\texto{138.6611}{583.6060}{11.9552}{\normalfont}{manual}
\texto{88.9980}{569.1600}{11.9552}{\normalfont}{de}
\texto{104.6019}{569.1600}{11.9552}{\normalfont}{tours}
\texto{88.9980}{554.7140}{11.9552}{\normalfont}{Comisiones}
\texto{88.9980}{540.2680}{11.9552}{\normalfont}{elevadas}
\texto{142.1939}{540.2680}{11.9552}{\normalfont}{(15}
\texto{160.4494}{540.2680}{11.9552}{\normalfont}{\%--}
\texto{88.9980}{525.8230}{11.9552}{\normalfont}{25}
\texto{102.6891}{525.8230}{11.9552}{\normalfont}{\%)}
\texto{88.9980}{511.3770}{11.9552}{\normalfont}{Falta}
\texto{119.5914}{511.3770}{11.9552}{\normalfont}{de}
\texto{135.7333}{511.3770}{11.9552}{\normalfont}{confian-}
\texto{88.9980}{496.9310}{11.9552}{\normalfont}{za}
\texto{103.9516}{496.9310}{11.9552}{\normalfont}{del}
\texto{122.8073}{496.9310}{11.9552}{\normalfont}{turista}
\texto{200.1660}{583.6060}{11.9552}{\normalfont}{Marketplace}
\texto{266.4743}{583.6060}{11.9552}{\normalfont}{con}
\texto{287.5968}{583.6060}{11.9552}{\normalfont}{comi-}
\texto{200.1660}{569.1600}{11.9552}{\normalfont}{sión}
\texto{224.2892}{569.1600}{11.9552}{\normalfont}{del}
\texto{243.1449}{569.1600}{11.9552}{\normalfont}{8}
\texto{250.9947}{569.1600}{11.9552}{\normalfont}{\%}
\texto{200.1660}{554.7140}{11.9552}{\normalfont}{Backoffice}
\texto{254.0122}{554.7140}{11.9552}{\normalfont}{de}
\texto{268.6836}{554.7140}{11.9552}{\normalfont}{escritorio}
\texto{200.1660}{540.2680}{11.9552}{\normalfont}{Reseñas}
\texto{244.6752}{540.2680}{11.9552}{\normalfont}{verificadas}
\texto{341.8980}{583.6060}{11.9552}{\normalfont}{Tasa}
\texto{373.0963}{583.6060}{11.9552}{\normalfont}{de}
\texto{392.2031}{583.6060}{11.9552}{\normalfont}{conver-}
\texto{341.8980}{569.1600}{11.9552}{\normalfont}{sión}
\texto{341.8980}{554.7140}{11.9552}{\normalfont}{NPS}
\texto{341.8980}{540.2680}{11.9552}{\normalfont}{Churn}
\texto{377.9907}{540.2680}{11.9552}{\normalfont}{rate}
\texto{453.0660}{583.6060}{11.9552}{\normalfont}{Primera}
\texto{510.2071}{583.6060}{11.9552}{\normalfont}{plata-}
\texto{453.0660}{569.1600}{11.9552}{\normalfont}{forma}
\texto{505.0018}{569.1600}{11.9552}{\normalfont}{nicara-}
\texto{453.0660}{554.7140}{11.9552}{\normalfont}{güense}
\texto{490.8468}{554.7140}{11.9552}{\normalfont}{integrada}
\texto{453.0660}{540.2680}{11.9552}{\normalfont}{Modelo}
\texto{494.9642}{540.2680}{11.9552}{\normalfont}{de}
\texto{510.8550}{540.2680}{11.9552}{\normalfont}{verifi-}
\texto{453.0660}{525.8230}{11.9552}{\normalfont}{cación}
\texto{488.8312}{525.8230}{11.9552}{\normalfont}{única}
\texto{76.8440}{476.4950}{11.9552}{\normalfont\bfseries}{Propuesta}
\texto{139.9675}{476.4950}{11.9552}{\normalfont\bfseries}{de}
\texto{157.1387}{476.4950}{11.9552}{\normalfont\bfseries}{valor:}
\texto{192.8320}{476.4950}{11.9552}{\normalfont}{Marketplace}
\texto{258.6860}{476.4950}{11.9552}{\normalfont}{turístico}
\texto{304.1325}{476.4950}{11.9552}{\normalfont}{ni-}
\texto{76.8440}{462.0490}{11.9552}{\normalfont}{caragüense}
\texto{136.3737}{462.0490}{11.9552}{\normalfont}{con}
\texto{158.1298}{462.0490}{11.9552}{\normalfont}{comisión}
\texto{206.6033}{462.0490}{11.9552}{\normalfont}{reducida}
\texto{254.3739}{462.0490}{11.9552}{\normalfont}{(8}
\texto{266.7763}{462.0490}{11.9552}{\normalfont}{\%),}
\texto{288.5204}{462.0490}{11.9552}{\normalfont}{verifi-}
\texto{76.8440}{447.6040}{11.9552}{\normalfont}{cación}
\texto{112.6092}{447.6040}{11.9552}{\normalfont}{de}
\texto{128.2131}{447.6040}{11.9552}{\normalfont}{operadores}
\texto{187.1403}{447.6040}{11.9552}{\normalfont}{y}
\texto{197.2161}{447.6040}{11.9552}{\normalfont}{reseñas}
\texto{237.6749}{447.6040}{11.9552}{\normalfont}{autenticadas}
\texto{329.7440}{476.4950}{11.9552}{\normalfont\bfseries}{Segmentos}
\texto{395.2513}{476.4950}{11.9552}{\normalfont\bfseries}{de}
\texto{412.3269}{476.4950}{11.9552}{\normalfont\bfseries}{clientes:}
\texto{462.4500}{476.4950}{11.9552}{\normalfont}{Turistas,}
\texto{510.1489}{476.4950}{11.9552}{\normalfont}{opera-}
\texto{329.7440}{462.0490}{11.9552}{\normalfont}{dores}
\texto{360.3708}{462.0490}{11.9552}{\normalfont}{turísticos,}
\texto{414.4753}{462.0490}{11.9552}{\normalfont}{equipo}
\texto{452.2011}{462.0490}{11.9552}{\normalfont}{gestor}
\texto{76.8440}{427.1680}{11.9552}{\normalfont\bfseries}{Canales:}
\texto{130.2510}{427.1680}{11.9552}{\normalfont}{Prototipo}
\texto{183.7433}{427.1680}{11.9552}{\normalfont}{inicial:}
\texto{221.7944}{427.1680}{11.9552}{\normalfont}{backoffice}
\texto{275.1217}{427.1680}{11.9552}{\normalfont}{de}
\texto{291.0604}{427.1680}{11.9552}{\normalfont}{escri-}
\texto{76.8440}{412.7220}{11.9552}{\normalfont}{torio}
\texto{105.9286}{412.7220}{11.9552}{\normalfont}{(operador)}
\texto{165.2647}{412.7220}{11.9552}{\normalfont}{y}
\texto{176.4643}{412.7220}{11.9552}{\normalfont}{panel}
\texto{208.8127}{412.7220}{11.9552}{\normalfont}{de}
\texto{225.5523}{412.7220}{11.9552}{\normalfont}{administración}
\texto{306.0802}{412.7220}{11.9552}{\normalfont}{de}
\texto{76.8440}{398.2760}{11.9552}{\normalfont}{escritorio}
\texto{127.9907}{398.2760}{11.9552}{\normalfont}{(admin).}
\texto{176.4834}{398.2760}{11.9552}{\normalfont}{Visión}
\texto{213.0042}{398.2760}{11.9552}{\normalfont}{completa:}
\texto{266.6997}{398.2760}{11.9552}{\normalfont}{app}
\texto{289.8283}{398.2760}{11.9552}{\normalfont}{móvil}
\texto{76.8440}{383.8300}{11.9552}{\normalfont}{(turista)}
\texto{329.7440}{427.1680}{11.9552}{\normalfont\bfseries}{Estructura}
\texto{401.6796}{427.1680}{11.9552}{\normalfont\bfseries}{de}
\texto{424.1949}{427.1680}{11.9552}{\normalfont\bfseries}{costos:}
\texto{470.9610}{427.1680}{11.9552}{\normalfont}{Desarrollo}
\texto{530.6366}{427.1680}{11.9552}{\normalfont}{de}
\texto{329.7440}{412.7220}{11.9552}{\normalfont}{software,}
\texto{381.4215}{412.7220}{11.9552}{\normalfont}{infraestructura}
\texto{463.6733}{412.7220}{11.9552}{\normalfont}{cloud,}
\texto{500.6603}{412.7220}{11.9552}{\normalfont}{comisio-}
\texto{329.7440}{398.2760}{11.9552}{\normalfont}{nes}
\texto{349.9650}{398.2760}{11.9552}{\normalfont}{de}
\texto{365.5690}{398.2760}{11.9552}{\normalfont}{pasarela}
\texto{411.1661}{398.2760}{11.9552}{\normalfont}{de}
\texto{426.7700}{398.2760}{11.9552}{\normalfont}{pago}
\texto{76.8440}{363.3940}{11.9552}{\normalfont\bfseries}{Flujo}
\texto{112.8770}{363.3940}{11.9552}{\normalfont\bfseries}{de}
\texto{132.0328}{363.3940}{11.9552}{\normalfont\bfseries}{ingresos:}
\texto{187.4990}{363.3940}{11.9552}{\normalfont}{Comisión}
\texto{239.8580}{363.3940}{11.9552}{\normalfont}{del}
\texto{259.6462}{363.3940}{11.9552}{\normalfont}{8}
\texto{267.4841}{363.3940}{11.9552}{\normalfont}{\%}
\texto{282.0694}{363.3940}{11.9552}{\normalfont}{por}
\texto{304.1315}{363.3940}{11.9552}{\normalfont}{re-}
\texto{76.8440}{348.9490}{11.9552}{\normalfont}{serva}
\texto{106.4881}{348.9490}{11.9552}{\normalfont}{+}
\texto{119.5025}{348.9490}{11.9552}{\normalfont}{suscripción}
\texto{179.7831}{348.9490}{11.9552}{\normalfont}{Premium}
\texto{230.6787}{348.9490}{11.9552}{\normalfont}{opcional}
\texto{329.7440}{363.3940}{11.9552}{\normalfont\bfseries}{Estructura}
\texto{397.8898}{363.3940}{11.9552}{\normalfont\bfseries}{de}
\texto{416.6153}{363.3940}{11.9552}{\normalfont\bfseries}{ingresos:}
\texto{471.7160}{363.3940}{11.9552}{\normalfont}{Transacciona-}
\texto{329.7440}{348.9490}{11.9552}{\normalfont}{les}
\texto{346.7132}{348.9490}{11.9552}{\normalfont}{(comisión)}
\texto{404.0049}{348.9490}{11.9552}{\normalfont}{y}
\texto{414.0927}{348.9490}{11.9552}{\normalfont}{recurrente}
\texto{469.6964}{348.9490}{11.9552}{\normalfont}{(suscripción)}
\texto{88.7990}{315.5510}{11.9552}{\normalfont}{Fuente:}
\texto{130.2501}{315.5510}{11.9552}{\normalfont}{Elaboración}
\texto{195.1262}{315.5510}{11.9552}{\normalfont}{propia}
\texto{231.5417}{315.5510}{11.9552}{\normalfont}{con}
\texto{252.9989}{315.5510}{11.9552}{\normalfont}{base}
\texto{279.0851}{315.5510}{11.9552}{\normalfont}{en}
\texto{294.6891}{315.5510}{11.9552}{\normalfont}{Maurya}
\texto{337.9310}{315.5510}{11.9552}{\normalfont}{(2012).}
\end{pagina}
% Página original 45
\begin{pagina}{45}
\texto{70.8660}{746.3890}{11.9552}{\normalfont}{3.12.}
\texto{102.7242}{746.3890}{11.9552}{\normalfont}{CONSIDERACIONES}
\texto{222.6803}{746.3890}{11.9552}{\normalfont}{ÉTICAS,}
\texto{274.2168}{746.3890}{11.9552}{\normalfont}{SOCIALES}
\texto{336.9577}{746.3890}{11.9552}{\normalfont}{E}
\texto{348.8316}{746.3890}{11.9552}{\normalfont}{IDENTIDAD}
\texto{421.6411}{746.3890}{11.9552}{\normalfont}{NACIONAL}
\texto{529.4280}{746.3890}{11.9552}{\normalfont}{35}
\texto{70.8660}{705.1940}{17.2154}{\normalfont\bfseries}{3.12.}
\texto{130.0439}{705.1940}{17.2154}{\normalfont\bfseries}{Consideraciones}
\texto{274.3968}{705.1940}{17.2154}{\normalfont\bfseries}{éticas,}
\texto{336.2603}{705.1940}{17.2154}{\normalfont\bfseries}{sociales}
\texto{408.4669}{705.1940}{17.2154}{\normalfont\bfseries}{e}
\texto{426.3761}{705.1940}{17.2154}{\normalfont\bfseries}{identidad}
\texto{514.4965}{705.1940}{17.2154}{\normalfont\bfseries}{na-}
\texto{130.0440}{683.2760}{17.2154}{\normalfont\bfseries}{cional}
\texto{88.7990}{656.9910}{11.9552}{\normalfont}{El}
\texto{103.1118}{656.9910}{11.9552}{\normalfont}{desarrollo}
\texto{155.7003}{656.9910}{11.9552}{\normalfont}{de}
\texto{170.5090}{656.9910}{11.9552}{\normalfont\bfseries}{Explora-Nicaragua}
\texto{283.7770}{656.9910}{11.9552}{\normalfont}{considera}
\texto{334.4144}{656.9910}{11.9552}{\normalfont}{las}
\texto{351.2330}{656.9910}{11.9552}{\normalfont}{siguientes}
\texto{403.5633}{656.9910}{11.9552}{\normalfont}{dimensiones}
\texto{467.9229}{656.9910}{11.9552}{\normalfont}{éticas,}
\texto{502.9636}{656.9910}{11.9552}{\normalfont}{sociales}
\texto{70.8660}{642.5450}{11.9552}{\normalfont}{y}
\texto{80.9418}{642.5450}{11.9552}{\normalfont}{de}
\texto{96.5458}{642.5450}{11.9552}{\normalfont}{identidad}
\texto{148.2472}{642.5450}{11.9552}{\normalfont}{nacional:}
\texto{100.7540}{615.1480}{11.9552}{\normalfont\bfseries}{Ética}
\texto{136.3482}{615.1480}{11.9552}{\normalfont\bfseries}{de}
\texto{155.2530}{615.1480}{11.9552}{\normalfont\bfseries}{datos}
\texto{191.8084}{615.1480}{11.9552}{\normalfont\bfseries}{y}
\texto{204.2035}{615.1480}{11.9552}{\normalfont\bfseries}{privacidad:}
\texto{273.3400}{615.1480}{11.9552}{\normalfont}{el}
\texto{286.3975}{615.1480}{11.9552}{\normalfont}{ecosistema}
\texto{345.1095}{615.1480}{11.9552}{\normalfont}{maneja}
\texto{386.4577}{615.1480}{11.9552}{\normalfont}{datos}
\texto{418.4522}{615.1480}{11.9552}{\normalfont}{personales}
\texto{475.5359}{615.1480}{11.9552}{\normalfont}{de}
\texto{491.8452}{615.1480}{11.9552}{\normalfont}{turistas}
\texto{534.9604}{615.1480}{11.9552}{\normalfont}{y}
\texto{100.7540}{600.7020}{11.9552}{\normalfont}{operadores}
\texto{159.1432}{600.7020}{11.9552}{\normalfont}{(nombres,}
\texto{212.9846}{600.7020}{11.9552}{\normalfont}{correos,}
\texto{255.4423}{600.7020}{11.9552}{\normalfont}{datos}
\texto{286.1935}{600.7020}{11.9552}{\normalfont}{de}
\texto{301.2714}{600.7020}{11.9552}{\normalfont}{pago).}
\texto{336.5225}{600.7020}{11.9552}{\normalfont}{Se}
\texto{351.6004}{600.7020}{11.9552}{\normalfont}{aplican}
\texto{391.3921}{600.7020}{11.9552}{\normalfont}{protocolos}
\texto{447.1798}{600.7020}{11.9552}{\normalfont}{de}
\texto{462.2577}{600.7020}{11.9552}{\normalfont}{seguridad}
\texto{514.4709}{600.7020}{11.9552}{\normalfont}{como}
\texto{100.7540}{586.2560}{11.9552}{\normalfont}{SSL/TLS}
\texto{152.9671}{586.2560}{11.9552}{\normalfont}{para}
\texto{179.5077}{586.2560}{11.9552}{\normalfont}{el}
\texto{191.7283}{586.2560}{11.9552}{\normalfont}{cifrado}
\texto{230.3005}{586.2560}{11.9552}{\normalfont}{de}
\texto{245.7849}{586.2560}{11.9552}{\normalfont}{comunicaciones}
\texto{327.9865}{586.2560}{11.9552}{\normalfont}{y}
\texto{337.9428}{586.2560}{11.9552}{\normalfont}{JWT}
\texto{368.2038}{586.2560}{11.9552}{\normalfont}{para}
\texto{394.7443}{586.2560}{11.9552}{\normalfont}{la}
\texto{407.6153}{586.2560}{11.9552}{\normalfont}{autenticación,}
\texto{482.6103}{586.2560}{11.9552}{\normalfont}{cumpliendo}
\texto{100.7540}{571.8100}{11.9552}{\normalfont}{con}
\texto{121.3863}{571.8100}{11.9552}{\normalfont}{las}
\texto{138.1809}{571.8100}{11.9552}{\normalfont}{buenas}
\texto{176.4376}{571.8100}{11.9552}{\normalfont}{prácticas}
\texto{225.1000}{571.8100}{11.9552}{\normalfont}{de}
\texto{239.8791}{571.8100}{11.9552}{\normalfont}{protección}
\texto{295.6309}{571.8100}{11.9552}{\normalfont}{de}
\texto{310.4100}{571.8100}{11.9552}{\normalfont}{datos}
\texto{340.8622}{571.8100}{11.9552}{\normalfont}{[7].}
\texto{359.5401}{571.8100}{11.9552}{\normalfont}{Los}
\texto{380.3972}{571.8100}{11.9552}{\normalfont}{datos}
\texto{410.8494}{571.8100}{11.9552}{\normalfont}{de}
\texto{425.6285}{571.8100}{11.9552}{\normalfont}{pago}
\texto{452.7644}{571.8100}{11.9552}{\normalfont}{se}
\texto{465.6569}{571.8100}{11.9552}{\normalfont}{procesan}
\texto{513.3414}{571.8100}{11.9552}{\normalfont}{a}
\texto{522.2671}{571.8100}{11.9552}{\normalfont}{tra-}
\texto{100.7540}{557.3650}{11.9552}{\normalfont}{vés}
\texto{119.6815}{557.3650}{11.9552}{\normalfont}{de}
\texto{134.6398}{557.3650}{11.9552}{\normalfont}{Stripe,}
\texto{172.0333}{557.3650}{11.9552}{\normalfont}{que}
\texto{193.1820}{557.3650}{11.9552}{\normalfont}{cumple}
\texto{232.8542}{557.3650}{11.9552}{\normalfont}{con}
\texto{253.6658}{557.3650}{11.9552}{\normalfont}{la}
\texto{266.0227}{557.3650}{11.9552}{\normalfont}{certificación}
\texto{330.4206}{557.3650}{11.9552}{\normalfont}{PCI}
\texto{354.3142}{557.3650}{11.9552}{\normalfont}{DSS,}
\texto{382.7652}{557.3650}{11.9552}{\normalfont}{garantizando}
\texto{452.0313}{557.3650}{11.9552}{\normalfont}{que}
\texto{473.1680}{557.3650}{11.9552}{\normalfont}{la}
\texto{485.5369}{557.3650}{11.9552}{\normalfont}{plataforma}
\texto{100.7540}{542.9190}{11.9552}{\normalfont}{no}
\texto{117.0083}{542.9190}{11.9552}{\normalfont}{almacena}
\texto{168.3822}{542.9190}{11.9552}{\normalfont}{datos}
\texto{199.6713}{542.9190}{11.9552}{\normalfont}{sensibles}
\texto{247.3367}{542.9190}{11.9552}{\normalfont}{de}
\texto{262.9406}{542.9190}{11.9552}{\normalfont}{tarjetas}
\texto{305.6111}{542.9190}{11.9552}{\normalfont}{de}
\texto{321.2151}{542.9190}{11.9552}{\normalfont}{crédito.}
\texto{100.7540}{518.5100}{11.9552}{\normalfont\bfseries}{Inclusión}
\texto{157.6524}{518.5100}{11.9552}{\normalfont\bfseries}{y}
\texto{168.7444}{518.5100}{11.9552}{\normalfont\bfseries}{accesibilidad:}
\texto{249.7970}{518.5100}{11.9552}{\normalfont}{el}
\texto{261.7187}{518.5100}{11.9552}{\normalfont}{backoffice}
\texto{314.3049}{518.5100}{11.9552}{\normalfont}{de}
\texto{329.4904}{518.5100}{11.9552}{\normalfont}{escritorio}
\texto{379.8600}{518.5100}{11.9552}{\normalfont}{ofrece}
\texto{412.9305}{518.5100}{11.9552}{\normalfont}{una}
\texto{435.2580}{518.5100}{11.9552}{\normalfont}{interfaz}
\texto{477.1108}{518.5100}{11.9552}{\normalfont}{accesible}
\texto{524.8790}{518.5100}{11.9552}{\normalfont}{pa-}
\texto{100.7540}{504.0650}{11.9552}{\normalfont}{ra}
\texto{115.9658}{504.0650}{11.9552}{\normalfont}{operadores}
\texto{175.8016}{504.0650}{11.9552}{\normalfont}{que}
\texto{198.5045}{504.0650}{11.9552}{\normalfont}{pueden}
\texto{239.7428}{504.0650}{11.9552}{\normalfont}{no}
\texto{256.9176}{504.0650}{11.9552}{\normalfont}{tener}
\texto{287.7501}{504.0650}{11.9552}{\normalfont}{experiencia}
\texto{349.7976}{504.0650}{11.9552}{\normalfont}{avanzada}
\texto{401.4368}{504.0650}{11.9552}{\normalfont}{con}
\texto{423.8146}{504.0650}{11.9552}{\normalfont}{tecnología,}
\texto{483.2630}{504.0650}{11.9552}{\normalfont}{priorizando}
\texto{100.7540}{489.6190}{11.9552}{\normalfont}{la}
\texto{114.5097}{489.6190}{11.9552}{\normalfont}{usabilidad}
\texto{171.2538}{489.6190}{11.9552}{\normalfont}{y}
\texto{182.0828}{489.6190}{11.9552}{\normalfont}{la}
\texto{195.8385}{489.6190}{11.9552}{\normalfont}{productividad.}
\texto{275.2904}{489.6190}{11.9552}{\normalfont}{La}
\texto{293.1084}{489.6190}{11.9552}{\normalfont}{app}
\texto{316.6195}{489.6190}{11.9552}{\normalfont}{móvil}
\texto{349.2380}{489.6190}{11.9552}{\normalfont}{se}
\texto{363.7086}{489.6190}{11.9552}{\normalfont}{diseña}
\texto{400.2915}{489.6190}{11.9552}{\normalfont}{con}
\texto{422.5019}{489.6190}{11.9552}{\normalfont}{interfaz}
\texto{465.5263}{489.6190}{11.9552}{\normalfont}{bilingüe}
\texto{510.4994}{489.6190}{11.9552}{\normalfont}{(espa-}
\texto{100.7540}{475.1730}{11.9552}{\normalfont}{ñol/inglés)}
\texto{159.3464}{475.1730}{11.9552}{\normalfont}{para}
\texto{186.0065}{475.1730}{11.9552}{\normalfont}{atender}
\texto{228.2873}{475.1730}{11.9552}{\normalfont}{tanto}
\texto{259.1771}{475.1730}{11.9552}{\normalfont}{a}
\texto{268.9278}{475.1730}{11.9552}{\normalfont}{turistas}
\texto{311.3377}{475.1730}{11.9552}{\normalfont}{locales}
\texto{348.7909}{475.1730}{11.9552}{\normalfont}{como}
\texto{379.3652}{475.1730}{11.9552}{\normalfont}{internacionales.}
\texto{100.7540}{450.7640}{11.9552}{\normalfont\bfseries}{Impacto}
\texto{152.5607}{450.7640}{11.9552}{\normalfont\bfseries}{económico}
\texto{216.9012}{450.7640}{11.9552}{\normalfont\bfseries}{local:}
\texto{250.9160}{450.7640}{11.9552}{\normalfont}{al}
\texto{263.2012}{450.7640}{11.9552}{\normalfont}{reducir}
\texto{302.1512}{450.7640}{11.9552}{\normalfont}{las}
\texto{319.0535}{450.7640}{11.9552}{\normalfont}{comisiones}
\texto{376.3428}{450.7640}{11.9552}{\normalfont}{frente}
\texto{408.7916}{450.7640}{11.9552}{\normalfont}{a}
\texto{417.8369}{450.7640}{11.9552}{\normalfont}{plataformas}
\texto{481.2401}{450.7640}{11.9552}{\normalfont}{extranjeras,}
\texto{100.7540}{436.3190}{11.9552}{\normalfont}{el}
\texto{112.9866}{436.3190}{11.9552}{\normalfont}{dinero}
\texto{148.6322}{436.3190}{11.9552}{\normalfont}{de}
\texto{164.1166}{436.3190}{11.9552}{\normalfont}{las}
\texto{181.6166}{436.3190}{11.9552}{\normalfont}{transacciones}
\texto{253.1613}{436.3190}{11.9552}{\normalfont}{turísticas}
\texto{303.8944}{436.3190}{11.9552}{\normalfont}{se}
\texto{317.5042}{436.3190}{11.9552}{\normalfont}{retiene}
\texto{355.7512}{436.3190}{11.9552}{\normalfont}{en}
\texto{371.2356}{436.3190}{11.9552}{\normalfont}{mayor}
\texto{406.5608}{436.3190}{11.9552}{\normalfont}{proporción}
\texto{465.2920}{436.3190}{11.9552}{\normalfont}{en}
\texto{480.7763}{436.3190}{11.9552}{\normalfont}{la}
\texto{493.6593}{436.3190}{11.9552}{\normalfont}{economía}
\texto{100.7540}{421.8730}{11.9552}{\normalfont}{nicaragüense,}
\texto{173.3268}{421.8730}{11.9552}{\normalfont}{fortaleciendo}
\texto{242.9085}{421.8730}{11.9552}{\normalfont}{a}
\texto{252.9939}{421.8730}{11.9552}{\normalfont}{los}
\texto{270.9363}{421.8730}{11.9552}{\normalfont}{operadores}
\texto{330.1743}{421.8730}{11.9552}{\normalfont}{locales}
\texto{367.9623}{421.8730}{11.9552}{\normalfont}{y}
\texto{378.3609}{421.8730}{11.9552}{\normalfont}{contribuyendo}
\texto{455.1013}{421.8730}{11.9552}{\normalfont}{al}
\texto{468.4386}{421.8730}{11.9552}{\normalfont}{desarrollo}
\texto{522.1509}{421.8730}{11.9552}{\normalfont}{sos-}
\texto{100.7540}{407.4270}{11.9552}{\normalfont}{tenible}
\texto{139.1206}{407.4270}{11.9552}{\normalfont}{del}
\texto{157.9764}{407.4270}{11.9552}{\normalfont}{sector.}
\texto{100.7540}{383.0180}{11.9552}{\normalfont\bfseries}{Promoción}
\texto{167.6732}{383.0180}{11.9552}{\normalfont\bfseries}{de}
\texto{184.5815}{383.0180}{11.9552}{\normalfont\bfseries}{la}
\texto{198.1674}{383.0180}{11.9552}{\normalfont\bfseries}{identidad}
\texto{256.3617}{383.0180}{11.9552}{\normalfont\bfseries}{nacional:}
\texto{311.1720}{383.0180}{11.9552}{\normalfont}{la}
\texto{323.1463}{383.0180}{11.9552}{\normalfont}{plataforma}
\texto{381.6216}{383.0180}{11.9552}{\normalfont}{promueve}
\texto{433.5980}{383.0180}{11.9552}{\normalfont}{los}
\texto{450.2014}{383.0180}{11.9552}{\normalfont}{destinos}
\texto{494.1726}{383.0180}{11.9552}{\normalfont}{turísticos}
\texto{100.7540}{368.5730}{11.9552}{\normalfont}{de}
\texto{116.3699}{368.5730}{11.9552}{\normalfont}{Nicaragua}
\texto{171.9903}{368.5730}{11.9552}{\normalfont}{(Ometepe,}
\texto{229.5641}{368.5730}{11.9552}{\normalfont}{Somoto,}
\texto{275.1086}{368.5730}{11.9552}{\normalfont}{San}
\texto{297.8785}{368.5730}{11.9552}{\normalfont}{Juan}
\texto{326.6738}{368.5730}{11.9552}{\normalfont}{del}
\texto{345.5535}{368.5730}{11.9552}{\normalfont}{Sur,}
\texto{370.2744}{368.5730}{11.9552}{\normalfont}{Solentiname,}
\texto{439.5596}{368.5730}{11.9552}{\normalfont}{Río}
\texto{461.2009}{368.5730}{11.9552}{\normalfont}{San}
\texto{483.9708}{368.5730}{11.9552}{\normalfont}{Juan)}
\texto{517.3186}{368.5730}{11.9552}{\normalfont}{y}
\texto{527.4184}{368.5730}{11.9552}{\normalfont}{las}
\texto{100.7540}{354.1270}{11.9552}{\normalfont}{experiencias}
\texto{165.7329}{354.1270}{11.9552}{\normalfont}{auténticas}
\texto{220.6360}{354.1270}{11.9552}{\normalfont}{que}
\texto{241.6532}{354.1270}{11.9552}{\normalfont}{ofrece}
\texto{274.3889}{354.1270}{11.9552}{\normalfont}{el}
\texto{285.9759}{354.1270}{11.9552}{\normalfont}{país,}
\texto{312.5858}{354.1270}{11.9552}{\normalfont}{contribuyendo}
\texto{388.2383}{354.1270}{11.9552}{\normalfont}{a}
\texto{397.2238}{354.1270}{11.9552}{\normalfont}{la}
\texto{409.4731}{354.1270}{11.9552}{\normalfont}{proyección}
\texto{466.2532}{354.1270}{11.9552}{\normalfont}{de}
\texto{481.1039}{354.1270}{11.9552}{\normalfont}{la}
\texto{493.3412}{354.1270}{11.9552}{\normalfont}{identidad}
\texto{100.7540}{339.6810}{11.9552}{\normalfont}{cultural}
\texto{144.3235}{339.6810}{11.9552}{\normalfont}{y}
\texto{154.3994}{339.6810}{11.9552}{\normalfont}{natural}
\texto{195.3794}{339.6810}{11.9552}{\normalfont}{nicaragüense}
\texto{264.3776}{339.6810}{11.9552}{\normalfont}{ante}
\texto{290.0646}{339.6810}{11.9552}{\normalfont}{el}
\texto{302.4287}{339.6810}{11.9552}{\normalfont}{turismo}
\texto{345.4124}{339.6810}{11.9552}{\normalfont}{internacional.}
\texto{100.7540}{315.2720}{11.9552}{\normalfont\bfseries}{F}
\texto{108.0969}{315.2720}{11.9552}{\normalfont\bfseries}{ormalización}
\texto{187.3156}{315.2720}{11.9552}{\normalfont\bfseries}{del}
\texto{210.0401}{315.2720}{11.9552}{\normalfont\bfseries}{sector:}
\texto{253.3280}{315.2720}{11.9552}{\normalfont}{el}
\texto{266.4572}{315.2720}{11.9552}{\normalfont}{proceso}
\texto{309.2401}{315.2720}{11.9552}{\normalfont}{de}
\texto{325.6211}{315.2720}{11.9552}{\normalfont}{verificación}
\texto{387.5419}{315.2720}{11.9552}{\normalfont}{documental}
\texto{451.3994}{315.2720}{11.9552}{\normalfont}{de}
\texto{467.7804}{315.2720}{11.9552}{\normalfont}{operadores}
\texto{527.4846}{315.2720}{11.9552}{\normalfont}{in-}
\texto{100.7540}{300.8270}{11.9552}{\normalfont}{centiva}
\texto{140.7752}{300.8270}{11.9552}{\normalfont}{la}
\texto{154.1364}{300.8270}{11.9552}{\normalfont}{formalización}
\texto{226.3553}{300.8270}{11.9552}{\normalfont}{de}
\texto{242.3179}{300.8270}{11.9552}{\normalfont}{agencias}
\texto{288.9121}{300.8270}{11.9552}{\normalfont}{y}
\texto{299.3466}{300.8270}{11.9552}{\normalfont}{guías,}
\texto{332.9335}{300.8270}{11.9552}{\normalfont}{promoviendo}
\texto{402.8786}{300.8270}{11.9552}{\normalfont}{el}
\texto{415.5894}{300.8270}{11.9552}{\normalfont}{cumplimiento}
\texto{489.1115}{300.8270}{11.9552}{\normalfont}{normativo}
\texto{100.7540}{286.3810}{11.9552}{\normalfont}{y}
\texto{110.8298}{286.3810}{11.9552}{\normalfont}{la}
\texto{123.8323}{286.3810}{11.9552}{\normalfont}{calidad}
\texto{164.1620}{286.3810}{11.9552}{\normalfont}{del}
\texto{183.0177}{286.3810}{11.9552}{\normalfont}{servicio}
\texto{225.0259}{286.3810}{11.9552}{\normalfont}{turístico.}
\texto{100.7540}{261.9720}{11.9552}{\normalfont\bfseries}{Sostenibilidad:}
\texto{190.8730}{261.9720}{11.9552}{\normalfont}{el}
\texto{203.4642}{261.9720}{11.9552}{\normalfont}{modelo}
\texto{244.3558}{261.9720}{11.9552}{\normalfont}{de}
\texto{260.1988}{261.9720}{11.9552}{\normalfont}{negocio}
\texto{302.3911}{261.9720}{11.9552}{\normalfont}{se}
\texto{316.3596}{261.9720}{11.9552}{\normalfont}{orienta}
\texto{355.9432}{261.9720}{11.9552}{\normalfont}{hacia}
\texto{386.7566}{261.9720}{11.9552}{\normalfont}{un}
\texto{403.9003}{261.9720}{11.9552}{\normalfont}{turismo}
\texto{447.1351}{261.9720}{11.9552}{\normalfont}{sostenible}
\texto{500.8283}{261.9720}{11.9552}{\normalfont}{que}
\texto{522.8617}{261.9720}{11.9552}{\normalfont}{res-}
\texto{100.7540}{247.5270}{11.9552}{\normalfont}{peta}
\texto{128.7244}{247.5270}{11.9552}{\normalfont}{el}
\texto{142.7263}{247.5270}{11.9552}{\normalfont}{medio}
\texto{178.8406}{247.5270}{11.9552}{\normalfont}{ambiente}
\texto{230.5564}{247.5270}{11.9552}{\normalfont}{y}
\texto{242.2701}{247.5270}{11.9552}{\normalfont}{las}
\texto{261.5395}{247.5270}{11.9552}{\normalfont}{comunidades}
\texto{332.5032}{247.5270}{11.9552}{\normalfont}{locales,}
\texto{374.8580}{247.5270}{11.9552}{\normalfont}{evitando}
\texto{424.3048}{247.5270}{11.9552}{\normalfont}{la}
\texto{438.9570}{247.5270}{11.9552}{\normalfont}{sobreexplotación}
\texto{529.4292}{247.5270}{11.9552}{\normalfont}{de}
\texto{100.7540}{233.0810}{11.9552}{\normalfont}{destinos}
\texto{145.7534}{233.0810}{11.9552}{\normalfont}{y}
\texto{155.8292}{233.0810}{11.9552}{\normalfont}{promoviendo}
\texto{225.4157}{233.0810}{11.9552}{\normalfont}{experiencias}
\texto{291.1716}{233.0810}{11.9552}{\normalfont}{responsables.}
\texto{100.7540}{208.6720}{11.9552}{\normalfont\bfseries}{T}
\texto{108.9959}{208.6720}{11.9552}{\normalfont\bfseries}{ransparencia:}
\texto{190.3380}{208.6720}{11.9552}{\normalfont}{el}
\texto{202.2239}{208.6720}{11.9552}{\normalfont}{sistema}
\texto{243.5171}{208.6720}{11.9552}{\normalfont}{de}
\texto{258.6668}{208.6720}{11.9552}{\normalfont}{reseñas}
\texto{298.6593}{208.6720}{11.9552}{\normalfont}{verificadas}
\texto{355.4871}{208.6720}{11.9552}{\normalfont}{y}
\texto{365.1087}{208.6720}{11.9552}{\normalfont}{la}
\texto{377.6569}{208.6720}{11.9552}{\normalfont}{bitácora}
\texto{422.7232}{208.6720}{11.9552}{\normalfont}{de}
\texto{437.8609}{208.6720}{11.9552}{\normalfont}{auditoría}
\texto{487.4797}{208.6720}{11.9552}{\normalfont}{garantizan}
\texto{100.7540}{194.2260}{11.9552}{\normalfont}{transparencia}
\texto{174.5271}{194.2260}{11.9552}{\normalfont}{en}
\texto{191.0038}{194.2260}{11.9552}{\normalfont}{las}
\texto{209.5081}{194.2260}{11.9552}{\normalfont}{transacciones}
\texto{282.0450}{194.2260}{11.9552}{\normalfont}{y}
\texto{292.9936}{194.2260}{11.9552}{\normalfont}{en}
\texto{309.4822}{194.2260}{11.9552}{\normalfont}{las}
\texto{327.9745}{194.2260}{11.9552}{\normalfont}{interacciones}
\texto{398.1731}{194.2260}{11.9552}{\normalfont}{entre}
\texto{428.6349}{194.2260}{11.9552}{\normalfont}{turistas}
\texto{471.9175}{194.2260}{11.9552}{\normalfont}{y}
\texto{482.8661}{194.2260}{11.9552}{\normalfont}{operadores,}
\texto{100.7540}{179.7810}{11.9552}{\normalfont}{generando}
\texto{156.6804}{179.7810}{11.9552}{\normalfont}{un}
\texto{173.5851}{179.7810}{11.9552}{\normalfont}{ecosistema}
\texto{231.6037}{179.7810}{11.9552}{\normalfont}{de}
\texto{247.2076}{179.7810}{11.9552}{\normalfont}{confianza.}
\end{pagina}

% Texto editable; coordenadas en puntos PDF.
% Página original 46
\begin{pagina}{46}
\texto{70.8660}{625.4930}{24.7871}{\normalfont\bfseries}{Capítulo}
\texto{185.1197}{625.4930}{24.7871}{\normalfont\bfseries}{4}
\texto{70.8660}{575.6790}{24.7871}{\normalfont\bfseries}{Modelado}
\texto{200.2745}{575.6790}{24.7871}{\normalfont\bfseries}{y}
\texto{224.2882}{575.6790}{24.7871}{\normalfont\bfseries}{automatización}
\texto{418.3266}{575.6790}{24.7871}{\normalfont\bfseries}{del}
\texto{70.8660}{545.7910}{24.7871}{\normalfont\bfseries}{proceso}
\texto{88.7990}{491.4950}{11.9552}{\normalfont}{El}
\texto{103.3867}{491.4950}{11.9552}{\normalfont}{presente}
\texto{148.7854}{491.4950}{11.9552}{\normalfont}{capítulo}
\texto{193.1416}{491.4950}{11.9552}{\normalfont}{modela}
\texto{233.2680}{491.4950}{11.9552}{\normalfont}{el}
\texto{245.1061}{491.4950}{11.9552}{\normalfont}{proceso}
\texto{286.5978}{491.4950}{11.9552}{\normalfont}{de}
\texto{301.6876}{491.4950}{11.9552}{\normalfont}{promoción,}
\texto{361.9753}{491.4950}{11.9552}{\normalfont}{reserva}
\texto{400.8608}{491.4950}{11.9552}{\normalfont}{y}
\texto{410.4226}{491.4950}{11.9552}{\normalfont}{gestión}
\texto{449.6404}{491.4950}{11.9552}{\normalfont}{de}
\texto{464.7303}{491.4950}{11.9552}{\normalfont}{tours}
\texto{494.1807}{491.4950}{11.9552}{\normalfont}{turísticos}
\texto{70.8660}{477.0490}{11.9552}{\normalfont}{que}
\texto{93.6765}{477.0490}{11.9552}{\normalfont}{atiende}
\texto{135.6710}{477.0490}{11.9552}{\normalfont\bfseries}{Explora-Nicaragua}
\texto{245.5780}{477.0490}{11.9552}{\normalfont}{,}
\texto{253.7554}{477.0490}{11.9552}{\normalfont}{tanto}
\texto{285.6614}{477.0490}{11.9552}{\normalfont}{en}
\texto{302.2935}{477.0490}{11.9552}{\normalfont}{su}
\texto{318.3397}{477.0490}{11.9552}{\normalfont}{estado}
\texto{355.8480}{477.0490}{11.9552}{\normalfont}{actual}
\texto{391.9909}{477.0490}{11.9552}{\normalfont}{(caracterizado}
\texto{469.1068}{477.0490}{11.9552}{\normalfont}{por}
\texto{491.2765}{477.0490}{11.9552}{\normalfont}{la}
\texto{505.3071}{477.0490}{11.9552}{\normalfont}{gestión}
\texto{70.8660}{462.6030}{11.9552}{\normalfont}{manual)}
\texto{117.9217}{462.6030}{11.9552}{\normalfont}{como}
\texto{149.7034}{462.6030}{11.9552}{\normalfont}{en}
\texto{166.5267}{462.6030}{11.9552}{\normalfont}{el}
\texto{180.0983}{462.6030}{11.9552}{\normalfont}{estado}
\texto{217.7978}{462.6030}{11.9552}{\normalfont}{propuesto}
\texto{273.0571}{462.6030}{11.9552}{\normalfont}{(automatizado}
\texto{352.3153}{462.6030}{11.9552}{\normalfont}{mediante}
\texto{403.9355}{462.6030}{11.9552}{\normalfont}{TI).}
\texto{429.5340}{462.6030}{11.9552}{\normalfont}{Se}
\texto{446.3573}{462.6030}{11.9552}{\normalfont}{utiliza}
\texto{483.3419}{462.6030}{11.9552}{\normalfont}{la}
\texto{497.5638}{462.6030}{11.9552}{\normalfont}{notación}
\texto{70.8660}{448.1580}{11.9552}{\normalfont\bfseries}{BPMN}
\texto{118.6390}{448.1580}{11.9552}{\normalfont\bfseries}{2.0}
\texto{140.8480}{448.1580}{11.9552}{\normalfont}{[11]}
\texto{164.0670}{448.1580}{11.9552}{\normalfont}{como}
\texto{195.7412}{448.1580}{11.9552}{\normalfont}{lenguaje}
\texto{243.0335}{448.1580}{11.9552}{\normalfont}{estándar}
\texto{291.6816}{448.1580}{11.9552}{\normalfont}{de}
\texto{308.4093}{448.1580}{11.9552}{\normalfont}{modelado}
\texto{362.5186}{448.1580}{11.9552}{\normalfont}{de}
\texto{379.2463}{448.1580}{11.9552}{\normalfont}{procesos,}
\texto{430.2448}{448.1580}{11.9552}{\normalfont}{apoyándose}
\texto{493.8536}{448.1580}{11.9552}{\normalfont}{en}
\texto{510.5694}{448.1580}{11.9552}{\normalfont}{herra-}
\texto{70.8660}{433.7120}{11.9552}{\normalfont}{mientas}
\texto{114.1773}{433.7120}{11.9552}{\normalfont}{como}
\texto{144.7450}{433.7120}{11.9552}{\normalfont\bfseries}{Camunda}
\texto{205.2598}{433.7120}{11.9552}{\normalfont\bfseries}{Modeler}
\texto{258.0020}{433.7120}{11.9552}{\normalfont}{[12].}
\texto{300.1470}{36.9930}{11.9552}{\normalfont}{36}
\end{pagina}
% Página original 47
\begin{pagina}{47}
\texto{70.8660}{746.3890}{11.9552}{\normalfont}{4.1.}
\texto{96.8710}{746.3890}{11.9552}{\normalfont}{DESCRIPCIÓN}
\texto{184.0028}{746.3890}{11.9552}{\normalfont}{DEL}
\texto{212.1191}{746.3890}{11.9552}{\normalfont}{PROCESO}
\texto{273.4062}{746.3890}{11.9552}{\normalfont}{ACTUAL}
\texto{529.4280}{746.3890}{11.9552}{\normalfont}{37}
\texto{70.8660}{705.1940}{17.2154}{\normalfont\bfseries}{4.1.}
\texto{120.3603}{705.1940}{17.2154}{\normalfont\bfseries}{Descripción}
\texto{225.7598}{705.1940}{17.2154}{\normalfont\bfseries}{del}
\texto{257.1917}{705.1940}{17.2154}{\normalfont\bfseries}{proceso}
\texto{327.3032}{705.1940}{17.2154}{\normalfont\bfseries}{actual}
\texto{88.7990}{678.9090}{11.9552}{\normalfont}{El}
\texto{104.4149}{678.9090}{11.9552}{\normalfont}{proceso}
\texto{146.9347}{678.9090}{11.9552}{\normalfont}{actual}
\texto{182.5636}{678.9090}{11.9552}{\normalfont}{de}
\texto{198.6697}{678.9090}{11.9552}{\normalfont}{promoción}
\texto{256.7337}{678.9090}{11.9552}{\normalfont}{y}
\texto{267.3236}{678.9090}{11.9552}{\normalfont}{reserva}
\texto{307.2372}{678.9090}{11.9552}{\normalfont}{de}
\texto{323.3552}{678.9090}{11.9552}{\normalfont}{tours}
\texto{353.8338}{678.9090}{11.9552}{\normalfont}{turísticos}
\texto{405.2005}{678.9090}{11.9552}{\normalfont}{en}
\texto{421.3185}{678.9090}{11.9552}{\normalfont}{Nicaragua}
\texto{477.4171}{678.9090}{11.9552}{\normalfont}{se}
\texto{491.6486}{678.9090}{11.9552}{\normalfont}{desarrolla}
\texto{70.8660}{664.4630}{11.9552}{\normalfont}{de}
\texto{86.2786}{664.4630}{11.9552}{\normalfont}{forma}
\texto{119.5763}{664.4630}{11.9552}{\normalfont}{mayoritariamente}
\texto{213.0396}{664.4630}{11.9552}{\normalfont}{manual}
\texto{254.1440}{664.4630}{11.9552}{\normalfont}{y}
\texto{264.0286}{664.4630}{11.9552}{\normalfont}{fragmentada.}
\texto{334.7364}{664.4630}{11.9552}{\normalfont}{El}
\texto{349.6589}{664.4630}{11.9552}{\normalfont}{punto}
\texto{382.9589}{664.4630}{11.9552}{\normalfont}{de}
\texto{398.3715}{664.4630}{11.9552}{\normalfont}{partida}
\texto{439.1603}{664.4630}{11.9552}{\normalfont}{es}
\texto{452.6864}{664.4630}{11.9552}{\normalfont}{la}
\texto{465.4976}{664.4630}{11.9552}{\normalfont}{gestión}
\texto{505.0502}{664.4630}{11.9552}{\normalfont}{interna}
\texto{70.8660}{650.0170}{11.9552}{\normalfont}{que}
\texto{92.8157}{650.0170}{11.9552}{\normalfont}{realiza}
\texto{130.0610}{650.0170}{11.9552}{\normalfont}{cada}
\texto{157.5508}{650.0170}{11.9552}{\normalfont}{operador}
\texto{206.8253}{650.0170}{11.9552}{\normalfont}{(agencia,}
\texto{256.4274}{650.0170}{11.9552}{\normalfont}{guía}
\texto{281.9542}{650.0170}{11.9552}{\normalfont}{independiente}
\texto{356.9205}{650.0170}{11.9552}{\normalfont}{o}
\texto{366.8505}{650.0170}{11.9552}{\normalfont}{cooperativa}
\texto{429.7850}{650.0170}{11.9552}{\normalfont}{comunitaria)}
\texto{499.2136}{650.0170}{11.9552}{\normalfont}{sobre}
\texto{530.0198}{650.0170}{11.9552}{\normalfont}{su}
\texto{70.8660}{635.5710}{11.9552}{\normalfont}{catálogo}
\texto{117.0370}{635.5710}{11.9552}{\normalfont}{de}
\texto{132.6409}{635.5710}{11.9552}{\normalfont}{tours,}
\texto{165.8811}{635.5710}{11.9552}{\normalfont}{sin}
\texto{184.1511}{635.5710}{11.9552}{\normalfont}{herramientas}
\texto{254.1273}{635.5710}{11.9552}{\normalfont}{digitales}
\texto{300.3748}{635.5710}{11.9552}{\normalfont}{centralizadas.}
\texto{88.7990}{621.1260}{11.9552}{\normalfont}{Las}
\texto{110.4810}{621.1260}{11.9552}{\normalfont}{actividades}
\texto{171.3497}{621.1260}{11.9552}{\normalfont}{principales}
\texto{229.9540}{621.1260}{11.9552}{\normalfont}{del}
\texto{248.8098}{621.1260}{11.9552}{\normalfont}{proceso}
\texto{290.8275}{621.1260}{11.9552}{\normalfont}{actual}
\texto{325.9423}{621.1260}{11.9552}{\normalfont}{son:}
\texto{85.6720}{593.7280}{11.9552}{\normalfont}{1.}
\texto{100.7540}{593.7280}{11.9552}{\normalfont\bfseries}{Registro}
\texto{156.3182}{593.7280}{11.9552}{\normalfont\bfseries}{de}
\texto{176.1076}{593.7280}{11.9552}{\normalfont\bfseries}{la}
\texto{192.5628}{593.7280}{11.9552}{\normalfont\bfseries}{oferta:}
\texto{235.9260}{593.7280}{11.9552}{\normalfont}{el}
\texto{249.7606}{593.7280}{11.9552}{\normalfont}{operador}
\texto{300.3382}{593.7280}{11.9552}{\normalfont}{registra}
\texto{344.1540}{593.7280}{11.9552}{\normalfont}{sus}
\texto{365.2717}{593.7280}{11.9552}{\normalfont}{tours,}
\texto{399.9944}{593.7280}{11.9552}{\normalfont}{precios}
\texto{440.5584}{593.7280}{11.9552}{\normalfont}{y}
\texto{452.1167}{593.7280}{11.9552}{\normalfont}{disponibilidad}
\texto{529.4238}{593.7280}{11.9552}{\normalfont}{en}
\texto{100.7540}{579.2820}{11.9552}{\normalfont}{documentos}
\texto{165.1997}{579.2820}{11.9552}{\normalfont}{locales}
\texto{202.6530}{579.2820}{11.9552}{\normalfont}{(Word,}
\texto{242.3227}{579.2820}{11.9552}{\normalfont}{Excel)}
\texto{278.5852}{579.2820}{11.9552}{\normalfont}{o}
\texto{288.3359}{579.2820}{11.9552}{\normalfont}{libretas}
\texto{330.0189}{579.2820}{11.9552}{\normalfont}{físicas,}
\texto{367.5510}{579.2820}{11.9552}{\normalfont}{sin}
\texto{385.8210}{579.2820}{11.9552}{\normalfont}{un}
\texto{402.7376}{579.2820}{11.9552}{\normalfont}{repositorio}
\texto{461.0024}{579.2820}{11.9552}{\normalfont}{centralizado.}
\texto{85.6720}{554.8740}{11.9552}{\normalfont}{2.}
\texto{100.7540}{554.8740}{11.9552}{\normalfont\bfseries}{Promoción}
\texto{169.1198}{554.8740}{11.9552}{\normalfont\bfseries}{por}
\texto{193.9269}{554.8740}{11.9552}{\normalfont\bfseries}{canales}
\texto{240.3776}{554.8740}{11.9552}{\normalfont\bfseries}{tradicionales:}
\texto{322.3490}{554.8740}{11.9552}{\normalfont}{la}
\texto{335.5786}{554.8740}{11.9552}{\normalfont}{oferta}
\texto{369.3066}{554.8740}{11.9552}{\normalfont}{se}
\texto{383.2512}{554.8740}{11.9552}{\normalfont}{difunde}
\texto{425.4339}{554.8740}{11.9552}{\normalfont}{mediante}
\texto{476.0618}{554.8740}{11.9552}{\normalfont}{agencias}
\texto{522.5364}{554.8740}{11.9552}{\normalfont}{físi-}
\texto{100.7540}{540.4280}{11.9552}{\normalfont}{cas,}
\texto{124.0786}{540.4280}{11.9552}{\normalfont}{propaganda}
\texto{188.3115}{540.4280}{11.9552}{\normalfont}{radial,}
\texto{225.2411}{540.4280}{11.9552}{\normalfont}{redes}
\texto{255.7197}{540.4280}{11.9552}{\normalfont}{sociales}
\texto{298.3041}{540.4280}{11.9552}{\normalfont}{(sin}
\texto{321.6287}{540.4280}{11.9552}{\normalfont}{funcionalidad}
\texto{394.6535}{540.4280}{11.9552}{\normalfont}{transaccional)}
\texto{470.6575}{540.4280}{11.9552}{\normalfont}{o}
\texto{480.9102}{540.4280}{11.9552}{\normalfont}{plataformas}
\texto{100.7540}{525.9820}{11.9552}{\normalfont}{extranjeras.}
\texto{85.6720}{501.5740}{11.9552}{\normalfont}{3.}
\texto{100.7540}{501.5740}{11.9552}{\normalfont\bfseries}{Recepción}
\texto{165.4639}{501.5740}{11.9552}{\normalfont\bfseries}{de}
\texto{183.7590}{501.5740}{11.9552}{\normalfont\bfseries}{solicitudes:}
\texto{252.4040}{501.5740}{11.9552}{\normalfont}{el}
\texto{264.9235}{501.5740}{11.9552}{\normalfont}{turista}
\texto{302.8836}{501.5740}{11.9552}{\normalfont}{consulta}
\texto{349.2985}{501.5740}{11.9552}{\normalfont}{disponibilidad}
\texto{425.2905}{501.5740}{11.9552}{\normalfont}{por}
\texto{446.5995}{501.5740}{11.9552}{\normalfont}{WhatsApp,}
\texto{509.2639}{501.5740}{11.9552}{\normalfont}{llama-}
\texto{100.7540}{487.1280}{11.9552}{\normalfont}{das}
\texto{122.7133}{487.1280}{11.9552}{\normalfont}{telefónicas}
\texto{180.7797}{487.1280}{11.9552}{\normalfont}{o}
\texto{191.6183}{487.1280}{11.9552}{\normalfont}{mensajes}
\texto{242.2677}{487.1280}{11.9552}{\normalfont}{en}
\texto{258.9596}{487.1280}{11.9552}{\normalfont}{redes}
\texto{290.0359}{487.1280}{11.9552}{\normalfont}{sociales.}
\texto{336.4579}{487.1280}{11.9552}{\normalfont}{Estas}
\texto{369.0598}{487.1280}{11.9552}{\normalfont}{solicitudes}
\texto{426.8536}{487.1280}{11.9552}{\normalfont}{se}
\texto{441.6709}{487.1280}{11.9552}{\normalfont}{reciben}
\texto{483.3993}{487.1280}{11.9552}{\normalfont}{en}
\texto{500.1031}{487.1280}{11.9552}{\normalfont}{horarios}
\texto{100.7540}{472.6820}{11.9552}{\normalfont}{limitados}
\texto{151.5421}{472.6820}{11.9552}{\normalfont}{y}
\texto{161.6179}{472.6820}{11.9552}{\normalfont}{dependen}
\texto{213.9769}{472.6820}{11.9552}{\normalfont}{de}
\texto{229.5808}{472.6820}{11.9552}{\normalfont}{la}
\texto{242.5953}{472.6820}{11.9552}{\normalfont}{disponibilidad}
\texto{318.4199}{472.6820}{11.9552}{\normalfont}{del}
\texto{337.2876}{472.6820}{11.9552}{\normalfont}{personal.}
\texto{85.6720}{448.2740}{11.9552}{\normalfont}{4.}
\texto{100.7540}{448.2740}{11.9552}{\normalfont\bfseries}{Confirmación}
\texto{185.3956}{448.2740}{11.9552}{\normalfont\bfseries}{manual:}
\texto{237.0510}{448.2740}{11.9552}{\normalfont}{el}
\texto{250.8258}{448.2740}{11.9552}{\normalfont}{operador}
\texto{301.3437}{448.2740}{11.9552}{\normalfont}{verifica}
\texto{343.0865}{448.2740}{11.9552}{\normalfont}{manualmente}
\texto{416.6874}{448.2740}{11.9552}{\normalfont}{la}
\texto{431.1126}{448.2740}{11.9552}{\normalfont}{disponibilidad}
\texto{508.3719}{448.2740}{11.9552}{\normalfont}{en}
\texto{525.3985}{448.2740}{11.9552}{\normalfont}{sus}
\texto{100.7540}{433.8280}{11.9552}{\normalfont}{registros}
\texto{148.4457}{433.8280}{11.9552}{\normalfont}{locales,}
\texto{189.8920}{433.8280}{11.9552}{\normalfont}{lo}
\texto{203.6357}{433.8280}{11.9552}{\normalfont}{que}
\texto{226.1593}{433.8280}{11.9552}{\normalfont}{frecuentemente}
\texto{307.8707}{433.8280}{11.9552}{\normalfont}{provoca}
\texto{352.1838}{433.8280}{11.9552}{\normalfont}{errores,}
\texto{394.6080}{433.8280}{11.9552}{\normalfont}{dobles}
\texto{431.1909}{433.8280}{11.9552}{\normalfont}{reservas}
\texto{475.9488}{433.8280}{11.9552}{\normalfont}{y}
\texto{486.7659}{433.8280}{11.9552}{\normalfont}{sobreventa}
\texto{100.7540}{419.3820}{11.9552}{\normalfont}{de}
\texto{116.3579}{419.3820}{11.9552}{\normalfont}{cupos.}
\texto{85.6720}{394.9740}{11.9552}{\normalfont}{5.}
\texto{100.7540}{394.9740}{11.9552}{\normalfont\bfseries}{Cobro}
\texto{141.5691}{394.9740}{11.9552}{\normalfont\bfseries}{y}
\texto{153.3664}{394.9740}{11.9552}{\normalfont\bfseries}{confirmación:}
\texto{235.9830}{394.9740}{11.9552}{\normalfont}{el}
\texto{248.5144}{394.9740}{11.9552}{\normalfont}{pago}
\texto{276.6665}{394.9740}{11.9552}{\normalfont}{se}
\texto{290.5633}{394.9740}{11.9552}{\normalfont}{realiza}
\texto{327.8205}{394.9740}{11.9552}{\normalfont}{por}
\texto{349.1414}{394.9740}{11.9552}{\normalfont}{transferencia}
\texto{418.6441}{394.9740}{11.9552}{\normalfont}{bancaria,}
\texto{469.5589}{394.9740}{11.9552}{\normalfont}{efectivo}
\texto{512.3346}{394.9740}{11.9552}{\normalfont}{o}
\texto{522.2766}{394.9740}{11.9552}{\normalfont}{me-}
\texto{100.7540}{380.5280}{11.9552}{\normalfont}{diante}
\texto{136.4355}{380.5280}{11.9552}{\normalfont}{la}
\texto{149.6771}{380.5280}{11.9552}{\normalfont}{retención}
\texto{200.6517}{380.5280}{11.9552}{\normalfont}{del}
\texto{219.7465}{380.5280}{11.9552}{\normalfont}{pago}
\texto{247.9464}{380.5280}{11.9552}{\normalfont}{por}
\texto{269.3271}{380.5280}{11.9552}{\normalfont}{la}
\texto{282.5806}{380.5280}{11.9552}{\normalfont}{plataforma}
\texto{342.3232}{380.5280}{11.9552}{\normalfont}{extranjera}
\texto{398.4887}{380.5280}{11.9552}{\normalfont}{intermediaria}
\texto{470.6025}{380.5280}{11.9552}{\normalfont}{(con}
\texto{496.8513}{380.5280}{11.9552}{\normalfont}{comisión}
\texto{100.7540}{366.0820}{11.9552}{\normalfont}{del}
\texto{119.6097}{366.0820}{11.9552}{\normalfont}{15}
\texto{133.3128}{366.0820}{11.9552}{\normalfont}{\%--25}
\texto{162.6126}{366.0820}{11.9552}{\normalfont}{\%)}
\texto{180.8299}{366.0820}{11.9552}{\normalfont}{[6].}
\texto{85.6720}{341.6730}{11.9552}{\normalfont}{6.}
\texto{100.7540}{341.6730}{11.9552}{\normalfont\bfseries}{Entrega}
\texto{151.0149}{341.6730}{11.9552}{\normalfont\bfseries}{de}
\texto{168.8317}{341.6730}{11.9552}{\normalfont\bfseries}{constancia:}
\texto{237.2210}{341.6730}{11.9552}{\normalfont}{en}
\texto{252.5858}{341.6730}{11.9552}{\normalfont}{el}
\texto{264.7108}{341.6730}{11.9552}{\normalfont}{mejor}
\texto{297.3222}{341.6730}{11.9552}{\normalfont}{de}
\texto{312.6990}{341.6730}{11.9552}{\normalfont}{los}
\texto{330.0914}{341.6730}{11.9552}{\normalfont}{casos,}
\texto{363.1571}{341.6730}{11.9552}{\normalfont}{el}
\texto{375.2820}{341.6730}{11.9552}{\normalfont}{operador}
\texto{424.1501}{341.6730}{11.9552}{\normalfont}{entrega}
\texto{465.2186}{341.6730}{11.9552}{\normalfont}{una}
\texto{487.7374}{341.6730}{11.9552}{\normalfont}{constancia}
\texto{100.7540}{327.2280}{11.9552}{\normalfont}{informal}
\texto{146.9154}{327.2280}{11.9552}{\normalfont}{(mensaje,}
\texto{199.3294}{327.2280}{11.9552}{\normalfont}{correo}
\texto{234.1095}{327.2280}{11.9552}{\normalfont}{o}
\texto{243.5373}{327.2280}{11.9552}{\normalfont}{captura}
\texto{286.1218}{327.2280}{11.9552}{\normalfont}{de}
\texto{301.4029}{327.2280}{11.9552}{\normalfont}{pantalla);}
\texto{354.0703}{327.2280}{11.9552}{\normalfont}{no}
\texto{370.0018}{327.2280}{11.9552}{\normalfont}{existe}
\texto{402.5702}{327.2280}{11.9552}{\normalfont}{un}
\texto{419.1521}{327.2280}{11.9552}{\normalfont}{voucher}
\texto{462.0664}{327.2280}{11.9552}{\normalfont}{digital}
\texto{498.1472}{327.2280}{11.9552}{\normalfont}{estanda-}
\texto{100.7540}{312.7820}{11.9552}{\normalfont}{rizado.}
\texto{85.6720}{288.3730}{11.9552}{\normalfont}{7.}
\texto{100.7540}{288.3730}{11.9552}{\normalfont\bfseries}{Control}
\texto{148.8139}{288.3730}{11.9552}{\normalfont\bfseries}{del}
\texto{169.4940}{288.3730}{11.9552}{\normalfont\bfseries}{servicio:}
\texto{220.3510}{288.3730}{11.9552}{\normalfont}{el}
\texto{231.7108}{288.3730}{11.9552}{\normalfont}{día}
\texto{250.2247}{288.3730}{11.9552}{\normalfont}{del}
\texto{268.0881}{288.3730}{11.9552}{\normalfont}{tour,}
\texto{295.7070}{288.3730}{11.9552}{\normalfont}{el}
\texto{307.0668}{288.3730}{11.9552}{\normalfont}{operador}
\texto{355.1698}{288.3730}{11.9552}{\normalfont}{identifica}
\texto{404.5782}{288.3730}{11.9552}{\normalfont}{al}
\texto{416.5884}{288.3730}{11.9552}{\normalfont}{turista}
\texto{453.3770}{288.3730}{11.9552}{\normalfont}{de}
\texto{467.9886}{288.3730}{11.9552}{\normalfont}{forma}
\texto{500.4852}{288.3730}{11.9552}{\normalfont}{manual;}
\texto{100.7540}{273.9280}{11.9552}{\normalfont}{no}
\texto{117.0083}{273.9280}{11.9552}{\normalfont}{existe}
\texto{149.9114}{273.9280}{11.9552}{\normalfont}{un}
\texto{166.8280}{273.9280}{11.9552}{\normalfont}{mecanismo}
\texto{226.7212}{273.9280}{11.9552}{\normalfont}{estandarizado}
\texto{300.9223}{273.9280}{11.9552}{\normalfont}{de}
\texto{316.5382}{273.9280}{11.9552}{\normalfont}{validación}
\texto{371.4939}{273.9280}{11.9552}{\normalfont}{ni}
\texto{385.1467}{273.9280}{11.9552}{\normalfont}{de}
\texto{400.7626}{273.9280}{11.9552}{\normalfont}{registro}
\texto{443.0959}{273.9280}{11.9552}{\normalfont}{de}
\texto{458.6999}{273.9280}{11.9552}{\normalfont}{auditoría.}
\texto{85.6720}{249.5190}{11.9552}{\normalfont}{8.}
\texto{100.7540}{249.5190}{11.9552}{\normalfont\bfseries}{Retroalimentación:}
\texto{215.3590}{249.5190}{11.9552}{\normalfont}{las}
\texto{232.1895}{249.5190}{11.9552}{\normalfont}{reseñas}
\texto{271.8593}{249.5190}{11.9552}{\normalfont}{se}
\texto{284.7996}{249.5190}{11.9552}{\normalfont}{gestionan}
\texto{336.0993}{249.5190}{11.9552}{\normalfont}{por}
\texto{356.4519}{249.5190}{11.9552}{\normalfont}{canales}
\texto{396.0571}{249.5190}{11.9552}{\normalfont}{informales}
\texto{451.5961}{249.5190}{11.9552}{\normalfont}{sin}
\texto{469.0770}{249.5190}{11.9552}{\normalfont}{verificación}
\texto{529.4317}{249.5190}{11.9552}{\normalfont}{de}
\texto{100.7540}{235.0730}{11.9552}{\normalfont}{que}
\texto{122.5364}{235.0730}{11.9552}{\normalfont}{el}
\texto{134.8885}{235.0730}{11.9552}{\normalfont}{turista}
\texto{172.6813}{235.0730}{11.9552}{\normalfont}{realmente}
\texto{226.3338}{235.0730}{11.9552}{\normalfont}{realizó}
\texto{263.3997}{235.0730}{11.9552}{\normalfont}{la}
\texto{276.4141}{235.0730}{11.9552}{\normalfont}{experiencia,}
\texto{340.7929}{235.0730}{11.9552}{\normalfont}{lo}
\texto{353.8073}{235.0730}{11.9552}{\normalfont}{que}
\texto{375.5897}{235.0730}{11.9552}{\normalfont}{limita}
\texto{409.4038}{235.0730}{11.9552}{\normalfont}{la}
\texto{422.4182}{235.0730}{11.9552}{\normalfont}{confianza}
\texto{473.7921}{235.0730}{11.9552}{\normalfont}{del}
\texto{492.6598}{235.0730}{11.9552}{\normalfont}{mercado.}
\end{pagina}
% Página original 48
\begin{pagina}{48}
\texto{70.8660}{746.3890}{11.9552}{\normalfont}{38}
\texto{168.2490}{746.3890}{11.9552}{\normalfont}{CAPÍTULO}
\texto{235.2125}{746.3890}{11.9552}{\normalfont}{4.}
\texto{252.1243}{746.3890}{11.9552}{\normalfont}{MODELADO}
\texto{326.8921}{746.3890}{11.9552}{\normalfont}{Y}
\texto{339.5646}{746.3890}{11.9552}{\normalfont}{AUTOMATIZACIÓN}
\texto{455.6233}{746.3890}{11.9552}{\normalfont}{DEL}
\texto{483.7396}{746.3890}{11.9552}{\normalfont}{PROCESO}
\texto{70.8660}{705.1940}{17.2154}{\normalfont\bfseries}{4.2.}
\texto{120.3603}{705.1940}{17.2154}{\normalfont\bfseries}{Actores}
\texto{191.1001}{705.1940}{17.2154}{\normalfont\bfseries}{y}
\texto{207.7784}{705.1940}{17.2154}{\normalfont\bfseries}{elementos}
\texto{297.3449}{705.1940}{17.2154}{\normalfont\bfseries}{del}
\texto{328.7768}{705.1940}{17.2154}{\normalfont\bfseries}{proceso}
\texto{88.7990}{678.9090}{11.9552}{\normalfont}{Los}
\texto{110.4810}{678.9090}{11.9552}{\normalfont}{actores}
\texto{150.2129}{678.9090}{11.9552}{\normalfont}{que}
\texto{172.0072}{678.9090}{11.9552}{\normalfont}{participan}
\texto{227.9336}{678.9090}{11.9552}{\normalfont}{en}
\texto{243.5375}{678.9090}{11.9552}{\normalfont}{el}
\texto{255.9016}{678.9090}{11.9552}{\normalfont}{proceso}
\texto{297.9074}{678.9090}{11.9552}{\normalfont}{actual}
\texto{333.0342}{678.9090}{11.9552}{\normalfont}{y}
\texto{343.1100}{678.9090}{11.9552}{\normalfont}{propuesto}
\texto{397.1619}{678.9090}{11.9552}{\normalfont}{son:}
\texto{224.4830}{644.4380}{11.9552}{\normalfont}{Cuadro}
\texto{266.1014}{644.4380}{11.9552}{\normalfont}{4.1:}
\texto{288.2160}{644.4380}{11.9552}{\normalfont}{Actores}
\texto{330.5470}{644.4380}{11.9552}{\normalfont}{del}
\texto{349.4027}{644.4380}{11.9552}{\normalfont}{proceso}
\texto{108.6840}{626.7230}{11.9552}{\normalfont\bfseries}{Actor}
\texto{234.0250}{626.7230}{11.9552}{\normalfont\bfseries}{Rol}
\texto{258.9887}{626.7230}{11.9552}{\normalfont\bfseries}{en}
\texto{277.0805}{626.7230}{11.9552}{\normalfont\bfseries}{el}
\texto{291.4363}{626.7230}{11.9552}{\normalfont\bfseries}{proceso}
\texto{108.6840}{606.2870}{11.9552}{\normalfont\bfseries}{T}
\texto{116.9259}{606.2870}{11.9552}{\normalfont\bfseries}{urista}
\texto{234.0250}{606.2870}{11.9552}{\normalfont}{Demandante}
\texto{302.9133}{606.2870}{11.9552}{\normalfont}{de}
\texto{318.9715}{606.2870}{11.9552}{\normalfont}{experiencias}
\texto{385.1698}{606.2870}{11.9552}{\normalfont}{turísticas;}
\texto{439.7166}{606.2870}{11.9552}{\normalfont}{consulta,}
\texto{489.6583}{606.2870}{11.9552}{\normalfont}{re-}
\texto{234.0250}{591.8420}{11.9552}{\normalfont}{serva}
\texto{264.2071}{591.8420}{11.9552}{\normalfont}{y}
\texto{274.8329}{591.8420}{11.9552}{\normalfont}{paga}
\texto{303.3317}{591.8420}{11.9552}{\normalfont}{por}
\texto{324.9993}{591.8420}{11.9552}{\normalfont}{los}
\texto{343.1688}{591.8420}{11.9552}{\normalfont}{tours.}
\texto{376.9351}{591.8420}{11.9552}{\normalfont}{En}
\texto{395.8386}{591.8420}{11.9552}{\normalfont}{el}
\texto{408.7407}{591.8420}{11.9552}{\normalfont}{modelo}
\texto{449.9192}{591.8420}{11.9552}{\normalfont}{propuesto,}
\texto{234.0250}{577.3960}{11.9552}{\normalfont}{interactúa}
\texto{288.8157}{577.3960}{11.9552}{\normalfont}{con}
\texto{309.4599}{577.3960}{11.9552}{\normalfont}{la}
\texto{321.6494}{577.3960}{11.9552}{\normalfont}{app}
\texto{343.5944}{577.3960}{11.9552}{\normalfont}{móvil,}
\texto{377.8987}{577.3960}{11.9552}{\normalfont}{genera}
\texto{414.1636}{577.3960}{11.9552}{\normalfont}{su}
\texto{428.3687}{577.3960}{11.9552}{\normalfont}{voucher}
\texto{470.8049}{577.3960}{11.9552}{\normalfont}{digital}
\texto{234.0250}{562.9500}{11.9552}{\normalfont}{y}
\texto{244.1008}{562.9500}{11.9552}{\normalfont}{deja}
\texto{269.1350}{562.9500}{11.9552}{\normalfont}{reseñas}
\texto{309.5938}{562.9500}{11.9552}{\normalfont}{verificadas.}
\texto{108.6840}{542.5140}{11.9552}{\normalfont\bfseries}{Operador}
\texto{234.0250}{542.5140}{11.9552}{\normalfont}{Oferente}
\texto{280.6025}{542.5140}{11.9552}{\normalfont}{de}
\texto{295.3097}{542.5140}{11.9552}{\normalfont}{tours}
\texto{324.4015}{542.5140}{11.9552}{\normalfont}{(agencia,}
\texto{372.9396}{542.5140}{11.9552}{\normalfont}{guía}
\texto{397.4024}{542.5140}{11.9552}{\normalfont}{o}
\texto{406.2683}{542.5140}{11.9552}{\normalfont}{cooperativa);}
\texto{475.9433}{542.5140}{11.9552}{\normalfont}{regis-}
\texto{234.0250}{528.0680}{11.9552}{\normalfont}{tra}
\texto{252.7373}{528.0680}{11.9552}{\normalfont}{su}
\texto{267.6119}{528.0680}{11.9552}{\normalfont}{catálogo,}
\texto{316.9032}{528.0680}{11.9552}{\normalfont}{gestiona}
\texto{362.3449}{528.0680}{11.9552}{\normalfont}{disponibilidad,}
\texto{441.2899}{528.0680}{11.9552}{\normalfont}{recibe}
\texto{475.2833}{528.0680}{11.9552}{\normalfont}{reser-}
\texto{234.0250}{513.6220}{11.9552}{\normalfont}{vas}
\texto{253.1127}{513.6220}{11.9552}{\normalfont}{y}
\texto{262.3756}{513.6220}{11.9552}{\normalfont}{valida}
\texto{295.7067}{513.6220}{11.9552}{\normalfont}{vouchers.}
\texto{346.0237}{513.6220}{11.9552}{\normalfont}{En}
\texto{363.5763}{513.6220}{11.9552}{\normalfont}{el}
\texto{375.1274}{513.6220}{11.9552}{\normalfont}{modelo}
\texto{414.9670}{513.6220}{11.9552}{\normalfont}{propuesto,}
\texto{471.4457}{513.6220}{11.9552}{\normalfont}{utiliza}
\texto{234.0250}{499.1770}{11.9552}{\normalfont}{el}
\texto{246.3771}{499.1770}{11.9552}{\normalfont}{backoffice}
\texto{299.3817}{499.1770}{11.9552}{\normalfont}{de}
\texto{314.9856}{499.1770}{11.9552}{\normalfont}{escritorio.}
\texto{108.6840}{478.7410}{11.9552}{\normalfont\bfseries}{Equipo}
\texto{154.8693}{478.7410}{11.9552}{\normalfont\bfseries}{gestor}
\texto{234.0250}{478.7410}{11.9552}{\normalfont}{Administrador}
\texto{312.3650}{478.7410}{11.9552}{\normalfont}{de}
\texto{328.2200}{478.7410}{11.9552}{\normalfont}{la}
\texto{341.4736}{478.7410}{11.9552}{\normalfont}{plataforma;}
\texto{404.4799}{478.7410}{11.9552}{\normalfont}{verifica}
\texto{445.0510}{478.7410}{11.9552}{\normalfont}{operadores,}
\texto{234.0250}{464.2950}{11.9552}{\normalfont}{modera}
\texto{276.5401}{464.2950}{11.9552}{\normalfont}{la}
\texto{290.1284}{464.2950}{11.9552}{\normalfont}{oferta,}
\texto{327.4429}{464.2950}{11.9552}{\normalfont}{supervisa}
\texto{379.5293}{464.2950}{11.9552}{\normalfont}{reservas}
\texto{424.1318}{464.2950}{11.9552}{\normalfont}{y}
\texto{434.7935}{464.2950}{11.9552}{\normalfont}{pagos,}
\texto{471.1970}{464.2950}{11.9552}{\normalfont}{y}
\texto{481.8587}{464.2950}{11.9552}{\normalfont}{con-}
\texto{234.0250}{449.8490}{11.9552}{\normalfont}{sulta}
\texto{262.7007}{449.8490}{11.9552}{\normalfont}{reportes}
\texto{307.9583}{449.8490}{11.9552}{\normalfont}{desde}
\texto{339.8979}{449.8490}{11.9552}{\normalfont}{el}
\texto{352.2500}{449.8490}{11.9552}{\normalfont}{panel}
\texto{383.4746}{449.8490}{11.9552}{\normalfont}{de}
\texto{399.0785}{449.8490}{11.9552}{\normalfont}{administración.}
\texto{88.7990}{416.4510}{11.9552}{\normalfont}{Fuente:}
\texto{130.2501}{416.4510}{11.9552}{\normalfont}{Elaboración}
\texto{195.1262}{416.4510}{11.9552}{\normalfont}{propia}
\texto{231.5417}{416.4510}{11.9552}{\normalfont}{(2026).}
\texto{88.7990}{402.0050}{11.9552}{\normalfont}{Los}
\texto{110.4810}{402.0050}{11.9552}{\normalfont}{elementos}
\texto{164.1981}{402.0050}{11.9552}{\normalfont}{que}
\texto{185.9804}{402.0050}{11.9552}{\normalfont}{intervienen}
\texto{246.4737}{402.0050}{11.9552}{\normalfont}{en}
\texto{262.0777}{402.0050}{11.9552}{\normalfont}{el}
\texto{274.4298}{402.0050}{11.9552}{\normalfont}{proceso}
\texto{316.4475}{402.0050}{11.9552}{\normalfont}{son:}
\texto{100.7540}{374.6080}{11.9552}{\normalfont\bfseries}{Información:}
\texto{178.2120}{374.6080}{11.9552}{\normalfont}{catálogo}
\texto{223.7374}{374.6080}{11.9552}{\normalfont}{de}
\texto{238.6957}{374.6080}{11.9552}{\normalfont}{tours,}
\texto{271.2784}{374.6080}{11.9552}{\normalfont}{precios,}
\texto{312.9782}{374.6080}{11.9552}{\normalfont}{disponibilidad,}
\texto{391.4091}{374.6080}{11.9552}{\normalfont}{datos}
\texto{422.0527}{374.6080}{11.9552}{\normalfont}{del}
\texto{440.2628}{374.6080}{11.9552}{\normalfont}{turista,}
\texto{480.6499}{374.6080}{11.9552}{\normalfont}{comproban-}
\texto{100.7540}{360.1620}{11.9552}{\normalfont}{tes}
\texto{119.0239}{360.1620}{11.9552}{\normalfont}{de}
\texto{134.6279}{360.1620}{11.9552}{\normalfont}{pago.}
\texto{100.7540}{335.7540}{11.9552}{\normalfont\bfseries}{Recursos}
\texto{157.6847}{335.7540}{11.9552}{\normalfont\bfseries}{físicos:}
\texto{200.2340}{335.7540}{11.9552}{\normalfont}{documentos}
\texto{264.6917}{335.7540}{11.9552}{\normalfont}{impresos,}
\texto{316.2066}{335.7540}{11.9552}{\normalfont}{libretas,}
\texto{361.1534}{335.7540}{11.9552}{\normalfont}{hojas}
\texto{391.7898}{335.7540}{11.9552}{\normalfont}{de}
\texto{407.4057}{335.7540}{11.9552}{\normalfont}{cálculo,}
\texto{449.6864}{335.7540}{11.9552}{\normalfont}{WhatsApp}
\texto{508.9316}{335.7540}{11.9552}{\normalfont}{y}
\texto{519.0194}{335.7540}{11.9552}{\normalfont}{telé-}
\texto{100.7540}{321.3080}{11.9552}{\normalfont}{fono.}
\texto{100.7540}{296.8990}{11.9552}{\normalfont\bfseries}{Recursos}
\texto{157.3141}{296.8990}{11.9552}{\normalfont\bfseries}{digitales:}
\texto{213.2540}{296.8990}{11.9552}{\normalfont}{redes}
\texto{242.9077}{296.8990}{11.9552}{\normalfont}{sociales,}
\texto{287.9190}{296.8990}{11.9552}{\normalfont}{plataformas}
\texto{351.7287}{296.8990}{11.9552}{\normalfont}{extranjeras}
\texto{411.9614}{296.8990}{11.9552}{\normalfont}{de}
\texto{427.2425}{296.8990}{11.9552}{\normalfont}{reserva}
\texto{466.3432}{296.8990}{11.9552}{\normalfont}{(opcional),}
\texto{524.8830}{296.8990}{11.9552}{\normalfont}{pa-}
\texto{100.7540}{282.4540}{11.9552}{\normalfont}{sarela}
\texto{134.4366}{282.4540}{11.9552}{\normalfont}{de}
\texto{150.4948}{282.4540}{11.9552}{\normalfont}{pago,}
\texto{182.1617}{282.4540}{11.9552}{\normalfont}{app}
\texto{205.3859}{282.4540}{11.9552}{\normalfont}{móvil,}
\texto{240.9574}{282.4540}{11.9552}{\normalfont}{backoffice}
\texto{294.4043}{282.4540}{11.9552}{\normalfont}{de}
\texto{310.4745}{282.4540}{11.9552}{\normalfont}{escritorio}
\texto{361.7169}{282.4540}{11.9552}{\normalfont}{y}
\texto{372.2470}{282.4540}{11.9552}{\normalfont}{panel}
\texto{403.9140}{282.4540}{11.9552}{\normalfont}{de}
\texto{419.9841}{282.4540}{11.9552}{\normalfont}{administración}
\texto{499.8425}{282.4540}{11.9552}{\normalfont}{(modelo}
\texto{100.7540}{268.0080}{11.9552}{\normalfont}{propuesto).}
\texto{100.7540}{243.5990}{11.9552}{\normalfont\bfseries}{Sistemas}
\texto{156.4688}{243.5990}{11.9552}{\normalfont\bfseries}{de}
\texto{174.8715}{243.5990}{11.9552}{\normalfont\bfseries}{información:}
\texto{251.9740}{243.5990}{11.9552}{\normalfont}{la}
\texto{265.2395}{243.5990}{11.9552}{\normalfont}{API}
\texto{290.3741}{243.5990}{11.9552}{\normalfont}{REST}
\texto{326.0843}{243.5990}{11.9552}{\normalfont}{centralizada,}
\texto{394.9582}{243.5990}{11.9552}{\normalfont}{la}
\texto{408.2356}{243.5990}{11.9552}{\normalfont}{base}
\texto{434.5849}{243.5990}{11.9552}{\normalfont}{de}
\texto{450.4638}{243.5990}{11.9552}{\normalfont}{datos}
\texto{482.0040}{243.5990}{11.9552}{\normalfont}{relacional}
\texto{534.9536}{243.5990}{11.9552}{\normalfont}{y}
\texto{100.7540}{229.1530}{11.9552}{\normalfont}{el}
\texto{113.1061}{229.1530}{11.9552}{\normalfont}{mecanismo}
\texto{172.9993}{229.1530}{11.9552}{\normalfont}{de}
\texto{188.6152}{229.1530}{11.9552}{\normalfont}{autenticación}
\texto{260.4779}{229.1530}{11.9552}{\normalfont}{JWT}
\texto{290.8704}{229.1530}{11.9552}{\normalfont}{con}
\texto{312.3395}{229.1530}{11.9552}{\normalfont}{control}
\texto{351.6841}{229.1530}{11.9552}{\normalfont}{de}
\texto{367.2880}{229.1530}{11.9552}{\normalfont}{acceso}
\texto{403.1297}{229.1530}{11.9552}{\normalfont}{basado}
\texto{442.2112}{229.1530}{11.9552}{\normalfont}{en}
\texto{457.8152}{229.1530}{11.9552}{\normalfont}{roles}
\texto{485.2021}{229.1530}{11.9552}{\normalfont}{(RBAC).}
\end{pagina}
% Página original 49
\begin{pagina}{49}
\texto{70.8660}{746.3890}{11.9552}{\normalfont}{4.3.}
\texto{96.8710}{746.3890}{11.9552}{\normalfont}{BPMN}
\texto{136.5359}{746.3890}{11.9552}{\normalfont}{DEL}
\texto{164.6522}{746.3890}{11.9552}{\normalfont}{PROCESO}
\texto{225.9393}{746.3890}{11.9552}{\normalfont}{ACTUAL}
\texto{529.4280}{746.3890}{11.9552}{\normalfont}{39}
\texto{70.8660}{705.1940}{17.2154}{\normalfont\bfseries}{4.3.}
\texto{120.3603}{705.1940}{17.2154}{\normalfont\bfseries}{BPMN}
\texto{187.3282}{705.1940}{17.2154}{\normalfont\bfseries}{del}
\texto{218.7601}{705.1940}{17.2154}{\normalfont\bfseries}{proceso}
\texto{288.8715}{705.1940}{17.2154}{\normalfont\bfseries}{actual}
\texto{88.7990}{678.9090}{11.9552}{\normalfont}{La}
\texto{107.6571}{678.9090}{11.9552}{\normalfont}{figura}
\texto{142.6140}{678.9090}{11.9552}{\normalfont}{4.1}
\texto{163.2630}{678.9090}{11.9552}{\normalfont}{presenta}
\texto{211.6074}{678.9090}{11.9552}{\normalfont}{el}
\texto{225.7528}{678.9090}{11.9552}{\normalfont}{diagrama}
\texto{278.9200}{678.9090}{11.9552}{\normalfont}{BPMN}
\texto{320.3663}{678.9090}{11.9552}{\normalfont}{del}
\texto{341.0153}{678.9090}{11.9552}{\normalfont}{proceso}
\texto{384.8263}{678.9090}{11.9552}{\normalfont}{actual}
\texto{421.7344}{678.9090}{11.9552}{\normalfont}{de}
\texto{439.1316}{678.9090}{11.9552}{\normalfont}{reserva}
\texto{480.3245}{678.9090}{11.9552}{\normalfont}{de}
\texto{497.7217}{678.9090}{11.9552}{\normalfont}{un}
\texto{516.4196}{678.9090}{11.9552}{\normalfont}{tour,}
\texto{70.8660}{664.4630}{11.9552}{\normalfont}{evidenciando}
\texto{142.3533}{664.4630}{11.9552}{\normalfont}{la}
\texto{156.9339}{664.4630}{11.9552}{\normalfont}{gestión}
\texto{198.2558}{664.4630}{11.9552}{\normalfont}{manual,}
\texto{244.3814}{664.4630}{11.9552}{\normalfont}{los}
\texto{263.5790}{664.4630}{11.9552}{\normalfont}{canales}
\texto{305.5394}{664.4630}{11.9552}{\normalfont}{informales}
\texto{363.4456}{664.4630}{11.9552}{\normalfont}{de}
\texto{380.6276}{664.4630}{11.9552}{\normalfont}{comunicación}
\texto{454.7188}{664.4630}{11.9552}{\normalfont}{y}
\texto{466.3847}{664.4630}{11.9552}{\normalfont}{la}
\texto{480.9652}{664.4630}{11.9552}{\normalfont}{ausencia}
\texto{529.4292}{664.4630}{11.9552}{\normalfont}{de}
\texto{70.8660}{650.0170}{11.9552}{\normalfont}{automatización.}
\texto{150.5160}{626.2790}{7.9701}{\normalfont}{Turista}
\texto{224.7360}{621.6470}{7.9701}{\normalfont}{Consulta}
\texto{260.5377}{621.6470}{7.9701}{\normalfont}{por}
\texto{209.4580}{612.1830}{7.9701}{\normalfont}{WhatsApp}
\texto{252.2997}{612.1830}{7.9701}{\normalfont}{/}
\texto{259.3556}{612.1830}{7.9701}{\normalfont}{teléfono}
\texto{322.2420}{620.6500}{7.9701}{\normalfont}{Espera}
\texto{350.3972}{620.6500}{7.9701}{\normalfont}{confir-}
\texto{319.9770}{611.1860}{7.9701}{\normalfont}{mación}
\texto{349.1444}{611.1860}{7.9701}{\normalfont}{manual}
\texto{146.3020}{558.2470}{7.9701}{\normalfont}{Operador}
\texto{219.7400}{553.1490}{7.9701}{\normalfont}{Revisa}
\texto{246.9507}{553.1490}{7.9701}{\normalfont}{libreta}
\texto{273.7725}{553.1490}{7.9701}{\normalfont}{/}
\texto{221.1170}{543.6850}{7.9701}{\normalfont}{hoja}
\texto{239.9312}{543.6850}{7.9701}{\normalfont}{de}
\texto{251.2209}{543.6850}{7.9701}{\normalfont}{cálculo}
\texto{337.5100}{553.3930}{7.9701}{\normalfont}{£Hay}
\texto{337.2670}{543.9290}{7.9701}{\normalfont}{cupo?}
\texto{408.3050}{558.1250}{7.9701}{\normalfont}{Confirma}
\texto{445.4783}{558.1250}{7.9701}{\normalfont}{por}
\texto{403.2270}{548.6610}{7.9701}{\normalfont}{mensaje}
\texto{435.7307}{548.6610}{7.9701}{\normalfont}{/}
\texto{442.7866}{548.6610}{7.9701}{\normalfont}{cobra}
\texto{401.8080}{539.1960}{7.9701}{\normalfont}{por}
\texto{417.1002}{539.1960}{7.9701}{\normalfont}{transferencia}
\texto{143.1910}{484.5450}{7.9701}{\normalfont}{Plataforma}
\texto{145.0080}{475.0810}{7.9701}{\normalfont}{extranjera}
\texto{324.7960}{478.8960}{7.9701}{\normalfont}{Cobra}
\texto{350.2079}{478.8960}{7.9701}{\normalfont}{15}
\texto{360.0876}{478.8960}{7.9701}{\normalfont}{\%--}
\texto{316.5430}{469.4320}{7.9701}{\normalfont}{25}
\texto{326.4227}{469.4320}{7.9701}{\normalfont}{\%}
\texto{336.3009}{469.4320}{7.9701}{\normalfont}{de}
\texto{347.5905}{469.4320}{7.9701}{\normalfont}{comisión}
\texto{384.2560}{554.9310}{5.9776}{\normalfont}{sí}
\texto{114.9220}{435.3150}{11.9552}{\normalfont}{Figura}
\texto{152.4733}{435.3150}{11.9552}{\normalfont}{4.1:}
\texto{174.5890}{435.3150}{11.9552}{\normalfont}{BPMN}
\texto{214.2420}{435.3150}{11.9552}{\normalfont}{del}
\texto{233.0977}{435.3150}{11.9552}{\normalfont}{proceso}
\texto{275.1155}{435.3150}{11.9552}{\normalfont}{actual}
\texto{310.2303}{435.3150}{11.9552}{\normalfont}{de}
\texto{325.8462}{435.3150}{11.9552}{\normalfont}{reserva}
\texto{365.2577}{435.3150}{11.9552}{\normalfont}{de}
\texto{380.8616}{435.3150}{11.9552}{\normalfont}{tours}
\texto{410.8381}{435.3150}{11.9552}{\normalfont}{(gestión}
\texto{455.1345}{435.3150}{11.9552}{\normalfont}{manual)}
\texto{88.7990}{403.9330}{11.9552}{\normalfont}{El}
\texto{103.7454}{403.9330}{11.9552}{\normalfont}{diagrama}
\texto{154.9519}{403.9330}{11.9552}{\normalfont}{muestra}
\texto{199.3966}{403.9330}{11.9552}{\normalfont}{que}
\texto{221.0116}{403.9330}{11.9552}{\normalfont}{el}
\texto{233.2083}{403.9330}{11.9552}{\normalfont}{proceso}
\texto{275.0467}{403.9330}{11.9552}{\normalfont}{depende}
\texto{320.7347}{403.9330}{11.9552}{\normalfont}{de}
\texto{336.1712}{403.9330}{11.9552}{\normalfont}{comunicaciones}
\texto{418.3369}{403.9330}{11.9552}{\normalfont}{informales,}
\texto{477.7375}{403.9330}{11.9552}{\normalfont}{carece}
\texto{512.6969}{403.9330}{11.9552}{\normalfont}{de}
\texto{528.1335}{403.9330}{11.9552}{\normalfont}{un}
\texto{70.8660}{389.4870}{11.9552}{\normalfont}{registro}
\texto{112.9842}{389.4870}{11.9552}{\normalfont}{estandarizado}
\texto{186.9821}{389.4870}{11.9552}{\normalfont}{de}
\texto{202.3828}{389.4870}{11.9552}{\normalfont}{las}
\texto{219.7991}{389.4870}{11.9552}{\normalfont}{solicitudes}
\texto{276.3018}{389.4870}{11.9552}{\normalfont}{y}
\texto{286.1744}{389.4870}{11.9552}{\normalfont}{los}
\texto{303.5907}{389.4870}{11.9552}{\normalfont}{pagos,}
\texto{339.2052}{389.4870}{11.9552}{\normalfont}{y}
\texto{349.0779}{389.4870}{11.9552}{\normalfont}{opera}
\texto{381.0604}{389.4870}{11.9552}{\normalfont}{con}
\texto{402.3144}{389.4870}{11.9552}{\normalfont}{intermediarios}
\texto{478.5909}{389.4870}{11.9552}{\normalfont}{que}
\texto{500.1701}{389.4870}{11.9552}{\normalfont}{retienen}
\texto{70.8660}{375.0410}{11.9552}{\normalfont}{comisiones}
\texto{128.8726}{375.0410}{11.9552}{\normalfont}{elevadas.}
\texto{178.0396}{375.0410}{11.9552}{\normalfont}{La}
\texto{195.1044}{375.0410}{11.9552}{\normalfont}{validación}
\texto{250.0601}{375.0410}{11.9552}{\normalfont}{de}
\texto{265.6760}{375.0410}{11.9552}{\normalfont}{la}
\texto{278.6785}{375.0410}{11.9552}{\normalfont}{disponibilidad}
\texto{354.5151}{375.0410}{11.9552}{\normalfont}{es}
\texto{368.2325}{375.0410}{11.9552}{\normalfont}{manual}
\texto{409.5281}{375.0410}{11.9552}{\normalfont}{y}
\texto{419.6159}{375.0410}{11.9552}{\normalfont}{propensa}
\texto{469.4261}{375.0410}{11.9552}{\normalfont}{a}
\texto{479.1887}{375.0410}{11.9552}{\normalfont}{errores.}
\end{pagina}
% Página original 50
\begin{pagina}{50}
\texto{70.8660}{746.3890}{11.9552}{\normalfont}{40}
\texto{168.2490}{746.3890}{11.9552}{\normalfont}{CAPÍTULO}
\texto{235.2125}{746.3890}{11.9552}{\normalfont}{4.}
\texto{252.1243}{746.3890}{11.9552}{\normalfont}{MODELADO}
\texto{326.8921}{746.3890}{11.9552}{\normalfont}{Y}
\texto{339.5646}{746.3890}{11.9552}{\normalfont}{AUTOMATIZACIÓN}
\texto{455.6233}{746.3890}{11.9552}{\normalfont}{DEL}
\texto{483.7396}{746.3890}{11.9552}{\normalfont}{PROCESO}
\texto{70.8660}{705.1940}{17.2154}{\normalfont\bfseries}{4.4.}
\texto{120.3603}{705.1940}{17.2154}{\normalfont\bfseries}{Identificación}
\texto{239.9591}{705.1940}{17.2154}{\normalfont\bfseries}{de}
\texto{266.0112}{705.1940}{17.2154}{\normalfont\bfseries}{oportunidades}
\texto{392.8973}{705.1940}{17.2154}{\normalfont\bfseries}{de}
\texto{418.9493}{705.1940}{17.2154}{\normalfont\bfseries}{mejora}
\texto{88.7990}{678.9090}{11.9552}{\normalfont}{A}
\texto{103.2170}{678.9090}{11.9552}{\normalfont}{partir}
\texto{138.1262}{678.9090}{11.9552}{\normalfont}{del}
\texto{158.7393}{678.9090}{11.9552}{\normalfont}{análisis}
\texto{201.5939}{678.9090}{11.9552}{\normalfont}{del}
\texto{222.1951}{678.9090}{11.9552}{\normalfont}{proceso}
\texto{265.9583}{678.9090}{11.9552}{\normalfont}{actual}
\texto{302.8306}{678.9090}{11.9552}{\normalfont}{se}
\texto{318.2934}{678.9090}{11.9552}{\normalfont}{identifican}
\texto{376.9432}{678.9090}{11.9552}{\normalfont}{las}
\texto{396.3202}{678.9090}{11.9552}{\normalfont}{siguientes}
\texto{451.1970}{678.9090}{11.9552}{\normalfont}{oportunidades}
\texto{529.4294}{678.9090}{11.9552}{\normalfont}{de}
\texto{70.8660}{664.4630}{11.9552}{\normalfont}{mejora:}
\texto{100.7540}{637.0660}{11.9552}{\normalfont\bfseries}{Digitalización}
\texto{184.8600}{637.0660}{11.9552}{\normalfont\bfseries}{del}
\texto{206.1379}{637.0660}{11.9552}{\normalfont\bfseries}{catálogo:}
\texto{261.5020}{637.0660}{11.9552}{\normalfont}{sustituir}
\texto{307.3239}{637.0660}{11.9552}{\normalfont}{los}
\texto{324.4772}{637.0660}{11.9552}{\normalfont}{documentos}
\texto{388.4447}{637.0660}{11.9552}{\normalfont}{locales}
\texto{425.4317}{637.0660}{11.9552}{\normalfont}{por}
\texto{446.0951}{637.0660}{11.9552}{\normalfont}{un}
\texto{462.5215}{637.0660}{11.9552}{\normalfont}{repositorio}
\texto{520.3201}{637.0660}{11.9552}{\normalfont}{cen-}
\texto{100.7540}{622.6200}{11.9552}{\normalfont}{tralizado}
\texto{149.5264}{622.6200}{11.9552}{\normalfont}{donde}
\texto{183.9909}{622.6200}{11.9552}{\normalfont}{el}
\texto{196.3550}{622.6200}{11.9552}{\normalfont}{operador}
\texto{245.4502}{622.6200}{11.9552}{\normalfont}{registre}
\texto{287.1451}{622.6200}{11.9552}{\normalfont}{sus}
\texto{306.7804}{622.6200}{11.9552}{\normalfont}{tours,}
\texto{340.0086}{622.6200}{11.9552}{\normalfont}{precios}
\texto{379.1021}{622.6200}{11.9552}{\normalfont}{y}
\texto{389.1780}{622.6200}{11.9552}{\normalfont}{disponibilidad.}
\texto{100.7540}{598.2110}{11.9552}{\normalfont\bfseries}{Canal}
\texto{139.0561}{598.2110}{11.9552}{\normalfont\bfseries}{transaccional}
\texto{220.1123}{598.2110}{11.9552}{\normalfont\bfseries}{en}
\texto{238.0248}{598.2110}{11.9552}{\normalfont\bfseries}{línea:}
\texto{273.1390}{598.2110}{11.9552}{\normalfont}{habilitar}
\texto{320.4553}{598.2110}{11.9552}{\normalfont}{un}
\texto{337.2165}{598.2110}{11.9552}{\normalfont}{canal}
\texto{367.6353}{598.2110}{11.9552}{\normalfont}{web/móvil}
\texto{425.0370}{598.2110}{11.9552}{\normalfont}{donde}
\texto{459.3580}{598.2110}{11.9552}{\normalfont}{el}
\texto{471.5666}{598.2110}{11.9552}{\normalfont}{turista}
\texto{509.2040}{598.2110}{11.9552}{\normalfont}{descu-}
\texto{100.7540}{583.7660}{11.9552}{\normalfont}{bra,}
\texto{123.8682}{583.7660}{11.9552}{\normalfont}{compare}
\texto{169.7451}{583.7660}{11.9552}{\normalfont}{y}
\texto{178.8764}{583.7660}{11.9552}{\normalfont}{reserve}
\texto{217.0159}{583.7660}{11.9552}{\normalfont}{tours}
\texto{246.0479}{583.7660}{11.9552}{\normalfont}{de}
\texto{260.7074}{583.7660}{11.9552}{\normalfont}{forma}
\texto{293.2518}{583.7660}{11.9552}{\normalfont}{autónoma,}
\texto{350.1849}{583.7660}{11.9552}{\normalfont}{sin}
\texto{367.5104}{583.7660}{11.9552}{\normalfont}{depender}
\texto{416.9738}{583.7660}{11.9552}{\normalfont}{de}
\texto{431.6333}{583.7660}{11.9552}{\normalfont}{llamadas}
\texto{479.5258}{583.7660}{11.9552}{\normalfont}{ni}
\texto{492.2342}{583.7660}{11.9552}{\normalfont}{mensajes.}
\texto{100.7540}{559.3570}{11.9552}{\normalfont\bfseries}{Automatización}
\texto{197.5290}{559.3570}{11.9552}{\normalfont\bfseries}{de}
\texto{215.5729}{559.3570}{11.9552}{\normalfont\bfseries}{la}
\texto{230.2826}{559.3570}{11.9552}{\normalfont\bfseries}{disponibilidad:}
\texto{320.3380}{559.3570}{11.9552}{\normalfont}{gestionar}
\texto{370.4279}{559.3570}{11.9552}{\normalfont}{fechas,}
\texto{408.1752}{559.3570}{11.9552}{\normalfont}{horarios}
\texto{453.0742}{559.3570}{11.9552}{\normalfont}{y}
\texto{463.1142}{559.3570}{11.9552}{\normalfont}{cupos}
\texto{495.9790}{559.3570}{11.9552}{\normalfont}{de}
\texto{511.5471}{559.3570}{11.9552}{\normalfont}{forma}
\texto{100.7540}{544.9110}{11.9552}{\normalfont}{sistematizada,}
\texto{176.9708}{544.9110}{11.9552}{\normalfont}{evitando}
\texto{224.7677}{544.9110}{11.9552}{\normalfont}{dobles}
\texto{260.6094}{544.9110}{11.9552}{\normalfont}{reservas}
\texto{304.6380}{544.9110}{11.9552}{\normalfont}{y}
\texto{314.7138}{544.9110}{11.9552}{\normalfont}{sobreventa.}
\texto{100.7540}{520.5030}{11.9552}{\normalfont\bfseries}{Pago}
\texto{136.1223}{520.5030}{11.9552}{\normalfont\bfseries}{en}
\texto{156.2943}{520.5030}{11.9552}{\normalfont\bfseries}{línea}
\texto{190.4885}{520.5030}{11.9552}{\normalfont\bfseries}{seguro:}
\texto{237.8060}{520.5030}{11.9552}{\normalfont}{integrar}
\texto{283.5203}{520.5030}{11.9552}{\normalfont}{una}
\texto{308.0954}{520.5030}{11.9552}{\normalfont}{pasarela}
\texto{355.4978}{520.5030}{11.9552}{\normalfont}{de}
\texto{372.9189}{520.5030}{11.9552}{\normalfont}{pago}
\texto{402.6969}{520.5030}{11.9552}{\normalfont}{con}
\texto{425.9832}{520.5030}{11.9552}{\normalfont}{certificación}
\texto{492.8319}{520.5030}{11.9552}{\normalfont}{PCI}
\texto{519.1884}{520.5030}{11.9552}{\normalfont}{DSS}
\texto{100.7540}{506.0570}{11.9552}{\normalfont}{(Stripe)}
\texto{144.6463}{506.0570}{11.9552}{\normalfont}{que}
\texto{166.4287}{506.0570}{11.9552}{\normalfont}{permita}
\texto{210.3330}{506.0570}{11.9552}{\normalfont}{el}
\texto{222.6851}{506.0570}{11.9552}{\normalfont}{cobro}
\texto{254.5600}{506.0570}{11.9552}{\normalfont}{inmediato}
\texto{309.1857}{506.0570}{11.9552}{\normalfont}{y}
\texto{319.2735}{506.0570}{11.9552}{\normalfont}{reduzca}
\texto{362.1927}{506.0570}{11.9552}{\normalfont}{la}
\texto{375.1952}{506.0570}{11.9552}{\normalfont}{fricción}
\texto{416.5004}{506.0570}{11.9552}{\normalfont}{del}
\texto{435.3561}{506.0570}{11.9552}{\normalfont}{proceso.}
\texto{100.7540}{481.6490}{11.9552}{\normalfont\bfseries}{V}
\texto{109.7981}{481.6490}{11.9552}{\normalfont\bfseries}{oucher}
\texto{155.0402}{481.6490}{11.9552}{\normalfont\bfseries}{digital:}
\texto{201.4880}{481.6490}{11.9552}{\normalfont}{generar}
\texto{244.7323}{481.6490}{11.9552}{\normalfont}{un}
\texto{263.2629}{481.6490}{11.9552}{\normalfont}{voucher}
\texto{308.1379}{481.6490}{11.9552}{\normalfont}{digital}
\texto{346.1913}{481.6490}{11.9552}{\normalfont}{con}
\texto{369.2744}{481.6490}{11.9552}{\normalfont}{código}
\texto{407.6506}{481.6490}{11.9552}{\normalfont}{QR}
\texto{430.8939}{481.6490}{11.9552}{\normalfont}{al}
\texto{445.5343}{481.6490}{11.9552}{\normalfont}{confirmar}
\texto{499.8348}{481.6490}{11.9552}{\normalfont}{el}
\texto{513.8248}{481.6490}{11.9552}{\normalfont}{pago,}
\texto{100.7540}{467.2030}{11.9552}{\normalfont}{estandarizando}
\texto{181.4588}{467.2030}{11.9552}{\normalfont}{la}
\texto{194.4612}{467.2030}{11.9552}{\normalfont}{validación}
\texto{249.4169}{467.2030}{11.9552}{\normalfont}{del}
\texto{268.2846}{467.2030}{11.9552}{\normalfont}{servicio}
\texto{310.2928}{467.2030}{11.9552}{\normalfont}{y}
\texto{320.3686}{467.2030}{11.9552}{\normalfont}{habilitando}
\texto{382.1603}{467.2030}{11.9552}{\normalfont}{reseñas}
\texto{422.6071}{467.2030}{11.9552}{\normalfont}{verificadas.}
\texto{100.7540}{442.7940}{11.9552}{\normalfont\bfseries}{Supervisión}
\texto{175.3533}{442.7940}{11.9552}{\normalfont\bfseries}{y}
\texto{188.5614}{442.7940}{11.9552}{\normalfont\bfseries}{trazabilidad:}
\texto{267.5610}{442.7940}{11.9552}{\normalfont}{incorporar}
\texto{325.8713}{442.7940}{11.9552}{\normalfont}{un}
\texto{344.1986}{442.7940}{11.9552}{\normalfont}{panel}
\texto{376.8339}{442.7940}{11.9552}{\normalfont}{de}
\texto{393.8605}{442.7940}{11.9552}{\normalfont}{administración}
\texto{474.6872}{442.7940}{11.9552}{\normalfont}{que}
\texto{497.8803}{442.7940}{11.9552}{\normalfont}{verifique}
\texto{100.7540}{428.3480}{11.9552}{\normalfont}{operadores,}
\texto{163.5308}{428.3480}{11.9552}{\normalfont}{modere}
\texto{205.4313}{428.3480}{11.9552}{\normalfont}{la}
\texto{219.0316}{428.3480}{11.9552}{\normalfont}{oferta}
\texto{253.1302}{428.3480}{11.9552}{\normalfont}{y}
\texto{263.8158}{428.3480}{11.9552}{\normalfont}{registre}
\texto{306.1085}{428.3480}{11.9552}{\normalfont}{todas}
\texto{338.3182}{428.3480}{11.9552}{\normalfont}{las}
\texto{356.5475}{428.3480}{11.9552}{\normalfont}{acciones}
\texto{402.7424}{428.3480}{11.9552}{\normalfont}{críticas}
\texto{443.7343}{428.3480}{11.9552}{\normalfont}{en}
\texto{459.9360}{428.3480}{11.9552}{\normalfont}{una}
\texto{483.3037}{428.3480}{11.9552}{\normalfont}{bitácora}
\texto{529.4340}{428.3480}{11.9552}{\normalfont}{de}
\texto{100.7540}{413.9030}{11.9552}{\normalfont}{auditoría.}
\texto{100.7540}{389.4940}{11.9552}{\normalfont\bfseries}{Reducción}
\texto{165.7592}{389.4940}{11.9552}{\normalfont\bfseries}{de}
\texto{183.3728}{389.4940}{11.9552}{\normalfont\bfseries}{comisiones:}
\texto{252.9340}{389.4940}{11.9552}{\normalfont}{eliminar}
\texto{298.0362}{389.4940}{11.9552}{\normalfont}{o}
\texto{307.3804}{389.4940}{11.9552}{\normalfont}{reducir}
\texto{346.6412}{389.4940}{11.9552}{\normalfont}{la}
\texto{359.2372}{389.4940}{11.9552}{\normalfont}{dependencia}
\texto{425.4858}{389.4940}{11.9552}{\normalfont}{de}
\texto{440.6832}{389.4940}{11.9552}{\normalfont}{plataformas}
\texto{504.3853}{389.4940}{11.9552}{\normalfont}{extran-}
\texto{100.7540}{375.0480}{11.9552}{\normalfont}{jeras}
\texto{128.4542}{375.0480}{11.9552}{\normalfont}{mediante}
\texto{178.8549}{375.0480}{11.9552}{\normalfont}{una}
\texto{201.6128}{375.0480}{11.9552}{\normalfont}{comisión}
\texto{249.8114}{375.0480}{11.9552}{\normalfont}{propia}
\texto{286.2270}{375.0480}{11.9552}{\normalfont}{reducida}
\texto{333.7106}{375.0480}{11.9552}{\normalfont}{del}
\texto{352.5664}{375.0480}{11.9552}{\normalfont}{8}
\texto{360.4162}{375.0480}{11.9552}{\normalfont}{\%.}
\end{pagina}
% Página original 51
\begin{pagina}{51}
\texto{70.8660}{746.3890}{11.9552}{\normalfont}{4.5.}
\texto{96.8710}{746.3890}{11.9552}{\normalfont}{PROPUESTA}
\texto{173.5994}{746.3890}{11.9552}{\normalfont}{DE}
\texto{194.4015}{746.3890}{11.9552}{\normalfont}{AUTOMATIZACIÓN}
\texto{310.4602}{746.3890}{11.9552}{\normalfont}{MEDIANTE}
\texto{380.1805}{746.3890}{11.9552}{\normalfont}{TI}
\texto{529.4280}{746.3890}{11.9552}{\normalfont}{41}
\texto{70.8660}{705.1940}{17.2154}{\normalfont\bfseries}{4.5.}
\texto{120.3603}{705.1940}{17.2154}{\normalfont\bfseries}{Propuesta}
\texto{212.5832}{705.1940}{17.2154}{\normalfont\bfseries}{de}
\texto{238.6352}{705.1940}{17.2154}{\normalfont\bfseries}{automatización}
\texto{373.4008}{705.1940}{17.2154}{\normalfont\bfseries}{mediante}
\texto{456.9644}{705.1940}{17.2154}{\normalfont\bfseries}{TI}
\texto{88.7990}{678.9090}{11.9552}{\normalfont}{La}
\texto{105.6487}{678.9090}{11.9552}{\normalfont}{automatización}
\texto{187.3744}{678.9090}{11.9552}{\normalfont}{del}
\texto{206.0150}{678.9090}{11.9552}{\normalfont}{proceso}
\texto{247.8175}{678.9090}{11.9552}{\normalfont}{se}
\texto{261.3197}{678.9090}{11.9552}{\normalfont}{materializa}
\texto{321.5835}{678.9090}{11.9552}{\normalfont}{en}
\texto{336.9722}{678.9090}{11.9552}{\normalfont}{el}
\texto{349.1211}{678.9090}{11.9552}{\normalfont}{ecosistema}
\texto{407.0550}{678.9090}{11.9552}{\normalfont\bfseries}{Explora-Nicaragua}
\texto{520.9090}{678.9090}{11.9552}{\normalfont}{des-}
\texto{70.8660}{664.4630}{11.9552}{\normalfont}{crito}
\texto{99.4676}{664.4630}{11.9552}{\normalfont}{en}
\texto{116.3627}{664.4630}{11.9552}{\normalfont}{los}
\texto{135.2734}{664.4630}{11.9552}{\normalfont}{capítulos}
\texto{186.0520}{664.4630}{11.9552}{\normalfont}{anteriores.}
\texto{244.3240}{664.4630}{11.9552}{\normalfont}{En}
\texto{263.9807}{664.4630}{11.9552}{\normalfont}{la}
\texto{278.2744}{664.4630}{11.9552}{\normalfont}{visión}
\texto{313.1190}{664.4630}{11.9552}{\normalfont}{completa,}
\texto{367.7351}{664.4630}{11.9552}{\normalfont}{el}
\texto{381.3784}{664.4630}{11.9552}{\normalfont}{ecosistema}
\texto{440.6882}{664.4630}{11.9552}{\normalfont}{contempla}
\texto{498.2333}{664.4630}{11.9552}{\normalfont}{una}
\texto{522.2824}{664.4630}{11.9552}{\normalfont}{app}
\texto{70.8660}{650.0170}{11.9552}{\normalfont}{móvil}
\texto{102.8270}{650.0170}{11.9552}{\normalfont}{para}
\texto{129.5947}{650.0170}{11.9552}{\normalfont}{el}
\texto{142.0544}{650.0170}{11.9552}{\normalfont}{turista,}
\texto{183.1827}{650.0170}{11.9552}{\normalfont}{un}
\texto{200.1950}{650.0170}{11.9552}{\normalfont}{backoffice}
\texto{253.3071}{650.0170}{11.9552}{\normalfont}{de}
\texto{269.0067}{650.0170}{11.9552}{\normalfont}{escritorio}
\texto{319.9024}{650.0170}{11.9552}{\normalfont}{para}
\texto{346.6701}{650.0170}{11.9552}{\normalfont}{el}
\texto{359.1298}{650.0170}{11.9552}{\normalfont}{operador}
\texto{408.3326}{650.0170}{11.9552}{\normalfont}{y}
\texto{418.5041}{650.0170}{11.9552}{\normalfont}{un}
\texto{435.5163}{650.0170}{11.9552}{\normalfont}{panel}
\texto{466.8366}{650.0170}{11.9552}{\normalfont}{de}
\texto{482.5481}{650.0170}{11.9552}{\normalfont}{administra-}
\texto{70.8660}{635.5710}{11.9552}{\normalfont}{ción}
\texto{95.1805}{635.5710}{11.9552}{\normalfont}{para}
\texto{121.4461}{635.5710}{11.9552}{\normalfont}{el}
\texto{133.4156}{635.5710}{11.9552}{\normalfont}{equipo}
\texto{170.7350}{635.5710}{11.9552}{\normalfont}{gestor,}
\texto{208.1332}{635.5710}{11.9552}{\normalfont}{todos}
\texto{239.3387}{635.5710}{11.9552}{\normalfont}{integrados}
\texto{295.2627}{635.5710}{11.9552}{\normalfont}{sobre}
\texto{325.5070}{635.5710}{11.9552}{\normalfont}{una}
\texto{347.8704}{635.5710}{11.9552}{\normalfont}{API}
\texto{372.3475}{635.5710}{11.9552}{\normalfont}{REST}
\texto{407.3881}{635.5710}{11.9552}{\normalfont}{común}
\texto{444.3871}{635.5710}{11.9552}{\normalfont}{[14].}
\texto{469.3590}{635.5710}{11.9552}{\normalfont\bfseries}{El}
\texto{485.9516}{635.5710}{11.9552}{\normalfont\bfseries}{prototipo}
\texto{70.8660}{621.1260}{11.9552}{\normalfont\bfseries}{inicial}
\texto{111.3929}{621.1260}{11.9552}{\normalfont\bfseries}{de}
\texto{130.5966}{621.1260}{11.9552}{\normalfont\bfseries}{esta}
\texto{159.4146}{621.1260}{11.9552}{\normalfont\bfseries}{etapa}
\texto{196.9252}{621.1260}{11.9552}{\normalfont\bfseries}{implementa}
\texto{271.0235}{621.1260}{11.9552}{\normalfont\bfseries}{las}
\texto{292.2106}{621.1260}{11.9552}{\normalfont\bfseries}{dos}
\texto{317.3081}{621.1260}{11.9552}{\normalfont\bfseries}{aplicaciones}
\texto{392.2541}{621.1260}{11.9552}{\normalfont\bfseries}{de}
\texto{411.4696}{621.1260}{11.9552}{\normalfont\bfseries}{escritorio}
\texto{470.9020}{621.1260}{11.9552}{\normalfont}{(backoffice}
\texto{529.4275}{621.1260}{11.9552}{\normalfont}{de}
\texto{70.8660}{606.6800}{11.9552}{\normalfont}{operadores}
\texto{130.2355}{606.6800}{11.9552}{\normalfont}{y}
\texto{140.7657}{606.6800}{11.9552}{\normalfont}{panel}
\texto{172.4326}{606.6800}{11.9552}{\normalfont}{de}
\texto{188.4908}{606.6800}{11.9552}{\normalfont}{administración),}
\texto{276.1535}{606.6800}{11.9552}{\normalfont}{consumiendo}
\texto{346.2564}{606.6800}{11.9552}{\normalfont}{la}
\texto{359.7132}{606.6800}{11.9552}{\normalfont}{misma}
\texto{397.2980}{606.6800}{11.9552}{\normalfont}{API}
\texto{422.5999}{606.6800}{11.9552}{\normalfont}{REST}
\texto{458.4894}{606.6800}{11.9552}{\normalfont}{que}
\texto{480.7261}{606.6800}{11.9552}{\normalfont}{en}
\texto{496.7843}{606.6800}{11.9552}{\normalfont}{el}
\texto{509.5908}{606.6800}{11.9552}{\normalfont}{futuro}
\texto{70.8660}{592.2340}{11.9552}{\normalfont}{alimentará}
\texto{129.0711}{592.2340}{11.9552}{\normalfont}{también}
\texto{174.2689}{592.2340}{11.9552}{\normalfont}{la}
\texto{187.2714}{592.2340}{11.9552}{\normalfont}{app}
\texto{210.0413}{592.2340}{11.9552}{\normalfont}{móvil.}
\texto{88.7990}{577.7880}{11.9552}{\normalfont}{Las}
\texto{110.4810}{577.7880}{11.9552}{\normalfont}{principales}
\texto{169.0734}{577.7880}{11.9552}{\normalfont}{acciones}
\texto{214.6705}{577.7880}{11.9552}{\normalfont}{de}
\texto{230.2744}{577.7880}{11.9552}{\normalfont}{automatización}
\texto{312.2154}{577.7880}{11.9552}{\normalfont}{son:}
\texto{100.7540}{550.3910}{11.9552}{\normalfont\bfseries}{Registro}
\texto{154.2021}{550.3910}{11.9552}{\normalfont\bfseries}{y}
\texto{165.3659}{550.3910}{11.9552}{\normalfont\bfseries}{verificación}
\texto{235.4234}{550.3910}{11.9552}{\normalfont\bfseries}{de}
\texto{253.0967}{550.3910}{11.9552}{\normalfont\bfseries}{operadores:}
\texto{324.2480}{550.3910}{11.9552}{\normalfont}{el}
\texto{336.2415}{550.3910}{11.9552}{\normalfont}{operador}
\texto{384.9780}{550.3910}{11.9552}{\normalfont}{se}
\texto{398.3368}{550.3910}{11.9552}{\normalfont}{registra}
\texto{440.3115}{550.3910}{11.9552}{\normalfont}{y}
\texto{450.0287}{550.3910}{11.9552}{\normalfont}{sube}
\texto{476.7174}{550.3910}{11.9552}{\normalfont}{su}
\texto{491.3769}{550.3910}{11.9552}{\normalfont}{documen-}
\texto{100.7540}{535.9450}{11.9552}{\normalfont}{tación;}
\texto{139.1206}{535.9450}{11.9552}{\normalfont}{el}
\texto{151.4727}{535.9450}{11.9552}{\normalfont}{panel}
\texto{182.6973}{535.9450}{11.9552}{\normalfont}{de}
\texto{198.3013}{535.9450}{11.9552}{\normalfont}{administración}
\texto{277.7053}{535.9450}{11.9552}{\normalfont}{de}
\texto{293.3212}{535.9450}{11.9552}{\normalfont}{escritorio}
\texto{344.1093}{535.9450}{11.9552}{\normalfont}{aprueba}
\texto{388.9915}{535.9450}{11.9552}{\normalfont}{o}
\texto{398.7421}{535.9450}{11.9552}{\normalfont}{rechaza}
\texto{440.6882}{535.9450}{11.9552}{\normalfont}{con}
\texto{462.1573}{535.9450}{11.9552}{\normalfont}{trazabilidad.}
\texto{100.7540}{511.5360}{11.9552}{\normalfont\bfseries}{Gestión}
\texto{151.6939}{511.5360}{11.9552}{\normalfont\bfseries}{centralizada}
\texto{227.9442}{511.5360}{11.9552}{\normalfont\bfseries}{del}
\texto{251.0631}{511.5360}{11.9552}{\normalfont\bfseries}{catálogo:}
\texto{308.0210}{511.5370}{11.9552}{\normalfont}{el}
\texto{321.4969}{511.5370}{11.9552}{\normalfont}{backoffice}
\texto{375.6253}{511.5370}{11.9552}{\normalfont}{de}
\texto{392.3530}{511.5370}{11.9552}{\normalfont}{escritorio}
\texto{444.2649}{511.5370}{11.9552}{\normalfont}{permite}
\texto{488.6306}{511.5370}{11.9552}{\normalfont}{crear,}
\texto{522.2797}{511.5370}{11.9552}{\normalfont}{edi-}
\texto{100.7540}{497.0910}{11.9552}{\normalfont}{tar,}
\texto{123.5430}{497.0910}{11.9552}{\normalfont}{publicar}
\texto{169.7451}{497.0910}{11.9552}{\normalfont}{y}
\texto{180.5024}{497.0910}{11.9552}{\normalfont}{despublicar}
\texto{243.0281}{497.0910}{11.9552}{\normalfont}{tours}
\texto{273.6860}{497.0910}{11.9552}{\normalfont}{con}
\texto{295.8246}{497.0910}{11.9552}{\normalfont}{descripciones,}
\texto{370.4442}{497.0910}{11.9552}{\normalfont}{fotografías,}
\texto{431.6692}{497.0910}{11.9552}{\normalfont}{precios}
\texto{471.4321}{497.0910}{11.9552}{\normalfont}{y}
\texto{482.1894}{497.0910}{11.9552}{\normalfont}{políticas}
\texto{529.4292}{497.0910}{11.9552}{\normalfont}{de}
\texto{100.7540}{482.6450}{11.9552}{\normalfont}{cancelación.}
\texto{100.7540}{458.2360}{11.9552}{\normalfont\bfseries}{Disponibilidad}
\texto{191.4701}{458.2360}{11.9552}{\normalfont\bfseries}{en}
\texto{210.5063}{458.2360}{11.9552}{\normalfont\bfseries}{tiempo}
\texto{256.8244}{458.2360}{11.9552}{\normalfont\bfseries}{real:}
\texto{287.1940}{458.2360}{11.9552}{\normalfont}{el}
\texto{300.3710}{458.2360}{11.9552}{\normalfont}{sistema}
\texto{342.9435}{458.2360}{11.9552}{\normalfont}{gestiona}
\texto{389.3536}{458.2360}{11.9552}{\normalfont}{fechas,}
\texto{427.9737}{458.2360}{11.9552}{\normalfont}{horarios}
\texto{473.7334}{458.2360}{11.9552}{\normalfont}{y}
\texto{484.6341}{458.2360}{11.9552}{\normalfont}{cupos,}
\texto{521.6235}{458.2360}{11.9552}{\normalfont}{blo-}
\texto{100.7540}{443.7910}{11.9552}{\normalfont}{queando}
\texto{147.2502}{443.7910}{11.9552}{\normalfont}{automáticamente}
\texto{239.2741}{443.7910}{11.9552}{\normalfont}{la}
\texto{252.2766}{443.7910}{11.9552}{\normalfont}{reserva}
\texto{291.6881}{443.7910}{11.9552}{\normalfont}{cuando}
\texto{332.0178}{443.7910}{11.9552}{\normalfont}{se}
\texto{345.7352}{443.7910}{11.9552}{\normalfont}{agota}
\texto{377.6101}{443.7910}{11.9552}{\normalfont}{el}
\texto{389.9622}{443.7910}{11.9552}{\normalfont}{cupo.}
\texto{100.7540}{419.3820}{11.9552}{\normalfont\bfseries}{Reserva}
\texto{150.4601}{419.3820}{11.9552}{\normalfont\bfseries}{y}
\texto{161.2892}{419.3820}{11.9552}{\normalfont\bfseries}{pago}
\texto{192.4671}{419.3820}{11.9552}{\normalfont\bfseries}{en}
\texto{209.8057}{419.3820}{11.9552}{\normalfont\bfseries}{línea:}
\texto{244.4070}{419.3820}{11.9552}{\normalfont}{en}
\texto{259.3534}{419.3820}{11.9552}{\normalfont}{la}
\texto{271.6983}{419.3820}{11.9552}{\normalfont}{visión}
\texto{304.5943}{419.3820}{11.9552}{\normalfont}{completa,}
\texto{357.2617}{419.3820}{11.9552}{\normalfont}{la}
\texto{369.6066}{419.3820}{11.9552}{\normalfont}{app}
\texto{391.7070}{419.3820}{11.9552}{\normalfont}{móvil}
\texto{422.9149}{419.3820}{11.9552}{\normalfont}{consume}
\texto{469.7936}{419.3820}{11.9552}{\normalfont}{la}
\texto{482.1385}{419.3820}{11.9552}{\normalfont}{API}
\texto{506.3406}{419.3820}{11.9552}{\normalfont}{REST,}
\texto{100.7540}{404.9360}{11.9552}{\normalfont}{procesa}
\texto{142.7120}{404.9360}{11.9552}{\normalfont}{el}
\texto{155.0163}{404.9360}{11.9552}{\normalfont}{pago}
\texto{182.9412}{404.9360}{11.9552}{\normalfont}{mediante}
\texto{233.2941}{404.9360}{11.9552}{\normalfont}{Stripe}
\texto{268.0335}{404.9360}{11.9552}{\normalfont}{y}
\texto{278.0735}{404.9360}{11.9552}{\normalfont}{genera}
\texto{315.0916}{404.9360}{11.9552}{\normalfont}{el}
\texto{327.3959}{404.9360}{11.9552}{\normalfont}{voucher}
\texto{370.5972}{404.9360}{11.9552}{\normalfont}{digital;}
\texto{410.2287}{404.9360}{11.9552}{\normalfont}{en}
\texto{425.7848}{404.9360}{11.9552}{\normalfont}{el}
\texto{438.0891}{404.9360}{11.9552}{\normalfont}{prototipo}
\texto{489.7499}{404.9360}{11.9552}{\normalfont}{inicial,}
\texto{527.4184}{404.9360}{11.9552}{\normalfont}{las}
\texto{100.7540}{390.4900}{11.9552}{\normalfont}{reservas}
\texto{144.7707}{390.4900}{11.9552}{\normalfont}{se}
\texto{158.5000}{390.4900}{11.9552}{\normalfont}{administran}
\texto{223.5961}{390.4900}{11.9552}{\normalfont}{desde}
\texto{255.5356}{390.4900}{11.9552}{\normalfont}{el}
\texto{267.8877}{390.4900}{11.9552}{\normalfont}{backoffice}
\texto{320.8923}{390.4900}{11.9552}{\normalfont}{de}
\texto{336.4962}{390.4900}{11.9552}{\normalfont}{escritorio}
\texto{387.2962}{390.4900}{11.9552}{\normalfont}{y}
\texto{397.3721}{390.4900}{11.9552}{\normalfont}{el}
\texto{409.7361}{390.4900}{11.9552}{\normalfont}{panel}
\texto{440.9488}{390.4900}{11.9552}{\normalfont}{de}
\texto{456.5527}{390.4900}{11.9552}{\normalfont}{administración.}
\texto{100.7540}{366.0820}{11.9552}{\normalfont\bfseries}{V}
\texto{109.7981}{366.0820}{11.9552}{\normalfont\bfseries}{alidación}
\texto{165.7102}{366.0820}{11.9552}{\normalfont\bfseries}{del}
\texto{187.0478}{366.0820}{11.9552}{\normalfont\bfseries}{voucher:}
\texto{239.8950}{366.0820}{11.9552}{\normalfont}{el}
\texto{251.8167}{366.0820}{11.9552}{\normalfont}{operador}
\texto{300.4935}{366.0820}{11.9552}{\normalfont}{valida}
\texto{334.2191}{366.0820}{11.9552}{\normalfont}{el}
\texto{346.1409}{366.0820}{11.9552}{\normalfont}{código}
\texto{382.4608}{366.0820}{11.9552}{\normalfont}{QR}
\texto{403.6597}{366.0820}{11.9552}{\normalfont}{del}
\texto{422.0970}{366.0820}{11.9552}{\normalfont}{voucher}
\texto{464.9158}{366.0820}{11.9552}{\normalfont}{en}
\texto{480.1013}{366.0820}{11.9552}{\normalfont}{el}
\texto{492.0350}{366.0820}{11.9552}{\normalfont}{backoffice}
\texto{100.7540}{351.6360}{11.9552}{\normalfont}{de}
\texto{116.3579}{351.6360}{11.9552}{\normalfont}{escritorio}
\texto{167.1460}{351.6360}{11.9552}{\normalfont}{y}
\texto{177.2338}{351.6360}{11.9552}{\normalfont}{marca}
\texto{212.3486}{351.6360}{11.9552}{\normalfont}{la}
\texto{225.3511}{351.6360}{11.9552}{\normalfont}{reserva}
\texto{264.7626}{351.6360}{11.9552}{\normalfont}{como}
\texto{295.3368}{351.6360}{11.9552}{\normalfont}{completada.}
\texto{100.7540}{327.2280}{11.9552}{\normalfont\bfseries}{Reseñas}
\texto{152.8213}{327.2280}{11.9552}{\normalfont\bfseries}{verificadas:}
\texto{222.4450}{327.2280}{11.9552}{\normalfont}{solo}
\texto{246.4917}{327.2280}{11.9552}{\normalfont}{los}
\texto{264.6971}{327.2280}{11.9552}{\normalfont}{turistas}
\texto{307.6808}{327.2280}{11.9552}{\normalfont}{que}
\texto{330.0490}{327.2280}{11.9552}{\normalfont}{completaron}
\texto{397.6174}{327.2280}{11.9552}{\normalfont}{la}
\texto{411.2057}{327.2280}{11.9552}{\normalfont}{reserva}
\texto{451.2030}{327.2280}{11.9552}{\normalfont}{pueden}
\texto{492.1065}{327.2280}{11.9552}{\normalfont}{dejar}
\texto{522.2790}{327.2280}{11.9552}{\normalfont}{una}
\texto{100.7540}{312.7820}{11.9552}{\normalfont}{reseña,}
\texto{139.8355}{312.7820}{11.9552}{\normalfont}{la}
\texto{152.8380}{312.7820}{11.9552}{\normalfont}{cual}
\texto{177.5590}{312.7820}{11.9552}{\normalfont}{es}
\texto{191.2764}{312.7820}{11.9552}{\normalfont}{moderada}
\texto{245.5865}{312.7820}{11.9552}{\normalfont}{por}
\texto{266.7161}{312.7820}{11.9552}{\normalfont}{el}
\texto{279.0802}{312.7820}{11.9552}{\normalfont}{equipo}
\texto{316.7940}{312.7820}{11.9552}{\normalfont}{gestor}
\texto{351.3350}{312.7820}{11.9552}{\normalfont}{desde}
\texto{383.2625}{312.7820}{11.9552}{\normalfont}{el}
\texto{395.6147}{312.7820}{11.9552}{\normalfont}{panel}
\texto{426.8393}{312.7820}{11.9552}{\normalfont}{de}
\texto{442.4432}{312.7820}{11.9552}{\normalfont}{administración.}
\texto{100.7540}{288.3730}{11.9552}{\normalfont\bfseries}{Reportes}
\texto{159.8844}{288.3730}{11.9552}{\normalfont\bfseries}{financieros:}
\texto{231.3740}{288.3730}{11.9552}{\normalfont}{el}
\texto{245.2444}{288.3730}{11.9552}{\normalfont}{sistema}
\texto{288.5103}{288.3730}{11.9552}{\normalfont}{calcula}
\texto{329.0575}{288.3730}{11.9552}{\normalfont}{automáticamente}
\texto{422.5998}{288.3730}{11.9552}{\normalfont}{monto}
\texto{460.2109}{288.3730}{11.9552}{\normalfont}{bruto,}
\texto{496.8440}{288.3730}{11.9552}{\normalfont}{comisión}
\texto{100.7540}{273.9280}{11.9552}{\normalfont}{(8}
\texto{113.1444}{273.9280}{11.9552}{\normalfont}{\%)}
\texto{131.3617}{273.9280}{11.9552}{\normalfont}{y}
\texto{141.4375}{273.9280}{11.9552}{\normalfont}{monto}
\texto{177.5303}{273.9280}{11.9552}{\normalfont}{neto}
\texto{203.5520}{273.9280}{11.9552}{\normalfont}{por}
\texto{224.6816}{273.9280}{11.9552}{\normalfont}{reserva,}
\texto{267.3449}{273.9280}{11.9552}{\normalfont}{con}
\texto{288.8021}{273.9280}{11.9552}{\normalfont}{exportación}
\texto{352.5425}{273.9280}{11.9552}{\normalfont}{de}
\texto{368.1464}{273.9280}{11.9552}{\normalfont}{reportes.}
\texto{100.7540}{249.5190}{11.9552}{\normalfont\bfseries}{Seguridad}
\texto{164.2564}{249.5190}{11.9552}{\normalfont\bfseries}{y}
\texto{176.3527}{249.5190}{11.9552}{\normalfont\bfseries}{control}
\texto{222.3264}{249.5190}{11.9552}{\normalfont\bfseries}{de}
\texto{240.9323}{249.5190}{11.9552}{\normalfont\bfseries}{acceso:}
\texto{285.6740}{249.5190}{11.9552}{\normalfont}{autenticación}
\texto{357.9790}{249.5190}{11.9552}{\normalfont}{JWT}
\texto{388.8139}{249.5190}{11.9552}{\normalfont}{y}
\texto{399.3440}{249.5190}{11.9552}{\normalfont}{RBAC}
\texto{437.4955}{249.5190}{11.9552}{\normalfont}{garantizan}
\texto{495.4925}{249.5190}{11.9552}{\normalfont}{que}
\texto{517.7173}{249.5190}{11.9552}{\normalfont}{cada}
\texto{100.7540}{235.0730}{11.9552}{\normalfont}{actor}
\texto{130.6659}{235.0730}{11.9552}{\normalfont}{acceda}
\texto{168.3822}{235.0730}{11.9552}{\normalfont}{únicamente}
\texto{230.4894}{235.0730}{11.9552}{\normalfont}{a}
\texto{240.2521}{235.0730}{11.9552}{\normalfont}{las}
\texto{257.8716}{235.0730}{11.9552}{\normalfont}{funcionalidades}
\texto{340.2143}{235.0730}{11.9552}{\normalfont}{autorizadas.}
\end{pagina}
% Página original 52
\begin{pagina}{52}
\texto{70.8660}{746.3890}{11.9552}{\normalfont}{42}
\texto{168.2490}{746.3890}{11.9552}{\normalfont}{CAPÍTULO}
\texto{235.2125}{746.3890}{11.9552}{\normalfont}{4.}
\texto{252.1243}{746.3890}{11.9552}{\normalfont}{MODELADO}
\texto{326.8921}{746.3890}{11.9552}{\normalfont}{Y}
\texto{339.5646}{746.3890}{11.9552}{\normalfont}{AUTOMATIZACIÓN}
\texto{455.6233}{746.3890}{11.9552}{\normalfont}{DEL}
\texto{483.7396}{746.3890}{11.9552}{\normalfont}{PROCESO}
\texto{70.8660}{705.1940}{17.2154}{\normalfont\bfseries}{4.6.}
\texto{120.3603}{705.1940}{17.2154}{\normalfont\bfseries}{BPMN}
\texto{187.3282}{705.1940}{17.2154}{\normalfont\bfseries}{del}
\texto{218.7601}{705.1940}{17.2154}{\normalfont\bfseries}{proceso}
\texto{288.8715}{705.1940}{17.2154}{\normalfont\bfseries}{propuesto}
\texto{88.7990}{678.9090}{11.9552}{\normalfont}{La}
\texto{106.0193}{678.9090}{11.9552}{\normalfont}{figura}
\texto{139.3510}{678.9090}{11.9552}{\normalfont}{4.2}
\texto{158.3680}{678.9090}{11.9552}{\normalfont}{presenta}
\texto{205.0865}{678.9090}{11.9552}{\normalfont}{el}
\texto{217.5941}{678.9090}{11.9552}{\normalfont}{diagrama}
\texto{269.1353}{678.9090}{11.9552}{\normalfont}{BPMN}
\texto{308.9557}{678.9090}{11.9552}{\normalfont}{del}
\texto{327.9669}{678.9090}{11.9552}{\normalfont}{proceso}
\texto{370.1400}{678.9090}{11.9552}{\normalfont}{propuesto,}
\texto{427.5991}{678.9090}{11.9552}{\normalfont}{mostrando}
\texto{485.6966}{678.9090}{11.9552}{\normalfont}{la}
\texto{498.8664}{678.9090}{11.9552}{\normalfont}{automa-}
\texto{70.8660}{664.4630}{11.9552}{\normalfont}{tización}
\texto{114.4355}{664.4630}{11.9552}{\normalfont}{de}
\texto{130.0395}{664.4630}{11.9552}{\normalfont}{las}
\texto{147.6710}{664.4630}{11.9552}{\normalfont}{actividades}
\texto{208.5397}{664.4630}{11.9552}{\normalfont}{y}
\texto{218.6155}{664.4630}{11.9552}{\normalfont}{la}
\texto{231.6300}{664.4630}{11.9552}{\normalfont}{integración}
\texto{291.7861}{664.4630}{11.9552}{\normalfont}{entre}
\texto{321.3753}{664.4630}{11.9552}{\normalfont}{los}
\texto{339.0068}{664.4630}{11.9552}{\normalfont}{tres}
\texto{361.8293}{664.4630}{11.9552}{\normalfont}{actores.}
\texto{149.0230}{614.7020}{7.9701}{\normalfont}{Turista}
\texto{140.6990}{605.2370}{7.9701}{\normalfont}{(app}
\texto{160.4577}{605.2370}{7.9701}{\normalfont}{móvil)}
\texto{225.6690}{609.0730}{7.9701}{\normalfont}{Explora}
\texto{257.5454}{609.0730}{7.9701}{\normalfont}{y}
\texto{264.8365}{609.0730}{7.9701}{\normalfont}{se-}
\texto{221.9850}{599.6080}{7.9701}{\normalfont}{lecciona}
\texto{253.9746}{599.6080}{7.9701}{\normalfont}{el}
\texto{262.9115}{599.6080}{7.9701}{\normalfont}{tour}
\texto{317.3210}{609.8480}{7.9701}{\normalfont}{Realiza}
\texto{347.0718}{609.8480}{7.9701}{\normalfont}{el}
\texto{310.7920}{600.3830}{7.9701}{\normalfont}{pago}
\texto{331.0217}{600.3830}{7.9701}{\normalfont}{en}
\texto{342.3114}{600.3830}{7.9701}{\normalfont}{línea}
\texto{398.5970}{605.1160}{7.9701}{\normalfont}{Recibe}
\texto{426.2261}{605.1160}{7.9701}{\normalfont}{voucher}
\texto{457.5184}{605.1160}{7.9701}{\normalfont}{QR}
\texto{147.9680}{523.9920}{7.9701}{\normalfont}{Sistema}
\texto{138.6560}{514.5270}{7.9701}{\normalfont}{(API}
\texto{159.9218}{514.5270}{7.9701}{\normalfont}{REST)}
\texto{185.7960}{519.1380}{7.9701}{\normalfont}{Verifica}
\texto{216.3693}{519.1380}{7.9701}{\normalfont}{dis-}
\texto{186.8780}{509.6730}{7.9701}{\normalfont}{ponibilidad}
\texto{269.9270}{514.4060}{7.9701}{\normalfont}{£Hay}
\texto{293.8987}{514.4060}{7.9701}{\normalfont}{cupo?}
\texto{353.3460}{523.8710}{7.9701}{\normalfont}{Procesa}
\texto{384.7944}{523.8710}{7.9701}{\normalfont}{pago}
\texto{359.6600}{514.4060}{7.9701}{\normalfont}{(Stripe)}
\texto{391.4177}{514.4060}{7.9701}{\normalfont}{y}
\texto{350.1290}{504.9420}{7.9701}{\normalfont}{genera}
\texto{376.9516}{504.9420}{7.9701}{\normalfont}{voucher}
\texto{144.8090}{433.2830}{7.9701}{\normalfont}{Operador}
\texto{141.2870}{423.8180}{7.9701}{\normalfont}{(backoffice)}
\texto{222.8120}{433.1610}{7.9701}{\normalfont}{Valida}
\texto{249.1444}{433.1610}{7.9701}{\normalfont}{voucher}
\texto{235.2750}{423.6960}{7.9701}{\normalfont}{y}
\texto{242.5660}{423.6960}{7.9701}{\normalfont}{marca}
\texto{229.0450}{414.2320}{7.9701}{\normalfont}{completado}
\texto{320.0520}{422.9220}{7.9701}{\normalfont}{Reserva}
\texto{351.5044}{422.9220}{7.9701}{\normalfont}{cerrada}
\texto{332.2730}{520.6770}{5.9776}{\normalfont}{sí}
\texto{109.2310}{366.7580}{11.9552}{\normalfont}{Figura}
\texto{146.7823}{366.7580}{11.9552}{\normalfont}{4.2:}
\texto{168.8980}{366.7580}{11.9552}{\normalfont}{BPMN}
\texto{208.5510}{366.7580}{11.9552}{\normalfont}{del}
\texto{227.4067}{366.7580}{11.9552}{\normalfont}{proceso}
\texto{269.4245}{366.7580}{11.9552}{\normalfont}{propuesto}
\texto{323.4644}{366.7580}{11.9552}{\normalfont}{de}
\texto{339.0803}{366.7580}{11.9552}{\normalfont}{reserva}
\texto{378.4918}{366.7580}{11.9552}{\normalfont}{de}
\texto{394.0957}{366.7580}{11.9552}{\normalfont}{tours}
\texto{424.0722}{366.7580}{11.9552}{\normalfont}{(automatizado)}
\texto{88.7990}{337.7000}{11.9552}{\normalfont}{El}
\texto{102.7053}{337.7000}{11.9552}{\normalfont}{diagrama}
\texto{152.8717}{337.7000}{11.9552}{\normalfont}{muestra}
\texto{196.2643}{337.7000}{11.9552}{\normalfont}{cómo}
\texto{225.6191}{337.7000}{11.9552}{\normalfont}{el}
\texto{236.7757}{337.7000}{11.9552}{\normalfont}{proceso}
\texto{277.5740}{337.7000}{11.9552}{\normalfont}{se}
\texto{290.0839}{337.7000}{11.9552}{\normalfont}{automatiza:}
\texto{353.2576}{337.7000}{11.9552}{\normalfont}{la}
\texto{365.0526}{337.7000}{11.9552}{\normalfont}{disponibilidad}
\texto{439.6817}{337.7000}{11.9552}{\normalfont}{se}
\texto{452.1917}{337.7000}{11.9552}{\normalfont}{verifica}
\texto{491.2923}{337.7000}{11.9552}{\normalfont}{en}
\texto{505.6888}{337.7000}{11.9552}{\normalfont}{tiempo}
\texto{70.8660}{323.2540}{11.9552}{\normalfont}{real,}
\texto{97.6648}{323.2540}{11.9552}{\normalfont}{el}
\texto{110.8059}{323.2540}{11.9552}{\normalfont}{pago}
\texto{139.5678}{323.2540}{11.9552}{\normalfont}{se}
\texto{154.0742}{323.2540}{11.9552}{\normalfont}{procesa}
\texto{196.8810}{323.2540}{11.9552}{\normalfont}{en}
\texto{213.2740}{323.2540}{11.9552}{\normalfont}{línea,}
\texto{245.2876}{323.2540}{11.9552}{\normalfont}{el}
\texto{258.4287}{323.2540}{11.9552}{\normalfont}{voucher}
\texto{302.4669}{323.2540}{11.9552}{\normalfont}{se}
\texto{316.9853}{323.2540}{11.9552}{\normalfont}{genera}
\texto{354.8403}{323.2540}{11.9552}{\normalfont}{automáticamente}
\texto{447.6533}{323.2540}{11.9552}{\normalfont}{y}
\texto{458.5301}{323.2540}{11.9552}{\normalfont}{el}
\texto{471.6713}{323.2540}{11.9552}{\normalfont}{operador}
\texto{521.5675}{323.2540}{11.9552}{\normalfont}{solo}
\texto{70.8660}{308.8090}{11.9552}{\normalfont}{interviene}
\texto{124.8437}{308.8090}{11.9552}{\normalfont}{en}
\texto{140.4477}{308.8090}{11.9552}{\normalfont}{la}
\texto{153.4501}{308.8090}{11.9552}{\normalfont}{validación}
\texto{208.4058}{308.8090}{11.9552}{\normalfont}{del}
\texto{227.2735}{308.8090}{11.9552}{\normalfont}{servicio.}
\end{pagina}
% Página original 53
\begin{pagina}{53}
\texto{70.8660}{746.3890}{11.9552}{\normalfont}{4.7.}
\texto{96.8710}{746.3890}{11.9552}{\normalfont}{BPMN}
\texto{136.5359}{746.3890}{11.9552}{\normalfont}{DE}
\texto{157.3380}{746.3890}{11.9552}{\normalfont}{LOS}
\texto{184.1583}{746.3890}{11.9552}{\normalfont}{PROCESOS}
\texto{251.9610}{746.3890}{11.9552}{\normalfont}{DEL}
\texto{280.0772}{746.3890}{11.9552}{\normalfont}{PROTOTIPO}
\texto{356.6431}{746.3890}{11.9552}{\normalfont}{DE}
\texto{377.4571}{746.3890}{11.9552}{\normalfont}{ESCRITORIO}
\texto{529.4280}{746.3890}{11.9552}{\normalfont}{43}
\texto{70.8660}{705.1940}{17.2154}{\normalfont\bfseries}{4.7.}
\texto{120.3603}{705.1940}{17.2154}{\normalfont\bfseries}{BPMN}
\texto{187.1560}{705.1940}{17.2154}{\normalfont\bfseries}{de}
\texto{213.0532}{705.1940}{17.2154}{\normalfont\bfseries}{los}
\texto{242.0404}{705.1940}{17.2154}{\normalfont\bfseries}{procesos}
\texto{319.6371}{705.1940}{17.2154}{\normalfont\bfseries}{del}
\texto{350.9141}{705.1940}{17.2154}{\normalfont\bfseries}{prototipo}
\texto{436.6726}{705.1940}{17.2154}{\normalfont\bfseries}{de}
\texto{462.5525}{705.1940}{17.2154}{\normalfont\bfseries}{escritorio}
\texto{88.7990}{676.5510}{11.9552}{\normalfont}{Con}
\texto{112.6472}{676.5510}{11.9552}{\normalfont}{el}
\texto{124.1386}{676.5510}{11.9552}{\normalfont}{fin}
\texto{140.1944}{676.5510}{11.9552}{\normalfont}{de}
\texto{154.9376}{676.5510}{11.9552}{\normalfont}{cumplir}
\texto{196.9960}{676.5510}{11.9552}{\normalfont}{el}
\texto{208.4992}{676.5510}{11.9552}{\normalfont}{objetivo}
\texto{252.5087}{676.5510}{11.9552}{\normalfont}{de}
\texto{267.2519}{676.5510}{11.9552}{\normalfont}{modelar}
\texto{311.5961}{676.5510}{11.9552}{\normalfont}{los}
\texto{328.3549}{676.5510}{11.9552}{\normalfont}{procesos}
\texto{374.1290}{676.5510}{11.9552}{\normalfont}{que}
\texto{395.0506}{676.5510}{11.9552}{\normalfont}{el}
\texto{406.5419}{676.5510}{11.9552}{\normalfont}{prototipo}
\texto{457.3898}{676.5510}{11.9552}{\normalfont}{automatiza,}
\texto{520.9101}{676.5510}{11.9552}{\normalfont}{esta}
\texto{70.8660}{662.1050}{11.9552}{\normalfont}{sección}
\texto{110.6936}{662.1050}{11.9552}{\normalfont}{detalla}
\texto{149.1558}{662.1050}{11.9552}{\normalfont}{mediante}
\texto{199.6522}{662.1050}{11.9552}{\normalfont}{BPMN}
\texto{239.4008}{662.1050}{11.9552}{\normalfont}{los}
\texto{257.1280}{662.1050}{11.9552}{\normalfont}{procesos}
\texto{303.8465}{662.1050}{11.9552}{\normalfont}{de}
\texto{319.5461}{662.1050}{11.9552}{\normalfont}{negocio}
\texto{361.5949}{662.1050}{11.9552}{\normalfont}{soportados}
\texto{420.6704}{662.1050}{11.9552}{\normalfont}{por}
\texto{441.8956}{662.1050}{11.9552}{\normalfont}{las}
\texto{459.6228}{662.1050}{11.9552}{\normalfont}{dos}
\texto{480.5898}{662.1050}{11.9552}{\normalfont}{aplicaciones}
\texto{70.8660}{647.6600}{11.9552}{\normalfont}{de}
\texto{87.7133}{647.6600}{11.9552}{\normalfont}{escritorio}
\texto{139.7566}{647.6600}{11.9552}{\normalfont}{del}
\texto{159.8557}{647.6600}{11.9552}{\normalfont}{alcance}
\texto{202.0791}{647.6600}{11.9552}{\normalfont}{inicial:}
\texto{241.0507}{647.6600}{11.9552}{\normalfont}{la}
\texto{255.2965}{647.6600}{11.9552}{\normalfont}{verificación}
\texto{317.6715}{647.6600}{11.9552}{\normalfont}{de}
\texto{334.5308}{647.6600}{11.9552}{\normalfont}{operadores,}
\texto{397.9531}{647.6600}{11.9552}{\normalfont}{la}
\texto{412.1989}{647.6600}{11.9552}{\normalfont}{gestión}
\texto{453.1861}{647.6600}{11.9552}{\normalfont}{del}
\texto{473.2852}{647.6600}{11.9552}{\normalfont}{catálogo}
\texto{520.7115}{647.6600}{11.9552}{\normalfont}{y}
\texto{532.0426}{647.6600}{11.9552}{\normalfont}{la}
\texto{70.8660}{633.2140}{11.9552}{\normalfont}{disponibilidad,}
\texto{148.8546}{633.2140}{11.9552}{\normalfont}{la}
\texto{160.7571}{633.2140}{11.9552}{\normalfont}{reserva}
\texto{199.0807}{633.2140}{11.9552}{\normalfont}{y}
\texto{208.0567}{633.2140}{11.9552}{\normalfont}{validación}
\texto{261.9244}{633.2140}{11.9552}{\normalfont}{de}
\texto{276.4285}{633.2140}{11.9552}{\normalfont}{vouchers,}
\texto{326.4586}{633.2140}{11.9552}{\normalfont}{y}
\texto{335.4346}{633.2140}{11.9552}{\normalfont}{la}
\texto{347.3491}{633.2140}{11.9552}{\normalfont}{moderación}
\texto{409.0140}{633.2140}{11.9552}{\normalfont}{de}
\texto{423.5301}{633.2140}{11.9552}{\normalfont}{reseñas}
\texto{462.8770}{633.2140}{11.9552}{\normalfont}{con}
\texto{483.2463}{633.2140}{11.9552}{\normalfont}{los}
\texto{499.7779}{633.2140}{11.9552}{\normalfont}{reportes}
\texto{70.8660}{618.7680}{11.9552}{\normalfont}{y}
\texto{80.6310}{618.7680}{11.9552}{\normalfont}{la}
\texto{93.3346}{618.7680}{11.9552}{\normalfont}{auditoría.}
\texto{146.3607}{618.7680}{11.9552}{\normalfont}{Cada}
\texto{176.6121}{618.7680}{11.9552}{\normalfont}{diagrama}
\texto{227.6871}{618.7680}{11.9552}{\normalfont}{identifica}
\texto{277.7770}{618.7680}{11.9552}{\normalfont}{las}
\texto{295.0977}{618.7680}{11.9552}{\normalfont}{actividades}
\texto{355.6676}{618.7680}{11.9552}{\normalfont}{que}
\texto{377.1391}{618.7680}{11.9552}{\normalfont}{los}
\texto{394.4598}{618.7680}{11.9552}{\normalfont}{actores}
\texto{433.8928}{618.7680}{11.9552}{\normalfont}{realizan}
\texto{477.1515}{618.7680}{11.9552}{\normalfont}{sobre}
\texto{507.4795}{618.7680}{11.9552}{\normalfont}{la}
\texto{520.1711}{618.7680}{11.9552}{\normalfont}{API}
\texto{70.8660}{604.3220}{11.9552}{\normalfont}{REST}
\texto{106.0980}{604.3220}{11.9552}{\normalfont}{común,}
\texto{146.5281}{604.3220}{11.9552}{\normalfont}{evidenciando}
\texto{216.2460}{604.3220}{11.9552}{\normalfont}{las}
\texto{233.6623}{604.3220}{11.9552}{\normalfont}{oportunidades}
\texto{309.9341}{604.3220}{11.9552}{\normalfont}{de}
\texto{325.3348}{604.3220}{11.9552}{\normalfont}{automatización}
\texto{407.0725}{604.3220}{11.9552}{\normalfont}{cubiertas}
\texto{457.0070}{604.3220}{11.9552}{\normalfont}{por}
\texto{477.9334}{604.3220}{11.9552}{\normalfont}{el}
\texto{490.0823}{604.3220}{11.9552}{\normalfont}{prototipo.}
\texto{70.8660}{562.5100}{14.3462}{\normalfont\bfseries}{4.7.1.}
\texto{124.6642}{562.5100}{14.3462}{\normalfont\bfseries}{BPMN:}
\texto{184.9542}{562.5100}{14.3462}{\normalfont\bfseries}{Alta}
\texto{221.1267}{562.5100}{14.3462}{\normalfont\bfseries}{y}
\texto{235.0253}{562.5100}{14.3462}{\normalfont\bfseries}{verificación}
\texto{319.5818}{562.5100}{14.3462}{\normalfont\bfseries}{de}
\texto{341.2919}{562.5100}{14.3462}{\normalfont\bfseries}{operadores}
\texto{88.7990}{537.9850}{11.9552}{\normalfont}{El}
\texto{103.2313}{537.9850}{11.9552}{\normalfont}{proceso}
\texto{144.5557}{537.9850}{11.9552}{\normalfont}{de}
\texto{159.4781}{537.9850}{11.9552}{\normalfont}{alta}
\texto{182.2050}{537.9850}{11.9552}{\normalfont}{y}
\texto{191.5994}{537.9850}{11.9552}{\normalfont}{verificación}
\texto{252.0496}{537.9850}{11.9552}{\normalfont}{comienza}
\texto{302.1037}{537.9850}{11.9552}{\normalfont}{cuando}
\texto{341.7399}{537.9850}{11.9552}{\normalfont}{el}
\texto{353.4106}{537.9850}{11.9552}{\normalfont}{operador}
\texto{401.8244}{537.9850}{11.9552}{\normalfont}{se}
\texto{414.8603}{537.9850}{11.9552}{\normalfont}{registra}
\texto{456.5242}{537.9850}{11.9552}{\normalfont}{desde}
\texto{487.7703}{537.9850}{11.9552}{\normalfont}{el}
\texto{499.4780}{537.9850}{11.9552}{\normalfont\bfseries}{backof-}
\texto{70.8660}{523.5400}{11.9552}{\normalfont\bfseries}{fice}
\texto{95.4017}{523.5400}{11.9552}{\normalfont\bfseries}{de}
\texto{113.9717}{523.5400}{11.9552}{\normalfont\bfseries}{escritorio}
\texto{172.8520}{523.5400}{11.9552}{\normalfont}{y}
\texto{183.3343}{523.5400}{11.9552}{\normalfont}{culmina}
\texto{227.9726}{523.5400}{11.9552}{\normalfont}{cuando}
\texto{268.7088}{523.5400}{11.9552}{\normalfont}{el}
\texto{281.4793}{523.5400}{11.9552}{\normalfont}{equipo}
\texto{319.6236}{523.5400}{11.9552}{\normalfont}{gestor}
\texto{354.5591}{523.5400}{11.9552}{\normalfont}{aprueba}
\texto{399.8478}{523.5400}{11.9552}{\normalfont}{o}
\texto{410.0169}{523.5400}{11.9552}{\normalfont}{rechaza}
\texto{452.3813}{523.5400}{11.9552}{\normalfont}{la}
\texto{465.8022}{523.5400}{11.9552}{\normalfont}{solicitud}
\texto{513.1066}{523.5400}{11.9552}{\normalfont}{desde}
\texto{70.8660}{509.0940}{11.9552}{\normalfont}{el}
\texto{83.2680}{509.0940}{11.9552}{\normalfont\bfseries}{panel}
\texto{119.0941}{509.0940}{11.9552}{\normalfont\bfseries}{de}
\texto{137.1860}{509.0940}{11.9552}{\normalfont\bfseries}{administración}
\texto{228.3097}{509.0940}{11.9552}{\normalfont\bfseries}{de}
\texto{246.4015}{509.0940}{11.9552}{\normalfont\bfseries}{escritorio}
\texto{300.9710}{509.0940}{11.9552}{\normalfont}{.}
\texto{308.1202}{509.0940}{11.9552}{\normalfont}{La}
\texto{325.1851}{509.0940}{11.9552}{\normalfont}{figura}
\texto{358.3610}{509.0940}{11.9552}{\normalfont}{4.3}
\texto{377.2210}{509.0940}{11.9552}{\normalfont}{lo}
\texto{390.2235}{509.0940}{11.9552}{\normalfont}{modela.}
\texto{124.2730}{441.9710}{7.9701}{\normalfont}{Operador}
\texto{120.7510}{432.5060}{7.9701}{\normalfont}{(backoffice)}
\texto{228.4790}{441.0740}{7.9701}{\normalfont}{Registra}
\texto{214.8720}{431.6090}{7.9701}{\normalfont}{solicitud}
\texto{248.7904}{431.6090}{7.9701}{\normalfont}{y}
\texto{256.0814}{431.6090}{7.9701}{\normalfont}{sube}
\texto{216.0950}{422.1450}{7.9701}{\normalfont}{documentación}
\texto{326.1960}{436.3420}{7.9701}{\normalfont}{Espera}
\texto{354.3512}{436.3420}{7.9701}{\normalfont}{resultado}
\texto{330.8900}{426.8770}{7.9701}{\normalfont}{de}
\texto{342.1796}{426.8770}{7.9701}{\normalfont}{verificación}
\texto{393.4660}{431.6100}{7.9701}{\normalfont}{Cuenta}
\texto{422.8653}{431.6100}{7.9701}{\normalfont}{activa}
\texto{447.5639}{431.6100}{7.9701}{\normalfont}{o}
\texto{454.6198}{431.6100}{7.9701}{\normalfont}{rechazada}
\texto{127.4320}{356.9310}{7.9701}{\normalfont}{Sistema}
\texto{118.1200}{347.4660}{7.9701}{\normalfont}{(API}
\texto{139.3858}{347.4660}{7.9701}{\normalfont}{REST)}
\texto{217.1350}{356.8090}{7.9701}{\normalfont}{Valida}
\texto{243.4674}{356.8090}{7.9701}{\normalfont}{datos}
\texto{266.0961}{356.8090}{7.9701}{\normalfont}{y}
\texto{216.7470}{347.3440}{7.9701}{\normalfont}{registra}
\texto{247.3873}{347.3440}{7.9701}{\normalfont}{estado}
\texto{225.9750}{337.8800}{7.9701}{\normalfont}{pendiente}
\texto{342.7770}{356.8090}{7.9701}{\normalfont}{Notifica}
\texto{336.9860}{347.3440}{7.9701}{\normalfont}{solicitud}
\texto{370.9044}{347.3440}{7.9701}{\normalfont}{al}
\texto{332.5130}{337.8800}{7.9701}{\normalfont}{equipo}
\texto{359.7947}{337.8800}{7.9701}{\normalfont}{gestor}
\texto{424.8770}{356.8090}{7.9701}{\normalfont}{Actualiza}
\texto{411.7950}{347.3440}{7.9701}{\normalfont}{estado}
\texto{438.1872}{347.3440}{7.9701}{\normalfont}{y}
\texto{445.4782}{347.3440}{7.9701}{\normalfont}{notifica}
\texto{421.2210}{337.8800}{7.9701}{\normalfont}{al}
\texto{430.6289}{337.8800}{7.9701}{\normalfont}{operador}
\texto{128.5750}{271.8900}{7.9701}{\normalfont}{Equipo}
\texto{117.7840}{262.4250}{7.9701}{\normalfont}{gestor}
\texto{142.7719}{262.4250}{7.9701}{\normalfont}{(panel}
\texto{128.6320}{252.9610}{7.9701}{\normalfont}{admin)}
\texto{224.3660}{271.7680}{7.9701}{\normalfont}{Revisa}
\texto{251.5767}{271.7680}{7.9701}{\normalfont}{do-}
\texto{220.6820}{262.3030}{7.9701}{\normalfont}{cumentación}
\texto{220.6760}{252.8390}{7.9701}{\normalfont}{del}
\texto{234.3176}{252.8390}{7.9701}{\normalfont}{operador}
\texto{282.7350}{261.6640}{7.9701}{\normalfont}{£Documentación}
\texto{349.1738}{261.6640}{7.9701}{\normalfont}{válida?}
\texto{398.0560}{267.0360}{7.9701}{\normalfont}{Aprueba}
\texto{397.5780}{257.5710}{7.9701}{\normalfont}{operador}
\texto{313.6640}{181.2210}{7.9701}{\normalfont}{Rechaza}
\texto{308.4290}{171.7560}{7.9701}{\normalfont}{con}
\texto{323.9532}{171.7560}{7.9701}{\normalfont}{motivo}
\texto{379.7870}{268.5740}{5.9776}{\normalfont}{sí}
\texto{333.1940}{197.7350}{5.9776}{\normalfont}{no}
\texto{87.1840}{143.2500}{11.9552}{\normalfont}{Figura}
\texto{124.7353}{143.2500}{11.9552}{\normalfont}{4.3:}
\texto{146.8510}{143.2500}{11.9552}{\normalfont}{BPMN}
\texto{186.5040}{143.2500}{11.9552}{\normalfont}{del}
\texto{205.3597}{143.2500}{11.9552}{\normalfont}{alta}
\texto{228.7800}{143.2500}{11.9552}{\normalfont}{y}
\texto{238.8558}{143.2500}{11.9552}{\normalfont}{verificación}
\texto{299.9875}{143.2500}{11.9552}{\normalfont}{de}
\texto{315.6034}{143.2500}{11.9552}{\normalfont}{operadores}
\texto{374.5187}{143.2500}{11.9552}{\normalfont}{en}
\texto{390.1345}{143.2500}{11.9552}{\normalfont}{el}
\texto{402.4867}{143.2500}{11.9552}{\normalfont}{panel}
\texto{433.6993}{143.2500}{11.9552}{\normalfont}{de}
\texto{449.3152}{143.2500}{11.9552}{\normalfont}{administración}
\texto{88.7990}{110.2190}{11.9552}{\normalfont}{El}
\texto{103.4465}{110.2190}{11.9552}{\normalfont}{proceso}
\texto{144.9980}{110.2190}{11.9552}{\normalfont}{automatiza}
\texto{205.6611}{110.2190}{11.9552}{\normalfont}{la}
\texto{218.2093}{110.2190}{11.9552}{\normalfont}{comunicación}
\texto{290.2561}{110.2190}{11.9552}{\normalfont}{entre}
\texto{319.3789}{110.2190}{11.9552}{\normalfont}{el}
\texto{331.2768}{110.2190}{11.9552}{\normalfont}{operador}
\texto{379.9057}{110.2190}{11.9552}{\normalfont}{y}
\texto{389.5273}{110.2190}{11.9552}{\normalfont}{el}
\texto{401.4131}{110.2190}{11.9552}{\normalfont}{equipo}
\texto{438.6727}{110.2190}{11.9552}{\normalfont}{gestor:}
\texto{475.9992}{110.2190}{11.9552}{\normalfont}{la}
\texto{488.5355}{110.2190}{11.9552}{\normalfont}{solicitud}
\texto{534.9671}{110.2190}{11.9552}{\normalfont}{y}
\texto{70.8660}{95.7730}{11.9552}{\normalfont}{la}
\texto{83.9402}{95.7730}{11.9552}{\normalfont}{documentación}
\texto{164.6641}{95.7730}{11.9552}{\normalfont}{viajan}
\texto{200.1854}{95.7730}{11.9552}{\normalfont}{por}
\texto{221.4107}{95.7730}{11.9552}{\normalfont}{la}
\texto{234.4968}{95.7730}{11.9552}{\normalfont}{API}
\texto{259.4282}{95.7730}{11.9552}{\normalfont}{REST,}
\texto{298.1989}{95.7730}{11.9552}{\normalfont}{el}
\texto{310.6347}{95.7730}{11.9552}{\normalfont}{administrador}
\texto{386.2203}{95.7730}{11.9552}{\normalfont}{aprueba}
\texto{431.1742}{95.7730}{11.9552}{\normalfont}{o}
\texto{441.0086}{95.7730}{11.9552}{\normalfont}{rechaza}
\texto{483.0383}{95.7730}{11.9552}{\normalfont}{con}
\texto{504.5791}{95.7730}{11.9552}{\normalfont}{un}
\texto{521.5675}{95.7730}{11.9552}{\normalfont}{solo}
\texto{70.8660}{81.3270}{11.9552}{\normalfont}{clic}
\texto{91.2664}{81.3270}{11.9552}{\normalfont}{y}
\texto{100.9477}{81.3270}{11.9552}{\normalfont}{cada}
\texto{127.8517}{81.3270}{11.9552}{\normalfont}{decisión}
\texto{171.7416}{81.3270}{11.9552}{\normalfont}{queda}
\texto{205.4744}{81.3270}{11.9552}{\normalfont}{registrada}
\texto{259.7701}{81.3270}{11.9552}{\normalfont}{en}
\texto{274.9676}{81.3270}{11.9552}{\normalfont}{la}
\texto{287.5755}{81.3270}{11.9552}{\normalfont}{bitácora}
\texto{332.6897}{81.3270}{11.9552}{\normalfont}{de}
\texto{347.8991}{81.3270}{11.9552}{\normalfont}{auditoría,}
\texto{400.8176}{81.3270}{11.9552}{\normalfont}{con}
\texto{421.8802}{81.3270}{11.9552}{\normalfont}{notificación}
\texto{483.9158}{81.3270}{11.9552}{\normalfont}{automática}
\texto{70.8660}{66.8810}{11.9552}{\normalfont}{al}
\texto{83.8685}{66.8810}{11.9552}{\normalfont}{operador.}
\end{pagina}
% Página original 54
\begin{pagina}{54}
\texto{70.8660}{746.3890}{11.9552}{\normalfont}{44}
\texto{168.2490}{746.3890}{11.9552}{\normalfont}{CAPÍTULO}
\texto{235.2125}{746.3890}{11.9552}{\normalfont}{4.}
\texto{252.1243}{746.3890}{11.9552}{\normalfont}{MODELADO}
\texto{326.8921}{746.3890}{11.9552}{\normalfont}{Y}
\texto{339.5646}{746.3890}{11.9552}{\normalfont}{AUTOMATIZACIÓN}
\texto{455.6233}{746.3890}{11.9552}{\normalfont}{DEL}
\texto{483.7396}{746.3890}{11.9552}{\normalfont}{PROCESO}
\texto{70.8660}{705.1940}{14.3462}{\normalfont\bfseries}{4.7.2.}
\texto{124.6642}{705.1940}{14.3462}{\normalfont\bfseries}{BPMN:}
\texto{184.9542}{705.1940}{14.3462}{\normalfont\bfseries}{Gestión}
\texto{244.5468}{705.1940}{14.3462}{\normalfont\bfseries}{del}
\texto{270.7401}{705.1940}{14.3462}{\normalfont\bfseries}{catálogo}
\texto{333.9538}{705.1940}{14.3462}{\normalfont\bfseries}{y}
\texto{347.8524}{705.1940}{14.3462}{\normalfont\bfseries}{disponibilidad}
\texto{88.7990}{682.3960}{11.9552}{\normalfont}{El}
\texto{103.4585}{682.3960}{11.9552}{\normalfont}{proceso}
\texto{145.0100}{682.3960}{11.9552}{\normalfont}{de}
\texto{160.1715}{682.3960}{11.9552}{\normalfont}{gestión}
\texto{199.4492}{682.3960}{11.9552}{\normalfont}{del}
\texto{217.8626}{682.3960}{11.9552}{\normalfont}{catálogo}
\texto{263.5792}{682.3960}{11.9552}{\normalfont}{cubre}
\texto{294.9879}{682.3960}{11.9552}{\normalfont}{el}
\texto{306.8977}{682.3960}{11.9552}{\normalfont}{ciclo}
\texto{333.1035}{682.3960}{11.9552}{\normalfont}{de}
\texto{348.2531}{682.3960}{11.9552}{\normalfont}{vida}
\texto{373.4954}{682.3960}{11.9552}{\normalfont}{de}
\texto{388.6450}{682.3960}{11.9552}{\normalfont}{un}
\texto{405.1073}{682.3960}{11.9552}{\normalfont}{tour}
\texto{430.0124}{682.3960}{11.9552}{\normalfont}{dentro}
\texto{466.3012}{682.3960}{11.9552}{\normalfont}{del}
\texto{484.7490}{682.3960}{11.9552}{\normalfont\bfseries}{backoffice}
\texto{70.8660}{667.9500}{11.9552}{\normalfont\bfseries}{de}
\texto{88.5035}{667.9500}{11.9552}{\normalfont\bfseries}{escritorio}
\texto{143.0650}{667.9500}{11.9552}{\normalfont}{:}
\texto{149.8197}{667.9500}{11.9552}{\normalfont}{creación}
\texto{194.9458}{667.9500}{11.9552}{\normalfont}{o}
\texto{204.3139}{667.9500}{11.9552}{\normalfont}{edición,}
\texto{246.8385}{667.9500}{11.9552}{\normalfont}{definición}
\texto{299.1306}{667.9500}{11.9552}{\normalfont}{de}
\texto{314.3400}{667.9500}{11.9552}{\normalfont}{disponibilidad}
\texto{389.7821}{667.9500}{11.9552}{\normalfont}{y}
\texto{399.4753}{667.9500}{11.9552}{\normalfont}{publicación.}
\texto{464.1123}{667.9500}{11.9552}{\normalfont}{La}
\texto{480.7946}{667.9500}{11.9552}{\normalfont}{figura}
\texto{513.5640}{667.9500}{11.9552}{\normalfont}{4.4}
\texto{532.0300}{667.9500}{11.9552}{\normalfont}{lo}
\texto{70.8660}{653.5040}{11.9552}{\normalfont}{modela.}
\texto{97.6760}{587.0970}{7.9701}{\normalfont}{Operador}
\texto{94.1540}{577.6330}{7.9701}{\normalfont}{(backoffice)}
\texto{177.4270}{587.1970}{7.9701}{\normalfont}{Crea}
\texto{197.6631}{587.1970}{7.9701}{\normalfont}{o}
\texto{204.7190}{587.1970}{7.9701}{\normalfont}{edita}
\texto{173.8750}{577.7330}{7.9701}{\normalfont}{el}
\texto{182.8119}{577.7330}{7.9701}{\normalfont}{tour}
\texto{201.1662}{577.7330}{7.9701}{\normalfont}{(datos,}
\texto{176.1090}{568.2680}{7.9701}{\normalfont}{fotos,}
\texto{198.9728}{568.2680}{7.9701}{\normalfont}{precio)}
\texto{268.4240}{582.2430}{7.9701}{\normalfont}{Define}
\texto{294.6464}{582.2430}{7.9701}{\normalfont}{fechas,}
\texto{263.3930}{572.7790}{7.9701}{\normalfont}{horarios}
\texto{295.9158}{572.7790}{7.9701}{\normalfont}{y}
\texto{303.2068}{572.7790}{7.9701}{\normalfont}{cupos}
\texto{361.1460}{582.2430}{7.9701}{\normalfont}{Publica}
\texto{391.8381}{582.2430}{7.9701}{\normalfont}{o}
\texto{398.8940}{582.2430}{7.9701}{\normalfont}{des-}
\texto{360.2820}{572.7790}{7.9701}{\normalfont}{publica}
\texto{389.9196}{572.7790}{7.9701}{\normalfont}{el}
\texto{398.8565}{572.7790}{7.9701}{\normalfont}{tour}
\texto{410.0920}{577.5110}{7.9701}{\normalfont}{Tour}
\texto{430.5600}{577.5110}{7.9701}{\normalfont}{visible}
\texto{456.7155}{577.5110}{7.9701}{\normalfont}{en}
\texto{468.0051}{577.5110}{7.9701}{\normalfont}{la}
\texto{477.4210}{577.5110}{7.9701}{\normalfont}{plataforma}
\texto{100.8340}{502.0570}{7.9701}{\normalfont}{Sistema}
\texto{91.5220}{492.5930}{7.9701}{\normalfont}{(API}
\texto{112.7878}{492.5930}{7.9701}{\normalfont}{REST)}
\texto{182.1150}{496.4280}{7.9701}{\normalfont}{Valida}
\texto{208.4474}{496.4280}{7.9701}{\normalfont}{los}
\texto{174.3450}{486.9640}{7.9701}{\normalfont}{datos}
\texto{196.9737}{486.9640}{7.9701}{\normalfont}{del}
\texto{210.6153}{486.9640}{7.9701}{\normalfont}{tour}
\texto{264.7710}{497.2030}{7.9701}{\normalfont}{Registra}
\texto{298.3403}{497.2030}{7.9701}{\normalfont}{el}
\texto{307.2771}{497.2030}{7.9701}{\normalfont}{tour}
\texto{268.3870}{487.7390}{7.9701}{\normalfont}{en}
\texto{279.6766}{487.7390}{7.9701}{\normalfont}{el}
\texto{288.6135}{487.7390}{7.9701}{\normalfont}{catálogo}
\texto{365.2290}{501.9350}{7.9701}{\normalfont}{Actualiza}
\texto{402.8519}{501.9350}{7.9701}{\normalfont}{la}
\texto{361.3200}{492.4710}{7.9701}{\normalfont}{disponibilidad}
\texto{360.6340}{483.0060}{7.9701}{\normalfont}{en}
\texto{371.9236}{483.0060}{7.9701}{\normalfont}{tiempo}
\texto{400.3809}{483.0060}{7.9701}{\normalfont}{real}
\texto{70.8660}{452.8260}{11.9552}{\normalfont}{Figura}
\texto{108.0945}{452.8260}{11.9552}{\normalfont}{4.4:}
\texto{129.8910}{452.8260}{11.9552}{\normalfont}{BPMN}
\texto{169.2212}{452.8260}{11.9552}{\normalfont}{de}
\texto{184.5024}{452.8260}{11.9552}{\normalfont}{la}
\texto{197.1940}{452.8260}{11.9552}{\normalfont}{gestión}
\texto{236.6031}{452.8260}{11.9552}{\normalfont}{del}
\texto{255.1480}{452.8260}{11.9552}{\normalfont}{catálogo}
\texto{300.9962}{452.8260}{11.9552}{\normalfont}{y}
\texto{310.7612}{452.8260}{11.9552}{\normalfont}{la}
\texto{323.4409}{452.8260}{11.9552}{\normalfont}{disponibilidad}
\texto{398.9547}{452.8260}{11.9552}{\normalfont}{en}
\texto{414.2478}{452.8260}{11.9552}{\normalfont}{el}
\texto{426.2771}{452.8260}{11.9552}{\normalfont}{backoffice}
\texto{478.9589}{452.8260}{11.9552}{\normalfont}{de}
\texto{494.2520}{452.8260}{11.9552}{\normalfont}{escritorio}
\texto{88.7990}{420.8670}{11.9552}{\normalfont}{La}
\texto{105.6128}{420.8670}{11.9552}{\normalfont}{gestión}
\texto{145.1056}{420.8670}{11.9552}{\normalfont}{de}
\texto{160.4704}{420.8670}{11.9552}{\normalfont}{la}
\texto{173.2218}{420.8670}{11.9552}{\normalfont}{disponibilidad}
\texto{248.8074}{420.8670}{11.9552}{\normalfont}{es}
\texto{262.2857}{420.8670}{11.9552}{\normalfont}{uno}
\texto{284.8045}{420.8670}{11.9552}{\normalfont}{de}
\texto{300.1693}{420.8670}{11.9552}{\normalfont}{los}
\texto{317.5378}{420.8670}{11.9552}{\normalfont}{puntos}
\texto{355.4071}{420.8670}{11.9552}{\normalfont}{críticos}
\texto{395.5503}{420.8670}{11.9552}{\normalfont}{que}
\texto{417.0816}{420.8670}{11.9552}{\normalfont}{automatiza}
\texto{477.9718}{420.8670}{11.9552}{\normalfont}{el}
\texto{490.0849}{420.8670}{11.9552}{\normalfont}{prototipo:}
\texto{70.8660}{406.4210}{11.9552}{\normalfont}{al}
\texto{83.2827}{406.4210}{11.9552}{\normalfont}{definir}
\texto{119.1124}{406.4210}{11.9552}{\normalfont}{fechas,}
\texto{156.3218}{406.4210}{11.9552}{\normalfont}{horarios}
\texto{200.6708}{406.4210}{11.9552}{\normalfont}{y}
\texto{210.1728}{406.4210}{11.9552}{\normalfont}{cupos}
\texto{242.4877}{406.4210}{11.9552}{\normalfont}{desde}
\texto{273.8414}{406.4210}{11.9552}{\normalfont}{el}
\texto{285.6077}{406.4210}{11.9552}{\normalfont}{backoffice,}
\texto{341.2783}{406.4210}{11.9552}{\normalfont}{el}
\texto{353.0565}{406.4210}{11.9552}{\normalfont}{sistema}
\texto{394.2183}{406.4210}{11.9552}{\normalfont}{centraliza}
\texto{446.6466}{406.4210}{11.9552}{\normalfont}{la}
\texto{459.0633}{406.4210}{11.9552}{\normalfont}{información}
\texto{522.5454}{406.4210}{11.9552}{\normalfont}{y}
\texto{532.0354}{406.4210}{11.9552}{\normalfont}{la}
\texto{70.8660}{391.9760}{11.9552}{\normalfont}{actualiza}
\texto{119.2009}{391.9760}{11.9552}{\normalfont}{en}
\texto{133.7169}{391.9760}{11.9552}{\normalfont}{tiempo}
\texto{171.9687}{391.9760}{11.9552}{\normalfont}{real,}
\texto{196.9025}{391.9760}{11.9552}{\normalfont}{eliminando}
\texto{255.6432}{391.9760}{11.9552}{\normalfont}{los}
\texto{272.1748}{391.9760}{11.9552}{\normalfont}{errores}
\texto{309.5181}{391.9760}{11.9552}{\normalfont}{de}
\texto{324.0461}{391.9760}{11.9552}{\normalfont}{doble}
\texto{354.1708}{391.9760}{11.9552}{\normalfont}{reserva}
\texto{392.4944}{391.9760}{11.9552}{\normalfont}{y}
\texto{401.4823}{391.9760}{11.9552}{\normalfont}{sobreventa}
\texto{458.6664}{391.9760}{11.9552}{\normalfont}{que}
\texto{479.3608}{391.9760}{11.9552}{\normalfont}{caracterizan}
\texto{70.8660}{377.5300}{11.9552}{\normalfont}{al}
\texto{83.8685}{377.5300}{11.9552}{\normalfont}{proceso}
\texto{125.8743}{377.5300}{11.9552}{\normalfont}{actual.}
\texto{70.8660}{340.9550}{14.3462}{\normalfont\bfseries}{4.7.3.}
\texto{124.6642}{340.9550}{14.3462}{\normalfont\bfseries}{BPMN:}
\texto{184.9542}{340.9550}{14.3462}{\normalfont\bfseries}{Reserva,}
\texto{250.0027}{340.9550}{14.3462}{\normalfont\bfseries}{pago}
\texto{288.3343}{340.9550}{14.3462}{\normalfont\bfseries}{y}
\texto{302.2329}{340.9550}{14.3462}{\normalfont\bfseries}{validación}
\texto{377.5591}{340.9550}{14.3462}{\normalfont\bfseries}{de}
\texto{399.2692}{340.9550}{14.3462}{\normalfont\bfseries}{voucher}
\texto{88.7990}{318.1570}{11.9552}{\normalfont}{El}
\texto{103.5063}{318.1570}{11.9552}{\normalfont}{proceso}
\texto{145.1176}{318.1570}{11.9552}{\normalfont}{de}
\texto{160.3150}{318.1570}{11.9552}{\normalfont}{reserva,}
\texto{202.5719}{318.1570}{11.9552}{\normalfont}{pago}
\texto{230.1382}{318.1570}{11.9552}{\normalfont}{y}
\texto{239.8075}{318.1570}{11.9552}{\normalfont}{validación}
\texto{294.3567}{318.1570}{11.9552}{\normalfont}{de}
\texto{309.5661}{318.1570}{11.9552}{\normalfont}{voucher}
\texto{352.4088}{318.1570}{11.9552}{\normalfont}{constituye}
\texto{407.9957}{318.1570}{11.9552}{\normalfont}{el}
\texto{419.9413}{318.1570}{11.9552}{\normalfont}{núcleo}
\texto{455.6395}{318.1570}{11.9552}{\normalfont}{transaccional}
\texto{526.1824}{318.1570}{11.9552}{\normalfont}{del}
\texto{70.8660}{303.7110}{11.9552}{\normalfont}{prototipo.}
\texto{126.2449}{303.7110}{11.9552}{\normalfont}{En}
\texto{145.0408}{303.7110}{11.9552}{\normalfont}{esta}
\texto{169.5944}{303.7110}{11.9552}{\normalfont}{etapa,}
\texto{205.1516}{303.7110}{11.9552}{\normalfont}{la}
\texto{218.5845}{303.7110}{11.9552}{\normalfont}{reserva}
\texto{258.4264}{303.7110}{11.9552}{\normalfont}{se}
\texto{272.5741}{303.7110}{11.9552}{\normalfont}{registra}
\texto{315.3379}{303.7110}{11.9552}{\normalfont}{y}
\texto{325.8441}{303.7110}{11.9552}{\normalfont}{administra}
\texto{384.8669}{303.7110}{11.9552}{\normalfont}{desde}
\texto{417.2249}{303.7110}{11.9552}{\normalfont}{el}
\texto{430.0580}{303.7110}{11.9552}{\normalfont\bfseries}{backoffice}
\texto{491.4121}{303.7110}{11.9552}{\normalfont\bfseries}{de}
\texto{509.9941}{303.7110}{11.9552}{\normalfont\bfseries}{escri-}
\texto{70.8660}{289.2650}{11.9552}{\normalfont\bfseries}{torio}
\texto{98.7750}{289.2650}{11.9552}{\normalfont}{;}
\texto{106.2350}{289.2650}{11.9552}{\normalfont}{la}
\texto{119.5484}{289.2650}{11.9552}{\normalfont}{generación}
\texto{177.7367}{289.2650}{11.9552}{\normalfont}{del}
\texto{196.9033}{289.2650}{11.9552}{\normalfont}{voucher}
\texto{240.4513}{289.2650}{11.9552}{\normalfont}{y}
\texto{250.8380}{289.2650}{11.9552}{\normalfont}{el}
\texto{263.5129}{289.2650}{11.9552}{\normalfont}{descuento}
\texto{317.5289}{289.2650}{11.9552}{\normalfont}{de}
\texto{333.4436}{289.2650}{11.9552}{\normalfont}{cupo}
\texto{362.0500}{289.2650}{11.9552}{\normalfont}{los}
\texto{379.9804}{289.2650}{11.9552}{\normalfont}{realiza}
\texto{417.3572}{289.2650}{11.9552}{\normalfont}{el}
\texto{430.0201}{289.2650}{11.9552}{\normalfont}{sistema.}
\texto{475.3303}{289.2650}{11.9552}{\normalfont}{La}
\texto{492.7060}{289.2650}{11.9552}{\normalfont}{figura}
\texto{526.1770}{289.2650}{11.9552}{\normalfont}{4.5}
\texto{70.8660}{274.8200}{11.9552}{\normalfont}{lo}
\texto{83.8685}{274.8200}{11.9552}{\normalfont}{modela.}
\texto{88.7990}{260.0630}{11.9552}{\normalfont}{El}
\texto{103.2313}{260.0630}{11.9552}{\normalfont}{sistema}
\texto{144.3094}{260.0630}{11.9552}{\normalfont}{bloquea}
\texto{187.2070}{260.0630}{11.9552}{\normalfont}{el}
\texto{198.8777}{260.0630}{11.9552}{\normalfont}{cupo}
\texto{226.4918}{260.0630}{11.9552}{\normalfont}{durante}
\texto{269.0571}{260.0630}{11.9552}{\normalfont}{la}
\texto{281.3901}{260.0630}{11.9552}{\normalfont}{transacción,}
\texto{345.8167}{260.0630}{11.9552}{\normalfont}{confirma}
\texto{393.2573}{260.0630}{11.9552}{\normalfont}{la}
\texto{405.5903}{260.0630}{11.9552}{\normalfont}{reserva}
\texto{444.3203}{260.0630}{11.9552}{\normalfont}{y}
\texto{453.7267}{260.0630}{11.9552}{\normalfont}{genera}
\texto{490.1112}{260.0630}{11.9552}{\normalfont}{el}
\texto{501.7938}{260.0630}{11.9552}{\normalfont}{voucher}
\texto{70.8660}{245.6180}{11.9552}{\normalfont}{con}
\texto{92.4547}{245.6180}{11.9552}{\normalfont}{código}
\texto{129.3365}{245.6180}{11.9552}{\normalfont}{QR.}
\texto{154.3492}{245.6180}{11.9552}{\normalfont}{El}
\texto{169.5944}{245.6180}{11.9552}{\normalfont}{día}
\texto{189.2440}{245.6180}{11.9552}{\normalfont}{del}
\texto{208.2432}{245.6180}{11.9552}{\normalfont}{tour,}
\texto{236.9859}{245.6180}{11.9552}{\normalfont}{el}
\texto{249.4815}{245.6180}{11.9552}{\normalfont}{operador}
\texto{298.7202}{245.6180}{11.9552}{\normalfont}{valida}
\texto{332.9957}{245.6180}{11.9552}{\normalfont}{el}
\texto{345.4913}{245.6180}{11.9552}{\normalfont}{voucher}
\texto{388.8719}{245.6180}{11.9552}{\normalfont}{desde}
\texto{420.9310}{245.6180}{11.9552}{\normalfont}{el}
\texto{433.4266}{245.6180}{11.9552}{\normalfont}{backoffice}
\texto{486.5627}{245.6180}{11.9552}{\normalfont}{y}
\texto{496.7820}{245.6180}{11.9552}{\normalfont}{marca}
\texto{532.0402}{245.6180}{11.9552}{\normalfont}{la}
\texto{70.8660}{231.1720}{11.9552}{\normalfont}{reserva}
\texto{110.0025}{231.1720}{11.9552}{\normalfont}{como}
\texto{140.2898}{231.1720}{11.9552}{\normalfont}{completada;}
\texto{205.7087}{231.1720}{11.9552}{\normalfont}{si}
\texto{217.2000}{231.1720}{11.9552}{\normalfont}{no}
\texto{233.1913}{231.1720}{11.9552}{\normalfont}{existe}
\texto{265.8194}{231.1720}{11.9552}{\normalfont}{cupo,}
\texto{297.0919}{231.1720}{11.9552}{\normalfont}{el}
\texto{309.1690}{231.1720}{11.9552}{\normalfont}{sistema}
\texto{350.6535}{231.1720}{11.9552}{\normalfont}{lo}
\texto{363.3811}{231.1720}{11.9552}{\normalfont}{comunica}
\texto{414.8195}{231.1720}{11.9552}{\normalfont}{de}
\texto{430.1485}{231.1720}{11.9552}{\normalfont}{inmediato}
\texto{484.5111}{231.1720}{11.9552}{\normalfont}{evitando}
\texto{532.0331}{231.1720}{11.9552}{\normalfont}{la}
\texto{70.8660}{216.7260}{11.9552}{\normalfont}{sobreventa.}
\texto{70.8660}{180.1510}{14.3462}{\normalfont\bfseries}{4.7.4.}
\texto{124.6642}{180.1510}{14.3462}{\normalfont\bfseries}{BPMN:}
\texto{185.7145}{180.1510}{14.3462}{\normalfont\bfseries}{Moderación}
\texto{275.1530}{180.1510}{14.3462}{\normalfont\bfseries}{de}
\texto{297.6378}{180.1510}{14.3462}{\normalfont\bfseries}{reseñas,}
\texto{359.1271}{180.1510}{14.3462}{\normalfont\bfseries}{reportes}
\texto{423.3177}{180.1510}{14.3462}{\normalfont\bfseries}{financieros}
\texto{504.2661}{180.1510}{14.3462}{\normalfont\bfseries}{y}
\texto{518.9394}{180.1510}{14.3462}{\normalfont\bfseries}{au-}
\texto{124.6640}{162.2180}{14.3462}{\normalfont\bfseries}{ditoría}
\texto{88.7990}{139.4200}{11.9552}{\normalfont}{El}
\texto{104.7735}{139.4200}{11.9552}{\normalfont}{proceso}
\texto{147.6521}{139.4200}{11.9552}{\normalfont}{de}
\texto{164.1287}{139.4200}{11.9552}{\normalfont}{moderación}
\texto{227.7543}{139.4200}{11.9552}{\normalfont}{y}
\texto{238.7029}{139.4200}{11.9552}{\normalfont}{supervisión}
\texto{300.8316}{139.4200}{11.9552}{\normalfont}{combina}
\texto{348.2029}{139.4200}{11.9552}{\normalfont}{las}
\texto{366.6952}{139.4200}{11.9552}{\normalfont}{capacidades}
\texto{432.0017}{139.4200}{11.9552}{\normalfont}{del}
\texto{451.7770}{139.4200}{11.9552}{\normalfont\bfseries}{panel}
\texto{488.6074}{139.4200}{11.9552}{\normalfont\bfseries}{de}
\texto{507.6915}{139.4200}{11.9552}{\normalfont\bfseries}{admi-}
\texto{70.8660}{124.9750}{11.9552}{\normalfont\bfseries}{nistración}
\texto{134.0397}{124.9750}{11.9552}{\normalfont\bfseries}{de}
\texto{153.1477}{124.9750}{11.9552}{\normalfont\bfseries}{escritorio}
\texto{212.5020}{124.9750}{11.9552}{\normalfont}{(moderación}
\texto{280.6921}{124.9750}{11.9552}{\normalfont}{de}
\texto{297.1807}{124.9750}{11.9552}{\normalfont}{reseñas}
\texto{338.5242}{124.9750}{11.9552}{\normalfont}{y}
\texto{349.4847}{124.9750}{11.9552}{\normalfont}{auditoría)}
\texto{404.9951}{124.9750}{11.9552}{\normalfont}{con}
\texto{427.3489}{124.9750}{11.9552}{\normalfont}{las}
\texto{445.8532}{124.9750}{11.9552}{\normalfont}{del}
\texto{465.6410}{124.9750}{11.9552}{\normalfont\bfseries}{backoffice}
\texto{527.5211}{124.9750}{11.9552}{\normalfont\bfseries}{de}
\texto{70.8660}{110.5290}{11.9552}{\normalfont\bfseries}{escritorio}
\texto{129.3270}{110.5290}{11.9552}{\normalfont}{(consulta}
\texto{180.1151}{110.5290}{11.9552}{\normalfont}{de}
\texto{195.7190}{110.5290}{11.9552}{\normalfont}{reportes}
\texto{240.9886}{110.5290}{11.9552}{\normalfont}{financieros).}
\texto{306.7350}{110.5290}{11.9552}{\normalfont}{La}
\texto{323.8118}{110.5290}{11.9552}{\normalfont}{figura}
\texto{356.9750}{110.5290}{11.9552}{\normalfont}{4.6}
\texto{375.8350}{110.5290}{11.9552}{\normalfont}{lo}
\texto{388.8375}{110.5290}{11.9552}{\normalfont}{modela.}
\texto{88.7990}{95.7730}{11.9552}{\normalfont}{La}
\texto{106.5333}{95.7730}{11.9552}{\normalfont}{moderación}
\texto{169.9557}{95.7730}{11.9552}{\normalfont}{garantiza}
\texto{221.6763}{95.7730}{11.9552}{\normalfont}{la}
\texto{235.3482}{95.7730}{11.9552}{\normalfont}{calidad}
\texto{276.3354}{95.7730}{11.9552}{\normalfont}{del}
\texto{295.8607}{95.7730}{11.9552}{\normalfont}{marketplace:}
\texto{365.4662}{95.7730}{11.9552}{\normalfont}{el}
\texto{378.4878}{95.7730}{11.9552}{\normalfont}{equipo}
\texto{416.8832}{95.7730}{11.9552}{\normalfont}{gestor}
\texto{452.0817}{95.7730}{11.9552}{\normalfont}{revisa}
\texto{486.3046}{95.7730}{11.9552}{\normalfont}{las}
\texto{504.5937}{95.7730}{11.9552}{\normalfont}{reseñas}
\texto{70.8660}{81.3270}{11.9552}{\normalfont}{pendientes}
\texto{128.1625}{81.3270}{11.9552}{\normalfont}{y}
\texto{137.6047}{81.3270}{11.9552}{\normalfont}{publica}
\texto{177.9272}{81.3270}{11.9552}{\normalfont}{o}
\texto{187.0442}{81.3270}{11.9552}{\normalfont}{rechaza}
\texto{228.3447}{81.3270}{11.9552}{\normalfont}{dejando}
\texto{271.6058}{81.3270}{11.9552}{\normalfont}{constancia}
\texto{328.2638}{81.3270}{11.9552}{\normalfont}{en}
\texto{343.2222}{81.3270}{11.9552}{\normalfont}{la}
\texto{355.5910}{81.3270}{11.9552}{\normalfont}{bitácora}
\texto{400.4661}{81.3270}{11.9552}{\normalfont}{de}
\texto{415.4364}{81.3270}{11.9552}{\normalfont}{auditoría.}
\texto{468.1277}{81.3270}{11.9552}{\normalfont}{En}
\texto{485.8477}{81.3270}{11.9552}{\normalfont}{paralelo,}
\texto{532.6858}{81.3270}{11.9552}{\normalfont}{el}
\texto{70.8660}{66.8810}{11.9552}{\normalfont}{sistema}
\texto{112.5657}{66.8810}{11.9552}{\normalfont}{calcula}
\texto{151.5469}{66.8810}{11.9552}{\normalfont}{automáticamente}
\texto{243.5350}{66.8810}{11.9552}{\normalfont}{los}
\texto{261.1187}{66.8810}{11.9552}{\normalfont}{montos}
\texto{301.7807}{66.8810}{11.9552}{\normalfont}{brutos,}
\texto{341.4767}{66.8810}{11.9552}{\normalfont}{la}
\texto{354.4433}{66.8810}{11.9552}{\normalfont}{comisión}
\texto{402.5941}{66.8810}{11.9552}{\normalfont}{del}
\texto{421.4020}{66.8810}{11.9552}{\normalfont}{8}
\texto{429.2518}{66.8810}{11.9552}{\normalfont}{\%}
\texto{442.8688}{66.8810}{11.9552}{\normalfont}{y}
\texto{452.9088}{66.8810}{11.9552}{\normalfont}{los}
\texto{470.4925}{66.8810}{11.9552}{\normalfont}{montos}
\texto{511.1545}{66.8810}{11.9552}{\normalfont}{netos,}
\end{pagina}
% Página original 55
\begin{pagina}{55}
\texto{70.8660}{746.3890}{11.9552}{\normalfont}{4.7.}
\texto{96.8710}{746.3890}{11.9552}{\normalfont}{BPMN}
\texto{136.5359}{746.3890}{11.9552}{\normalfont}{DE}
\texto{157.3380}{746.3890}{11.9552}{\normalfont}{LOS}
\texto{184.1583}{746.3890}{11.9552}{\normalfont}{PROCESOS}
\texto{251.9610}{746.3890}{11.9552}{\normalfont}{DEL}
\texto{280.0772}{746.3890}{11.9552}{\normalfont}{PROTOTIPO}
\texto{356.6431}{746.3890}{11.9552}{\normalfont}{DE}
\texto{377.4571}{746.3890}{11.9552}{\normalfont}{ESCRITORIO}
\texto{529.4280}{746.3890}{11.9552}{\normalfont}{45}
\texto{97.5170}{669.7410}{7.9701}{\normalfont}{Operador}
\texto{93.9950}{660.2760}{7.9701}{\normalfont}{(backoffice)}
\texto{184.7140}{668.8440}{7.9701}{\normalfont}{Registra}
\texto{182.5330}{659.3790}{7.9701}{\normalfont}{la}
\texto{191.9409}{659.3790}{7.9701}{\normalfont}{reserva}
\texto{181.0080}{649.9150}{7.9701}{\normalfont}{del}
\texto{194.6496}{649.9150}{7.9701}{\normalfont}{turista}
\texto{264.4600}{664.8870}{7.9701}{\normalfont}{Selecciona}
\texto{304.9178}{664.8870}{7.9701}{\normalfont}{tour,}
\texto{265.9860}{655.4220}{7.9701}{\normalfont}{fecha}
\texto{287.6224}{655.4220}{7.9701}{\normalfont}{y}
\texto{294.9135}{655.4220}{7.9701}{\normalfont}{horario}
\texto{376.7810}{668.8440}{7.9701}{\normalfont}{Valida}
\texto{403.1134}{668.8440}{7.9701}{\normalfont}{voucher}
\texto{386.2510}{659.3790}{7.9701}{\normalfont}{QR}
\texto{401.8883}{659.3790}{7.9701}{\normalfont}{el}
\texto{410.8252}{659.3790}{7.9701}{\normalfont}{día}
\texto{389.5980}{649.9150}{7.9701}{\normalfont}{del}
\texto{403.2396}{649.9150}{7.9701}{\normalfont}{tour}
\texto{443.8240}{660.1550}{7.9701}{\normalfont}{Reserva}
\texto{475.2764}{660.1550}{7.9701}{\normalfont}{completada}
\texto{100.6750}{584.7010}{7.9701}{\normalfont}{Sistema}
\texto{91.3630}{575.2360}{7.9701}{\normalfont}{(API}
\texto{112.6288}{575.2360}{7.9701}{\normalfont}{REST)}
\texto{181.5060}{584.5790}{7.9701}{\normalfont}{Verifica}
\texto{212.0793}{584.5790}{7.9701}{\normalfont}{la}
\texto{174.0710}{575.1140}{7.9701}{\normalfont}{disponibilidad}
\texto{173.3860}{565.6500}{7.9701}{\normalfont}{en}
\texto{184.6756}{565.6500}{7.9701}{\normalfont}{tiempo}
\texto{213.1329}{565.6500}{7.9701}{\normalfont}{real}
\texto{270.8230}{575.1140}{7.9701}{\normalfont}{£Hay}
\texto{294.7947}{575.1140}{7.9701}{\normalfont}{cupo?}
\texto{387.0100}{589.3110}{7.9701}{\normalfont}{Confirma}
\texto{374.8380}{579.8460}{7.9701}{\normalfont}{reserva,}
\texto{405.7134}{579.8460}{7.9701}{\normalfont}{registra}
\texto{378.4240}{570.3820}{7.9701}{\normalfont}{pago}
\texto{398.6537}{570.3820}{7.9701}{\normalfont}{y}
\texto{405.9448}{570.3820}{7.9701}{\normalfont}{genera}
\texto{382.1310}{560.9170}{7.9701}{\normalfont}{voucher}
\texto{413.4232}{560.9170}{7.9701}{\normalfont}{QR}
\texto{461.8230}{583.8040}{7.9701}{\normalfont}{Descuenta}
\texto{458.0220}{574.3390}{7.9701}{\normalfont}{el}
\texto{466.9589}{574.3390}{7.9701}{\normalfont}{cupo}
\texto{487.4189}{574.3390}{7.9701}{\normalfont}{defi-}
\texto{459.0810}{564.8750}{7.9701}{\normalfont}{nitivamente}
\texto{238.8160}{507.0820}{7.9701}{\normalfont}{Sin}
\texto{253.3989}{507.0820}{7.9701}{\normalfont}{cupo:}
\texto{276.2109}{507.0820}{7.9701}{\normalfont}{notifica}
\texto{306.3275}{507.0820}{7.9701}{\normalfont}{al}
\texto{315.7354}{507.0820}{7.9701}{\normalfont}{operador}
\texto{345.9260}{581.3850}{5.9776}{\normalfont}{sí}
\texto{297.9330}{554.3060}{5.9776}{\normalfont}{no}
\texto{70.8660}{424.7680}{11.9552}{\normalfont}{Figura}
\texto{108.9911}{424.7680}{11.9552}{\normalfont}{4.5:}
\texto{131.6780}{424.7680}{11.9552}{\normalfont}{BPMN}
\texto{171.9049}{424.7680}{11.9552}{\normalfont}{del}
\texto{191.3344}{424.7680}{11.9552}{\normalfont}{proceso}
\texto{233.9141}{424.7680}{11.9552}{\normalfont}{de}
\texto{250.0919}{424.7680}{11.9552}{\normalfont}{reserva,}
\texto{293.3290}{424.7680}{11.9552}{\normalfont}{pago}
\texto{321.8637}{424.7680}{11.9552}{\normalfont}{y}
\texto{332.5254}{424.7680}{11.9552}{\normalfont}{validación}
\texto{388.0429}{424.7680}{11.9552}{\normalfont}{de}
\texto{404.2326}{424.7680}{11.9552}{\normalfont}{voucher}
\texto{448.0437}{424.7680}{11.9552}{\normalfont}{en}
\texto{464.2334}{424.7680}{11.9552}{\normalfont}{el}
\texto{477.1594}{424.7680}{11.9552}{\normalfont}{prototipo}
\texto{529.4299}{424.7680}{11.9552}{\normalfont}{de}
\texto{70.8660}{410.3220}{11.9552}{\normalfont}{escritorio}
\texto{129.6180}{356.4410}{7.9701}{\normalfont}{Sistema}
\texto{120.3060}{346.9770}{7.9701}{\normalfont}{(API}
\texto{141.5718}{346.9770}{7.9701}{\normalfont}{REST)}
\texto{231.7240}{365.0090}{7.9701}{\normalfont}{Reserva}
\texto{223.6920}{355.5450}{7.9701}{\normalfont}{completada:}
\texto{218.9610}{346.0800}{7.9701}{\normalfont}{habilita}
\texto{250.0101}{346.0800}{7.9701}{\normalfont}{reseña}
\texto{228.4840}{336.6160}{7.9701}{\normalfont}{y}
\texto{235.7750}{336.6160}{7.9701}{\normalfont}{registra}
\texto{229.3350}{327.1510}{7.9701}{\normalfont}{auditoría}
\texto{333.6300}{361.2730}{7.9701}{\normalfont}{Genera}
\texto{362.8611}{361.2730}{7.9701}{\normalfont}{repor-}
\texto{336.6050}{351.8090}{7.9701}{\normalfont}{te}
\texto{346.4831}{351.8090}{7.9701}{\normalfont}{financiero}
\texto{339.6640}{342.3440}{7.9701}{\normalfont}{(bruto,}
\texto{368.3683}{342.3440}{7.9701}{\normalfont}{co-}
\texto{335.1740}{332.8800}{7.9701}{\normalfont}{misión,}
\texto{364.3876}{332.8800}{7.9701}{\normalfont}{neto)}
\texto{443.2670}{360.2770}{7.9701}{\normalfont}{Registra}
\texto{434.8350}{350.8130}{7.9701}{\normalfont}{decisiones}
\texto{473.9754}{350.8130}{7.9701}{\normalfont}{en}
\texto{438.8780}{341.3480}{7.9701}{\normalfont}{la}
\texto{448.2859}{341.3480}{7.9701}{\normalfont}{bitácora}
\texto{436.2900}{331.8840}{7.9701}{\normalfont}{de}
\texto{447.5796}{331.8840}{7.9701}{\normalfont}{auditoría}
\texto{130.7620}{271.4010}{7.9701}{\normalfont}{Equipo}
\texto{119.9710}{261.9370}{7.9701}{\normalfont}{gestor}
\texto{144.9589}{261.9370}{7.9701}{\normalfont}{(panel}
\texto{130.8180}{252.4720}{7.9701}{\normalfont}{admin)}
\texto{217.5630}{266.5470}{7.9701}{\normalfont}{Revisa}
\texto{244.7737}{266.5470}{7.9701}{\normalfont}{las}
\texto{257.5219}{266.5470}{7.9701}{\normalfont}{rese-}
\texto{218.9400}{257.0830}{7.9701}{\normalfont}{ñas}
\texto{234.0409}{257.0830}{7.9701}{\normalfont}{pendientes}
\texto{288.3890}{261.1750}{7.9701}{\normalfont}{£Contenido}
\texto{335.4222}{261.1750}{7.9701}{\normalfont}{adecuado?}
\texto{402.1850}{265.7720}{7.9701}{\normalfont}{Publica}
\texto{399.8620}{256.3080}{7.9701}{\normalfont}{la}
\texto{409.2699}{256.3080}{7.9701}{\normalfont}{reseña}
\texto{275.3320}{193.7820}{7.9701}{\normalfont}{Reseña}
\texto{304.1893}{193.7820}{7.9701}{\normalfont}{rechazada}
\texto{343.4756}{193.7820}{7.9701}{\normalfont}{con}
\texto{359.0077}{193.7820}{7.9701}{\normalfont}{registro}
\texto{126.4600}{186.3600}{7.9701}{\normalfont}{Operador}
\texto{122.9380}{176.8960}{7.9701}{\normalfont}{(backoffice)}
\texto{229.5500}{185.4640}{7.9701}{\normalfont}{Consulta}
\texto{223.9680}{176.0000}{7.9701}{\normalfont}{sus}
\texto{238.1747}{176.0000}{7.9701}{\normalfont}{reportes}
\texto{226.4880}{166.5350}{7.9701}{\normalfont}{financieros}
\texto{263.7050}{176.7740}{7.9701}{\normalfont}{Exporta}
\texto{296.7538}{176.7740}{7.9701}{\normalfont}{reporte}
\texto{326.1746}{176.7740}{7.9701}{\normalfont}{por}
\texto{341.4748}{176.7740}{7.9701}{\normalfont}{rango}
\texto{365.0050}{176.7740}{7.9701}{\normalfont}{de}
\texto{376.2946}{176.7740}{7.9701}{\normalfont}{fechas}
\texto{380.2390}{268.0850}{5.9776}{\normalfont}{sí}
\texto{335.3800}{231.7360}{5.9776}{\normalfont}{no}
\texto{85.9490}{82.0850}{11.9552}{\normalfont}{Figura}
\texto{123.5003}{82.0850}{11.9552}{\normalfont}{4.6:}
\texto{145.6160}{82.0850}{11.9552}{\normalfont}{BPMN}
\texto{185.2690}{82.0850}{11.9552}{\normalfont}{de}
\texto{200.8729}{82.0850}{11.9552}{\normalfont}{la}
\texto{213.8874}{82.0850}{11.9552}{\normalfont}{moderación}
\texto{276.6402}{82.0850}{11.9552}{\normalfont}{de}
\texto{292.2561}{82.0850}{11.9552}{\normalfont}{reseñas,}
\texto{335.9547}{82.0850}{11.9552}{\normalfont}{los}
\texto{353.5743}{82.0850}{11.9552}{\normalfont}{reportes}
\texto{398.8439}{82.0850}{11.9552}{\normalfont}{financieros}
\texto{456.7859}{82.0850}{11.9552}{\normalfont}{y}
\texto{466.8737}{82.0850}{11.9552}{\normalfont}{la}
\texto{479.8762}{82.0850}{11.9552}{\normalfont}{auditoría}
\end{pagina}
% Página original 56
\begin{pagina}{56}
\texto{70.8660}{746.3890}{11.9552}{\normalfont}{46}
\texto{168.2490}{746.3890}{11.9552}{\normalfont}{CAPÍTULO}
\texto{235.2125}{746.3890}{11.9552}{\normalfont}{4.}
\texto{252.1243}{746.3890}{11.9552}{\normalfont}{MODELADO}
\texto{326.8921}{746.3890}{11.9552}{\normalfont}{Y}
\texto{339.5646}{746.3890}{11.9552}{\normalfont}{AUTOMATIZACIÓN}
\texto{455.6233}{746.3890}{11.9552}{\normalfont}{DEL}
\texto{483.7396}{746.3890}{11.9552}{\normalfont}{PROCESO}
\texto{70.8660}{705.1940}{11.9552}{\normalfont}{que}
\texto{91.7876}{705.1940}{11.9552}{\normalfont}{el}
\texto{103.2909}{705.1940}{11.9552}{\normalfont}{operador}
\texto{151.5253}{705.1940}{11.9552}{\normalfont}{consulta}
\texto{196.9121}{705.1940}{11.9552}{\normalfont}{y}
\texto{206.1271}{705.1940}{11.9552}{\normalfont}{exporta}
\texto{248.1951}{705.1940}{11.9552}{\normalfont}{desde}
\texto{279.2738}{705.1940}{11.9552}{\normalfont}{el}
\texto{290.7652}{705.1940}{11.9552}{\normalfont}{backoffice,}
\texto{346.1608}{705.1940}{11.9552}{\normalfont}{y}
\texto{355.3878}{705.1940}{11.9552}{\normalfont}{que}
\texto{376.3214}{705.1940}{11.9552}{\normalfont}{el}
\texto{387.8127}{705.1940}{11.9552}{\normalfont}{administrador}
\texto{462.4657}{705.1940}{11.9552}{\normalfont}{supervisa}
\texto{513.1056}{705.1940}{11.9552}{\normalfont}{desde}
\texto{70.8660}{690.7480}{11.9552}{\normalfont}{el}
\texto{83.2181}{690.7480}{11.9552}{\normalfont}{panel.}
\end{pagina}
% Página original 57
\begin{pagina}{57}
\texto{70.8660}{746.3890}{11.9552}{\normalfont}{4.8.}
\texto{96.8710}{746.3890}{11.9552}{\normalfont}{DESCRIPCIÓN}
\texto{184.0028}{746.3890}{11.9552}{\normalfont}{FUNCIONAL}
\texto{259.7366}{746.3890}{11.9552}{\normalfont}{DEL}
\texto{287.8529}{746.3890}{11.9552}{\normalfont}{PROCESO}
\texto{349.1520}{746.3890}{11.9552}{\normalfont}{PROPUESTO}
\texto{529.4280}{746.3890}{11.9552}{\normalfont}{47}
\texto{70.8660}{705.1940}{17.2154}{\normalfont\bfseries}{4.8.}
\texto{120.3603}{705.1940}{17.2154}{\normalfont\bfseries}{Descripción}
\texto{225.7598}{705.1940}{17.2154}{\normalfont\bfseries}{funcional}
\texto{308.8792}{705.1940}{17.2154}{\normalfont\bfseries}{del}
\texto{340.3111}{705.1940}{17.2154}{\normalfont\bfseries}{proceso}
\texto{410.4225}{705.1940}{17.2154}{\normalfont\bfseries}{propuesto}
\texto{88.7990}{678.9090}{11.9552}{\normalfont}{La}
\texto{106.3182}{678.9090}{11.9552}{\normalfont}{descripción}
\texto{167.6507}{678.9090}{11.9552}{\normalfont}{funcional}
\texto{218.5034}{678.9090}{11.9552}{\normalfont}{del}
\texto{237.8134}{678.9090}{11.9552}{\normalfont}{proceso}
\texto{280.2854}{678.9090}{11.9552}{\normalfont}{propuesto}
\texto{334.7796}{678.9090}{11.9552}{\normalfont}{detalla}
\texto{373.6125}{678.9090}{11.9552}{\normalfont}{cada}
\texto{401.3773}{678.9090}{11.9552}{\normalfont}{actividad,}
\texto{456.1440}{678.9090}{11.9552}{\normalfont}{su}
\texto{471.6165}{678.9090}{11.9552}{\normalfont}{responsable}
\texto{534.9647}{678.9090}{11.9552}{\normalfont}{y}
\texto{70.8660}{664.4630}{11.9552}{\normalfont}{el}
\texto{83.3138}{664.4630}{11.9552}{\normalfont}{mecanismo}
\texto{143.3145}{664.4630}{11.9552}{\normalfont}{de}
\texto{159.0260}{664.4630}{11.9552}{\normalfont}{automatización,}
\texto{244.3144}{664.4630}{11.9552}{\normalfont}{complementando}
\texto{334.4949}{664.4630}{11.9552}{\normalfont}{los}
\texto{352.2221}{664.4630}{11.9552}{\normalfont}{diagramas}
\texto{408.3087}{664.4630}{11.9552}{\normalfont}{BPMN}
\texto{448.0693}{664.4630}{11.9552}{\normalfont}{presentados}
\texto{511.7020}{664.4630}{11.9552}{\normalfont}{en}
\texto{527.4136}{664.4630}{11.9552}{\normalfont}{las}
\texto{70.8660}{650.0170}{11.9552}{\normalfont}{figuras}
\texto{109.4620}{650.0170}{11.9552}{\normalfont}{4.3}
\texto{129.1330}{650.0170}{11.9552}{\normalfont}{a}
\texto{139.6990}{650.0170}{11.9552}{\normalfont}{4.6.}
\texto{162.6192}{650.0170}{11.9552}{\normalfont}{Las}
\texto{185.1141}{650.0170}{11.9552}{\normalfont}{actividades}
\texto{246.8077}{650.0170}{11.9552}{\normalfont}{que}
\texto{269.4030}{650.0170}{11.9552}{\normalfont}{dependen}
\texto{322.5630}{650.0170}{11.9552}{\normalfont}{de}
\texto{338.9799}{650.0170}{11.9552}{\normalfont}{la}
\texto{352.8073}{650.0170}{11.9552}{\normalfont}{app}
\texto{376.3781}{650.0170}{11.9552}{\normalfont}{móvil}
\texto{409.0565}{650.0170}{11.9552}{\normalfont}{forman}
\texto{449.8620}{650.0170}{11.9552}{\normalfont}{parte}
\texto{481.2372}{650.0170}{11.9552}{\normalfont}{de}
\texto{497.6660}{650.0170}{11.9552}{\normalfont}{la}
\texto{511.4815}{650.0170}{11.9552}{\normalfont}{visión}
\texto{70.8660}{635.5710}{11.9552}{\normalfont}{completa}
\texto{120.7479}{635.5710}{11.9552}{\normalfont}{del}
\texto{139.4243}{635.5710}{11.9552}{\normalfont}{ecosistema}
\texto{197.2396}{635.5710}{11.9552}{\normalfont}{y}
\texto{207.1362}{635.5710}{11.9552}{\normalfont}{se}
\texto{220.6623}{635.5710}{11.9552}{\normalfont}{indican}
\texto{261.4510}{635.5710}{11.9552}{\normalfont}{como}
\texto{291.8220}{635.5710}{11.9552}{\normalfont}{trabajo}
\texto{332.6084}{635.5710}{11.9552}{\normalfont}{futuro;}
\texto{371.1208}{635.5710}{11.9552}{\normalfont}{las}
\texto{388.5491}{635.5710}{11.9552}{\normalfont}{actividades}
\texto{449.2385}{635.5710}{11.9552}{\normalfont}{implementadas}
\texto{529.4292}{635.5710}{11.9552}{\normalfont}{en}
\texto{70.8660}{621.1260}{11.9552}{\normalfont}{el}
\texto{83.2181}{621.1260}{11.9552}{\normalfont}{prototipo}
\texto{134.9148}{621.1260}{11.9552}{\normalfont}{inicial}
\texto{169.3912}{621.1260}{11.9552}{\normalfont}{corresponden}
\texto{240.6633}{621.1260}{11.9552}{\normalfont}{a}
\texto{250.4259}{621.1260}{11.9552}{\normalfont}{las}
\texto{268.0455}{621.1260}{11.9552}{\normalfont}{dos}
\texto{288.9288}{621.1260}{11.9552}{\normalfont}{aplicaciones}
\texto{353.3745}{621.1260}{11.9552}{\normalfont}{de}
\texto{368.9785}{621.1260}{11.9552}{\normalfont}{escritorio:}
\texto{85.6720}{593.7280}{11.9552}{\normalfont}{1.}
\texto{100.7540}{593.7280}{11.9552}{\normalfont\bfseries}{Exploración}
\texto{173.5803}{593.7280}{11.9552}{\normalfont\bfseries}{y}
\texto{183.7159}{593.7280}{11.9552}{\normalfont\bfseries}{selección}
\texto{237.9435}{593.7280}{11.9552}{\normalfont\bfseries}{(turista,}
\texto{288.9528}{593.7280}{11.9552}{\normalfont\bfseries}{app}
\texto{313.4597}{593.7280}{11.9552}{\normalfont\bfseries}{móvil}
\texto{348.6295}{593.7280}{11.9552}{\normalfont\bfseries}{---}
\texto{365.1038}{593.7280}{11.9552}{\normalfont\bfseries}{trabajo}
\texto{410.6184}{593.7280}{11.9552}{\normalfont\bfseries}{futuro):}
\texto{458.7280}{593.7280}{11.9552}{\normalfont}{el}
\texto{469.8129}{593.7280}{11.9552}{\normalfont}{turista}
\texto{506.3384}{593.7280}{11.9552}{\normalfont}{navega}
\texto{100.7540}{579.2820}{11.9552}{\normalfont}{por}
\texto{121.8597}{579.2820}{11.9552}{\normalfont}{el}
\texto{134.1879}{579.2820}{11.9552}{\normalfont}{catálogo,}
\texto{183.5988}{579.2820}{11.9552}{\normalfont}{filtra}
\texto{212.1860}{579.2820}{11.9552}{\normalfont}{por}
\texto{233.3037}{579.2820}{11.9552}{\normalfont}{destino}
\texto{273.6621}{579.2820}{11.9552}{\normalfont}{y}
\texto{283.7259}{579.2820}{11.9552}{\normalfont}{categoría,}
\texto{337.0270}{579.2820}{11.9552}{\normalfont}{consulta}
\texto{383.2386}{579.2820}{11.9552}{\normalfont}{reseñas}
\texto{423.6735}{579.2820}{11.9552}{\normalfont}{verificadas}
\texto{480.9437}{579.2820}{11.9552}{\normalfont}{y}
\texto{490.9956}{579.2820}{11.9552}{\normalfont}{selecciona}
\texto{100.7540}{564.8370}{11.9552}{\normalfont}{un}
\texto{117.0728}{564.8370}{11.9552}{\normalfont}{tour}
\texto{141.8464}{564.8370}{11.9552}{\normalfont}{con}
\texto{162.7178}{564.8370}{11.9552}{\normalfont}{su}
\texto{177.1501}{564.8370}{11.9552}{\normalfont}{fecha}
\texto{206.4786}{564.8370}{11.9552}{\normalfont}{y}
\texto{215.9687}{564.8370}{11.9552}{\normalfont}{horario.}
\texto{258.9524}{564.8370}{11.9552}{\normalfont}{En}
\texto{276.7322}{564.8370}{11.9552}{\normalfont}{el}
\texto{288.4985}{564.8370}{11.9552}{\normalfont}{prototipo}
\texto{339.6093}{564.8370}{11.9552}{\normalfont}{inicial,}
\texto{376.7517}{564.8370}{11.9552}{\normalfont}{esta}
\texto{400.2891}{564.8370}{11.9552}{\normalfont}{exploración}
\texto{461.8082}{564.8370}{11.9552}{\normalfont}{se}
\texto{474.9398}{564.8370}{11.9552}{\normalfont}{sustituye}
\texto{523.9083}{564.8370}{11.9552}{\normalfont}{por}
\texto{100.7540}{550.3910}{11.9552}{\normalfont}{la}
\texto{113.7565}{550.3910}{11.9552}{\normalfont}{gestión}
\texto{153.4884}{550.3910}{11.9552}{\normalfont}{directa}
\texto{192.5173}{550.3910}{11.9552}{\normalfont}{de}
\texto{208.1213}{550.3910}{11.9552}{\normalfont}{las}
\texto{225.7408}{550.3910}{11.9552}{\normalfont}{reservas}
\texto{269.7694}{550.3910}{11.9552}{\normalfont}{desde}
\texto{301.7090}{550.3910}{11.9552}{\normalfont}{el}
\texto{314.0611}{550.3910}{11.9552}{\normalfont}{backoffice}
\texto{367.0656}{550.3910}{11.9552}{\normalfont}{de}
\texto{382.6815}{550.3910}{11.9552}{\normalfont}{escritorio.}
\texto{85.6720}{525.9820}{11.9552}{\normalfont}{2.}
\texto{100.7540}{525.9820}{11.9552}{\normalfont\bfseries}{V}
\texto{109.7981}{525.9820}{11.9552}{\normalfont\bfseries}{erificación}
\texto{172.9240}{525.9820}{11.9552}{\normalfont\bfseries}{de}
\texto{190.4060}{525.9820}{11.9552}{\normalfont\bfseries}{disponibilidad}
\texto{276.7310}{525.9820}{11.9552}{\normalfont\bfseries}{(sistema,}
\texto{333.0316}{525.9820}{11.9552}{\normalfont\bfseries}{API}
\texto{361.2578}{525.9820}{11.9552}{\normalfont\bfseries}{REST):}
\texto{409.2860}{525.9820}{11.9552}{\normalfont}{el}
\texto{421.1121}{525.9820}{11.9552}{\normalfont}{backend}
\texto{465.7743}{525.9820}{11.9552}{\normalfont}{valida}
\texto{499.3923}{525.9820}{11.9552}{\normalfont}{en}
\texto{514.4702}{525.9820}{11.9552}{\normalfont}{tiem-}
\texto{100.7540}{511.5370}{11.9552}{\normalfont}{po}
\texto{119.1722}{511.5370}{11.9552}{\normalfont}{real}
\texto{143.7832}{511.5370}{11.9552}{\normalfont}{que}
\texto{167.4066}{511.5370}{11.9552}{\normalfont}{exista}
\texto{202.8132}{511.5370}{11.9552}{\normalfont}{cupo}
\texto{232.9498}{511.5370}{11.9552}{\normalfont}{para}
\texto{261.4510}{511.5370}{11.9552}{\normalfont}{la}
\texto{276.3065}{511.5370}{11.9552}{\normalfont}{fecha}
\texto{308.0620}{511.5370}{11.9552}{\normalfont}{y}
\texto{319.9908}{511.5370}{11.9552}{\normalfont}{horario}
\texto{362.1616}{511.5370}{11.9552}{\normalfont}{seleccionados,}
\texto{438.2684}{511.5370}{11.9552}{\normalfont}{bloqueando}
\texto{502.5491}{511.5370}{11.9552}{\normalfont}{el}
\texto{516.7543}{511.5370}{11.9552}{\normalfont}{cupo}
\texto{100.7540}{497.0910}{11.9552}{\normalfont}{provisionalmente}
\texto{190.8819}{497.0910}{11.9552}{\normalfont}{durante}
\texto{234.1406}{497.0910}{11.9552}{\normalfont}{la}
\texto{247.1430}{497.0910}{11.9552}{\normalfont}{transacción.}
\texto{85.6720}{472.6820}{11.9552}{\normalfont}{3.}
\texto{100.7540}{472.6820}{11.9552}{\normalfont\bfseries}{Pago}
\texto{134.9148}{472.6820}{11.9552}{\normalfont\bfseries}{en}
\texto{153.8913}{472.6820}{11.9552}{\normalfont\bfseries}{línea}
\texto{186.8781}{472.6820}{11.9552}{\normalfont\bfseries}{(sistema,}
\texto{244.6611}{472.6820}{11.9552}{\normalfont\bfseries}{Stripe):}
\texto{294.2240}{472.6820}{11.9552}{\normalfont}{el}
\texto{307.3412}{472.6820}{11.9552}{\normalfont}{sistema}
\texto{349.8659}{472.6820}{11.9552}{\normalfont}{procesa}
\texto{392.6488}{472.6820}{11.9552}{\normalfont}{el}
\texto{405.7660}{472.6820}{11.9552}{\normalfont}{pago}
\texto{434.5039}{472.6820}{11.9552}{\normalfont}{mediante}
\texto{485.6698}{472.6820}{11.9552}{\normalfont}{la}
\texto{499.4494}{472.6820}{11.9552}{\normalfont}{pasarela}
\texto{100.7540}{458.2360}{11.9552}{\normalfont}{Stripe}
\texto{136.3303}{458.2360}{11.9552}{\normalfont}{con}
\texto{158.5885}{458.2360}{11.9552}{\normalfont}{certificación}
\texto{224.4090}{458.2360}{11.9552}{\normalfont}{PCI}
\texto{249.7493}{458.2360}{11.9552}{\normalfont}{DSS.}
\texto{279.6349}{458.2360}{11.9552}{\normalfont}{Al}
\texto{296.3602}{458.2360}{11.9552}{\normalfont}{confirmarse,}
\texto{362.8957}{458.2360}{11.9552}{\normalfont}{la}
\texto{376.6991}{458.2360}{11.9552}{\normalfont}{reserva}
\texto{416.8997}{458.2360}{11.9552}{\normalfont}{adquiere}
\texto{464.8472}{458.2360}{11.9552}{\normalfont}{estado}
\texto{502.1164}{458.2360}{11.9552}{\normalfont}{«confir-}
\texto{100.7540}{443.7910}{11.9552}{\normalfont}{mada»}
\texto{139.1206}{443.7910}{11.9552}{\normalfont}{y}
\texto{149.1965}{443.7910}{11.9552}{\normalfont}{se}
\texto{162.9258}{443.7910}{11.9552}{\normalfont}{descuenta}
\texto{216.6429}{443.7910}{11.9552}{\normalfont}{el}
\texto{228.9950}{443.7910}{11.9552}{\normalfont}{cupo}
\texto{257.2906}{443.7910}{11.9552}{\normalfont}{definitivamente.}
\texto{85.6720}{419.3820}{11.9552}{\normalfont}{4.}
\texto{100.7540}{419.3820}{11.9552}{\normalfont\bfseries}{Generación}
\texto{171.7476}{419.3820}{11.9552}{\normalfont\bfseries}{del}
\texto{193.8264}{419.3820}{11.9552}{\normalfont\bfseries}{voucher}
\texto{244.1949}{419.3820}{11.9552}{\normalfont\bfseries}{(sistema):}
\texto{305.9830}{419.3820}{11.9552}{\normalfont}{se}
\texto{319.9275}{419.3820}{11.9552}{\normalfont}{genera}
\texto{357.2206}{419.3820}{11.9552}{\normalfont}{un}
\texto{374.3404}{419.3820}{11.9552}{\normalfont}{voucher}
\texto{417.8167}{419.3820}{11.9552}{\normalfont}{digital}
\texto{454.4594}{419.3820}{11.9552}{\normalfont}{con}
\texto{476.1438}{419.3820}{11.9552}{\normalfont}{código}
\texto{513.1092}{419.3820}{11.9552}{\normalfont}{QR}
\texto{534.9538}{419.3820}{11.9552}{\normalfont}{y}
\texto{100.7540}{404.9360}{11.9552}{\normalfont}{los}
\texto{118.3736}{404.9360}{11.9552}{\normalfont}{datos}
\texto{149.6508}{404.9360}{11.9552}{\normalfont}{de}
\texto{165.2667}{404.9360}{11.9552}{\normalfont}{la}
\texto{178.2691}{404.9360}{11.9552}{\normalfont}{reserva}
\texto{217.6806}{404.9360}{11.9552}{\normalfont}{(tour,}
\texto{250.8444}{404.9360}{11.9552}{\normalfont}{fecha,}
\texto{284.0224}{404.9360}{11.9552}{\normalfont}{horario,}
\texto{327.5920}{404.9360}{11.9552}{\normalfont}{monto).}
\texto{85.6720}{380.5280}{11.9552}{\normalfont}{5.}
\texto{100.7540}{380.5280}{11.9552}{\normalfont\bfseries}{V}
\texto{109.7981}{380.5280}{11.9552}{\normalfont\bfseries}{alidación}
\texto{166.2123}{380.5280}{11.9552}{\normalfont\bfseries}{y}
\texto{177.8065}{380.5280}{11.9552}{\normalfont\bfseries}{cierre}
\texto{215.2753}{380.5280}{11.9552}{\normalfont\bfseries}{(operador,}
\texto{281.1628}{380.5280}{11.9552}{\normalfont\bfseries}{backoffice}
\texto{342.0506}{380.5280}{11.9552}{\normalfont\bfseries}{de}
\texto{360.1543}{380.5280}{11.9552}{\normalfont\bfseries}{escritorio):}
\texto{427.5950}{380.5280}{11.9552}{\normalfont}{el}
\texto{439.9591}{380.5280}{11.9552}{\normalfont}{día}
\texto{459.4771}{380.5280}{11.9552}{\normalfont}{del}
\texto{478.3448}{380.5280}{11.9552}{\normalfont}{tour,}
\texto{506.9799}{380.5280}{11.9552}{\normalfont}{el}
\texto{519.3440}{380.5280}{11.9552}{\normalfont}{ope-}
\texto{100.7540}{366.0820}{11.9552}{\normalfont}{rador}
\texto{131.6199}{366.0820}{11.9552}{\normalfont}{valida}
\texto{165.4053}{366.0820}{11.9552}{\normalfont}{el}
\texto{177.4107}{366.0820}{11.9552}{\normalfont}{voucher}
\texto{220.3132}{366.0820}{11.9552}{\normalfont}{desde}
\texto{251.8940}{366.0820}{11.9552}{\normalfont}{el}
\texto{263.8994}{366.0820}{11.9552}{\normalfont}{backoffice,}
\texto{319.8091}{366.0820}{11.9552}{\normalfont}{confirma}
\texto{367.5845}{366.0820}{11.9552}{\normalfont}{la}
\texto{380.2403}{366.0820}{11.9552}{\normalfont}{identidad}
\texto{431.5950}{366.0820}{11.9552}{\normalfont}{del}
\texto{450.1041}{366.0820}{11.9552}{\normalfont}{turista}
\texto{487.5382}{366.0820}{11.9552}{\normalfont}{y}
\texto{497.2673}{366.0820}{11.9552}{\normalfont}{marca}
\texto{532.0354}{366.0820}{11.9552}{\normalfont}{la}
\texto{100.7540}{351.6360}{11.9552}{\normalfont}{reserva}
\texto{140.1536}{351.6360}{11.9552}{\normalfont}{como}
\texto{170.7278}{351.6360}{11.9552}{\normalfont}{«completada».}
\texto{248.4460}{351.6360}{11.9552}{\normalfont\itshape}{Implementado}
\texto{324.3842}{351.6360}{11.9552}{\normalfont\itshape}{en}
\texto{340.5238}{351.6360}{11.9552}{\normalfont\itshape}{el}
\texto{353.0767}{351.6360}{11.9552}{\normalfont\itshape}{prototipo}
\texto{402.2425}{351.6360}{11.9552}{\normalfont\itshape}{inicial.}
\texto{85.6720}{327.2280}{11.9552}{\normalfont}{6.}
\texto{100.7540}{327.2280}{11.9552}{\normalfont\bfseries}{Reseña}
\texto{148.7111}{327.2280}{11.9552}{\normalfont\bfseries}{verificada}
\texto{211.1399}{327.2280}{11.9552}{\normalfont\bfseries}{(turista,}
\texto{265.4488}{327.2280}{11.9552}{\normalfont\bfseries}{app}
\texto{293.2674}{327.2280}{11.9552}{\normalfont\bfseries}{móvil}
\texto{331.7368}{327.2280}{11.9552}{\normalfont\bfseries}{---}
\texto{351.5346}{327.2280}{11.9552}{\normalfont\bfseries}{trabajo}
\texto{400.3489}{327.2280}{11.9552}{\normalfont\bfseries}{futuro):}
\texto{451.3430}{327.2280}{11.9552}{\normalfont}{al}
\texto{465.9594}{327.2280}{11.9552}{\normalfont}{completarse}
\texto{532.0310}{327.2280}{11.9552}{\normalfont}{la}
\texto{100.7540}{312.7820}{11.9552}{\normalfont}{reserva,}
\texto{143.9912}{312.7820}{11.9552}{\normalfont}{el}
\texto{156.9291}{312.7820}{11.9552}{\normalfont}{turista}
\texto{195.2838}{312.7820}{11.9552}{\normalfont}{puede}
\texto{229.6837}{312.7820}{11.9552}{\normalfont}{escribir}
\texto{271.3021}{312.7820}{11.9552}{\normalfont}{una}
\texto{294.6458}{312.7820}{11.9552}{\normalfont}{reseña}
\texto{331.0494}{312.7820}{11.9552}{\normalfont}{que}
\texto{353.4176}{312.7820}{11.9552}{\normalfont}{el}
\texto{366.3555}{312.7820}{11.9552}{\normalfont}{sistema}
\texto{408.6769}{312.7820}{11.9552}{\normalfont}{marca}
\texto{444.3775}{312.7820}{11.9552}{\normalfont}{como}
\texto{475.5256}{312.7820}{11.9552}{\normalfont}{verificada;}
\texto{532.0402}{312.7820}{11.9552}{\normalfont}{la}
\texto{100.7540}{298.3360}{11.9552}{\normalfont}{reseña}
\texto{136.1653}{298.3360}{11.9552}{\normalfont}{es}
\texto{149.4643}{298.3360}{11.9552}{\normalfont}{moderada}
\texto{203.3440}{298.3360}{11.9552}{\normalfont}{por}
\texto{224.0552}{298.3360}{11.9552}{\normalfont}{el}
\texto{235.9888}{298.3360}{11.9552}{\normalfont}{equipo}
\texto{273.2843}{298.3360}{11.9552}{\normalfont}{gestor}
\texto{307.3948}{298.3360}{11.9552}{\normalfont}{desde}
\texto{338.9159}{298.3360}{11.9552}{\normalfont}{el}
\texto{350.8496}{298.3360}{11.9552}{\normalfont}{panel}
\texto{381.6438}{298.3360}{11.9552}{\normalfont}{de}
\texto{396.8293}{298.3360}{11.9552}{\normalfont}{administración}
\texto{475.8149}{298.3360}{11.9552}{\normalfont}{de}
\texto{491.0004}{298.3360}{11.9552}{\normalfont}{escritorio.}
\texto{100.7540}{283.8900}{11.9552}{\normalfont\itshape}{La}
\texto{117.6515}{283.8900}{11.9552}{\normalfont\itshape}{moderación}
\texto{179.9680}{283.8900}{11.9552}{\normalfont\itshape}{se}
\texto{194.3022}{283.8900}{11.9552}{\normalfont\itshape}{implementa}
\texto{257.3659}{283.8900}{11.9552}{\normalfont\itshape}{en}
\texto{273.5054}{283.8900}{11.9552}{\normalfont\itshape}{el}
\texto{286.0584}{283.8900}{11.9552}{\normalfont\itshape}{prototipo}
\texto{335.2242}{283.8900}{11.9552}{\normalfont\itshape}{inicial.}
\texto{85.6720}{259.4820}{11.9552}{\normalfont}{7.}
\texto{100.7540}{259.4820}{11.9552}{\normalfont\bfseries}{Reporte}
\texto{153.2039}{259.4820}{11.9552}{\normalfont\bfseries}{financiero}
\texto{215.0804}{259.4820}{11.9552}{\normalfont\bfseries}{(sistema}
\texto{268.6373}{259.4820}{11.9552}{\normalfont\bfseries}{y}
\texto{280.5901}{259.4820}{11.9552}{\normalfont\bfseries}{operador):}
\texto{346.2230}{259.4820}{11.9552}{\normalfont}{el}
\texto{358.8979}{259.4820}{11.9552}{\normalfont}{sistema}
\texto{400.9802}{259.4820}{11.9552}{\normalfont}{calcula}
\texto{440.3319}{259.4820}{11.9552}{\normalfont}{automáticamente}
\texto{532.6787}{259.4820}{11.9552}{\normalfont}{el}
\texto{100.7540}{245.0360}{11.9552}{\normalfont}{monto}
\texto{137.5880}{245.0360}{11.9552}{\normalfont}{bruto,}
\texto{173.4679}{245.0360}{11.9552}{\normalfont}{la}
\texto{187.2236}{245.0360}{11.9552}{\normalfont}{comisión}
\texto{236.1634}{245.0360}{11.9552}{\normalfont}{del}
\texto{255.7723}{245.0360}{11.9552}{\normalfont}{8}
\texto{263.6101}{245.0360}{11.9552}{\normalfont}{\%}
\texto{278.0161}{245.0360}{11.9552}{\normalfont}{y}
\texto{288.8571}{245.0360}{11.9552}{\normalfont}{el}
\texto{301.9624}{245.0360}{11.9552}{\normalfont}{monto}
\texto{338.8083}{245.0360}{11.9552}{\normalfont}{neto,}
\texto{368.8231}{245.0360}{11.9552}{\normalfont}{y}
\texto{379.6521}{245.0360}{11.9552}{\normalfont}{el}
\texto{392.7574}{245.0360}{11.9552}{\normalfont}{operador}
\texto{442.6058}{245.0360}{11.9552}{\normalfont}{puede}
\texto{477.1730}{245.0360}{11.9552}{\normalfont}{exportar}
\texto{525.3955}{245.0360}{11.9552}{\normalfont}{sus}
\texto{100.7540}{230.5900}{11.9552}{\normalfont}{reportes}
\texto{146.0116}{230.5900}{11.9552}{\normalfont}{por}
\texto{167.1532}{230.5900}{11.9552}{\normalfont}{rango}
\texto{199.6665}{230.5900}{11.9552}{\normalfont}{de}
\texto{215.2705}{230.5900}{11.9552}{\normalfont}{fechas}
\texto{249.8138}{230.5900}{11.9552}{\normalfont}{desde}
\texto{281.7414}{230.5900}{11.9552}{\normalfont}{el}
\texto{294.0935}{230.5900}{11.9552}{\normalfont}{backoffice.}
\texto{350.3530}{230.5900}{11.9552}{\normalfont\itshape}{Implementado}
\texto{426.2793}{230.5900}{11.9552}{\normalfont\itshape}{en}
\texto{442.4188}{230.5900}{11.9552}{\normalfont\itshape}{el}
\texto{454.9718}{230.5900}{11.9552}{\normalfont\itshape}{prototipo}
\texto{504.1375}{230.5900}{11.9552}{\normalfont\itshape}{inicial.}
\end{pagina}
% Página original 58
\begin{pagina}{58}
\texto{70.8660}{746.3890}{11.9552}{\normalfont}{48}
\texto{168.2490}{746.3890}{11.9552}{\normalfont}{CAPÍTULO}
\texto{235.2125}{746.3890}{11.9552}{\normalfont}{4.}
\texto{252.1243}{746.3890}{11.9552}{\normalfont}{MODELADO}
\texto{326.8921}{746.3890}{11.9552}{\normalfont}{Y}
\texto{339.5646}{746.3890}{11.9552}{\normalfont}{AUTOMATIZACIÓN}
\texto{455.6233}{746.3890}{11.9552}{\normalfont}{DEL}
\texto{483.7396}{746.3890}{11.9552}{\normalfont}{PROCESO}
\texto{70.8660}{705.1940}{17.2154}{\normalfont\bfseries}{4.9.}
\texto{120.3603}{705.1940}{17.2154}{\normalfont\bfseries}{Beneficios}
\texto{211.0889}{705.1940}{17.2154}{\normalfont\bfseries}{esperados}
\texto{88.7990}{678.9090}{11.9552}{\normalfont}{A}
\texto{100.1684}{678.9090}{11.9552}{\normalfont}{partir}
\texto{132.0410}{678.9090}{11.9552}{\normalfont}{de}
\texto{146.3418}{678.9090}{11.9552}{\normalfont}{los}
\texto{162.6702}{678.9090}{11.9552}{\normalfont}{procesos}
\texto{208.0019}{678.9090}{11.9552}{\normalfont}{modelados}
\texto{264.3133}{678.9090}{11.9552}{\normalfont}{en}
\texto{278.6261}{678.9090}{11.9552}{\normalfont}{la}
\texto{290.3254}{678.9090}{11.9552}{\normalfont}{sección}
\texto{328.7680}{678.9090}{11.9552}{\normalfont}{4.7,}
\texto{349.5721}{678.9090}{11.9552}{\normalfont}{la}
\texto{361.2834}{678.9090}{11.9552}{\normalfont}{automatización}
\texto{441.9212}{678.9090}{11.9552}{\normalfont}{mediante}
\texto{491.0300}{678.9090}{11.9552}{\normalfont\bfseries}{Explora-}
\texto{70.8660}{664.4630}{11.9552}{\normalfont\bfseries}{Nicaragua}
\texto{134.8100}{664.4630}{11.9552}{\normalfont}{genera}
\texto{172.1150}{664.4630}{11.9552}{\normalfont}{beneficios}
\texto{225.4400}{664.4630}{11.9552}{\normalfont}{directamente}
\texto{295.6027}{664.4630}{11.9552}{\normalfont}{asociados}
\texto{347.6795}{664.4630}{11.9552}{\normalfont}{a}
\texto{357.6812}{664.4630}{11.9552}{\normalfont}{cada}
\texto{385.2427}{664.4630}{11.9552}{\normalfont}{proceso}
\texto{427.4996}{664.4630}{11.9552}{\normalfont}{del}
\texto{446.6064}{664.4630}{11.9552}{\normalfont}{prototipo}
\texto{498.5541}{664.4630}{11.9552}{\normalfont}{de}
\texto{514.4091}{664.4630}{11.9552}{\normalfont}{escri-}
\texto{70.8660}{650.0170}{11.9552}{\normalfont}{torio:}
\texto{100.7540}{622.6200}{11.9552}{\normalfont\bfseries}{Reducción}
\texto{166.3809}{622.6200}{11.9552}{\normalfont\bfseries}{de}
\texto{184.6281}{622.6200}{11.9552}{\normalfont\bfseries}{errores:}
\texto{233.1860}{622.6200}{11.9552}{\normalfont}{la}
\texto{246.3200}{622.6200}{11.9552}{\normalfont}{gestión}
\texto{286.1834}{622.6200}{11.9552}{\normalfont}{sistematizada}
\texto{359.2918}{622.6200}{11.9552}{\normalfont}{de}
\texto{375.0273}{622.6200}{11.9552}{\normalfont}{la}
\texto{388.1613}{622.6200}{11.9552}{\normalfont}{disponibilidad}
\texto{464.1294}{622.6200}{11.9552}{\normalfont}{(figura}
\texto{501.9830}{622.6200}{11.9552}{\normalfont}{4.4)}
\texto{525.5224}{622.6200}{11.9552}{\normalfont}{eli-}
\texto{100.7540}{608.1740}{11.9552}{\normalfont}{mina}
\texto{130.0992}{608.1740}{11.9552}{\normalfont}{las}
\texto{147.8144}{608.1740}{11.9552}{\normalfont}{dobles}
\texto{183.7398}{608.1740}{11.9552}{\normalfont}{reservas}
\texto{227.8521}{608.1740}{11.9552}{\normalfont}{y}
\texto{238.0236}{608.1740}{11.9552}{\normalfont}{la}
\texto{251.1217}{608.1740}{11.9552}{\normalfont}{sobreventa}
\texto{309.4894}{608.1740}{11.9552}{\normalfont}{de}
\texto{325.1770}{608.1740}{11.9552}{\normalfont}{cupos,}
\texto{361.4252}{608.1740}{11.9552}{\normalfont}{y}
\texto{371.5967}{608.1740}{11.9552}{\normalfont}{la}
\texto{384.6948}{608.1740}{11.9552}{\normalfont}{verificación}
\texto{445.9221}{608.1740}{11.9552}{\normalfont}{de}
\texto{461.6217}{608.1740}{11.9552}{\normalfont}{cupo}
\texto{490.0010}{608.1740}{11.9552}{\normalfont}{en}
\texto{505.6886}{608.1740}{11.9552}{\normalfont}{tiempo}
\texto{100.7540}{593.7280}{11.9552}{\normalfont}{real}
\texto{123.5119}{593.7280}{11.9552}{\normalfont}{(figura}
\texto{161.2350}{593.7280}{11.9552}{\normalfont}{4.5)}
\texto{184.6429}{593.7280}{11.9552}{\normalfont}{evita}
\texto{213.5793}{593.7280}{11.9552}{\normalfont}{reservas}
\texto{257.6079}{593.7280}{11.9552}{\normalfont}{sin}
\texto{275.8898}{593.7280}{11.9552}{\normalfont}{capacidad.}
\texto{100.7540}{569.3200}{11.9552}{\normalfont\bfseries}{Eficiencia}
\texto{162.4141}{569.3200}{11.9552}{\normalfont\bfseries}{operativa:}
\texto{225.8050}{569.3200}{11.9552}{\normalfont}{el}
\texto{239.3048}{569.3200}{11.9552}{\normalfont}{operador}
\texto{289.5477}{569.3200}{11.9552}{\normalfont}{reduce}
\texto{327.7613}{569.3200}{11.9552}{\normalfont}{el}
\texto{341.2611}{569.3200}{11.9552}{\normalfont}{tiempo}
\texto{381.7486}{569.3200}{11.9552}{\normalfont}{dedicado}
\texto{431.6688}{569.3200}{11.9552}{\normalfont}{a}
\texto{442.5671}{569.3200}{11.9552}{\normalfont}{registros}
\texto{490.6653}{569.3200}{11.9552}{\normalfont}{manuales,}
\texto{100.7540}{554.8740}{11.9552}{\normalfont}{llamadas}
\texto{149.5910}{554.8740}{11.9552}{\normalfont}{y}
\texto{159.6668}{554.8740}{11.9552}{\normalfont}{mensajes,}
\texto{212.4801}{554.8740}{11.9552}{\normalfont}{concentrándose}
\texto{294.1628}{554.8740}{11.9552}{\normalfont}{en}
\texto{309.7668}{554.8740}{11.9552}{\normalfont}{la}
\texto{322.7812}{554.8740}{11.9552}{\normalfont}{operación}
\texto{375.7786}{554.8740}{11.9552}{\normalfont}{del}
\texto{394.6463}{554.8740}{11.9552}{\normalfont}{servicio.}
\texto{100.7540}{530.4660}{11.9552}{\normalfont\bfseries}{Mejora}
\texto{146.5077}{530.4660}{11.9552}{\normalfont\bfseries}{de}
\texto{164.1213}{530.4660}{11.9552}{\normalfont\bfseries}{la}
\texto{178.4006}{530.4660}{11.9552}{\normalfont\bfseries}{experiencia}
\texto{248.7104}{530.4660}{11.9552}{\normalfont\bfseries}{del}
\texto{270.0600}{530.4660}{11.9552}{\normalfont\bfseries}{turista:}
\texto{316.2870}{530.4660}{11.9552}{\normalfont}{reserva}
\texto{355.2681}{530.4660}{11.9552}{\normalfont}{en}
\texto{370.4656}{530.4660}{11.9552}{\normalfont}{línea}
\texto{398.0080}{530.4660}{11.9552}{\normalfont}{en}
\texto{413.2054}{530.4660}{11.9552}{\normalfont}{cualquier}
\texto{463.1853}{530.4660}{11.9552}{\normalfont}{momento,}
\texto{517.0698}{530.4660}{11.9552}{\normalfont}{pago}
\texto{100.7540}{516.0200}{11.9552}{\normalfont}{seguro}
\texto{137.2341}{516.0200}{11.9552}{\normalfont}{y}
\texto{147.3099}{516.0200}{11.9552}{\normalfont}{voucher}
\texto{190.5591}{516.0200}{11.9552}{\normalfont}{digital}
\texto{226.9866}{516.0200}{11.9552}{\normalfont}{estandarizado.}
\texto{100.7540}{491.6110}{11.9552}{\normalfont\bfseries}{Incremento}
\texto{172.1576}{491.6110}{11.9552}{\normalfont\bfseries}{de}
\texto{190.6918}{491.6110}{11.9552}{\normalfont\bfseries}{la}
\texto{205.8916}{491.6110}{11.9552}{\normalfont\bfseries}{confianza:}
\texto{268.0960}{491.6110}{11.9552}{\normalfont}{el}
\texto{280.8307}{491.6110}{11.9552}{\normalfont}{proceso}
\texto{323.2310}{491.6110}{11.9552}{\normalfont}{de}
\texto{339.2175}{491.6110}{11.9552}{\normalfont}{verificación}
\texto{400.7437}{491.6110}{11.9552}{\normalfont}{de}
\texto{416.7302}{491.6110}{11.9552}{\normalfont}{operadores}
\texto{476.0400}{491.6110}{11.9552}{\normalfont}{modelado}
\texto{529.4319}{491.6110}{11.9552}{\normalfont}{en}
\texto{100.7540}{477.1650}{11.9552}{\normalfont}{la}
\texto{114.0195}{477.1650}{11.9552}{\normalfont}{figura}
\texto{147.4540}{477.1650}{11.9552}{\normalfont}{4.3}
\texto{166.5760}{477.1650}{11.9552}{\normalfont}{y}
\texto{176.9149}{477.1650}{11.9552}{\normalfont}{las}
\texto{194.7974}{477.1650}{11.9552}{\normalfont}{reseñas}
\texto{235.5192}{477.1650}{11.9552}{\normalfont}{verificadas}
\texto{293.0764}{477.1650}{11.9552}{\normalfont}{moderadas}
\texto{352.2546}{477.1650}{11.9552}{\normalfont}{en}
\texto{368.1335}{477.1650}{11.9552}{\normalfont}{la}
\texto{381.3990}{477.1650}{11.9552}{\normalfont}{figura}
\texto{414.8310}{477.1650}{11.9552}{\normalfont}{4.6}
\texto{433.9540}{477.1650}{11.9552}{\normalfont}{generan}
\texto{477.7865}{477.1650}{11.9552}{\normalfont}{confianza}
\texto{529.4234}{477.1650}{11.9552}{\normalfont}{en}
\texto{100.7540}{462.7200}{11.9552}{\normalfont}{el}
\texto{113.1061}{462.7200}{11.9552}{\normalfont}{marketplace.}
\texto{100.7540}{438.3110}{11.9552}{\normalfont\bfseries}{Mayores}
\texto{153.4872}{438.3110}{11.9552}{\normalfont\bfseries}{ingresos}
\texto{203.8054}{438.3110}{11.9552}{\normalfont\bfseries}{para}
\texto{233.2786}{438.3110}{11.9552}{\normalfont\bfseries}{los}
\texto{252.4643}{438.3110}{11.9552}{\normalfont\bfseries}{operadores:}
\texto{323.0610}{438.3110}{11.9552}{\normalfont}{la}
\texto{335.1429}{438.3110}{11.9552}{\normalfont}{comisión}
\texto{382.4090}{438.3110}{11.9552}{\normalfont}{reducida}
\texto{428.9721}{438.3110}{11.9552}{\normalfont}{del}
\texto{446.9073}{438.3110}{11.9552}{\normalfont}{8}
\texto{454.7571}{438.3110}{11.9552}{\normalfont}{\%}
\texto{467.4894}{438.3110}{11.9552}{\normalfont}{frente}
\texto{499.7349}{438.3110}{11.9552}{\normalfont}{al}
\texto{511.8288}{438.3110}{11.9552}{\normalfont}{15}
\texto{525.5199}{438.3110}{11.9552}{\normalfont}{\%--}
\texto{100.7540}{423.8650}{11.9552}{\normalfont}{25}
\texto{114.4451}{423.8650}{11.9552}{\normalfont}{\%}
\texto{129.2576}{423.8650}{11.9552}{\normalfont}{de}
\texto{146.0331}{423.8650}{11.9552}{\normalfont}{las}
\texto{164.8124}{423.8650}{11.9552}{\normalfont}{plataformas}
\texto{230.0925}{423.8650}{11.9552}{\normalfont}{internacionales}
\texto{310.9838}{423.8650}{11.9552}{\normalfont}{[6]}
\texto{328.4029}{423.8650}{11.9552}{\normalfont}{incrementa}
\texto{389.7187}{423.8650}{11.9552}{\normalfont}{el}
\texto{403.2305}{423.8650}{11.9552}{\normalfont}{margen}
\texto{446.0086}{423.8650}{11.9552}{\normalfont}{de}
\texto{462.7841}{423.8650}{11.9552}{\normalfont}{los}
\texto{481.5633}{423.8650}{11.9552}{\normalfont}{proveedores}
\texto{100.7540}{409.4190}{11.9552}{\normalfont}{locales.}
\texto{100.7540}{385.0110}{11.9552}{\normalfont\bfseries}{Supervisión}
\texto{174.6718}{385.0110}{11.9552}{\normalfont\bfseries}{y}
\texto{187.1865}{385.0110}{11.9552}{\normalfont\bfseries}{trazabilidad:}
\texto{265.5940}{385.0110}{11.9552}{\normalfont}{el}
\texto{278.7591}{385.0110}{11.9552}{\normalfont}{panel}
\texto{310.7966}{385.0110}{11.9552}{\normalfont}{de}
\texto{327.2254}{385.0110}{11.9552}{\normalfont}{administración}
\texto{407.4544}{385.0110}{11.9552}{\normalfont}{y}
\texto{418.3552}{385.0110}{11.9552}{\normalfont}{la}
\texto{432.1825}{385.0110}{11.9552}{\normalfont}{bitácora}
\texto{478.5281}{385.0110}{11.9552}{\normalfont}{de}
\texto{494.9569}{385.0110}{11.9552}{\normalfont}{auditoría}
\texto{100.7540}{370.5650}{11.9552}{\normalfont}{garantizan}
\texto{158.3087}{370.5650}{11.9552}{\normalfont}{la}
\texto{171.3112}{370.5650}{11.9552}{\normalfont}{supervisión}
\texto{232.5792}{370.5650}{11.9552}{\normalfont}{de}
\texto{248.1831}{370.5650}{11.9552}{\normalfont}{la}
\texto{261.1856}{370.5650}{11.9552}{\normalfont}{plataforma}
\texto{320.7010}{370.5650}{11.9552}{\normalfont}{y}
\texto{330.7768}{370.5650}{11.9552}{\normalfont}{la}
\texto{343.7913}{370.5650}{11.9552}{\normalfont}{trazabilidad}
\texto{408.8228}{370.5650}{11.9552}{\normalfont}{de}
\texto{424.4267}{370.5650}{11.9552}{\normalfont}{las}
\texto{442.0582}{370.5650}{11.9552}{\normalfont}{acciones}
\texto{487.6434}{370.5650}{11.9552}{\normalfont}{críticas.}
\texto{100.7540}{346.1570}{11.9552}{\normalfont\bfseries}{T}
\texto{108.9959}{346.1570}{11.9552}{\normalfont\bfseries}{ransparencia}
\texto{186.8637}{346.1570}{11.9552}{\normalfont\bfseries}{financiera:}
\texto{250.6740}{346.1570}{11.9552}{\normalfont}{los}
\texto{267.6241}{346.1570}{11.9552}{\normalfont}{reportes}
\texto{312.2122}{346.1570}{11.9552}{\normalfont}{automáticos}
\texto{377.2891}{346.1570}{11.9552}{\normalfont}{de}
\texto{392.2236}{346.1570}{11.9552}{\normalfont}{monto}
\texto{427.6468}{346.1570}{11.9552}{\normalfont}{bruto,}
\texto{462.0921}{346.1570}{11.9552}{\normalfont}{comisión}
\texto{509.6093}{346.1570}{11.9552}{\normalfont}{y}
\texto{519.0276}{346.1570}{11.9552}{\normalfont}{neto}
\texto{100.7540}{331.7110}{11.9552}{\normalfont}{permiten}
\texto{150.2246}{331.7110}{11.9552}{\normalfont}{a}
\texto{159.7003}{331.7110}{11.9552}{\normalfont}{operadores}
\texto{218.3406}{331.7110}{11.9552}{\normalfont}{y}
\texto{228.1414}{331.7110}{11.9552}{\normalfont}{administradores}
\texto{313.1883}{331.7110}{11.9552}{\normalfont}{conocer}
\texto{355.5050}{331.7110}{11.9552}{\normalfont}{con}
\texto{376.6872}{331.7110}{11.9552}{\normalfont}{exactitud}
\texto{428.1113}{331.7110}{11.9552}{\normalfont}{el}
\texto{440.1765}{331.7110}{11.9552}{\normalfont}{estado}
\texto{476.3816}{331.7110}{11.9552}{\normalfont}{financiero}
\texto{529.4316}{331.7110}{11.9552}{\normalfont}{de}
\texto{100.7540}{317.2650}{11.9552}{\normalfont}{la}
\texto{113.7565}{317.2650}{11.9552}{\normalfont}{plataforma.}
\texto{100.7540}{292.8570}{11.9552}{\normalfont\bfseries}{Escalabilidad:}
\texto{187.1640}{292.8570}{11.9552}{\normalfont}{el}
\texto{201.2018}{292.8570}{11.9552}{\normalfont}{modelo}
\texto{243.5280}{292.8570}{11.9552}{\normalfont}{de}
\texto{260.8176}{292.8570}{11.9552}{\normalfont}{negocio}
\texto{304.4445}{292.8570}{11.9552}{\normalfont}{automatizado}
\texto{379.6284}{292.8570}{11.9552}{\normalfont}{es}
\texto{395.0315}{292.8570}{11.9552}{\normalfont}{escalable}
\texto{445.5541}{292.8570}{11.9552}{\normalfont}{a}
\texto{456.9905}{292.8570}{11.9552}{\normalfont}{nuevos}
\texto{496.7869}{292.8570}{11.9552}{\normalfont}{destinos,}
\texto{100.7540}{278.4110}{11.9552}{\normalfont}{operadores}
\texto{159.6692}{278.4110}{11.9552}{\normalfont}{y}
\texto{169.7451}{278.4110}{11.9552}{\normalfont}{servicios}
\texto{216.3823}{278.4110}{11.9552}{\normalfont}{complementarios}
\texto{305.8694}{278.4110}{11.9552}{\normalfont}{(hospedaje,}
\texto{368.0460}{278.4110}{11.9552}{\normalfont}{transporte)}
\texto{429.5626}{278.4110}{11.9552}{\normalfont}{en}
\texto{445.1785}{278.4110}{11.9552}{\normalfont}{el}
\texto{457.5306}{278.4110}{11.9552}{\normalfont}{futuro.}
\end{pagina}

\clearpage\restoregeometry
\pagestyle{plain}
\setlength{\parindent}{1.2em}
\setcounter{chapter}{4}
\chapter{Desarrollo e implementación del prototipo}
\label{cap:implementacion}

ExploraNica es un sistema de gestión turística desarrollado con .NET 10, Windows Forms, ASP.NET Core, Entity Framework Core y PostgreSQL. Se implementaron dos aplicaciones de escritorio sobre una API común: el backoffice del operador y el panel administrativo. El sistema centraliza el registro de agencias, la gestión documental, el catálogo de tours, la disponibilidad de salidas y la administración de reservas.

La arquitectura distribuye las responsabilidades entre dominio, aplicación, infraestructura y presentación. Las reglas de reserva y validación se ejecutan en el servidor, mientras los clientes presentan formularios, tablas e indicadores para la operación diaria. Las integraciones de pago, reseñas y acceso móvil forman parte del plan de ampliación del producto.

\section{Arquitectura y estructura}
\subsection{Diagrama de arquitectura}
La solución comprende dos clientes de escritorio y una API central. El cliente del operador permite gestionar la actividad de la agencia; el cliente administrativo concentra la supervisión. Ambos intercambian datos con la API mediante HTTP y JSON. La API resuelve autenticación, autorización y operaciones del negocio, mientras PostgreSQL conserva los registros. Windows Forms proporciona los controles de interfaz utilizados por los clientes \cite{ref15}.

La organización distingue dominio, aplicación, infraestructura y presentación. El dominio define entidades y estados; aplicación contiene contratos, DTO y validadores; infraestructura implementa servicios, acceso a datos y almacenamiento documental; presentación contiene los controladores HTTP y los formularios de escritorio. La inyección de dependencias conecta las interfaces con sus implementaciones y permite mantener las responsabilidades de cada proyecto.

\begin{figure}[htbp]
\centering
\begin{tikzpicture}[node distance=11mm,box/.style={draw,rounded corners,align=center,text width=5cm,minimum height=12mm,font=\small},>=stealth]
\node[box,fill=blue!7] (op) {Cliente del operador\\Windows Forms / .NET 10};
\node[box,fill=blue!7,right=8mm of op] (ad) {Cliente de administración\\Windows Forms / .NET 10};
\coordinate (medio) at ($(op)!0.5!(ad)$);
\node[box,fill=green!8,text width=10.8cm,below=17mm of medio] (api) {API ASP.NET Core: controladores, JWT, autorización\\Contratos y validadores de aplicación / entidades del dominio};
\node[box,fill=orange!10,text width=10.8cm,below=of api] (infra) {Servicios de infraestructura / Entity Framework Core\\Registro de errores, auditoría y almacenamiento documental};
\node[box,fill=gray!8,below left=12mm and 3mm of infra.south] (db) {PostgreSQL 16\\Usuarios, operadores, tours, sesiones, reservas y vouchers};
\node[box,fill=gray!8,below right=12mm and 3mm of infra.south] (files) {Almacenamiento local\\Documentos de operadores};
\draw[->] (op.south) -- (api.north west);
\draw[->] (ad.south) -- (api.north east);
\draw[->] (api) -- (infra);
\draw[->] (infra) -- (db);
\draw[->] (infra) -- (files);
\end{tikzpicture}
\caption{Arquitectura de ExploraNica y comunicación entre sus componentes.}
\label{fig:arquitectura-real}
\end{figure}

JWT identifica al usuario y transporta los atributos de su sesión. La API valida el token y comprueba los permisos antes de ejecutar las operaciones protegidas \cite{ref20}. Los clientes concentran la interacción con el usuario y consumen los servicios mediante solicitudes HTTP con cuerpos JSON.

\clearpage
\subsection{Estructura de carpetas y archivos}
La raíz contiene \archivo{ExploraNica.slnx}, \archivo{README.md}, \archivo{docker-compose.yml}, los scripts de preparación y el directorio de documentación \archivo{odd/tasks}. La estructura principal es la siguiente:
\begin{Verbatim}[fontsize=\small]
ExploraNica/
  ExploraNica.slnx
  README.md
  docker-compose.yml
  scripts/
  odd/tasks/
  src/
    Core/
      ExploraNica.Domain/
        Entities/  Enums/  Exceptions/  Builders/
      ExploraNica.Application/
        DTOs/  Interfaces/  Validators/
    Infrastructure/
      ExploraNica.Infrastructure/
        Services/  Persistence/  Middleware/
      ExploraNica.Shared/
        Configuration/  Services/  UI/
    Presentation/
      ExploraNica.Api/
        Controllers/  Middleware/  Health/  OpenApi/
      ExploraNica.OperatorApp/
        Forms/  Services/  Resources/
      ExploraNica.AdminApp/
        Forms/  Controls/  Resources/
\end{Verbatim}

\archivo{ExploraNicaDbContext.cs} coordina el acceso a datos y las migraciones conservan la evolución del esquema. \archivo{Program.cs} de la API registra servicios y configura el procesamiento de solicitudes. Los archivos \archivo{appsettings.json} contienen las opciones de ejecución, que pueden sobrescribirse con variables de entorno. En los clientes, \archivo{Api:BaseUrl} define la dirección del servidor.

\section{Módulos implementados}
\subsection{Descripción de cada módulo: propósito, entradas y salidas}
\begin{longtable}{@{}P{2.5cm}P{3.25cm}P{3.25cm}P{3.3cm}@{}}
\caption{Módulos de ExploraNica.}\label{tab:modulos}\\
\toprule Módulo & Propósito & Entradas & Salidas \\ \midrule\endfirsthead
\toprule Módulo & Propósito & Entradas & Salidas \\ \midrule\endhead
Autenticación & Comprobar identidad y permitir acceso según rol. & Correo y contraseña; configuración JWT. & Token firmado, vencimiento o error de credenciales. \\
Operadores y documentos & Registrar agencias y mantener su expediente. & Datos personales y comerciales; RUC, identificación y documento institucional. & Cuenta pendiente; documentos registrados; estado y motivo de rechazo. \\
Verificación administrativa & Revisar agencias y resolver su solicitud. & Operador seleccionado, documentos y decisión; motivo si se rechaza. & Estado aprobado o rechazado y actualización visible al operador. \\
Catálogo de tours & Crear y actualizar las ofertas de la agencia; supervisar su activación. & Título, descripción, destino, precio y duración. & Tour registrado o actualizado; indicador de actividad. \\
Sesiones y disponibilidad & Definir salidas y consultar los cupos. & Tour, fecha, hora, cupos y modificador de precio. & Sesión asociada al tour y disponibilidad consultable. \\
Reservas y vouchers & Registrar reservas y validar su utilización. & Sesión, cantidad de cupos e identidad autenticada; reserva a validar. & Reserva confirmada, cupo descontado; reserva completada y voucher usado. \\
Paneles y seguimiento & Mostrar información operativa y registrar acciones. & Listas de operadores, tours y reservas; eventos del sistema. & Conteos de control, errores correlacionados y registros de auditoría. \\
\bottomrule
\end{longtable}

\archivo{AuthService} verifica las contraseñas con \archivo{IPasswordHasher<User>} y genera un token JWT con vencimiento. El sistema utiliza tres roles: \archivo{Admin}, \archivo{TourOperator} y \archivo{Tourist}. Los controladores restringen las operaciones por rol y los servicios resuelven la identidad autenticada para consultar o modificar los recursos asociados al usuario.

\archivo{OperatorRegistrationService}, \archivo{OperatorDocumentService} y \archivo{AdminService} implementan el registro, la carga documental y la verificación administrativa. Los estados son \archivo{Pending}, \archivo{Approved} y \archivo{Rejected}. Las pantallas \archivo{UcVendors}, \archivo{VendorDetailForm}, \archivo{VendorEditForm} y \archivo{VendorDocumentsForm} permiten consultar, editar y resolver los expedientes de las agencias.

\archivo{OperatorService} gestiona los tours de la agencia autenticada. \archivo{BookingService} calcula el importe de la reserva, descuenta los cupos disponibles y registra la confirmación dentro de una transacción. Durante la validación del servicio, actualiza o crea el voucher, registra su utilización y cambia la reserva a completada.

\subsection{Interacción entre módulos}
El inicio de sesión habilita las consultas del cliente y la API obtiene la identidad desde los atributos del token. La agencia se resuelve en el servidor para consultar sus tours y documentos. La creación de un tour desde el servicio del operador exige una cuenta aprobada; posteriormente se asocian sesiones y se consultan las reservas de cada salida.

La validación conecta reserva, sesión, tour y operador. Antes de completar la reserva, el servicio comprueba que el tour pertenece al operador autenticado. El cambio de estado, la actualización del voucher y el evento de auditoría se guardan dentro de una transacción. Los paneles administrativos consultan operadores, tours y reservas para presentar los indicadores de seguimiento.

\section{Funcionalidades implementadas}
\subsection{Funcionalidades principales}
Se implementaron inicio de sesión; registro de operadores por API con documentos; consulta, sustitución y reenvío de documentación; aprobación y rechazo administrativo; consulta y edición de agencias; creación y edición de tours; activación administrativa de ofertas; registro de sesiones; creación de reservas por API; consulta de reservas por sesión y validación de vouchers desde el cliente del operador.

El cliente administrativo incorpora listados de operadores, tours y reservas, además de un tablero con operadores pendientes y verificados, tours activos y reservas confirmadas. El backoffice del operador organiza sus opciones en panel, tours, documentos y reservas, con navegación desde un formulario principal.

\subsection{Correspondencia entre funcionalidades y el proceso propuesto}
\begin{longtable}{@{}P{3.05cm}P{5.05cm}P{4.2cm}@{}}
\caption{Correspondencia entre las funciones y el proceso de gestión.}\label{tab:trazabilidad}\\
\toprule Proceso & Funcionalidad & Resultado operativo \\ \midrule\endfirsthead
\toprule Proceso & Funcionalidad & Resultado operativo \\ \midrule\endhead
Alta y verificación & Registro de agencias, carga documental y decisión administrativa. & Expediente con estado y motivo de rechazo cuando corresponde. \\
Catálogo & Creación y edición de tours; activación administrativa de ofertas. & Oferta asociada a la agencia responsable. \\
Disponibilidad & Registro de sesiones con fecha, hora y cupos. & Salidas organizadas por tour. \\
Reserva & Creación transaccional y actualización condicionada de disponibilidad. & Reserva confirmada y cupos descontados. \\
Validación y cierre & Comprobación de propiedad, estado de la reserva y uso del voucher. & Reserva completada y utilización registrada. \\
Seguimiento & Conteos de operadores, tours y reservas; eventos de auditoría. & Indicadores de operación y trazabilidad de acciones. \\
\bottomrule
\end{longtable}

Las funcionalidades siguen una secuencia operativa: registrar y verificar la agencia, mantener el catálogo, programar salidas, registrar reservas y confirmar la prestación del servicio. La API relaciona estas actividades mediante identificadores de operador, tour, sesión y reserva. Los controles de propiedad y estado conservan la continuidad del proceso entre los dos clientes de escritorio.

\section{Validaciones implementadas}
\subsection{Validaciones de datos de entrada}
Los validadores de aplicación comprueban que las reservas indiquen una sesión y una cantidad de cupos mayor que cero. Para los tours exigen título y ubicación, ambos de hasta 250 caracteres, precio base no negativo y duración positiva. El validador de sesiones exige tour, cupos positivos y una fecha posterior a la hora UTC de referencia con tolerancia de un minuto.

En el registro de operadores se normalizan los textos y se revisan campos obligatorios, formato de correo, longitudes compatibles con el modelo, contraseña de al menos ocho caracteres y unicidad de correo y RUC. El expediente requiere tres tipos documentales: RUC, identidad y constitución institucional. Los archivos se limitan a 5 MiB, se admiten PDF, PNG y JPEG y se verifica su firma binaria, además del tipo declarado. La sustitución documental no se permite si la cuenta ya está aprobada.

La edición administrativa limita el nombre de agencia a 200 caracteres, el RUC a 50 y el teléfono a 50 cuando se proporciona. Un rechazo requiere motivo. Para validar un voucher, la reserva debe pertenecer al operador, estar confirmada y no haber sido utilizada. En las descargas se filtra por propietario y se devuelven nombres controlados, evitando exponer rutas internas.

\subsection{Manejo de errores y mensajes de usuario}
\archivo{ErrorHandlingMiddleware} transforma las excepciones de dominio en respuestas JSON con código, mensaje e identificador de correlación. Los errores de validación de dominio incluyen un mapa por campo. La validación automática de modelos se procesa mediante ASP.NET Core. El cliente utiliza estas respuestas para comunicar el problema y orientar la corrección de los datos.

\begin{longtable}{@{}P{1.25cm}P{3.3cm}P{7.75cm}@{}}
\caption{Códigos y mensajes de error del sistema.}\\
\toprule HTTP & Situación & Mensaje o respuesta representativa \\ \midrule\endfirsthead
\toprule HTTP & Situación & Mensaje o respuesta representativa \\ \midrule\endhead
400 & Regla de negocio & ``No hay cupos suficientes para la cantidad solicitada.'' \\
401 & Credenciales inválidas & ``Credenciales inválidas.'' \\
403 & Operación prohibida & ``Tu cuenta debe estar aprobada para publicar tours.'' \\
404 & Recurso inexistente & ``No se encontró el documento solicitado.'' \\
409 & Conflicto & ``El RUC ya está registrado por otro operador.'' \\
422 & Validación de dominio & ``El motivo de rechazo es obligatorio.'' y errores por campo. \\
500 & Fallo no previsto & ``Ocurrió un error inesperado.'' con correlación. \\
\bottomrule
\end{longtable}

Los clientes incorporan \archivo{UiErrorHandler} y utilidades para actualizar controles desde el hilo de interfaz. Estos componentes centralizan los mensajes de error de las operaciones asíncronas y permiten recuperar el control de la pantalla después de una excepción. La API relaciona los registros técnicos con la respuesta mediante el identificador de correlación y limita los detalles de excepciones al entorno de desarrollo.

\section{Estructuras de datos y lógica de programación}
\subsection{Estructuras de datos usadas}
Las entidades principales son \archivo{User}, \archivo{Operator}, \archivo{OperatorDocument}, \archivo{Tour}, \archivo{TourSession}, \archivo{Booking}, \archivo{Voucher} y \archivo{AuditLog}. Los identificadores son de tipo \archivo{Guid}; los importes se calculan con valores decimales y las fechas de negocio emplean \archivo{DateTimeOffset}. Los enumerados \archivo{UserRole}, \archivo{BookingStatus} y \archivo{DocumentType} delimitan opciones conocidas.

Las listas y colecciones de DTO transportan resultados a las tablas de Windows Forms. Las consultas LINQ filtran y proyectan información, con \archivo{AsNoTracking} en varias consultas de lectura. Los diccionarios \archivo{Dictionary<string,string>} asocian un campo con su error; las listas de documentos preparados permiten verificar archivos antes de escribirlos. Las relaciones entre entidades se conservan con claves foráneas y configuraciones de Entity Framework Core.

La relación operador--tour y tour--sesión es de uno a muchos. La reserva referencia una sesión y un usuario. El voucher se relaciona con una reserva y registra el código, el estado de uso y la fecha de validación. Los documentos se vinculan a la agencia y se distinguen por su tipo, tamaño, formato y fecha de carga.

\subsection{Lógica y algoritmos principales}
\textbf{Creación de reserva.} El servicio abre una transacción dentro de la estrategia de reintentos de EF Core. Ejecuta una actualización condicionada que descuenta cupos únicamente cuando la sesión conserva disponibilidad suficiente. Si ninguna fila cambia, informa falta de cupo. Después carga la sesión y su tour, calcula el total, crea una reserva confirmada y confirma la transacción. Ante un fallo previo a la confirmación, el descuento no debe persistir independientemente de la reserva.

El importe se obtiene mediante:
\[
T=\operatorname{redondear}\bigl((P_b+M_s)\,q,2\bigr),
\]
donde $P_b$ es el precio base, $M_s$ el modificador de la sesión y $q$ la cantidad de cupos. Por ejemplo, con precio base de 30 unidades monetarias, modificador de 5 y dos cupos, el total calculado es 70.

\textbf{Validación de voucher.} Se consulta la reserva junto con su sesión, tour, operador y voucher. Se comprueba la propiedad y el estado; después se realiza un cambio condicionado de confirmado a completado. Si otra solicitud ya realizó el cambio, se rechaza una segunda validación. El voucher usado y la bitácora se guardan en la misma transacción. El control condicional evita depender únicamente de una lectura previa del estado.

\textbf{Registro documental.} Primero se validan datos y archivos; después se crea usuario, operador y documentos dentro de una transacción. Si falla la persistencia, se intentan retirar los archivos escritos para no conservar documentos sin expediente. La sustitución conserva el archivo anterior hasta confirmar la nueva referencia. El servicio coordina la escritura de archivos con la persistencia del expediente para conservar su integridad.

\textbf{Selección en cascada.} La pantalla de reservas carga tours, luego sesiones y finalmente reservas. Al cambiar el nivel superior se limpian las selecciones dependientes y se suprimen eventos disparados durante la recarga. El objetivo es evitar que una tabla muestre reservas de una salida distinta a la seleccionada.

\section{Evidencia y funcionamiento}
\subsection{Capturas de pantalla por módulo}
La interfaz se organiza alrededor de las tareas del operador y del administrador. Cada pantalla mantiene un propósito definido y presenta los datos de la actividad seleccionada:
\begin{description}
\item[Inicio de sesión.] Formulario de correo y contraseña, validación de acceso y comunicación de errores.
\item[Operadores.] Listado administrativo con filtros por estado, datos comerciales y acceso al expediente.
\item[Documentos.] Consulta de RUC, identificación y documento institucional, carga de archivos y reenvío de la solicitud.
\item[Tours.] Listado de ofertas y formulario con título, descripción, destino, precio y duración.
\item[Sesiones.] Selección de fecha, horario, disponibilidad y modificador de precio de cada salida.
\item[Reservas.] Selección de tour y sesión, consulta de reservas y acción de validación del voucher.
\item[Panel administrativo.] Indicadores de operadores pendientes y verificados, tours activos y reservas confirmadas.
\end{description}
\evidencia{01-login.png}{Inicio de sesión}{}
\evidencia{02-operadores.png}{Gestión administrativa de operadores}{}
\evidencia{03-documentos.png}{Gestión documental y verificación}{}
\evidencia{04-tours.png}{Catálogo de tours}{}
\evidencia{05-sesiones.png}{Sesiones y disponibilidad}{}
\evidencia{06-reservas.png}{Reservas y validación de vouchers}{}
\evidencia{07-panel.png}{Panel administrativo}{}

\subsection{Flujos de ejemplo}
\textbf{Flujo 1: alta y revisión de una agencia.} El operador registra sus datos y los tres documentos requeridos. El sistema crea el expediente con estado pendiente. El administrador consulta los documentos y aprueba la cuenta o registra el motivo de rechazo. Si se requiere una corrección, el operador sustituye los archivos y reenvía la solicitud de verificación.

\textbf{Flujo 2: publicación y programación.} Una agencia aprobada inicia sesión, crea un tour y registra una salida con diez cupos. El catálogo conserva la oferta y la sesión queda disponible para consulta. Si una cuenta pendiente intenta crear un tour desde el servicio del operador, el sistema deniega la operación y comunica que necesita aprobación.

\textbf{Flujo 3: reserva y cierre.} Un usuario autenticado solicita por API dos cupos de una sesión con diez disponibles. La transacción registra la reserva confirmada y actualiza el cupo a ocho. El operador selecciona el tour y su sesión en el backoffice, consulta la reserva y valida la atención del servicio. La reserva pasa a completada y el voucher queda utilizado. Una segunda validación produce un mensaje de reutilización.

\textbf{Flujo 4: control de disponibilidad.} Cuando una sesión conserva un cupo y se solicitan dos, el servicio responde que no hay disponibilidad suficiente. La actualización condicionada impide descontar una cantidad superior al cupo restante. La reserva y el descuento se ejecutan dentro de la misma transacción para conservar la consistencia de la operación.

\section{Repositorio GitHub}
\subsection{Enlace del repositorio}
El repositorio oficial es \url{https://github.com/AJcano09/explora-nicaragua}. Contiene la solución, los proyectos de dominio y aplicación, los servicios de infraestructura, la API y ambos clientes Windows Forms \cite{ref19}. Git conserva la evolución del código y permite recuperar el estado de cada componente.

La raíz reúne las instrucciones de instalación, la configuración de servicios y los scripts de preparación. El código de la aplicación se concentra en \archivo{src}, mientras \archivo{odd/tasks} organiza las tareas técnicas. Los archivos de configuración definen las opciones de cada entorno y las credenciales se suministran mediante variables de entorno o secretos de desarrollo.

\subsection{Organización del repositorio e historial de cambios}
La organización por proyectos permite localizar interfaces, entidades, servicios y clientes sin mezclar sus responsabilidades. El historial utiliza mensajes que identifican el componente y el tipo de cambio: funcionalidades, correcciones, refactorizaciones y documentación. La siguiente tabla resume los hitos principales del desarrollo.

\begin{longtable}{@{}P{2.25cm}P{2cm}P{8.3cm}@{}}
\caption{Historial de cambios de ExploraNica.}\\
\toprule Fecha & Commit & Cambio registrado \\ \midrule\endfirsthead
\toprule Fecha & Commit & Cambio registrado \\ \midrule\endhead
10/10/2026 & \texttt{fe8c449} & Indicadores de tours activos y reservas confirmadas en administración. \\
11/10/2026 & \texttt{0b21e52} & Endpoint y acción de escritorio para validar vouchers. \\
11/10/2026 & \texttt{9f4d173} & Control de concurrencia en validación de vouchers. \\
11/10/2026 & \texttt{96824c0} & Transacción de reserva dentro de la estrategia de ejecución. \\
12/10/2026 & \texttt{a63af9c} & Registro público de operadores con documentación. \\
12/10/2026 & \texttt{54ea7af} & Actualización de operadores y endpoints documentales. \\
12/10/2026 & \texttt{88236ac} & Vistas de detalle, edición y documentos en administración. \\
12/10/2026 & \texttt{ea75b35} & Carga y visualización documental para operadores. \\
12/10/2026 & \texttt{75eccef} & Centralización del tratamiento de errores asíncronos de interfaz. \\
12/10/2026 & \texttt{33ce0d8} & Traducción de mensajes visibles pendientes al español. \\
\bottomrule
\end{longtable}

El desarrollo se realizó de forma incremental. Primero se incorporaron la gestión de tours y los indicadores administrativos; después se añadieron la validación de vouchers, los controles de concurrencia, el registro documental y las pantallas de verificación. La centralización de errores asíncronos y la traducción de mensajes completaron la integración de los clientes de escritorio.

\section{Documentación técnica básica}
\subsection{Guía de instalación y ejecución}
\textbf{Requisitos.} Los clientes requieren Windows. El entorno de desarrollo utiliza SDK .NET 10, Git y PostgreSQL. Docker Compose permite iniciar PostgreSQL 16 y los servicios auxiliares definidos en el proyecto. La API y los dos clientes se ejecutan como procesos independientes.

\textbf{Preparación.} Para una copia nueva:
\begin{Verbatim}[fontsize=\small,breaklines=true]
git clone https://github.com/AJcano09/explora-nicaragua.git
cd explora-nicaragua
dotnet restore ExploraNica.slnx
dotnet build ExploraNica.slnx -c Release
\end{Verbatim}
Para trabajar con una copia existente, abrir su carpeta raíz. Con Docker Compose, copiar \archivo{.env.example} a \archivo{.env}, definir las variables de los servicios y ejecutar \archivo{docker compose up -d explora-db}. Con PostgreSQL instalado directamente, crear la base de datos y configurar la cadena de conexión. La API recibe su configuración mediante \archivo{appsettings.json} y variables de entorno; el archivo \archivo{.env} pertenece a la configuración de Docker.

\textbf{Configuración de API.} Establecer \archivo{ConnectionStrings:DefaultConnection}, \archivo{Jwt:Key}, \archivo{Jwt:Issuer}, \archivo{Jwt:Audience} y \archivo{Jwt:ExpiryMinutes}. Usar secretos de desarrollo o variables de entorno con doble guion bajo, por ejemplo \archivo{ConnectionStrings__DefaultConnection}. Mantener los secretos fuera del control de versiones. La carga inicial puede crear usuarios de demostración; el hash de sus contraseñas sólo se habilita en desarrollo cuando se suministra \archivo{Seed:DevPassword}.

\textbf{Base de datos.} Instalar una versión 10.x de \archivo{dotnet-ef} compatible con los paquetes Entity Framework Core 10 del proyecto. Ejecutar desde la raíz:
\begin{Verbatim}[fontsize=\small,breaklines=true]
dotnet ef database update --project src/Infrastructure/ExploraNica.Infrastructure --startup-project src/Presentation/ExploraNica.Api
\end{Verbatim}
La API también intenta migrar e inicializar datos al arrancar. Si la base de datos falla, puede mantener el proceso vivo aunque no esté listo para atender operaciones. Por eso se deben comprobar los dos endpoints de salud.

\textbf{Ejecución de API y clientes.} En terminales independientes:
\begin{Verbatim}[fontsize=\small,breaklines=true]
dotnet run --project src/Presentation/ExploraNica.Api --launch-profile https
dotnet run --project src/Presentation/ExploraNica.OperatorApp
dotnet run --project src/Presentation/ExploraNica.AdminApp
\end{Verbatim}
El perfil HTTPS declara \url{https://localhost:7152} y el perfil HTTP \url{http://localhost:5241}. Configurar \archivo{Api:BaseUrl} de ambos clientes con la dirección elegida. En desarrollo puede ser necesario preparar el certificado local con \archivo{dotnet dev-certs https --trust}. Revisar \archivo{/health/live}, \archivo{/health/ready} y, en desarrollo, \archivo{/swagger}. El endpoint de disponibilidad debe responder satisfactoriamente antes de probar operaciones que dependan de PostgreSQL.

\textbf{Verificación de instalación.} Comprobar la compilación de la solución, la aplicación de las migraciones y la disponibilidad de la base de datos. Iniciar la API y ambos clientes, verificar la dirección del servidor e iniciar sesión con el rol correspondiente. Los endpoints de salud permiten distinguir entre proceso activo y servicio listo para atender solicitudes.

\subsection{Manual de usuario}
\textbf{Acceso.} Abrir el cliente correspondiente, escribir correo y contraseña y pulsar iniciar sesión. Si aparece un error de credenciales, corregir los datos; si el problema es de conexión, comprobar la API. No compartir cuentas. Al concluir las tareas, utilizar la opción de cerrar sesión.

\textbf{Operador: panel y tours.} Consultar el estado de la cuenta y el resumen disponible. En ``Mis Tours'', revisar el listado y abrir la creación o edición. Completar título, descripción, ubicación, precio y duración; guardar y comprobar que el listado refleja la modificación. La creación de ofertas exige aprobación de la agencia en el servicio correspondiente.

\textbf{Operador: documentos.} Abrir la gestión documental desde el acceso disponible en el panel. Consultar los tres tipos requeridos y adjuntar el archivo correcto para cada uno. En caso de rechazo, revisar el motivo, sustituir los documentos permitidos y reenviar la solicitud. Una cuenta aprobada no admite sustitución documental en el flujo actual.

\textbf{Operador: sesiones y reservas.} Seleccionar el tour y registrar la salida mediante el diálogo de sesiones. En ``Reservas'', seleccionar primero el tour y después su sesión; comprobar fecha, cupos y reserva antes de pulsar la validación. Validar sólo al verificar la atención del servicio, porque esta acción completa la reserva. El registro de reservas se realiza mediante la API.

\textbf{Administrador.} Consultar el panel de control, abrir operadores, filtrar por estado y revisar datos y documentos. Aprobar tras la revisión o rechazar indicando un motivo comprensible. En tours, activar o desactivar la oferta seleccionada y confirmar el cambio. En reservas, consultar los registros globales y su estado.

\textbf{Incidencias.} Anotar operación, hora e identificador de correlación cuando esté disponible. Antes de reintentar, consultar si la operación fue guardada. No adjuntar documentos personales, claves o tokens a los reportes.

\chapter{Resultados del proyecto}
\label{cap:resultados}
\section{Resultados obtenidos}
Se desarrolló una solución de gestión turística organizada en siete proyectos: dominio, aplicación, infraestructura, componentes compartidos, API, cliente del operador y cliente administrativo. Esta organización separa las entidades, los contratos de servicio, la persistencia y las interfaces de usuario. Los dos clientes de escritorio comparten el mismo servidor y la misma base de datos.

Se implementaron autenticación JWT, verificación documental de operadores, gestión de tours y sesiones, creación de reservas, validación de vouchers e indicadores operativos. La persistencia utiliza Entity Framework Core y PostgreSQL. Las migraciones mantienen el esquema de datos, el middleware centraliza las excepciones y los endpoints de salud informan sobre el estado del servidor y la conexión a la base de datos.

El registro de agencias crea expedientes con sus tres tipos documentales y estado de verificación. La gestión del catálogo vincula cada tour con su operador y sus salidas. El proceso de reserva calcula el importe y descuenta disponibilidad dentro de una transacción. La validación del voucher completa la reserva, registra su utilización y genera el evento de auditoría correspondiente.

Se incorporaron dos mecanismos de integridad para las operaciones concurrentes: actualización condicionada del cupo y cambio condicionado de la reserva confirmada a completada. La primera operación evita descontar más plazas de las disponibles; la segunda impide que dos validaciones completen repetidamente la misma reserva. El plan de validación funcional incluye estos recorridos y sus casos de excepción.

\section{Beneficios del sistema}
El beneficio funcional esperado es concentrar información antes dispersa entre mensajes, registros manuales y archivos independientes. La asociación de operadores, tours, sesiones y reservas permite consultar la operación con una estructura uniforme. La revisión documental facilita un proceso administrativo explícito y los estados de reserva permiten identificar qué servicios siguen confirmados y cuáles fueron atendidos.

Desde el punto de vista técnico, los contratos de aplicación y la API común ofrecen una base para reutilizar servicios desde ambos clientes. La centralización de reglas facilita corregirlas en el servidor. Los mensajes de dominio y la correlación de errores favorecen el diagnóstico; la bitácora existente permite seguir determinadas acciones críticas, aunque su cobertura debe ampliarse.

Para las agencias, la herramienta organiza la oferta, las salidas y las reservas en una interfaz de escritorio. Para la administración, los expedientes y los filtros por estado facilitan la distribución de tareas de verificación. La consulta centralizada permite localizar registros sin recorrer conversaciones o archivos independientes.

\section{Limitaciones del prototipo}
Los clientes de escritorio funcionan en Windows y necesitan conectividad con la API para las operaciones remotas. Los documentos se almacenan en el servidor mediante un servicio de archivos local. La continuidad operativa requiere respaldar tanto PostgreSQL como los archivos documentales y conservar su correspondencia.

Los pagos en línea, la conciliación, los reembolsos, las reseñas verificadas, los reportes financieros y la aplicación móvil pertenecen a la ampliación funcional del sistema. La reserva actual se confirma desde la API y no ejecuta un cargo de pasarela. El voucher registra el código y su utilización durante el cierre de la reserva; la emisión anticipada y lectura de una imagen QR requieren un flujo específico.

La autenticación utiliza tokens con vencimiento y el servidor no expone un servicio de renovación. El modelo reúne al personal de la agencia bajo el rol TourOperator. La analítica se concentra en conteos de control; el análisis financiero necesita sus propias entidades y reglas. La cancelación de reservas y la cobertura de permisos de todas las rutas requieren fortalecer sus controles de estado y propiedad.

La evolución del prototipo comprende pruebas automatizadas, evaluación de concurrencia, accesibilidad y recuperación ante fallos. La operación comercial también requiere procedimientos de soporte, despliegue, monitoreo y atención de incidencias.

\section{Propuestas de mejora}
Se propone priorizar cuatro etapas. La primera debe cerrar la coherencia del núcleo: permisos por propiedad en todas las rutas, validación uniforme, cancelación idempotente, renovación segura de sesión y pruebas de disponibilidad y validación concurrentes. La segunda debe completar evidencias: pruebas automatizadas, capturas por módulo y un piloto con datos de demostración y operadores voluntarios.

La tercera etapa puede incorporar pagos en entorno de prueba, estados de pago separados de los de reserva, procesamiento idempotente de eventos, conciliación y manejo de cancelaciones. Los eventos de la pasarela deben validarse y no asumirse únicos ni necesariamente ordenados \cite{ref18}. La selección de proveedor deberá confirmar su disponibilidad comercial y operativa para el proyecto antes de efectuar cobros reales.

La cuarta etapa puede desarrollar reseñas vinculadas a servicios atendidos, reportes de bruto, comisión y neto, interfaz móvil y mejoras de accesibilidad. Cada ampliación necesita criterios de aceptación, responsable y evidencia. El orden prioriza integridad de datos y demostración del núcleo antes de ampliar el catálogo de funciones.

\chapter{Impacto, ética e identidad nacional}
\label{cap:impacto}
\section{Impacto tecnológico}
ExploraNica aplica una arquitectura de servicios centralizados a la gestión operativa del turismo. Los clientes de escritorio comparten una API y una base de datos con contratos, estados y relaciones definidos. Esta organización concentra las operaciones sobre una fuente común y permite ampliar los servicios sin duplicar la persistencia en cada cliente.

El proyecto integra programación orientada a objetos, persistencia relacional, autenticación, consumo de servicios, programación asíncrona y control de versiones. La combinación de estas tecnologías permite desarrollar interfaces de gestión y coordinar operaciones sobre datos compartidos. La separación de componentes facilita mantener el producto y ampliar sus servicios.

La sostenibilidad tecnológica se apoya en dependencias mantenidas, decisiones de arquitectura documentadas y separación entre clientes y servidor. Los contratos permiten evolucionar los componentes de presentación y mantener las reglas de negocio en los servicios de la API.

\section{Impacto económico}
ExploraNica organiza la oferta de las agencias y facilita su incorporación a procesos digitales. La Cuenta Satélite de Turismo del Banco Central de Nicaragua describe la producción, el consumo y el aporte económico del sector y sirve como referencia para analizar su contexto \cite{ref2}. El proyecto se orienta a las necesidades de gestión de los prestadores de servicios turísticos.

El modelo de negocio plantea una comisión del 8\,\% por reserva. Para una base de 100 unidades monetarias, esta regla corresponde a una comisión de 8 y un saldo preliminar de 92, antes de costos de pasarela, impuestos y devoluciones. Su aplicación requiere incorporar el módulo financiero, registrar las transacciones y definir las reglas de conciliación.

La evaluación económica futura debe registrar costo de infraestructura, mantenimiento, capacitación y soporte; reservas atendidas; ingresos efectivamente conciliados y tasa de cancelación. Sólo con estos datos podrá valorarse si la herramienta aporta ahorro operativo o rentabilidad suficiente para sostener su servicio.

\section{Impacto social}
La digitalización facilita que las agencias de menor tamaño organicen su oferta con una estructura común de destinos, precios, fechas y responsables. La administración de expedientes y estados establece un procedimiento de revisión compartido. La claridad de la información y la identificación del operador contribuyen a una relación más ordenada entre prestadores y usuarios.

El acceso también puede verse limitado por conectividad, disponibilidad de equipos Windows y experiencia digital. El piloto debe incluir capacitación breve, tareas comprensibles y mecanismos de soporte. La inclusión requiere probar legibilidad, navegación por teclado, contraste y claridad de los mensajes, además de considerar agencias fuera de los centros urbanos.

El sistema complementa la atención humana mediante la organización del servicio y el registro de las decisiones. La atención de incidentes de viaje, disputas y necesidades de accesibilidad permanece a cargo de las personas responsables. La capacitación y el soporte facilitan su incorporación a la operación de las agencias.

\section{Ética profesional en el desarrollo tecnológico}
La ética profesional orienta el diseño de reglas transparentes, la protección de datos y el tratamiento responsable de la información turística. El sistema diferencia las responsabilidades del administrador, el operador y el turista. Las operaciones críticas utilizan controles de identidad y estado para proteger la información y mantener la integridad del servicio.

La protección de datos comprende acceso restringido a los expedientes, uso de información necesaria para prestar el servicio y gestión de su conservación. Las interfaces de consulta limitan los datos documentales expuestos. Las credenciales, los tokens y las claves de firma se administran como secretos de configuración y se excluyen del material de capacitación.

La autorización debe aplicarse en el servidor y comprobar la relación entre usuario y recurso. El equipo debe revisar los accesos por rol y documentar incidentes sin exponer datos personales. La bitácora debe identificar al responsable real, proteger su integridad y evitar almacenar secretos. Los registros actuales con responsable genérico requieren mejorarse para ofrecer trazabilidad individual.

La verificación de agencias comunica su estado y el motivo de rechazo para permitir la corrección del expediente. La incorporación de reseñas y pagos requiere reglas de moderación, cargos, cancelación y reclamación comprensibles para el usuario. La operación del servicio debe integrar los requisitos normativos aplicables al tratamiento de datos y a la actividad turística.

\section{Historia e identidad nacional en el proyecto}
La identidad nacional se expresa en la selección de Nicaragua como territorio de referencia y en la vinculación de sus destinos con operadores locales. Ometepe, Somoto y San Juan del Sur representan escenarios diversos de paisaje y actividad turística. El catálogo permite organizar experiencias asociadas con destinos y prestadores del país.

El catálogo debe presentar la historia y cultura de cada destino con fuentes identificadas y revisión local. Las descripciones de patrimonio, acontecimientos y tradiciones deben evitar errores, estereotipos o apropiación de expresiones comunitarias. Cuando se incorporen fotografías, relatos o conocimientos de una comunidad, debe verificarse su autorización de uso y reconocerse su procedencia.

La utilización del español en los mensajes y de una identidad visual vinculada al país contribuye a la familiaridad del producto. No obstante, los símbolos y nombres nacionales por sí solos no garantizan pertinencia cultural. Esta debe evaluarse mediante participación de operadores y comunidades, respeto por sus formas de representación y actualización responsable de la información.

La proyección turística debe equilibrar promoción y cuidado del entorno: describir recomendaciones de visita, capacidad de atención y condiciones del servicio, sin prometer acceso irrestricto a espacios sensibles. El sistema puede incorporar campos de procedencia y fecha de revisión del contenido para mantener responsabilidad editorial.

\section{Proyección}
La proyección inmediata es consolidar un prototipo demostrable para agencias y administración, con instalación reproducible, pruebas del núcleo y evidencias de uso. La siguiente etapa es un piloto acotado en el que se midan tareas concretas: crear una oferta, programar una salida, consultar reservas y validar atención. Cada participante debe conocer el propósito de la evaluación y el tratamiento de sus datos.

Si el piloto acredita utilidad, puede avanzarse hacia pagos de prueba, reportes y acceso móvil, sin comprometer de antemano una operación comercial a escala nacional. La expansión requiere responsables de soporte, presupuesto de mantenimiento, respaldo y recuperación, acuerdos con operadores y políticas de atención al usuario.

La proyección de largo plazo es consolidar una plataforma de gestión turística con participación local y mejora continua. La expansión del producto se apoya en estabilidad técnica, uso sostenido y sostenibilidad económica. Los nuevos módulos deben responder a necesidades concretas de operadores y turistas, con criterios de aceptación definidos.

\chapter{Conclusiones, recomendaciones y anexos}
\label{cap:cierre}
\section{Conclusiones}
El proyecto traduce la necesidad de organizar la oferta turística y la gestión de agencias en un prototipo con clientes de escritorio y servicios compartidos. La relación entre operadores, tours, sesiones y reservas ofrece una base concreta para digitalizar tareas internas descritas en el diagnóstico y los procesos propuestos.

Se implementaron autenticación, verificación documental, catálogo, disponibilidad, reservas y validación de vouchers. La solución utiliza .NET 10, Windows Forms, ASP.NET Core, Entity Framework Core y PostgreSQL. La organización por proyectos y las migraciones permiten mantener los componentes y la evolución del esquema de datos.

La gestión interna constituye el núcleo del sistema. Las integraciones de pago, reseñas, reportes financieros y aplicación móvil conforman su ruta de ampliación. La arquitectura de servicios permite incorporar estos componentes sobre las entidades y procesos existentes, manteniendo las reglas de identidad, propiedad y estado.

Las transacciones de reserva y los cambios condicionados de disponibilidad y estado aportan mecanismos concretos de integridad. La validación del voucher relaciona al operador con la reserva y registra el cierre del servicio. La correlación de errores y los eventos de auditoría facilitan el seguimiento de incidencias y operaciones críticas.

Finalmente, el valor del proyecto depende de combinar desarrollo técnico con responsabilidad sobre los datos, transparencia de las decisiones e información cultural verificable. La pertinencia nacional se fortalecerá mediante participación de operadores y comunidades, documentación completa y evaluación de las necesidades reales del sector.

\section{Recomendaciones}
\begin{enumerate}
\item Consolidar una versión de entrega identificada por commit y etiqueta. Separar configuración local de código publicado y documentar el procedimiento de actualización.
\item Revisar permisos y pertenencia de recursos en todos los endpoints; homogeneizar validaciones y cerrar la renovación de sesión y la cancelación repetida de reservas.
\item Incorporar pruebas automatizadas de reserva, disponibilidad y voucher. Incluir fallos de base de datos, reintentos y solicitudes concurrentes, y conservar los resultados.
\item Mantener el manual de usuario y el registro de pruebas alineados con las funcionalidades del sistema. Utilizar datos de demostración anonimizados en la capacitación.
\item Realizar un piloto con tareas y métricas definidas antes de valorar ahorro, satisfacción o aceptación. Recoger dificultades de conectividad y aprendizaje, además de resultados positivos.
\item Incorporar pagos primero en entorno de prueba, con estados separados, validación de eventos, idempotencia y conciliación. Evaluar costos y disponibilidad del proveedor antes de operar.
\item Desarrollar reportes financieros y reseñas después de estabilizar el núcleo, definiendo fuentes, reglas de cálculo, moderación y derechos de reclamación.
\item Fortalecer respaldos de PostgreSQL y archivos documentales, recuperación y trazabilidad del usuario real. Mantener secretos fuera del repositorio y limitar la retención de documentos.
\item Revisar accesibilidad, mensajes y representación cultural con usuarios locales. Registrar procedencia y derechos de uso de fotografías y textos de los destinos.
\item Establecer un ciclo periódico de mantenimiento, actualización de dependencias y revisión de la información de los destinos.
\end{enumerate}

\section{Referencias bibliográficas}
La fundamentación del sistema reúne fuentes de arquitectura de software, desarrollo de escritorio, autenticación, modelado de procesos y gestión turística.
\begingroup
\renewcommand{\chapter}[2]{}%
\renewcommand{\bibname}{}%
\begin{thebibliography}{20}
\bibitem{ref1} Instituto Nicaragüense de Turismo. \emph{Oferta turística: agencias de viajes}. Mapa Nacional de Turismo. \url{https://www.mapanicaragua.com/oferta-turistica-agencias-de-viajes/}.
\bibitem{ref2} Banco Central de Nicaragua. \emph{Informe de resultados de la Cuenta Satélite de Turismo de Nicaragua en 2023}. \url{https://www.bcn.gob.ni/publicaciones/informe-de-resultados-de-la-cuenta-satelite-de-turismo-de-nicaragua-en-2023}.
\bibitem{ref3} Ries, E. (2011). \emph{The Lean Startup}. Crown Business.
\bibitem{ref4} Osterwalder, A. y Pigneur, Y. (2010). \emph{Business Model Generation}. Wiley.
\bibitem{ref5} React Native. \emph{Documentación oficial}. \url{https://reactnative.dev/}. 
\bibitem{ref6} Civitatis. \emph{Proveedores: preguntas frecuentes}. \url{https://www.civitatis.com/es/proveedores/landing-preguntas-frecuentes/}.
\bibitem{ref7} Hardt, D. (ed.) (2012). \emph{The OAuth 2.0 Authorization Framework}. RFC 6749. \url{https://www.rfc-editor.org/info/rfc6749/}.
\bibitem{ref8} Tripadvisor (2023). \emph{Review Transparency Report}. \url{https://www.tripadvisor.com/TransparencyReport2023}.
\bibitem{ref9} Mozilla Developer Network. \emph{Offline and background operation}. \url{https://developer.mozilla.org/en-US/docs/Web/Progressive_web_apps/Guides/Offline_and_background_operation}.
\bibitem{ref10} Firebase. \emph{Firebase Cloud Messaging}. \url{https://firebase.google.com/docs/cloud-messaging}.
\bibitem{ref11} Object Management Group (2011). \emph{Business Process Model and Notation, Version 2.0}. \url{https://www.omg.org/spec/BPMN/2.0/}.
\bibitem{ref12} Camunda. \emph{Desktop Modeler}. \url{https://docs.camunda.io/docs/components/modeler/desktop-modeler/}.
\bibitem{ref13} Maurya, A. (2012). \emph{Running Lean}. O'Reilly Media.
\bibitem{ref14} Fielding, R. T. (2000). \emph{Architectural Styles and the Design of Network-based Software Architectures}. Tesis doctoral, University of California, Irvine. \url{https://ics.uci.edu/~fielding/pubs/dissertation/top.htm}.
\bibitem{ref15} Microsoft. \emph{Windows Forms overview}. \url{https://learn.microsoft.com/en-us/dotnet/desktop/winforms/overview/}.
\bibitem{ref16} Microsoft. \emph{Documentación de .NET}. \url{https://learn.microsoft.com/dotnet/}.
\bibitem{ref17} NIST. \emph{Role Based Access Control}. \url{https://csrc.nist.gov/projects/role-based-access-control/faqs}.
\bibitem{ref18} Stripe. \emph{Receive Stripe events in your webhook endpoint}. \url{https://docs.stripe.com/webhooks}.
\bibitem{ref19} ExploraNica. \emph{Repositorio oficial del sistema}. \url{https://github.com/AJcano09/explora-nicaragua}.
\bibitem{ref20} Jones, M., Bradley, J. y Sakimura, N. (2015). \emph{JSON Web Token (JWT)}. RFC 7519. \url{https://www.rfc-editor.org/info/rfc7519/}.
\end{thebibliography}
\endgroup

\section{Anexos}
\subsection*{Anexo A. Matriz de verificación funcional}
El plan de validación funcional define las operaciones principales y sus condiciones de aceptación. Cada caso establece una acción, la respuesta del sistema y la comprobación de integridad correspondiente.
\begin{longtable}{@{}P{1cm}P{4.5cm}P{6.8cm}@{}}
\toprule Caso & Acción & Condición de aceptación \\ \midrule\endfirsthead
\toprule Caso & Acción & Condición de aceptación \\ \midrule\endhead
A01 & Iniciar sesión con datos válidos e inválidos. & Acceso por rol o rechazo; captura sin credenciales y respuesta HTTP. \\
A02 & Registrar operador con tres documentos válidos. & Cuenta pendiente y expediente completo; identificador anonimizado. \\
A03 & Registrar correo o RUC repetido. & Conflicto controlado; sin duplicados ni archivos huérfanos. \\
A04 & Adjuntar archivo inválido o mayor a 5 MiB. & Rechazo por campo; sin persistencia parcial. \\
A05 & Aprobar, rechazar y reenviar una solicitud. & Estados y motivo coherentes; capturas antes y después. \\
A06 & Crear un tour con cuenta aprobada y pendiente. & Creación permitida o denegada según ruta y permiso. \\
A07 & Consultar y editar recursos de otra agencia. & Denegación sin exposición de datos. \\
A08 & Reservar dos cupos de una sesión con diez. & Reserva confirmada y disponibilidad igual a ocho. \\
A09 & Disputar concurrentemente el último cupo. & Una confirmación; otra solicitud rechazada; cupo no negativo. \\
A10 & Validar dos veces la misma reserva. & Una transición a completada y rechazo de reutilización. \\
A11 & Validar una reserva ajena. & Sin cambio de estado ni acceso a datos de otro operador. \\
A12 & Cancelar repetidamente una reserva. & Probar idempotencia y cupos; registrar cualquier desviación. \\
A13 & Interrumpir base de datos o conectividad. & Error comprensible, correlación y recuperación documentada. \\
A14 & Comprobar vencimiento de sesión. & Solicitud de autenticación al expirar el token y acceso restablecido tras iniciar sesión. \\
\bottomrule
\end{longtable}
El registro de ejecución incluye fecha, entorno, datos iniciales, pasos, resultado, responsable y evidencia. Las incidencias se relacionan con el caso que las detecta y con el cambio que las resuelve.

\subsection*{Anexo B. Registro de evidencia visual}
La documentación visual se organiza por tarea: acceso al sistema, gestión de operadores, expediente documental, catálogo de tours, sesiones, reservas y panel administrativo. Los identificadores internos de estas pantallas son FormLogin, UcVendors, OperatorDocumentsForm, OperatorToursForm, SessionDialog, OperatorBookingsForm y UcDashboard. Las imágenes de soporte se conservan en \archivo{assets/evidencias} y se vinculan con el manual de usuario.

\subsection*{Anexo C. Rutas de referencia de la API}
\begin{longtable}{@{}P{2.35cm}P{7.1cm}P{3.8cm}@{}}
\toprule Método & Ruta & Propósito \\ \midrule\endfirsthead
\toprule Método & Ruta & Propósito \\ \midrule\endhead
POST & \archivo{/api/auth/login} & Inicio de sesión. \\
POST & \archivo{/api/operators} & Registro con documentos. \\
GET & \archivo{/api/operators/me/documents} & Expediente propio. \\
POST & \archivo{/api/operators/me/submit-verification} & Reenvío a revisión. \\
GET & \archivo{/api/admin/operators} & Listado administrativo. \\
POST & \archivo{/api/admin/operators/{id}/approve} & Aprobar operador. \\
POST & \archivo{/api/admin/operators/{id}/reject} & Rechazar con motivo. \\
GET/POST & \archivo{/api/operator/tours} & Consultar o crear tours propios. \\
POST & \archivo{/api/bookings} & Crear reserva autenticada. \\
POST & \archivo{/api/operator/tours/bookings/{id}/validate} & Validar una reserva del operador. \\
GET & \archivo{/health/live} y \archivo{/health/ready} & Salud de proceso y disponibilidad. \\
\bottomrule
\end{longtable}
OpenAPI publica los contratos de la API y Swagger permite consultar las rutas durante el desarrollo. Las solicitudes de los recursos protegidos incluyen el token de acceso, los identificadores de recurso y los campos definidos por cada DTO.

\subsection*{Anexo D. Glosario técnico básico}
\begingroup\small\setlength{\parskip}{0pt}
\begin{description}
\setlength{\itemsep}{0.1em}
\item[API REST.] Interfaz de servicios que expone recursos mediante operaciones HTTP.
\item[DTO.] Objeto para transferir datos entre componentes sin exponer toda la entidad.
\item[JWT.] Token firmado con atributos de identidad y vigencia verificables.
\item[Idempotencia.] Propiedad por la que repetir una operación no duplica su efecto de negocio.
\item[Migración.] Cambio versionado del esquema de base de datos.
\item[Transacción.] Conjunto de escrituras que se confirma o revierte como una unidad.
\item[Voucher.] Comprobante de reserva con código y registro de utilización.
\item[Verificación documental.] Evaluación administrativa del expediente de una agencia.
\end{description}
\endgroup

\end{document}
